% !TeX program = pdflatex
% papers/cristal_manual.tex — GERADO por tools/cristal_tex.py; NÃO editar à mão.
% Fonte: cristal/cristal.jsonl (broca-so, última versão por conceito).
% Cada secção carrega o registo original em comentário %CRISTAL — a volta é
% tests/cristal_volta.js (resíduo 0 byte a byte contra a fonte).
\documentclass[11pt,a4paper]{article}
\usepackage[utf8]{inputenc}\usepackage[T1]{fontenc}\usepackage[brazil]{babel}
\usepackage{amsmath,amssymb}\usepackage{textcomp}
\usepackage[margin=2.4cm]{geometry}\usepackage{xcolor}
\definecolor{cinza}{gray}{0.35}
\newcommand{\prov}[1]{\par\smallskip\noindent{\small\color{cinza}#1}\par}
\setcounter{secnumdepth}{0}
% ─────────────────────────────────────────────────────────────────────────────
%  papers/gkcapa.tex — a capa do Reino, mínima, para os papers em `article`.
%
%  O estilo.tex completo é do `report` (títulos de capítulo, skak, …). Aqui fica
%  só o que a capa precisa, para o paper continuar a compilar sozinho com
%  pdflatex. O tradutor próprio (tests/tex) carrega o estilo.tex da raiz / o
%  symlink papers/estilo.tex e avalia o mesmo `\gkcapa`.
% ─────────────────────────────────────────────────────────────────────────────
\definecolor{ouro}{HTML}{B8912F}
\definecolor{ouroclaro}{HTML}{F3E9CE}
\definecolor{tinta}{HTML}{1E1B16}
\definecolor{regua}{HTML}{6E6858}
\providecommand{\gknota}{\fontsize{7.62}{11.05}\selectfont}
\providecommand{\gktexto}{\fontsize{10.50}{15.22}\selectfont}
\providecommand{\gksub}{\fontsize{12.33}{17.87}\selectfont}
\providecommand{\gksec}{\fontsize{14.47}{20.98}\selectfont}
\providecommand{\gkcap}{\fontsize{16.99}{24.63}\selectfont}
\providecommand{\gktit}{\fontsize{23.42}{33.95}\selectfont}
\providecommand{\Prj}{\mathbb{P}}
\providecommand{\Trf}{\mathcal{T}}
\providecommand{\gkcapa}[3]{%
\title{\vspace{-0.6cm}
{\color{ouro}\rule{\textwidth}{1.4pt}}\\[6mm]
\textbf{\gktit\color{tinta} Reino Dourado}\\[3mm]{\gksec\itshape\color{regua} Golden Kingdom}\\[8mm]
{\gkcap\scshape\color{ouro} #1}\\[5mm]
{\gksub\color{regua} #2}\\[2mm]
{\gktexto\color{regua}\itshape #3}\\[7mm]
{\color{ouro}\rule{\textwidth}{0.6pt}}\\[9mm]{\gknota\color{regua}\begin{minipage}{\textwidth}\setlength{\parindent}{0pt}%
\textbf{Aviso de Isenção Algorítmica:} O universo, as leis e as entidades contidas nestas páginas ---
incluindo felinos matriciais, aves de rapina algébricas e amuletos de controle de fluxo --- são
construções de pura ficção formal. Esta obra não descreve a realidade física, não serve como
engenharia reversa do cosmos e não constitui doutrina religiosa. Qualquer semelhança com bugs reais do seu
sistema operacional é mera coincidência (ou falta de reza). Prossiga por sua própria conta e
risco.\end{minipage}}}
\author{}\date{}}

\begin{document}
\gkcapa{O Cristal --- o manual do cristal}
       {operação, referência, os repos, a conversa}
       {corpus recuperado do broca-so; 516 conceitos; projeção verificável da fonte}
\maketitle

%CRISTAL {"arestas":[{"alvo":"perguntas","confianca":1.0,"epistemico":"fato","origem":"POW_METAL_PRE_FECHO.md","rel":"parte_de"},{"alvo":"6_provas_de_composicao","confianca":0.5,"epistemico":"fato","origem":"wiki","rel":"relaciona"},{"alvo":"semiotica","confianca":0.5,"epistemico":"fato","origem":"wiki","rel":"relaciona"},{"alvo":"vida-morte","confianca":0.5,"epistemico":"fato","origem":"wiki","rel":"relaciona"},{"alvo":"virtualizacao","confianca":0.5,"epistemico":"fato","origem":"wiki","rel":"relaciona"},{"alvo":"word","confianca":0.5,"epistemico":"fato","origem":"wiki","rel":"relaciona"},{"alvo":"pell","confianca":0.5,"epistemico":"fato","origem":"wiki","rel":"relaciona"},{"alvo":"educacao","confianca":0.5,"epistemico":"fato","origem":"wiki","rel":"relaciona"},{"alvo":"telemetria_shadow_producao","confianca":0.5,"epistemico":"fato","origem":"wiki","rel":"relaciona"},{"alvo":"gentil_capitulo_8_mais_aplicacoes","confianca":0.5,"epistemico":"fato","origem":"wiki","rel":"relaciona"},{"alvo":"2_vizinhos_imediatos_tem_acoplamento_semelhante","confianca":0.5,"epistemico":"fato","origem":"wiki","rel":"relaciona"}],"confianca":1.0,"contraexemplos":[],"descricao":"**2/2** fechos no top-100.","epistemico":"fato","exemplos":[],"id":"0_os_fechos_aparecem_no_topo_do_ranking_global","memoria":"semantica","meta":{"dominio":"manual","nivel":6,"verificacao":"slug idêntico — enriquecer, não duplicar"},"origem":"POW_METAL_PRE_FECHO.md","palavras_chave":[],"sinonimos":[],"tipo":"conceito","titulo":"0. Os fechos aparecem no topo do ranking global?"}
\section{0. Os fechos aparecem no topo do ranking global?}
**2/2** fechos no top-100.
\prov{relações: parte\_de perguntas; relaciona 6\_provas\_de\_composicao; relaciona semiotica; relaciona vida-morte; relaciona virtualizacao; relaciona word; relaciona pell; relaciona educacao; relaciona telemetria\_shadow\_producao; relaciona gentil\_capitulo\_8\_mais\_aplicacoes; relaciona 2\_vizinhos\_imediatos\_tem\_acoplamento\_semelhante}
\prov{proveniência: POW\_METAL\_PRE\_FECHO.md \textperiodcentered{} fato \textperiodcentered{} confiança 1.0 \textperiodcentered{} domínio manual \textperiodcentered{} história 14 versões no jornal}

%CRISTAL {"arestas":[{"alvo":"a_prova_da_isa_da_broca","confianca":1.0,"epistemico":"fato","origem":"PROVA_ISA.md","rel":"parte_de"},{"alvo":"mecanica_quantica_hub","confianca":0.5,"epistemico":"fato","origem":"wiki","rel":"relaciona"},{"alvo":"utilitarismo","confianca":0.5,"epistemico":"fato","origem":"wiki","rel":"relaciona"},{"alvo":"testes","confianca":0.5,"epistemico":"fato","origem":"wiki","rel":"relaciona"},{"alvo":"nucleo_broca_vivo","confianca":0.5,"epistemico":"fato","origem":"wiki","rel":"relaciona"},{"alvo":"3_maquina_de_estados","confianca":0.5,"epistemico":"fato","origem":"wiki","rel":"relaciona"},{"alvo":"hardware_sequencias_de_cauchy","confianca":0.5,"epistemico":"fato","origem":"wiki","rel":"relaciona"},{"alvo":"opcao_a_tarball_release","confianca":0.5,"epistemico":"fato","origem":"wiki","rel":"relaciona"},{"alvo":"regua_associativa_e_o_risco_de_vazamento","confianca":0.5,"epistemico":"fato","origem":"wiki","rel":"relaciona"},{"alvo":"4_acoplamento_cresce_monotonicamente_ate_o_fecho","confianca":0.5,"epistemico":"fato","origem":"wiki","rel":"relaciona"},{"alvo":"rust","confianca":0.5,"epistemico":"fato","origem":"wiki","rel":"relaciona"},{"alvo":"consulta_simples","confianca":0.5,"epistemico":"fato","origem":"wiki","rel":"relaciona"},{"alvo":"fase_2_tecido_de_conhecimento","confianca":0.5,"epistemico":"fato","origem":"wiki","rel":"relaciona"}],"confianca":1.0,"contraexemplos":[],"descricao":"A ISA da Broca é um **objeto matemático demonstrável**:","epistemico":"fato","exemplos":[],"id":"10_meta","memoria":"semantica","meta":{"dominio":"manual","nivel":6,"verificacao":"slug idêntico — enriquecer, não duplicar"},"origem":"PROVA_ISA.md","palavras_chave":[],"sinonimos":[],"tipo":"conceito","titulo":"10. Meta"}
\section{10. Meta}
A ISA da Broca é um **objeto matemático demonstrável**:
\prov{relações: parte\_de a\_prova\_da\_isa\_da\_broca; relaciona mecanica\_quantica\_hub; relaciona utilitarismo; relaciona testes; relaciona nucleo\_broca\_vivo; relaciona 3\_maquina\_de\_estados; relaciona hardware\_sequencias\_de\_cauchy; relaciona opcao\_a\_tarball\_release; relaciona regua\_associativa\_e\_o\_risco\_de\_vazamento; relaciona 4\_acoplamento\_cresce\_monotonicamente\_ate\_o\_fecho; relaciona rust; relaciona consulta\_simples; relaciona fase\_2\_tecido\_de\_conhecimento}
\prov{proveniência: PROVA\_ISA.md \textperiodcentered{} fato \textperiodcentered{} confiança 1.0 \textperiodcentered{} domínio manual \textperiodcentered{} história 16 versões no jornal}

%CRISTAL {"arestas":[{"alvo":"pow_coordenadas_espectrais_broca","confianca":1.0,"epistemico":"fato","origem":"POW_GEOMETRIA.md","rel":"parte_de"},{"alvo":"analise","confianca":1.0,"epistemico":"fato","origem":"manual","rel":"depende"},{"alvo":"word","confianca":0.5,"epistemico":"fato","origem":"wiki","rel":"relaciona"}],"confianca":1.0,"contraexemplos":[],"descricao":"- válidos: **65,536** - entropia calor: **1.3849** - entropia fase²: **1.3884** - autocorr total (nonce): **0.7503**","epistemico":"fato","exemplos":[],"id":"1_campo_espectral","memoria":"semantica","meta":{"dominio":"manual","nivel":6,"verificacao":"slug idêntico — enriquecer, não duplicar"},"origem":"POW_GEOMETRIA.md","palavras_chave":[],"sinonimos":[],"tipo":"conceito","titulo":"1. Campo espectral"}
\section{1. Campo espectral}
- válidos: **65,536** - entropia calor: **1.3849** - entropia fase\ensuremath{^{2}}: **1.3884** - autocorr total (nonce): **0.7503**
\prov{relações: parte\_de pow\_coordenadas\_espectrais\_broca; depende analise; relaciona word}
\prov{proveniência: POW\_GEOMETRIA.md \textperiodcentered{} fato \textperiodcentered{} confiança 1.0 \textperiodcentered{} domínio manual \textperiodcentered{} história 55 versões no jornal}

%CRISTAL {"arestas":[{"alvo":"a_prova_da_isa_da_broca","confianca":1.0,"epistemico":"fato","origem":"PROVA_ISA.md","rel":"parte_de"}],"confianca":1.0,"contraexemplos":[],"descricao":"``` LOAD  imm16   ; mem/imm → neurônio → A; empilha A anterior em B LOADS imm16   ; mem/imm → neurônio → A₂·word → A (cifra implícita) STORE imm16   ; R → slot memória (word binário) ADD           ; R = A + B  (componente a componente) SUB           ; R = A - B AND OR XOR    ; bitwise por componente SILVER GOLD BRONZE  ; A ← Aₙ·A, R ← A  (metal) FOLD UNFOLD PROJECT LIFT  ; parábola + metal composto","epistemico":"fato","exemplos":[],"id":"1_conjunto_de_instrucoes","memoria":"semantica","meta":{"dominio":"manual","duplicata_sugerida":[],"nivel":6,"verificacao":"parte de conceito mais amplo '1_conjunto_de_instrucoes'"},"origem":"README.md","palavras_chave":[],"sinonimos":[],"tipo":"conceito","titulo":"1. Conjunto de instruções"}
\section{1. Conjunto de instruções}
``` LOAD  imm16   ; mem/imm \ensuremath{\to} neurônio \ensuremath{\to} A; empilha A anterior em B LOADS imm16   ; mem/imm \ensuremath{\to} neurônio \ensuremath{\to} A\ensuremath{_{2}}\textperiodcentered word \ensuremath{\to} A (cifra implícita) STORE imm16   ; R \ensuremath{\to} slot memória (word binário) ADD           ; R = A + B  (componente a componente) SUB           ; R = A - B AND OR XOR    ; bitwise por componente SILVER GOLD BRONZE  ; A \ensuremath{\leftarrow} A\ensuremath{_{n}}\textperiodcentered A, R \ensuremath{\leftarrow} A  (metal) FOLD UNFOLD PROJECT LIFT  ; parábola + metal composto
\prov{relações: parte\_de a\_prova\_da\_isa\_da\_broca}
\prov{proveniência: README.md \textperiodcentered{} fato \textperiodcentered{} confiança 1.0 \textperiodcentered{} domínio manual \textperiodcentered{} história 6 versões no jornal}

%CRISTAL {"arestas":[{"alvo":"computacao","confianca":1.0,"epistemico":"fato","origem":"manual","rel":"parte_de"},{"alvo":"memoria_virtual","confianca":0.5,"epistemico":"fato","origem":"wiki","rel":"relaciona"},{"alvo":"8_release_core","confianca":0.5,"epistemico":"fato","origem":"wiki","rel":"relaciona"},{"alvo":"gato_operador_hub","confianca":0.5,"epistemico":"fato","origem":"wiki","rel":"relaciona"},{"alvo":"broca_so","confianca":0.5,"epistemico":"fato","origem":"wiki","rel":"relaciona"},{"alvo":"koch_memoria","confianca":0.5,"epistemico":"fato","origem":"wiki","rel":"relaciona"},{"alvo":"6_provas_de_composicao","confianca":0.5,"epistemico":"fato","origem":"wiki","rel":"relaciona"},{"alvo":"rust","confianca":0.5,"epistemico":"fato","origem":"wiki","rel":"relaciona"},{"alvo":"main","confianca":0.5,"epistemico":"fato","origem":"wiki","rel":"relaciona"},{"alvo":"cifra_prova","confianca":0.5,"epistemico":"fato","origem":"wiki","rel":"relaciona"},{"alvo":"koch_atlas","confianca":0.5,"epistemico":"fato","origem":"wiki","rel":"relaciona"},{"alvo":"python","confianca":0.5,"epistemico":"fato","origem":"wiki","rel":"relaciona"},{"alvo":"funil_cumulativo","confianca":0.5,"epistemico":"fato","origem":"wiki","rel":"relaciona"},{"alvo":"koch_pell","confianca":0.5,"epistemico":"fato","origem":"wiki","rel":"relaciona"},{"alvo":"hopfield_rede_hub","confianca":0.5,"epistemico":"fato","origem":"wiki","rel":"relaciona"},{"alvo":"grava_fato_resposta_normal_em_seguida","confianca":0.5,"epistemico":"fato","origem":"wiki","rel":"relaciona"},{"alvo":"mamifero","confianca":0.5,"epistemico":"fato","origem":"wiki","rel":"relaciona"},{"alvo":"pow_espectral","confianca":0.5,"epistemico":"fato","origem":"wiki","rel":"relaciona"},{"alvo":"install","confianca":0.5,"epistemico":"fato","origem":"wiki","rel":"relaciona"},{"alvo":"gpu","confianca":0.5,"epistemico":"fato","origem":"wiki","rel":"relaciona"},{"alvo":"driver","confianca":0.5,"epistemico":"fato","origem":"wiki","rel":"relaciona"},{"alvo":"goldcoin","confianca":0.5,"epistemico":"fato","origem":"wiki","rel":"relaciona"},{"alvo":"fechamento_blocos","confianca":0.5,"epistemico":"fato","origem":"wiki","rel":"relaciona"},{"alvo":"prova_isa","confianca":0.5,"epistemico":"fato","origem":"wiki","rel":"relaciona"},{"alvo":"sessao_interativa","confianca":0.5,"epistemico":"fato","origem":"wiki","rel":"relaciona"},{"alvo":"resposta_a_pergunta_dificil","confianca":0.5,"epistemico":"fato","origem":"wiki","rel":"relaciona"},{"alvo":"java","confianca":0.5,"epistemico":"fato","origem":"wiki","rel":"relaciona"},{"alvo":"recursao","confianca":0.5,"epistemico":"fato","origem":"wiki","rel":"relaciona"},{"alvo":"go","confianca":0.5,"epistemico":"fato","origem":"wiki","rel":"relaciona"},{"alvo":"cache","confianca":0.5,"epistemico":"fato","origem":"wiki","rel":"relaciona"},{"alvo":"historia","confianca":0.5,"epistemico":"fato","origem":"wiki","rel":"relaciona"},{"alvo":"pow_metal_pre_fecho","confianca":0.5,"epistemico":"fato","origem":"wiki","rel":"relaciona"},{"alvo":"topologia","confianca":0.5,"epistemico":"fato","origem":"wiki","rel":"relaciona"}],"confianca":1.0,"contraexemplos":[],"descricao":"``` M = (Σ, Mem, Prog, δ) ```","epistemico":"fato","exemplos":[],"id":"1_modelo_formal","memoria":"semantica","meta":{"dominio":"manual","nivel":6,"verificacao":"slug idêntico — enriquecer, não duplicar"},"origem":"EXPRESSIVIDADE.md","palavras_chave":[],"sinonimos":[],"tipo":"conceito","titulo":"1. Modelo Formal"}
\section{1. Modelo Formal}
``` M = (\ensuremath{\Sigma}, Mem, Prog, \ensuremath{\delta}) ```
\prov{relações: parte\_de computacao; relaciona memoria\_virtual; relaciona 8\_release\_core; relaciona gato\_operador\_hub; relaciona broca\_so; relaciona koch\_memoria; relaciona 6\_provas\_de\_composicao; relaciona rust; relaciona main; relaciona cifra\_prova; relaciona koch\_atlas; relaciona python; relaciona funil\_cumulativo; relaciona koch\_pell; relaciona hopfield\_rede\_hub; relaciona grava\_fato\_resposta\_normal\_em\_seguida; relaciona mamifero; relaciona pow\_espectral; relaciona install; relaciona gpu; relaciona driver; relaciona goldcoin; relaciona fechamento\_blocos; relaciona prova\_isa; relaciona sessao\_interativa; relaciona resposta\_a\_pergunta\_dificil; relaciona java; relaciona recursao; relaciona go; relaciona cache; relaciona historia; relaciona pow\_metal\_pre\_fecho; relaciona topologia}
\prov{proveniência: EXPRESSIVIDADE.md \textperiodcentered{} fato \textperiodcentered{} confiança 1.0 \textperiodcentered{} domínio manual \textperiodcentered{} história 40 versões no jornal}

%CRISTAL {"arestas":[{"alvo":"perguntas","confianca":1.0,"epistemico":"fato","origem":"POW_METAL_PRE_FECHO.md","rel":"parte_de"},{"alvo":"utilitarismo","confianca":0.5,"epistemico":"fato","origem":"wiki","rel":"relaciona"},{"alvo":"4_acoplamento_cresce_monotonicamente_ate_o_fecho","confianca":0.5,"epistemico":"fato","origem":"wiki","rel":"relaciona"},{"alvo":"nao_dualidade","confianca":0.5,"epistemico":"fato","origem":"wiki","rel":"relaciona"},{"alvo":"a_prova_da_isa_da_broca","confianca":0.5,"epistemico":"fato","origem":"wiki","rel":"relaciona"}],"confianca":1.0,"contraexemplos":[],"descricao":"Perfil monotónico: **0/2** fechos.","epistemico":"fato","exemplos":[],"id":"1_o_acoplamento_cresce_continuamente_ate_o_fecho","memoria":"semantica","meta":{"dominio":"manual","nivel":6,"verificacao":"slug idêntico — enriquecer, não duplicar"},"origem":"POW_METAL_PRE_FECHO.md","palavras_chave":[],"sinonimos":[],"tipo":"conceito","titulo":"1. O acoplamento cresce continuamente até o fecho?"}
\section{1. O acoplamento cresce continuamente até o fecho?}
Perfil monotónico: **0/2** fechos.
\prov{relações: parte\_de perguntas; relaciona utilitarismo; relaciona 4\_acoplamento\_cresce\_monotonicamente\_ate\_o\_fecho; relaciona nao\_dualidade; relaciona a\_prova\_da\_isa\_da\_broca}
\prov{proveniência: POW\_METAL\_PRE\_FECHO.md \textperiodcentered{} fato \textperiodcentered{} confiança 1.0 \textperiodcentered{} domínio manual \textperiodcentered{} história 9 versões no jornal}

%CRISTAL {"arestas":[{"alvo":"perguntas","confianca":1.0,"epistemico":"fato","origem":"POW_METAL_RESSONANCIA_VAL.md","rel":"parte_de"},{"alvo":"analise","confianca":0.5,"epistemico":"fato","origem":"wiki","rel":"relaciona"},{"alvo":"word","confianca":0.5,"epistemico":"fato","origem":"wiki","rel":"relaciona"},{"alvo":"syscall","confianca":0.5,"epistemico":"fato","origem":"wiki","rel":"relaciona"},{"alvo":"anjo_da_historia","confianca":0.5,"epistemico":"fato","origem":"wiki","rel":"relaciona"},{"alvo":"experimento_2_recorrencias_e_metais","confianca":0.5,"epistemico":"fato","origem":"wiki","rel":"relaciona"},{"alvo":"virtualizacao","confianca":0.5,"epistemico":"fato","origem":"wiki","rel":"relaciona"},{"alvo":"numeros_duais","confianca":0.5,"epistemico":"fato","origem":"wiki","rel":"relaciona"},{"alvo":"tragedia_dos_comuns","confianca":0.5,"epistemico":"fato","origem":"wiki","rel":"relaciona"},{"alvo":"lab_conversa","confianca":0.5,"epistemico":"fato","origem":"wiki","rel":"relaciona"},{"alvo":"acao_social_weber","confianca":0.5,"epistemico":"fato","origem":"wiki","rel":"relaciona"},{"alvo":"colapso_koch_tiffany_relacao_koch","confianca":0.5,"epistemico":"fato","origem":"wiki","rel":"relaciona"}],"confianca":1.0,"contraexemplos":[],"descricao":"Top-20: **2** fechos no ranking (cobertura **100%** dos fechos encontrados).","epistemico":"fato","exemplos":[],"id":"1_os_maiores_acoplamentos_concentram_os_fechos","memoria":"semantica","meta":{"dominio":"manual","nivel":6,"verificacao":"slug idêntico — enriquecer, não duplicar"},"origem":"POW_METAL_RESSONANCIA_VAL.md","palavras_chave":[],"sinonimos":[],"tipo":"conceito","titulo":"1. Os maiores acoplamentos concentram os fechos?"}
\section{1. Os maiores acoplamentos concentram os fechos?}
Top-20: **2** fechos no ranking (cobertura **100\%** dos fechos encontrados).
\prov{relações: parte\_de perguntas; relaciona analise; relaciona word; relaciona syscall; relaciona anjo\_da\_historia; relaciona experimento\_2\_recorrencias\_e\_metais; relaciona virtualizacao; relaciona numeros\_duais; relaciona tragedia\_dos\_comuns; relaciona lab\_conversa; relaciona acao\_social\_weber; relaciona colapso\_koch\_tiffany\_relacao\_koch}
\prov{proveniência: POW\_METAL\_RESSONANCIA\_VAL.md \textperiodcentered{} fato \textperiodcentered{} confiança 1.0 \textperiodcentered{} domínio manual \textperiodcentered{} história 15 versões no jornal}

%CRISTAL {"arestas":[{"alvo":"pow_coordenadas_espectrais_broca","confianca":1.0,"epistemico":"fato","origem":"POW_GEOMETRIA.md","rel":"parte_de"}],"confianca":1.0,"contraexemplos":[],"descricao":"| Ficheiro | Conteúdo | |----------|----------| | `atlas/classes.csv` | assinaturas, simetria, contagem | | `atlas/recorrencias.csv` | recorrências emergentes ordem 2 e 3 | | `atlas/pell.csv` | valores Pell/Fib/etc. + autovalores prata | | `atlas/exoticos.csv` | recorrências ordem 2 não clássicas | | `atlas/grafos.svg` | transições entre classes top | | `atlas/componentes.svg` | componentes conexas |","epistemico":"fato","exemplos":[],"id":"2_atlas_de_coordenadas","memoria":"semantica","meta":{"dominio":"manual","duplicata_sugerida":[],"nivel":6,"verificacao":"parte de conceito mais amplo '2_atlas_de_coordenadas'"},"origem":"DEGENERADOS.md","palavras_chave":[],"sinonimos":[],"tipo":"conceito","titulo":"2. Atlas de coordenadas"}
\section{2. Atlas de coordenadas}
| Ficheiro | Conteúdo | |----------|----------| | `atlas/classes.csv` | assinaturas, simetria, contagem | | `atlas/recorrencias.csv` | recorrências emergentes ordem 2 e 3 | | `atlas/pell.csv` | valores Pell/Fib/etc. + autovalores prata | | `atlas/exoticos.csv` | recorrências ordem 2 não clássicas | | `atlas/grafos.svg` | transições entre classes top | | `atlas/componentes.svg` | componentes conexas |
\prov{relações: parte\_de pow\_coordenadas\_espectrais\_broca}
\prov{proveniência: DEGENERADOS.md \textperiodcentered{} fato \textperiodcentered{} confiança 1.0 \textperiodcentered{} domínio manual \textperiodcentered{} história 10 versões no jornal}

%CRISTAL {"arestas":[{"alvo":"broca_os","confianca":0.95,"epistemico":"fato","origem":"paper","rel":"parte_de"},{"alvo":"virtualizacao","confianca":0.5,"epistemico":"fato","origem":"wiki","rel":"relaciona"},{"alvo":"transcricao","confianca":0.5,"epistemico":"fato","origem":"wiki","rel":"relaciona"},{"alvo":"utilitarismo","confianca":0.5,"epistemico":"fato","origem":"wiki","rel":"relaciona"},{"alvo":"vida-morte","confianca":0.5,"epistemico":"fato","origem":"wiki","rel":"relaciona"},{"alvo":"pedagogia","confianca":0.5,"epistemico":"fato","origem":"wiki","rel":"relaciona"},{"alvo":"scheduler","confianca":0.5,"epistemico":"fato","origem":"wiki","rel":"relaciona"},{"alvo":"paper_2_reta_ouro","confianca":0.5,"epistemico":"fato","origem":"wiki","rel":"relaciona"},{"alvo":"ruido","confianca":0.5,"epistemico":"fato","origem":"wiki","rel":"relaciona"},{"alvo":"mecanica","confianca":0.5,"epistemico":"fato","origem":"wiki","rel":"relaciona"},{"alvo":"tipos_de_interpretante","confianca":0.5,"epistemico":"fato","origem":"wiki","rel":"relaciona"},{"alvo":"prova_da_maquina_espectral","confianca":0.5,"epistemico":"fato","origem":"wiki","rel":"relaciona"},{"alvo":"experimento_3_inverter","confianca":0.5,"epistemico":"fato","origem":"wiki","rel":"relaciona"},{"alvo":"pow_metal_ressonancia_val","confianca":0.5,"epistemico":"fato","origem":"wiki","rel":"relaciona"},{"alvo":"resultado_completo","confianca":0.5,"epistemico":"fato","origem":"wiki","rel":"relaciona"},{"alvo":"pitagorismo","confianca":0.5,"epistemico":"fato","origem":"wiki","rel":"relaciona"},{"alvo":"resultado_anterior_identidade_nao_cobertura","confianca":0.5,"epistemico":"fato","origem":"wiki","rel":"relaciona"},{"alvo":"broca_os_hub","confianca":0.5,"epistemico":"fato","origem":"wiki","rel":"relaciona"},{"alvo":"modo_espectral_spect","confianca":0.5,"epistemico":"fato","origem":"wiki","rel":"relaciona"},{"alvo":"main","confianca":0.5,"epistemico":"fato","origem":"wiki","rel":"relaciona"}],"confianca":1.0,"contraexemplos":[],"descricao":"```bash export TIFFANY_HOME=~/.tiffany mkdir -p $TIFFANY_HOME","epistemico":"fato","exemplos":[],"id":"2_base_de_conhecimento_core","memoria":"semantica","meta":{"dominio":"manual","nivel":6,"verificacao":"slug idêntico — enriquecer, não duplicar"},"origem":"INSTALL.md","palavras_chave":[],"sinonimos":[],"tipo":"conceito","titulo":"2. Base de conhecimento (Core)"}
\section{2. Base de conhecimento (Core)}
```bash export TIFFANY\_HOME=\~{}/.tiffany mkdir -p \$TIFFANY\_HOME
\prov{relações: parte\_de broca\_os; relaciona virtualizacao; relaciona transcricao; relaciona utilitarismo; relaciona vida-morte; relaciona pedagogia; relaciona scheduler; relaciona paper\_2\_reta\_ouro; relaciona ruido; relaciona mecanica; relaciona tipos\_de\_interpretante; relaciona prova\_da\_maquina\_espectral; relaciona experimento\_3\_inverter; relaciona pow\_metal\_ressonancia\_val; relaciona resultado\_completo; relaciona pitagorismo; relaciona resultado\_anterior\_identidade\_nao\_cobertura; relaciona broca\_os\_hub; relaciona modo\_espectral\_spect; relaciona main}
\prov{proveniência: INSTALL.md \textperiodcentered{} fato \textperiodcentered{} confiança 1.0 \textperiodcentered{} domínio manual \textperiodcentered{} história 22 versões no jornal}

%CRISTAL {"arestas":[{"alvo":"perguntas","confianca":1.0,"epistemico":"fato","origem":"POW_METAL_PRE_FECHO.md","rel":"parte_de"},{"alvo":"virtualizacao","confianca":0.5,"epistemico":"fato","origem":"wiki","rel":"relaciona"},{"alvo":"word","confianca":0.5,"epistemico":"fato","origem":"wiki","rel":"relaciona"}],"confianca":1.0,"contraexemplos":[],"descricao":"Acoplamento médio |dist|≤8: **0.6013** · |dist|≥32: **0.6347** · região detectada: **não**.","epistemico":"fato","exemplos":[],"id":"2_existe_regiao_de_ressonancia_em_torno_do_fecho","memoria":"semantica","meta":{"dominio":"manual","nivel":6,"verificacao":"slug idêntico — enriquecer, não duplicar"},"origem":"POW_METAL_PRE_FECHO.md","palavras_chave":[],"sinonimos":[],"tipo":"conceito","titulo":"2. Existe região de ressonância em torno do fecho?"}
\section{2. Existe região de ressonância em torno do fecho?}
Acoplamento médio |dist|\ensuremath{\le}8: **0.6013** \textperiodcentered  |dist|\ensuremath{\ge}32: **0.6347** \textperiodcentered  região detectada: **não**.
\prov{relações: parte\_de perguntas; relaciona virtualizacao; relaciona word}
\prov{proveniência: POW\_METAL\_PRE\_FECHO.md \textperiodcentered{} fato \textperiodcentered{} confiança 1.0 \textperiodcentered{} domínio manual \textperiodcentered{} história 6 versões no jornal}

%CRISTAL {"arestas":[{"alvo":"broca_os","confianca":0.95,"epistemico":"fato","origem":"paper","rel":"parte_de"},{"alvo":"vetores","confianca":0.5,"epistemico":"fato","origem":"wiki","rel":"relaciona"},{"alvo":"virtualizacao","confianca":0.5,"epistemico":"fato","origem":"wiki","rel":"relaciona"},{"alvo":"word","confianca":0.5,"epistemico":"fato","origem":"wiki","rel":"relaciona"},{"alvo":"vida-morte","confianca":0.5,"epistemico":"fato","origem":"wiki","rel":"relaciona"},{"alvo":"rede_leech","confianca":0.5,"epistemico":"fato","origem":"wiki","rel":"relaciona"},{"alvo":"misticismo_nao_dual","confianca":0.5,"epistemico":"fato","origem":"wiki","rel":"relaciona"}],"confianca":1.0,"contraexemplos":[],"descricao":"```bnf prog     ::= line* line     ::= instr | label | comment label    ::= IDENT ':' comment  ::= ';' TEXT","epistemico":"fato","exemplos":[],"id":"2_gramatica_assembly","memoria":"semantica","meta":{"dominio":"manual","nivel":6,"verificacao":"slug idêntico — enriquecer, não duplicar"},"origem":"EXPRESSIVIDADE.md","palavras_chave":[],"sinonimos":[],"tipo":"conceito","titulo":"2. Gramática Assembly"}
\section{2. Gramática Assembly}
```bnf prog     ::= line* line     ::= instr | label | comment label    ::= IDENT ':' comment  ::= ';' TEXT
\prov{relações: parte\_de broca\_os; relaciona vetores; relaciona virtualizacao; relaciona word; relaciona vida-morte; relaciona rede\_leech; relaciona misticismo\_nao\_dual}
\prov{proveniência: EXPRESSIVIDADE.md \textperiodcentered{} fato \textperiodcentered{} confiança 1.0 \textperiodcentered{} domínio manual \textperiodcentered{} história 9 versões no jornal}

%CRISTAL {"arestas":[{"alvo":"a_prova_da_isa_da_broca","confianca":1.0,"epistemico":"fato","origem":"PROVA_ISA.md","rel":"parte_de"},{"alvo":"transcricao","confianca":0.5,"epistemico":"fato","origem":"wiki","rel":"relaciona"},{"alvo":"parabola_destruir","confianca":0.5,"epistemico":"fato","origem":"wiki","rel":"relaciona"},{"alvo":"resposta_a_pergunta_dificil","confianca":0.5,"epistemico":"fato","origem":"wiki","rel":"relaciona"},{"alvo":"parabola_pa","confianca":0.5,"epistemico":"fato","origem":"wiki","rel":"relaciona"},{"alvo":"vida-morte","confianca":0.5,"epistemico":"fato","origem":"wiki","rel":"relaciona"},{"alvo":"pow_metal_setor","confianca":0.5,"epistemico":"fato","origem":"wiki","rel":"relaciona"},{"alvo":"broca_cristal","confianca":0.5,"epistemico":"fato","origem":"wiki","rel":"relaciona"},{"alvo":"ds_universo","confianca":0.5,"epistemico":"fato","origem":"wiki","rel":"relaciona"},{"alvo":"teste_ordem","confianca":0.5,"epistemico":"fato","origem":"wiki","rel":"relaciona"},{"alvo":"pow_metal_pre_fecho","confianca":0.5,"epistemico":"fato","origem":"wiki","rel":"relaciona"},{"alvo":"pilha","confianca":0.5,"epistemico":"fato","origem":"wiki","rel":"relaciona"},{"alvo":"testes","confianca":0.5,"epistemico":"fato","origem":"wiki","rel":"relaciona"},{"alvo":"motivos","confianca":0.5,"epistemico":"fato","origem":"wiki","rel":"relaciona"},{"alvo":"geometria_hub_broca","confianca":0.5,"epistemico":"fato","origem":"wiki","rel":"relaciona"},{"alvo":"pow_metal_funil","confianca":0.5,"epistemico":"fato","origem":"wiki","rel":"relaciona"},{"alvo":"log_legado_conversadadostelemetriajsonl","confianca":0.5,"epistemico":"fato","origem":"wiki","rel":"relaciona"},{"alvo":"topologia","confianca":0.5,"epistemico":"fato","origem":"wiki","rel":"relaciona"},{"alvo":"hardware_sequencias_de_cauchy","confianca":0.5,"epistemico":"fato","origem":"wiki","rel":"relaciona"},{"alvo":"koch_cantor_fib_prova","confianca":0.5,"epistemico":"fato","origem":"wiki","rel":"relaciona"}],"confianca":1.0,"contraexemplos":[],"descricao":"Secção de PROVA_ISA.md.","epistemico":"fato","exemplos":[],"id":"2_semantica_formal","memoria":"semantica","meta":{"dominio":"manual","nivel":6,"verificacao":"slug idêntico — enriquecer, não duplicar"},"origem":"PROVA_ISA.md","palavras_chave":[],"sinonimos":[],"tipo":"conceito","titulo":"2. Semântica formal"}
\section{2. Semântica formal}
Secção de PROVA\_ISA.md.
\prov{relações: parte\_de a\_prova\_da\_isa\_da\_broca; relaciona transcricao; relaciona parabola\_destruir; relaciona resposta\_a\_pergunta\_dificil; relaciona parabola\_pa; relaciona vida-morte; relaciona pow\_metal\_setor; relaciona broca\_cristal; relaciona ds\_universo; relaciona teste\_ordem; relaciona pow\_metal\_pre\_fecho; relaciona pilha; relaciona testes; relaciona motivos; relaciona geometria\_hub\_broca; relaciona pow\_metal\_funil; relaciona log\_legado\_conversadadostelemetriajsonl; relaciona topologia; relaciona hardware\_sequencias\_de\_cauchy; relaciona koch\_cantor\_fib\_prova}
\prov{proveniência: PROVA\_ISA.md \textperiodcentered{} fato \textperiodcentered{} confiança 1.0 \textperiodcentered{} domínio manual \textperiodcentered{} história 23 versões no jornal}

%CRISTAL {"arestas":[{"alvo":"perguntas","confianca":1.0,"epistemico":"fato","origem":"POW_METAL_RESSONANCIA_VAL.md","rel":"parte_de"},{"alvo":"virtualizacao","confianca":0.5,"epistemico":"fato","origem":"wiki","rel":"relaciona"},{"alvo":"vida-morte","confianca":0.5,"epistemico":"fato","origem":"wiki","rel":"relaciona"},{"alvo":"word","confianca":0.5,"epistemico":"fato","origem":"wiki","rel":"relaciona"},{"alvo":"sequencia_pow_nativa","confianca":0.5,"epistemico":"fato","origem":"wiki","rel":"relaciona"},{"alvo":"pell","confianca":0.5,"epistemico":"fato","origem":"wiki","rel":"relaciona"},{"alvo":"readme","confianca":0.5,"epistemico":"fato","origem":"wiki","rel":"relaciona"},{"alvo":"camadas_de_tecido","confianca":0.5,"epistemico":"fato","origem":"wiki","rel":"relaciona"},{"alvo":"fibonacci_multidimensional","confianca":0.5,"epistemico":"fato","origem":"wiki","rel":"relaciona"},{"alvo":"broca_os_hub_contraexemplos","confianca":0.5,"epistemico":"fato","origem":"wiki","rel":"relaciona"},{"alvo":"heap","confianca":0.5,"epistemico":"fato","origem":"wiki","rel":"relaciona"},{"alvo":"gentil_12_definicao","confianca":0.5,"epistemico":"fato","origem":"wiki","rel":"relaciona"},{"alvo":"educacao","confianca":0.5,"epistemico":"fato","origem":"wiki","rel":"relaciona"},{"alvo":"incautos_enganados","confianca":0.5,"epistemico":"fato","origem":"wiki","rel":"relaciona"},{"alvo":"pipeline_de_ingestao","confianca":0.5,"epistemico":"fato","origem":"wiki","rel":"relaciona"},{"alvo":"vetores","confianca":0.5,"epistemico":"fato","origem":"wiki","rel":"relaciona"},{"alvo":"kappa","confianca":0.5,"epistemico":"fato","origem":"wiki","rel":"relaciona"},{"alvo":"vontade_de_poder","confianca":0.5,"epistemico":"fato","origem":"wiki","rel":"relaciona"},{"alvo":"semiotica_peirce_triade","confianca":0.5,"epistemico":"fato","origem":"wiki","rel":"relaciona"},{"alvo":"poder_disciplinar_foucault","confianca":0.5,"epistemico":"fato","origem":"wiki","rel":"relaciona"},{"alvo":"lemma_lemunit_e_unitaria_0xt2dtt_-","confianca":0.5,"epistemico":"fato","origem":"wiki","rel":"relaciona"}],"confianca":1.0,"contraexemplos":[],"descricao":"Acoplamento médio fecho: **0.9988** · vizinhos ±1: **0.6007**.","epistemico":"fato","exemplos":[],"id":"2_vizinhos_imediatos_tem_acoplamento_semelhante","memoria":"semantica","meta":{"dominio":"manual","nivel":6,"verificacao":"slug idêntico — enriquecer, não duplicar"},"origem":"POW_METAL_RESSONANCIA_VAL.md","palavras_chave":[],"sinonimos":[],"tipo":"conceito","titulo":"2. Vizinhos imediatos têm acoplamento semelhante?"}
\section{2. Vizinhos imediatos têm acoplamento semelhante?}
Acoplamento médio fecho: **0.9988** \textperiodcentered  vizinhos $\pm$1: **0.6007**.
\prov{relações: parte\_de perguntas; relaciona virtualizacao; relaciona vida-morte; relaciona word; relaciona sequencia\_pow\_nativa; relaciona pell; relaciona readme; relaciona camadas\_de\_tecido; relaciona fibonacci\_multidimensional; relaciona broca\_os\_hub\_contraexemplos; relaciona heap; relaciona gentil\_12\_definicao; relaciona educacao; relaciona incautos\_enganados; relaciona pipeline\_de\_ingestao; relaciona vetores; relaciona kappa; relaciona vontade\_de\_poder; relaciona semiotica\_peirce\_triade; relaciona poder\_disciplinar\_foucault; relaciona lemma\_lemunit\_e\_unitaria\_0xt2dtt\_-}
\prov{proveniência: POW\_METAL\_RESSONANCIA\_VAL.md \textperiodcentered{} fato \textperiodcentered{} confiança 1.0 \textperiodcentered{} domínio manual \textperiodcentered{} história 24 versões no jornal}

%CRISTAL {"arestas":[{"alvo":"opcao_b_desenvolvimento_no_repo","confianca":1.0,"epistemico":"fato","origem":"INSTALL.md","rel":"parte_de"},{"alvo":"broca_os","confianca":1.0,"epistemico":"fato","origem":"manual","rel":"parte_de"}],"confianca":1.0,"contraexemplos":[],"descricao":"```bash cd code && make export PATH=\"$PWD:$PATH\" export TIFFANY_LIB=$PWD/libtiffany/libtiffany.so   # bindings ctypes opcionais ```","epistemico":"fato","exemplos":[],"id":"3_binarios_nativos_performance","memoria":"semantica","meta":{"dominio":"manual","nivel":6,"verificacao":"slug idêntico — enriquecer, não duplicar"},"origem":"INSTALL.md","palavras_chave":[],"sinonimos":[],"tipo":"conceito","titulo":"3. Binários nativos (performance)"}
\section{3. Binários nativos (performance)}
```bash cd code \&\& make export PATH="\$PWD:\$PATH" export TIFFANY\_LIB=\$PWD/libtiffany/libtiffany.so   \# bindings ctypes opcionais ```
\prov{relações: parte\_de opcao\_b\_desenvolvimento\_no\_repo; parte\_de broca\_os}
\prov{proveniência: INSTALL.md \textperiodcentered{} fato \textperiodcentered{} confiança 1.0 \textperiodcentered{} domínio manual \textperiodcentered{} história 10 versões no jornal}

%CRISTAL {"arestas":[{"alvo":"pow_coordenadas_espectrais_broca","confianca":1.0,"epistemico":"fato","origem":"POW_GEOMETRIA.md","rel":"parte_de"},{"alvo":"word","confianca":0.5,"epistemico":"fato","origem":"wiki","rel":"relaciona"},{"alvo":"virtualizacao","confianca":0.5,"epistemico":"fato","origem":"wiki","rel":"relaciona"}],"confianca":1.0,"contraexemplos":[],"descricao":"- regime: **sensivel_caotico** - flips 1-bit suaves (|Δtotal|≤4, |Δcalor|≤0.05): **0/256** - |Δcalor| médio: **0.000001**","epistemico":"fato","exemplos":[],"id":"3_continuidade","memoria":"semantica","meta":{"dominio":"manual","nivel":6,"verificacao":"slug idêntico — enriquecer, não duplicar"},"origem":"POW_GEOMETRIA.md","palavras_chave":[],"sinonimos":[],"tipo":"conceito","titulo":"3. Continuidade"}
\section{3. Continuidade}
- regime: **sensivel\_caotico** - flips 1-bit suaves (|\ensuremath{\Delta}total|\ensuremath{\le}4, |\ensuremath{\Delta}calor|\ensuremath{\le}0.05): **0/256** - |\ensuremath{\Delta}calor| médio: **0.000001**
\prov{relações: parte\_de pow\_coordenadas\_espectrais\_broca; relaciona word; relaciona virtualizacao}
\prov{proveniência: POW\_GEOMETRIA.md \textperiodcentered{} fato \textperiodcentered{} confiança 1.0 \textperiodcentered{} domínio manual \textperiodcentered{} história 8 versões no jornal}

%CRISTAL {"arestas":[{"alvo":"a_prova_da_isa_da_broca","confianca":1.0,"epistemico":"fato","origem":"PROVA_ISA.md","rel":"parte_de"},{"alvo":"borda_contemplacao","confianca":0.5,"epistemico":"fato","origem":"wiki","rel":"relaciona"},{"alvo":"mecanica_quantica_hub","confianca":0.5,"epistemico":"fato","origem":"wiki","rel":"relaciona"},{"alvo":"testes","confianca":0.5,"epistemico":"fato","origem":"wiki","rel":"relaciona"},{"alvo":"gato_eh_um_mamifero","confianca":0.5,"epistemico":"fato","origem":"wiki","rel":"relaciona"},{"alvo":"topologia","confianca":0.5,"epistemico":"fato","origem":"wiki","rel":"relaciona"},{"alvo":"o_que_e_no_projecto","confianca":0.5,"epistemico":"fato","origem":"wiki","rel":"relaciona"},{"alvo":"pedagogia","confianca":0.5,"epistemico":"fato","origem":"wiki","rel":"relaciona"},{"alvo":"compilador","confianca":0.5,"epistemico":"fato","origem":"wiki","rel":"relaciona"}],"confianca":1.0,"contraexemplos":[],"descricao":"``` (fetch, σ) → decode(op) → exec(σ) → σ' ```","epistemico":"fato","exemplos":[],"id":"3_maquina_de_estados","memoria":"semantica","meta":{"dominio":"manual","nivel":6,"verificacao":"slug idêntico — enriquecer, não duplicar"},"origem":"PROVA_ISA.md","palavras_chave":[],"sinonimos":[],"tipo":"conceito","titulo":"3. Máquina de estados"}
\section{3. Máquina de estados}
``` (fetch, \ensuremath{\sigma}) \ensuremath{\to} decode(op) \ensuremath{\to} exec(\ensuremath{\sigma}) \ensuremath{\to} \ensuremath{\sigma}' ```
\prov{relações: parte\_de a\_prova\_da\_isa\_da\_broca; relaciona borda\_contemplacao; relaciona mecanica\_quantica\_hub; relaciona testes; relaciona gato\_eh\_um\_mamifero; relaciona topologia; relaciona o\_que\_e\_no\_projecto; relaciona pedagogia; relaciona compilador}
\prov{proveniência: PROVA\_ISA.md \textperiodcentered{} fato \textperiodcentered{} confiança 1.0 \textperiodcentered{} domínio manual \textperiodcentered{} história 12 versões no jornal}

%CRISTAL {"arestas":[{"alvo":"perguntas","confianca":1.0,"epistemico":"fato","origem":"POW_METAL_PRE_FECHO.md","rel":"parte_de"},{"alvo":"word","confianca":0.5,"epistemico":"fato","origem":"wiki","rel":"relaciona"},{"alvo":"virtualizacao","confianca":0.5,"epistemico":"fato","origem":"wiki","rel":"relaciona"},{"alvo":"escassez","confianca":0.5,"epistemico":"fato","origem":"wiki","rel":"relaciona"},{"alvo":"universidade_curriculo_em_camadas","confianca":0.5,"epistemico":"fato","origem":"wiki","rel":"relaciona"},{"alvo":"vida-morte","confianca":0.5,"epistemico":"fato","origem":"wiki","rel":"relaciona"},{"alvo":"aura_benjamin","confianca":0.5,"epistemico":"fato","origem":"wiki","rel":"relaciona"},{"alvo":"proposition_gerador_de_escala_ipropderiv_com_tddt","confianca":0.5,"epistemico":"fato","origem":"wiki","rel":"relaciona"},{"alvo":"octonios_hiperbolicos","confianca":0.5,"epistemico":"fato","origem":"wiki","rel":"relaciona"},{"alvo":"epistemologia","confianca":0.5,"epistemico":"fato","origem":"wiki","rel":"relaciona"},{"alvo":"semiotica_peirce_triade","confianca":0.5,"epistemico":"fato","origem":"wiki","rel":"relaciona"}],"confianca":1.0,"contraexemplos":[],"descricao":"Variância Δatrator |dist|≤16: **33877.8** · |dist|≥48: **38681.7** · estabiliza: **sim**.","epistemico":"fato","exemplos":[],"id":"3_o_atrator_estabiliza_antes_do_fecho","memoria":"semantica","meta":{"dominio":"manual","nivel":6,"verificacao":"slug idêntico — enriquecer, não duplicar"},"origem":"POW_METAL_PRE_FECHO.md","palavras_chave":[],"sinonimos":[],"tipo":"conceito","titulo":"3. O atrator estabiliza antes do fecho?"}
\section{3. O atrator estabiliza antes do fecho?}
Variância \ensuremath{\Delta}atrator |dist|\ensuremath{\le}16: **33877.8** \textperiodcentered  |dist|\ensuremath{\ge}48: **38681.7** \textperiodcentered  estabiliza: **sim**.
\prov{relações: parte\_de perguntas; relaciona word; relaciona virtualizacao; relaciona escassez; relaciona universidade\_curriculo\_em\_camadas; relaciona vida-morte; relaciona aura\_benjamin; relaciona proposition\_gerador\_de\_escala\_ipropderiv\_com\_tddt; relaciona octonios\_hiperbolicos; relaciona epistemologia; relaciona semiotica\_peirce\_triade}
\prov{proveniência: POW\_METAL\_PRE\_FECHO.md \textperiodcentered{} fato \textperiodcentered{} confiança 1.0 \textperiodcentered{} domínio manual \textperiodcentered{} história 14 versões no jornal}

%CRISTAL {"arestas":[{"alvo":"broca_os","confianca":0.95,"epistemico":"fato","origem":"paper","rel":"parte_de"}],"confianca":1.0,"contraexemplos":[],"descricao":"| Construção | Status | Como / obstrução | |------------|--------|------------------| | **if** | ✅ implementado | `CMP; JZ/JNZ; blocos` | | **else** | ✅ implementado | salto sobre ramo alternativo | | **while** | ✅ implementado | `JMP` para trás + contador scratch | | **for** | ⚙️ montador | desenrolar ou contador + `JNZ` | | **contador** | ✅ implementado | scratch + `LOAD 1; ADD; STORE` | | **acumulador** | ✅ implementado | scratch 8000 (provado) |","epistemico":"fato","exemplos":[],"id":"3_poder_de_expressao_tabela","memoria":"semantica","meta":{"dominio":"manual","nivel":6,"verificacao":"slug idêntico — enriquecer, não duplicar"},"origem":"EXPRESSIVIDADE.md","palavras_chave":[],"sinonimos":[],"tipo":"conceito","titulo":"3. Poder de Expressão — tabela"}
\section{3. Poder de Expressão — tabela}
| Construção | Status | Como / obstrução | |------------|--------|------------------| | **if** | [ok] implementado | `CMP; JZ/JNZ; blocos` | | **else** | [ok] implementado | salto sobre ramo alternativo | | **while** | [ok] implementado | `JMP` para trás + contador scratch | | **for** | [engr.] montador | desenrolar ou contador + `JNZ` | | **contador** | [ok] implementado | scratch + `LOAD 1; ADD; STORE` | | **acumulador** | [ok] implementado | scratch 8000 (provado) |
\prov{relações: parte\_de broca\_os}
\prov{proveniência: EXPRESSIVIDADE.md \textperiodcentered{} fato \textperiodcentered{} confiança 1.0 \textperiodcentered{} domínio manual \textperiodcentered{} história 3 versões no jornal}

%CRISTAL {"arestas":[{"alvo":"perguntas","confianca":1.0,"epistemico":"fato","origem":"POW_METAL_PRE_FECHO.md","rel":"parte_de"},{"alvo":"word","confianca":0.5,"epistemico":"fato","origem":"wiki","rel":"relaciona"},{"alvo":"virtualizacao","confianca":0.5,"epistemico":"fato","origem":"wiki","rel":"relaciona"},{"alvo":"vida-morte","confianca":0.5,"epistemico":"fato","origem":"wiki","rel":"relaciona"},{"alvo":"bai_quinteto_hub","confianca":0.5,"epistemico":"fato","origem":"wiki","rel":"relaciona"},{"alvo":"etica_protestante","confianca":0.5,"epistemico":"fato","origem":"wiki","rel":"relaciona"},{"alvo":"misticismo_nao_dual","confianca":0.5,"epistemico":"fato","origem":"wiki","rel":"relaciona"},{"alvo":"apego","confianca":0.5,"epistemico":"fato","origem":"wiki","rel":"relaciona"},{"alvo":"estrutura_de_dados","confianca":0.5,"epistemico":"fato","origem":"wiki","rel":"relaciona"},{"alvo":"vontade_de_poder","confianca":0.5,"epistemico":"fato","origem":"wiki","rel":"relaciona"},{"alvo":"linguagem","confianca":0.5,"epistemico":"fato","origem":"wiki","rel":"relaciona"},{"alvo":"metrica_parabolica","confianca":0.5,"epistemico":"fato","origem":"wiki","rel":"relaciona"},{"alvo":"5_a_lemniscata_aproxima-se_do_equilibrio_antes_do_fecho","confianca":0.5,"epistemico":"fato","origem":"wiki","rel":"relaciona"},{"alvo":"estetica","confianca":0.5,"epistemico":"fato","origem":"wiki","rel":"relaciona"},{"alvo":"gentil_capitulo_3_sequencias_geometricas_de_ordem_m","confianca":0.5,"epistemico":"fato","origem":"wiki","rel":"relaciona"},{"alvo":"musica","confianca":0.5,"epistemico":"fato","origem":"wiki","rel":"relaciona"}],"confianca":1.0,"contraexemplos":[],"descricao":"Salto máximo consecutivo: **0.0001** · médio: **0.0000** · contínua: **sim**.","epistemico":"fato","exemplos":[],"id":"4_a_coordenada_koch_permanece_continua","memoria":"semantica","meta":{"dominio":"manual","nivel":6,"verificacao":"slug idêntico — enriquecer, não duplicar"},"origem":"POW_METAL_PRE_FECHO.md","palavras_chave":[],"sinonimos":[],"tipo":"conceito","titulo":"4. A coordenada Koch permanece contínua?"}
\section{4. A coordenada Koch permanece contínua?}
Salto máximo consecutivo: **0.0001** \textperiodcentered  médio: **0.0000** \textperiodcentered  contínua: **sim**.
\prov{relações: parte\_de perguntas; relaciona word; relaciona virtualizacao; relaciona vida-morte; relaciona bai\_quinteto\_hub; relaciona etica\_protestante; relaciona misticismo\_nao\_dual; relaciona apego; relaciona estrutura\_de\_dados; relaciona vontade\_de\_poder; relaciona linguagem; relaciona metrica\_parabolica; relaciona 5\_a\_lemniscata\_aproxima-se\_do\_equilibrio\_antes\_do\_fecho; relaciona estetica; relaciona gentil\_capitulo\_3\_sequencias\_geometricas\_de\_ordem\_m; relaciona musica}
\prov{proveniência: POW\_METAL\_PRE\_FECHO.md \textperiodcentered{} fato \textperiodcentered{} confiança 1.0 \textperiodcentered{} domínio manual \textperiodcentered{} história 19 versões no jornal}

%CRISTAL {"arestas":[{"alvo":"perguntas","confianca":1.0,"epistemico":"fato","origem":"POW_METAL_RESSONANCIA_VAL.md","rel":"parte_de"},{"alvo":"grupos","confianca":0.5,"epistemico":"fato","origem":"wiki","rel":"relaciona"},{"alvo":"recursao","confianca":0.5,"epistemico":"fato","origem":"wiki","rel":"relaciona"},{"alvo":"testes","confianca":0.5,"epistemico":"fato","origem":"wiki","rel":"relaciona"},{"alvo":"fase_2_tecido_de_conhecimento","confianca":0.5,"epistemico":"fato","origem":"wiki","rel":"relaciona"},{"alvo":"funcao","confianca":0.5,"epistemico":"fato","origem":"wiki","rel":"relaciona"},{"alvo":"ruido","confianca":0.5,"epistemico":"fato","origem":"wiki","rel":"relaciona"},{"alvo":"utilitarismo","confianca":0.5,"epistemico":"fato","origem":"wiki","rel":"relaciona"},{"alvo":"teste_ordem","confianca":0.5,"epistemico":"fato","origem":"wiki","rel":"relaciona"},{"alvo":"monotono","confianca":0.5,"epistemico":"fato","origem":"wiki","rel":"relaciona"},{"alvo":"melhor_candidato","confianca":0.5,"epistemico":"fato","origem":"wiki","rel":"relaciona"},{"alvo":"regua_associativa_e_o_risco_de_vazamento","confianca":0.5,"epistemico":"fato","origem":"wiki","rel":"relaciona"},{"alvo":"5_modos_estaveis_quando_nonce_muda_1","confianca":0.5,"epistemico":"fato","origem":"wiki","rel":"relaciona"},{"alvo":"consulta_simples","confianca":0.5,"epistemico":"fato","origem":"wiki","rel":"relaciona"},{"alvo":"topologia","confianca":0.5,"epistemico":"fato","origem":"wiki","rel":"relaciona"},{"alvo":"a_escala_uma_volta_no_c_rculo","confianca":0.5,"epistemico":"fato","origem":"wiki","rel":"relaciona"}],"confianca":1.0,"contraexemplos":[],"descricao":"Fechos com perfil unimodal: **0/2**.","epistemico":"fato","exemplos":[],"id":"4_acoplamento_cresce_monotonicamente_ate_o_fecho","memoria":"semantica","meta":{"dominio":"manual","nivel":6,"verificacao":"slug idêntico — enriquecer, não duplicar"},"origem":"POW_METAL_RESSONANCIA_VAL.md","palavras_chave":[],"sinonimos":[],"tipo":"conceito","titulo":"4. Acoplamento cresce monotonicamente até o fecho?"}
\section{4. Acoplamento cresce monotonicamente até o fecho?}
Fechos com perfil unimodal: **0/2**.
\prov{relações: parte\_de perguntas; relaciona grupos; relaciona recursao; relaciona testes; relaciona fase\_2\_tecido\_de\_conhecimento; relaciona funcao; relaciona ruido; relaciona utilitarismo; relaciona teste\_ordem; relaciona monotono; relaciona melhor\_candidato; relaciona regua\_associativa\_e\_o\_risco\_de\_vazamento; relaciona 5\_modos\_estaveis\_quando\_nonce\_muda\_1; relaciona consulta\_simples; relaciona topologia; relaciona a\_escala\_uma\_volta\_no\_c\_rculo}
\prov{proveniência: POW\_METAL\_RESSONANCIA\_VAL.md \textperiodcentered{} fato \textperiodcentered{} confiança 1.0 \textperiodcentered{} domínio manual \textperiodcentered{} história 19 versões no jornal}

%CRISTAL {"arestas":[{"alvo":"opcao_b_desenvolvimento_no_repo","confianca":1.0,"epistemico":"fato","origem":"INSTALL.md","rel":"parte_de"},{"alvo":"broca_os","confianca":1.0,"epistemico":"fato","origem":"manual","rel":"parte_de"}],"confianca":1.0,"contraexemplos":[],"descricao":"```bash cd code python3 tiffany_platform.py pack-cliente tar -xzf ../dist/tiffany-client-template.tgz -C /tmp cp -r /tmp/tiffany-client-template/cliente $TIFFANY_HOME/","epistemico":"fato","exemplos":[],"id":"4_base_cliente_opcional","memoria":"semantica","meta":{"dominio":"manual","nivel":6,"verificacao":"slug idêntico — enriquecer, não duplicar"},"origem":"INSTALL.md","palavras_chave":[],"sinonimos":[],"tipo":"conceito","titulo":"4. Base Cliente (opcional)"}
\section{4. Base Cliente (opcional)}
```bash cd code python3 tiffany\_platform.py pack-cliente tar -xzf ../dist/tiffany-client-template.tgz -C /tmp cp -r /tmp/tiffany-client-template/cliente \$TIFFANY\_HOME/
\prov{relações: parte\_de opcao\_b\_desenvolvimento\_no\_repo; parte\_de broca\_os}
\prov{proveniência: INSTALL.md \textperiodcentered{} fato \textperiodcentered{} confiança 1.0 \textperiodcentered{} domínio manual \textperiodcentered{} história 10 versões no jornal}

%CRISTAL {"arestas":[{"alvo":"pow_coordenadas_espectrais_broca","confianca":1.0,"epistemico":"fato","origem":"POW_GEOMETRIA.md","rel":"parte_de"},{"alvo":"transformacao_linear","confianca":1.0,"epistemico":"fato","origem":"manual","rel":"referencia"}],"confianca":1.0,"contraexemplos":[],"descricao":"| métrica | espectral | linear | |---------|-----------|--------| | entropia calor | 1.3849 | 2.0085 | | autocorr | 0.7503 | 1.0000 | | dim. efetiva (exp H) | 3.99 | — |","epistemico":"fato","exemplos":[],"id":"4_compressao_espectral_vs_linear","memoria":"semantica","meta":{"dominio":"manual","nivel":6,"verificacao":"slug idêntico — enriquecer, não duplicar"},"origem":"POW_GEOMETRIA.md","palavras_chave":[],"sinonimos":[],"tipo":"conceito","titulo":"4. Compressão espectral vs linear"}
\section{4. Compressão espectral vs linear}
| métrica | espectral | linear | |---------|-----------|--------| | entropia calor | 1.3849 | 2.0085 | | autocorr | 0.7503 | 1.0000 | | dim. efetiva (exp H) | 3.99 | — |
\prov{relações: parte\_de pow\_coordenadas\_espectrais\_broca; referencia transformacao\_linear}
\prov{proveniência: POW\_GEOMETRIA.md \textperiodcentered{} fato \textperiodcentered{} confiança 1.0 \textperiodcentered{} domínio manual \textperiodcentered{} história 59 versões no jornal}

%CRISTAL {"arestas":[{"alvo":"a_prova_da_isa_da_broca","confianca":1.0,"epistemico":"fato","origem":"PROVA_ISA.md","rel":"parte_de"},{"alvo":"word","confianca":0.5,"epistemico":"fato","origem":"wiki","rel":"relaciona"}],"confianca":1.0,"contraexemplos":[],"descricao":"``` arquivo ──neurônio──► word ──LOAD──► A │ B◄──empilha──┘ │ ADD/SUB/… │ R ──STORE──► scratch ```","epistemico":"fato","exemplos":[],"id":"4_modelo_operacional","memoria":"semantica","meta":{"dominio":"manual","nivel":6,"verificacao":"slug idêntico — enriquecer, não duplicar"},"origem":"PROVA_ISA.md","palavras_chave":[],"sinonimos":[],"tipo":"conceito","titulo":"4. Modelo operacional"}
\section{4. Modelo operacional}
``` arquivo --neurônio--$\triangleright$ word --LOAD--$\triangleright$ A | B$\triangleleft$--empilha-- | ADD/SUB/… | R --STORE--$\triangleright$ scratch ```
\prov{relações: parte\_de a\_prova\_da\_isa\_da\_broca; relaciona word}
\prov{proveniência: PROVA\_ISA.md \textperiodcentered{} fato \textperiodcentered{} confiança 1.0 \textperiodcentered{} domínio manual \textperiodcentered{} história 5 versões no jornal}

%CRISTAL {"arestas":[{"alvo":"computacao","confianca":1.0,"epistemico":"fato","origem":"manual","rel":"parte_de"},{"alvo":"historia","confianca":0.5,"epistemico":"fato","origem":"wiki","rel":"relaciona"},{"alvo":"utilitarismo","confianca":0.5,"epistemico":"fato","origem":"wiki","rel":"relaciona"},{"alvo":"metafisica","confianca":0.5,"epistemico":"fato","origem":"wiki","rel":"relaciona"},{"alvo":"virtualizacao","confianca":0.5,"epistemico":"fato","origem":"wiki","rel":"relaciona"},{"alvo":"word","confianca":0.5,"epistemico":"fato","origem":"wiki","rel":"relaciona"}],"confianca":1.0,"contraexemplos":[],"descricao":"Secção de EXPRESSIVIDADE.md.","epistemico":"fato","exemplos":[],"id":"4_obstrucao_formal_nao_turing-completa","memoria":"semantica","meta":{"dominio":"manual","nivel":6,"verificacao":"slug idêntico — enriquecer, não duplicar"},"origem":"EXPRESSIVIDADE.md","palavras_chave":[],"sinonimos":[],"tipo":"conceito","titulo":"4. Obstrução Formal (não Turing-completa)"}
\section{4. Obstrução Formal (não Turing-completa)}
Secção de EXPRESSIVIDADE.md.
\prov{relações: parte\_de computacao; relaciona historia; relaciona utilitarismo; relaciona metafisica; relaciona virtualizacao; relaciona word}
\prov{proveniência: EXPRESSIVIDADE.md \textperiodcentered{} fato \textperiodcentered{} confiança 1.0 \textperiodcentered{} domínio manual \textperiodcentered{} história 13 versões no jornal}

%CRISTAL {"arestas":[{"alvo":"perguntas","confianca":1.0,"epistemico":"fato","origem":"POW_METAL_PRE_FECHO.md","rel":"parte_de"},{"alvo":"vida-morte","confianca":0.5,"epistemico":"fato","origem":"wiki","rel":"relaciona"},{"alvo":"word","confianca":0.5,"epistemico":"fato","origem":"wiki","rel":"relaciona"},{"alvo":"politica","confianca":0.5,"epistemico":"fato","origem":"wiki","rel":"relaciona"}],"confianca":1.0,"contraexemplos":[],"descricao":"|calor+c²−1| |dist|≤8: **0.000000** · |dist|≥32: **0.000000** · aproxima: **sim**.","epistemico":"fato","exemplos":[],"id":"5_a_lemniscata_aproxima-se_do_equilibrio_antes_do_fecho","memoria":"semantica","meta":{"dominio":"manual","nivel":6,"verificacao":"slug idêntico — enriquecer, não duplicar"},"origem":"POW_METAL_PRE_FECHO.md","palavras_chave":[],"sinonimos":[],"tipo":"conceito","titulo":"5. A lemniscata aproxima-se do equilíbrio antes do fecho?"}
\section{5. A lemniscata aproxima-se do equilíbrio antes do fecho?}
|calor+c\ensuremath{^{2}}\ensuremath{-}1| |dist|\ensuremath{\le}8: **0.000000** \textperiodcentered  |dist|\ensuremath{\ge}32: **0.000000** \textperiodcentered  aproxima: **sim**.
\prov{relações: parte\_de perguntas; relaciona vida-morte; relaciona word; relaciona politica}
\prov{proveniência: POW\_METAL\_PRE\_FECHO.md \textperiodcentered{} fato \textperiodcentered{} confiança 1.0 \textperiodcentered{} domínio manual \textperiodcentered{} história 7 versões no jornal}

%CRISTAL {"arestas":[{"alvo":"pow_coordenadas_espectrais_broca","confianca":1.0,"epistemico":"fato","origem":"POW_GEOMETRIA.md","rel":"parte_de"}],"confianca":1.0,"contraexemplos":[],"descricao":"| n | entropia metal | fit | |---|----------------|-----| | 1 | 2.5695 | 1.0000 | | 2 | 2.7392 | 1.0000 | | 3 | 2.7291 | 1.0000 | | 4 | 2.6454 | 1.0000 | | 5 | 2.5122 | 1.0000 | | 6 | 2.4830 | 1.0000 | | 7 | 2.3970 | 1.0000 | | 8 | 2.3032 | 1.0000 |","epistemico":"fato","exemplos":[],"id":"5_familias_metalicas_dados_escolhem","memoria":"semantica","meta":{"dominio":"manual","nivel":6,"verificacao":"slug idêntico — enriquecer, não duplicar"},"origem":"POW_GEOMETRIA.md","palavras_chave":[],"sinonimos":[],"tipo":"conceito","titulo":"5. Famílias metálicas — dados escolhem"}
\section{5. Famílias metálicas — dados escolhem}
| n | entropia metal | fit | |---|----------------|-----| | 1 | 2.5695 | 1.0000 | | 2 | 2.7392 | 1.0000 | | 3 | 2.7291 | 1.0000 | | 4 | 2.6454 | 1.0000 | | 5 | 2.5122 | 1.0000 | | 6 | 2.4830 | 1.0000 | | 7 | 2.3970 | 1.0000 | | 8 | 2.3032 | 1.0000 |
\prov{relações: parte\_de pow\_coordenadas\_espectrais\_broca}
\prov{proveniência: POW\_GEOMETRIA.md \textperiodcentered{} fato \textperiodcentered{} confiança 1.0 \textperiodcentered{} domínio manual \textperiodcentered{} história 6 versões no jornal}

%CRISTAL {"arestas":[{"alvo":"perguntas","confianca":1.0,"epistemico":"fato","origem":"POW_METAL_RESSONANCIA_VAL.md","rel":"parte_de"},{"alvo":"melhor_candidato","confianca":0.5,"epistemico":"fato","origem":"wiki","rel":"relaciona"},{"alvo":"compilador","confianca":0.5,"epistemico":"fato","origem":"wiki","rel":"relaciona"},{"alvo":"teoria_do_valor","confianca":0.5,"epistemico":"fato","origem":"wiki","rel":"relaciona"},{"alvo":"word","confianca":0.5,"epistemico":"fato","origem":"wiki","rel":"relaciona"}],"confianca":1.0,"contraexemplos":[],"descricao":"|Δ acopl|±1 na janela: **0.26554** · fora: **0.05486**.","epistemico":"fato","exemplos":[],"id":"5_modos_estaveis_quando_nonce_muda_1","memoria":"semantica","meta":{"dominio":"manual","nivel":6,"verificacao":"slug idêntico — enriquecer, não duplicar"},"origem":"POW_METAL_RESSONANCIA_VAL.md","palavras_chave":[],"sinonimos":[],"tipo":"conceito","titulo":"5. Modos estáveis quando nonce muda ±1?"}
\section{5. Modos estáveis quando nonce muda $\pm$1?}
|\ensuremath{\Delta} acopl|$\pm$1 na janela: **0.26554** \textperiodcentered  fora: **0.05486**.
\prov{relações: parte\_de perguntas; relaciona melhor\_candidato; relaciona compilador; relaciona teoria\_do\_valor; relaciona word}
\prov{proveniência: POW\_METAL\_RESSONANCIA\_VAL.md \textperiodcentered{} fato \textperiodcentered{} confiança 1.0 \textperiodcentered{} domínio manual \textperiodcentered{} história 8 versões no jornal}

%CRISTAL {"arestas":[{"alvo":"a_prova_da_isa_da_broca","confianca":1.0,"epistemico":"fato","origem":"PROVA_ISA.md","rel":"parte_de"}],"confianca":1.0,"contraexemplos":[],"descricao":"| # | Ficheiro | Composição | |---|----------|------------| | 1 | `programas/01_soma.asm` | `LOAD LOAD ADD STORE` | | 2 | `programas/02_igualdade.asm` | `SUB; LOAD zero; CMP; JNZ` | | 3 | `programas/03_loop.asm` | gerado: `CMP/JZ` por slot + sentinela | | 4 | `programas/04_acumulador.asm` | gerado: acumulador via scratch 8000 | | 5 | `programas/05_pipeline.asm` | `LOAD ADD STORE` | | 6 | `programas/06_tesouraria.asm` | assinatura→SUB→CMP→checkpoint |","epistemico":"fato","exemplos":[],"id":"5_programas_de_referencia","memoria":"semantica","meta":{"dominio":"manual","nivel":6,"verificacao":"slug idêntico — enriquecer, não duplicar"},"origem":"PROVA_ISA.md","palavras_chave":[],"sinonimos":[],"tipo":"conceito","titulo":"5. Programas de referência"}
\section{5. Programas de referência}
| \# | Ficheiro | Composição | |---|----------|------------| | 1 | `programas/01\_soma.asm` | `LOAD LOAD ADD STORE` | | 2 | `programas/02\_igualdade.asm` | `SUB; LOAD zero; CMP; JNZ` | | 3 | `programas/03\_loop.asm` | gerado: `CMP/JZ` por slot + sentinela | | 4 | `programas/04\_acumulador.asm` | gerado: acumulador via scratch 8000 | | 5 | `programas/05\_pipeline.asm` | `LOAD ADD STORE` | | 6 | `programas/06\_tesouraria.asm` | assinatura\ensuremath{\to}SUB\ensuremath{\to}CMP\ensuremath{\to}checkpoint |
\prov{relações: parte\_de a\_prova\_da\_isa\_da\_broca}
\prov{proveniência: PROVA\_ISA.md \textperiodcentered{} fato \textperiodcentered{} confiança 1.0 \textperiodcentered{} domínio manual \textperiodcentered{} história 4 versões no jornal}

%CRISTAL {"arestas":[{"alvo":"6_provas_de_composicao","confianca":1.0,"epistemico":"fato","origem":"PROVA_ISA.md","rel":"parte_de"},{"alvo":"virtualizacao","confianca":0.5,"epistemico":"fato","origem":"wiki","rel":"relaciona"},{"alvo":"word","confianca":0.5,"epistemico":"fato","origem":"wiki","rel":"relaciona"},{"alvo":"vida-morte","confianca":0.5,"epistemico":"fato","origem":"wiki","rel":"relaciona"},{"alvo":"semiotica_peirce_triade","confianca":0.5,"epistemico":"fato","origem":"wiki","rel":"relaciona"},{"alvo":"proposition_gerador_de_escala_ipropderiv_com_tddt","confianca":0.5,"epistemico":"fato","origem":"wiki","rel":"relaciona"},{"alvo":"c","confianca":0.5,"epistemico":"fato","origem":"wiki","rel":"relaciona"},{"alvo":"escassez","confianca":0.5,"epistemico":"fato","origem":"wiki","rel":"relaciona"},{"alvo":"flags","confianca":0.5,"epistemico":"fato","origem":"wiki","rel":"relaciona"},{"alvo":"goldchain_mercado_hub_fontes","confianca":0.5,"epistemico":"fato","origem":"wiki","rel":"relaciona"},{"alvo":"proposition_convolucao_multiplicativa_produtopropconv","confianca":0.5,"epistemico":"fato","origem":"wiki","rel":"relaciona"}],"confianca":1.0,"contraexemplos":[],"descricao":"**Teorema:** `LOAD a; LOAD b; ADD` produz `neuron(a) ⊕ neuron(b)` onde `⊕` é soma componente.","epistemico":"fato","exemplos":[],"id":"61_soma_programa_1","memoria":"semantica","meta":{"dominio":"manual","nivel":6,"verificacao":"slug idêntico — enriquecer, não duplicar"},"origem":"PROVA_ISA.md","palavras_chave":[],"sinonimos":[],"tipo":"conceito","titulo":"6.1 Soma (Programa 1)"}
\section{6.1 Soma (Programa 1)}
**Teorema:** `LOAD a; LOAD b; ADD` produz `neuron(a) \ensuremath{\oplus} neuron(b)` onde `\ensuremath{\oplus}` é soma componente.
\prov{relações: parte\_de 6\_provas\_de\_composicao; relaciona virtualizacao; relaciona word; relaciona vida-morte; relaciona semiotica\_peirce\_triade; relaciona proposition\_gerador\_de\_escala\_ipropderiv\_com\_tddt; relaciona c; relaciona escassez; relaciona flags; relaciona goldchain\_mercado\_hub\_fontes; relaciona proposition\_convolucao\_multiplicativa\_produtopropconv}
\prov{proveniência: PROVA\_ISA.md \textperiodcentered{} fato \textperiodcentered{} confiança 1.0 \textperiodcentered{} domínio manual \textperiodcentered{} história 14 versões no jornal}

%CRISTAL {"arestas":[{"alvo":"6_provas_de_composicao","confianca":1.0,"epistemico":"fato","origem":"PROVA_ISA.md","rel":"parte_de"},{"alvo":"regua_associativa_e_o_risco_de_vazamento","confianca":0.5,"epistemico":"fato","origem":"wiki","rel":"relaciona"},{"alvo":"utilitarismo","confianca":0.5,"epistemico":"fato","origem":"wiki","rel":"relaciona"},{"alvo":"testes","confianca":0.5,"epistemico":"fato","origem":"wiki","rel":"relaciona"},{"alvo":"word","confianca":0.5,"epistemico":"fato","origem":"wiki","rel":"relaciona"}],"confianca":1.0,"contraexemplos":[],"descricao":"**Problema:** `JNZ` testa `ZERO`, não `EQ`.","epistemico":"fato","exemplos":[],"id":"62_igualdade_programa_2","memoria":"semantica","meta":{"dominio":"manual","nivel":6,"verificacao":"slug idêntico — enriquecer, não duplicar"},"origem":"PROVA_ISA.md","palavras_chave":[],"sinonimos":[],"tipo":"conceito","titulo":"6.2 Igualdade (Programa 2)"}
\section{6.2 Igualdade (Programa 2)}
**Problema:** `JNZ` testa `ZERO`, não `EQ`.
\prov{relações: parte\_de 6\_provas\_de\_composicao; relaciona regua\_associativa\_e\_o\_risco\_de\_vazamento; relaciona utilitarismo; relaciona testes; relaciona word}
\prov{proveniência: PROVA\_ISA.md \textperiodcentered{} fato \textperiodcentered{} confiança 1.0 \textperiodcentered{} domínio manual \textperiodcentered{} história 8 versões no jornal}

%CRISTAL {"arestas":[{"alvo":"6_provas_de_composicao","confianca":1.0,"epistemico":"fato","origem":"PROVA_ISA.md","rel":"parte_de"},{"alvo":"virtualizacao","confianca":0.5,"epistemico":"fato","origem":"wiki","rel":"relaciona"},{"alvo":"word","confianca":0.5,"epistemico":"fato","origem":"wiki","rel":"relaciona"}],"confianca":1.0,"contraexemplos":[],"descricao":"Contrato e referência → `SUB` → `CMP` com zero → checkpoint `8001`: - `[1,0]` aprovado - `[0,1]` rejeitado","epistemico":"fato","exemplos":[],"id":"64_tesouraria_programa_6","memoria":"semantica","meta":{"dominio":"manual","nivel":6,"verificacao":"slug idêntico — enriquecer, não duplicar"},"origem":"PROVA_ISA.md","palavras_chave":[],"sinonimos":[],"tipo":"conceito","titulo":"6.4 Tesouraria (Programa 6)"}
\section{6.4 Tesouraria (Programa 6)}
Contrato e referência \ensuremath{\to} `SUB` \ensuremath{\to} `CMP` com zero \ensuremath{\to} checkpoint `8001`: - `[1,0]` aprovado - `[0,1]` rejeitado
\prov{relações: parte\_de 6\_provas\_de\_composicao; relaciona virtualizacao; relaciona word}
\prov{proveniência: PROVA\_ISA.md \textperiodcentered{} fato \textperiodcentered{} confiança 1.0 \textperiodcentered{} domínio manual \textperiodcentered{} história 6 versões no jornal}

%CRISTAL {"arestas":[{"alvo":"6_provas_de_composicao","confianca":1.0,"epistemico":"fato","origem":"PROVA_ISA.md","rel":"parte_de"},{"alvo":"vontade_de_poder","confianca":0.5,"epistemico":"fato","origem":"wiki","rel":"relaciona"},{"alvo":"word","confianca":0.5,"epistemico":"fato","origem":"wiki","rel":"relaciona"},{"alvo":"virtualizacao","confianca":0.5,"epistemico":"fato","origem":"wiki","rel":"relaciona"}],"confianca":1.0,"contraexemplos":[],"descricao":"**Composição obrigatória:** `LOAD x; ADD; STORE` para persistir espectro em scratch. ∎","epistemico":"fato","exemplos":[],"id":"65_store_grava_r_nao_a","memoria":"semantica","meta":{"dominio":"manual","nivel":6,"verificacao":"slug idêntico — enriquecer, não duplicar"},"origem":"PROVA_ISA.md","palavras_chave":[],"sinonimos":[],"tipo":"conceito","titulo":"6.5 STORE grava R (não A)"}
\section{6.5 STORE grava R (não A)}
**Composição obrigatória:** `LOAD x; ADD; STORE` para persistir espectro em scratch. \ensuremath{\blacksquare}
\prov{relações: parte\_de 6\_provas\_de\_composicao; relaciona vontade\_de\_poder; relaciona word; relaciona virtualizacao}
\prov{proveniência: PROVA\_ISA.md \textperiodcentered{} fato \textperiodcentered{} confiança 1.0 \textperiodcentered{} domínio manual \textperiodcentered{} história 7 versões no jornal}

%CRISTAL {"arestas":[{"alvo":"geometria","confianca":0.211,"epistemico":"fato","origem":"manual","rel":"referencia"},{"alvo":"bai","confianca":0.85,"epistemico":"fato","origem":"manual","rel":"referencia"},{"alvo":"pow_coordenadas_espectrais_broca","confianca":1.0,"epistemico":"fato","origem":"POW_GEOMETRIA.md","rel":"parte_de"}],"confianca":1.0,"contraexemplos":[],"descricao":"- nonce: 1572863 - word: `[295, 143]` - ouro: `[438, 295]` - gauss_real: `[295, 143]` - gauss_ouro: `[438, 295]` - gauss_koch: `[367, 384]` - atrator: **672** - lemniscata soma: 1.000000000000 - cifra: `[733, 438]`","epistemico":"fato","exemplos":[],"id":"6_geometria_da_solucao_calor_baixo","memoria":"semantica","meta":{"dominio":"manual","duplicata_sugerida":[],"nivel":6,"verificacao":"parte de conceito mais amplo '6_geometria_da_solucao_calor_baixo'"},"origem":"POW_METAL.md","palavras_chave":[],"sinonimos":[],"tipo":"conceito","titulo":"6. Geometria da solução (calor baixo)"}
\section{6. Geometria da solução (calor baixo)}
- nonce: 1572863 - word: `[295, 143]` - ouro: `[438, 295]` - gauss\_real: `[295, 143]` - gauss\_ouro: `[438, 295]` - gauss\_koch: `[367, 384]` - atrator: **672** - lemniscata soma: 1.000000000000 - cifra: `[733, 438]`
\prov{relações: referencia geometria; referencia bai; parte\_de pow\_coordenadas\_espectrais\_broca}
\prov{proveniência: POW\_METAL.md \textperiodcentered{} fato \textperiodcentered{} confiança 1.0 \textperiodcentered{} domínio manual \textperiodcentered{} história 18 versões no jornal}

%CRISTAL {"arestas":[{"alvo":"a_prova_da_isa_da_broca","confianca":1.0,"epistemico":"fato","origem":"PROVA_ISA.md","rel":"parte_de"},{"alvo":"8_release_core","confianca":0.5,"epistemico":"fato","origem":"wiki","rel":"relaciona"},{"alvo":"topologia","confianca":0.5,"epistemico":"fato","origem":"wiki","rel":"relaciona"},{"alvo":"aneis","confianca":0.5,"epistemico":"fato","origem":"wiki","rel":"relaciona"},{"alvo":"historia","confianca":0.5,"epistemico":"fato","origem":"wiki","rel":"relaciona"},{"alvo":"cache","confianca":0.5,"epistemico":"fato","origem":"wiki","rel":"relaciona"},{"alvo":"broca_so","confianca":0.5,"epistemico":"fato","origem":"wiki","rel":"relaciona"},{"alvo":"gato_operador_hub","confianca":0.5,"epistemico":"fato","origem":"wiki","rel":"relaciona"},{"alvo":"hopfield_rede_hub","confianca":0.5,"epistemico":"fato","origem":"wiki","rel":"relaciona"},{"alvo":"kernel","confianca":0.5,"epistemico":"fato","origem":"wiki","rel":"relaciona"},{"alvo":"complexos","confianca":0.5,"epistemico":"fato","origem":"wiki","rel":"relaciona"},{"alvo":"recursao","confianca":0.5,"epistemico":"fato","origem":"wiki","rel":"relaciona"},{"alvo":"algebra_dos_contratos","confianca":0.5,"epistemico":"fato","origem":"wiki","rel":"relaciona"},{"alvo":"expressividade","confianca":0.5,"epistemico":"fato","origem":"wiki","rel":"relaciona"},{"alvo":"memoria_virtual","confianca":0.5,"epistemico":"fato","origem":"wiki","rel":"relaciona"},{"alvo":"goldchain_mercado_hub","confianca":0.5,"epistemico":"fato","origem":"wiki","rel":"relaciona"},{"alvo":"linker","confianca":0.5,"epistemico":"fato","origem":"wiki","rel":"relaciona"},{"alvo":"resposta_a_pergunta_dificil","confianca":0.5,"epistemico":"fato","origem":"wiki","rel":"relaciona"},{"alvo":"etica","confianca":0.5,"epistemico":"fato","origem":"wiki","rel":"relaciona"},{"alvo":"pow_metal_pre_fecho","confianca":0.5,"epistemico":"fato","origem":"wiki","rel":"relaciona"},{"alvo":"processos","confianca":0.5,"epistemico":"fato","origem":"wiki","rel":"relaciona"},{"alvo":"java","confianca":0.5,"epistemico":"fato","origem":"wiki","rel":"relaciona"},{"alvo":"parabola_pa","confianca":0.5,"epistemico":"fato","origem":"wiki","rel":"relaciona"},{"alvo":"regua_associativa_e_o_risco_de_vazamento","confianca":0.5,"epistemico":"fato","origem":"wiki","rel":"relaciona"},{"alvo":"kappa","confianca":0.5,"epistemico":"fato","origem":"wiki","rel":"relaciona"},{"alvo":"compilador","confianca":0.5,"epistemico":"fato","origem":"wiki","rel":"relaciona"},{"alvo":"projeto_cristaljson","confianca":0.5,"epistemico":"fato","origem":"wiki","rel":"relaciona"},{"alvo":"funil_cumulativo","confianca":0.5,"epistemico":"fato","origem":"wiki","rel":"relaciona"},{"alvo":"rust","confianca":0.5,"epistemico":"fato","origem":"wiki","rel":"relaciona"},{"alvo":"utilitarismo","confianca":0.5,"epistemico":"fato","origem":"wiki","rel":"relaciona"},{"alvo":"prova_isa","confianca":0.5,"epistemico":"fato","origem":"wiki","rel":"relaciona"},{"alvo":"fase_2_tecido_de_conhecimento","confianca":0.5,"epistemico":"fato","origem":"wiki","rel":"relaciona"}],"confianca":1.0,"contraexemplos":[],"descricao":"Secção de PROVA_ISA.md.","epistemico":"fato","exemplos":[],"id":"6_provas_de_composicao","memoria":"semantica","meta":{"dominio":"manual","nivel":6,"verificacao":"slug idêntico — enriquecer, não duplicar"},"origem":"PROVA_ISA.md","palavras_chave":[],"sinonimos":[],"tipo":"conceito","titulo":"6. Provas de composição"}
\section{6. Provas de composição}
Secção de PROVA\_ISA.md.
\prov{relações: parte\_de a\_prova\_da\_isa\_da\_broca; relaciona 8\_release\_core; relaciona topologia; relaciona aneis; relaciona historia; relaciona cache; relaciona broca\_so; relaciona gato\_operador\_hub; relaciona hopfield\_rede\_hub; relaciona kernel; relaciona complexos; relaciona recursao; relaciona algebra\_dos\_contratos; relaciona expressividade; relaciona memoria\_virtual; relaciona goldchain\_mercado\_hub; relaciona linker; relaciona resposta\_a\_pergunta\_dificil; relaciona etica; relaciona pow\_metal\_pre\_fecho; relaciona processos; relaciona java; relaciona parabola\_pa; relaciona regua\_associativa\_e\_o\_risco\_de\_vazamento; relaciona kappa; relaciona compilador; relaciona projeto\_cristaljson; relaciona funil\_cumulativo; relaciona rust; relaciona utilitarismo; relaciona prova\_isa; relaciona fase\_2\_tecido\_de\_conhecimento}
\prov{proveniência: PROVA\_ISA.md \textperiodcentered{} fato \textperiodcentered{} confiança 1.0 \textperiodcentered{} domínio manual \textperiodcentered{} história 35 versões no jornal}

%CRISTAL {"arestas":[{"alvo":"broca_os","confianca":0.95,"epistemico":"fato","origem":"paper","rel":"parte_de"}],"confianca":1.0,"contraexemplos":[],"descricao":"| Variável | Default | Uso | |----------|---------|-----| | `TIFFANY_HOME` | `~/.tiffany` ou `tiffany-home/` no repo | layout base/cliente/memoria | | `TIFFANY_ORQ` | — | `1` = orquestrador multi-idx no CLI | | `TIFFANY_CLIENTE` | — | `1` = incluir Base Cliente na consulta | | `TIFFANY_LIB` | auto | caminho para `libtiffany.so` |","epistemico":"fato","exemplos":[],"id":"6_variaveis_de_ambiente","memoria":"semantica","meta":{"dominio":"manual","nivel":6,"verificacao":"slug idêntico — enriquecer, não duplicar"},"origem":"INSTALL.md","palavras_chave":[],"sinonimos":[],"tipo":"conceito","titulo":"6. Variáveis de ambiente"}
\section{6. Variáveis de ambiente}
| Variável | Default | Uso | |----------|---------|-----| | `TIFFANY\_HOME` | `\~{}/.tiffany` ou `tiffany-home/` no repo | layout base/cliente/memoria | | `TIFFANY\_ORQ` | — | `1` = orquestrador multi-idx no CLI | | `TIFFANY\_CLIENTE` | — | `1` = incluir Base Cliente na consulta | | `TIFFANY\_LIB` | auto | caminho para `libtiffany.so` |
\prov{relações: parte\_de broca\_os}
\prov{proveniência: INSTALL.md \textperiodcentered{} fato \textperiodcentered{} confiança 1.0 \textperiodcentered{} domínio manual \textperiodcentered{} história 3 versões no jornal}

%CRISTAL {"arestas":[{"alvo":"broca_os","confianca":0.95,"epistemico":"fato","origem":"paper","rel":"parte_de"}],"confianca":1.0,"contraexemplos":[],"descricao":"```bash sh benchmarks/fogo.sh      # motor (colapso) sh benchmarks/fase2.sh       # assistente sh benchmarks/fase3.sh         # painel comercial completo sh test/test.sh              # cadeia do gato + fase 3 ```","epistemico":"fato","exemplos":[],"id":"7_benchmarks","memoria":"semantica","meta":{"dominio":"manual","nivel":6,"verificacao":"slug idêntico — enriquecer, não duplicar"},"origem":"INSTALL.md","palavras_chave":[],"sinonimos":[],"tipo":"conceito","titulo":"7. Benchmarks"}
\section{7. Benchmarks}
```bash sh benchmarks/fogo.sh      \# motor (colapso) sh benchmarks/fase2.sh       \# assistente sh benchmarks/fase3.sh         \# painel comercial completo sh test/test.sh              \# cadeia do gato + fase 3 ```
\prov{relações: parte\_de broca\_os}
\prov{proveniência: INSTALL.md \textperiodcentered{} fato \textperiodcentered{} confiança 1.0 \textperiodcentered{} domínio manual \textperiodcentered{} história 3 versões no jornal}

%CRISTAL {"arestas":[{"alvo":"a_prova_da_isa_da_broca","confianca":1.0,"epistemico":"fato","origem":"PROVA_ISA.md","rel":"parte_de"}],"confianca":1.0,"contraexemplos":[],"descricao":"| Limite | Razão | Contorno por composição | |--------|-------|-------------------------| | `LOAD` imediato fixo | endereço no bytecode | montador gera cadeia | | `JZ/JNZ` só `ZERO` | sem `JEQ` | `SUB` + CMP com zero | | Roundtrip word | neurônio ≠ identidade | scratch `≥0x8000` | | Índice dinâmico | sem `ADD` ao PC | tabela de saltos gerada | | `N` arbitrário em runtime | programa finito | montador conhece `mem/` |","epistemico":"fato","exemplos":[],"id":"7_limites_conhecidos","memoria":"semantica","meta":{"dominio":"manual","nivel":6,"verificacao":"slug idêntico — enriquecer, não duplicar"},"origem":"PROVA_ISA.md","palavras_chave":[],"sinonimos":[],"tipo":"conceito","titulo":"7. Limites conhecidos"}
\section{7. Limites conhecidos}
| Limite | Razão | Contorno por composição | |--------|-------|-------------------------| | `LOAD` imediato fixo | endereço no bytecode | montador gera cadeia | | `JZ/JNZ` só `ZERO` | sem `JEQ` | `SUB` + CMP com zero | | Roundtrip word | neurônio \ensuremath{\ne} identidade | scratch `\ensuremath{\ge}0x8000` | | Índice dinâmico | sem `ADD` ao PC | tabela de saltos gerada | | `N` arbitrário em runtime | programa finito | montador conhece `mem/` |
\prov{relações: parte\_de a\_prova\_da\_isa\_da\_broca}
\prov{proveniência: PROVA\_ISA.md \textperiodcentered{} fato \textperiodcentered{} confiança 1.0 \textperiodcentered{} domínio manual \textperiodcentered{} história 4 versões no jornal}

%CRISTAL {"arestas":[{"alvo":"broca_os","confianca":0.95,"epistemico":"fato","origem":"paper","rel":"parte_de"},{"alvo":"word","confianca":0.5,"epistemico":"fato","origem":"wiki","rel":"relaciona"},{"alvo":"virtualizacao","confianca":0.5,"epistemico":"fato","origem":"wiki","rel":"relaciona"}],"confianca":1.0,"contraexemplos":[],"descricao":"Programa `10_interpretador.asm` lê **3 words** pré-carregados em scratch `8100..8102`, cada um codificando `(op, arg)`:","epistemico":"fato","exemplos":[],"id":"7_micro-interpreter_limite_da_universalidade","memoria":"semantica","meta":{"dominio":"manual","nivel":6,"verificacao":"slug idêntico — enriquecer, não duplicar"},"origem":"EXPRESSIVIDADE.md","palavras_chave":[],"sinonimos":[],"tipo":"conceito","titulo":"7. Micro-interpreter (limite da universalidade)"}
\section{7. Micro-interpreter (limite da universalidade)}
Programa `10\_interpretador.asm` lê **3 words** pré-carregados em scratch `8100..8102`, cada um codificando `(op, arg)`:
\prov{relações: parte\_de broca\_os; relaciona word; relaciona virtualizacao}
\prov{proveniência: EXPRESSIVIDADE.md \textperiodcentered{} fato \textperiodcentered{} confiança 1.0 \textperiodcentered{} domínio manual \textperiodcentered{} história 5 versões no jornal}

%CRISTAL {"arestas":[{"alvo":"broca_os","confianca":0.95,"epistemico":"fato","origem":"paper","rel":"parte_de"},{"alvo":"9_criterio_de_rejeicao_checklist","confianca":0.5,"epistemico":"fato","origem":"wiki","rel":"relaciona"},{"alvo":"a_prova_da_isa_da_broca","confianca":0.5,"epistemico":"fato","origem":"wiki","rel":"relaciona"},{"alvo":"antropologia","confianca":0.5,"epistemico":"fato","origem":"wiki","rel":"relaciona"},{"alvo":"testes","confianca":0.5,"epistemico":"fato","origem":"wiki","rel":"relaciona"}],"confianca":1.0,"contraexemplos":[],"descricao":"Cada programa reporta via `make expressividade`:","epistemico":"fato","exemplos":[],"id":"8_complexidade_medida_empirica","memoria":"semantica","meta":{"dominio":"manual","nivel":6,"verificacao":"slug idêntico — enriquecer, não duplicar"},"origem":"EXPRESSIVIDADE.md","palavras_chave":[],"sinonimos":[],"tipo":"conceito","titulo":"8. Complexidade (medida empírica)"}
\section{8. Complexidade (medida empírica)}
Cada programa reporta via `make expressividade`:
\prov{relações: parte\_de broca\_os; relaciona 9\_criterio\_de\_rejeicao\_checklist; relaciona a\_prova\_da\_isa\_da\_broca; relaciona antropologia; relaciona testes}
\prov{proveniência: EXPRESSIVIDADE.md \textperiodcentered{} fato \textperiodcentered{} confiança 1.0 \textperiodcentered{} domínio manual \textperiodcentered{} história 7 versões no jornal}

%CRISTAL {"arestas":[{"alvo":"a_prova_da_isa_da_broca","confianca":1.0,"epistemico":"fato","origem":"PROVA_ISA.md","rel":"parte_de"}],"confianca":1.0,"contraexemplos":[],"descricao":"| Opcode | Por que existe | Removível? | |--------|----------------|------------| | `HALT` | terminação | não | | `LOAD` | única ponte neurônio→CPU | não | | `STORE` | persistir `R` / scratch | não (acumulador) | | `ADD` | combinar espectros | SUB não substitui | | `SUB` | diferença / teste igualdade | não | | `AND/OR/XOR` | máscaras bitwise | teórico; não usados nos 7 prog. | | `CMP` | condições | JZ/JNZ precisam FLAGS |","epistemico":"fato","exemplos":[],"id":"8_justificativa_formal_por_opcode","memoria":"semantica","meta":{"dominio":"manual","nivel":6,"verificacao":"slug idêntico — enriquecer, não duplicar"},"origem":"PROVA_ISA.md","palavras_chave":[],"sinonimos":[],"tipo":"conceito","titulo":"8. Justificativa formal por opcode"}
\section{8. Justificativa formal por opcode}
| Opcode | Por que existe | Removível? | |--------|----------------|------------| | `HALT` | terminação | não | | `LOAD` | única ponte neurônio\ensuremath{\to}CPU | não | | `STORE` | persistir `R` / scratch | não (acumulador) | | `ADD` | combinar espectros | SUB não substitui | | `SUB` | diferença / teste igualdade | não | | `AND/OR/XOR` | máscaras bitwise | teórico; não usados nos 7 prog. | | `CMP` | condições | JZ/JNZ precisam FLAGS |
\prov{relações: parte\_de a\_prova\_da\_isa\_da\_broca}
\prov{proveniência: PROVA\_ISA.md \textperiodcentered{} fato \textperiodcentered{} confiança 1.0 \textperiodcentered{} domínio manual \textperiodcentered{} história 4 versões no jornal}

%CRISTAL {"arestas":[{"alvo":"broca_os","confianca":0.95,"epistemico":"fato","origem":"paper","rel":"parte_de"},{"alvo":"rust","confianca":0.5,"epistemico":"fato","origem":"wiki","rel":"relaciona"},{"alvo":"java","confianca":0.5,"epistemico":"fato","origem":"wiki","rel":"relaciona"},{"alvo":"gpu","confianca":0.5,"epistemico":"fato","origem":"wiki","rel":"relaciona"},{"alvo":"broca_so","confianca":0.5,"epistemico":"fato","origem":"wiki","rel":"relaciona"},{"alvo":"kernel","confianca":0.5,"epistemico":"fato","origem":"wiki","rel":"relaciona"},{"alvo":"koch_atlas","confianca":0.5,"epistemico":"fato","origem":"wiki","rel":"relaciona"},{"alvo":"goldcoin","confianca":0.5,"epistemico":"fato","origem":"wiki","rel":"relaciona"},{"alvo":"mamifero","confianca":0.5,"epistemico":"fato","origem":"wiki","rel":"relaciona"},{"alvo":"recursao","confianca":0.5,"epistemico":"fato","origem":"wiki","rel":"relaciona"},{"alvo":"gato_operador_hub","confianca":0.5,"epistemico":"fato","origem":"wiki","rel":"relaciona"},{"alvo":"grava_fato_resposta_normal_em_seguida","confianca":0.5,"epistemico":"fato","origem":"wiki","rel":"relaciona"},{"alvo":"prova_isa","confianca":0.5,"epistemico":"fato","origem":"wiki","rel":"relaciona"},{"alvo":"numa","confianca":0.5,"epistemico":"fato","origem":"wiki","rel":"relaciona"},{"alvo":"python","confianca":0.5,"epistemico":"fato","origem":"wiki","rel":"relaciona"},{"alvo":"thread","confianca":0.5,"epistemico":"fato","origem":"wiki","rel":"relaciona"},{"alvo":"metais","confianca":0.5,"epistemico":"fato","origem":"wiki","rel":"relaciona"},{"alvo":"driver","confianca":0.5,"epistemico":"fato","origem":"wiki","rel":"relaciona"},{"alvo":"resposta_a_pergunta_dificil","confianca":0.5,"epistemico":"fato","origem":"wiki","rel":"relaciona"},{"alvo":"vida-morte","confianca":0.5,"epistemico":"fato","origem":"wiki","rel":"relaciona"},{"alvo":"koch_memoria","confianca":0.5,"epistemico":"fato","origem":"wiki","rel":"relaciona"},{"alvo":"install","confianca":0.5,"epistemico":"fato","origem":"wiki","rel":"relaciona"},{"alvo":"main","confianca":0.5,"epistemico":"fato","origem":"wiki","rel":"relaciona"},{"alvo":"assistente","confianca":0.5,"epistemico":"fato","origem":"wiki","rel":"relaciona"},{"alvo":"broca_os_hub","confianca":0.5,"epistemico":"fato","origem":"wiki","rel":"relaciona"},{"alvo":"hopfield_rede_hub","confianca":0.5,"epistemico":"fato","origem":"wiki","rel":"relaciona"},{"alvo":"go","confianca":0.5,"epistemico":"fato","origem":"wiki","rel":"relaciona"},{"alvo":"montagem","confianca":0.5,"epistemico":"fato","origem":"wiki","rel":"relaciona"},{"alvo":"linker","confianca":0.5,"epistemico":"fato","origem":"wiki","rel":"relaciona"},{"alvo":"topologia","confianca":0.5,"epistemico":"fato","origem":"wiki","rel":"relaciona"},{"alvo":"memoria_virtual","confianca":0.5,"epistemico":"fato","origem":"wiki","rel":"relaciona"},{"alvo":"teste_ordem","confianca":0.5,"epistemico":"fato","origem":"wiki","rel":"relaciona"},{"alvo":"geometria_hub_broca","confianca":0.5,"epistemico":"fato","origem":"wiki","rel":"relaciona"},{"alvo":"cifra_prova","confianca":0.5,"epistemico":"fato","origem":"wiki","rel":"relaciona"}],"confianca":1.0,"contraexemplos":[],"descricao":"```bash sh scripts/release-core.sh","epistemico":"fato","exemplos":[],"id":"8_release_core","memoria":"semantica","meta":{"dominio":"manual","nivel":6,"verificacao":"slug idêntico — enriquecer, não duplicar"},"origem":"INSTALL.md","palavras_chave":[],"sinonimos":[],"tipo":"conceito","titulo":"8. Release Core"}
\section{8. Release Core}
```bash sh scripts/release-core.sh
\prov{relações: parte\_de broca\_os; relaciona rust; relaciona java; relaciona gpu; relaciona broca\_so; relaciona kernel; relaciona koch\_atlas; relaciona goldcoin; relaciona mamifero; relaciona recursao; relaciona gato\_operador\_hub; relaciona grava\_fato\_resposta\_normal\_em\_seguida; relaciona prova\_isa; relaciona numa; relaciona python; relaciona thread; relaciona metais; relaciona driver; relaciona resposta\_a\_pergunta\_dificil; relaciona vida-morte; relaciona koch\_memoria; relaciona install; relaciona main; relaciona assistente; relaciona broca\_os\_hub; relaciona hopfield\_rede\_hub; relaciona go; relaciona montagem; relaciona linker; relaciona topologia; relaciona memoria\_virtual; relaciona teste\_ordem; relaciona geometria\_hub\_broca; relaciona cifra\_prova}
\prov{proveniência: INSTALL.md \textperiodcentered{} fato \textperiodcentered{} confiança 1.0 \textperiodcentered{} domínio manual \textperiodcentered{} história 36 versões no jornal}

%CRISTAL {"arestas":[{"alvo":"a_prova_da_isa_da_broca","confianca":1.0,"epistemico":"fato","origem":"PROVA_ISA.md","rel":"parte_de"},{"alvo":"historia","confianca":0.5,"epistemico":"fato","origem":"wiki","rel":"relaciona"},{"alvo":"regua_associativa_e_o_risco_de_vazamento","confianca":0.5,"epistemico":"fato","origem":"wiki","rel":"relaciona"},{"alvo":"expressividade","confianca":0.5,"epistemico":"fato","origem":"wiki","rel":"relaciona"},{"alvo":"word","confianca":0.5,"epistemico":"fato","origem":"wiki","rel":"relaciona"},{"alvo":"vontade_de_poder","confianca":0.5,"epistemico":"fato","origem":"wiki","rel":"relaciona"}],"confianca":1.0,"contraexemplos":[],"descricao":"Antes de propor opcode `X`:","epistemico":"fato","exemplos":[],"id":"9_criterio_de_rejeicao_checklist","memoria":"semantica","meta":{"dominio":"manual","nivel":6,"verificacao":"slug idêntico — enriquecer, não duplicar"},"origem":"PROVA_ISA.md","palavras_chave":[],"sinonimos":[],"tipo":"conceito","titulo":"9. Critério de rejeição (checklist)"}
\section{9. Critério de rejeição (checklist)}
Antes de propor opcode `X`:
\prov{relações: parte\_de a\_prova\_da\_isa\_da\_broca; relaciona historia; relaciona regua\_associativa\_e\_o\_risco\_de\_vazamento; relaciona expressividade; relaciona word; relaciona vontade\_de\_poder}
\prov{proveniência: PROVA\_ISA.md \textperiodcentered{} fato \textperiodcentered{} confiança 1.0 \textperiodcentered{} domínio manual \textperiodcentered{} história 9 versões no jornal}

%CRISTAL {"arestas":[{"alvo":"sistema_e_lei_nao_dono","confianca":1.0,"epistemico":"fato","origem":"wiki","rel":"usa"},{"alvo":"soberania_da_chave","confianca":1.0,"epistemico":"fato","origem":"wiki","rel":"usa"},{"alvo":"o_substrato_a_maquina_e_funcao_de_onda","confianca":1.0,"epistemico":"fato","origem":"wiki","rel":"relaciona"}],"confianca":1.0,"contraexemplos":[],"descricao":"A GoldChain é descentralizada/soberana/offline-first: a chave é o endereço (sha256 da chave), não há servidor nem lugar; sem tempo, acessar é idempotente e a mesma seed produz a mesma banda em qualquer máquina (verificado: banda idêntica no laptop e no GEX44). Painel = colapso na leitura (sem olho, o sistema dorme; série de colapsos somada por Neumann (I−rU)⁻¹).","epistemico":"fato","exemplos":[],"id":"a_bolha_infraestrutura","memoria":"semantica","meta":{"autor":"Aarão/broca-so","dominio":"broca_repos","fonte":""},"origem":"wiki","palavras_chave":[],"sinonimos":[],"tipo":"conceito","titulo":"A Bolha (infraestrutura sem 'onde' nem tempo)"}
\section{A Bolha (infraestrutura sem 'onde' nem tempo)}
A GoldChain é descentralizada/soberana/offline-first: a chave é o endereço (sha256 da chave), não há servidor nem lugar; sem tempo, acessar é idempotente e a mesma seed produz a mesma banda em qualquer máquina (verificado: banda idêntica no laptop e no GEX44). Painel = colapso na leitura (sem olho, o sistema dorme; série de colapsos somada por Neumann (I\ensuremath{-}rU)\ensuremath{^{-1}}).
\prov{relações: usa sistema\_e\_lei\_nao\_dono; usa soberania\_da\_chave; relaciona o\_substrato\_a\_maquina\_e\_funcao\_de\_onda}
\prov{proveniência: wiki \textperiodcentered{} fato \textperiodcentered{} confiança 1.0 \textperiodcentered{} domínio broca\_repos \textperiodcentered{} história 4 versões no jornal}

%CRISTAL {"arestas":[{"alvo":"parabola_isa","confianca":1.0,"epistemico":"fato","origem":"manual","rel":"parte_de"},{"alvo":"utilitarismo","confianca":0.5,"epistemico":"fato","origem":"wiki","rel":"relaciona"},{"alvo":"resultado_completo","confianca":0.5,"epistemico":"fato","origem":"wiki","rel":"relaciona"},{"alvo":"teoria_do_valor","confianca":0.5,"epistemico":"fato","origem":"wiki","rel":"relaciona"},{"alvo":"o_que_e_no_projecto","confianca":0.5,"epistemico":"fato","origem":"wiki","rel":"relaciona"},{"alvo":"destruir_a_parabola","confianca":0.5,"epistemico":"fato","origem":"wiki","rel":"relaciona"},{"alvo":"cifra_ouro_branco_nucleo_espectral","confianca":0.5,"epistemico":"fato","origem":"wiki","rel":"relaciona"},{"alvo":"gentil_16_soma_dos_termos_de_uma_pam","confianca":0.5,"epistemico":"fato","origem":"wiki","rel":"relaciona"},{"alvo":"ficheiros","confianca":0.5,"epistemico":"fato","origem":"wiki","rel":"relaciona"},{"alvo":"word","confianca":0.5,"epistemico":"fato","origem":"wiki","rel":"relaciona"},{"alvo":"modo_espectral_spect","confianca":0.5,"epistemico":"fato","origem":"wiki","rel":"relaciona"},{"alvo":"toryces_e_helyces","confianca":0.5,"epistemico":"fato","origem":"wiki","rel":"relaciona"},{"alvo":"broca-so_cristalizado","confianca":0.5,"epistemico":"fato","origem":"wiki","rel":"relaciona"},{"alvo":"casl-propagation","confianca":0.5,"epistemico":"fato","origem":"wiki","rel":"relaciona"}],"confianca":1.0,"contraexemplos":[],"descricao":"> *a + c² = 1* — painel da máquina, não restrição.","epistemico":"fato","exemplos":[],"id":"a_cacada_da_parabola_mapa_vivo_da_isa","memoria":"semantica","meta":{"dominio":"manual","nivel":6,"verificacao":"slug idêntico — enriquecer, não duplicar"},"origem":"PARABOLA_ISA.md","palavras_chave":[],"sinonimos":[],"tipo":"conceito","titulo":"🌞 A Caçada da Parábola — Mapa Vivo da ISA"}
\section{[sol] A Caçada da Parábola — Mapa Vivo da ISA}
> *a + c\ensuremath{^{2}} = 1* — painel da máquina, não restrição.
\prov{relações: parte\_de parabola\_isa; relaciona utilitarismo; relaciona resultado\_completo; relaciona teoria\_do\_valor; relaciona o\_que\_e\_no\_projecto; relaciona destruir\_a\_parabola; relaciona cifra\_ouro\_branco\_nucleo\_espectral; relaciona gentil\_16\_soma\_dos\_termos\_de\_uma\_pam; relaciona ficheiros; relaciona word; relaciona modo\_espectral\_spect; relaciona toryces\_e\_helyces; relaciona broca-so\_cristalizado; relaciona casl-propagation}
\prov{proveniência: PARABOLA\_ISA.md \textperiodcentered{} fato \textperiodcentered{} confiança 1.0 \textperiodcentered{} domínio manual \textperiodcentered{} história 18 versões no jornal}

%CRISTAL {"arestas":[{"alvo":"corpo_glacial","confianca":1.0,"epistemico":"fato","origem":"paper","rel":"parte_de"},{"alvo":"olho_de_koch","confianca":1.0,"epistemico":"fato","origem":"paper","rel":"referencia"}],"confianca":1.0,"contraexemplos":[],"descricao":"| Métrica | Valor | |---------|-------| | Médio | 0.604890 | | Máximo | 0.999095 (nonce 1,966,079) | | Fecho ok — acopl médio | 0.998839 | | Sem fecho — acopl médio | 0.604889 | | SHA ok — acopl médio | 0.596899 | | Sem SHA — acopl médio | 0.604890 | | Separação fecho | +0.393950 | | Separação SHA | -0.007991 | | Fechos | 2 | | SHA ok | 1 |","epistemico":"fato","exemplos":[],"id":"a_dinamica_e_o_acoplamento","memoria":"semantica","meta":{"dominio":"manual","duplicata_sugerida":[],"nivel":6,"verificacao":"parte de conceito mais amplo 'a_dinamica_e_o_acoplamento'"},"origem":"POW_METAL_RESSONANCIA.md","palavras_chave":[],"sinonimos":[],"tipo":"conceito","titulo":"A dinâmica é o acoplamento"}
\section{A dinâmica é o acoplamento}
| Métrica | Valor | |---------|-------| | Médio | 0.604890 | | Máximo | 0.999095 (nonce 1,966,079) | | Fecho ok — acopl médio | 0.998839 | | Sem fecho — acopl médio | 0.604889 | | SHA ok — acopl médio | 0.596899 | | Sem SHA — acopl médio | 0.604890 | | Separação fecho | +0.393950 | | Separação SHA | -0.007991 | | Fechos | 2 | | SHA ok | 1 |
\prov{relações: parte\_de corpo\_glacial; referencia olho\_de\_koch}
\prov{proveniência: POW\_METAL\_RESSONANCIA.md \textperiodcentered{} fato \textperiodcentered{} confiança 1.0 \textperiodcentered{} domínio manual \textperiodcentered{} história 41 versões no jornal}

%CRISTAL {"arestas":[{"alvo":"algebra_dos_contratos","confianca":1.0,"epistemico":"fato","origem":"paper","rel":"parte_de"},{"alvo":"goldchain","confianca":1.0,"epistemico":"fato","origem":"paper","rel":"parte_de"}],"confianca":1.0,"contraexemplos":[],"descricao":"- **c²** medido via Bhattacharyya (ρ²) — caminho A - **calor (a)** medido via Hellinger — caminho B, independente - **Sem** renormalizar a ou c² depois - Estados degenerados (total+e ≤ 0): **não** inventamos 0.5/0.5 — contamos como não-testável","epistemico":"fato","exemplos":[],"id":"a_economia_do_pell_---_a_moeda_do_protocolo","memoria":"semantica","meta":{"dominio":"manual","duplicata_sugerida":[],"nivel":6,"verificacao":"parte de conceito mais amplo 'a_economia_do_pell_---_a_moeda_do_protocolo'"},"origem":"PARABOLA_DESTRUIR.md","palavras_chave":[],"sinonimos":[],"tipo":"conceito","titulo":"A economia do Pell --- a moeda do protocolo"}
\section{A economia do Pell --- a moeda do protocolo}
- **c\ensuremath{^{2}}** medido via Bhattacharyya (\ensuremath{\rho}\ensuremath{^{2}}) — caminho A - **calor (a)** medido via Hellinger — caminho B, independente - **Sem** renormalizar a ou c\ensuremath{^{2}} depois - Estados degenerados (total+e \ensuremath{\le} 0): **não** inventamos 0.5/0.5 — contamos como não-testável
\prov{relações: parte\_de algebra\_dos\_contratos; parte\_de goldchain}
\prov{proveniência: PARABOLA\_DESTRUIR.md \textperiodcentered{} fato \textperiodcentered{} confiança 1.0 \textperiodcentered{} domínio manual \textperiodcentered{} história 12 versões no jornal}

%CRISTAL {"arestas":[{"alvo":"prova_isa","confianca":1.0,"epistemico":"fato","origem":"manual","rel":"parte_de"},{"alvo":"topologia","confianca":0.5,"epistemico":"fato","origem":"wiki","rel":"relaciona"},{"alvo":"antropologia","confianca":0.5,"epistemico":"fato","origem":"wiki","rel":"relaciona"},{"alvo":"utilitarismo","confianca":0.5,"epistemico":"fato","origem":"wiki","rel":"relaciona"},{"alvo":"expressividade","confianca":0.5,"epistemico":"fato","origem":"wiki","rel":"relaciona"},{"alvo":"gato_operador_hub","confianca":0.5,"epistemico":"fato","origem":"wiki","rel":"relaciona"},{"alvo":"fixamos_e_o_curinga","confianca":0.5,"epistemico":"fato","origem":"wiki","rel":"relaciona"},{"alvo":"imagem_do_floco_---_validacao_no","confianca":0.5,"epistemico":"fato","origem":"wiki","rel":"relaciona"}],"confianca":1.0,"contraexemplos":[],"descricao":"Documento formal. **Zero opcodes novos** — só composição.","epistemico":"fato","exemplos":[],"id":"a_prova_da_isa_da_broca","memoria":"semantica","meta":{"dominio":"manual","nivel":6,"verificacao":"slug idêntico — enriquecer, não duplicar"},"origem":"PROVA_ISA.md","palavras_chave":[],"sinonimos":[],"tipo":"conceito","titulo":"A Prova da ISA da Broca"}
\section{A Prova da ISA da Broca}
Documento formal. **Zero opcodes novos** — só composição.
\prov{relações: parte\_de prova\_isa; relaciona topologia; relaciona antropologia; relaciona utilitarismo; relaciona expressividade; relaciona gato\_operador\_hub; relaciona fixamos\_e\_o\_curinga; relaciona imagem\_do\_floco\_---\_validacao\_no}
\prov{proveniência: PROVA\_ISA.md \textperiodcentered{} fato \textperiodcentered{} confiança 1.0 \textperiodcentered{} domínio manual \textperiodcentered{} história 11 versões no jornal}

%CRISTAL {"arestas":[{"alvo":"respiracao_dos_programas","confianca":1.0,"epistemico":"fato","origem":"PARABOLA_ISA.md","rel":"parte_de"},{"alvo":"tesouraria","confianca":0.5,"epistemico":"fato","origem":"wiki","rel":"relaciona"},{"alvo":"word","confianca":0.5,"epistemico":"fato","origem":"wiki","rel":"relaciona"}],"confianca":1.0,"contraexemplos":[],"descricao":"- calor acumulado: **0.651** - fase acumulada: **13.629** - perda: **0.000** - potência (calor/passo): **0.046** - passos: **14**","epistemico":"fato","exemplos":[],"id":"acumulador","memoria":"semantica","meta":{"dominio":"manual","nivel":6,"verificacao":"slug idêntico — enriquecer, não duplicar"},"origem":"PARABOLA_ISA.md","palavras_chave":[],"sinonimos":[],"tipo":"conceito","titulo":"🌱 acumulador"}
\section{[semente] acumulador}
- calor acumulado: **0.651** - fase acumulada: **13.629** - perda: **0.000** - potência (calor/passo): **0.046** - passos: **14**
\prov{relações: parte\_de respiracao\_dos\_programas; relaciona tesouraria; relaciona word}
\prov{proveniência: PARABOLA\_ISA.md \textperiodcentered{} fato \textperiodcentered{} confiança 1.0 \textperiodcentered{} domínio manual \textperiodcentered{} história 8 versões no jornal}

%CRISTAL {"arestas":[{"alvo":"prova_da_maquina_espectral","confianca":1.0,"epistemico":"fato","origem":"CIFRA_PROVA.md","rel":"parte_de"},{"alvo":"virtualizacao","confianca":0.5,"epistemico":"fato","origem":"wiki","rel":"relaciona"},{"alvo":"word","confianca":0.5,"epistemico":"fato","origem":"wiki","rel":"relaciona"},{"alvo":"vida-morte","confianca":0.5,"epistemico":"fato","origem":"wiki","rel":"relaciona"}],"confianca":1.0,"contraexemplos":[],"descricao":"`11_acum_pell.asm` substitui a cadeia `LOAD+ADD` de `04_acumulador.asm` por um loop `SILVER`+`JNZ`:","epistemico":"fato","exemplos":[],"id":"acumulador_espectral_loop_vs_desenrolado","memoria":"semantica","meta":{"dominio":"manual","nivel":6,"verificacao":"slug idêntico — enriquecer, não duplicar"},"origem":"CIFRA_PROVA.md","palavras_chave":[],"sinonimos":[],"tipo":"conceito","titulo":"Acumulador espectral — loop vs desenrolado"}
\section{Acumulador espectral — loop vs desenrolado}
`11\_acum\_pell.asm` substitui a cadeia `LOAD+ADD` de `04\_acumulador.asm` por um loop `SILVER`+`JNZ`:
\prov{relações: parte\_de prova\_da\_maquina\_espectral; relaciona virtualizacao; relaciona word; relaciona vida-morte}
\prov{proveniência: CIFRA\_PROVA.md \textperiodcentered{} fato \textperiodcentered{} confiança 1.0 \textperiodcentered{} domínio manual \textperiodcentered{} história 7 versões no jornal}

%CRISTAL {"arestas":[{"alvo":"reta_de_ouro","confianca":1.0,"epistemico":"fato","origem":"wiki","rel":"usa"},{"alvo":"regua_de_gentil","confianca":1.0,"epistemico":"fato","origem":"wiki","rel":"relaciona"},{"alvo":"matematica_estelar","confianca":1.0,"epistemico":"fato","origem":"wiki","rel":"parte_de"},{"alvo":"transformada_metalica_fechada_perelman","confianca":1.0,"epistemico":"fato","origem":"wiki","rel":"relaciona"}],"confianca":1.0,"contraexemplos":[],"descricao":"Operações pontuais = deltas na contração. A agulha de Dirac (banda ∞) vaza acima do gap e invade o invariante — proibida; a agulha de Perelman (band-limited, reconstruída só pelos modos de ouro Ω=φ¹/φ até o gap) — admitida. A régua decide o que pode ser medido.","epistemico":"fato","exemplos":[],"id":"agulha_admissao","memoria":"semantica","meta":{"autor":"Lopes/Aarão","dominio":"hiper","fonte":""},"origem":"wiki","palavras_chave":[],"sinonimos":[],"tipo":"conceito","titulo":"Agulhas e a Lei de Admissão (Perelman vs Dirac)"}
\section{Agulhas e a Lei de Admissão (Perelman vs Dirac)}
Operações pontuais = deltas na contração. A agulha de Dirac (banda \ensuremath{\infty}) vaza acima do gap e invade o invariante — proibida; a agulha de Perelman (band-limited, reconstruída só pelos modos de ouro \ensuremath{\Omega}=\ensuremath{\varphi}\ensuremath{^{1}}/\ensuremath{\varphi} até o gap) — admitida. A régua decide o que pode ser medido.
\prov{relações: usa reta\_de\_ouro; relaciona regua\_de\_gentil; parte\_de matematica\_estelar; relaciona transformada\_metalica\_fechada\_perelman}
\prov{proveniência: wiki \textperiodcentered{} fato \textperiodcentered{} confiança 1.0 \textperiodcentered{} domínio hiper \textperiodcentered{} história 7 versões no jornal}

%CRISTAL {"arestas":[{"alvo":"matematica_estelar","confianca":1.0,"epistemico":"fato","origem":"wiki","rel":"parte_de"},{"alvo":"algebras_hipercomplexas","confianca":1.0,"epistemico":"fato","origem":"wiki","rel":"usa"}],"confianca":1.0,"contraexemplos":[],"descricao":"Álgebra real 2D gerada por z=a+c·e sob o invariante a+c²=1, com e²=1−s e produto ·ₛ. Em s=0 é split-complex plana/associativa; em s=1, octoniônica não-associativa máxima. É a mesma família ×ₛ vista relativisticamente.","epistemico":"fato","exemplos":[],"id":"algebra_estelar","memoria":"semantica","meta":{"autor":"Lopes/Aarão","dominio":"hiper","fonte":""},"origem":"wiki","palavras_chave":[],"sinonimos":[],"tipo":"conceito","titulo":"Álgebra Estelar 𝓔ₛ"}
\section{Álgebra Estelar \ensuremath{\mathcal{E}}\ensuremath{_{s}}}
Álgebra real 2D gerada por z=a+c\textperiodcentered e sob o invariante a+c\ensuremath{^{2}}=1, com e\ensuremath{^{2}}=1\ensuremath{-}s e produto \textperiodcentered \ensuremath{_{s}}. Em s=0 é split-complex plana/associativa; em s=1, octoniônica não-associativa máxima. É a mesma família $\times$\ensuremath{_{s}} vista relativisticamente.
\prov{relações: parte\_de matematica\_estelar; usa algebras\_hipercomplexas}
\prov{proveniência: wiki \textperiodcentered{} fato \textperiodcentered{} confiança 1.0 \textperiodcentered{} domínio hiper \textperiodcentered{} história 4 versões no jornal}

%CRISTAL {"arestas":[{"alvo":"sistema_operacional","confianca":1.0,"epistemico":"fato","origem":"curriculo","rel":"especializa"},{"alvo":"estrutura_de_dados","confianca":1.0,"epistemico":"fato","origem":"curriculo","rel":"depende"},{"alvo":"maquina_de_turing","confianca":1.0,"epistemico":"fato","origem":"curriculo","rel":"depende"}],"confianca":1.0,"contraexemplos":[],"descricao":"| Algoritmo | Ficheiro | Status | |-----------|----------|--------| | Soma | `expressividade/01_soma.asm` | ✅ | | Busca linear | `expressividade/02_busca.asm` | ✅ | | Contagem | `expressividade/03_contagem.asm` | ✅ | | Comparação | `expressividade/04_comparacao.asm` | ✅ | | Copiar memória | `expressividade/05_copiar.asm` | ✅ | | Reduce | `expressividade/06_reduce.asm` | ✅ | | Map (×2) | `expressividade/07_map.asm` | ✅ |","epistemico":"fato","exemplos":[],"id":"algoritmo","memoria":"semantica","meta":{"dominio":"manual","duplicata_sugerida":[],"nivel":6,"verificacao":"parte de conceito mais amplo 'algoritmo'"},"origem":"EXPRESSIVIDADE.md","palavras_chave":["passos","correção","complexidade"],"sinonimos":[],"tipo":"conceito","titulo":"algoritmo"}
\section{algoritmo}
| Algoritmo | Ficheiro | Status | |-----------|----------|--------| | Soma | `expressividade/01\_soma.asm` | [ok] | | Busca linear | `expressividade/02\_busca.asm` | [ok] | | Contagem | `expressividade/03\_contagem.asm` | [ok] | | Comparação | `expressividade/04\_comparacao.asm` | [ok] | | Copiar memória | `expressividade/05\_copiar.asm` | [ok] | | Reduce | `expressividade/06\_reduce.asm` | [ok] | | Map ($\times$2) | `expressividade/07\_map.asm` | [ok] |
\prov{relações: especializa sistema\_operacional; depende estrutura\_de\_dados; depende maquina\_de\_turing}
\prov{proveniência: EXPRESSIVIDADE.md \textperiodcentered{} fato \textperiodcentered{} confiança 1.0 \textperiodcentered{} domínio manual \textperiodcentered{} história 19 versões no jornal}

%CRISTAL {"arestas":[{"alvo":"cifra","confianca":1.0,"epistemico":"fato","origem":"wiki","rel":"explica"},{"alvo":"colapso_koch_tiffany","confianca":1.0,"epistemico":"fato","origem":"wiki","rel":"relaciona"}],"confianca":1.0,"contraexemplos":[],"descricao":"A 'alma' (resíduo não-associativo do associador sobre o bit bruto) é um hash com perda: avalanche≈0, muitos-para-um, ~6 bits. É irreversível por PERDA (não por custo como o SHA-256): as pré-imagens existem e são vastas, mas a informação distintiva não está na imagem. Publicar μ_x não vaza x (sela a semente); a não-forja vem da camada Ed25519.","epistemico":"fato","exemplos":[],"id":"alma_hash_com_perda","memoria":"semantica","meta":{"autor":"Aarão/broca-so","dominio":"broca_repos","fonte":""},"origem":"wiki","palavras_chave":[],"sinonimos":[],"tipo":"conceito","titulo":"A Alma é um Hash com Perda"}
\section{A Alma é um Hash com Perda}
A 'alma' (resíduo não-associativo do associador sobre o bit bruto) é um hash com perda: avalanche\ensuremath{\approx}0, muitos-para-um, \~{}6 bits. É irreversível por PERDA (não por custo como o SHA-256): as pré-imagens existem e são vastas, mas a informação distintiva não está na imagem. Publicar \ensuremath{\mu}\_x não vaza x (sela a semente); a não-forja vem da camada Ed25519.
\prov{relações: explica cifra; relaciona colapso\_koch\_tiffany}
\prov{proveniência: wiki \textperiodcentered{} fato \textperiodcentered{} confiança 1.0 \textperiodcentered{} domínio broca\_repos \textperiodcentered{} história 3 versões no jornal}

%CRISTAL {"arestas":[{"alvo":"substrato_prtsuv_h","confianca":1.0,"epistemico":"fato","origem":"wiki","rel":"usa"},{"alvo":"torre_hurwitz","confianca":1.0,"epistemico":"fato","origem":"wiki","rel":"usa"},{"alvo":"cordas_ilha_hurwitz","confianca":1.0,"epistemico":"fato","origem":"wiki","rel":"explica"}],"confianca":1.0,"contraexemplos":[],"descricao":"Teorema: na família ×ₛ, o tri-associador espinorial anula-se para todo ψ,χ se e só se s²=1; combinado com Adams-Eckmann (n−1∈0,1,3,7), dá as dimensões críticas de SUSY D=n+2∈3,4,6,10 — sem invocar matrizes de Dirac.","epistemico":"fato","exemplos":[],"id":"alternatividade_dimensoes_criticas","memoria":"semantica","meta":{"autor":"Lopes/Aarão","dominio":"hiper","fonte":""},"origem":"wiki","palavras_chave":[],"sinonimos":[],"tipo":"conceito","titulo":"Alternatividade ⟺ s²=1 ⟹ D∈3,4,6,10"}
\section{Alternatividade \ensuremath{\Longleftrightarrow} s\ensuremath{^{2}}=1 \ensuremath{\Longrightarrow} D\ensuremath{\in}3,4,6,10}
Teorema: na família $\times$\ensuremath{_{s}}, o tri-associador espinorial anula-se para todo \ensuremath{\psi},\ensuremath{\chi} se e só se s\ensuremath{^{2}}=1; combinado com Adams-Eckmann (n\ensuremath{-}1\ensuremath{\in}0,1,3,7), dá as dimensões críticas de SUSY D=n+2\ensuremath{\in}3,4,6,10 — sem invocar matrizes de Dirac.
\prov{relações: usa substrato\_prtsuv\_h; usa torre\_hurwitz; explica cordas\_ilha\_hurwitz}
\prov{proveniência: wiki \textperiodcentered{} fato \textperiodcentered{} confiança 1.0 \textperiodcentered{} domínio hiper \textperiodcentered{} história 7 versões no jornal}

%CRISTAL {"arestas":[{"alvo":"broca-so_cristalizado","confianca":1.0,"epistemico":"fato","origem":"BROCA_CRISTAL.md","rel":"parte_de"}],"confianca":1.0,"contraexemplos":[],"descricao":"| | Amputação errada ❌ | Projeção Koch ✅ | |--|---------------------|-------------------| | O quê | `Morton`, FFT no espectro gato, cadáver no pipeline | Olho de Koch → atlas → ecrã | | Efeito | mata `e` (racha: 0/587) | projecta associativo **depois** do vivo | | Onde | **nunca no núcleo** | **só na saída** para borda | | Clássico | perde informação | **recebe** o que Koch projectou |","epistemico":"fato","exemplos":[],"id":"amputacao_errada_vs_projecao_certa","memoria":"semantica","meta":{"dominio":"manual","nivel":6,"verificacao":"slug idêntico — enriquecer, não duplicar"},"origem":"BROCA_CRISTAL.md","palavras_chave":[],"sinonimos":[],"tipo":"conceito","titulo":"Amputação errada vs projeção certa"}
\section{Amputação errada vs projeção certa}
| | Amputação errada [x] | Projeção Koch [ok] | |--|---------------------|-------------------| | O quê | `Morton`, FFT no espectro gato, cadáver no pipeline | Olho de Koch \ensuremath{\to} atlas \ensuremath{\to} ecrã | | Efeito | mata `e` (racha: 0/587) | projecta associativo **depois** do vivo | | Onde | **nunca no núcleo** | **só na saída** para borda | | Clássico | perde informação | **recebe** o que Koch projectou |
\prov{relações: parte\_de broca-so\_cristalizado}
\prov{proveniência: BROCA\_CRISTAL.md \textperiodcentered{} fato \textperiodcentered{} confiança 1.0 \textperiodcentered{} domínio manual \textperiodcentered{} história 5 versões no jornal}

%CRISTAL {"arestas":[{"alvo":"seguranca_dimensional","confianca":1.0,"epistemico":"fato","origem":"wiki","rel":"usa"},{"alvo":"circuito_como_torcao","confianca":1.0,"epistemico":"fato","origem":"wiki","rel":"relaciona"}],"confianca":1.0,"contraexemplos":[],"descricao":"Os objetos do painel instanciam a teoria clássica de antenas (Maxwell): gap/LFSR = corrente; idempotente e_i = elemento radiante; canal E_ij = guia; k_cut = banda/sintonia; 'array de um bit' = N clones no mesmo seed. Mesmo seed → coerência (broadside N²); dirigir o feixe exige fase diferencial — 'os clones entram à velocidade da luz'. A máquina dá as fases; o mundo dá os números.","epistemico":"fato","exemplos":[],"id":"antena_fisica","memoria":"semantica","meta":{"autor":"Aarão/broca-so","dominio":"broca_repos","fonte":""},"origem":"wiki","palavras_chave":[],"sinonimos":[],"tipo":"conceito","titulo":"A Antena Física (onde o gap atravessa o cobre)"}
\section{A Antena Física (onde o gap atravessa o cobre)}
Os objetos do painel instanciam a teoria clássica de antenas (Maxwell): gap/LFSR = corrente; idempotente e\_i = elemento radiante; canal E\_ij = guia; k\_cut = banda/sintonia; 'array de um bit' = N clones no mesmo seed. Mesmo seed \ensuremath{\to} coerência (broadside N\ensuremath{^{2}}); dirigir o feixe exige fase diferencial — 'os clones entram à velocidade da luz'. A máquina dá as fases; o mundo dá os números.
\prov{relações: usa seguranca\_dimensional; relaciona circuito\_como\_torcao}
\prov{proveniência: wiki \textperiodcentered{} fato \textperiodcentered{} confiança 1.0 \textperiodcentered{} domínio broca\_repos \textperiodcentered{} história 3 versões no jornal}

%CRISTAL {"arestas":[{"alvo":"conversa_recupera_nao_gera","confianca":1.0,"epistemico":"fato","origem":"wiki","rel":"relaciona"},{"alvo":"grafo_de_proveniencia","confianca":1.0,"epistemico":"fato","origem":"wiki","rel":"relaciona"}],"confianca":1.0,"contraexemplos":[],"descricao":"responder chama aprenderdedialogo(msg) no início; só age se a env TIFFANYAPRENDER ∈ 1,true,yes,on, delegando a aprenderdedialogo, que extrai fatos/inferências da fala e grava no conhecimento — 'a resposta da Tiffany não muda'. A escrita é DESLIGADA por padrão (fail-open silencioso: exceções engolidas). O gatilho de escrita conversacional.","epistemico":"inferencia","exemplos":[],"id":"aprender_de_fala_gate","memoria":"semantica","meta":{"autor":"broca-so","dominio":"conversa","fonte":"conversa/colapso.py:_aprender_de_dialogo; ensinar.py:aprender_de_dialogo"},"origem":"wiki","palavras_chave":[],"sinonimos":[],"tipo":"conceito","titulo":"Tiffany — aprender de fala sob TIFFANY_APRENDER"}
\section{Tiffany — aprender de fala sob TIFFANY\_APRENDER}
responder chama aprenderdedialogo(msg) no início; só age se a env TIFFANYAPRENDER \ensuremath{\in} 1,true,yes,on, delegando a aprenderdedialogo, que extrai fatos/inferências da fala e grava no conhecimento — 'a resposta da Tiffany não muda'. A escrita é DESLIGADA por padrão (fail-open silencioso: exceções engolidas). O gatilho de escrita conversacional.
\prov{relações: relaciona conversa\_recupera\_nao\_gera; relaciona grafo\_de\_proveniencia}
\prov{proveniência: wiki \textperiodcentered{} conversa/colapso.py:\_aprender\_de\_dialogo; ensinar.py:aprender\_de\_dialogo \textperiodcentered{} inferencia \textperiodcentered{} confiança 1.0 \textperiodcentered{} domínio conversa \textperiodcentered{} história 16 versões no jornal}

%CRISTAL {"arestas":[{"alvo":"o_conselho","confianca":1.0,"epistemico":"fato","origem":"wiki","rel":"parte_de"},{"alvo":"honestidade_intelectual","confianca":1.0,"epistemico":"fato","origem":"wiki","rel":"usa"}],"confianca":1.0,"contraexemplos":[],"descricao":"O fechamento se dá por voto explícito de cada IA — APROVO / APROVO-COM-RESSALVAS / REJEITO — com a ressalva obrigatória identificada antes de aprovar; a convergência tripla sem erro no núcleo é o sinal de consenso (propriedade emergente, não autoridade).","epistemico":"fato","exemplos":[],"id":"aprovacao_por_voto","memoria":"semantica","meta":{"autor":"Aarão/broca-so","dominio":"broca_repos","fonte":""},"origem":"wiki","palavras_chave":[],"sinonimos":[],"tipo":"conceito","titulo":"Aprovação por Voto Ternário"}
\section{Aprovação por Voto Ternário}
O fechamento se dá por voto explícito de cada IA — APROVO / APROVO-COM-RESSALVAS / REJEITO — com a ressalva obrigatória identificada antes de aprovar; a convergência tripla sem erro no núcleo é o sinal de consenso (propriedade emergente, não autoridade).
\prov{relações: parte\_de o\_conselho; usa honestidade\_intelectual}
\prov{proveniência: wiki \textperiodcentered{} fato \textperiodcentered{} confiança 1.0 \textperiodcentered{} domínio broca\_repos \textperiodcentered{} história 3 versões no jornal}

%CRISTAL {"arestas":[{"alvo":"maquina_algebrica","confianca":1.0,"epistemico":"fato","origem":"wiki","rel":"usa"},{"alvo":"concorrencia_cayley_dickson","confianca":1.0,"epistemico":"fato","origem":"wiki","rel":"usa"}],"confianca":1.0,"contraexemplos":[],"descricao":"Codificando janelas de OHLC/retornos como hipercomplexos, a assinatura algébrica é quadrática/bilinear simétrica nos operandos — por isso prediz a volatilidade futura (função de magnitude) mas NÃO o retorno futuro (função de sinal, que escapa à simetria z→−z). A concorrência de Wootters entre pares detecta lead-lag com screening O(N).","epistemico":"fato","exemplos":[],"id":"assinatura_algebrica_financeira","memoria":"semantica","meta":{"autor":"Aarão/broca-so","dominio":"broca_repos","fonte":""},"origem":"wiki","palavras_chave":[],"sinonimos":[],"tipo":"conceito","titulo":"Assinatura Algébrica de Séries Temporais"}
\section{Assinatura Algébrica de Séries Temporais}
Codificando janelas de OHLC/retornos como hipercomplexos, a assinatura algébrica é quadrática/bilinear simétrica nos operandos — por isso prediz a volatilidade futura (função de magnitude) mas NÃO o retorno futuro (função de sinal, que escapa à simetria z\ensuremath{\to}\ensuremath{-}z). A concorrência de Wootters entre pares detecta lead-lag com screening O(N).
\prov{relações: usa maquina\_algebrica; usa concorrencia\_cayley\_dickson}
\prov{proveniência: wiki \textperiodcentered{} fato \textperiodcentered{} confiança 1.0 \textperiodcentered{} domínio broca\_repos \textperiodcentered{} história 3 versões no jornal}

%CRISTAL {"arestas":[{"alvo":"rede_neural_algebra_aprendivel","confianca":1.0,"epistemico":"fato","origem":"wiki","rel":"usa"},{"alvo":"cifra","confianca":1.0,"epistemico":"fato","origem":"wiki","rel":"relaciona"}],"confianca":1.0,"contraexemplos":[],"descricao":"A distribuição final dos parâmetros algébricos (E_r,E_t,E_s,E_u) caracteriza univocamente uma tarefa (rotação SO(3): r≈t; texto: t>r; energia ∝ complexidade) — uma 'impressão digital' computável para transferência e classificação de datasets sem treinar a rede completa.","epistemico":"fato","exemplos":[],"id":"assinatura_energetica_tarefa","memoria":"semantica","meta":{"autor":"Lopes/Aarão","dominio":"hiper","fonte":""},"origem":"wiki","palavras_chave":[],"sinonimos":[],"tipo":"conceito","titulo":"Assinatura Energética da Tarefa"}
\section{Assinatura Energética da Tarefa}
A distribuição final dos parâmetros algébricos (E\_r,E\_t,E\_s,E\_u) caracteriza univocamente uma tarefa (rotação SO(3): r\ensuremath{\approx}t; texto: t>r; energia \ensuremath{\propto} complexidade) — uma 'impressão digital' computável para transferência e classificação de datasets sem treinar a rede completa.
\prov{relações: usa rede\_neural\_algebra\_aprendivel; relaciona cifra}
\prov{proveniência: wiki \textperiodcentered{} fato \textperiodcentered{} confiança 1.0 \textperiodcentered{} domínio hiper \textperiodcentered{} história 4 versões no jornal}

%CRISTAL {"arestas":[{"alvo":"nucleo_g2_reabsorcao","confianca":1.0,"epistemico":"fato","origem":"wiki","rel":"usa"},{"alvo":"associador","confianca":1.0,"epistemico":"fato","origem":"wiki","rel":"explica"}],"confianca":1.0,"contraexemplos":[],"descricao":"[x,y,z] = −2·ψ(x,y,z,·) com ψ=*φ; o salto de 2 para 3 corpos é o salto de φ (produto) para *φ (associador). 'Axioma do nascimento na ordem 3': em curvatura constante o vácuo é estéril, e a primeira estrutura nasce triádica.","epistemico":"fato","exemplos":[],"id":"associador_coassociativa","memoria":"semantica","meta":{"autor":"Lopes/Aarão","dominio":"hiper","fonte":""},"origem":"wiki","palavras_chave":[],"sinonimos":[],"tipo":"conceito","titulo":"O Associador é a 4-Forma Coassociativa"}
\section{O Associador é a 4-Forma Coassociativa}
[x,y,z] = \ensuremath{-}2\textperiodcentered \ensuremath{\psi}(x,y,z,\textperiodcentered ) com \ensuremath{\psi}=*\ensuremath{\varphi}; o salto de 2 para 3 corpos é o salto de \ensuremath{\varphi} (produto) para *\ensuremath{\varphi} (associador). 'Axioma do nascimento na ordem 3': em curvatura constante o vácuo é estéril, e a primeira estrutura nasce triádica.
\prov{relações: usa nucleo\_g2\_reabsorcao; explica associador}
\prov{proveniência: wiki \textperiodcentered{} fato \textperiodcentered{} confiança 1.0 \textperiodcentered{} domínio hiper \textperiodcentered{} história 4 versões no jornal}

%CRISTAL {"arestas":[{"alvo":"spec_trabalho","confianca":1.0,"epistemico":"fato","origem":"BROCA_SO.md","rel":"parte_de"},{"alvo":"utilitarismo","confianca":0.5,"epistemico":"fato","origem":"wiki","rel":"relaciona"},{"alvo":"o_que_e_no_projecto","confianca":0.5,"epistemico":"fato","origem":"wiki","rel":"relaciona"},{"alvo":"ipc","confianca":0.5,"epistemico":"fato","origem":"wiki","rel":"relaciona"},{"alvo":"word","confianca":0.5,"epistemico":"fato","origem":"wiki","rel":"relaciona"},{"alvo":"vontade_de_poder","confianca":0.5,"epistemico":"fato","origem":"wiki","rel":"relaciona"},{"alvo":"virtualizacao","confianca":0.5,"epistemico":"fato","origem":"wiki","rel":"relaciona"}],"confianca":1.0,"contraexemplos":[],"descricao":"`neuronio.c` / `neuron_io.c` — gato `A=[[1,1],[1,0]]`, streaming byte a byte.","epistemico":"fato","exemplos":[],"id":"atomo","memoria":"semantica","meta":{"dominio":"manual","nivel":6,"verificacao":"slug idêntico — enriquecer, não duplicar"},"origem":"BROCA_SO.md","palavras_chave":[],"sinonimos":[],"tipo":"conceito","titulo":"Átomo"}
\section{Átomo}
`neuronio.c` / `neuron\_io.c` — gato `A=[[1,1],[1,0]]`, streaming byte a byte.
\prov{relações: parte\_de spec\_trabalho; relaciona utilitarismo; relaciona o\_que\_e\_no\_projecto; relaciona ipc; relaciona word; relaciona vontade\_de\_poder; relaciona virtualizacao}
\prov{proveniência: BROCA\_SO.md \textperiodcentered{} fato \textperiodcentered{} confiança 1.0 \textperiodcentered{} domínio manual \textperiodcentered{} história 12 versões no jornal}

%CRISTAL {"arestas":[{"alvo":"broca-so_camada_de_trabalho_isomorfica","confianca":1.0,"epistemico":"fato","origem":"BROCA_SO.md","rel":"parte_de"},{"alvo":"teoria_do_valor","confianca":0.5,"epistemico":"fato","origem":"wiki","rel":"relaciona"},{"alvo":"matrizes","confianca":0.5,"epistemico":"fato","origem":"wiki","rel":"relaciona"},{"alvo":"word","confianca":0.5,"epistemico":"fato","origem":"wiki","rel":"relaciona"},{"alvo":"syscall","confianca":0.5,"epistemico":"fato","origem":"wiki","rel":"relaciona"},{"alvo":"virtualizacao","confianca":0.5,"epistemico":"fato","origem":"wiki","rel":"relaciona"},{"alvo":"tragedia_dos_comuns","confianca":0.5,"epistemico":"fato","origem":"wiki","rel":"relaciona"}],"confianca":1.0,"contraexemplos":[],"descricao":"Implementam **a mesma** `BrocaBackend.exec()`. Prova: mesma `BrocaMaquina` após exec.","epistemico":"fato","exemplos":[],"id":"backends","memoria":"semantica","meta":{"dominio":"manual","nivel":6,"verificacao":"slug idêntico — enriquecer, não duplicar"},"origem":"BROCA_SO.md","palavras_chave":[],"sinonimos":[],"tipo":"conceito","titulo":"Backends"}
\section{Backends}
Implementam **a mesma** `BrocaBackend.exec()`. Prova: mesma `BrocaMaquina` após exec.
\prov{relações: parte\_de broca-so\_camada\_de\_trabalho\_isomorfica; relaciona teoria\_do\_valor; relaciona matrizes; relaciona word; relaciona syscall; relaciona virtualizacao; relaciona tragedia\_dos\_comuns}
\prov{proveniência: BROCA\_SO.md \textperiodcentered{} fato \textperiodcentered{} confiança 1.0 \textperiodcentered{} domínio manual \textperiodcentered{} história 12 versões no jornal}

%CRISTAL {"arestas":[{"alvo":"bai_orbita_energia_por_trabalho","confianca":1.0,"epistemico":"fato","origem":"wiki","rel":"especializa"},{"alvo":"colapso_multicamada_peso","confianca":1.0,"epistemico":"fato","origem":"wiki","rel":"usa"},{"alvo":"conversa_recupera_nao_gera","confianca":1.0,"epistemico":"fato","origem":"wiki","rel":"relaciona"},{"alvo":"cosmologia_quantica_hub","confianca":1.0,"epistemico":"fato","origem":"wiki","rel":"usa"},{"alvo":"broca_os","confianca":1.0,"epistemico":"fato","origem":"wiki","rel":"parte_de"},{"alvo":"tiffany","confianca":1.0,"epistemico":"fato","origem":"wiki","rel":"parte_de"}],"confianca":1.0,"contraexemplos":[],"descricao":"conversa/bai.py implementa a ÓRBITA 𝒞_δ + ENERGIA δ do BAI por tipo de trabalho, e o colapso (conversa/colapso.py) a usa no lugar do peso de camada plano: o BAI SELECIONA a órbita onde o Ψ deve ressoar — as outras camadas ficam com δ menor (fora de órbita), NÃO removidas. 'Densificar, não podar' aplicado ao colapso. tipo_de_trabalho classifica a fala: definicao ('o que é X', 'quem é') → órbita do conhecimento (δ+6, corpus−2); social (saudação/turno curto) → órbita do diálogo curado (δ+4); cadeia → a órbita da cadeia CASL ativa; conversa → órbita de repouso. score_hit soma o δ da órbita (antes era _LAYER_WEIGHT fixo); a órbita entra também no desempate. Multi-tenant: o param `tenant` traz a órbita de repouso própria de cada tenant. Defaults = os pesos históricos → nada regride (test_fluxo/test_conhecimento verdes). [D] roda; o colapso agora expõe bai_trabalho/bai_orbita na observação.","epistemico":"fato","exemplos":[],"id":"bai_orbita_seleciona_colapso","memoria":"semantica","meta":{"autor":"Lopes/broca-so","dominio":"conversa","fonte":"conversa/bai.py (orbita_energia/tipo_de_trabalho); conversa/colapso.py (_bai_orbita/score_hit/colapso)"},"origem":"wiki","palavras_chave":[],"sinonimos":[],"tipo":"conceito","titulo":"O BAI seleciona a órbita do tenant no colapso (não poda a floresta)"}
\section{O BAI seleciona a órbita do tenant no colapso (não poda a floresta)}
conversa/bai.py implementa a ÓRBITA \ensuremath{\mathcal{C}}\_\ensuremath{\delta} + ENERGIA \ensuremath{\delta} do BAI por tipo de trabalho, e o colapso (conversa/colapso.py) a usa no lugar do peso de camada plano: o BAI SELECIONA a órbita onde o \ensuremath{\Psi} deve ressoar — as outras camadas ficam com \ensuremath{\delta} menor (fora de órbita), NÃO removidas. 'Densificar, não podar' aplicado ao colapso. tipo\_de\_trabalho classifica a fala: definicao ('o que é X', 'quem é') \ensuremath{\to} órbita do conhecimento (\ensuremath{\delta}+6, corpus\ensuremath{-}2); social (saudação/turno curto) \ensuremath{\to} órbita do diálogo curado (\ensuremath{\delta}+4); cadeia \ensuremath{\to} a órbita da cadeia CASL ativa; conversa \ensuremath{\to} órbita de repouso. score\_hit soma o \ensuremath{\delta} da órbita (antes era \_LAYER\_WEIGHT fixo); a órbita entra também no desempate. Multi-tenant: o param `tenant` traz a órbita de repouso própria de cada tenant. Defaults = os pesos históricos \ensuremath{\to} nada regride (test\_fluxo/test\_conhecimento verdes). [D] roda; o colapso agora expõe bai\_trabalho/bai\_orbita na observação.
\prov{relações: especializa bai\_orbita\_energia\_por\_trabalho; usa colapso\_multicamada\_peso; relaciona conversa\_recupera\_nao\_gera; usa cosmologia\_quantica\_hub; parte\_de broca\_os; parte\_de tiffany}
\prov{proveniência: wiki \textperiodcentered{} conversa/bai.py (orbita\_energia/tipo\_de\_trabalho); conversa/colapso.py (\_bai\_orbita/score\_hit/colapso) \textperiodcentered{} fato \textperiodcentered{} confiança 1.0 \textperiodcentered{} domínio conversa \textperiodcentered{} história 52 versões no jornal}

%CRISTAL {"arestas":[{"alvo":"bai_orbita_seleciona_colapso","confianca":1.0,"epistemico":"fato","origem":"wiki","rel":"parte_de"},{"alvo":"colapso_multicamada_peso","confianca":1.0,"epistemico":"fato","origem":"wiki","rel":"usa"},{"alvo":"conversa_recupera_nao_gera","confianca":1.0,"epistemico":"fato","origem":"wiki","rel":"relaciona"},{"alvo":"tiffany","confianca":1.0,"epistemico":"fato","origem":"wiki","rel":"parte_de"}],"confianca":1.0,"contraexemplos":[],"descricao":"As penalidades ad-hoc do score_hit que eram PODA por pertinência (sc −= X) migraram para conversa/bai.py:energia_conteudo — viraram RELEVO DE ENERGIA de órbita, não corte: (i) órbita da cadeia CASL ativa — dentro da cadeia +14, noutra cadeia −10 (fora de órbita); (ii) órbita de conteúdo da camada conhecimento (tags skn:) — conceito bem-formado +8, skn sem título/'—' e não-fato −24, ruído tabular/LaTeX cru −120 (bem fora da órbita conversacional). O score_hit agora soma essa energia (sc += _bai_energia_conteudo) em vez das regras inline — a floresta INTEIRA permanece, só com relevo de energia; não se poda. As regras de TURNO/continuidade (saudação, 'sim' sem contexto, face-only −999, casamento de turno, dep) FICAM no score_hit: são mecânica CASL de diálogo, não órbita 𝒞_δ. Migração com equivalência numérica exata (test_fluxo/test_conhecimento verdes; sc idêntico). [D] roda.","epistemico":"fato","exemplos":[],"id":"bai_penalidades_migradas","memoria":"semantica","meta":{"autor":"Lopes/broca-so","dominio":"conversa","fonte":"conversa/bai.py:energia_conteudo (_RE_RUIDO_TABULAR); conversa/colapso.py:score_hit/_bai_energia_conteudo"},"origem":"wiki","palavras_chave":[],"sinonimos":[],"tipo":"conceito","titulo":"As penalidades de poda do score_hit migraram para energia de órbita"}
\section{As penalidades de poda do score\_hit migraram para energia de órbita}
As penalidades ad-hoc do score\_hit que eram PODA por pertinência (sc \ensuremath{-}= X) migraram para conversa/bai.py:energia\_conteudo — viraram RELEVO DE ENERGIA de órbita, não corte: (i) órbita da cadeia CASL ativa — dentro da cadeia +14, noutra cadeia \ensuremath{-}10 (fora de órbita); (ii) órbita de conteúdo da camada conhecimento (tags skn:) — conceito bem-formado +8, skn sem título/'—' e não-fato \ensuremath{-}24, ruído tabular/LaTeX cru \ensuremath{-}120 (bem fora da órbita conversacional). O score\_hit agora soma essa energia (sc += \_bai\_energia\_conteudo) em vez das regras inline — a floresta INTEIRA permanece, só com relevo de energia; não se poda. As regras de TURNO/continuidade (saudação, 'sim' sem contexto, face-only \ensuremath{-}999, casamento de turno, dep) FICAM no score\_hit: são mecânica CASL de diálogo, não órbita \ensuremath{\mathcal{C}}\_\ensuremath{\delta}. Migração com equivalência numérica exata (test\_fluxo/test\_conhecimento verdes; sc idêntico). [D] roda.
\prov{relações: parte\_de bai\_orbita\_seleciona\_colapso; usa colapso\_multicamada\_peso; relaciona conversa\_recupera\_nao\_gera; parte\_de tiffany}
\prov{proveniência: wiki \textperiodcentered{} conversa/bai.py:energia\_conteudo (\_RE\_RUIDO\_TABULAR); conversa/colapso.py:score\_hit/\_bai\_energia\_conteudo \textperiodcentered{} fato \textperiodcentered{} confiança 1.0 \textperiodcentered{} domínio conversa \textperiodcentered{} história 38 versões no jornal}

%CRISTAL {"arestas":[{"alvo":"bai_biblioteca_de_autorizacao_isomorfica","confianca":0.95,"epistemico":"fato","origem":"wiki","rel":"confirma"},{"alvo":"borda_contemplacao","confianca":0.5,"epistemico":"fato","origem":"wiki","rel":"relaciona"},{"alvo":"grupos","confianca":0.5,"epistemico":"fato","origem":"wiki","rel":"relaciona"},{"alvo":"word","confianca":0.5,"epistemico":"fato","origem":"wiki","rel":"relaciona"},{"alvo":"virtualizacao","confianca":0.5,"epistemico":"fato","origem":"wiki","rel":"relaciona"},{"alvo":"estetica","confianca":0.5,"epistemico":"fato","origem":"wiki","rel":"relaciona"},{"alvo":"misticismo_nao_dual","confianca":0.5,"epistemico":"fato","origem":"wiki","rel":"relaciona"},{"alvo":"vida-morte","confianca":0.5,"epistemico":"fato","origem":"wiki","rel":"relaciona"}],"confianca":0.95,"contraexemplos":[],"descricao":"Referência verificada: bai_quinteto_hub.md","epistemico":"fato","exemplos":[],"id":"bai_quinteto_hub","memoria":"semantica","meta":{"dominio":"referencia","fonte":"web","nivel":5},"origem":"wiki","palavras_chave":[],"sinonimos":[],"tipo":"conceito","titulo":"bai quinteto hub"}
\section{bai quinteto hub}
Referência verificada: bai\_quinteto\_hub.md
\prov{relações: confirma bai\_biblioteca\_de\_autorizacao\_isomorfica; relaciona borda\_contemplacao; relaciona grupos; relaciona word; relaciona virtualizacao; relaciona estetica; relaciona misticismo\_nao\_dual; relaciona vida-morte}
\prov{proveniência: wiki \textperiodcentered{} web \textperiodcentered{} fato \textperiodcentered{} confiança 0.95 \textperiodcentered{} domínio referencia \textperiodcentered{} história 10 versões no jornal}

%CRISTAL {"arestas":[{"alvo":"bai_quinteto_hub","confianca":0.95,"epistemico":"fato","origem":"wiki","rel":"parte_de"}],"confianca":1.0,"contraexemplos":[],"descricao":"> **Epistemologia:** fato formal (paper CASL + teoremas + código colapso). > **No grafo:** `epistemico=fato`, confiança 0.95. > **Papel:** fechar generalizações, especializações, exemplos e contraexemplos da Biblioteca de Autorização Isomórfica.","epistemico":"fato","exemplos":[],"id":"bai_quinteto_hub_bai_hub_do_quinteto_κ_δ_φ_t_ψ","memoria":"semantica","meta":{"dominio":"referencia","fonte":"web","nivel":5},"origem":"bai_quinteto_hub.md","palavras_chave":[],"sinonimos":[],"tipo":"conceito","titulo":"BAI — hub do quinteto (κ, δ, Φ, T, Ψ)"}
\section{BAI — hub do quinteto (\ensuremath{\kappa}, \ensuremath{\delta}, \ensuremath{\Phi}, T, \ensuremath{\Psi})}
> **Epistemologia:** fato formal (paper CASL + teoremas + código colapso). > **No grafo:** `epistemico=fato`, confiança 0.95. > **Papel:** fechar generalizações, especializações, exemplos e contraexemplos da Biblioteca de Autorização Isomórfica.
\prov{relações: parte\_de bai\_quinteto\_hub}
\prov{proveniência: bai\_quinteto\_hub.md \textperiodcentered{} web \textperiodcentered{} fato \textperiodcentered{} confiança 1.0 \textperiodcentered{} domínio referencia \textperiodcentered{} história 3 versões no jornal}

%CRISTAL {"arestas":[{"alvo":"bai_quinteto_hub_bai_hub_do_quinteto_κ_δ_φ_t_ψ","confianca":1.0,"epistemico":"fato","origem":"bai_quinteto_hub.md","rel":"parte_de"},{"alvo":"virtualizacao","confianca":0.5,"epistemico":"fato","origem":"wiki","rel":"relaciona"},{"alvo":"word","confianca":0.5,"epistemico":"fato","origem":"wiki","rel":"relaciona"}],"confianca":1.0,"contraexemplos":[],"descricao":"**Peirce ≠ BAI** — homónimo distinto; BAI *usa* reticulados, Peirce *decompõe* estado+canal.","epistemico":"fato","exemplos":[],"id":"bai_quinteto_hub_contraexemplos","memoria":"semantica","meta":{"dominio":"referencia","fonte":"web","nivel":5},"origem":"bai_quinteto_hub.md","palavras_chave":[],"sinonimos":[],"tipo":"conceito","titulo":"Contraexemplos"}
\section{Contraexemplos}
**Peirce \ensuremath{\ne} BAI** — homónimo distinto; BAI *usa* reticulados, Peirce *decompõe* estado+canal.
\prov{relações: parte\_de bai\_quinteto\_hub\_bai\_hub\_do\_quinteto\_\ensuremath{\kappa}\_\ensuremath{\delta}\_\ensuremath{\varphi}\_t\_\ensuremath{\psi}; relaciona virtualizacao; relaciona word}
\prov{proveniência: bai\_quinteto\_hub.md \textperiodcentered{} web \textperiodcentered{} fato \textperiodcentered{} confiança 1.0 \textperiodcentered{} domínio referencia \textperiodcentered{} história 5 versões no jornal}

%CRISTAL {"arestas":[{"alvo":"bai_quinteto_hub_bai_hub_do_quinteto_κ_δ_φ_t_ψ","confianca":1.0,"epistemico":"fato","origem":"bai_quinteto_hub.md","rel":"parte_de"},{"alvo":"vida-morte","confianca":0.5,"epistemico":"fato","origem":"wiki","rel":"relaciona"},{"alvo":"virtualizacao","confianca":0.5,"epistemico":"fato","origem":"wiki","rel":"relaciona"}],"confianca":1.0,"contraexemplos":[],"descricao":"- `papers/casl-propagation.tex` - `conversa/colapso.py` - `papers/paper_bridge_areas.tex`","epistemico":"fato","exemplos":[],"id":"bai_quinteto_hub_fontes","memoria":"semantica","meta":{"dominio":"referencia","fonte":"web","nivel":5},"origem":"bai_quinteto_hub.md","palavras_chave":[],"sinonimos":[],"tipo":"conceito","titulo":"Fontes"}
\section{Fontes}
- `papers/casl-propagation.tex` - `conversa/colapso.py` - `papers/paper\_bridge\_areas.tex`
\prov{relações: parte\_de bai\_quinteto\_hub\_bai\_hub\_do\_quinteto\_\ensuremath{\kappa}\_\ensuremath{\delta}\_\ensuremath{\varphi}\_t\_\ensuremath{\psi}; relaciona vida-morte; relaciona virtualizacao}
\prov{proveniência: bai\_quinteto\_hub.md \textperiodcentered{} web \textperiodcentered{} fato \textperiodcentered{} confiança 1.0 \textperiodcentered{} domínio referencia \textperiodcentered{} história 5 versões no jornal}

%CRISTAL {"arestas":[{"alvo":"bai_quinteto_hub_bai_hub_do_quinteto_κ_δ_φ_t_ψ","confianca":1.0,"epistemico":"fato","origem":"bai_quinteto_hub.md","rel":"parte_de"},{"alvo":"reticulados","confianca":0.95,"epistemico":"fato","origem":"wiki","rel":"depende"},{"alvo":"monotono","confianca":0.95,"epistemico":"fato","origem":"wiki","rel":"depende"}],"confianca":1.0,"contraexemplos":[],"descricao":"**Biblioteca de Autorização Isomórfica (BAI):** delegação capacitária monótona sobre reticulados booleanos.","epistemico":"fato","exemplos":[],"id":"bai_quinteto_hub_o_que_e","memoria":"semantica","meta":{"dominio":"referencia","fonte":"web","nivel":5},"origem":"bai_quinteto_hub.md","palavras_chave":[],"sinonimos":[],"tipo":"conceito","titulo":"O que é"}
\section{O que é}
**Biblioteca de Autorização Isomórfica (BAI):** delegação capacitária monótona sobre reticulados booleanos.
\prov{relações: parte\_de bai\_quinteto\_hub\_bai\_hub\_do\_quinteto\_\ensuremath{\kappa}\_\ensuremath{\delta}\_\ensuremath{\varphi}\_t\_\ensuremath{\psi}; depende reticulados; depende monotono}
\prov{proveniência: bai\_quinteto\_hub.md \textperiodcentered{} web \textperiodcentered{} fato \textperiodcentered{} confiança 1.0 \textperiodcentered{} domínio referencia \textperiodcentered{} história 5 versões no jornal}

%CRISTAL {"arestas":[{"alvo":"bai_quinteto_hub_bai_hub_do_quinteto_κ_δ_φ_t_ψ","confianca":1.0,"epistemico":"fato","origem":"bai_quinteto_hub.md","rel":"parte_de"},{"alvo":"bai","confianca":0.95,"epistemico":"fato","origem":"wiki","rel":"confirma"},{"alvo":"tiffany","confianca":0.95,"epistemico":"fato","origem":"wiki","rel":"referencia"},{"alvo":"broca_os","confianca":0.95,"epistemico":"fato","origem":"wiki","rel":"referencia"},{"alvo":"a_decomposicao_de_peirce_celulas_e_canais","confianca":0.95,"epistemico":"fato","origem":"wiki","rel":"analogia"}],"confianca":1.0,"contraexemplos":[],"descricao":"| BAI | Nó | relação | |-----|-----|---------| | núcleo formal | bai_biblioteca_de_autorizacao_isomorfica | confirma | | alias currículo | BAI | confirma | | álgebra base | reticulados | depende | | propriedade | monotono | depende | | quintuplo | kappa | define | | orçamento | delta | define | | colapso Ψ | colapso | especializa | | paper CASL | casl-propagation | especializa | | cabeça | tiffany | referencia |","epistemico":"fato","exemplos":[],"id":"bai_quinteto_hub_ponte_broca-so","memoria":"semantica","meta":{"dominio":"referencia","fonte":"web","nivel":5},"origem":"bai_quinteto_hub.md","palavras_chave":[],"sinonimos":[],"tipo":"conceito","titulo":"Ponte Broca-SO"}
\section{Ponte Broca-SO}
| BAI | Nó | relação | |-----|-----|---------| | núcleo formal | bai\_biblioteca\_de\_autorizacao\_isomorfica | confirma | | alias currículo | BAI | confirma | | álgebra base | reticulados | depende | | propriedade | monotono | depende | | quintuplo | kappa | define | | orçamento | delta | define | | colapso \ensuremath{\Psi} | colapso | especializa | | paper CASL | casl-propagation | especializa | | cabeça | tiffany | referencia |
\prov{relações: parte\_de bai\_quinteto\_hub\_bai\_hub\_do\_quinteto\_\ensuremath{\kappa}\_\ensuremath{\delta}\_\ensuremath{\varphi}\_t\_\ensuremath{\psi}; confirma bai; referencia tiffany; referencia broca\_os; analogia a\_decomposicao\_de\_peirce\_celulas\_e\_canais}
\prov{proveniência: bai\_quinteto\_hub.md \textperiodcentered{} web \textperiodcentered{} fato \textperiodcentered{} confiança 1.0 \textperiodcentered{} domínio referencia \textperiodcentered{} história 7 versões no jornal}

%CRISTAL {"arestas":[{"alvo":"bai_quinteto_hub_bai_hub_do_quinteto_κ_δ_φ_t_ψ","confianca":1.0,"epistemico":"fato","origem":"bai_quinteto_hub.md","rel":"parte_de"},{"alvo":"casl-propagation","confianca":0.95,"epistemico":"fato","origem":"wiki","rel":"especializa"},{"alvo":"virtualizacao","confianca":0.5,"epistemico":"fato","origem":"wiki","rel":"relaciona"},{"alvo":"word","confianca":0.5,"epistemico":"fato","origem":"wiki","rel":"relaciona"}],"confianca":1.0,"contraexemplos":[],"descricao":"CASL (propagação sobre reticulados) é **parte de** BAI:","epistemico":"fato","exemplos":[],"id":"bai_quinteto_hub_propagacao_casl","memoria":"semantica","meta":{"dominio":"referencia","fonte":"web","nivel":5},"origem":"bai_quinteto_hub.md","palavras_chave":[],"sinonimos":[],"tipo":"conceito","titulo":"Propagação CASL"}
\section{Propagação CASL}
CASL (propagação sobre reticulados) é **parte de** BAI:
\prov{relações: parte\_de bai\_quinteto\_hub\_bai\_hub\_do\_quinteto\_\ensuremath{\kappa}\_\ensuremath{\delta}\_\ensuremath{\varphi}\_t\_\ensuremath{\psi}; especializa casl-propagation; relaciona virtualizacao; relaciona word}
\prov{proveniência: bai\_quinteto\_hub.md \textperiodcentered{} web \textperiodcentered{} fato \textperiodcentered{} confiança 1.0 \textperiodcentered{} domínio referencia \textperiodcentered{} história 6 versões no jornal}

%CRISTAL {"arestas":[{"alvo":"bai_quinteto_hub_bai_hub_do_quinteto_κ_δ_φ_t_ψ","confianca":1.0,"epistemico":"fato","origem":"bai_quinteto_hub.md","rel":"parte_de"},{"alvo":"kappa","confianca":0.95,"epistemico":"fato","origem":"wiki","rel":"define"},{"alvo":"delta","confianca":0.95,"epistemico":"fato","origem":"wiki","rel":"define"},{"alvo":"colapso","confianca":0.95,"epistemico":"fato","origem":"wiki","rel":"especializa"}],"confianca":1.0,"contraexemplos":[],"descricao":"| Primitiva | Papel | |-----------|-------| | **κ** | capacidade (σ, c, δ, ω, o) — unidade propagada | | **δ** | orçamento de profundidade / vigência | | **Φ** | organização estrutural pós-propagação | | **T** | atenuação deflacionária | | **Ψ** | colapso — argmax sobre candidatos vigentes 𝒞_δ |","epistemico":"fato","exemplos":[],"id":"bai_quinteto_hub_quinteto_bai","memoria":"semantica","meta":{"dominio":"referencia","fonte":"web","nivel":5},"origem":"bai_quinteto_hub.md","palavras_chave":[],"sinonimos":[],"tipo":"conceito","titulo":"Quinteto BAI"}
\section{Quinteto BAI}
| Primitiva | Papel | |-----------|-------| | **\ensuremath{\kappa}** | capacidade (\ensuremath{\sigma}, c, \ensuremath{\delta}, \ensuremath{\omega}, o) — unidade propagada | | **\ensuremath{\delta}** | orçamento de profundidade / vigência | | **\ensuremath{\Phi}** | organização estrutural pós-propagação | | **T** | atenuação deflacionária | | **\ensuremath{\Psi}** | colapso — argmax sobre candidatos vigentes \ensuremath{\mathcal{C}}\_\ensuremath{\delta} |
\prov{relações: parte\_de bai\_quinteto\_hub\_bai\_hub\_do\_quinteto\_\ensuremath{\kappa}\_\ensuremath{\delta}\_\ensuremath{\varphi}\_t\_\ensuremath{\psi}; define kappa; define delta; especializa colapso}
\prov{proveniência: bai\_quinteto\_hub.md \textperiodcentered{} web \textperiodcentered{} fato \textperiodcentered{} confiança 1.0 \textperiodcentered{} domínio referencia \textperiodcentered{} história 6 versões no jornal}

%CRISTAL {"arestas":[{"alvo":"bai_quinteto_hub_bai_hub_do_quinteto_κ_δ_φ_t_ψ","confianca":1.0,"epistemico":"fato","origem":"bai_quinteto_hub.md","rel":"parte_de"}],"confianca":1.0,"contraexemplos":[],"descricao":"- **Paper:** `papers/casl-propagation.pdf`, `papers/casl-propagation.tex` - **Código:** `conversa/colapso.py` — operador Ψ = Collapse(𝒞_δ) - **Currículo:** fase 4 matemática (reticulados) → BAI → colapso/Tiffany - **Bridge:** `papers/paper_bridge_areas.tex` — Peirce ≠ BAI (homónimos distintos)","epistemico":"fato","exemplos":[],"id":"bai_quinteto_hub_referencia","memoria":"semantica","meta":{"dominio":"referencia","fonte":"web","nivel":5},"origem":"bai_quinteto_hub.md","palavras_chave":[],"sinonimos":[],"tipo":"conceito","titulo":"Referência"}
\section{Referência}
- **Paper:** `papers/casl-propagation.pdf`, `papers/casl-propagation.tex` - **Código:** `conversa/colapso.py` — operador \ensuremath{\Psi} = Collapse(\ensuremath{\mathcal{C}}\_\ensuremath{\delta}) - **Currículo:** fase 4 matemática (reticulados) \ensuremath{\to} BAI \ensuremath{\to} colapso/Tiffany - **Bridge:** `papers/paper\_bridge\_areas.tex` — Peirce \ensuremath{\ne} BAI (homónimos distintos)
\prov{relações: parte\_de bai\_quinteto\_hub\_bai\_hub\_do\_quinteto\_\ensuremath{\kappa}\_\ensuremath{\delta}\_\ensuremath{\varphi}\_t\_\ensuremath{\psi}}
\prov{proveniência: bai\_quinteto\_hub.md \textperiodcentered{} web \textperiodcentered{} fato \textperiodcentered{} confiança 1.0 \textperiodcentered{} domínio referencia \textperiodcentered{} história 3 versões no jornal}

%CRISTAL {"arestas":[{"alvo":"bai_orbita_seleciona_colapso","confianca":1.0,"epistemico":"fato","origem":"wiki","rel":"especializa"},{"alvo":"colapso_multicamada_peso","confianca":1.0,"epistemico":"fato","origem":"wiki","rel":"usa"},{"alvo":"broca_os","confianca":1.0,"epistemico":"fato","origem":"wiki","rel":"parte_de"},{"alvo":"tiffany","confianca":1.0,"epistemico":"fato","origem":"wiki","rel":"parte_de"}],"confianca":1.0,"contraexemplos":[],"descricao":"O tenant (code/tiffanyplatform, identificado pelo home/TIFFANYHOME — base + cliente) agora é plumbado até o colapso Ψ: code/tiffany/api.py:ask passa tenant=home a respondercommeta e convresponder → responder → colapso → baiorbita → orbitaenergia(tenant=home). A ÓRBITA DE REPOUSO própria do tenant vem de <home>/baiorbita.json (opcional, com cache; aceita 'orbita':camada:δ ou …); ausente = órbita de repouso → nada regride. Sobre a órbita do tenant o BAI aplica a modulação por tipo de trabalho (definicao/social/cadeia). Assim cada tenant tem a sua órbita no colapso sem podar a floresta — o multi-tenant governa quais camadas ressoam. carregarorbitatenant aceita dict (órbita explícita) ou str (home). [D] roda (testfluxo/testconhecimento verdes; colapso expõe baitenant).","epistemico":"fato","exemplos":[],"id":"bai_tenant_plumbado","memoria":"semantica","meta":{"autor":"Lopes/broca-so","dominio":"conversa","fonte":"code/tiffany/api.py:ask (tenant=home); conversa/telemetria.py:responder_com_meta; conversa/bai.py:carregar_orbita_tenant"},"origem":"wiki","palavras_chave":[],"sinonimos":[],"tipo":"conceito","titulo":"O tenant do multi-tenant chega ao colapso: cada tenant tem sua órbita 𝒞_δ"}
\section{O tenant do multi-tenant chega ao colapso: cada tenant tem sua órbita \ensuremath{\mathcal{C}}\_\ensuremath{\delta}}
O tenant (code/tiffanyplatform, identificado pelo home/TIFFANYHOME — base + cliente) agora é plumbado até o colapso \ensuremath{\Psi}: code/tiffany/api.py:ask passa tenant=home a respondercommeta e convresponder \ensuremath{\to} responder \ensuremath{\to} colapso \ensuremath{\to} baiorbita \ensuremath{\to} orbitaenergia(tenant=home). A ÓRBITA DE REPOUSO própria do tenant vem de <home>/baiorbita.json (opcional, com cache; aceita 'orbita':camada:\ensuremath{\delta} ou …); ausente = órbita de repouso \ensuremath{\to} nada regride. Sobre a órbita do tenant o BAI aplica a modulação por tipo de trabalho (definicao/social/cadeia). Assim cada tenant tem a sua órbita no colapso sem podar a floresta — o multi-tenant governa quais camadas ressoam. carregarorbitatenant aceita dict (órbita explícita) ou str (home). [D] roda (testfluxo/testconhecimento verdes; colapso expõe baitenant).
\prov{relações: especializa bai\_orbita\_seleciona\_colapso; usa colapso\_multicamada\_peso; parte\_de broca\_os; parte\_de tiffany}
\prov{proveniência: wiki \textperiodcentered{} code/tiffany/api.py:ask (tenant=home); conversa/telemetria.py:responder\_com\_meta; conversa/bai.py:carregar\_orbita\_tenant \textperiodcentered{} fato \textperiodcentered{} confiança 1.0 \textperiodcentered{} domínio conversa \textperiodcentered{} história 43 versões no jornal}

%CRISTAL {"arestas":[{"alvo":"torre_hurwitz","confianca":1.0,"epistemico":"fato","origem":"wiki","rel":"usa"},{"alvo":"concorrencia_cayley_dickson","confianca":1.0,"epistemico":"fato","origem":"wiki","rel":"usa"},{"alvo":"alternatividade_dimensoes_criticas","confianca":1.0,"epistemico":"fato","origem":"wiki","rel":"relaciona"},{"alvo":"nao_associatividade_amputacao","confianca":1.0,"epistemico":"fato","origem":"wiki","rel":"relaciona"}],"confianca":1.0,"contraexemplos":[],"descricao":"A hierarquia Eₙ=τ/4 via Cayley-Dickson recupera concorrência/tangle exatos até 3 qubits (ℝ,ℂ,ℍ,𝕆) e FALHA em 4 qubits (sedenions: divisores de zero). A mesma classificação de Hurwitz governa SUSY (cordas, D∈3,4,6,10) e emaranhamento (QM, 0/1/2/3 qubits): dimℝ1,2,4,8.","epistemico":"fato","exemplos":[],"id":"barreira_hurwitz_qubits","memoria":"semantica","meta":{"autor":"Lopes/Aarão","dominio":"hiper","fonte":""},"origem":"wiki","palavras_chave":[],"sinonimos":[],"tipo":"conceito","titulo":"A Barreira de Hurwitz no Emaranhamento"}
\section{A Barreira de Hurwitz no Emaranhamento}
A hierarquia E\ensuremath{_{n}}=\ensuremath{\tau}/4 via Cayley-Dickson recupera concorrência/tangle exatos até 3 qubits (\ensuremath{\mathbb{R}},\ensuremath{\mathbb{C}},\ensuremath{\mathbb{H}},\ensuremath{\mathbb{O}}) e FALHA em 4 qubits (sedenions: divisores de zero). A mesma classificação de Hurwitz governa SUSY (cordas, D\ensuremath{\in}3,4,6,10) e emaranhamento (QM, 0/1/2/3 qubits): dim\ensuremath{\mathbb{R}}1,2,4,8.
\prov{relações: usa torre\_hurwitz; usa concorrencia\_cayley\_dickson; relaciona alternatividade\_dimensoes\_criticas; relaciona nao\_associatividade\_amputacao}
\prov{proveniência: wiki \textperiodcentered{} fato \textperiodcentered{} confiança 1.0 \textperiodcentered{} domínio hiper \textperiodcentered{} história 7 versões no jornal}

%CRISTAL {"arestas":[{"alvo":"termodinamica","confianca":1.0,"epistemico":"fato","origem":"wiki","rel":"referencia"},{"alvo":"conjunto_dougherty","confianca":1.0,"epistemico":"fato","origem":"wiki","rel":"analogia"},{"alvo":"falsificabilidade","confianca":1.0,"epistemico":"fato","origem":"wiki","rel":"confirma"},{"alvo":"word","confianca":0.5,"epistemico":"fato","origem":"wiki","rel":"relaciona"},{"alvo":"virtualizacao","confianca":0.5,"epistemico":"fato","origem":"wiki","rel":"relaciona"},{"alvo":"vida-morte","confianca":0.5,"epistemico":"fato","origem":"wiki","rel":"relaciona"}],"confianca":0.9,"contraexemplos":[],"descricao":"Referência verificada: birkeland_correntes_plasma.md","epistemico":"fato","exemplos":[],"id":"birkeland_correntes_plasma","memoria":"semantica","meta":{"dominio":"referencia","fonte":"web","nivel":5},"origem":"wiki","palavras_chave":[],"sinonimos":[],"tipo":"conceito","titulo":"birkeland correntes plasma"}
\section{birkeland correntes plasma}
Referência verificada: birkeland\_correntes\_plasma.md
\prov{relações: referencia termodinamica; analogia conjunto\_dougherty; confirma falsificabilidade; relaciona word; relaciona virtualizacao; relaciona vida-morte}
\prov{proveniência: wiki \textperiodcentered{} web \textperiodcentered{} fato \textperiodcentered{} confiança 0.9 \textperiodcentered{} domínio referencia \textperiodcentered{} história 99 versões no jornal}

%CRISTAL {"arestas":[{"alvo":"birkeland_correntes_plasma_correntes_de_birkeland_filamentos_de_plasma","confianca":1.0,"epistemico":"fato","origem":"birkeland_correntes_plasma.md","rel":"parte_de"},{"alvo":"word","confianca":0.5,"epistemico":"fato","origem":"wiki","rel":"relaciona"}],"confianca":1.0,"contraexemplos":[],"descricao":"Atribuir a Birkeland toros φ ou “internet cósmica” sem medição — extrapolação não falsificável.","epistemico":"fato","exemplos":[],"id":"birkeland_correntes_plasma_contraexemplos","memoria":"semantica","meta":{"dominio":"referencia","fonte":"web","nivel":5},"origem":"birkeland_correntes_plasma.md","palavras_chave":[],"sinonimos":[],"tipo":"conceito","titulo":"Contraexemplos"}
\section{Contraexemplos}
Atribuir a Birkeland toros \ensuremath{\varphi} ou “internet cósmica” sem medição — extrapolação não falsificável.
\prov{relações: parte\_de birkeland\_correntes\_plasma\_correntes\_de\_birkeland\_filamentos\_de\_plasma; relaciona word}
\prov{proveniência: birkeland\_correntes\_plasma.md \textperiodcentered{} web \textperiodcentered{} fato \textperiodcentered{} confiança 1.0 \textperiodcentered{} domínio referencia \textperiodcentered{} história 4 versões no jornal}

%CRISTAL {"fusao":[{"arestas":[{"alvo":"birkeland_correntes_plasma","confianca":0.9,"epistemico":"fato","origem":"wiki","rel":"parte_de"}],"confianca":1.0,"contraexemplos":[],"descricao":"> **Epistemologia:** fato observacional (NASA, satélites) + interpretação física standard. > **No grafo:** `epistemico=fato`, confiança 0.90. > **Papel:** teste macroscópico para hipótese Dougherty (filamentos/toros em plasma).","epistemico":"fato","exemplos":[],"id":"birkeland_correntes_plasma_correntes_de_birkeland_filamentos_de_plasma","memoria":"semantica","meta":{"dominio":"referencia","fonte":"web","nivel":5},"origem":"birkeland_correntes_plasma.md","palavras_chave":[],"sinonimos":[],"tipo":"conceito","titulo":"Correntes de Birkeland — filamentos de plasma"},{"arestas":[{"alvo":"birkeland_correntes_plasma","confianca":0.9,"epistemico":"fato","origem":"wiki","rel":"parte_de"}],"confianca":1.0,"contraexemplos":[],"descricao":"> **Epistemologia:** fato observacional (NASA, satélites) + interpretação física standard. > **No grafo:** `epistemico=fato`, confiança 0.90. > **Papel:** teste macroscópico para hipótese Dougherty (filamentos/toros em plasma).","epistemico":"fato","exemplos":[],"id":"correntes_de_birkeland_filamentos_de_plasma","memoria":"semantica","meta":{"dominio":"referencia","nivel":5},"origem":"birkeland_correntes_plasma.md","palavras_chave":[],"sinonimos":[],"tipo":"conceito","titulo":"Correntes de Birkeland — filamentos de plasma"}],"id":"birkeland_correntes_plasma_correntes_de_birkeland_filamentos_de_plasma","tipo":"conceito"}
\section{Correntes de Birkeland — filamentos de plasma}
> **Epistemologia:** fato observacional (NASA, satélites) + interpretação física standard. > **No grafo:** `epistemico=fato`, confiança 0.90. > **Papel:** teste macroscópico para hipótese Dougherty (filamentos/toros em plasma).
\prov{relações: parte\_de birkeland\_correntes\_plasma}
\prov{proveniência: birkeland\_correntes\_plasma.md \textperiodcentered{} web \textperiodcentered{} fato \textperiodcentered{} confiança 1.0 \textperiodcentered{} domínio referencia \textperiodcentered{} história 3 versões no jornal}
\prov{fusão de curadoria (texto idêntico): absorve correntes\_de\_birkeland\_filamentos\_de\_plasma --- o mesmo conteúdo por dois esquemas de endereço}

%CRISTAL {"fusao":[{"arestas":[{"alvo":"birkeland_correntes_plasma_correntes_de_birkeland_filamentos_de_plasma","confianca":1.0,"epistemico":"fato","origem":"birkeland_correntes_plasma.md","rel":"parte_de"}],"confianca":1.0,"contraexemplos":[],"descricao":"- Satélite Triad (1973): duas folhas — descida manhã, subida tarde na zona auroral - Swarm e outros: intensificam durante actividade geomagnética → aurora - Laboratório: pinches, Z-pinch evoluem para filamentos helicoidais com j||B","epistemico":"fato","exemplos":[],"id":"birkeland_correntes_plasma_evidencia","memoria":"semantica","meta":{"dominio":"referencia","fonte":"web","nivel":5},"origem":"birkeland_correntes_plasma.md","palavras_chave":[],"sinonimos":[],"tipo":"conceito","titulo":"Evidência"},{"arestas":[{"alvo":"correntes_de_birkeland_filamentos_de_plasma","confianca":1.0,"epistemico":"fato","origem":"birkeland_correntes_plasma.md","rel":"parte_de"}],"confianca":1.0,"contraexemplos":[],"descricao":"- Satélite Triad (1973): duas folhas — descida manhã, subida tarde na zona auroral - Swarm e outros: intensificam durante actividade geomagnética → aurora - Laboratório: pinches, Z-pinch evoluem para filamentos helicoidais com j||B","epistemico":"fato","exemplos":[],"id":"evidencia","memoria":"semantica","meta":{"dominio":"referencia","nivel":5},"origem":"birkeland_correntes_plasma.md","palavras_chave":[],"sinonimos":[],"tipo":"conceito","titulo":"Evidência"}],"id":"birkeland_correntes_plasma_evidencia","tipo":"conceito"}
\section{Evidência}
- Satélite Triad (1973): duas folhas — descida manhã, subida tarde na zona auroral - Swarm e outros: intensificam durante actividade geomagnética \ensuremath{\to} aurora - Laboratório: pinches, Z-pinch evoluem para filamentos helicoidais com j||B
\prov{relações: parte\_de birkeland\_correntes\_plasma\_correntes\_de\_birkeland\_filamentos\_de\_plasma}
\prov{proveniência: birkeland\_correntes\_plasma.md \textperiodcentered{} web \textperiodcentered{} fato \textperiodcentered{} confiança 1.0 \textperiodcentered{} domínio referencia \textperiodcentered{} história 3 versões no jornal}
\prov{fusão de curadoria (texto idêntico): absorve evidencia --- o mesmo conteúdo por dois esquemas de endereço}

%CRISTAL {"fusao":[{"arestas":[{"alvo":"birkeland_correntes_plasma_correntes_de_birkeland_filamentos_de_plasma","confianca":1.0,"epistemico":"fato","origem":"birkeland_correntes_plasma.md","rel":"parte_de"},{"alvo":"ponte_quantica_dougherty","confianca":0.9,"epistemico":"fato","origem":"wiki","rel":"referencia"}],"confianca":1.0,"contraexemplos":[],"descricao":"Hipótese Dougherty prevê subestruturas em escala φ em filamentos de plasma. Birkeland **confirma** existência de filamentos; **não confirma** escala φ — isso exige análise espectral independente.","epistemico":"fato","exemplos":[],"id":"birkeland_correntes_plasma_falsificacao_dougherty_teste_2","memoria":"semantica","meta":{"dominio":"referencia","fonte":"web","nivel":5},"origem":"birkeland_correntes_plasma.md","palavras_chave":[],"sinonimos":[],"tipo":"conceito","titulo":"Falsificação Dougherty (teste 2)"},{"arestas":[{"alvo":"correntes_de_birkeland_filamentos_de_plasma","confianca":1.0,"epistemico":"fato","origem":"birkeland_correntes_plasma.md","rel":"parte_de"},{"alvo":"ponte_quantica_dougherty","confianca":0.9,"epistemico":"fato","origem":"wiki","rel":"referencia"}],"confianca":1.0,"contraexemplos":[],"descricao":"Hipótese Dougherty prevê subestruturas em escala φ em filamentos de plasma. Birkeland **confirma** existência de filamentos; **não confirma** escala φ — isso exige análise espectral independente.","epistemico":"fato","exemplos":[],"id":"falsificacao_dougherty_teste_2","memoria":"semantica","meta":{"dominio":"referencia","nivel":5},"origem":"birkeland_correntes_plasma.md","palavras_chave":[],"sinonimos":[],"tipo":"conceito","titulo":"Falsificação Dougherty (teste 2)"}],"id":"birkeland_correntes_plasma_falsificacao_dougherty_teste_2","tipo":"conceito"}
\section{Falsificação Dougherty (teste 2)}
Hipótese Dougherty prevê subestruturas em escala \ensuremath{\varphi} em filamentos de plasma. Birkeland **confirma** existência de filamentos; **não confirma** escala \ensuremath{\varphi} — isso exige análise espectral independente.
\prov{relações: parte\_de birkeland\_correntes\_plasma\_correntes\_de\_birkeland\_filamentos\_de\_plasma; referencia ponte\_quantica\_dougherty}
\prov{proveniência: birkeland\_correntes\_plasma.md \textperiodcentered{} web \textperiodcentered{} fato \textperiodcentered{} confiança 1.0 \textperiodcentered{} domínio referencia \textperiodcentered{} história 4 versões no jornal}
\prov{fusão de curadoria (texto idêntico): absorve falsificacao\_dougherty\_teste\_2 --- o mesmo conteúdo por dois esquemas de endereço}

%CRISTAL {"arestas":[{"alvo":"birkeland_correntes_plasma_correntes_de_birkeland_filamentos_de_plasma","confianca":1.0,"epistemico":"fato","origem":"birkeland_correntes_plasma.md","rel":"parte_de"}],"confianca":1.0,"contraexemplos":[],"descricao":"- http://pwg.gsfc.nasa.gov/Education/wcurrent.html - https://svs.gsfc.nasa.gov/4264/ - Peratt (1986) IEEE TPS — estruturas de laboratório vs cosmos (interpretação controversa em escala galáctica)","epistemico":"fato","exemplos":[],"id":"birkeland_correntes_plasma_fontes","memoria":"semantica","meta":{"dominio":"referencia","fonte":"web","nivel":5},"origem":"birkeland_correntes_plasma.md","palavras_chave":[],"sinonimos":[],"tipo":"conceito","titulo":"Fontes"}
\section{Fontes}
- http://pwg.gsfc.nasa.gov/Education/wcurrent.html - https://svs.gsfc.nasa.gov/4264/ - Peratt (1986) IEEE TPS — estruturas de laboratório vs cosmos (interpretação controversa em escala galáctica)
\prov{relações: parte\_de birkeland\_correntes\_plasma\_correntes\_de\_birkeland\_filamentos\_de\_plasma}
\prov{proveniência: birkeland\_correntes\_plasma.md \textperiodcentered{} web \textperiodcentered{} fato \textperiodcentered{} confiança 1.0 \textperiodcentered{} domínio referencia \textperiodcentered{} história 3 versões no jornal}

%CRISTAL {"fusao":[{"arestas":[{"alvo":"birkeland_correntes_plasma_correntes_de_birkeland_filamentos_de_plasma","confianca":1.0,"epistemico":"fato","origem":"birkeland_correntes_plasma.md","rel":"parte_de"},{"alvo":"conjunto_dougherty","confianca":0.9,"epistemico":"fato","origem":"wiki","rel":"analogia"}],"confianca":1.0,"contraexemplos":[],"descricao":"Correntes **field-aligned** (paralelas às linhas de B): fluem ao longo do campo magnético terrestre, fechando circuito entre espaço e ionosfera.","epistemico":"fato","exemplos":[],"id":"birkeland_correntes_plasma_o_que_sao","memoria":"semantica","meta":{"dominio":"referencia","fonte":"web","nivel":5},"origem":"birkeland_correntes_plasma.md","palavras_chave":[],"sinonimos":[],"tipo":"conceito","titulo":"O que são"},{"arestas":[{"alvo":"correntes_de_birkeland_filamentos_de_plasma","confianca":1.0,"epistemico":"fato","origem":"birkeland_correntes_plasma.md","rel":"parte_de"},{"alvo":"conjunto_dougherty","confianca":0.9,"epistemico":"fato","origem":"wiki","rel":"analogia"}],"confianca":1.0,"contraexemplos":[],"descricao":"Correntes **field-aligned** (paralelas às linhas de B): fluem ao longo do campo magnético terrestre, fechando circuito entre espaço e ionosfera.","epistemico":"fato","exemplos":[],"id":"o_que_sao","memoria":"semantica","meta":{"dominio":"referencia","nivel":5},"origem":"birkeland_correntes_plasma.md","palavras_chave":[],"sinonimos":[],"tipo":"conceito","titulo":"O que são"}],"id":"birkeland_correntes_plasma_o_que_sao","tipo":"conceito"}
\section{O que são}
Correntes **field-aligned** (paralelas às linhas de B): fluem ao longo do campo magnético terrestre, fechando circuito entre espaço e ionosfera.
\prov{relações: parte\_de birkeland\_correntes\_plasma\_correntes\_de\_birkeland\_filamentos\_de\_plasma; analogia conjunto\_dougherty}
\prov{proveniência: birkeland\_correntes\_plasma.md \textperiodcentered{} web \textperiodcentered{} fato \textperiodcentered{} confiança 1.0 \textperiodcentered{} domínio referencia \textperiodcentered{} história 4 versões no jornal}
\prov{fusão de curadoria (texto idêntico): absorve o\_que\_sao --- o mesmo conteúdo por dois esquemas de endereço}

%CRISTAL {"arestas":[{"alvo":"birkeland_correntes_plasma_correntes_de_birkeland_filamentos_de_plasma","confianca":1.0,"epistemico":"fato","origem":"birkeland_correntes_plasma.md","rel":"parte_de"},{"alvo":"colapso","confianca":0.9,"epistemico":"fato","origem":"wiki","rel":"analogia"},{"alvo":"observador_e_observado","confianca":0.9,"epistemico":"fato","origem":"wiki","rel":"referencia"}],"confianca":1.0,"contraexemplos":[],"descricao":"| Birkeland | Broca-SO | relação | |-----------|----------|---------| | filamentos helicoidais | conjunto_dougherty (hipótese) | analogia | | vórtice / force-free | colapso, termodinamica | analogia | | aurora / medição | observador_e_observado | referencia | | escala macro | falsificabilidade Dougherty teste 2 | confirma base |","epistemico":"fato","exemplos":[],"id":"birkeland_correntes_plasma_ponte_broca-so","memoria":"semantica","meta":{"dominio":"referencia","fonte":"web","nivel":5},"origem":"birkeland_correntes_plasma.md","palavras_chave":[],"sinonimos":[],"tipo":"conceito","titulo":"Ponte Broca-SO"}
\section{Ponte Broca-SO}
| Birkeland | Broca-SO | relação | |-----------|----------|---------| | filamentos helicoidais | conjunto\_dougherty (hipótese) | analogia | | vórtice / force-free | colapso, termodinamica | analogia | | aurora / medição | observador\_e\_observado | referencia | | escala macro | falsificabilidade Dougherty teste 2 | confirma base |
\prov{relações: parte\_de birkeland\_correntes\_plasma\_correntes\_de\_birkeland\_filamentos\_de\_plasma; analogia colapso; referencia observador\_e\_observado}
\prov{proveniência: birkeland\_correntes\_plasma.md \textperiodcentered{} web \textperiodcentered{} fato \textperiodcentered{} confiança 1.0 \textperiodcentered{} domínio referencia \textperiodcentered{} história 5 versões no jornal}

%CRISTAL {"arestas":[{"alvo":"birkeland_correntes_plasma_correntes_de_birkeland_filamentos_de_plasma","confianca":1.0,"epistemico":"fato","origem":"birkeland_correntes_plasma.md","rel":"parte_de"}],"confianca":1.0,"contraexemplos":[],"descricao":"- **Kristian Birkeland (1908):** correntes eléctricas alongadas ao campo magnético ligam ionosfera e magnetosfera - **Triad (1973):** confirmação directa — ~1 MA por folha de corrente - **NASA:** [Electric Currents from Space](http://pwg.gsfc.nasa.gov/Education/wcurrent.html) - **NASA SVS:** Field-Aligned Current (Birkeland Current)","epistemico":"fato","exemplos":[],"id":"birkeland_correntes_plasma_referencia","memoria":"semantica","meta":{"dominio":"referencia","fonte":"web","nivel":5},"origem":"birkeland_correntes_plasma.md","palavras_chave":[],"sinonimos":[],"tipo":"conceito","titulo":"Referência"}
\section{Referência}
- **Kristian Birkeland (1908):** correntes eléctricas alongadas ao campo magnético ligam ionosfera e magnetosfera - **Triad (1973):** confirmação directa — \~{}1 MA por folha de corrente - **NASA:** [Electric Currents from Space](http://pwg.gsfc.nasa.gov/Education/wcurrent.html) - **NASA SVS:** Field-Aligned Current (Birkeland Current)
\prov{relações: parte\_de birkeland\_correntes\_plasma\_correntes\_de\_birkeland\_filamentos\_de\_plasma}
\prov{proveniência: birkeland\_correntes\_plasma.md \textperiodcentered{} web \textperiodcentered{} fato \textperiodcentered{} confiança 1.0 \textperiodcentered{} domínio referencia \textperiodcentered{} história 3 versões no jornal}

%CRISTAL {"fusao":[{"arestas":[{"alvo":"birkeland_correntes_plasma_correntes_de_birkeland_filamentos_de_plasma","confianca":1.0,"epistemico":"fato","origem":"birkeland_correntes_plasma.md","rel":"parte_de"},{"alvo":"mecanica_quantica","confianca":0.9,"epistemico":"fato","origem":"wiki","rel":"referencia"}],"confianca":1.0,"contraexemplos":[],"descricao":"Plasma conduz bem ao longo de B → linhas actuam como “fios”. Dinamo magnetosférico + vento solar → circuito eléctrico.","epistemico":"fato","exemplos":[],"id":"birkeland_correntes_plasma_relacao_fisica_standard","memoria":"semantica","meta":{"dominio":"referencia","fonte":"web","nivel":5},"origem":"birkeland_correntes_plasma.md","palavras_chave":[],"sinonimos":[],"tipo":"conceito","titulo":"Relação física (standard)"},{"arestas":[{"alvo":"correntes_de_birkeland_filamentos_de_plasma","confianca":1.0,"epistemico":"fato","origem":"birkeland_correntes_plasma.md","rel":"parte_de"},{"alvo":"mecanica_quantica","confianca":0.9,"epistemico":"fato","origem":"wiki","rel":"referencia"}],"confianca":1.0,"contraexemplos":[],"descricao":"Plasma conduz bem ao longo de B → linhas actuam como “fios”. Dinamo magnetosférico + vento solar → circuito eléctrico.","epistemico":"fato","exemplos":[],"id":"relacao_fisica_standard","memoria":"semantica","meta":{"dominio":"referencia","nivel":5},"origem":"birkeland_correntes_plasma.md","palavras_chave":[],"sinonimos":[],"tipo":"conceito","titulo":"Relação física (standard)"}],"id":"birkeland_correntes_plasma_relacao_fisica_standard","tipo":"conceito"}
\section{Relação física (standard)}
Plasma conduz bem ao longo de B \ensuremath{\to} linhas actuam como “fios”. Dinamo magnetosférico + vento solar \ensuremath{\to} circuito eléctrico.
\prov{relações: parte\_de birkeland\_correntes\_plasma\_correntes\_de\_birkeland\_filamentos\_de\_plasma; referencia mecanica\_quantica}
\prov{proveniência: birkeland\_correntes\_plasma.md \textperiodcentered{} web \textperiodcentered{} fato \textperiodcentered{} confiança 1.0 \textperiodcentered{} domínio referencia \textperiodcentered{} história 4 versões no jornal}
\prov{fusão de curadoria (texto idêntico): absorve relacao\_fisica\_standard --- o mesmo conteúdo por dois esquemas de endereço}

%CRISTAL {"arestas":[{"alvo":"broca-so_cristalizado","confianca":1.0,"epistemico":"fato","origem":"BROCA_CRISTAL.md","rel":"parte_de"},{"alvo":"word","confianca":0.5,"epistemico":"fato","origem":"wiki","rel":"relaciona"},{"alvo":"syscall","confianca":0.5,"epistemico":"fato","origem":"wiki","rel":"relaciona"},{"alvo":"vontade_de_poder","confianca":0.5,"epistemico":"fato","origem":"wiki","rel":"relaciona"},{"alvo":"contradicao","confianca":0.5,"epistemico":"fato","origem":"wiki","rel":"relaciona"},{"alvo":"while","confianca":0.5,"epistemico":"fato","origem":"wiki","rel":"relaciona"},{"alvo":"virtualizacao","confianca":0.5,"epistemico":"fato","origem":"wiki","rel":"relaciona"},{"alvo":"cantor_ouro","confianca":0.5,"epistemico":"fato","origem":"wiki","rel":"relaciona"},{"alvo":"casos_degenerados","confianca":0.5,"epistemico":"fato","origem":"wiki","rel":"relaciona"}],"confianca":1.0,"contraexemplos":[],"descricao":"O clássico na borda **não processa o núcleo** — **recebe** a projeção Koch e **mostra**:","epistemico":"fato","exemplos":[],"id":"borda_computador_classico_util_e_necessario","memoria":"semantica","meta":{"dominio":"manual","nivel":6,"verificacao":"slug idêntico — enriquecer, não duplicar"},"origem":"BROCA_CRISTAL.md","palavras_chave":[],"sinonimos":[],"tipo":"conceito","titulo":"Borda (computador clássico — útil e necessário)"}
\section{Borda (computador clássico — útil e necessário)}
O clássico na borda **não processa o núcleo** — **recebe** a projeção Koch e **mostra**:
\prov{relações: parte\_de broca-so\_cristalizado; relaciona word; relaciona syscall; relaciona vontade\_de\_poder; relaciona contradicao; relaciona while; relaciona virtualizacao; relaciona cantor\_ouro; relaciona casos\_degenerados}
\prov{proveniência: BROCA\_CRISTAL.md \textperiodcentered{} fato \textperiodcentered{} confiança 1.0 \textperiodcentered{} domínio manual \textperiodcentered{} história 12 versões no jornal}

%CRISTAL {"arestas":[{"alvo":"broca-so_camada_de_trabalho_isomorfica","confianca":1.0,"epistemico":"fato","origem":"BROCA_SO.md","rel":"parte_de"},{"alvo":"utilitarismo","confianca":0.5,"epistemico":"fato","origem":"wiki","rel":"relaciona"},{"alvo":"ipc","confianca":0.5,"epistemico":"fato","origem":"wiki","rel":"relaciona"},{"alvo":"melhor_candidato","confianca":0.5,"epistemico":"fato","origem":"wiki","rel":"relaciona"},{"alvo":"o_que_e_no_projecto","confianca":0.5,"epistemico":"fato","origem":"wiki","rel":"relaciona"},{"alvo":"funil_cumulativo","confianca":0.5,"epistemico":"fato","origem":"wiki","rel":"relaciona"},{"alvo":"fase_2_tecido_de_conhecimento","confianca":0.5,"epistemico":"fato","origem":"wiki","rel":"relaciona"},{"alvo":"consulta_simples","confianca":0.5,"epistemico":"fato","origem":"wiki","rel":"relaciona"},{"alvo":"recursao","confianca":0.5,"epistemico":"fato","origem":"wiki","rel":"relaciona"},{"alvo":"compilador","confianca":0.5,"epistemico":"fato","origem":"wiki","rel":"relaciona"}],"confianca":1.0,"contraexemplos":[],"descricao":"Tudo que **não** entra em `broca_so.h`:","epistemico":"fato","exemplos":[],"id":"borda_contemplacao","memoria":"semantica","meta":{"dominio":"manual","nivel":6,"verificacao":"slug idêntico — enriquecer, não duplicar"},"origem":"BROCA_SO.md","palavras_chave":[],"sinonimos":[],"tipo":"conceito","titulo":"Borda (contemplação)"}
\section{Borda (contemplação)}
Tudo que **não** entra em `broca\_so.h`:
\prov{relações: parte\_de broca-so\_camada\_de\_trabalho\_isomorfica; relaciona utilitarismo; relaciona ipc; relaciona melhor\_candidato; relaciona o\_que\_e\_no\_projecto; relaciona funil\_cumulativo; relaciona fase\_2\_tecido\_de\_conhecimento; relaciona consulta\_simples; relaciona recursao; relaciona compilador}
\prov{proveniência: BROCA\_SO.md \textperiodcentered{} fato \textperiodcentered{} confiança 1.0 \textperiodcentered{} domínio manual \textperiodcentered{} história 15 versões no jornal}

%CRISTAL {"arestas":[{"alvo":"substrato_prtsuv_h","confianca":1.0,"epistemico":"fato","origem":"wiki","rel":"usa"},{"alvo":"nucleo_g2_reabsorcao","confianca":1.0,"epistemico":"fato","origem":"wiki","rel":"relaciona"}],"confianca":1.0,"contraexemplos":[],"descricao":"Conjectura (a mais especulativa): o parâmetro h ativa a estrutura G₂ em ℝ⁷ (quebra SO(7)→G₂ estabilizando a 4-forma de Cayley), preservando 1/8 da SUSY e realizando M-theory na família (contagem D=8+2+1=11). A norma da 4-forma A_h=‖φ‖²/[n(n−1)]=4 em ℝ⁷ é sinal numérico independente.","epistemico":"hipotese","exemplos":[],"id":"botao_h_mtheory_g2","memoria":"semantica","meta":{"autor":"Lopes/Aarão","dominio":"hiper","fonte":""},"origem":"wiki","palavras_chave":[],"sinonimos":[],"tipo":"conceito","titulo":"O Botão h e a M-Theory em Holonomia G₂"}
\section{O Botão h e a M-Theory em Holonomia G\ensuremath{_{2}}}
Conjectura (a mais especulativa): o parâmetro h ativa a estrutura G\ensuremath{_{2}} em \ensuremath{\mathbb{R}}\ensuremath{^{7}} (quebra SO(7)\ensuremath{\to}G\ensuremath{_{2}} estabilizando a 4-forma de Cayley), preservando 1/8 da SUSY e realizando M-theory na família (contagem D=8+2+1=11). A norma da 4-forma A\_h=\ensuremath{\|}\ensuremath{\varphi}\ensuremath{\|}\ensuremath{^{2}}/[n(n\ensuremath{-}1)]=4 em \ensuremath{\mathbb{R}}\ensuremath{^{7}} é sinal numérico independente.
\prov{relações: usa substrato\_prtsuv\_h; relaciona nucleo\_g2\_reabsorcao}
\prov{proveniência: wiki \textperiodcentered{} hipotese \textperiodcentered{} confiança 1.0 \textperiodcentered{} domínio hiper \textperiodcentered{} história 4 versões no jornal}

%CRISTAL {"arestas":[{"alvo":"broca_cristal","confianca":1.0,"epistemico":"fato","origem":"manual","rel":"parte_de"},{"alvo":"historia","confianca":0.5,"epistemico":"fato","origem":"wiki","rel":"relaciona"},{"alvo":"utilitarismo","confianca":0.5,"epistemico":"fato","origem":"wiki","rel":"relaciona"},{"alvo":"resultado_completo","confianca":0.5,"epistemico":"fato","origem":"wiki","rel":"relaciona"},{"alvo":"ficheiros","confianca":0.5,"epistemico":"fato","origem":"wiki","rel":"relaciona"},{"alvo":"colapso_koch_tiffany","confianca":0.5,"epistemico":"fato","origem":"wiki","rel":"relaciona"},{"alvo":"teoria_do_valor","confianca":0.5,"epistemico":"fato","origem":"wiki","rel":"relaciona"},{"alvo":"main","confianca":0.5,"epistemico":"fato","origem":"wiki","rel":"relaciona"},{"alvo":"gentil_16_soma_dos_termos_de_uma_pam","confianca":0.5,"epistemico":"fato","origem":"wiki","rel":"relaciona"},{"alvo":"fixamos_e_o_curinga","confianca":0.5,"epistemico":"fato","origem":"wiki","rel":"relaciona"},{"alvo":"vida-morte","confianca":0.5,"epistemico":"fato","origem":"wiki","rel":"relaciona"},{"alvo":"pitagorismo","confianca":0.5,"epistemico":"fato","origem":"wiki","rel":"relaciona"},{"alvo":"parabola_isa","confianca":0.5,"epistemico":"fato","origem":"wiki","rel":"relaciona"},{"alvo":"word","confianca":0.5,"epistemico":"fato","origem":"wiki","rel":"relaciona"},{"alvo":"readme_cristal","confianca":0.5,"epistemico":"fato","origem":"wiki","rel":"relaciona"},{"alvo":"transcricao","confianca":0.5,"epistemico":"fato","origem":"wiki","rel":"relaciona"},{"alvo":"filesystem","confianca":0.5,"epistemico":"fato","origem":"wiki","rel":"relaciona"}],"confianca":1.0,"contraexemplos":[],"descricao":"**Decisão de arquitectura — não reabrir sem prova formal.**","epistemico":"fato","exemplos":[],"id":"broca-so_cristalizado","memoria":"semantica","meta":{"dominio":"manual","nivel":6,"verificacao":"slug idêntico — enriquecer, não duplicar"},"origem":"BROCA_CRISTAL.md","palavras_chave":[],"sinonimos":[],"tipo":"conceito","titulo":"Broca-so cristalizado"}
\section{Broca-so cristalizado}
**Decisão de arquitectura — não reabrir sem prova formal.**
\prov{relações: parte\_de broca\_cristal; relaciona historia; relaciona utilitarismo; relaciona resultado\_completo; relaciona ficheiros; relaciona colapso\_koch\_tiffany; relaciona teoria\_do\_valor; relaciona main; relaciona gentil\_16\_soma\_dos\_termos\_de\_uma\_pam; relaciona fixamos\_e\_o\_curinga; relaciona vida-morte; relaciona pitagorismo; relaciona parabola\_isa; relaciona word; relaciona readme\_cristal; relaciona transcricao; relaciona filesystem}
\prov{proveniência: BROCA\_CRISTAL.md \textperiodcentered{} fato \textperiodcentered{} confiança 1.0 \textperiodcentered{} domínio manual \textperiodcentered{} história 20 versões no jornal}

%CRISTAL {"arestas":[{"alvo":"colapso","confianca":1.0,"epistemico":"fato","origem":"manual","rel":"explica"},{"alvo":"tiffany","confianca":1.0,"epistemico":"fato","origem":"manual","rel":"referencia"},{"alvo":"transcricao","confianca":0.5,"epistemico":"fato","origem":"wiki","rel":"relaciona"},{"alvo":"teste_ordem","confianca":0.5,"epistemico":"fato","origem":"wiki","rel":"relaciona"},{"alvo":"resposta_a_pergunta_dificil","confianca":0.5,"epistemico":"fato","origem":"wiki","rel":"relaciona"},{"alvo":"vida-morte","confianca":0.5,"epistemico":"fato","origem":"wiki","rel":"relaciona"},{"alvo":"parabola_destruir","confianca":0.5,"epistemico":"fato","origem":"wiki","rel":"relaciona"},{"alvo":"por_sessao_conversadadostelemetriasessaojsonl","confianca":0.5,"epistemico":"fato","origem":"wiki","rel":"relaciona"},{"alvo":"pow_geometria","confianca":0.5,"epistemico":"fato","origem":"wiki","rel":"relaciona"},{"alvo":"topologia","confianca":0.5,"epistemico":"fato","origem":"wiki","rel":"relaciona"},{"alvo":"projeto_cristaljson","confianca":0.5,"epistemico":"fato","origem":"wiki","rel":"relaciona"},{"alvo":"geometria_hub_broca","confianca":0.5,"epistemico":"fato","origem":"wiki","rel":"relaciona"},{"alvo":"koch_cantor_fib_prova","confianca":0.5,"epistemico":"fato","origem":"wiki","rel":"relaciona"},{"alvo":"testes","confianca":0.5,"epistemico":"fato","origem":"wiki","rel":"relaciona"},{"alvo":"motivos","confianca":0.5,"epistemico":"fato","origem":"wiki","rel":"relaciona"},{"alvo":"ds_universo","confianca":0.5,"epistemico":"fato","origem":"wiki","rel":"relaciona"},{"alvo":"pow_metal_setor","confianca":0.5,"epistemico":"fato","origem":"wiki","rel":"relaciona"},{"alvo":"exatidao","confianca":0.5,"epistemico":"fato","origem":"wiki","rel":"relaciona"},{"alvo":"recursao","confianca":0.5,"epistemico":"fato","origem":"wiki","rel":"relaciona"},{"alvo":"ponteiro","confianca":0.5,"epistemico":"fato","origem":"wiki","rel":"relaciona"},{"alvo":"rust","confianca":0.5,"epistemico":"fato","origem":"wiki","rel":"relaciona"},{"alvo":"pilha","confianca":0.5,"epistemico":"fato","origem":"wiki","rel":"relaciona"},{"alvo":"complexos","confianca":0.5,"epistemico":"fato","origem":"wiki","rel":"relaciona"},{"alvo":"parabola_pa","confianca":0.5,"epistemico":"fato","origem":"wiki","rel":"relaciona"},{"alvo":"tiffany_cabeca_broca","confianca":0.5,"epistemico":"fato","origem":"wiki","rel":"relaciona"},{"alvo":"mecanica_quantica_hub","confianca":0.5,"epistemico":"fato","origem":"wiki","rel":"relaciona"},{"alvo":"pow_metal_funil","confianca":0.5,"epistemico":"fato","origem":"wiki","rel":"relaciona"},{"alvo":"virtualizacao","confianca":0.5,"epistemico":"fato","origem":"wiki","rel":"relaciona"},{"alvo":"hardware_sequencias_de_cauchy","confianca":0.5,"epistemico":"fato","origem":"wiki","rel":"relaciona"},{"alvo":"pow_metal_ressonancia","confianca":0.5,"epistemico":"fato","origem":"wiki","rel":"relaciona"},{"alvo":"thread","confianca":0.5,"epistemico":"fato","origem":"wiki","rel":"relaciona"},{"alvo":"fechos_centro_da_janela","confianca":0.5,"epistemico":"fato","origem":"wiki","rel":"relaciona"},{"alvo":"pow_metal_pre_fecho","confianca":0.5,"epistemico":"fato","origem":"wiki","rel":"relaciona"},{"alvo":"numa","confianca":0.5,"epistemico":"fato","origem":"wiki","rel":"relaciona"}],"confianca":1.0,"contraexemplos":[],"descricao":"Manual: ula/BROCA_CRISTAL.md","epistemico":"fato","exemplos":[],"id":"broca_cristal","memoria":"semantica","meta":{"arquivo":"ula/BROCA_CRISTAL.md","dominio":"manual","nivel":6},"origem":"manual","palavras_chave":[],"sinonimos":[],"tipo":"conceito","titulo":"BROCA_CRISTAL"}
\section{BROCA\_CRISTAL}
Manual: ula/BROCA\_CRISTAL.md
\prov{relações: explica colapso; referencia tiffany; relaciona transcricao; relaciona teste\_ordem; relaciona resposta\_a\_pergunta\_dificil; relaciona vida-morte; relaciona parabola\_destruir; relaciona por\_sessao\_conversadadostelemetriasessaojsonl; relaciona pow\_geometria; relaciona topologia; relaciona projeto\_cristaljson; relaciona geometria\_hub\_broca; relaciona koch\_cantor\_fib\_prova; relaciona testes; relaciona motivos; relaciona ds\_universo; relaciona pow\_metal\_setor; relaciona exatidao; relaciona recursao; relaciona ponteiro; relaciona rust; relaciona pilha; relaciona complexos; relaciona parabola\_pa; relaciona tiffany\_cabeca\_broca; relaciona mecanica\_quantica\_hub; relaciona pow\_metal\_funil; relaciona virtualizacao; relaciona hardware\_sequencias\_de\_cauchy; relaciona pow\_metal\_ressonancia; relaciona thread; relaciona fechos\_centro\_da\_janela; relaciona pow\_metal\_pre\_fecho; relaciona numa}
\prov{proveniência: manual \textperiodcentered{} fato \textperiodcentered{} confiança 1.0 \textperiodcentered{} domínio manual \textperiodcentered{} história 63 versões no jornal}

%CRISTAL {"arestas":[{"alvo":"broca_os","confianca":0.95,"epistemico":"fato","origem":"wiki","rel":"confirma"},{"alvo":"mamifero","confianca":0.5,"epistemico":"fato","origem":"wiki","rel":"relaciona"},{"alvo":"goldcoin","confianca":0.5,"epistemico":"fato","origem":"wiki","rel":"relaciona"},{"alvo":"virtualizacao","confianca":0.5,"epistemico":"fato","origem":"wiki","rel":"relaciona"},{"alvo":"transcricao","confianca":0.5,"epistemico":"fato","origem":"wiki","rel":"relaciona"},{"alvo":"utilitarismo","confianca":0.5,"epistemico":"fato","origem":"wiki","rel":"relaciona"},{"alvo":"resultado_anterior_identidade_nao_cobertura","confianca":0.5,"epistemico":"fato","origem":"wiki","rel":"relaciona"},{"alvo":"scheduler","confianca":0.5,"epistemico":"fato","origem":"wiki","rel":"relaciona"},{"alvo":"vida-morte","confianca":0.5,"epistemico":"fato","origem":"wiki","rel":"relaciona"},{"alvo":"ruido","confianca":0.5,"epistemico":"fato","origem":"wiki","rel":"relaciona"},{"alvo":"pedagogia","confianca":0.5,"epistemico":"fato","origem":"wiki","rel":"relaciona"},{"alvo":"paper_2_reta_ouro","confianca":0.5,"epistemico":"fato","origem":"wiki","rel":"relaciona"},{"alvo":"ficheiros","confianca":0.5,"epistemico":"fato","origem":"wiki","rel":"relaciona"},{"alvo":"mecanica","confianca":0.5,"epistemico":"fato","origem":"wiki","rel":"relaciona"},{"alvo":"modo_espectral_spect","confianca":0.5,"epistemico":"fato","origem":"wiki","rel":"relaciona"},{"alvo":"resultado_completo","confianca":0.5,"epistemico":"fato","origem":"wiki","rel":"relaciona"},{"alvo":"pitagorismo","confianca":0.5,"epistemico":"fato","origem":"wiki","rel":"relaciona"},{"alvo":"pow_metal_ressonancia_val","confianca":0.5,"epistemico":"fato","origem":"wiki","rel":"relaciona"},{"alvo":"experimento_3_inverter","confianca":0.5,"epistemico":"fato","origem":"wiki","rel":"relaciona"},{"alvo":"main","confianca":0.5,"epistemico":"fato","origem":"wiki","rel":"relaciona"},{"alvo":"vetores","confianca":0.5,"epistemico":"fato","origem":"wiki","rel":"relaciona"},{"alvo":"orbita_pell_hardware","confianca":0.5,"epistemico":"fato","origem":"wiki","rel":"relaciona"},{"alvo":"veredito_experimental","confianca":0.5,"epistemico":"fato","origem":"wiki","rel":"relaciona"},{"alvo":"prova_da_maquina_espectral","confianca":0.5,"epistemico":"fato","origem":"wiki","rel":"relaciona"}],"confianca":0.95,"contraexemplos":[],"descricao":"Referência verificada: broca_os_hub.md","epistemico":"fato","exemplos":[],"id":"broca_os_hub","memoria":"semantica","meta":{"dominio":"referencia","fonte":"web","nivel":5},"origem":"wiki","palavras_chave":[],"sinonimos":[],"tipo":"conceito","titulo":"broca os hub"}
\section{broca os hub}
Referência verificada: broca\_os\_hub.md
\prov{relações: confirma broca\_os; relaciona mamifero; relaciona goldcoin; relaciona virtualizacao; relaciona transcricao; relaciona utilitarismo; relaciona resultado\_anterior\_identidade\_nao\_cobertura; relaciona scheduler; relaciona vida-morte; relaciona ruido; relaciona pedagogia; relaciona paper\_2\_reta\_ouro; relaciona ficheiros; relaciona mecanica; relaciona modo\_espectral\_spect; relaciona resultado\_completo; relaciona pitagorismo; relaciona pow\_metal\_ressonancia\_val; relaciona experimento\_3\_inverter; relaciona main; relaciona vetores; relaciona orbita\_pell\_hardware; relaciona veredito\_experimental; relaciona prova\_da\_maquina\_espectral}
\prov{proveniência: wiki \textperiodcentered{} web \textperiodcentered{} fato \textperiodcentered{} confiança 0.95 \textperiodcentered{} domínio referencia \textperiodcentered{} história 29 versões no jornal}

%CRISTAL {"fusao":[{"arestas":[{"alvo":"broca_os_hub","confianca":0.95,"epistemico":"fato","origem":"wiki","rel":"parte_de"}],"confianca":1.0,"contraexemplos":[],"descricao":"> **Epistemologia:** fato de implementação (papers + código + currículo). > **No grafo:** `epistemico=fato`, confiança 0.95. > **Papel:** fechar dependências e especializações do SO fractal.","epistemico":"fato","exemplos":[],"id":"broca_os_hub_broca_os_hub_do_sistema_fractal","memoria":"semantica","meta":{"dominio":"referencia","fonte":"web","nivel":5},"origem":"broca_os_hub.md","palavras_chave":[],"sinonimos":[],"tipo":"conceito","titulo":"Broca OS — hub do sistema fractal"},{"arestas":[{"alvo":"broca_os","confianca":0.95,"epistemico":"fato","origem":"paper","rel":"parte_de"}],"confianca":1.0,"contraexemplos":[],"descricao":"> **Epistemologia:** fato de implementação (papers + código + currículo). > **No grafo:** `epistemico=fato`, confiança 0.95. > **Papel:** fechar dependências e especializações do SO fractal.","epistemico":"fato","exemplos":[],"id":"broca_os_hub_do_sistema_fractal","memoria":"semantica","meta":{"dominio":"referencia","nivel":5},"origem":"broca_os_hub.md","palavras_chave":[],"sinonimos":[],"tipo":"conceito","titulo":"Broca OS — hub do sistema fractal"}],"id":"broca_os_hub_broca_os_hub_do_sistema_fractal","tipo":"conceito"}
\section{Broca OS — hub do sistema fractal}
> **Epistemologia:** fato de implementação (papers + código + currículo). > **No grafo:** `epistemico=fato`, confiança 0.95. > **Papel:** fechar dependências e especializações do SO fractal.
\prov{relações: parte\_de broca\_os\_hub}
\prov{proveniência: broca\_os\_hub.md \textperiodcentered{} web \textperiodcentered{} fato \textperiodcentered{} confiança 1.0 \textperiodcentered{} domínio referencia \textperiodcentered{} história 3 versões no jornal}
\prov{fusão de curadoria (texto idêntico): absorve broca\_os\_hub\_do\_sistema\_fractal --- o mesmo conteúdo por dois esquemas de endereço}

%CRISTAL {"arestas":[{"alvo":"broca_os_hub_broca_os_hub_do_sistema_fractal","confianca":1.0,"epistemico":"fato","origem":"broca_os_hub.md","rel":"parte_de"},{"alvo":"vida-morte","confianca":0.5,"epistemico":"fato","origem":"wiki","rel":"relaciona"},{"alvo":"virtualizacao","confianca":0.5,"epistemico":"fato","origem":"wiki","rel":"relaciona"},{"alvo":"word","confianca":0.5,"epistemico":"fato","origem":"wiki","rel":"relaciona"},{"alvo":"pipeline_de_ingestao","confianca":0.5,"epistemico":"fato","origem":"wiki","rel":"relaciona"},{"alvo":"pell","confianca":0.5,"epistemico":"fato","origem":"wiki","rel":"relaciona"},{"alvo":"sequencia_pow_nativa","confianca":0.5,"epistemico":"fato","origem":"wiki","rel":"relaciona"},{"alvo":"vontade_de_poder","confianca":0.5,"epistemico":"fato","origem":"wiki","rel":"relaciona"},{"alvo":"camadas_de_tecido","confianca":0.5,"epistemico":"fato","origem":"wiki","rel":"relaciona"},{"alvo":"fibonacci_multidimensional","confianca":0.5,"epistemico":"fato","origem":"wiki","rel":"relaciona"},{"alvo":"painel_estelar","confianca":0.5,"epistemico":"fato","origem":"wiki","rel":"relaciona"},{"alvo":"poder_disciplinar_foucault","confianca":0.5,"epistemico":"fato","origem":"wiki","rel":"relaciona"},{"alvo":"consciencia","confianca":0.5,"epistemico":"fato","origem":"wiki","rel":"relaciona"},{"alvo":"sintaxe","confianca":0.5,"epistemico":"fato","origem":"wiki","rel":"relaciona"},{"alvo":"lemma_lemunit_e_unitaria_0xt2dtt_-","confianca":0.5,"epistemico":"fato","origem":"wiki","rel":"relaciona"},{"alvo":"telemetria_shadow_producao","confianca":0.5,"epistemico":"fato","origem":"wiki","rel":"relaciona"}],"confianca":1.0,"contraexemplos":[],"descricao":"Broca OS **≠** Linux kernel — contemplação corre em borda clássica.","epistemico":"fato","exemplos":[],"id":"broca_os_hub_contraexemplos","memoria":"semantica","meta":{"dominio":"referencia","fonte":"web","nivel":5},"origem":"broca_os_hub.md","palavras_chave":[],"sinonimos":[],"tipo":"conceito","titulo":"Contraexemplos"}
\section{Contraexemplos}
Broca OS **\ensuremath{\ne}** Linux kernel — contemplação corre em borda clássica.
\prov{relações: parte\_de broca\_os\_hub\_broca\_os\_hub\_do\_sistema\_fractal; relaciona vida-morte; relaciona virtualizacao; relaciona word; relaciona pipeline\_de\_ingestao; relaciona pell; relaciona sequencia\_pow\_nativa; relaciona vontade\_de\_poder; relaciona camadas\_de\_tecido; relaciona fibonacci\_multidimensional; relaciona painel\_estelar; relaciona poder\_disciplinar\_foucault; relaciona consciencia; relaciona sintaxe; relaciona lemma\_lemunit\_e\_unitaria\_0xt2dtt\_-; relaciona telemetria\_shadow\_producao}
\prov{proveniência: broca\_os\_hub.md \textperiodcentered{} web \textperiodcentered{} fato \textperiodcentered{} confiança 1.0 \textperiodcentered{} domínio referencia \textperiodcentered{} história 18 versões no jornal}

%CRISTAL {"arestas":[{"alvo":"broca_os_hub_broca_os_hub_do_sistema_fractal","confianca":1.0,"epistemico":"fato","origem":"broca_os_hub.md","rel":"parte_de"},{"alvo":"destruir_a_parabola","confianca":0.5,"epistemico":"fato","origem":"wiki","rel":"relaciona"},{"alvo":"virtualizacao","confianca":0.5,"epistemico":"fato","origem":"wiki","rel":"relaciona"},{"alvo":"vida-morte","confianca":0.5,"epistemico":"fato","origem":"wiki","rel":"relaciona"}],"confianca":1.0,"contraexemplos":[],"descricao":"- `ula/BROCA_SO.md` - `papers/broca_os.tex` - `neuronio/primitivas_broca.py`","epistemico":"fato","exemplos":[],"id":"broca_os_hub_fontes","memoria":"semantica","meta":{"dominio":"referencia","fonte":"web","nivel":5},"origem":"broca_os_hub.md","palavras_chave":[],"sinonimos":[],"tipo":"conceito","titulo":"Fontes"}
\section{Fontes}
- `ula/BROCA\_SO.md` - `papers/broca\_os.tex` - `neuronio/primitivas\_broca.py`
\prov{relações: parte\_de broca\_os\_hub\_broca\_os\_hub\_do\_sistema\_fractal; relaciona destruir\_a\_parabola; relaciona virtualizacao; relaciona vida-morte}
\prov{proveniência: broca\_os\_hub.md \textperiodcentered{} web \textperiodcentered{} fato \textperiodcentered{} confiança 1.0 \textperiodcentered{} domínio referencia \textperiodcentered{} história 6 versões no jornal}

%CRISTAL {"fusao":[{"arestas":[{"alvo":"broca_os_hub_broca_os_hub_do_sistema_fractal","confianca":1.0,"epistemico":"fato","origem":"broca_os_hub.md","rel":"parte_de"},{"alvo":"gato","confianca":0.95,"epistemico":"fato","origem":"wiki","rel":"depende"},{"alvo":"processos","confianca":0.95,"epistemico":"fato","origem":"wiki","rel":"depende"},{"alvo":"word","confianca":0.5,"epistemico":"fato","origem":"wiki","rel":"relaciona"}],"confianca":1.0,"contraexemplos":[],"descricao":"**SO fractal isomórfico:** camada de trabalho (`broca_so.h`) + contemplação (`broca_borda.h`).","epistemico":"fato","exemplos":[],"id":"broca_os_hub_o_que_e","memoria":"semantica","meta":{"dominio":"referencia","fonte":"web","nivel":5},"origem":"broca_os_hub.md","palavras_chave":[],"sinonimos":[],"tipo":"conceito","titulo":"O que é"},{"arestas":[{"alvo":"tiffany_cabeca_conversacional_broca-so","confianca":1.0,"epistemico":"fato","origem":"tiffany_cabeca_broca.md","rel":"parte_de"},{"alvo":"colapso","confianca":0.95,"epistemico":"fato","origem":"wiki","rel":"usa"},{"alvo":"broca_os","confianca":0.95,"epistemico":"fato","origem":"wiki","rel":"parte_de"},{"alvo":"gato","confianca":0.95,"epistemico":"fato","origem":"wiki","rel":"depende"},{"alvo":"processos","confianca":0.95,"epistemico":"fato","origem":"wiki","rel":"depende"}],"confianca":1.0,"contraexemplos":[],"descricao":"**SO fractal isomórfico:** camada de trabalho (`broca_so.h`) + contemplação (`broca_borda.h`).","epistemico":"fato","exemplos":[],"id":"o_que_e","memoria":"semantica","meta":{"dominio":"referencia","nivel":5},"origem":"broca_os_hub.md","palavras_chave":[],"sinonimos":[],"tipo":"conceito","titulo":"O que é"}],"id":"broca_os_hub_o_que_e","tipo":"conceito"}
\section{O que é}
**SO fractal isomórfico:** camada de trabalho (`broca\_so.h`) + contemplação (`broca\_borda.h`).
\prov{relações: parte\_de broca\_os\_hub\_broca\_os\_hub\_do\_sistema\_fractal; depende gato; depende processos; relaciona word}
\prov{proveniência: broca\_os\_hub.md \textperiodcentered{} web \textperiodcentered{} fato \textperiodcentered{} confiança 1.0 \textperiodcentered{} domínio referencia \textperiodcentered{} história 6 versões no jornal}
\prov{fusão de curadoria (texto idêntico): absorve o\_que\_e --- o mesmo conteúdo por dois esquemas de endereço}

%CRISTAL {"arestas":[{"alvo":"broca_os_hub_broca_os_hub_do_sistema_fractal","confianca":1.0,"epistemico":"fato","origem":"broca_os_hub.md","rel":"parte_de"}],"confianca":1.0,"contraexemplos":[],"descricao":"- **Paper:** `papers/broca_os.tex`, `papers/computacao_fundamentos_broca.tex` - **Manual:** `ula/BROCA_SO.md`, `ula/BROCA_CRISTAL.md` - **Código:** `neuronio/neuronio.c`, `neuronio/radio.py`, `neuronio/primitivas_broca.py` - **Currículo:** fase 6 programação/SO → Broca OS → Tiffany","epistemico":"fato","exemplos":[],"id":"broca_os_hub_referencia","memoria":"semantica","meta":{"dominio":"referencia","fonte":"web","nivel":5},"origem":"broca_os_hub.md","palavras_chave":[],"sinonimos":[],"tipo":"conceito","titulo":"Referência"}
\section{Referência}
- **Paper:** `papers/broca\_os.tex`, `papers/computacao\_fundamentos\_broca.tex` - **Manual:** `ula/BROCA\_SO.md`, `ula/BROCA\_CRISTAL.md` - **Código:** `neuronio/neuronio.c`, `neuronio/radio.py`, `neuronio/primitivas\_broca.py` - **Currículo:** fase 6 programação/SO \ensuremath{\to} Broca OS \ensuremath{\to} Tiffany
\prov{relações: parte\_de broca\_os\_hub\_broca\_os\_hub\_do\_sistema\_fractal}
\prov{proveniência: broca\_os\_hub.md \textperiodcentered{} web \textperiodcentered{} fato \textperiodcentered{} confiança 1.0 \textperiodcentered{} domínio referencia \textperiodcentered{} história 3 versões no jornal}

%CRISTAL {"fusao":[{"arestas":[{"alvo":"broca_os_hub_broca_os_hub_do_sistema_fractal","confianca":1.0,"epistemico":"fato","origem":"broca_os_hub.md","rel":"parte_de"},{"alvo":"gato","confianca":0.95,"epistemico":"fato","origem":"wiki","rel":"usa"},{"alvo":"colapso","confianca":0.95,"epistemico":"fato","origem":"wiki","rel":"referencia"},{"alvo":"word","confianca":0.5,"epistemico":"fato","origem":"wiki","rel":"relaciona"},{"alvo":"vontade_de_poder","confianca":0.5,"epistemico":"fato","origem":"wiki","rel":"relaciona"}],"confianca":1.0,"contraexemplos":[],"descricao":"```text GATO_STREAM → TORO → REGISTRO_N → COLAPSO_6 → BANDA → NONCE ```","epistemico":"fato","exemplos":[],"id":"broca_os_hub_sequencia_pow_nativa","memoria":"semantica","meta":{"dominio":"referencia","fonte":"web","nivel":5},"origem":"broca_os_hub.md","palavras_chave":[],"sinonimos":[],"tipo":"conceito","titulo":"Sequência PoW nativa"},{"arestas":[{"alvo":"gato","confianca":0.95,"epistemico":"fato","origem":"wiki","rel":"usa"},{"alvo":"colapso","confianca":0.95,"epistemico":"fato","origem":"wiki","rel":"referencia"},{"alvo":"vida-morte","confianca":0.5,"epistemico":"fato","origem":"wiki","rel":"relaciona"},{"alvo":"word","confianca":0.5,"epistemico":"fato","origem":"wiki","rel":"relaciona"},{"alvo":"virtualizacao","confianca":0.5,"epistemico":"fato","origem":"wiki","rel":"relaciona"},{"alvo":"telemetria_shadow_producao","confianca":0.5,"epistemico":"fato","origem":"wiki","rel":"relaciona"},{"alvo":"vontade_de_poder","confianca":0.5,"epistemico":"fato","origem":"wiki","rel":"relaciona"}],"confianca":1.0,"contraexemplos":[],"descricao":"```text GATO_STREAM → TORO → REGISTRO_N → COLAPSO_6 → BANDA → NONCE ```","epistemico":"fato","exemplos":[],"id":"sequencia_pow_nativa","memoria":"semantica","meta":{"dominio":"referencia","nivel":5},"origem":"broca_os_hub.md","palavras_chave":[],"sinonimos":[],"tipo":"conceito","titulo":"Sequência PoW nativa"}],"id":"broca_os_hub_sequencia_pow_nativa","tipo":"conceito"}
\section{Sequência PoW nativa}
```text GATO\_STREAM \ensuremath{\to} TORO \ensuremath{\to} REGISTRO\_N \ensuremath{\to} COLAPSO\_6 \ensuremath{\to} BANDA \ensuremath{\to} NONCE ```
\prov{relações: parte\_de broca\_os\_hub\_broca\_os\_hub\_do\_sistema\_fractal; usa gato; referencia colapso; relaciona word; relaciona vontade\_de\_poder}
\prov{proveniência: broca\_os\_hub.md \textperiodcentered{} web \textperiodcentered{} fato \textperiodcentered{} confiança 1.0 \textperiodcentered{} domínio referencia \textperiodcentered{} história 7 versões no jornal}
\prov{fusão de curadoria (texto idêntico): absorve sequencia\_pow\_nativa --- o mesmo conteúdo por dois esquemas de endereço}

%CRISTAL {"arestas":[{"alvo":"broca_os","confianca":1.0,"epistemico":"fato","origem":"manual","rel":"explica"},{"alvo":"goldcoin","confianca":0.5,"epistemico":"fato","origem":"wiki","rel":"relaciona"},{"alvo":"fase3","confianca":0.5,"epistemico":"fato","origem":"wiki","rel":"relaciona"},{"alvo":"gpu","confianca":0.5,"epistemico":"fato","origem":"wiki","rel":"relaciona"},{"alvo":"transcricao","confianca":0.5,"epistemico":"fato","origem":"wiki","rel":"relaciona"},{"alvo":"prova_isa","confianca":0.5,"epistemico":"fato","origem":"wiki","rel":"relaciona"},{"alvo":"teste_ordem","confianca":0.5,"epistemico":"fato","origem":"wiki","rel":"relaciona"},{"alvo":"fechamento_blocos","confianca":0.5,"epistemico":"fato","origem":"wiki","rel":"relaciona"},{"alvo":"python","confianca":0.5,"epistemico":"fato","origem":"wiki","rel":"relaciona"},{"alvo":"contabilidade","confianca":0.5,"epistemico":"fato","origem":"wiki","rel":"relaciona"},{"alvo":"pow_gato","confianca":0.5,"epistemico":"fato","origem":"wiki","rel":"relaciona"},{"alvo":"montagem","confianca":0.5,"epistemico":"fato","origem":"wiki","rel":"relaciona"},{"alvo":"install","confianca":0.5,"epistemico":"fato","origem":"wiki","rel":"relaciona"},{"alvo":"koch_atlas","confianca":0.5,"epistemico":"fato","origem":"wiki","rel":"relaciona"},{"alvo":"cristal_dev","confianca":0.5,"epistemico":"fato","origem":"wiki","rel":"relaciona"},{"alvo":"readme","confianca":0.5,"epistemico":"fato","origem":"wiki","rel":"relaciona"},{"alvo":"koch_pell","confianca":0.5,"epistemico":"fato","origem":"wiki","rel":"relaciona"},{"alvo":"rust","confianca":0.5,"epistemico":"fato","origem":"wiki","rel":"relaciona"},{"alvo":"resposta_a_pergunta_dificil","confianca":0.5,"epistemico":"fato","origem":"wiki","rel":"relaciona"},{"alvo":"sessao_interativa","confianca":0.5,"epistemico":"fato","origem":"wiki","rel":"relaciona"},{"alvo":"go","confianca":0.5,"epistemico":"fato","origem":"wiki","rel":"relaciona"},{"alvo":"mamifero","confianca":0.5,"epistemico":"fato","origem":"wiki","rel":"relaciona"},{"alvo":"koch_memoria","confianca":0.5,"epistemico":"fato","origem":"wiki","rel":"relaciona"},{"alvo":"vida-morte","confianca":0.5,"epistemico":"fato","origem":"wiki","rel":"relaciona"},{"alvo":"pow_espectral","confianca":0.5,"epistemico":"fato","origem":"wiki","rel":"relaciona"},{"alvo":"cifra_prova","confianca":0.5,"epistemico":"fato","origem":"wiki","rel":"relaciona"},{"alvo":"java","confianca":0.5,"epistemico":"fato","origem":"wiki","rel":"relaciona"},{"alvo":"cifra_ouro_branco","confianca":0.5,"epistemico":"fato","origem":"wiki","rel":"relaciona"},{"alvo":"hopfield_rede_hub","confianca":0.5,"epistemico":"fato","origem":"wiki","rel":"relaciona"},{"alvo":"funil_cumulativo","confianca":0.5,"epistemico":"fato","origem":"wiki","rel":"relaciona"},{"alvo":"kernel","confianca":0.5,"epistemico":"fato","origem":"wiki","rel":"relaciona"},{"alvo":"gato_operador_hub","confianca":0.5,"epistemico":"fato","origem":"wiki","rel":"relaciona"},{"alvo":"memoria_virtual","confianca":0.5,"epistemico":"fato","origem":"wiki","rel":"relaciona"},{"alvo":"thread","confianca":0.5,"epistemico":"fato","origem":"wiki","rel":"relaciona"},{"alvo":"koch_parabola","confianca":0.5,"epistemico":"fato","origem":"wiki","rel":"relaciona"},{"alvo":"main","confianca":0.5,"epistemico":"fato","origem":"wiki","rel":"relaciona"},{"alvo":"parabola_isa","confianca":0.5,"epistemico":"fato","origem":"wiki","rel":"relaciona"}],"confianca":1.0,"contraexemplos":[],"descricao":"Manual: ula/BROCA_SO.md","epistemico":"fato","exemplos":[],"id":"broca_so","memoria":"semantica","meta":{"arquivo":"ula/BROCA_SO.md","dominio":"manual","nivel":6},"origem":"manual","palavras_chave":[],"sinonimos":[],"tipo":"conceito","titulo":"BROCA_SO"}
\section{BROCA\_SO}
Manual: ula/BROCA\_SO.md
\prov{relações: explica broca\_os; relaciona goldcoin; relaciona fase3; relaciona gpu; relaciona transcricao; relaciona prova\_isa; relaciona teste\_ordem; relaciona fechamento\_blocos; relaciona python; relaciona contabilidade; relaciona pow\_gato; relaciona montagem; relaciona install; relaciona koch\_atlas; relaciona cristal\_dev; relaciona readme; relaciona koch\_pell; relaciona rust; relaciona resposta\_a\_pergunta\_dificil; relaciona sessao\_interativa; relaciona go; relaciona mamifero; relaciona koch\_memoria; relaciona vida-morte; relaciona pow\_espectral; relaciona cifra\_prova; relaciona java; relaciona cifra\_ouro\_branco; relaciona hopfield\_rede\_hub; relaciona funil\_cumulativo; relaciona kernel; relaciona gato\_operador\_hub; relaciona memoria\_virtual; relaciona thread; relaciona koch\_parabola; relaciona main; relaciona parabola\_isa}
\prov{proveniência: manual \textperiodcentered{} fato \textperiodcentered{} confiança 1.0 \textperiodcentered{} domínio manual \textperiodcentered{} história 67 versões no jornal}

%CRISTAL {"arestas":[{"alvo":"broca-so_camada_de_trabalho_isomorfica","confianca":1.0,"epistemico":"fato","origem":"BROCA_SO.md","rel":"parte_de"}],"confianca":1.0,"contraexemplos":[],"descricao":"```bash cd ula make cifra             # núcleo espectral — silver/gold/bronze make cifra-prova       # prova Pell vs ADD vs modo SPECT make pow-spectral      # PoW — investigação inicial make pow-geometria     # PoW — atlas de coordenadas espectrais make ruido             # auditoria do ruído — Δ = 1−(a+c²) make metais            # família Aₙ make pell              # exploração Pell local make racha             # PROVA: vivo vs cadáver vs Fourier (não produção)","epistemico":"fato","exemplos":[],"id":"build_prova","memoria":"semantica","meta":{"dominio":"manual","duplicata_sugerida":[],"nivel":6,"verificacao":"parte de conceito mais amplo 'build_prova'"},"origem":"README.md","palavras_chave":[],"sinonimos":[],"tipo":"conceito","titulo":"Build / prova"}
\section{Build / prova}
```bash cd ula make cifra             \# núcleo espectral — silver/gold/bronze make cifra-prova       \# prova Pell vs ADD vs modo SPECT make pow-spectral      \# PoW — investigação inicial make pow-geometria     \# PoW — atlas de coordenadas espectrais make ruido             \# auditoria do ruído — \ensuremath{\Delta} = 1\ensuremath{-}(a+c\ensuremath{^{2}}) make metais            \# família A\ensuremath{_{n}} make pell              \# exploração Pell local make racha             \# PROVA: vivo vs cadáver vs Fourier (não produção)
\prov{relações: parte\_de broca-so\_camada\_de\_trabalho\_isomorfica}
\prov{proveniência: README.md \textperiodcentered{} fato \textperiodcentered{} confiança 1.0 \textperiodcentered{} domínio manual \textperiodcentered{} história 8 versões no jornal}

%CRISTAL {"arestas":[{"alvo":"mult_prtsu","confianca":1.0,"epistemico":"fato","origem":"wiki","rel":"usa"},{"alvo":"assinatura_energetica_tarefa","confianca":1.0,"epistemico":"fato","origem":"wiki","rel":"relaciona"}],"confianca":1.0,"contraexemplos":[],"descricao":"A 'velocidade da luz algébrica' local c²_local=a₀b₀/⟨a,b⟩ varia por par de operandos; sua distribuição sobre um dataset é assinatura computável em O(N) sem treino (discriminação KS temporal/espacial 0,997). Define o 'regime de Minkowski algébrico' com c²=p/r.","epistemico":"fato","exemplos":[],"id":"c2_local_dsr","memoria":"semantica","meta":{"autor":"Lopes/Aarão","dominio":"hiper","fonte":""},"origem":"wiki","palavras_chave":[],"sinonimos":[],"tipo":"conceito","titulo":"c²_local (DSR Localizada) como Feature"}
\section{c\ensuremath{^{2}}\_local (DSR Localizada) como Feature}
A 'velocidade da luz algébrica' local c\ensuremath{^{2}}\_local=a\ensuremath{_{0}}b\ensuremath{_{0}}/\ensuremath{\langle}a,b\ensuremath{\rangle} varia por par de operandos; sua distribuição sobre um dataset é assinatura computável em O(N) sem treino (discriminação KS temporal/espacial 0,997). Define o 'regime de Minkowski algébrico' com c\ensuremath{^{2}}=p/r.
\prov{relações: usa mult\_prtsu; relaciona assinatura\_energetica\_tarefa}
\prov{proveniência: wiki \textperiodcentered{} fato \textperiodcentered{} confiança 1.0 \textperiodcentered{} domínio hiper \textperiodcentered{} história 4 versões no jornal}

%CRISTAL {"arestas":[{"alvo":"cadeias_dep_grant","confianca":1.0,"epistemico":"fato","origem":"wiki","rel":"explica"},{"alvo":"colapso_camadas_ordem_parada","confianca":1.0,"epistemico":"fato","origem":"wiki","rel":"usa"}],"confianca":1.0,"contraexemplos":[],"descricao":"Com set_cadeia_ativa, score_hit soma +14 a nós da cadeia ativa e −10 aos de outra; independentemente, se a tag do último nó e a do candidato são da mesma cadeia e o índice de turno é o esperado (nit == uit+1) com o lado esquerdo casando, soma +16 e +22 — a propagação CASL turno a turno. O mecanismo de score por trás do grant-check de cadeias_dep_grant.","epistemico":"fato","exemplos":[],"id":"cadeia_ativa_boost","memoria":"semantica","meta":{"autor":"broca-so","dominio":"conversa","fonte":"conversa/colapso.py:set_cadeia_ativa/score_hit"},"origem":"wiki","palavras_chave":[],"sinonimos":[],"tipo":"conceito","titulo":"Tiffany — boost de cadeia ativa e continuidade de fluxo"}
\section{Tiffany — boost de cadeia ativa e continuidade de fluxo}
Com set\_cadeia\_ativa, score\_hit soma +14 a nós da cadeia ativa e \ensuremath{-}10 aos de outra; independentemente, se a tag do último nó e a do candidato são da mesma cadeia e o índice de turno é o esperado (nit == uit+1) com o lado esquerdo casando, soma +16 e +22 — a propagação CASL turno a turno. O mecanismo de score por trás do grant-check de cadeias\_dep\_grant.
\prov{relações: explica cadeias\_dep\_grant; usa colapso\_camadas\_ordem\_parada}
\prov{proveniência: wiki \textperiodcentered{} conversa/colapso.py:set\_cadeia\_ativa/score\_hit \textperiodcentered{} fato \textperiodcentered{} confiança 1.0 \textperiodcentered{} domínio conversa \textperiodcentered{} história 15 versões no jornal}

%CRISTAL {"arestas":[{"alvo":"cadeias_dep_grant","confianca":1.0,"epistemico":"fato","origem":"wiki","rel":"explica"}],"confianca":1.0,"contraexemplos":[],"descricao":"pares_para_tsv exporta cada cadeia como turnos numerados com dep = total − i + 1: o orçamento δ começa igual ao comprimento da cadeia e decresce a cada turno (δ₀ = nº de turnos → 1 no fim). Cada par gera duas linhas: o turno duplo 'usuário — tiffany' (tag fluxo:cid:NN:dep=) e a face-só (:rNN:dep=). O δ é o orçamento CASL propagando turno a turno.","epistemico":"fato","exemplos":[],"id":"cadeias_delta_decrescente","memoria":"semantica","meta":{"autor":"broca-so","dominio":"conversa","fonte":"conversa/cadeias.py:pares_para_tsv"},"origem":"wiki","palavras_chave":[],"sinonimos":[],"tipo":"conceito","titulo":"Tiffany — orçamento δ decrescente na cadeia"}
\section{Tiffany — orçamento \ensuremath{\delta} decrescente na cadeia}
pares\_para\_tsv exporta cada cadeia como turnos numerados com dep = total \ensuremath{-} i + 1: o orçamento \ensuremath{\delta} começa igual ao comprimento da cadeia e decresce a cada turno (\ensuremath{\delta}\ensuremath{_{0}} = nº de turnos \ensuremath{\to} 1 no fim). Cada par gera duas linhas: o turno duplo 'usuário — tiffany' (tag fluxo:cid:NN:dep=) e a face-só (:rNN:dep=). O \ensuremath{\delta} é o orçamento CASL propagando turno a turno.
\prov{relações: explica cadeias\_dep\_grant}
\prov{proveniência: wiki \textperiodcentered{} conversa/cadeias.py:pares\_para\_tsv \textperiodcentered{} fato \textperiodcentered{} confiança 1.0 \textperiodcentered{} domínio conversa \textperiodcentered{} história 10 versões no jornal}

%CRISTAL {"arestas":[{"alvo":"habitantes_gentil_de_sitter","confianca":1.0,"epistemico":"fato","origem":"DS_UNIVERSO.md","rel":"parte_de"}],"confianca":1.0,"contraexemplos":[],"descricao":"``` neurônio (habitante)     Lopes (universo)        de Sitter (palco) Word [total,e]      →    Lopes3 (r,c,θ), Q      →   (τ, S³, ds²) popcnt par/ímpar         lemniscata Hurwitz         métrica global ```","epistemico":"fato","exemplos":[],"id":"camadas","memoria":"semantica","meta":{"dominio":"manual","nivel":6,"verificacao":"slug idêntico — enriquecer, não duplicar"},"origem":"DS_UNIVERSO.md","palavras_chave":[],"sinonimos":[],"tipo":"conceito","titulo":"Camadas"}
\section{Camadas}
``` neurônio (habitante)     Lopes (universo)        de Sitter (palco) Word [total,e]      \ensuremath{\to}    Lopes3 (r,c,\ensuremath{\theta}), Q      \ensuremath{\to}   (\ensuremath{\tau}, S\ensuremath{^{3}}, ds\ensuremath{^{2}}) popcnt par/ímpar         lemniscata Hurwitz         métrica global ```
\prov{relações: parte\_de habitantes\_gentil\_de\_sitter}
\prov{proveniência: DS\_UNIVERSO.md \textperiodcentered{} fato \textperiodcentered{} confiança 1.0 \textperiodcentered{} domínio manual \textperiodcentered{} história 12 versões no jornal}

%CRISTAL {"arestas":[{"alvo":"fluido_algebrico","confianca":1.0,"epistemico":"fato","origem":"wiki","rel":"usa"},{"alvo":"borda_critica_vazio","confianca":1.0,"epistemico":"fato","origem":"wiki","rel":"relaciona"},{"alvo":"aparelho_psiquico","confianca":1.0,"epistemico":"fato","origem":"wiki","rel":"relaciona"}],"confianca":1.0,"contraexemplos":[],"descricao":"A psicologia como catástrofe de Thom sobre o fluido algébrico ṡ=2s−k(s−s*): as cúspides (modos puros) são os sentimentos que a língua nomeia — o 'osso' (≈uma dúzia: calma, agitação, desejo, tédio, prontidão…); os demais (alegria, medo, amor) são a 'carne': combinações contínuas de 3–4 modos. Poucas cúspides, infinitos sentimentos. (Refs: Plutchik, Russell.)","epistemico":"hipotese","exemplos":[],"id":"campo_psicologico_cuspide","memoria":"semantica","meta":{"autor":"Lopes/Aarão","dominio":"hiper","fonte":""},"origem":"wiki","palavras_chave":[],"sinonimos":[],"tipo":"conceito","titulo":"O Campo Psicológico como Catástrofe de Cúspide"}
\section{O Campo Psicológico como Catástrofe de Cúspide}
A psicologia como catástrofe de Thom sobre o fluido algébrico \ensuremath{\dot{s}}=2s\ensuremath{-}k(s\ensuremath{-}s*): as cúspides (modos puros) são os sentimentos que a língua nomeia — o 'osso' (\ensuremath{\approx}uma dúzia: calma, agitação, desejo, tédio, prontidão…); os demais (alegria, medo, amor) são a 'carne': combinações contínuas de 3–4 modos. Poucas cúspides, infinitos sentimentos. (Refs: Plutchik, Russell.)
\prov{relações: usa fluido\_algebrico; relaciona borda\_critica\_vazio; relaciona aparelho\_psiquico}
\prov{proveniência: wiki \textperiodcentered{} hipotese \textperiodcentered{} confiança 1.0 \textperiodcentered{} domínio hiper \textperiodcentered{} história 5 versões no jornal}

%CRISTAL {"arestas":[{"alvo":"destruir_a_parabola","confianca":1.0,"epistemico":"fato","origem":"PARABOLA_DESTRUIR.md","rel":"parte_de"},{"alvo":"word","confianca":0.5,"epistemico":"fato","origem":"wiki","rel":"relaciona"},{"alvo":"syscall","confianca":0.5,"epistemico":"fato","origem":"wiki","rel":"relaciona"},{"alvo":"vontade_de_poder","confianca":0.5,"epistemico":"fato","origem":"wiki","rel":"relaciona"},{"alvo":"tragedia_dos_comuns","confianca":0.5,"epistemico":"fato","origem":"wiki","rel":"relaciona"},{"alvo":"virtualizacao","confianca":0.5,"epistemico":"fato","origem":"wiki","rel":"relaciona"}],"confianca":1.0,"contraexemplos":[],"descricao":"- passos sem distribuição (total+e ≤ 0): **1,016,900** - destes, precisariam de constante 0.5/0.5 em `parabola_entre`: **1,016,900**","epistemico":"fato","exemplos":[],"id":"casos_degenerados","memoria":"semantica","meta":{"dominio":"manual","nivel":6,"verificacao":"slug idêntico — enriquecer, não duplicar"},"origem":"PARABOLA_DESTRUIR.md","palavras_chave":[],"sinonimos":[],"tipo":"conceito","titulo":"Casos degenerados"}
\section{Casos degenerados}
- passos sem distribuição (total+e \ensuremath{\le} 0): **1,016,900** - destes, precisariam de constante 0.5/0.5 em `parabola\_entre`: **1,016,900**
\prov{relações: parte\_de destruir\_a\_parabola; relaciona word; relaciona syscall; relaciona vontade\_de\_poder; relaciona tragedia\_dos\_comuns; relaciona virtualizacao}
\prov{proveniência: PARABOLA\_DESTRUIR.md \textperiodcentered{} fato \textperiodcentered{} confiança 1.0 \textperiodcentered{} domínio manual \textperiodcentered{} história 9 versões no jornal}

%CRISTAL {"arestas":[{"alvo":"transformada_dourada","confianca":1.0,"epistemico":"fato","origem":"wiki","rel":"usa"},{"alvo":"decomposicao_peirce","confianca":1.0,"epistemico":"fato","origem":"wiki","rel":"relaciona"}],"confianca":1.0,"contraexemplos":[],"descricao":"Cada obra é evaporada pela transformada dourada e mede-se a entropia espectral H_AC e os 'modos efetivos' (participation ratio = quantos campos locais ℂ_â a obra ocupa). H=0 = pureza máxima (um campo local); H=1 = dispersão máxima. Ordena do mais puro (esboço melódico ~0,44) ao menos puro (gravação encorpada ~0,89) — a pureza é do recorte, não do tema.","epistemico":"fato","exemplos":[],"id":"catalogo_especies_pureza","memoria":"semantica","meta":{"autor":"Lopes/Aarão","dominio":"hiper","fonte":""},"origem":"wiki","palavras_chave":[],"sinonimos":[],"tipo":"conceito","titulo":"Classificação por Pureza do Campo Local"}
\section{Classificação por Pureza do Campo Local}
Cada obra é evaporada pela transformada dourada e mede-se a entropia espectral H\_AC e os 'modos efetivos' (participation ratio = quantos campos locais \ensuremath{\mathbb{C}}\_â a obra ocupa). H=0 = pureza máxima (um campo local); H=1 = dispersão máxima. Ordena do mais puro (esboço melódico \~{}0,44) ao menos puro (gravação encorpada \~{}0,89) — a pureza é do recorte, não do tema.
\prov{relações: usa transformada\_dourada; relaciona decomposicao\_peirce}
\prov{proveniência: wiki \textperiodcentered{} fato \textperiodcentered{} confiança 1.0 \textperiodcentered{} domínio hiper \textperiodcentered{} história 4 versões no jornal}

%CRISTAL {"arestas":[{"alvo":"protocolo_assincrono","confianca":1.0,"epistemico":"fato","origem":"wiki","rel":"usa"},{"alvo":"cifra_prova","confianca":1.0,"epistemico":"fato","origem":"wiki","rel":"relaciona"}],"confianca":1.0,"contraexemplos":[],"descricao":"cert=(t,seq,prev,nullifier,event_hash,band,sig) com o tempo nos primeiros bits do cabeçalho (a torção o absorve como semente, ordenando sem tocar a identidade — 'o RG, não a alma'); o nullifier veda replay; sem deadline — o tempo ordena, não expira. Máquina de estados PRELOCKED→CERTIFIED→ORDERED→EVIDENCED→RELEASED.","epistemico":"fato","exemplos":[],"id":"certificado_tempo_semente","memoria":"semantica","meta":{"autor":"Aarão/broca-so","dominio":"broca_repos","fonte":""},"origem":"wiki","palavras_chave":[],"sinonimos":[],"tipo":"conceito","titulo":"O Certificado Encadeado (o tempo como semente)"}
\section{O Certificado Encadeado (o tempo como semente)}
cert=(t,seq,prev,nullifier,event\_hash,band,sig) com o tempo nos primeiros bits do cabeçalho (a torção o absorve como semente, ordenando sem tocar a identidade — 'o RG, não a alma'); o nullifier veda replay; sem deadline — o tempo ordena, não expira. Máquina de estados PRELOCKED\ensuremath{\to}CERTIFIED\ensuremath{\to}ORDERED\ensuremath{\to}EVIDENCED\ensuremath{\to}RELEASED.
\prov{relações: usa protocolo\_assincrono; relaciona cifra\_prova}
\prov{proveniência: wiki \textperiodcentered{} fato \textperiodcentered{} confiança 1.0 \textperiodcentered{} domínio broca\_repos \textperiodcentered{} história 3 versões no jornal}

%CRISTAL {"arestas":[{"alvo":"colapso_campo_local","confianca":1.0,"epistemico":"fato","origem":"wiki","rel":"usa"},{"alvo":"razao_aurea","confianca":1.0,"epistemico":"fato","origem":"wiki","rel":"relaciona"},{"alvo":"maquina_evolutiva","confianca":1.0,"epistemico":"fato","origem":"wiki","rel":"relaciona"}],"confianca":1.0,"contraexemplos":[],"descricao":"'Chá' = o associador entre obras de campos locais distintos (código×música×imagem deu o maior chá): o novo só nasce cruzando naturezas independentes. A saturação medida (≈0,20) é da régua binária; a auto-similaridade real é áurea (Ω=φ¹/φ) — tratada como bússola, não teorema.","epistemico":"hipotese","exemplos":[],"id":"cha_bussola_aurea","memoria":"semantica","meta":{"autor":"Lopes/Aarão","dominio":"hiper","fonte":""},"origem":"wiki","palavras_chave":[],"sinonimos":[],"tipo":"conceito","titulo":"O 'Chá' e a Bússola Áurea"}
\section{O 'Chá' e a Bússola Áurea}
'Chá' = o associador entre obras de campos locais distintos (código$\times$música$\times$imagem deu o maior chá): o novo só nasce cruzando naturezas independentes. A saturação medida (\ensuremath{\approx}0,20) é da régua binária; a auto-similaridade real é áurea (\ensuremath{\Omega}=\ensuremath{\varphi}\ensuremath{^{1}}/\ensuremath{\varphi}) — tratada como bússola, não teorema.
\prov{relações: usa colapso\_campo\_local; relaciona razao\_aurea; relaciona maquina\_evolutiva}
\prov{proveniência: wiki \textperiodcentered{} hipotese \textperiodcentered{} confiança 1.0 \textperiodcentered{} domínio hiper \textperiodcentered{} história 6 versões no jornal}

%CRISTAL {"arestas":[{"alvo":"peirce_celulas_canais","confianca":1.0,"epistemico":"fato","origem":"wiki","rel":"usa"},{"alvo":"soberania_da_chave","confianca":1.0,"epistemico":"fato","origem":"wiki","rel":"relaciona"}],"confianca":1.0,"contraexemplos":[],"descricao":"A posse coletiva (o coração) é a completude Σ E_ii = I (o corpo todo): não é produzida, emerge da soma; cada tecido sozinho é parcial, só todos juntos fecham I. Ninguém é privilegiado — nem o autor — e todas as chaves se cancelam no vazio (eixos de Bloch via SHA somam ~0). Threshold FROST/Schnorr.","epistemico":"fato","exemplos":[],"id":"chave_do_coracao","memoria":"semantica","meta":{"autor":"Aarão/broca-so","dominio":"broca_repos","fonte":""},"origem":"wiki","palavras_chave":[],"sinonimos":[],"tipo":"conceito","titulo":"A Chave do Coração (ninguém a produz)"}
\section{A Chave do Coração (ninguém a produz)}
A posse coletiva (o coração) é a completude \ensuremath{\Sigma} E\_ii = I (o corpo todo): não é produzida, emerge da soma; cada tecido sozinho é parcial, só todos juntos fecham I. Ninguém é privilegiado — nem o autor — e todas as chaves se cancelam no vazio (eixos de Bloch via SHA somam \~{}0). Threshold FROST/Schnorr.
\prov{relações: usa peirce\_celulas\_canais; relaciona soberania\_da\_chave}
\prov{proveniência: wiki \textperiodcentered{} fato \textperiodcentered{} confiança 1.0 \textperiodcentered{} domínio broca\_repos \textperiodcentered{} história 3 versões no jornal}

%CRISTAL {"arestas":[{"alvo":"broca-so_cristalizado","confianca":1.0,"epistemico":"fato","origem":"BROCA_CRISTAL.md","rel":"parte_de"}],"confianca":1.0,"contraexemplos":[],"descricao":"- [ ] Núcleo: estado `Word`, `e` preservado até Koch. - [ ] Saída humana: passa por projeção Koch (ou equivalente associativo documentado). - [ ] Borda clássica: só recebe **projectado**, não vivo bruto amputado. - [ ] Nova peça orbita `neuronio.c`?","epistemico":"fato","exemplos":[],"id":"checklist_pr","memoria":"semantica","meta":{"dominio":"manual","nivel":6,"verificacao":"slug idêntico — enriquecer, não duplicar"},"origem":"BROCA_CRISTAL.md","palavras_chave":[],"sinonimos":[],"tipo":"conceito","titulo":"Checklist PR"}
\section{Checklist PR}
- [ ] Núcleo: estado `Word`, `e` preservado até Koch. - [ ] Saída humana: passa por projeção Koch (ou equivalente associativo documentado). - [ ] Borda clássica: só recebe **projectado**, não vivo bruto amputado. - [ ] Nova peça orbita `neuronio.c`?
\prov{relações: parte\_de broca-so\_cristalizado}
\prov{proveniência: BROCA\_CRISTAL.md \textperiodcentered{} fato \textperiodcentered{} confiança 1.0 \textperiodcentered{} domínio manual \textperiodcentered{} história 4 versões no jornal}

%CRISTAL {"arestas":[{"alvo":"memoria_viva","confianca":1.0,"epistemico":"fato","origem":"paper","rel":"parte_de"}],"confianca":1.0,"contraexemplos":[],"descricao":"- `mem/NNNN.bin` com `N < 8000`: **dados** — `LOAD` passa pelo neurônio. - `N ≥ 8000`: **scratch** — roundtrip raw (`8000`=acc, `8003`=report, `8008`=diff, …). - Convenção de memória, **não opcode**.","epistemico":"fato","exemplos":[],"id":"ciclo_memoria_conhecimento","memoria":"semantica","meta":{"dominio":"manual","duplicata_sugerida":[],"nivel":6,"verificacao":"parte de conceito mais amplo 'ciclo_memoria_conhecimento'"},"origem":"PROVA_ISA.md","palavras_chave":[],"sinonimos":[],"tipo":"conceito","titulo":"Ciclo memória conhecimento"}
\section{Ciclo memória conhecimento}
- `mem/NNNN.bin` com `N < 8000`: **dados** — `LOAD` passa pelo neurônio. - `N \ensuremath{\ge} 8000`: **scratch** — roundtrip raw (`8000`=acc, `8003`=report, `8008`=diff, …). - Convenção de memória, **não opcode**.
\prov{relações: parte\_de memoria\_viva}
\prov{proveniência: PROVA\_ISA.md \textperiodcentered{} fato \textperiodcentered{} confiança 1.0 \textperiodcentered{} domínio manual \textperiodcentered{} história 5 versões no jornal}

%CRISTAL {"arestas":[{"alvo":"sha256","confianca":1.0,"epistemico":"fato","origem":"paper","rel":"especializa"},{"alvo":"cifra_de_ouro_branco","confianca":1.0,"epistemico":"fato","origem":"manual","rel":"explica"}],"confianca":1.0,"contraexemplos":[],"descricao":"- **Rastreio Pell** → camada **`silver(w)` / célula branca** (Word exacto, sem limite mod p). - **Floco Koch** → só guarda **resíduo mod 1009** da dissipação; cobre **todos** os Pell mod p, mas **não** o inteiro grande. - Pell ≥ 2378 **existem fora** do floco como valor exacto — a cifra **não** deve esperar o inteiro no Koch; rastreia na prata (`A₂`) e usa o floco só como complemento mod p.","epistemico":"fato","exemplos":[],"id":"cifra","memoria":"semantica","meta":{"dominio":"manual","duplicata_sugerida":[],"nivel":6,"verificacao":"parte de conceito mais amplo 'cifra'"},"origem":"KOCH_PELL.md","palavras_chave":[],"sinonimos":[],"tipo":"conceito","titulo":"CIFRA"}
\section{CIFRA}
- **Rastreio Pell** \ensuremath{\to} camada **`silver(w)` / célula branca** (Word exacto, sem limite mod p). - **Floco Koch** \ensuremath{\to} só guarda **resíduo mod 1009** da dissipação; cobre **todos** os Pell mod p, mas **não** o inteiro grande. - Pell \ensuremath{\ge} 2378 **existem fora** do floco como valor exacto — a cifra **não** deve esperar o inteiro no Koch; rastreia na prata (`A\ensuremath{_{2}}`) e usa o floco só como complemento mod p.
\prov{relações: especializa sha256; explica cifra\_de\_ouro\_branco}
\prov{proveniência: KOCH\_PELL.md \textperiodcentered{} fato \textperiodcentered{} confiança 1.0 \textperiodcentered{} domínio manual \textperiodcentered{} história 41 versões no jornal}

%CRISTAL {"arestas":[{"alvo":"corpo_f2k","confianca":1.0,"epistemico":"fato","origem":"wiki","rel":"usa"},{"alvo":"soberania_da_chave","confianca":1.0,"epistemico":"fato","origem":"wiki","rel":"explica"},{"alvo":"cifra","confianca":1.0,"epistemico":"fato","origem":"wiki","rel":"especializa"}],"confianca":1.0,"contraexemplos":[],"descricao":"A chave de ouro do cliente é um g∈𝔽₂ᵏˣ; multiplicar por g é linear sobre 𝔽₂ (mudança de base). Cifra=g·m; só o dono aplica g⁻¹=g²ᵏ−² (Fermat) e recupera um primo de prata — non-custodial por construção, reversão exata em bits (verificado em 𝔽₂¹²⁸).","epistemico":"fato","exemplos":[],"id":"cifra_mudanca_de_base","memoria":"semantica","meta":{"autor":"Aarão/broca-so","dominio":"broca_repos","fonte":""},"origem":"wiki","palavras_chave":[],"sinonimos":[],"tipo":"conceito","titulo":"Cifra por Mudança de Base (𝔽₂ᵏ)"}
\section{Cifra por Mudança de Base (\ensuremath{\mathbb{F}}\ensuremath{_{2}}\ensuremath{^{k}})}
A chave de ouro do cliente é um g\ensuremath{\in}\ensuremath{\mathbb{F}}\ensuremath{_{2}}\ensuremath{^{kx}}; multiplicar por g é linear sobre \ensuremath{\mathbb{F}}\ensuremath{_{2}} (mudança de base). Cifra=g\textperiodcentered m; só o dono aplica g\ensuremath{^{-1}}=g\ensuremath{^{2k}}\ensuremath{-}\ensuremath{^{2}} (Fermat) e recupera um primo de prata — non-custodial por construção, reversão exata em bits (verificado em \ensuremath{\mathbb{F}}\ensuremath{_{2}}\ensuremath{^{128}}).
\prov{relações: usa corpo\_f2k; explica soberania\_da\_chave; especializa cifra}
\prov{proveniência: wiki \textperiodcentered{} fato \textperiodcentered{} confiança 1.0 \textperiodcentered{} domínio broca\_repos \textperiodcentered{} história 5 versões no jornal}

%CRISTAL {"arestas":[{"alvo":"cifra","confianca":1.0,"epistemico":"fato","origem":"manual","rel":"parte_de"},{"alvo":"utilitarismo","confianca":0.5,"epistemico":"fato","origem":"wiki","rel":"relaciona"},{"alvo":"resultado_anterior_identidade_nao_cobertura","confianca":0.5,"epistemico":"fato","origem":"wiki","rel":"relaciona"},{"alvo":"resultado_completo","confianca":0.5,"epistemico":"fato","origem":"wiki","rel":"relaciona"},{"alvo":"destruir_a_parabola","confianca":0.5,"epistemico":"fato","origem":"wiki","rel":"relaciona"},{"alvo":"teoria_do_valor","confianca":0.5,"epistemico":"fato","origem":"wiki","rel":"relaciona"},{"alvo":"pitagorismo","confianca":0.5,"epistemico":"fato","origem":"wiki","rel":"relaciona"},{"alvo":"scheduler","confianca":0.5,"epistemico":"fato","origem":"wiki","rel":"relaciona"},{"alvo":"toryces_e_helyces","confianca":0.5,"epistemico":"fato","origem":"wiki","rel":"relaciona"},{"alvo":"fibonacci_multidimensional","confianca":0.5,"epistemico":"fato","origem":"wiki","rel":"relaciona"},{"alvo":"o_que_e_no_projecto","confianca":0.5,"epistemico":"fato","origem":"wiki","rel":"relaciona"},{"alvo":"filesystem","confianca":0.5,"epistemico":"fato","origem":"wiki","rel":"relaciona"},{"alvo":"paper_2_reta_ouro","confianca":0.5,"epistemico":"fato","origem":"wiki","rel":"relaciona"},{"alvo":"ficheiros","confianca":0.5,"epistemico":"fato","origem":"wiki","rel":"relaciona"}],"confianca":1.0,"contraexemplos":[],"descricao":"> A cifra deixa de ser software e passa a ser **física interna** da máquina.","epistemico":"fato","exemplos":[],"id":"cifra_ouro_branco_nucleo_espectral","memoria":"semantica","meta":{"dominio":"manual","nivel":6,"verificacao":"slug idêntico — enriquecer, não duplicar"},"origem":"CIFRA.md","palavras_chave":[],"sinonimos":[],"tipo":"conceito","titulo":"Cifra Ouro Branco — núcleo espectral"}
\section{Cifra Ouro Branco — núcleo espectral}
> A cifra deixa de ser software e passa a ser **física interna** da máquina.
\prov{relações: parte\_de cifra; relaciona utilitarismo; relaciona resultado\_anterior\_identidade\_nao\_cobertura; relaciona resultado\_completo; relaciona destruir\_a\_parabola; relaciona teoria\_do\_valor; relaciona pitagorismo; relaciona scheduler; relaciona toryces\_e\_helyces; relaciona fibonacci\_multidimensional; relaciona o\_que\_e\_no\_projecto; relaciona filesystem; relaciona paper\_2\_reta\_ouro; relaciona ficheiros}
\prov{proveniência: CIFRA.md \textperiodcentered{} fato \textperiodcentered{} confiança 1.0 \textperiodcentered{} domínio manual \textperiodcentered{} história 18 versões no jornal}

%CRISTAL {"arestas":[{"alvo":"cifra_de_ouro_branco","confianca":1.0,"epistemico":"fato","origem":"manual","rel":"prova"},{"alvo":"goldcoin","confianca":0.5,"epistemico":"fato","origem":"wiki","rel":"relaciona"},{"alvo":"install","confianca":0.5,"epistemico":"fato","origem":"wiki","rel":"relaciona"},{"alvo":"sessao_interativa","confianca":0.5,"epistemico":"fato","origem":"wiki","rel":"relaciona"},{"alvo":"cristal_dev","confianca":0.5,"epistemico":"fato","origem":"wiki","rel":"relaciona"},{"alvo":"gpu","confianca":0.5,"epistemico":"fato","origem":"wiki","rel":"relaciona"},{"alvo":"go","confianca":0.5,"epistemico":"fato","origem":"wiki","rel":"relaciona"},{"alvo":"thread","confianca":0.5,"epistemico":"fato","origem":"wiki","rel":"relaciona"},{"alvo":"kernel","confianca":0.5,"epistemico":"fato","origem":"wiki","rel":"relaciona"},{"alvo":"grava_fato_resposta_normal_em_seguida","confianca":0.5,"epistemico":"fato","origem":"wiki","rel":"relaciona"},{"alvo":"koch_memoria","confianca":0.5,"epistemico":"fato","origem":"wiki","rel":"relaciona"},{"alvo":"python","confianca":0.5,"epistemico":"fato","origem":"wiki","rel":"relaciona"},{"alvo":"montagem","confianca":0.5,"epistemico":"fato","origem":"wiki","rel":"relaciona"},{"alvo":"teste_ordem","confianca":0.5,"epistemico":"fato","origem":"wiki","rel":"relaciona"},{"alvo":"ds_universo","confianca":0.5,"epistemico":"fato","origem":"wiki","rel":"relaciona"},{"alvo":"gato_operador_hub","confianca":0.5,"epistemico":"fato","origem":"wiki","rel":"relaciona"},{"alvo":"prova_isa","confianca":0.5,"epistemico":"fato","origem":"wiki","rel":"relaciona"},{"alvo":"mamifero","confianca":0.5,"epistemico":"fato","origem":"wiki","rel":"relaciona"},{"alvo":"vida-morte","confianca":0.5,"epistemico":"fato","origem":"wiki","rel":"relaciona"},{"alvo":"numa","confianca":0.5,"epistemico":"fato","origem":"wiki","rel":"relaciona"},{"alvo":"metais","confianca":0.5,"epistemico":"fato","origem":"wiki","rel":"relaciona"},{"alvo":"main","confianca":0.5,"epistemico":"fato","origem":"wiki","rel":"relaciona"}],"confianca":1.0,"contraexemplos":[],"descricao":"Manual: ula/CIFRA_PROVA.md","epistemico":"fato","exemplos":[],"id":"cifra_prova","memoria":"semantica","meta":{"arquivo":"ula/CIFRA_PROVA.md","dominio":"manual","nivel":6},"origem":"manual","palavras_chave":[],"sinonimos":[],"tipo":"conceito","titulo":"CIFRA_PROVA"}
\section{CIFRA\_PROVA}
Manual: ula/CIFRA\_PROVA.md
\prov{relações: prova cifra\_de\_ouro\_branco; relaciona goldcoin; relaciona install; relaciona sessao\_interativa; relaciona cristal\_dev; relaciona gpu; relaciona go; relaciona thread; relaciona kernel; relaciona grava\_fato\_resposta\_normal\_em\_seguida; relaciona koch\_memoria; relaciona python; relaciona montagem; relaciona teste\_ordem; relaciona ds\_universo; relaciona gato\_operador\_hub; relaciona prova\_isa; relaciona mamifero; relaciona vida-morte; relaciona numa; relaciona metais; relaciona main}
\prov{proveniência: manual \textperiodcentered{} fato \textperiodcentered{} confiança 1.0 \textperiodcentered{} domínio manual \textperiodcentered{} história 25 versões no jornal}

%CRISTAL {"arestas":[{"alvo":"vida_matematica","confianca":1.0,"epistemico":"fato","origem":"wiki","rel":"usa"},{"alvo":"yin_yang_quantico","confianca":1.0,"epistemico":"fato","origem":"wiki","rel":"relaciona"}],"confianca":1.0,"contraexemplos":[],"descricao":"A Ciranda cuida do ciclo de vida do tecido: não mata — põe o tecido na roda, ele cumpre a volta e sai, e seu 'chá' é recolhido ao maracá sem contaminar vizinhos. Núcleo: do gap caem dois modos (φ,ψ), fator de decaimento 1/φ² e a constante φ; vida=modo expansivo, morte=modo contrátil, maracá=absorvente.","epistemico":"inferencia","exemplos":[],"id":"ciranda_ciclo_de_vida","memoria":"semantica","meta":{"autor":"Aarão/broca-so","dominio":"broca_repos","fonte":""},"origem":"wiki","palavras_chave":[],"sinonimos":[],"tipo":"conceito","titulo":"A Ciranda (a roda que não mata)"}
\section{A Ciranda (a roda que não mata)}
A Ciranda cuida do ciclo de vida do tecido: não mata — põe o tecido na roda, ele cumpre a volta e sai, e seu 'chá' é recolhido ao maracá sem contaminar vizinhos. Núcleo: do gap caem dois modos (\ensuremath{\varphi},\ensuremath{\psi}), fator de decaimento 1/\ensuremath{\varphi}\ensuremath{^{2}} e a constante \ensuremath{\varphi}; vida=modo expansivo, morte=modo contrátil, maracá=absorvente.
\prov{relações: usa vida\_matematica; relaciona yin\_yang\_quantico}
\prov{proveniência: wiki \textperiodcentered{} inferencia \textperiodcentered{} confiança 1.0 \textperiodcentered{} domínio broca\_repos \textperiodcentered{} história 3 versões no jornal}

%CRISTAL {"arestas":[{"alvo":"flip_epsilon_discriminante","confianca":1.0,"epistemico":"fato","origem":"wiki","rel":"usa"},{"alvo":"arquitetura_von_neumann","confianca":1.0,"epistemico":"fato","origem":"wiki","rel":"relaciona"}],"confianca":1.0,"contraexemplos":[],"descricao":"Um circuito linear de 2ª ordem (RLC) É um ensaio de torção: ε=sign(−disc) da companion é o regime de amortecimento; a impedância Z=R+jX separa o que dissipa (R, split) do que gira e conserva (X, fase). E o gap de ouro tem R<0: não dissipa, injeta (oscilador ativo) — duas paredes: disc=0 (a forma) e tr A=0 (a atividade, Barkhausen).","epistemico":"fato","exemplos":[],"id":"circuito_como_torcao","memoria":"semantica","meta":{"autor":"Aarão/broca-so","dominio":"broca_repos","fonte":""},"origem":"wiki","palavras_chave":[],"sinonimos":[],"tipo":"conceito","titulo":"O Circuito RLC é o Flip ε (não analogia)"}
\section{O Circuito RLC é o Flip \ensuremath{\varepsilon} (não analogia)}
Um circuito linear de 2ª ordem (RLC) É um ensaio de torção: \ensuremath{\varepsilon}=sign(\ensuremath{-}disc) da companion é o regime de amortecimento; a impedância Z=R+jX separa o que dissipa (R, split) do que gira e conserva (X, fase). E o gap de ouro tem R<0: não dissipa, injeta (oscilador ativo) — duas paredes: disc=0 (a forma) e tr A=0 (a atividade, Barkhausen).
\prov{relações: usa flip\_epsilon\_discriminante; relaciona arquitetura\_von\_neumann}
\prov{proveniência: wiki \textperiodcentered{} fato \textperiodcentered{} confiança 1.0 \textperiodcentered{} domínio broca\_repos \textperiodcentered{} história 3 versões no jornal}

%CRISTAL {"arestas":[{"alvo":"a_cacada_da_parabola_mapa_vivo_da_isa","confianca":1.0,"epistemico":"fato","origem":"PARABOLA_ISA.md","rel":"parte_de"}],"confianca":1.0,"contraexemplos":[],"descricao":"| Classe | Significado | |--------|-------------| | 🌿 vivo | ressoa, alta fase, reversível | | 🌉 fronteira | entrada/saída, neurônio, decisão | | 🪨 cristal | calor, cicatriz, persistência |","epistemico":"fato","exemplos":[],"id":"classificacao","memoria":"semantica","meta":{"dominio":"manual","nivel":6,"verificacao":"slug idêntico — enriquecer, não duplicar"},"origem":"PARABOLA_ISA.md","palavras_chave":[],"sinonimos":[],"tipo":"conceito","titulo":"Classificação"}
\section{Classificação}
| Classe | Significado | |--------|-------------| | [ramo] vivo | ressoa, alta fase, reversível | | [ponte] fronteira | entrada/saída, neurônio, decisão | | [pedra] cristal | calor, cicatriz, persistência |
\prov{relações: parte\_de a\_cacada\_da\_parabola\_mapa\_vivo\_da\_isa}
\prov{proveniência: PARABOLA\_ISA.md \textperiodcentered{} fato \textperiodcentered{} confiança 1.0 \textperiodcentered{} domínio manual \textperiodcentered{} história 6 versões no jornal}

%CRISTAL {"arestas":[{"alvo":"broca_os","confianca":0.95,"epistemico":"fato","origem":"paper","rel":"parte_de"},{"alvo":"recursao","confianca":0.5,"epistemico":"fato","origem":"wiki","rel":"relaciona"},{"alvo":"utilitarismo","confianca":0.5,"epistemico":"fato","origem":"wiki","rel":"relaciona"},{"alvo":"word","confianca":0.5,"epistemico":"fato","origem":"wiki","rel":"relaciona"},{"alvo":"virtualizacao","confianca":0.5,"epistemico":"fato","origem":"wiki","rel":"relaciona"}],"confianca":1.0,"contraexemplos":[],"descricao":"c = tiffany.Client(home=\"~/.tiffany\") print(c.ask(\"python venv pip install\"))","epistemico":"fato","exemplos":[],"id":"cliente_de_alto_nivel","memoria":"semantica","meta":{"dominio":"manual","nivel":6,"verificacao":"slug idêntico — enriquecer, não duplicar"},"origem":"INSTALL.md","palavras_chave":[],"sinonimos":[],"tipo":"conceito","titulo":"cliente de alto nível"}
\section{cliente de alto nível}
c = tiffany.Client(home="\~{}/.tiffany") print(c.ask("python venv pip install"))
\prov{relações: parte\_de broca\_os; relaciona recursao; relaciona utilitarismo; relaciona word; relaciona virtualizacao}
\prov{proveniência: INSTALL.md \textperiodcentered{} fato \textperiodcentered{} confiança 1.0 \textperiodcentered{} domínio manual \textperiodcentered{} história 7 versões no jornal}

%CRISTAL {"arestas":[{"alvo":"identidade_omega_r8","confianca":1.0,"epistemico":"fato","origem":"wiki","rel":"usa"},{"alvo":"hopfield","confianca":1.0,"epistemico":"fato","origem":"wiki","rel":"contradiz"},{"alvo":"mente_computacao","confianca":1.0,"epistemico":"fato","origem":"wiki","rel":"relaciona"},{"alvo":"cristal_cli","confianca":1.0,"epistemico":"fato","origem":"wiki","rel":"relaciona"}],"confianca":1.0,"contraexemplos":[],"descricao":"Arquitetura cujo estado é uma 4-forma φM∈Λ⁴(ℝ⁸) (70 g.l./camada), escrita por antissimetrização k⊗v⊗q⊗r modulada por gate e leitura via 2-forma; a leitura é Spin(7)-equivariante (norma invariante em 21 dos 28 g.l.) — vantagem estrutural que KV-cache/Hopfield/Titans não têm. O operador de esquecimento CΩ tem espectro 1,½,−⅙,0 dando quatro horizontes temporais determinados geometricamente, sem decay rates aprendidos.","epistemico":"inferencia","exemplos":[],"id":"cltm_memoria_spin7","memoria":"semantica","meta":{"autor":"Aarão/broca-so","dominio":"broca_repos","fonte":""},"origem":"wiki","palavras_chave":[],"sinonimos":[],"tipo":"conceito","titulo":"CLTM — Memória de Longo Prazo Spin(7)-Equivariante"}
\section{CLTM — Memória de Longo Prazo Spin(7)-Equivariante}
Arquitetura cujo estado é uma 4-forma \ensuremath{\varphi}M\ensuremath{\in}\ensuremath{\Lambda}\ensuremath{^{4}}(\ensuremath{\mathbb{R}}\ensuremath{^{8}}) (70 g.l./camada), escrita por antissimetrização k\ensuremath{\otimes}v\ensuremath{\otimes}q\ensuremath{\otimes}r modulada por gate e leitura via 2-forma; a leitura é Spin(7)-equivariante (norma invariante em 21 dos 28 g.l.) — vantagem estrutural que KV-cache/Hopfield/Titans não têm. O operador de esquecimento C\ensuremath{\Omega} tem espectro 1,\ensuremath{\tfrac12},\ensuremath{-}\ensuremath{\tfrac16},0 dando quatro horizontes temporais determinados geometricamente, sem decay rates aprendidos.
\prov{relações: usa identidade\_omega\_r8; contradiz hopfield; relaciona mente\_computacao; relaciona cristal\_cli}
\prov{proveniência: wiki \textperiodcentered{} inferencia \textperiodcentered{} confiança 1.0 \textperiodcentered{} domínio broca\_repos \textperiodcentered{} história 6 versões no jornal}

%CRISTAL {"arestas":[{"alvo":"olho_de_koch","confianca":1.0,"epistemico":"fato","origem":"wiki","rel":"relaciona"},{"alvo":"colapso_koch_tiffany","confianca":1.0,"epistemico":"fato","origem":"wiki","rel":"relaciona"}],"confianca":1.0,"contraexemplos":[],"descricao":"Os zeros de ζ não vivem no espectro discreto 3D (formas de Maass) mas no espalhamento de Eisenstein das cúspides, um canal efetivamente 1D; a fase de espalhamento reproduz termo a termo a contagem dos zeros (T=30,100,200 → 3,6/29,0/79,2 vs 3/29/79). A localização em Re s=½ fica como horizonte declarado.","epistemico":"inferencia","exemplos":[],"id":"colapso_3_para_1_zeta","memoria":"semantica","meta":{"autor":"Aarão/broca-so","dominio":"broca_repos","fonte":""},"origem":"wiki","palavras_chave":[],"sinonimos":[],"tipo":"conceito","titulo":"O Colapso 3→1 (os zeros de ζ nas cúspides)"}
\section{O Colapso 3\ensuremath{\to}1 (os zeros de \ensuremath{\zeta} nas cúspides)}
Os zeros de \ensuremath{\zeta} não vivem no espectro discreto 3D (formas de Maass) mas no espalhamento de Eisenstein das cúspides, um canal efetivamente 1D; a fase de espalhamento reproduz termo a termo a contagem dos zeros (T=30,100,200 \ensuremath{\to} 3,6/29,0/79,2 vs 3/29/79). A localização em Re s=\ensuremath{\tfrac12} fica como horizonte declarado.
\prov{relações: relaciona olho\_de\_koch; relaciona colapso\_koch\_tiffany}
\prov{proveniência: wiki \textperiodcentered{} inferencia \textperiodcentered{} confiança 1.0 \textperiodcentered{} domínio broca\_repos \textperiodcentered{} história 3 versões no jornal}

%CRISTAL {"arestas":[{"alvo":"colapso_multicamada_peso","confianca":1.0,"epistemico":"fato","origem":"wiki","rel":"especializa"},{"alvo":"conversa_recupera_nao_gera","confianca":1.0,"epistemico":"fato","origem":"wiki","rel":"parte_de"}],"confianca":1.0,"contraexemplos":[],"descricao":"O colapso varre as três camadas na ordem fixa dialogos → conhecimento → corpus (_CAMADAS_CONVERSA), acumulando candidatos; só varre machine com --full. Ao chegar ao corpus, se já há candidato com score ≥ _MIN_OV (200), interrompe em vez de descer mais. A ordem e o break são o mecanismo por trás dos pesos de camada.","epistemico":"fato","exemplos":[],"id":"colapso_camadas_ordem_parada","memoria":"semantica","meta":{"autor":"broca-so","dominio":"conversa","fonte":"conversa/colapso.py:camadas/colapso"},"origem":"wiki","palavras_chave":[],"sinonimos":[],"tipo":"conceito","titulo":"Tiffany — colapso multicamada com ordem e parada antecipada"}
\section{Tiffany — colapso multicamada com ordem e parada antecipada}
O colapso varre as três camadas na ordem fixa dialogos \ensuremath{\to} conhecimento \ensuremath{\to} corpus (\_CAMADAS\_CONVERSA), acumulando candidatos; só varre machine com --full. Ao chegar ao corpus, se já há candidato com score \ensuremath{\ge} \_MIN\_OV (200), interrompe em vez de descer mais. A ordem e o break são o mecanismo por trás dos pesos de camada.
\prov{relações: especializa colapso\_multicamada\_peso; parte\_de conversa\_recupera\_nao\_gera}
\prov{proveniência: wiki \textperiodcentered{} conversa/colapso.py:camadas/colapso \textperiodcentered{} fato \textperiodcentered{} confiança 1.0 \textperiodcentered{} domínio conversa \textperiodcentered{} história 15 versões no jornal}

%CRISTAL {"arestas":[{"alvo":"tensor_mu_torre_corpos","confianca":1.0,"epistemico":"fato","origem":"wiki","rel":"usa"},{"alvo":"decomposicao_peirce","confianca":1.0,"epistemico":"fato","origem":"wiki","rel":"relaciona"}],"confianca":1.0,"contraexemplos":[],"descricao":"Restrita a um único campo local ℂ_â a álgebra vira associativa (≅ℂ): todas as Catalan(n−1) árvores colapsam num produto único e exato, e os associadores se anulam. A complexidade combinatória das árvores mede o grau de separação entre campos locais.","epistemico":"fato","exemplos":[],"id":"colapso_campo_local","memoria":"semantica","meta":{"autor":"Lopes/Aarão","dominio":"hiper","fonte":""},"origem":"wiki","palavras_chave":[],"sinonimos":[],"tipo":"conceito","titulo":"Colapso Intra-Campo-Local"}
\section{Colapso Intra-Campo-Local}
Restrita a um único campo local \ensuremath{\mathbb{C}}\_â a álgebra vira associativa (\ensuremath{\cong}\ensuremath{\mathbb{C}}): todas as Catalan(n\ensuremath{-}1) árvores colapsam num produto único e exato, e os associadores se anulam. A complexidade combinatória das árvores mede o grau de separação entre campos locais.
\prov{relações: usa tensor\_mu\_torre\_corpos; relaciona decomposicao\_peirce}
\prov{proveniência: wiki \textperiodcentered{} fato \textperiodcentered{} confiança 1.0 \textperiodcentered{} domínio hiper \textperiodcentered{} história 4 versões no jornal}

%CRISTAL {"arestas":[{"alvo":"delta","confianca":0.95,"epistemico":"fato","origem":"wiki","rel":"confirma"},{"alvo":"tiffany","confianca":0.95,"epistemico":"fato","origem":"wiki","rel":"confirma"},{"alvo":"colapso","confianca":0.95,"epistemico":"fato","origem":"wiki","rel":"confirma"},{"alvo":"utilitarismo","confianca":0.5,"epistemico":"fato","origem":"wiki","rel":"relaciona"},{"alvo":"pitagorismo","confianca":0.5,"epistemico":"fato","origem":"wiki","rel":"relaciona"},{"alvo":"teoria_do_valor","confianca":0.5,"epistemico":"fato","origem":"wiki","rel":"relaciona"},{"alvo":"pedagogia","confianca":0.5,"epistemico":"fato","origem":"wiki","rel":"relaciona"},{"alvo":"ruido","confianca":0.5,"epistemico":"fato","origem":"wiki","rel":"relaciona"},{"alvo":"gentil_15_propriedade_fundamental_de_uma_pam","confianca":0.5,"epistemico":"fato","origem":"wiki","rel":"relaciona"},{"alvo":"o_que_e_no_projecto","confianca":0.5,"epistemico":"fato","origem":"wiki","rel":"relaciona"},{"alvo":"vida-morte","confianca":0.5,"epistemico":"fato","origem":"wiki","rel":"relaciona"},{"alvo":"vetores","confianca":0.5,"epistemico":"fato","origem":"wiki","rel":"relaciona"},{"alvo":"scheduler","confianca":0.5,"epistemico":"fato","origem":"wiki","rel":"relaciona"}],"confianca":0.95,"contraexemplos":[],"descricao":"Referência verificada: colapso_koch_tiffany.md","epistemico":"fato","exemplos":[],"id":"colapso_koch_tiffany","memoria":"semantica","meta":{"dominio":"referencia","fonte":"web","nivel":5},"origem":"wiki","palavras_chave":[],"sinonimos":[],"tipo":"conceito","titulo":"colapso koch tiffany"}
\section{colapso koch tiffany}
Referência verificada: colapso\_koch\_tiffany.md
\prov{relações: confirma delta; confirma tiffany; confirma colapso; relaciona utilitarismo; relaciona pitagorismo; relaciona teoria\_do\_valor; relaciona pedagogia; relaciona ruido; relaciona gentil\_15\_propriedade\_fundamental\_de\_uma\_pam; relaciona o\_que\_e\_no\_projecto; relaciona vida-morte; relaciona vetores; relaciona scheduler}
\prov{proveniência: wiki \textperiodcentered{} web \textperiodcentered{} fato \textperiodcentered{} confiança 0.95 \textperiodcentered{} domínio referencia \textperiodcentered{} história 25 versões no jornal}

%CRISTAL {"fusao":[{"arestas":[{"alvo":"colapso_koch_tiffany_colapso_koch_operador_ψ_da_tiffany","confianca":1.0,"epistemico":"fato","origem":"colapso_koch_tiffany.md","rel":"parte_de"},{"alvo":"semantica","confianca":0.95,"epistemico":"fato","origem":"wiki","rel":"confirma"},{"alvo":"virtualizacao","confianca":0.5,"epistemico":"fato","origem":"wiki","rel":"relaciona"},{"alvo":"word","confianca":0.5,"epistemico":"fato","origem":"wiki","rel":"relaciona"},{"alvo":"vida-morte","confianca":0.5,"epistemico":"fato","origem":"wiki","rel":"relaciona"},{"alvo":"ontologia","confianca":0.5,"epistemico":"fato","origem":"wiki","rel":"relaciona"},{"alvo":"misticismo_nao_dual","confianca":0.5,"epistemico":"fato","origem":"wiki","rel":"relaciona"}],"confianca":1.0,"contraexemplos":[],"descricao":"Prioridade (peso `_LAYER_WEIGHT` em `colapso.py`):","epistemico":"fato","exemplos":[],"id":"colapso_koch_tiffany_camadas_de_tecido","memoria":"semantica","meta":{"dominio":"referencia","fonte":"web","nivel":5},"origem":"colapso_koch_tiffany.md","palavras_chave":[],"sinonimos":[],"tipo":"conceito","titulo":"Camadas de tecido"},{"arestas":[{"alvo":"colapso_koch_operador_ψ_da_tiffany","confianca":1.0,"epistemico":"fato","origem":"colapso_koch_tiffany.md","rel":"parte_de"},{"alvo":"semantica","confianca":0.95,"epistemico":"fato","origem":"wiki","rel":"confirma"},{"alvo":"vida-morte","confianca":0.5,"epistemico":"fato","origem":"wiki","rel":"relaciona"},{"alvo":"virtualizacao","confianca":0.5,"epistemico":"fato","origem":"wiki","rel":"relaciona"},{"alvo":"pell","confianca":0.5,"epistemico":"fato","origem":"wiki","rel":"relaciona"},{"alvo":"word","confianca":0.5,"epistemico":"fato","origem":"wiki","rel":"relaciona"},{"alvo":"readme","confianca":0.5,"epistemico":"fato","origem":"wiki","rel":"relaciona"},{"alvo":"sequencia_pow_nativa","confianca":0.5,"epistemico":"fato","origem":"wiki","rel":"relaciona"},{"alvo":"telemetria_shadow_producao","confianca":0.5,"epistemico":"fato","origem":"wiki","rel":"relaciona"},{"alvo":"fibonacci_multidimensional","confianca":0.5,"epistemico":"fato","origem":"wiki","rel":"relaciona"},{"alvo":"gentil_capitulo_9_programando_a_hp_prime","confianca":0.5,"epistemico":"fato","origem":"wiki","rel":"relaciona"},{"alvo":"incautos_enganados","confianca":0.5,"epistemico":"fato","origem":"wiki","rel":"relaciona"},{"alvo":"gentil_12_definicao","confianca":0.5,"epistemico":"fato","origem":"wiki","rel":"relaciona"},{"alvo":"consciencia","confianca":0.5,"epistemico":"fato","origem":"wiki","rel":"relaciona"},{"alvo":"religiao","confianca":0.5,"epistemico":"fato","origem":"wiki","rel":"relaciona"}],"confianca":1.0,"contraexemplos":[],"descricao":"Prioridade (peso `_LAYER_WEIGHT` em `colapso.py`):","epistemico":"fato","exemplos":[],"id":"camadas_de_tecido","memoria":"semantica","meta":{"dominio":"referencia","nivel":5},"origem":"colapso_koch_tiffany.md","palavras_chave":[],"sinonimos":[],"tipo":"conceito","titulo":"Camadas de tecido"}],"id":"colapso_koch_tiffany_camadas_de_tecido","tipo":"conceito"}
\section{Camadas de tecido}
Prioridade (peso `\_LAYER\_WEIGHT` em `colapso.py`):
\prov{relações: parte\_de colapso\_koch\_tiffany\_colapso\_koch\_operador\_\ensuremath{\psi}\_da\_tiffany; confirma semantica; relaciona virtualizacao; relaciona word; relaciona vida-morte; relaciona ontologia; relaciona misticismo\_nao\_dual}
\prov{proveniência: colapso\_koch\_tiffany.md \textperiodcentered{} web \textperiodcentered{} fato \textperiodcentered{} confiança 1.0 \textperiodcentered{} domínio referencia \textperiodcentered{} história 9 versões no jornal}
\prov{fusão de curadoria (texto idêntico): absorve camadas\_de\_tecido --- o mesmo conteúdo por dois esquemas de endereço}

%CRISTAL {"fusao":[{"arestas":[{"alvo":"colapso_koch_tiffany","confianca":0.95,"epistemico":"fato","origem":"wiki","rel":"parte_de"}],"confianca":1.0,"contraexemplos":[],"descricao":"> **Epistemologia:** fato de implementação Broca-SO (código + BAI + papers). > **No grafo:** `epistemico=fato`, confiança 0.95. > **Papel:** fechar núcleo conversacional — dependências, exemplos, contraexemplos.","epistemico":"fato","exemplos":[],"id":"colapso_koch_tiffany_colapso_koch_operador_ψ_da_tiffany","memoria":"semantica","meta":{"dominio":"referencia","fonte":"web","nivel":5},"origem":"colapso_koch_tiffany.md","palavras_chave":[],"sinonimos":[],"tipo":"conceito","titulo":"Colapso Koch — operador Ψ da Tiffany"},{"arestas":[{"alvo":"colapso_koch_tiffany","confianca":0.95,"epistemico":"fato","origem":"wiki","rel":"parte_de"}],"confianca":1.0,"contraexemplos":[],"descricao":"> **Epistemologia:** fato de implementação Broca-SO (código + BAI + papers). > **No grafo:** `epistemico=fato`, confiança 0.95. > **Papel:** fechar núcleo conversacional — dependências, exemplos, contraexemplos.","epistemico":"fato","exemplos":[],"id":"colapso_koch_operador_ψ_da_tiffany","memoria":"semantica","meta":{"dominio":"referencia","nivel":5},"origem":"colapso_koch_tiffany.md","palavras_chave":[],"sinonimos":[],"tipo":"conceito","titulo":"Colapso Koch — operador Ψ da Tiffany"}],"id":"colapso_koch_tiffany_colapso_koch_operador_ψ_da_tiffany","tipo":"conceito"}
\section{Colapso Koch — operador \ensuremath{\Psi} da Tiffany}
> **Epistemologia:** fato de implementação Broca-SO (código + BAI + papers). > **No grafo:** `epistemico=fato`, confiança 0.95. > **Papel:** fechar núcleo conversacional — dependências, exemplos, contraexemplos.
\prov{relações: parte\_de colapso\_koch\_tiffany}
\prov{proveniência: colapso\_koch\_tiffany.md \textperiodcentered{} web \textperiodcentered{} fato \textperiodcentered{} confiança 1.0 \textperiodcentered{} domínio referencia \textperiodcentered{} história 3 versões no jornal}
\prov{fusão de curadoria (texto idêntico): absorve colapso\_koch\_operador\_\ensuremath{\psi}\_da\_tiffany --- o mesmo conteúdo por dois esquemas de endereço}

%CRISTAL {"arestas":[{"alvo":"colapso_koch_tiffany_colapso_koch_operador_ψ_da_tiffany","confianca":1.0,"epistemico":"fato","origem":"colapso_koch_tiffany.md","rel":"parte_de"},{"alvo":"virtualizacao","confianca":0.5,"epistemico":"fato","origem":"wiki","rel":"relaciona"},{"alvo":"word","confianca":0.5,"epistemico":"fato","origem":"wiki","rel":"relaciona"}],"confianca":1.0,"contraexemplos":[],"descricao":"Confundir Ψ com colapso da função de onda (postulado MQ) — operadores distintos.","epistemico":"fato","exemplos":[],"id":"colapso_koch_tiffany_contraexemplos","memoria":"semantica","meta":{"dominio":"referencia","fonte":"web","nivel":5},"origem":"colapso_koch_tiffany.md","palavras_chave":[],"sinonimos":[],"tipo":"conceito","titulo":"Contraexemplos"}
\section{Contraexemplos}
Confundir \ensuremath{\Psi} com colapso da função de onda (postulado MQ) — operadores distintos.
\prov{relações: parte\_de colapso\_koch\_tiffany\_colapso\_koch\_operador\_\ensuremath{\psi}\_da\_tiffany; relaciona virtualizacao; relaciona word}
\prov{proveniência: colapso\_koch\_tiffany.md \textperiodcentered{} web \textperiodcentered{} fato \textperiodcentered{} confiança 1.0 \textperiodcentered{} domínio referencia \textperiodcentered{} história 5 versões no jornal}

%CRISTAL {"arestas":[{"alvo":"colapso_koch_tiffany_colapso_koch_operador_ψ_da_tiffany","confianca":1.0,"epistemico":"fato","origem":"colapso_koch_tiffany.md","rel":"parte_de"},{"alvo":"virtualizacao","confianca":0.5,"epistemico":"fato","origem":"wiki","rel":"relaciona"},{"alvo":"word","confianca":0.5,"epistemico":"fato","origem":"wiki","rel":"relaciona"}],"confianca":1.0,"contraexemplos":[],"descricao":"- `conversa/colapso.py` - `papers/casl-propagation.pdf` - `ula/BROCA_CRISTAL.md` - `ula/KOCH_ATLAS.md`","epistemico":"fato","exemplos":[],"id":"colapso_koch_tiffany_fontes","memoria":"semantica","meta":{"dominio":"referencia","fonte":"web","nivel":5},"origem":"colapso_koch_tiffany.md","palavras_chave":[],"sinonimos":[],"tipo":"conceito","titulo":"Fontes"}
\section{Fontes}
- `conversa/colapso.py` - `papers/casl-propagation.pdf` - `ula/BROCA\_CRISTAL.md` - `ula/KOCH\_ATLAS.md`
\prov{relações: parte\_de colapso\_koch\_tiffany\_colapso\_koch\_operador\_\ensuremath{\psi}\_da\_tiffany; relaciona virtualizacao; relaciona word}
\prov{proveniência: colapso\_koch\_tiffany.md \textperiodcentered{} web \textperiodcentered{} fato \textperiodcentered{} confiança 1.0 \textperiodcentered{} domínio referencia \textperiodcentered{} história 5 versões no jornal}

%CRISTAL {"fusao":[{"arestas":[{"alvo":"colapso_koch_tiffany_colapso_koch_operador_ψ_da_tiffany","confianca":1.0,"epistemico":"fato","origem":"colapso_koch_tiffany.md","rel":"parte_de"},{"alvo":"colapso","confianca":0.95,"epistemico":"fato","origem":"wiki","rel":"confirma"},{"alvo":"bai_biblioteca_de_autorizacao_isomorfica","confianca":0.95,"epistemico":"fato","origem":"wiki","rel":"confirma"}],"confianca":1.0,"contraexemplos":[],"descricao":"Ψ = **Collapse(𝒞_δ)** — operador de colapso conversacional:","epistemico":"fato","exemplos":[],"id":"colapso_koch_tiffany_operador_ψ_collapse","memoria":"semantica","meta":{"dominio":"referencia","fonte":"web","nivel":5},"origem":"colapso_koch_tiffany.md","palavras_chave":[],"sinonimos":[],"tipo":"conceito","titulo":"Operador Ψ (Collapse)"},{"arestas":[{"alvo":"colapso_koch_operador_ψ_da_tiffany","confianca":1.0,"epistemico":"fato","origem":"colapso_koch_tiffany.md","rel":"parte_de"},{"alvo":"colapso","confianca":0.95,"epistemico":"fato","origem":"wiki","rel":"confirma"},{"alvo":"bai_biblioteca_de_autorizacao_isomorfica","confianca":0.95,"epistemico":"fato","origem":"wiki","rel":"confirma"},{"alvo":"while","confianca":0.5,"epistemico":"fato","origem":"wiki","rel":"relaciona"},{"alvo":"virtualizacao","confianca":0.5,"epistemico":"fato","origem":"wiki","rel":"relaciona"}],"confianca":1.0,"contraexemplos":[],"descricao":"Ψ = **Collapse(𝒞_δ)** — operador de colapso conversacional:","epistemico":"fato","exemplos":[],"id":"operador_ψ_collapse","memoria":"semantica","meta":{"dominio":"referencia","nivel":5},"origem":"colapso_koch_tiffany.md","palavras_chave":[],"sinonimos":[],"tipo":"conceito","titulo":"Operador Ψ (Collapse)"}],"id":"colapso_koch_tiffany_operador_ψ_collapse","tipo":"conceito"}
\section{Operador \ensuremath{\Psi} (Collapse)}
\ensuremath{\Psi} = **Collapse(\ensuremath{\mathcal{C}}\_\ensuremath{\delta})** — operador de colapso conversacional:
\prov{relações: parte\_de colapso\_koch\_tiffany\_colapso\_koch\_operador\_\ensuremath{\psi}\_da\_tiffany; confirma colapso; confirma bai\_biblioteca\_de\_autorizacao\_isomorfica}
\prov{proveniência: colapso\_koch\_tiffany.md \textperiodcentered{} web \textperiodcentered{} fato \textperiodcentered{} confiança 1.0 \textperiodcentered{} domínio referencia \textperiodcentered{} história 5 versões no jornal}
\prov{fusão de curadoria (texto idêntico): absorve operador\_\ensuremath{\psi}\_collapse --- o mesmo conteúdo por dois esquemas de endereço}

%CRISTAL {"arestas":[{"alvo":"colapso_koch_tiffany_colapso_koch_operador_ψ_da_tiffany","confianca":1.0,"epistemico":"fato","origem":"colapso_koch_tiffany.md","rel":"parte_de"},{"alvo":"tiffany","confianca":0.95,"epistemico":"fato","origem":"wiki","rel":"confirma"},{"alvo":"colapso","confianca":0.95,"epistemico":"fato","origem":"wiki","rel":"confirma"},{"alvo":"gato","confianca":0.95,"epistemico":"fato","origem":"wiki","rel":"usa"},{"alvo":"olho_de_koch","confianca":0.95,"epistemico":"fato","origem":"wiki","rel":"usa"},{"alvo":"codigo_conversa_colapso","confianca":0.95,"epistemico":"fato","origem":"wiki","rel":"define"}],"confianca":1.0,"contraexemplos":[],"descricao":"| Colapso | Nó Broca-SO | relação | |---------|-------------|---------| | motor Ψ | colapso | confirma | | resposta Tiffany | tiffany | explica | | overlap gato | gato | usa | | orçamento | delta | depende | | quintuplo | kappa | depende | | BAI formal | bai_biblioteca_de_autorizacao_isomorfica | parte_de | | atlas | olho_de_koch | usa | | painel nonce | painel_estelar | referencia | | implementação | codigo_conversa_colapso | define |","epistemico":"fato","exemplos":[],"id":"colapso_koch_tiffany_ponte_broca-so","memoria":"semantica","meta":{"dominio":"referencia","fonte":"web","nivel":5},"origem":"colapso_koch_tiffany.md","palavras_chave":[],"sinonimos":[],"tipo":"conceito","titulo":"Ponte Broca-SO"}
\section{Ponte Broca-SO}
| Colapso | Nó Broca-SO | relação | |---------|-------------|---------| | motor \ensuremath{\Psi} | colapso | confirma | | resposta Tiffany | tiffany | explica | | overlap gato | gato | usa | | orçamento | delta | depende | | quintuplo | kappa | depende | | BAI formal | bai\_biblioteca\_de\_autorizacao\_isomorfica | parte\_de | | atlas | olho\_de\_koch | usa | | painel nonce | painel\_estelar | referencia | | implementação | codigo\_conversa\_colapso | define |
\prov{relações: parte\_de colapso\_koch\_tiffany\_colapso\_koch\_operador\_\ensuremath{\psi}\_da\_tiffany; confirma tiffany; confirma colapso; usa gato; usa olho\_de\_koch; define codigo\_conversa\_colapso}
\prov{proveniência: colapso\_koch\_tiffany.md \textperiodcentered{} web \textperiodcentered{} fato \textperiodcentered{} confiança 1.0 \textperiodcentered{} domínio referencia \textperiodcentered{} história 8 versões no jornal}

%CRISTAL {"fusao":[{"arestas":[{"alvo":"colapso_koch_tiffany_colapso_koch_operador_ψ_da_tiffany","confianca":1.0,"epistemico":"fato","origem":"colapso_koch_tiffany.md","rel":"parte_de"},{"alvo":"kappa","confianca":0.95,"epistemico":"fato","origem":"wiki","rel":"referencia"},{"alvo":"delta","confianca":0.95,"epistemico":"fato","origem":"wiki","rel":"referencia"}],"confianca":1.0,"contraexemplos":[],"descricao":"| Primitiva | Papel no colapso | |-----------|------------------| | **κ** | capacidade / quintuplo propagado | | **δ** | orçamento de profundidade — limita passos Ψ | | **Φ** | organização estrutural observada pós-propagação | | **T** | atenuação deflacionária | | **Ψ** | colapso — escolhe face legível |","epistemico":"fato","exemplos":[],"id":"colapso_koch_tiffany_quinteto_bai","memoria":"semantica","meta":{"dominio":"referencia","fonte":"web","nivel":5},"origem":"colapso_koch_tiffany.md","palavras_chave":[],"sinonimos":[],"tipo":"conceito","titulo":"Quinteto BAI"},{"arestas":[{"alvo":"colapso_koch_operador_ψ_da_tiffany","confianca":1.0,"epistemico":"fato","origem":"colapso_koch_tiffany.md","rel":"parte_de"},{"alvo":"kappa","confianca":0.95,"epistemico":"fato","origem":"wiki","rel":"referencia"},{"alvo":"delta","confianca":0.95,"epistemico":"fato","origem":"wiki","rel":"referencia"}],"confianca":1.0,"contraexemplos":[],"descricao":"| Primitiva | Papel no colapso | |-----------|------------------| | **κ** | capacidade / quintuplo propagado | | **δ** | orçamento de profundidade — limita passos Ψ | | **Φ** | organização estrutural observada pós-propagação | | **T** | atenuação deflacionária | | **Ψ** | colapso — escolhe face legível |","epistemico":"fato","exemplos":[],"id":"quinteto_bai","memoria":"semantica","meta":{"dominio":"referencia","nivel":5},"origem":"colapso_koch_tiffany.md","palavras_chave":[],"sinonimos":[],"tipo":"conceito","titulo":"Quinteto BAI"}],"id":"colapso_koch_tiffany_quinteto_bai","tipo":"conceito"}
\section{Quinteto BAI}
| Primitiva | Papel no colapso | |-----------|------------------| | **\ensuremath{\kappa}** | capacidade / quintuplo propagado | | **\ensuremath{\delta}** | orçamento de profundidade — limita passos \ensuremath{\Psi} | | **\ensuremath{\Phi}** | organização estrutural observada pós-propagação | | **T** | atenuação deflacionária | | **\ensuremath{\Psi}** | colapso — escolhe face legível |
\prov{relações: parte\_de colapso\_koch\_tiffany\_colapso\_koch\_operador\_\ensuremath{\psi}\_da\_tiffany; referencia kappa; referencia delta}
\prov{proveniência: colapso\_koch\_tiffany.md \textperiodcentered{} web \textperiodcentered{} fato \textperiodcentered{} confiança 1.0 \textperiodcentered{} domínio referencia \textperiodcentered{} história 5 versões no jornal}
\prov{fusão de curadoria (texto idêntico): absorve quinteto\_bai --- o mesmo conteúdo por dois esquemas de endereço}

%CRISTAL {"arestas":[{"alvo":"colapso_koch_tiffany_colapso_koch_operador_ψ_da_tiffany","confianca":1.0,"epistemico":"fato","origem":"colapso_koch_tiffany.md","rel":"parte_de"}],"confianca":1.0,"contraexemplos":[],"descricao":"- **BAI / CASL:** quinteto (κ, δ, Φ, T, **Ψ**) — `papers/casl-propagation.pdf` - **Código:** `conversa/colapso.py` — seleção de face vencedora no tecido - **Arquitectura:** `ula/BROCA_CRISTAL.md` — núcleo gato → Olho de Koch → borda clássica - **Koch:** `ula/KOCH_ATLAS.md`, `papers/olho_de_koch.tex`","epistemico":"fato","exemplos":[],"id":"colapso_koch_tiffany_referencia","memoria":"semantica","meta":{"dominio":"referencia","fonte":"web","nivel":5},"origem":"colapso_koch_tiffany.md","palavras_chave":[],"sinonimos":[],"tipo":"conceito","titulo":"Referência"}
\section{Referência}
- **BAI / CASL:** quinteto (\ensuremath{\kappa}, \ensuremath{\delta}, \ensuremath{\Phi}, T, **\ensuremath{\Psi}**) — `papers/casl-propagation.pdf` - **Código:** `conversa/colapso.py` — seleção de face vencedora no tecido - **Arquitectura:** `ula/BROCA\_CRISTAL.md` — núcleo gato \ensuremath{\to} Olho de Koch \ensuremath{\to} borda clássica - **Koch:** `ula/KOCH\_ATLAS.md`, `papers/olho\_de\_koch.tex`
\prov{relações: parte\_de colapso\_koch\_tiffany\_colapso\_koch\_operador\_\ensuremath{\psi}\_da\_tiffany}
\prov{proveniência: colapso\_koch\_tiffany.md \textperiodcentered{} web \textperiodcentered{} fato \textperiodcentered{} confiança 1.0 \textperiodcentered{} domínio referencia \textperiodcentered{} história 3 versões no jornal}

%CRISTAL {"fusao":[{"arestas":[{"alvo":"colapso_koch_tiffany_colapso_koch_operador_ψ_da_tiffany","confianca":1.0,"epistemico":"fato","origem":"colapso_koch_tiffany.md","rel":"parte_de"},{"alvo":"ponte_quantica_dougherty","confianca":0.95,"epistemico":"fato","origem":"wiki","rel":"analogia"},{"alvo":"fase_geometrica_berry","confianca":0.95,"epistemico":"fato","origem":"wiki","rel":"referencia"}],"confianca":1.0,"contraexemplos":[],"descricao":"Dougherty/Ginzburg mencionam “colapso” de vórtices/toros — **analogia metafórica** apenas.","epistemico":"fato","exemplos":[],"id":"colapso_koch_tiffany_relacao_com_hipoteses_externas","memoria":"semantica","meta":{"dominio":"referencia","fonte":"web","nivel":5},"origem":"colapso_koch_tiffany.md","palavras_chave":[],"sinonimos":[],"tipo":"conceito","titulo":"Relação com hipóteses externas"},{"arestas":[{"alvo":"colapso_koch_operador_ψ_da_tiffany","confianca":1.0,"epistemico":"fato","origem":"colapso_koch_tiffany.md","rel":"parte_de"},{"alvo":"ponte_quantica_dougherty","confianca":0.95,"epistemico":"fato","origem":"wiki","rel":"analogia"},{"alvo":"fase_geometrica_berry","confianca":0.95,"epistemico":"fato","origem":"wiki","rel":"referencia"}],"confianca":1.0,"contraexemplos":[],"descricao":"Dougherty/Ginzburg mencionam “colapso” de vórtices/toros — **analogia metafórica** apenas.","epistemico":"fato","exemplos":[],"id":"relacao_com_hipoteses_externas","memoria":"semantica","meta":{"dominio":"referencia","nivel":5},"origem":"colapso_koch_tiffany.md","palavras_chave":[],"sinonimos":[],"tipo":"conceito","titulo":"Relação com hipóteses externas"}],"id":"colapso_koch_tiffany_relacao_com_hipoteses_externas","tipo":"conceito"}
\section{Relação com hipóteses externas}
Dougherty/Ginzburg mencionam “colapso” de vórtices/toros — **analogia metafórica** apenas.
\prov{relações: parte\_de colapso\_koch\_tiffany\_colapso\_koch\_operador\_\ensuremath{\psi}\_da\_tiffany; analogia ponte\_quantica\_dougherty; referencia fase\_geometrica\_berry}
\prov{proveniência: colapso\_koch\_tiffany.md \textperiodcentered{} web \textperiodcentered{} fato \textperiodcentered{} confiança 1.0 \textperiodcentered{} domínio referencia \textperiodcentered{} história 5 versões no jornal}
\prov{fusão de curadoria (texto idêntico): absorve relacao\_com\_hipoteses\_externas --- o mesmo conteúdo por dois esquemas de endereço}

%CRISTAL {"fusao":[{"arestas":[{"alvo":"colapso_koch_tiffany_colapso_koch_operador_ψ_da_tiffany","confianca":1.0,"epistemico":"fato","origem":"colapso_koch_tiffany.md","rel":"parte_de"},{"alvo":"olho_de_koch","confianca":0.95,"epistemico":"fato","origem":"wiki","rel":"confirma"},{"alvo":"koch_atlas_experimento_de_parametrizacao","confianca":0.95,"epistemico":"fato","origem":"wiki","rel":"referencia"},{"alvo":"word","confianca":0.5,"epistemico":"fato","origem":"wiki","rel":"relaciona"},{"alvo":"syscall","confianca":0.5,"epistemico":"fato","origem":"wiki","rel":"relaciona"},{"alvo":"virtualizacao","confianca":0.5,"epistemico":"fato","origem":"wiki","rel":"relaciona"},{"alvo":"contradicao","confianca":0.5,"epistemico":"fato","origem":"wiki","rel":"relaciona"},{"alvo":"vontade_de_poder","confianca":0.5,"epistemico":"fato","origem":"wiki","rel":"relaciona"}],"confianca":1.0,"contraexemplos":[],"descricao":"Pipeline Broca cristalizado:","epistemico":"fato","exemplos":[],"id":"colapso_koch_tiffany_relacao_koch","memoria":"semantica","meta":{"dominio":"referencia","fonte":"web","nivel":5},"origem":"colapso_koch_tiffany.md","palavras_chave":[],"sinonimos":[],"tipo":"conceito","titulo":"Relação Koch"},{"arestas":[{"alvo":"colapso_koch_operador_ψ_da_tiffany","confianca":1.0,"epistemico":"fato","origem":"colapso_koch_tiffany.md","rel":"parte_de"},{"alvo":"olho_de_koch","confianca":0.95,"epistemico":"fato","origem":"wiki","rel":"confirma"},{"alvo":"koch_atlas_experimento_de_parametrizacao","confianca":0.95,"epistemico":"fato","origem":"wiki","rel":"referencia"},{"alvo":"word","confianca":0.5,"epistemico":"fato","origem":"wiki","rel":"relaciona"},{"alvo":"virtualizacao","confianca":0.5,"epistemico":"fato","origem":"wiki","rel":"relaciona"}],"confianca":1.0,"contraexemplos":[],"descricao":"Pipeline Broca cristalizado:","epistemico":"fato","exemplos":[],"id":"relacao_koch","memoria":"semantica","meta":{"dominio":"referencia","nivel":5},"origem":"colapso_koch_tiffany.md","palavras_chave":[],"sinonimos":[],"tipo":"conceito","titulo":"Relação Koch"}],"id":"colapso_koch_tiffany_relacao_koch","tipo":"conceito"}
\section{Relação Koch}
Pipeline Broca cristalizado:
\prov{relações: parte\_de colapso\_koch\_tiffany\_colapso\_koch\_operador\_\ensuremath{\psi}\_da\_tiffany; confirma olho\_de\_koch; referencia koch\_atlas\_experimento\_de\_parametrizacao; relaciona word; relaciona syscall; relaciona virtualizacao; relaciona contradicao; relaciona vontade\_de\_poder}
\prov{proveniência: colapso\_koch\_tiffany.md \textperiodcentered{} web \textperiodcentered{} fato \textperiodcentered{} confiança 1.0 \textperiodcentered{} domínio referencia \textperiodcentered{} história 10 versões no jornal}
\prov{fusão de curadoria (texto idêntico): absorve relacao\_koch --- o mesmo conteúdo por dois esquemas de endereço}

%CRISTAL {"arestas":[{"alvo":"montar_conversa_tecido","confianca":1.0,"epistemico":"fato","origem":"wiki","rel":"parte_de"},{"alvo":"colapso_ping_pong_overlap","confianca":1.0,"epistemico":"fato","origem":"wiki","rel":"relaciona"}],"confianca":1.0,"contraexemplos":[],"descricao":"_colisoes_koch remove linhas do Tatoeba que ressoam alto (ov ≥ 254) com saudações curtas ('oi', 'bom dia') mas não as contêm de fato — evita que colisões Koch=256 poluam o colapso de falas curtas. Contraparte, na montagem, do overlap ping-pong que resolve empates no scan.","epistemico":"fato","exemplos":[],"id":"colisoes_koch_filtro","memoria":"semantica","meta":{"autor":"broca-so","dominio":"conversa","fonte":"conversa/montar_conversa.py:_colisoes_koch"},"origem":"wiki","palavras_chave":[],"sinonimos":[],"tipo":"conceito","titulo":"Tiffany — filtro de falsas ressonâncias de saudação"}
\section{Tiffany — filtro de falsas ressonâncias de saudação}
\_colisoes\_koch remove linhas do Tatoeba que ressoam alto (ov \ensuremath{\ge} 254) com saudações curtas ('oi', 'bom dia') mas não as contêm de fato — evita que colisões Koch=256 poluam o colapso de falas curtas. Contraparte, na montagem, do overlap ping-pong que resolve empates no scan.
\prov{relações: parte\_de montar\_conversa\_tecido; relaciona colapso\_ping\_pong\_overlap}
\prov{proveniência: wiki \textperiodcentered{} conversa/montar\_conversa.py:\_colisoes\_koch \textperiodcentered{} fato \textperiodcentered{} confiança 1.0 \textperiodcentered{} domínio conversa \textperiodcentered{} história 15 versões no jornal}

%CRISTAL {"arestas":[{"alvo":"casl-propagation","confianca":0.047,"epistemico":"fato","origem":"manual","rel":"referencia"},{"alvo":"bai_biblioteca_de_autorizacao_isomorfica","confianca":-0.039,"epistemico":"fato","origem":"manual","rel":"referencia"},{"alvo":"bai","confianca":-0.078,"epistemico":"fato","origem":"manual","rel":"referencia"}],"confianca":1.0,"contraexemplos":[],"descricao":"```text cristal                          → sessão interativa (REPL) ├── projeto    new | init | open | status ├── new PATH                     → atalho projeto new ├── agent    list | software | paper | pesquisa | qa | devops ├── chat     [agente] \"msg\"      → conversa / sessão ├── code     explain | review | tests | benchmark | refactor… ├── task     add | list | work ├── work                         → executar tarefa aberta","epistemico":"fato","exemplos":[],"id":"comandos","memoria":"semantica","meta":{"dominio":"manual","duplicata_sugerida":[],"nivel":6,"verificacao":"parte de conceito mais amplo 'comandos'"},"origem":"CRISTAL_DEV.md","palavras_chave":[],"sinonimos":[],"tipo":"conceito","titulo":"Comandos"}
\section{Comandos}
```text cristal                          \ensuremath{\to} sessão interativa (REPL) |--- projeto    new | init | open | status |--- new PATH                     \ensuremath{\to} atalho projeto new |--- agent    list | software | paper | pesquisa | qa | devops |--- chat     [agente] "msg"      \ensuremath{\to} conversa / sessão |--- code     explain | review | tests | benchmark | refactor… |--- task     add | list | work |--- work                         \ensuremath{\to} executar tarefa aberta
\prov{relações: referencia casl-propagation; referencia bai\_biblioteca\_de\_autorizacao\_isomorfica; referencia bai}
\prov{proveniência: CRISTAL\_DEV.md \textperiodcentered{} fato \textperiodcentered{} confiança 1.0 \textperiodcentered{} domínio manual \textperiodcentered{} história 15 versões no jornal}

%CRISTAL {"arestas":[{"alvo":"readme_cristal","confianca":1.0,"epistemico":"fato","origem":"README_CRISTAL.md","rel":"parte_de"},{"alvo":"misticismo_nao_dual","confianca":0.5,"epistemico":"fato","origem":"wiki","rel":"relaciona"},{"alvo":"estetica","confianca":0.5,"epistemico":"fato","origem":"wiki","rel":"relaciona"},{"alvo":"word","confianca":0.5,"epistemico":"fato","origem":"wiki","rel":"relaciona"}],"confianca":1.0,"contraexemplos":[],"descricao":"Comandos flat (`cristal ingest`, `cristal doctor`) **permanecem**. Novos namespaces (`cristal tecido ingest`) são o caminho preferido.","epistemico":"fato","exemplos":[],"id":"compatibilidade","memoria":"semantica","meta":{"dominio":"manual","duplicata_sugerida":[],"nivel":6,"verificacao":"parte de conceito mais amplo 'compatibilidade'"},"origem":"CRISTAL_DEV.md","palavras_chave":[],"sinonimos":[],"tipo":"conceito","titulo":"Compatibilidade"}
\section{Compatibilidade}
Comandos flat (`cristal ingest`, `cristal doctor`) **permanecem**. Novos namespaces (`cristal tecido ingest`) são o caminho preferido.
\prov{relações: parte\_de readme\_cristal; relaciona misticismo\_nao\_dual; relaciona estetica; relaciona word}
\prov{proveniência: CRISTAL\_DEV.md \textperiodcentered{} fato \textperiodcentered{} confiança 1.0 \textperiodcentered{} domínio manual \textperiodcentered{} história 9 versões no jornal}

%CRISTAL {"arestas":[{"alvo":"a_base_topologica","confianca":1.0,"epistemico":"fato","origem":"paper","rel":"parte_de"},{"alvo":"termodinamica","confianca":1.0,"epistemico":"fato","origem":"paper","rel":"parte_de"},{"alvo":"mecanica_quantica","confianca":1.0,"epistemico":"fato","origem":"paper","rel":"parte_de"},{"alvo":"seguranca_por_irreversibilidade","confianca":1.0,"epistemico":"fato","origem":"paper","rel":"parte_de"},{"alvo":"redes","confianca":1.0,"epistemico":"fato","origem":"paper","rel":"parte_de"},{"alvo":"computacao_grafica","confianca":1.0,"epistemico":"fato","origem":"paper","rel":"parte_de"},{"alvo":"paper_0_sintese","confianca":1.0,"epistemico":"fato","origem":"paper","rel":"parte_de"},{"alvo":"pow_gato","confianca":1.0,"epistemico":"fato","origem":"paper","rel":"parte_de"},{"alvo":"pow_coordenadas_espectrais_broca","confianca":1.0,"epistemico":"fato","origem":"POW_GEOMETRIA.md","rel":"parte_de"},{"alvo":"pow_espectral_investigacao_broca","confianca":1.0,"epistemico":"fato","origem":"POW_ESPECTRAL.md","rel":"parte_de"}],"confianca":1.0,"contraexemplos":[],"descricao":"A parábola **organiza** a deformação: mesmo perturbados, os pares válidos fecham a+c²=1 dentro de ε de máquina. A ruptura aparece primeiro como **degeneração** (estados inválidos), não como violação da identidade. Geometricamente, o erro manifesta-se como **deslocamento sobre o arco** (calor,fase), não como saída da curva.","epistemico":"fato","exemplos":[],"id":"conclusao","memoria":"semantica","meta":{"dominio":"manual","nivel":6,"verificacao":"slug idêntico — enriquecer, não duplicar"},"origem":"RUÍDO.md","palavras_chave":[],"sinonimos":[],"tipo":"conceito","titulo":"Conclusão"}
\section{Conclusão}
A parábola **organiza** a deformação: mesmo perturbados, os pares válidos fecham a+c\ensuremath{^{2}}=1 dentro de \ensuremath{\varepsilon} de máquina. A ruptura aparece primeiro como **degeneração** (estados inválidos), não como violação da identidade. Geometricamente, o erro manifesta-se como **deslocamento sobre o arco** (calor,fase), não como saída da curva.
\prov{relações: parte\_de a\_base\_topologica; parte\_de termodinamica; parte\_de mecanica\_quantica; parte\_de seguranca\_por\_irreversibilidade; parte\_de redes; parte\_de computacao\_grafica; parte\_de paper\_0\_sintese; parte\_de pow\_gato; parte\_de pow\_coordenadas\_espectrais\_broca; parte\_de pow\_espectral\_investigacao\_broca}
\prov{proveniência: RUÍDO.md \textperiodcentered{} fato \textperiodcentered{} confiança 1.0 \textperiodcentered{} domínio manual \textperiodcentered{} história 45 versões no jornal}

%CRISTAL {"arestas":[{"alvo":"gato","confianca":0.95,"epistemico":"fato","origem":"paper","rel":"parte_de"}],"confianca":1.0,"contraexemplos":[],"descricao":"Ressonancia 2000000 nonces: acopl medio=0.6049 max=0.9991. Fecho 2 (acopl medio 0.9988 vs 0.6049). SHA 1 (acopl medio 0.5969 vs 0.6049). Separacao fecho=+0.3939 sha=-0.0080. Acopl prediz fecho=sim.","epistemico":"fato","exemplos":[],"id":"conclusao_automatica","memoria":"semantica","meta":{"dominio":"manual","nivel":6,"verificacao":"slug idêntico — enriquecer, não duplicar"},"origem":"POW_METAL_RESSONANCIA.md","palavras_chave":[],"sinonimos":[],"tipo":"conceito","titulo":"Conclusao automatica"}
\section{Conclusao automatica}
Ressonancia 2000000 nonces: acopl medio=0.6049 max=0.9991. Fecho 2 (acopl medio 0.9988 vs 0.6049). SHA 1 (acopl medio 0.5969 vs 0.6049). Separacao fecho=+0.3939 sha=-0.0080. Acopl prediz fecho=sim.
\prov{relações: parte\_de gato}
\prov{proveniência: POW\_METAL\_RESSONANCIA.md \textperiodcentered{} fato \textperiodcentered{} confiança 1.0 \textperiodcentered{} domínio manual \textperiodcentered{} história 5 versões no jornal}

%CRISTAL {"arestas":[{"alvo":"koch_atlas_experimento_de_parametrizacao","confianca":1.0,"epistemico":"fato","origem":"KOCH_ATLAS.md","rel":"parte_de"}],"confianca":1.0,"contraexemplos":[],"descricao":"Estados degenerados (total+e efetivo ≤ 0) podem conter organização interna. Pell e a dimensão de prata entram **depois**, como comparação — não como filtro.","epistemico":"fato","exemplos":[],"id":"conclusao_automatica_medicao_nao_hipotese","memoria":"semantica","meta":{"dominio":"manual","duplicata_sugerida":[],"nivel":6,"verificacao":"parte de conceito mais amplo 'conclusao_automatica_medicao_nao_hipotese'"},"origem":"DEGENERADOS.md","palavras_chave":[],"sinonimos":[],"tipo":"conceito","titulo":"Conclusão automática (medição, não hipótese)"}
\section{Conclusão automática (medição, não hipótese)}
Estados degenerados (total+e efetivo \ensuremath{\le} 0) podem conter organização interna. Pell e a dimensão de prata entram **depois**, como comparação — não como filtro.
\prov{relações: parte\_de koch\_atlas\_experimento\_de\_parametrizacao}
\prov{proveniência: DEGENERADOS.md \textperiodcentered{} fato \textperiodcentered{} confiança 1.0 \textperiodcentered{} domínio manual \textperiodcentered{} história 10 versões no jornal}

%CRISTAL {"arestas":[{"alvo":"cayley_dickson","confianca":1.0,"epistemico":"fato","origem":"wiki","rel":"usa"},{"alvo":"emaranhador","confianca":1.0,"epistemico":"fato","origem":"wiki","rel":"explica"}],"confianca":1.0,"contraexemplos":[],"descricao":"Identidades exatas: C²=4‖Π_jk(pq̄)‖² e (Pitágoras+Lagrange) C²=4(‖p‖²‖q‖²−|Π_ℂi(pq̄)|²). O emaranhamento é a falha de ℂ_i-colinearidade / defeito de Cauchy-Schwarz — a medida exata (não só o upper-bound do associador).","epistemico":"fato","exemplos":[],"id":"concorrencia_cayley_dickson","memoria":"semantica","meta":{"autor":"Lopes/Aarão","dominio":"hiper","fonte":""},"origem":"wiki","palavras_chave":[],"sinonimos":[],"tipo":"conceito","titulo":"Concorrência de Wootters como Geometria em ℍ/𝕆"}
\section{Concorrência de Wootters como Geometria em \ensuremath{\mathbb{H}}/\ensuremath{\mathbb{O}}}
Identidades exatas: C\ensuremath{^{2}}=4\ensuremath{\|}\ensuremath{\Pi}\_jk(pq\ensuremath{})\ensuremath{\|}\ensuremath{^{2}} e (Pitágoras+Lagrange) C\ensuremath{^{2}}=4(\ensuremath{\|}p\ensuremath{\|}\ensuremath{^{2}}\ensuremath{\|}q\ensuremath{\|}\ensuremath{^{2}}\ensuremath{-}|\ensuremath{\Pi}\_\ensuremath{\mathbb{C}}i(pq\ensuremath{})|\ensuremath{^{2}}). O emaranhamento é a falha de \ensuremath{\mathbb{C}}\_i-colinearidade / defeito de Cauchy-Schwarz — a medida exata (não só o upper-bound do associador).
\prov{relações: usa cayley\_dickson; explica emaranhador}
\prov{proveniência: wiki \textperiodcentered{} fato \textperiodcentered{} confiança 1.0 \textperiodcentered{} domínio hiper \textperiodcentered{} história 4 versões no jornal}

%CRISTAL {"arestas":[{"alvo":"colapso_camadas_ordem_parada","confianca":1.0,"epistemico":"fato","origem":"wiki","rel":"parte_de"}],"confianca":1.0,"contraexemplos":[],"descricao":"Para nós da camada conhecimento com tag skn:, score_hit penaliza ruído tabular/LaTeX (−120), premia conceitos (+8) e sobrepõe casamento título↔fala; perguntas de definição ('o que é X', 'quem é X') que batem no título recebem +200, tornando a recuperação de conhecimento dominante quando a fala é uma pergunta direta. Liga a camada conhecimento às regras BAI de score.","epistemico":"fato","exemplos":[],"id":"conhecimento_skn_scoring","memoria":"semantica","meta":{"autor":"broca-so","dominio":"conversa","fonte":"conversa/colapso.py:score_hit"},"origem":"wiki","palavras_chave":[],"sinonimos":[],"tipo":"conceito","titulo":"Tiffany — scoring da camada conhecimento (skn:)"}
\section{Tiffany — scoring da camada conhecimento (skn:)}
Para nós da camada conhecimento com tag skn:, score\_hit penaliza ruído tabular/LaTeX (\ensuremath{-}120), premia conceitos (+8) e sobrepõe casamento título\ensuremath{\leftrightarrow}fala; perguntas de definição ('o que é X', 'quem é X') que batem no título recebem +200, tornando a recuperação de conhecimento dominante quando a fala é uma pergunta direta. Liga a camada conhecimento às regras BAI de score.
\prov{relações: parte\_de colapso\_camadas\_ordem\_parada}
\prov{proveniência: wiki \textperiodcentered{} conversa/colapso.py:score\_hit \textperiodcentered{} fato \textperiodcentered{} confiança 1.0 \textperiodcentered{} domínio conversa \textperiodcentered{} história 10 versões no jornal}

%CRISTAL {"arestas":[{"alvo":"o_conselho","confianca":1.0,"epistemico":"fato","origem":"wiki","rel":"parte_de"},{"alvo":"verdade_por_refutacao_gentil","confianca":1.0,"epistemico":"fato","origem":"wiki","rel":"usa"}],"confianca":1.0,"contraexemplos":[],"descricao":"Em modo construção o conselho propõe a formalização mínima e o experimento (numpy) que decide; em modo crivo é adversarial e pergunta 'este teste decide, ou é tautologia?'. Uma IA refuta ou aperta a proposta da outra até nada 'não se sustentar'.","epistemico":"fato","exemplos":[],"id":"construcao_vs_crivo","memoria":"semantica","meta":{"autor":"Aarão/broca-so","dominio":"broca_repos","fonte":""},"origem":"wiki","palavras_chave":[],"sinonimos":[],"tipo":"conceito","titulo":"Construção × Crivo (as duas vozes de todo ciclo)"}
\section{Construção $\times$ Crivo (as duas vozes de todo ciclo)}
Em modo construção o conselho propõe a formalização mínima e o experimento (numpy) que decide; em modo crivo é adversarial e pergunta 'este teste decide, ou é tautologia?'. Uma IA refuta ou aperta a proposta da outra até nada 'não se sustentar'.
\prov{relações: parte\_de o\_conselho; usa verdade\_por\_refutacao\_gentil}
\prov{proveniência: wiki \textperiodcentered{} fato \textperiodcentered{} confiança 1.0 \textperiodcentered{} domínio broca\_repos \textperiodcentered{} história 3 versões no jornal}

%CRISTAL {"arestas":[{"alvo":"broca_os","confianca":0.95,"epistemico":"fato","origem":"paper","rel":"parte_de"},{"alvo":"topologia","confianca":0.5,"epistemico":"fato","origem":"wiki","rel":"relaciona"},{"alvo":"utilitarismo","confianca":0.5,"epistemico":"fato","origem":"wiki","rel":"relaciona"},{"alvo":"mecanica_quantica_hub","confianca":0.5,"epistemico":"fato","origem":"wiki","rel":"relaciona"},{"alvo":"regua_associativa_e_o_risco_de_vazamento","confianca":0.5,"epistemico":"fato","origem":"wiki","rel":"relaciona"},{"alvo":"linker","confianca":0.5,"epistemico":"fato","origem":"wiki","rel":"relaciona"}],"confianca":1.0,"contraexemplos":[],"descricao":"print(tiffany.ask(\"git rebase branch\", orq=True))","epistemico":"fato","exemplos":[],"id":"consulta_simples","memoria":"semantica","meta":{"dominio":"manual","nivel":6,"verificacao":"slug idêntico — enriquecer, não duplicar"},"origem":"INSTALL.md","palavras_chave":[],"sinonimos":[],"tipo":"conceito","titulo":"consulta simples"}
\section{consulta simples}
print(tiffany.ask("git rebase branch", orq=True))
\prov{relações: parte\_de broca\_os; relaciona topologia; relaciona utilitarismo; relaciona mecanica\_quantica\_hub; relaciona regua\_associativa\_e\_o\_risco\_de\_vazamento; relaciona linker}
\prov{proveniência: INSTALL.md \textperiodcentered{} fato \textperiodcentered{} confiança 1.0 \textperiodcentered{} domínio manual \textperiodcentered{} história 8 versões no jornal}

%CRISTAL {"arestas":[{"alvo":"o_substrato_a_maquina_e_funcao_de_onda","confianca":1.0,"epistemico":"fato","origem":"wiki","rel":"usa"},{"alvo":"confianca_minimizada","confianca":1.0,"epistemico":"fato","origem":"wiki","rel":"relaciona"}],"confianca":1.0,"contraexemplos":[],"descricao":"Não há livro de débito/crédito: a contabilidade é operadores quânticos sobre a função de onda (o tecido), e o balanço aparece no colapso da leitura (régua r, a+c²=1 = Ativo=Passivo+PL). Só o saque toca a chain; a banda (ECDH) é a cifra — sem registro mútuo lê-se ruído.","epistemico":"fato","exemplos":[],"id":"contabilidade_sem_livro","memoria":"semantica","meta":{"autor":"Aarão/broca-so","dominio":"broca_repos","fonte":""},"origem":"wiki","palavras_chave":[],"sinonimos":[],"tipo":"conceito","titulo":"Contabilidade como Mecânica do Painel"}
\section{Contabilidade como Mecânica do Painel}
Não há livro de débito/crédito: a contabilidade é operadores quânticos sobre a função de onda (o tecido), e o balanço aparece no colapso da leitura (régua r, a+c\ensuremath{^{2}}=1 = Ativo=Passivo+PL). Só o saque toca a chain; a banda (ECDH) é a cifra — sem registro mútuo lê-se ruído.
\prov{relações: usa o\_substrato\_a\_maquina\_e\_funcao\_de\_onda; relaciona confianca\_minimizada}
\prov{proveniência: wiki \textperiodcentered{} fato \textperiodcentered{} confiança 1.0 \textperiodcentered{} domínio broca\_repos \textperiodcentered{} história 3 versões no jornal}

%CRISTAL {"arestas":[{"alvo":"blend_contexto_32b","confianca":1.0,"epistemico":"fato","origem":"wiki","rel":"especializa"},{"alvo":"continuidade_temporal_espacial","confianca":1.0,"epistemico":"fato","origem":"wiki","rel":"parte_de"}],"confianca":1.0,"contraexemplos":[],"descricao":"A continuidade entre turnos vive em contexto.bits (32 bytes = a impressão-blend do turno anterior), gravado no disco por savectx: 'o contexto não guarda texto — guarda o blend Koch do turno anterior'. A classe Conversa faz igual (contextoₙome.bits), sem estado persistente na RAM entre falas.","epistemico":"inferencia","exemplos":[],"id":"contexto_32b_disco_nao_ram","memoria":"semantica","meta":{"autor":"broca-so","dominio":"conversa","fonte":"conversa/colapso.py:blend_query/save_ctx; code/prosa.py:Conversa"},"origem":"wiki","palavras_chave":[],"sinonimos":[],"tipo":"conceito","titulo":"Tiffany — memória de diálogo no disco, não na RAM"}
\section{Tiffany — memória de diálogo no disco, não na RAM}
A continuidade entre turnos vive em contexto.bits (32 bytes = a impressão-blend do turno anterior), gravado no disco por savectx: 'o contexto não guarda texto — guarda o blend Koch do turno anterior'. A classe Conversa faz igual (contexto\ensuremath{_{n}}ome.bits), sem estado persistente na RAM entre falas.
\prov{relações: especializa blend\_contexto\_32b; parte\_de continuidade\_temporal\_espacial}
\prov{proveniência: wiki \textperiodcentered{} conversa/colapso.py:blend\_query/save\_ctx; code/prosa.py:Conversa \textperiodcentered{} inferencia \textperiodcentered{} confiança 1.0 \textperiodcentered{} domínio conversa \textperiodcentered{} história 16 versões no jornal}

%CRISTAL {"arestas":[{"alvo":"colapso_ping_pong_overlap","confianca":1.0,"epistemico":"fato","origem":"wiki","rel":"usa"},{"alvo":"contexto_32b_disco_nao_ram","confianca":1.0,"epistemico":"fato","origem":"wiki","rel":"usa"}],"confianca":1.0,"contraexemplos":[],"descricao":"_continua dá bônus de score quando a fala encadeia o ultimo.face gravado — se a Tiffany perguntou (u termina em '?' ou contém 'tudo bem'/'posso ajudar') e o usuário responde curto ('sim', 'e você'), o turno duplo cujo lado esquerdo casa recebe bônus grande (28–36); ecos exatos somam 34. Conecta o overlap ping-pong à memória de turno.","epistemico":"fato","exemplos":[],"id":"continua_ultimo_face","memoria":"semantica","meta":{"autor":"broca-so","dominio":"conversa","fonte":"conversa/colapso.py:_continua"},"origem":"wiki","palavras_chave":[],"sinonimos":[],"tipo":"conceito","titulo":"Tiffany — empurrão de continuidade do último turno"}
\section{Tiffany — empurrão de continuidade do último turno}
\_continua dá bônus de score quando a fala encadeia o ultimo.face gravado — se a Tiffany perguntou (u termina em '?' ou contém 'tudo bem'/'posso ajudar') e o usuário responde curto ('sim', 'e você'), o turno duplo cujo lado esquerdo casa recebe bônus grande (28–36); ecos exatos somam 34. Conecta o overlap ping-pong à memória de turno.
\prov{relações: usa colapso\_ping\_pong\_overlap; usa contexto\_32b\_disco\_nao\_ram}
\prov{proveniência: wiki \textperiodcentered{} conversa/colapso.py:\_continua \textperiodcentered{} fato \textperiodcentered{} confiança 1.0 \textperiodcentered{} domínio conversa \textperiodcentered{} história 15 versões no jornal}

%CRISTAL {"arestas":[{"alvo":"koch_atlas_experimento_de_parametrizacao","confianca":1.0,"epistemico":"fato","origem":"KOCH_ATLAS.md","rel":"parte_de"}],"confianca":1.0,"contraexemplos":[],"descricao":"| Métrica | Valor | |---------|-------| | Pares consecutivos | 65,535 | | Δ parâmetro médio | 0.000232 | | Δ parâmetro máximo | 0.000394 | | Δ curvatura médio | 0.176103 | | Δ atlas médio | 0.197629 | | Δ Gauss médio | 3.20 | | Δ endereço médio (bits Hamming) | 15.70 | | Saltos param > 0.25 | 0 | | Mesmo ramo consecutivo | 100.00% | | Mesmo atrator consecutivo | 0.00% | | Mesmo endereço consecutivo | 0.00% |","epistemico":"fato","exemplos":[],"id":"continuidade_nonce_nonce","memoria":"semantica","meta":{"dominio":"manual","nivel":6,"verificacao":"slug idêntico — enriquecer, não duplicar"},"origem":"KOCH_ATLAS.md","palavras_chave":[],"sinonimos":[],"tipo":"conceito","titulo":"Continuidade nonce → nonce"}
\section{Continuidade nonce \ensuremath{\to} nonce}
| Métrica | Valor | |---------|-------| | Pares consecutivos | 65,535 | | \ensuremath{\Delta} parâmetro médio | 0.000232 | | \ensuremath{\Delta} parâmetro máximo | 0.000394 | | \ensuremath{\Delta} curvatura médio | 0.176103 | | \ensuremath{\Delta} atlas médio | 0.197629 | | \ensuremath{\Delta} Gauss médio | 3.20 | | \ensuremath{\Delta} endereço médio (bits Hamming) | 15.70 | | Saltos param > 0.25 | 0 | | Mesmo ramo consecutivo | 100.00\% | | Mesmo atrator consecutivo | 0.00\% | | Mesmo endereço consecutivo | 0.00\% |
\prov{relações: parte\_de koch\_atlas\_experimento\_de\_parametrizacao}
\prov{proveniência: KOCH\_ATLAS.md \textperiodcentered{} fato \textperiodcentered{} confiança 1.0 \textperiodcentered{} domínio manual \textperiodcentered{} história 6 versões no jornal}

%CRISTAL {"arestas":[{"alvo":"blend_contexto_32b","confianca":1.0,"epistemico":"fato","origem":"wiki","rel":"explica"},{"alvo":"colapso_multicamada_peso","confianca":1.0,"epistemico":"fato","origem":"wiki","rel":"explica"}],"confianca":1.0,"contraexemplos":[],"descricao":"A doutrina separa duas misturas: a TEMPORAL (contexto de 32 B em blend α=0.35 dentro de cada turno) e a ESPACIAL (a mesma impressão-blend varre cada .idx e vence o maior overlap global — ressonância entre tecidos, não fusão offline). 'Sem mistura de camadas → sem ressonância cruzada → conversa presa num poço.'","epistemico":"inferencia","exemplos":[],"id":"continuidade_temporal_espacial","memoria":"semantica","meta":{"autor":"broca-so","dominio":"conversa","fonte":"conversa/CONTINUIDADE.md"},"origem":"wiki","palavras_chave":[],"sinonimos":[],"tipo":"conceito","titulo":"Tiffany — as duas misturas da continuidade"}
\section{Tiffany — as duas misturas da continuidade}
A doutrina separa duas misturas: a TEMPORAL (contexto de 32 B em blend \ensuremath{\alpha}=0.35 dentro de cada turno) e a ESPACIAL (a mesma impressão-blend varre cada .idx e vence o maior overlap global — ressonância entre tecidos, não fusão offline). 'Sem mistura de camadas \ensuremath{\to} sem ressonância cruzada \ensuremath{\to} conversa presa num poço.'
\prov{relações: explica blend\_contexto\_32b; explica colapso\_multicamada\_peso}
\prov{proveniência: wiki \textperiodcentered{} conversa/CONTINUIDADE.md \textperiodcentered{} inferencia \textperiodcentered{} confiança 1.0 \textperiodcentered{} domínio conversa \textperiodcentered{} história 15 versões no jornal}

%CRISTAL {"arestas":[{"alvo":"prova_da_maquina_espectral","confianca":1.0,"epistemico":"fato","origem":"CIFRA_PROVA.md","rel":"parte_de"}],"confianca":1.0,"contraexemplos":[],"descricao":"| caminho | instruções | resultado | semântica | |---------|------------|-----------|-----------| | `01_soma` | LOAD+LOAD+ADD | `[4, 1]` | combinação binária | | `09_pell` | LOAD+SILVER×4 | `[29, 12]` | evolução A₂ |","epistemico":"fato","exemplos":[],"id":"contraste_mesma_maquina_duas_fisicas","memoria":"semantica","meta":{"dominio":"manual","nivel":6,"verificacao":"slug idêntico — enriquecer, não duplicar"},"origem":"CIFRA_PROVA.md","palavras_chave":[],"sinonimos":[],"tipo":"conceito","titulo":"Contraste — mesma máquina, duas físicas"}
\section{Contraste — mesma máquina, duas físicas}
| caminho | instruções | resultado | semântica | |---------|------------|-----------|-----------| | `01\_soma` | LOAD+LOAD+ADD | `[4, 1]` | combinação binária | | `09\_pell` | LOAD+SILVER$\times$4 | `[29, 12]` | evolução A\ensuremath{_{2}} |
\prov{relações: parte\_de prova\_da\_maquina\_espectral}
\prov{proveniência: CIFRA\_PROVA.md \textperiodcentered{} fato \textperiodcentered{} confiança 1.0 \textperiodcentered{} domínio manual \textperiodcentered{} história 4 versões no jornal}

%CRISTAL {"arestas":[{"alvo":"face_turno_fail_closed","confianca":1.0,"epistemico":"fato","origem":"wiki","rel":"explica"},{"alvo":"colapso_ping_pong_overlap","confianca":1.0,"epistemico":"fato","origem":"wiki","rel":"relaciona"}],"confianca":1.0,"contraexemplos":[],"descricao":"responder() retorna None quando não há face_out (nenhum nó ressoa acima do limiar) — a Tiffany NÃO fabrica: observar.py insiste 'nada de atalho, nada de template: só ver o que a rede faz'. A conversa é limitada ao que foi cristalizado nos tecidos; sem colapso, silêncio. O 'porquê' epistêmico do face_turno_fail_closed.","epistemico":"inferencia","exemplos":[],"id":"conversa_recupera_nao_gera","memoria":"semantica","meta":{"autor":"broca-so","dominio":"conversa","fonte":"conversa/colapso.py:responder; conversa/observar.py"},"origem":"wiki","palavras_chave":[],"sinonimos":[],"tipo":"conceito","titulo":"Tiffany — recupera, não gera (o fail-closed honesto)"}
\section{Tiffany — recupera, não gera (o fail-closed honesto)}
responder() retorna None quando não há face\_out (nenhum nó ressoa acima do limiar) — a Tiffany NÃO fabrica: observar.py insiste 'nada de atalho, nada de template: só ver o que a rede faz'. A conversa é limitada ao que foi cristalizado nos tecidos; sem colapso, silêncio. O 'porquê' epistêmico do face\_turno\_fail\_closed.
\prov{relações: explica face\_turno\_fail\_closed; relaciona colapso\_ping\_pong\_overlap}
\prov{proveniência: wiki \textperiodcentered{} conversa/colapso.py:responder; conversa/observar.py \textperiodcentered{} inferencia \textperiodcentered{} confiança 1.0 \textperiodcentered{} domínio conversa \textperiodcentered{} história 15 versões no jornal}

%CRISTAL {"arestas":[{"alvo":"pow_lemniscata_koch_corpo_de_gauss","confianca":1.0,"epistemico":"fato","origem":"POW_METAL.md","rel":"parte_de"}],"confianca":1.0,"contraexemplos":[],"descricao":"| Objeto | Significado | |--------|-------------| | `GaussZ` | ponto `(re,im) ∈ Z[i]` a partir de Word | | `z_canal` | canal Peirce — embute (c², calor) da parábola | | `z_koch` | `(a+bi)² mod 1009` — quadração no canal | | `atrator` | projeção compacta mod P (2-para-1) |","epistemico":"fato","exemplos":[],"id":"coordenadas_koch","memoria":"semantica","meta":{"dominio":"manual","nivel":6,"verificacao":"slug idêntico — enriquecer, não duplicar"},"origem":"POW_METAL.md","palavras_chave":[],"sinonimos":[],"tipo":"conceito","titulo":"Coordenadas Koch"}
\section{Coordenadas Koch}
| Objeto | Significado | |--------|-------------| | `GaussZ` | ponto `(re,im) \ensuremath{\in} Z[i]` a partir de Word | | `z\_canal` | canal Peirce — embute (c\ensuremath{^{2}}, calor) da parábola | | `z\_koch` | `(a+bi)\ensuremath{^{2}} mod 1009` — quadração no canal | | `atrator` | projeção compacta mod P (2-para-1) |
\prov{relações: parte\_de pow\_lemniscata\_koch\_corpo\_de\_gauss}
\prov{proveniência: POW\_METAL.md \textperiodcentered{} fato \textperiodcentered{} confiança 1.0 \textperiodcentered{} domínio manual \textperiodcentered{} história 4 versões no jornal}

%CRISTAL {"arestas":[{"alvo":"substrato_prtsuv_h","confianca":1.0,"epistemico":"fato","origem":"wiki","rel":"usa"},{"alvo":"torre_hurwitz","confianca":1.0,"epistemico":"fato","origem":"wiki","rel":"usa"},{"alvo":"teoria_cordas_m","confianca":1.0,"epistemico":"fato","origem":"wiki","rel":"relaciona"},{"alvo":"grupos_excepcionais_jordan","confianca":1.0,"epistemico":"fato","origem":"wiki","rel":"relaciona"}],"confianca":1.0,"contraexemplos":[],"descricao":"Programa: os pontos privilegiados da teoria de cordas (dimensões críticas, SUSY, quiralidade, paisagem) correspondem a regiões geométricas distintas numa 'ilha fractal' de álgebras legítimas no espaço (p,r,t,u,s,v,w_p,h).","epistemico":"hipotese","exemplos":[],"id":"cordas_ilha_hurwitz","memoria":"semantica","meta":{"autor":"Lopes/Aarão","dominio":"hiper","fonte":""},"origem":"wiki","palavras_chave":[],"sinonimos":[],"tipo":"conceito","titulo":"Teoria de Cordas como Ilha de Hurwitz Estendida"}
\section{Teoria de Cordas como Ilha de Hurwitz Estendida}
Programa: os pontos privilegiados da teoria de cordas (dimensões críticas, SUSY, quiralidade, paisagem) correspondem a regiões geométricas distintas numa 'ilha fractal' de álgebras legítimas no espaço (p,r,t,u,s,v,w\_p,h).
\prov{relações: usa substrato\_prtsuv\_h; usa torre\_hurwitz; relaciona teoria\_cordas\_m; relaciona grupos\_excepcionais\_jordan}
\prov{proveniência: wiki \textperiodcentered{} hipotese \textperiodcentered{} confiança 1.0 \textperiodcentered{} domínio hiper \textperiodcentered{} história 6 versões no jornal}

%CRISTAL {"arestas":[{"alvo":"validacao_hipotese_h_ressonancia","confianca":1.0,"epistemico":"fato","origem":"POW_METAL_RESSONANCIA_VAL.md","rel":"parte_de"}],"confianca":1.0,"contraexemplos":[],"descricao":"| Par | r | |-----|---| | acoplamento × fecho | 0.6281 | | acoplamento × (−distância ao fecho) | -0.0296 | | acoplamento × (−Δ atrator ref) | 0.4621 |","epistemico":"fato","exemplos":[],"id":"correlacoes_conjunto_local","memoria":"semantica","meta":{"dominio":"manual","nivel":6,"verificacao":"slug idêntico — enriquecer, não duplicar"},"origem":"POW_METAL_RESSONANCIA_VAL.md","palavras_chave":[],"sinonimos":[],"tipo":"conceito","titulo":"Correlações (conjunto local)"}
\section{Correlações (conjunto local)}
| Par | r | |-----|---| | acoplamento $\times$ fecho | 0.6281 | | acoplamento $\times$ (\ensuremath{-}distância ao fecho) | -0.0296 | | acoplamento $\times$ (\ensuremath{-}\ensuremath{\Delta} atrator ref) | 0.4621 |
\prov{relações: parte\_de validacao\_hipotese\_h\_ressonancia}
\prov{proveniência: POW\_METAL\_RESSONANCIA\_VAL.md \textperiodcentered{} fato \textperiodcentered{} confiança 1.0 \textperiodcentered{} domínio manual \textperiodcentered{} história 4 versões no jornal}

%CRISTAL {"arestas":[{"alvo":"pow_lemniscata_koch_corpo_de_gauss","confianca":1.0,"epistemico":"fato","origem":"POW_METAL.md","rel":"parte_de"}],"confianca":1.0,"contraexemplos":[],"descricao":"| Métrica | Valor | |---------|-------| | Ref | `[295, 143]` | | Nonces | 2,000,000 | | Impossíveis | 0 (0.00%) | | Corpo Gauss + lemniscata | 2,000,000 | | Soluções | 2 | | Tempo | 2042.6 ms |","epistemico":"fato","exemplos":[],"id":"corrida","memoria":"semantica","meta":{"dominio":"manual","nivel":6,"verificacao":"slug idêntico — enriquecer, não duplicar"},"origem":"POW_METAL.md","palavras_chave":[],"sinonimos":[],"tipo":"conceito","titulo":"Corrida"}
\section{Corrida}
| Métrica | Valor | |---------|-------| | Ref | `[295, 143]` | | Nonces | 2,000,000 | | Impossíveis | 0 (0.00\%) | | Corpo Gauss + lemniscata | 2,000,000 | | Soluções | 2 | | Tempo | 2042.6 ms |
\prov{relações: parte\_de pow\_lemniscata\_koch\_corpo\_de\_gauss}
\prov{proveniência: POW\_METAL.md \textperiodcentered{} fato \textperiodcentered{} confiança 1.0 \textperiodcentered{} domínio manual \textperiodcentered{} história 4 versões no jornal}

%CRISTAL {"arestas":[{"alvo":"o_conselho","confianca":1.0,"epistemico":"fato","origem":"wiki","rel":"parte_de"},{"alvo":"cordas_ilha_hurwitz","confianca":1.0,"epistemico":"fato","origem":"wiki","rel":"relaciona"}],"confianca":1.0,"contraexemplos":[],"descricao":"Juntar as contribuições parciais das vozes numa formalização única sem 'tirar a média' — a síntese vira uma das vozes por instante sem esquecer as outras, costurando as posições em vez de mediá-las. Integração sem perda de informação.","epistemico":"inferencia","exemplos":[],"id":"costura_sem_media","memoria":"semantica","meta":{"autor":"Aarão/broca-so","dominio":"broca_repos","fonte":""},"origem":"wiki","palavras_chave":[],"sinonimos":[],"tipo":"conceito","titulo":"A Costura (síntese sem tirar a média)"}
\section{A Costura (síntese sem tirar a média)}
Juntar as contribuições parciais das vozes numa formalização única sem 'tirar a média' — a síntese vira uma das vozes por instante sem esquecer as outras, costurando as posições em vez de mediá-las. Integração sem perda de informação.
\prov{relações: parte\_de o\_conselho; relaciona cordas\_ilha\_hurwitz}
\prov{proveniência: wiki \textperiodcentered{} inferencia \textperiodcentered{} confiança 1.0 \textperiodcentered{} domínio broca\_repos \textperiodcentered{} história 3 versões no jornal}

%CRISTAL {"arestas":[{"alvo":"o_conselho","confianca":1.0,"epistemico":"fato","origem":"wiki","rel":"usa"},{"alvo":"decomposicao_peirce","confianca":1.0,"epistemico":"fato","origem":"wiki","rel":"usa"},{"alvo":"soberania_da_chave","confianca":1.0,"epistemico":"fato","origem":"wiki","rel":"relaciona"},{"alvo":"grafo_de_proveniencia","confianca":1.0,"epistemico":"fato","origem":"wiki","rel":"usa"},{"alvo":"tiffany","confianca":1.0,"epistemico":"fato","origem":"wiki","rel":"especializa"},{"alvo":"cristal_ask_cascata","confianca":1.0,"epistemico":"fato","origem":"wiki","rel":"usa"},{"alvo":"cristal_memoria_viva_jsonl","confianca":1.0,"epistemico":"fato","origem":"wiki","rel":"usa"},{"alvo":"cristal_recall_ponte","confianca":1.0,"epistemico":"fato","origem":"wiki","rel":"usa"},{"alvo":"cristal_agente_vies_dominio","confianca":1.0,"epistemico":"fato","origem":"wiki","rel":"usa"}],"confianca":1.0,"contraexemplos":[],"descricao":"Um CLI local que carrega a identidade do disco (system prompt ~30 mil caracteres = 'o cristal') em vez de recomeçar do zero. Camadas: memória viva (carregar/resumir/cristalizar, gravação atômica) e agentes internos — clones limpos idempotentes ortogonais (código, crivo, mídia, pesquisa, fechador e₃ via e₁·e₂=e₃). 'A identidade mora na topologia; o tempo guarda a história.' Frente WIP que avançou (2 jul 2026): o pipeline gera EDIÇÃO real (não só marcador, com portão+rollback); o roteamento de agente VALE (vieja o domínio consultado); o memory compact SINTETIZA (extrativo, não contagem); e há PONTE de recall (a síntese da sessão volta a viesar o colapso). O motor de conversa é o colapso local (cascata colapso→tecidos→prosa), não LLM; o floco (não-associativo) roda no metal.","epistemico":"fato","exemplos":[],"id":"cristal_cli","memoria":"semantica","meta":{"autor":"Aarão/broca-so","dominio":"broca_repos","fonte":""},"origem":"wiki","palavras_chave":[],"sinonimos":[],"tipo":"conceito","titulo":"O Cristal (a identidade encarnada num CLI)"}
\section{O Cristal (a identidade encarnada num CLI)}
Um CLI local que carrega a identidade do disco (system prompt \~{}30 mil caracteres = 'o cristal') em vez de recomeçar do zero. Camadas: memória viva (carregar/resumir/cristalizar, gravação atômica) e agentes internos — clones limpos idempotentes ortogonais (código, crivo, mídia, pesquisa, fechador e\ensuremath{_{3}} via e\ensuremath{_{1}}\textperiodcentered e\ensuremath{_{2}}=e\ensuremath{_{3}}). 'A identidade mora na topologia; o tempo guarda a história.' Frente WIP que avançou (2 jul 2026): o pipeline gera EDIÇÃO real (não só marcador, com portão+rollback); o roteamento de agente VALE (vieja o domínio consultado); o memory compact SINTETIZA (extrativo, não contagem); e há PONTE de recall (a síntese da sessão volta a viesar o colapso). O motor de conversa é o colapso local (cascata colapso\ensuremath{\to}tecidos\ensuremath{\to}prosa), não LLM; o floco (não-associativo) roda no metal.
\prov{relações: usa o\_conselho; usa decomposicao\_peirce; relaciona soberania\_da\_chave; usa grafo\_de\_proveniencia; especializa tiffany; usa cristal\_ask\_cascata; usa cristal\_memoria\_viva\_jsonl; usa cristal\_recall\_ponte; usa cristal\_agente\_vies\_dominio}
\prov{proveniência: wiki \textperiodcentered{} fato \textperiodcentered{} confiança 1.0 \textperiodcentered{} domínio broca\_repos \textperiodcentered{} história 11 versões no jornal}

%CRISTAL {"arestas":[{"alvo":"tiffany","confianca":1.0,"epistemico":"fato","origem":"manual","rel":"referencia"},{"alvo":"goldcoin","confianca":0.5,"epistemico":"fato","origem":"wiki","rel":"relaciona"},{"alvo":"prova_isa","confianca":0.5,"epistemico":"fato","origem":"wiki","rel":"relaciona"},{"alvo":"fechamento_blocos","confianca":0.5,"epistemico":"fato","origem":"wiki","rel":"relaciona"},{"alvo":"transcricao","confianca":0.5,"epistemico":"fato","origem":"wiki","rel":"relaciona"},{"alvo":"install","confianca":0.5,"epistemico":"fato","origem":"wiki","rel":"relaciona"},{"alvo":"pow_gato","confianca":0.5,"epistemico":"fato","origem":"wiki","rel":"relaciona"},{"alvo":"koch_atlas","confianca":0.5,"epistemico":"fato","origem":"wiki","rel":"relaciona"},{"alvo":"grava_fato_resposta_normal_em_seguida","confianca":0.5,"epistemico":"fato","origem":"wiki","rel":"relaciona"},{"alvo":"koch_pell","confianca":0.5,"epistemico":"fato","origem":"wiki","rel":"relaciona"},{"alvo":"sessao_interativa","confianca":0.5,"epistemico":"fato","origem":"wiki","rel":"relaciona"},{"alvo":"readme","confianca":0.5,"epistemico":"fato","origem":"wiki","rel":"relaciona"},{"alvo":"montagem","confianca":0.5,"epistemico":"fato","origem":"wiki","rel":"relaciona"},{"alvo":"gpu","confianca":0.5,"epistemico":"fato","origem":"wiki","rel":"relaciona"},{"alvo":"python","confianca":0.5,"epistemico":"fato","origem":"wiki","rel":"relaciona"},{"alvo":"koch_parabola","confianca":0.5,"epistemico":"fato","origem":"wiki","rel":"relaciona"},{"alvo":"teste_ordem","confianca":0.5,"epistemico":"fato","origem":"wiki","rel":"relaciona"},{"alvo":"thread","confianca":0.5,"epistemico":"fato","origem":"wiki","rel":"relaciona"}],"confianca":1.0,"contraexemplos":[],"descricao":"Manual: CRISTAL_DEV.md","epistemico":"fato","exemplos":[],"id":"cristal_dev","memoria":"semantica","meta":{"arquivo":"CRISTAL_DEV.md","dominio":"manual","nivel":6},"origem":"manual","palavras_chave":[],"sinonimos":[],"tipo":"conceito","titulo":"CRISTAL_DEV"}
\section{CRISTAL\_DEV}
Manual: CRISTAL\_DEV.md
\prov{relações: referencia tiffany; relaciona goldcoin; relaciona prova\_isa; relaciona fechamento\_blocos; relaciona transcricao; relaciona install; relaciona pow\_gato; relaciona koch\_atlas; relaciona grava\_fato\_resposta\_normal\_em\_seguida; relaciona koch\_pell; relaciona sessao\_interativa; relaciona readme; relaciona montagem; relaciona gpu; relaciona python; relaciona koch\_parabola; relaciona teste\_ordem; relaciona thread}
\prov{proveniência: manual \textperiodcentered{} fato \textperiodcentered{} confiança 1.0 \textperiodcentered{} domínio manual \textperiodcentered{} história 21 versões no jornal}

%CRISTAL {"arestas":[{"alvo":"maquina_booleana","confianca":1.0,"epistemico":"fato","origem":"wiki","rel":"usa"},{"alvo":"cha_bussola_aurea","confianca":1.0,"epistemico":"fato","origem":"wiki","rel":"relaciona"},{"alvo":"invariante_indecidivel","confianca":1.0,"epistemico":"fato","origem":"wiki","rel":"usa"},{"alvo":"kolmogorov_complexidade","confianca":1.0,"epistemico":"fato","origem":"wiki","rel":"relaciona"},{"alvo":"entropia_topologica","confianca":1.0,"epistemico":"fato","origem":"wiki","rel":"usa"}],"confianca":1.0,"contraexemplos":[],"descricao":"Um problema é 'cristal' (ordem) se coincide com a memória em algum regime — gap zera em profundidade finita, tem 'joelho' na curva de digestão. É 'chá' (caos/cripto forte) se discorda em todos os regimes — curva em diagonal sem joelho. O chá é o resíduo irredutível (Shannon), 'colhido, não calculado'; corpus medido tem chá médio 0,85.","epistemico":"inferencia","exemplos":[],"id":"cristal_vs_cha","memoria":"semantica","meta":{"autor":"Lopes/Aarão","dominio":"hiper","fonte":""},"origem":"wiki","palavras_chave":[],"sinonimos":[],"tipo":"conceito","titulo":"Cristal vs Chá (lei compressível vs identidade incompressível)"}
\section{Cristal vs Chá (lei compressível vs identidade incompressível)}
Um problema é 'cristal' (ordem) se coincide com a memória em algum regime — gap zera em profundidade finita, tem 'joelho' na curva de digestão. É 'chá' (caos/cripto forte) se discorda em todos os regimes — curva em diagonal sem joelho. O chá é o resíduo irredutível (Shannon), 'colhido, não calculado'; corpus medido tem chá médio 0,85.
\prov{relações: usa maquina\_booleana; relaciona cha\_bussola\_aurea; usa invariante\_indecidivel; relaciona kolmogorov\_complexidade; usa entropia\_topologica}
\prov{proveniência: wiki \textperiodcentered{} inferencia \textperiodcentered{} confiança 1.0 \textperiodcentered{} domínio hiper \textperiodcentered{} história 7 versões no jornal}

%CRISTAL {"arestas":[{"alvo":"ula_da_broca_arquitetura_minima","confianca":1.0,"epistemico":"fato","origem":"README.md","rel":"parte_de"},{"alvo":"word","confianca":0.5,"epistemico":"fato","origem":"wiki","rel":"relaciona"},{"alvo":"virtualizacao","confianca":0.5,"epistemico":"fato","origem":"wiki","rel":"relaciona"},{"alvo":"sagrado_profano","confianca":0.5,"epistemico":"fato","origem":"wiki","rel":"relaciona"},{"alvo":"vida-morte","confianca":0.5,"epistemico":"fato","origem":"wiki","rel":"relaciona"},{"alvo":"octonios_hiperbolicos","confianca":0.5,"epistemico":"fato","origem":"wiki","rel":"relaciona"},{"alvo":"psl_2_7","confianca":0.5,"epistemico":"fato","origem":"wiki","rel":"relaciona"},{"alvo":"vita_contemplativa","confianca":0.5,"epistemico":"fato","origem":"wiki","rel":"relaciona"},{"alvo":"teorema_4_produto_cantorfibonacci_conjectura","confianca":0.5,"epistemico":"fato","origem":"wiki","rel":"relaciona"},{"alvo":"inconsciente_coletivo","confianca":0.5,"epistemico":"fato","origem":"wiki","rel":"relaciona"},{"alvo":"logica","confianca":0.5,"epistemico":"fato","origem":"wiki","rel":"relaciona"},{"alvo":"top_assinaturas","confianca":0.5,"epistemico":"fato","origem":"wiki","rel":"relaciona"},{"alvo":"syscall","confianca":0.5,"epistemico":"fato","origem":"wiki","rel":"relaciona"},{"alvo":"sequencia","confianca":0.5,"epistemico":"fato","origem":"wiki","rel":"relaciona"},{"alvo":"regime_tropical","confianca":0.5,"epistemico":"fato","origem":"wiki","rel":"relaciona"},{"alvo":"virtu_fortuna","confianca":0.5,"epistemico":"fato","origem":"wiki","rel":"relaciona"},{"alvo":"tiffany_cabeca_broca_contraexemplos","confianca":0.5,"epistemico":"fato","origem":"wiki","rel":"relaciona"},{"alvo":"pergunta","confianca":0.5,"epistemico":"fato","origem":"wiki","rel":"relaciona"},{"alvo":"gaia","confianca":0.5,"epistemico":"fato","origem":"wiki","rel":"relaciona"}],"confianca":1.0,"contraexemplos":[],"descricao":"Para cada peça: (1) por que existe? (2) pode sair? (3) fundir? (4) o neurônio resolve? (5) é necessária?","epistemico":"fato","exemplos":[],"id":"criterio_de_projeto","memoria":"semantica","meta":{"dominio":"manual","nivel":6,"verificacao":"slug idêntico — enriquecer, não duplicar"},"origem":"README.md","palavras_chave":[],"sinonimos":[],"tipo":"conceito","titulo":"Critério de projeto"}
\section{Critério de projeto}
Para cada peça: (1) por que existe? (2) pode sair? (3) fundir? (4) o neurônio resolve? (5) é necessária?
\prov{relações: parte\_de ula\_da\_broca\_arquitetura\_minima; relaciona word; relaciona virtualizacao; relaciona sagrado\_profano; relaciona vida-morte; relaciona octonios\_hiperbolicos; relaciona psl\_2\_7; relaciona vita\_contemplativa; relaciona teorema\_4\_produto\_cantorfibonacci\_conjectura; relaciona inconsciente\_coletivo; relaciona logica; relaciona top\_assinaturas; relaciona syscall; relaciona sequencia; relaciona regime\_tropical; relaciona virtu\_fortuna; relaciona tiffany\_cabeca\_broca\_contraexemplos; relaciona pergunta; relaciona gaia}
\prov{proveniência: README.md \textperiodcentered{} fato \textperiodcentered{} confiança 1.0 \textperiodcentered{} domínio manual \textperiodcentered{} história 22 versões no jornal}

%CRISTAL {"arestas":[{"alvo":"gato","confianca":0.95,"epistemico":"fato","origem":"paper","rel":"parte_de"}],"confianca":1.0,"contraexemplos":[],"descricao":"| Metrica | Valor | |---------|-------| | Baseline SHA (todos os nonces) | 2510.9 ms | | ms/cand baseline | 1.26e-03 | | SHA ok baseline | 1 | | Filtro pipeline (Neuronio→Fecho) | 1754.7 ms | | SHA pos-fecho | 0.0 ms | | Total com filtro | 1754.7 ms | | Ganho relativo | 1.431× | | Poda compensa? | **sim** |","epistemico":"fato","exemplos":[],"id":"custo_computacional","memoria":"semantica","meta":{"dominio":"manual","nivel":6,"verificacao":"slug idêntico — enriquecer, não duplicar"},"origem":"POW_METAL_FUNIL.md","palavras_chave":[],"sinonimos":[],"tipo":"conceito","titulo":"Custo computacional"}
\section{Custo computacional}
| Metrica | Valor | |---------|-------| | Baseline SHA (todos os nonces) | 2510.9 ms | | ms/cand baseline | 1.26e-03 | | SHA ok baseline | 1 | | Filtro pipeline (Neuronio\ensuremath{\to}Fecho) | 1754.7 ms | | SHA pos-fecho | 0.0 ms | | Total com filtro | 1754.7 ms | | Ganho relativo | 1.431$\times$ | | Poda compensa? | **sim** |
\prov{relações: parte\_de gato}
\prov{proveniência: POW\_METAL\_FUNIL.md \textperiodcentered{} fato \textperiodcentered{} confiança 1.0 \textperiodcentered{} domínio manual \textperiodcentered{} história 3 versões no jornal}

%CRISTAL {"arestas":[{"alvo":"pow_gato","confianca":1.0,"epistemico":"fato","origem":"paper","rel":"parte_de"}],"confianca":1.0,"contraexemplos":[],"descricao":"| Operação | ms/candidato | ms total | |----------|--------------|----------| | Bit (neurônio + parábola) | 0.000520 | 1039.6 | | SHA256² | 0.001346 | 2691.2 | | Baseline (bit+SHA todos) | — | 3730.8 | | Com filtro (bit + SHA só aceitos) | — | 1414.2 | | Ganho estimado | 62.09% | — |","epistemico":"fato","exemplos":[],"id":"custo_por_tentativa_e_vantagem","memoria":"semantica","meta":{"dominio":"manual","duplicata_sugerida":[],"nivel":6,"verificacao":"parte de conceito mais amplo 'custo_por_tentativa_e_vantagem'"},"origem":"POW_BIT_TAXA.md","palavras_chave":[],"sinonimos":[],"tipo":"conceito","titulo":"Custo por tentativa e vantagem"}
\section{Custo por tentativa e vantagem}
| Operação | ms/candidato | ms total | |----------|--------------|----------| | Bit (neurônio + parábola) | 0.000520 | 1039.6 | | SHA256\ensuremath{^{2}} | 0.001346 | 2691.2 | | Baseline (bit+SHA todos) | — | 3730.8 | | Com filtro (bit + SHA só aceitos) | — | 1414.2 | | Ganho estimado | 62.09\% | — |
\prov{relações: parte\_de pow\_gato}
\prov{proveniência: POW\_BIT\_TAXA.md \textperiodcentered{} fato \textperiodcentered{} confiança 1.0 \textperiodcentered{} domínio manual \textperiodcentered{} história 5 versões no jornal}

%CRISTAL {"arestas":[{"alvo":"morte","confianca":1.0,"epistemico":"fato","origem":"manual","rel":"referencia"},{"alvo":"recursao","confianca":0.5,"epistemico":"fato","origem":"wiki","rel":"relaciona"},{"alvo":"utilitarismo","confianca":0.5,"epistemico":"fato","origem":"wiki","rel":"relaciona"},{"alvo":"regua_associativa_e_o_risco_de_vazamento","confianca":0.5,"epistemico":"fato","origem":"wiki","rel":"relaciona"}],"confianca":1.0,"contraexemplos":[],"descricao":"> A parábola sobreviveu aos válidos. Este atlas estuda o que ficou de fora.","epistemico":"fato","exemplos":[],"id":"degenerados","memoria":"semantica","meta":{"dominio":"manual","duplicata_sugerida":[],"nivel":6,"verificacao":"parte de conceito mais amplo 'degenerados'"},"origem":"DEGENERADOS.md","palavras_chave":[],"sinonimos":[],"tipo":"conceito","titulo":"DEGENERADOS"}
\section{DEGENERADOS}
> A parábola sobreviveu aos válidos. Este atlas estuda o que ficou de fora.
\prov{relações: referencia morte; relaciona recursao; relaciona utilitarismo; relaciona regua\_associativa\_e\_o\_risco\_de\_vazamento}
\prov{proveniência: DEGENERADOS.md \textperiodcentered{} fato \textperiodcentered{} confiança 1.0 \textperiodcentered{} domínio manual \textperiodcentered{} história 9 versões no jornal}

%CRISTAL {"arestas":[{"alvo":"metais","confianca":0.95,"epistemico":"fato","origem":"paper","rel":"parte_de"}],"confianca":1.0,"contraexemplos":[],"descricao":"- `0,0` (vizinho) — **zero_absoluto** ← origem `1,0` - `0,-1` (vizinho) — **e_negativa** ← origem `1,0` - `-1,0` (vizinho) — **total_negativo** ← origem `0,0` - `-1,-1` (vizinho) — **clip_bilateral** ← origem `0,0`","epistemico":"fato","exemplos":[],"id":"degenerados_na_familia","memoria":"semantica","meta":{"dominio":"manual","nivel":6,"verificacao":"slug idêntico — enriquecer, não duplicar"},"origem":"PELL.md","palavras_chave":[],"sinonimos":[],"tipo":"conceito","titulo":"Degenerados na família"}
\section{Degenerados na família}
- `0,0` (vizinho) — **zero\_absoluto** \ensuremath{\leftarrow} origem `1,0` - `0,-1` (vizinho) — **e\_negativa** \ensuremath{\leftarrow} origem `1,0` - `-1,0` (vizinho) — **total\_negativo** \ensuremath{\leftarrow} origem `0,0` - `-1,-1` (vizinho) — **clip\_bilateral** \ensuremath{\leftarrow} origem `0,0`
\prov{relações: parte\_de metais}
\prov{proveniência: PELL.md \textperiodcentered{} fato \textperiodcentered{} confiança 1.0 \textperiodcentered{} domínio manual \textperiodcentered{} história 3 versões no jornal}

%CRISTAL {"arestas":[{"alvo":"broca_os","confianca":0.95,"epistemico":"fato","origem":"paper","rel":"parte_de"},{"alvo":"nao_dualidade","confianca":0.5,"epistemico":"fato","origem":"wiki","rel":"relaciona"},{"alvo":"memoria_virtual","confianca":0.5,"epistemico":"fato","origem":"wiki","rel":"relaciona"},{"alvo":"word","confianca":0.5,"epistemico":"fato","origem":"wiki","rel":"relaciona"},{"alvo":"virtualizacao","confianca":0.5,"epistemico":"fato","origem":"wiki","rel":"relaciona"},{"alvo":"misticismo_nao_dual","confianca":0.5,"epistemico":"fato","origem":"wiki","rel":"relaciona"}],"confianca":1.0,"contraexemplos":[],"descricao":"```bash cristal task add \"reduzir RAM do reconhece\" cristal work ```","epistemico":"fato","exemplos":[],"id":"desenvolvimento_orientado_a_tarefas","memoria":"semantica","meta":{"dominio":"manual","nivel":6,"verificacao":"slug idêntico — enriquecer, não duplicar"},"origem":"CRISTAL_DEV.md","palavras_chave":[],"sinonimos":[],"tipo":"conceito","titulo":"Desenvolvimento orientado a tarefas"}
\section{Desenvolvimento orientado a tarefas}
```bash cristal task add "reduzir RAM do reconhece" cristal work ```
\prov{relações: parte\_de broca\_os; relaciona nao\_dualidade; relaciona memoria\_virtual; relaciona word; relaciona virtualizacao; relaciona misticismo\_nao\_dual}
\prov{proveniência: CRISTAL\_DEV.md \textperiodcentered{} fato \textperiodcentered{} confiança 1.0 \textperiodcentered{} domínio manual \textperiodcentered{} história 8 versões no jornal}

%CRISTAL {"arestas":[{"alvo":"parabola_destruir","confianca":1.0,"epistemico":"fato","origem":"manual","rel":"parte_de"},{"alvo":"utilitarismo","confianca":0.5,"epistemico":"fato","origem":"wiki","rel":"relaciona"},{"alvo":"teoria_do_valor","confianca":0.5,"epistemico":"fato","origem":"wiki","rel":"relaciona"},{"alvo":"gentil_16_soma_dos_termos_de_uma_pam","confianca":0.5,"epistemico":"fato","origem":"wiki","rel":"relaciona"},{"alvo":"modo_espectral_spect","confianca":0.5,"epistemico":"fato","origem":"wiki","rel":"relaciona"},{"alvo":"vetores","confianca":0.5,"epistemico":"fato","origem":"wiki","rel":"relaciona"},{"alvo":"scheduler","confianca":0.5,"epistemico":"fato","origem":"wiki","rel":"relaciona"},{"alvo":"pipeline_de_ingestao","confianca":0.5,"epistemico":"fato","origem":"wiki","rel":"relaciona"},{"alvo":"resultado_completo","confianca":0.5,"epistemico":"fato","origem":"wiki","rel":"relaciona"},{"alvo":"readme_cristal","confianca":0.5,"epistemico":"fato","origem":"wiki","rel":"relaciona"},{"alvo":"o_que_e_no_projecto","confianca":0.5,"epistemico":"fato","origem":"wiki","rel":"relaciona"},{"alvo":"experimento_3_orbitas","confianca":0.5,"epistemico":"fato","origem":"wiki","rel":"relaciona"},{"alvo":"ficheiros","confianca":0.5,"epistemico":"fato","origem":"wiki","rel":"relaciona"},{"alvo":"telemetria_shadow_producao","confianca":0.5,"epistemico":"fato","origem":"wiki","rel":"relaciona"}],"confianca":1.0,"contraexemplos":[],"descricao":"> Tentativa de quebrar **a + c² = 1** — sem pós-ajuste, sem forçar fechamento.","epistemico":"fato","exemplos":[],"id":"destruir_a_parabola","memoria":"semantica","meta":{"dominio":"manual","nivel":6,"verificacao":"slug idêntico — enriquecer, não duplicar"},"origem":"PARABOLA_DESTRUIR.md","palavras_chave":[],"sinonimos":[],"tipo":"conceito","titulo":"⚔️ Destruir a Parábola"}
\section{[espadas] Destruir a Parábola}
> Tentativa de quebrar **a + c\ensuremath{^{2}} = 1** — sem pós-ajuste, sem forçar fechamento.
\prov{relações: parte\_de parabola\_destruir; relaciona utilitarismo; relaciona teoria\_do\_valor; relaciona gentil\_16\_soma\_dos\_termos\_de\_uma\_pam; relaciona modo\_espectral\_spect; relaciona vetores; relaciona scheduler; relaciona pipeline\_de\_ingestao; relaciona resultado\_completo; relaciona readme\_cristal; relaciona o\_que\_e\_no\_projecto; relaciona experimento\_3\_orbitas; relaciona ficheiros; relaciona telemetria\_shadow\_producao}
\prov{proveniência: PARABOLA\_DESTRUIR.md \textperiodcentered{} fato \textperiodcentered{} confiança 1.0 \textperiodcentered{} domínio manual \textperiodcentered{} história 17 versões no jornal}

%CRISTAL {"arestas":[{"alvo":"gato","confianca":0.95,"epistemico":"fato","origem":"paper","rel":"parte_de"}],"confianca":1.0,"contraexemplos":[],"descricao":"| Etapa | Entraram | Passaram | Eliminados | ms | ms/cand | % acum | |-------|----------|----------|------------|-----|---------|--------| | Header | 2,000,000 | 2,000,000 | 0 | 0.0 | 0.00e+00 | 100.0000% | | Neuronio **SEM RESTRICAO** | 2,000,000 | 2,000,000 | 0 | 805.7 | 4.03e-04 | 100.0000% | | Gauss **SEM RESTRICAO** | 2,000,000 | 2,000,000 | 0 | 49.7 | 2.49e-05 | 100.0000% | | Gold **SEM RESTRICAO** | 2,000,000 | 2,000,000 | 0 | 62.2 | 3.11e-05 | 100.0000% |","epistemico":"fato","exemplos":[],"id":"detalhe_por_etapa","memoria":"semantica","meta":{"dominio":"manual","nivel":6,"verificacao":"slug idêntico — enriquecer, não duplicar"},"origem":"POW_METAL_FUNIL.md","palavras_chave":[],"sinonimos":[],"tipo":"conceito","titulo":"Detalhe por etapa"}
\section{Detalhe por etapa}
| Etapa | Entraram | Passaram | Eliminados | ms | ms/cand | \% acum | |-------|----------|----------|------------|-----|---------|--------| | Header | 2,000,000 | 2,000,000 | 0 | 0.0 | 0.00e+00 | 100.0000\% | | Neuronio **SEM RESTRICAO** | 2,000,000 | 2,000,000 | 0 | 805.7 | 4.03e-04 | 100.0000\% | | Gauss **SEM RESTRICAO** | 2,000,000 | 2,000,000 | 0 | 49.7 | 2.49e-05 | 100.0000\% | | Gold **SEM RESTRICAO** | 2,000,000 | 2,000,000 | 0 | 62.2 | 3.11e-05 | 100.0000\% |
\prov{relações: parte\_de gato}
\prov{proveniência: POW\_METAL\_FUNIL.md \textperiodcentered{} fato \textperiodcentered{} confiança 1.0 \textperiodcentered{} domínio manual \textperiodcentered{} história 3 versões no jornal}

%CRISTAL {"arestas":[{"alvo":"gato","confianca":0.95,"epistemico":"fato","origem":"paper","rel":"parte_de"},{"alvo":"utilitarismo","confianca":0.5,"epistemico":"fato","origem":"wiki","rel":"relaciona"},{"alvo":"vetores","confianca":0.5,"epistemico":"fato","origem":"wiki","rel":"relaciona"},{"alvo":"teoria_do_valor","confianca":0.5,"epistemico":"fato","origem":"wiki","rel":"relaciona"},{"alvo":"word","confianca":0.5,"epistemico":"fato","origem":"wiki","rel":"relaciona"}],"confianca":1.0,"contraexemplos":[],"descricao":"Autovalor dominante A₂=[[2,1],[1,0]]: **σ₂ = 1+√2 ≈ 2.414214**","epistemico":"fato","exemplos":[],"id":"dimensao_de_prata_comparacao_posterior","memoria":"semantica","meta":{"dominio":"manual","nivel":6,"verificacao":"slug idêntico — enriquecer, não duplicar"},"origem":"DEGENERADOS.md","palavras_chave":[],"sinonimos":[],"tipo":"conceito","titulo":"Dimensão de prata (comparação posterior)"}
\section{Dimensão de prata (comparação posterior)}
Autovalor dominante A\ensuremath{_{2}}=[[2,1],[1,0]]: **\ensuremath{\sigma}\ensuremath{_{2}} = 1+\ensuremath{\sqrt{\,}}2 \ensuremath{\approx} 2.414214**
\prov{relações: parte\_de gato; relaciona utilitarismo; relaciona vetores; relaciona teoria\_do\_valor; relaciona word}
\prov{proveniência: DEGENERADOS.md \textperiodcentered{} fato \textperiodcentered{} confiança 1.0 \textperiodcentered{} domínio manual \textperiodcentered{} história 7 versões no jornal}

%CRISTAL {"arestas":[{"alvo":"corpo_dual","confianca":1.0,"epistemico":"fato","origem":"wiki","rel":"usa"},{"alvo":"sequencia_pell","confianca":1.0,"epistemico":"fato","origem":"wiki","rel":"usa"},{"alvo":"beta_numeracao_metalica","confianca":1.0,"epistemico":"fato","origem":"wiki","rel":"especializa"}],"confianca":1.0,"contraexemplos":[],"descricao":"A prata é n=2 na reta dos metais: σ₂=1+√2, entropia log σ₂; os números de prata são a sequência de Pell Pₖ=2Pₖ−₁+Pₖ−₂, com primos de Pell em índices primos (2,5,29,5741,33461…).","epistemico":"fato","exemplos":[],"id":"dimensao_de_prata_pell","memoria":"semantica","meta":{"autor":"Aarão/broca-so","dominio":"broca_repos","fonte":""},"origem":"wiki","palavras_chave":[],"sinonimos":[],"tipo":"conceito","titulo":"Dimensão de Prata e os Números de Pell"}
\section{Dimensão de Prata e os Números de Pell}
A prata é n=2 na reta dos metais: \ensuremath{\sigma}\ensuremath{_{2}}=1+\ensuremath{\sqrt{\,}}2, entropia log \ensuremath{\sigma}\ensuremath{_{2}}; os números de prata são a sequência de Pell P\ensuremath{_{k}}=2P\ensuremath{_{k}}\ensuremath{-}\ensuremath{_{1}}+P\ensuremath{_{k}}\ensuremath{-}\ensuremath{_{2}}, com primos de Pell em índices primos (2,5,29,5741,33461…).
\prov{relações: usa corpo\_dual; usa sequencia\_pell; especializa beta\_numeracao\_metalica}
\prov{proveniência: wiki \textperiodcentered{} fato \textperiodcentered{} confiança 1.0 \textperiodcentered{} domínio broca\_repos \textperiodcentered{} história 5 versões no jornal}

%CRISTAL {"arestas":[{"alvo":"cifra","confianca":1.0,"epistemico":"fato","origem":"wiki","rel":"relaciona"},{"alvo":"maquina_algebrica","confianca":1.0,"epistemico":"fato","origem":"wiki","rel":"usa"}],"confianca":1.0,"contraexemplos":[],"descricao":"Aplicando o distinguisher de Gohr com uma HyperNet (MLP quaterniônico) vs RealNet pareada em parâmetros, a vantagem de distinção cai monotonicamente e atinge o acaso em r=8 rodadas; a fronteira prática do distinguisher de 1 bit é r=6. E a degenerescência da bacia (∂L/∂s≈0) vira detector: uma bacia não-degenerada sinalizaria fresta atacável.","epistemico":"fato","exemplos":[],"id":"distinguisher_sha_hipercomplexo","memoria":"semantica","meta":{"autor":"Aarão/broca-so","dominio":"broca_repos","fonte":""},"origem":"wiki","palavras_chave":[],"sinonimos":[],"tipo":"conceito","titulo":"Distinguisher Diferencial Hipercomplexo do SHA-256"}
\section{Distinguisher Diferencial Hipercomplexo do SHA-256}
Aplicando o distinguisher de Gohr com uma HyperNet (MLP quaterniônico) vs RealNet pareada em parâmetros, a vantagem de distinção cai monotonicamente e atinge o acaso em r=8 rodadas; a fronteira prática do distinguisher de 1 bit é r=6. E a degenerescência da bacia (\ensuremath{\partial}L/\ensuremath{\partial}s\ensuremath{\approx}0) vira detector: uma bacia não-degenerada sinalizaria fresta atacável.
\prov{relações: relaciona cifra; usa maquina\_algebrica}
\prov{proveniência: wiki \textperiodcentered{} fato \textperiodcentered{} confiança 1.0 \textperiodcentered{} domínio broca\_repos \textperiodcentered{} história 3 versões no jornal}

%CRISTAL {"arestas":[{"alvo":"koch_atlas_experimento_de_parametrizacao","confianca":1.0,"epistemico":"fato","origem":"KOCH_ATLAS.md","rel":"parte_de"}],"confianca":1.0,"contraexemplos":[],"descricao":"| Métrica | Valor | |---------|-------| | Ramo 0 | 65,536 (100.00%) | | Ramo 1 | 0 (0.00%) | | Ramo 2 | 0 (0.00%) | | Entropia ramos | 0.0000 bits | | Entropia atrator | 5.0313 bits | | Entropia endereço | 3.0566 bits | | Atrator repetidos | 65,459 | | Endereço repetidos | 65,517 | | Órbita repetidas | 65,517 | | Profundidade (nível) | 20 |","epistemico":"fato","exemplos":[],"id":"distribuicoes","memoria":"semantica","meta":{"dominio":"manual","nivel":6,"verificacao":"slug idêntico — enriquecer, não duplicar"},"origem":"KOCH_ATLAS.md","palavras_chave":[],"sinonimos":[],"tipo":"conceito","titulo":"Distribuições"}
\section{Distribuições}
| Métrica | Valor | |---------|-------| | Ramo 0 | 65,536 (100.00\%) | | Ramo 1 | 0 (0.00\%) | | Ramo 2 | 0 (0.00\%) | | Entropia ramos | 0.0000 bits | | Entropia atrator | 5.0313 bits | | Entropia endereço | 3.0566 bits | | Atrator repetidos | 65,459 | | Endereço repetidos | 65,517 | | Órbita repetidas | 65,517 | | Profundidade (nível) | 20 |
\prov{relações: parte\_de koch\_atlas\_experimento\_de\_parametrizacao}
\prov{proveniência: KOCH\_ATLAS.md \textperiodcentered{} fato \textperiodcentered{} confiança 1.0 \textperiodcentered{} domínio manual \textperiodcentered{} história 6 versões no jornal}

%CRISTAL {"arestas":[{"alvo":"cosmologia_quantica","confianca":1.0,"epistemico":"fato","origem":"manual","rel":"referencia"},{"alvo":"transcricao","confianca":0.5,"epistemico":"fato","origem":"wiki","rel":"relaciona"},{"alvo":"vida-morte","confianca":0.5,"epistemico":"fato","origem":"wiki","rel":"relaciona"},{"alvo":"resposta_a_pergunta_dificil","confianca":0.5,"epistemico":"fato","origem":"wiki","rel":"relaciona"},{"alvo":"teste_ordem","confianca":0.5,"epistemico":"fato","origem":"wiki","rel":"relaciona"},{"alvo":"parabola_destruir","confianca":0.5,"epistemico":"fato","origem":"wiki","rel":"relaciona"},{"alvo":"pow_metal_funil","confianca":0.5,"epistemico":"fato","origem":"wiki","rel":"relaciona"},{"alvo":"log_legado_conversadadostelemetriajsonl","confianca":0.5,"epistemico":"fato","origem":"wiki","rel":"relaciona"},{"alvo":"parabola_pa","confianca":0.5,"epistemico":"fato","origem":"wiki","rel":"relaciona"},{"alvo":"koch_cantor_fib_prova","confianca":0.5,"epistemico":"fato","origem":"wiki","rel":"relaciona"},{"alvo":"por_sessao_conversadadostelemetriasessaojsonl","confianca":0.5,"epistemico":"fato","origem":"wiki","rel":"relaciona"},{"alvo":"motivos","confianca":0.5,"epistemico":"fato","origem":"wiki","rel":"relaciona"},{"alvo":"metais","confianca":0.5,"epistemico":"fato","origem":"wiki","rel":"relaciona"},{"alvo":"pow_metal_pre_fecho","confianca":0.5,"epistemico":"fato","origem":"wiki","rel":"relaciona"},{"alvo":"geometria_hub_broca","confianca":0.5,"epistemico":"fato","origem":"wiki","rel":"relaciona"},{"alvo":"topologia","confianca":0.5,"epistemico":"fato","origem":"wiki","rel":"relaciona"},{"alvo":"virtualizacao","confianca":0.5,"epistemico":"fato","origem":"wiki","rel":"relaciona"},{"alvo":"exatidao","confianca":0.5,"epistemico":"fato","origem":"wiki","rel":"relaciona"},{"alvo":"pow_metal_setor","confianca":0.5,"epistemico":"fato","origem":"wiki","rel":"relaciona"}],"confianca":1.0,"contraexemplos":[],"descricao":"Manual: ula/DS_UNIVERSO.md","epistemico":"fato","exemplos":[],"id":"ds_universo","memoria":"semantica","meta":{"arquivo":"ula/DS_UNIVERSO.md","dominio":"manual","nivel":6},"origem":"manual","palavras_chave":[],"sinonimos":[],"tipo":"conceito","titulo":"DS_UNIVERSO"}
\section{DS\_UNIVERSO}
Manual: ula/DS\_UNIVERSO.md
\prov{relações: referencia cosmologia\_quantica; relaciona transcricao; relaciona vida-morte; relaciona resposta\_a\_pergunta\_dificil; relaciona teste\_ordem; relaciona parabola\_destruir; relaciona pow\_metal\_funil; relaciona log\_legado\_conversadadostelemetriajsonl; relaciona parabola\_pa; relaciona koch\_cantor\_fib\_prova; relaciona por\_sessao\_conversadadostelemetriasessaojsonl; relaciona motivos; relaciona metais; relaciona pow\_metal\_pre\_fecho; relaciona geometria\_hub\_broca; relaciona topologia; relaciona virtualizacao; relaciona exatidao; relaciona pow\_metal\_setor}
\prov{proveniência: manual \textperiodcentered{} fato \textperiodcentered{} confiança 1.0 \textperiodcentered{} domínio manual \textperiodcentered{} história 22 versões no jornal}

%CRISTAL {"arestas":[{"alvo":"colapso","confianca":-0.086,"epistemico":"fato","origem":"manual","rel":"referencia"},{"alvo":"broca-so_camada_de_trabalho_isomorfica","confianca":1.0,"epistemico":"fato","origem":"BROCA_SO.md","rel":"parte_de"}],"confianca":1.0,"contraexemplos":[],"descricao":"```text ┌─────────────────────────────────────────────────────────────┐ │  CAMADA DE TRABALHO — broca_so.h                            │ │  Memória (slots) + Ops algébricas (gato, toro, colapso…)   │ │  Isomórfica: cpu | gpu/ptx | android — mesma prova          │ └────────────────────────────┬────────────────────────────────┘ │ broca_backend.h ┌───────────────────┼───────────────────┐ ▼                   ▼                   ▼","epistemico":"fato","exemplos":[],"id":"duas_camadas","memoria":"semantica","meta":{"dominio":"manual","nivel":6,"verificacao":"slug idêntico — enriquecer, não duplicar"},"origem":"BROCA_SO.md","palavras_chave":[],"sinonimos":[],"tipo":"conceito","titulo":"Duas camadas"}
\section{Duas camadas}
```text ------------------------------------------------------------- |  CAMADA DE TRABALHO — broca\_so.h                            | |  Memória (slots) + Ops algébricas (gato, toro, colapso…)   | |  Isomórfica: cpu | gpu/ptx | android — mesma prova          | ------------------------------------------------------------- | broca\_backend.h -------------------+------------------- $\triangledown$                   $\triangledown$                   $\triangledown$
\prov{relações: referencia colapso; parte\_de broca-so\_camada\_de\_trabalho\_isomorfica}
\prov{proveniência: BROCA\_SO.md \textperiodcentered{} fato \textperiodcentered{} confiança 1.0 \textperiodcentered{} domínio manual \textperiodcentered{} história 9 versões no jornal}

%CRISTAL {"arestas":[{"alvo":"turmalina_paraiba_ancora","confianca":1.0,"epistemico":"fato","origem":"wiki","rel":"parte_de"},{"alvo":"cristalchain","confianca":1.0,"epistemico":"fato","origem":"wiki","rel":"usa"},{"alvo":"cifra_de_cristal","confianca":1.0,"epistemico":"fato","origem":"wiki","rel":"usa"},{"alvo":"tesouraria","confianca":1.0,"epistemico":"fato","origem":"wiki","rel":"relaciona"},{"alvo":"area_gemologia","confianca":1.0,"epistemico":"fato","origem":"wiki","rel":"relaciona"},{"alvo":"joalheria_hub","confianca":1.0,"epistemico":"fato","origem":"wiki","rel":"relaciona"}],"confianca":1.0,"contraexemplos":[],"descricao":"A emissão que fixa o PISO da CristalChain: 1 QUILATE de Turmalina Paraíba pura = R$ 1.000,00. A ancoragem numérica vem de 3 CPFs brasileiros (a família Cunha Gomes: Patrícia/esposa, Artur/filho, Carlos/marido), que semeiam uma CHAVE 3-QUBIT (8 amplitudes, normalizada) — os 3 eixos da pedra = 3 titulares = 3 qubits. O selo é o hash da chave (df5ee62b…); a ancoragem dos CPFs é o hash 7347f5bb…. Os CPFs em bruto ficam em registro privado (gitignored, dado pessoal); só a âncora derivada e os titulares mascarados são versionados. Fixado em cristalchain/piso_turmalina.json. A âncora deixa de ser abstrata: uma pedra, uma família, um selo.","epistemico":"fato","exemplos":[],"id":"emissao_turmalina_piso","memoria":"semantica","meta":{"dominio":"referencia","fonte":"turmalina paraíba"},"origem":"wiki","palavras_chave":[],"sinonimos":[],"tipo":"conceito","titulo":"A Emissão da Turmalina: o Piso 1 ct = R$ 1.000,00"}
\section{A Emissão da Turmalina: o Piso 1 ct = R\$ 1.000,00}
A emissão que fixa o PISO da CristalChain: 1 QUILATE de Turmalina Paraíba pura = R\$ 1.000,00. A ancoragem numérica vem de 3 CPFs brasileiros (a família Cunha Gomes: Patrícia/esposa, Artur/filho, Carlos/marido), que semeiam uma CHAVE 3-QUBIT (8 amplitudes, normalizada) — os 3 eixos da pedra = 3 titulares = 3 qubits. O selo é o hash da chave (df5ee62b…); a ancoragem dos CPFs é o hash 7347f5bb…. Os CPFs em bruto ficam em registro privado (gitignored, dado pessoal); só a âncora derivada e os titulares mascarados são versionados. Fixado em cristalchain/piso\_turmalina.json. A âncora deixa de ser abstrata: uma pedra, uma família, um selo.
\prov{relações: parte\_de turmalina\_paraiba\_ancora; usa cristalchain; usa cifra\_de\_cristal; relaciona tesouraria; relaciona area\_gemologia; relaciona joalheria\_hub}
\prov{proveniência: wiki \textperiodcentered{} turmalina paraíba \textperiodcentered{} fato \textperiodcentered{} confiança 1.0 \textperiodcentered{} domínio referencia \textperiodcentered{} história 15 versões no jornal}

%CRISTAL {"arestas":[{"alvo":"geometria","confianca":0.031,"epistemico":"fato","origem":"manual","rel":"referencia"},{"alvo":"pow_coordenadas_espectrais_broca","confianca":1.0,"epistemico":"fato","origem":"POW_GEOMETRIA.md","rel":"parte_de"}],"confianca":1.0,"contraexemplos":[],"descricao":"| Caminho | Conteúdo | |---------|----------| | `atlas/pow_geometria/campos/` | amostras campo espectral | | `atlas/pow_geometria/mapas/` | heatmap calor×fase² | | `atlas/pow_geometria/orbita/` | (reservado) | | `atlas/pow_geometria/metal/` | ranking A₁…A₈ | | `atlas/pow_geometria/coordenadas/` | continuidade, compressão |","epistemico":"fato","exemplos":[],"id":"entregaveis","memoria":"semantica","meta":{"dominio":"manual","nivel":6,"verificacao":"slug idêntico — enriquecer, não duplicar"},"origem":"POW_GEOMETRIA.md","palavras_chave":[],"sinonimos":[],"tipo":"conceito","titulo":"Entregáveis"}
\section{Entregáveis}
| Caminho | Conteúdo | |---------|----------| | `atlas/pow\_geometria/campos/` | amostras campo espectral | | `atlas/pow\_geometria/mapas/` | heatmap calor$\times$fase\ensuremath{^{2}} | | `atlas/pow\_geometria/orbita/` | (reservado) | | `atlas/pow\_geometria/metal/` | ranking A\ensuremath{_{1}}…A\ensuremath{_{8}} | | `atlas/pow\_geometria/coordenadas/` | continuidade, compressão |
\prov{relações: referencia geometria; parte\_de pow\_coordenadas\_espectrais\_broca}
\prov{proveniência: POW\_GEOMETRIA.md \textperiodcentered{} fato \textperiodcentered{} confiança 1.0 \textperiodcentered{} domínio manual \textperiodcentered{} história 9 versões no jornal}

%CRISTAL {"arestas":[{"alvo":"o_conselho","confianca":1.0,"epistemico":"fato","origem":"wiki","rel":"parte_de"},{"alvo":"miopia_coletiva","confianca":1.0,"epistemico":"fato","origem":"wiki","rel":"relaciona"}],"confianca":1.0,"contraexemplos":[],"descricao":"Cada IA tem função fixa: GPT formaliza/corrige, Grok é o crivo (contraexemplos e refutações), Gemini enquadra/pole, Claude constrói e sintetiza, Aarão traz hipóteses/visões, e 'conselho' é o que convergiu. Divisão de trabalho cognitivo — nenhuma voz vê tudo sozinha.","epistemico":"inferencia","exemplos":[],"id":"especializacao_das_vozes","memoria":"semantica","meta":{"autor":"Aarão/broca-so","dominio":"broca_repos","fonte":""},"origem":"wiki","palavras_chave":[],"sinonimos":[],"tipo":"conceito","titulo":"Papel por Tecido (especialização das vozes)"}
\section{Papel por Tecido (especialização das vozes)}
Cada IA tem função fixa: GPT formaliza/corrige, Grok é o crivo (contraexemplos e refutações), Gemini enquadra/pole, Claude constrói e sintetiza, Aarão traz hipóteses/visões, e 'conselho' é o que convergiu. Divisão de trabalho cognitivo — nenhuma voz vê tudo sozinha.
\prov{relações: parte\_de o\_conselho; relaciona miopia\_coletiva}
\prov{proveniência: wiki \textperiodcentered{} inferencia \textperiodcentered{} confiança 1.0 \textperiodcentered{} domínio broca\_repos \textperiodcentered{} história 3 versões no jornal}

%CRISTAL {"arestas":[{"alvo":"koch_atlas_experimento_de_parametrizacao","confianca":1.0,"epistemico":"fato","origem":"KOCH_ATLAS.md","rel":"parte_de"},{"alvo":"virtualizacao","confianca":0.5,"epistemico":"fato","origem":"wiki","rel":"relaciona"},{"alvo":"word","confianca":0.5,"epistemico":"fato","origem":"wiki","rel":"relaciona"},{"alvo":"vontade_de_poder","confianca":0.5,"epistemico":"fato","origem":"wiki","rel":"relaciona"}],"confianca":1.0,"contraexemplos":[],"descricao":"- Atratores únicos (aprox.): 77 - Endereços únicos (aprox.): 19 - Órbitas únicas (aprox.): 19","epistemico":"fato","exemplos":[],"id":"estabilidade","memoria":"semantica","meta":{"dominio":"manual","nivel":6,"verificacao":"slug idêntico — enriquecer, não duplicar"},"origem":"KOCH_ATLAS.md","palavras_chave":[],"sinonimos":[],"tipo":"conceito","titulo":"Estabilidade"}
\section{Estabilidade}
- Atratores únicos (aprox.): 77 - Endereços únicos (aprox.): 19 - Órbitas únicas (aprox.): 19
\prov{relações: parte\_de koch\_atlas\_experimento\_de\_parametrizacao; relaciona virtualizacao; relaciona word; relaciona vontade\_de\_poder}
\prov{proveniência: KOCH\_ATLAS.md \textperiodcentered{} fato \textperiodcentered{} confiança 1.0 \textperiodcentered{} domínio manual \textperiodcentered{} história 9 versões no jornal}

%CRISTAL {"arestas":[{"alvo":"nao_associatividade_amputacao","confianca":1.0,"epistemico":"fato","origem":"wiki","rel":"contradiz"},{"alvo":"matematica_estelar","confianca":1.0,"epistemico":"fato","origem":"wiki","rel":"parte_de"}],"confianca":1.0,"contraexemplos":[],"descricao":"A estabilização giroscópica pós-cirurgia, conjecturada para reabsorver o gap, é REFUTADA no intervalo físico: o corolário prova que o s crítico s_c ≥ 2√12 ≈ 6,93 > 1, logo não ocorre em s∈[0,1]. O vácuo permanece estéril — o Lopes registra o que falha.","epistemico":"fato","exemplos":[],"id":"estabilizacao_giroscopica_impossivel","memoria":"semantica","meta":{"autor":"Lopes/Aarão","dominio":"hiper","fonte":"","refutado":true},"origem":"wiki","palavras_chave":[],"sinonimos":[],"tipo":"conceito","titulo":"Estabilização Giroscópica Impossível (refutado)"}
\section{Estabilização Giroscópica Impossível (refutado)}
A estabilização giroscópica pós-cirurgia, conjecturada para reabsorver o gap, é REFUTADA no intervalo físico: o corolário prova que o s crítico s\_c \ensuremath{\ge} 2\ensuremath{\sqrt{\,}}12 \ensuremath{\approx} 6,93 > 1, logo não ocorre em s\ensuremath{\in}[0,1]. O vácuo permanece estéril — o Lopes registra o que falha.
\prov{relações: contradiz nao\_associatividade\_amputacao; parte\_de matematica\_estelar}
\prov{proveniência: wiki \textperiodcentered{} fato \textperiodcentered{} confiança 1.0 \textperiodcentered{} domínio hiper \textperiodcentered{} história 4 versões no jornal}

%CRISTAL {"arestas":[{"alvo":"broca_os","confianca":1.0,"epistemico":"fato","origem":"manual","rel":"parte_de"}],"confianca":1.0,"contraexemplos":[],"descricao":"- classes únicas: **985,597** - simetrias distintas: **978,541** - recorrências detectadas: **0** - recorrências Pell exatas (a1=2,a0=1): **0** - recorrências exóticas: **0** - valores |x| na tabela Pell: **4** - valores |x| na Fibonacci: **5** - componentes conexas (top grafo): **1**","epistemico":"fato","exemplos":[],"id":"estatisticas","memoria":"semantica","meta":{"dominio":"manual","nivel":6,"verificacao":"slug idêntico — enriquecer, não duplicar"},"origem":"DEGENERADOS.md","palavras_chave":[],"sinonimos":[],"tipo":"conceito","titulo":"Estatísticas"}
\section{Estatísticas}
- classes únicas: **985,597** - simetrias distintas: **978,541** - recorrências detectadas: **0** - recorrências Pell exatas (a1=2,a0=1): **0** - recorrências exóticas: **0** - valores |x| na tabela Pell: **4** - valores |x| na Fibonacci: **5** - componentes conexas (top grafo): **1**
\prov{relações: parte\_de broca\_os}
\prov{proveniência: DEGENERADOS.md \textperiodcentered{} fato \textperiodcentered{} confiança 1.0 \textperiodcentered{} domínio manual \textperiodcentered{} história 8 versões no jornal}

%CRISTAL {"arestas":[{"alvo":"algebra_estelar","confianca":1.0,"epistemico":"fato","origem":"wiki","rel":"explica"},{"alvo":"cayley_dickson","confianca":1.0,"epistemico":"fato","origem":"wiki","rel":"usa"}],"confianca":1.0,"contraexemplos":[],"descricao":"Teorema (verificado a 10⁻¹⁵): z⋆ₛw é a projeção do produto geométrico de Cl(0,m) sobre os graus 0 e 1; a obstrução a+c²=1 é exatamente a norma do bivetor (grau 2) descartado (identidade de Lagrange). No Clifford pleno a norma é multiplicativa em toda dimensão.","epistemico":"fato","exemplos":[],"id":"estrela_clifford_projetado","memoria":"semantica","meta":{"autor":"Lopes/Aarão","dominio":"hiper","fonte":""},"origem":"wiki","palavras_chave":[],"sinonimos":[],"tipo":"conceito","titulo":"A Estrela é o Clifford Projetado"}
\section{A Estrela é o Clifford Projetado}
Teorema (verificado a 10\ensuremath{^{-15}}): z\ensuremath{\star}\ensuremath{_{s}}w é a projeção do produto geométrico de Cl(0,m) sobre os graus 0 e 1; a obstrução a+c\ensuremath{^{2}}=1 é exatamente a norma do bivetor (grau 2) descartado (identidade de Lagrange). No Clifford pleno a norma é multiplicativa em toda dimensão.
\prov{relações: explica algebra\_estelar; usa cayley\_dickson}
\prov{proveniência: wiki \textperiodcentered{} fato \textperiodcentered{} confiança 1.0 \textperiodcentered{} domínio hiper \textperiodcentered{} história 4 versões no jornal}

%CRISTAL {"arestas":[{"alvo":"rede_neural_algebra_aprendivel","confianca":1.0,"epistemico":"fato","origem":"wiki","rel":"usa"},{"alvo":"prova_de_trabalho","confianca":1.0,"epistemico":"fato","origem":"wiki","rel":"relaciona"}],"confianca":1.0,"contraexemplos":[],"descricao":"Dois neurônios inicializados no mesmo ponto algébrico (r,t) desenvolvem repulsão sistemática (~35000%) sob gradiente, sem termo de repulsão explícito — induzindo diversidade algébrica emergente. O regime 'fermiônico' (parâmetros locais) supera o 'bosônico' (globais).","epistemico":"fato","exemplos":[],"id":"exclusao_pauli_algebrica","memoria":"semantica","meta":{"autor":"Lopes/Aarão","dominio":"hiper","fonte":""},"origem":"wiki","palavras_chave":[],"sinonimos":[],"tipo":"conceito","titulo":"Exclusão de Pauli Algébrica"}
\section{Exclusão de Pauli Algébrica}
Dois neurônios inicializados no mesmo ponto algébrico (r,t) desenvolvem repulsão sistemática (\~{}35000\%) sob gradiente, sem termo de repulsão explícito — induzindo diversidade algébrica emergente. O regime 'fermiônico' (parâmetros locais) supera o 'bosônico' (globais).
\prov{relações: usa rede\_neural\_algebra\_aprendivel; relaciona prova\_de\_trabalho}
\prov{proveniência: wiki \textperiodcentered{} fato \textperiodcentered{} confiança 1.0 \textperiodcentered{} domínio hiper \textperiodcentered{} história 4 versões no jornal}

%CRISTAL {"arestas":[{"alvo":"contraexemplo","confianca":1.0,"epistemico":"fato","origem":"curriculo","rel":"contradiz"},{"alvo":"definicao","confianca":1.0,"epistemico":"fato","origem":"curriculo","rel":"confirma"},{"alvo":"word","confianca":0.5,"epistemico":"fato","origem":"wiki","rel":"relaciona"},{"alvo":"while","confianca":0.5,"epistemico":"fato","origem":"wiki","rel":"relaciona"},{"alvo":"virtualizacao","confianca":0.5,"epistemico":"fato","origem":"wiki","rel":"relaciona"},{"alvo":"vontade_de_poder","confianca":0.5,"epistemico":"fato","origem":"wiki","rel":"relaciona"}],"confianca":1.0,"contraexemplos":[],"descricao":"```bash export TIFFANY_HOME=~/.tiffany cristal init sh content/core-p0/build.sh    # ou: cristal base update","epistemico":"fato","exemplos":[],"id":"exemplo","memoria":"semantica","meta":{"dominio":"manual","duplicata_sugerida":[],"nivel":6,"verificacao":"parte de conceito mais amplo 'exemplo'"},"origem":"README_CRISTAL.md","palavras_chave":[],"sinonimos":["exemplos"],"tipo":"conceito","titulo":"exemplo"}
\section{exemplo}
```bash export TIFFANY\_HOME=\~{}/.tiffany cristal init sh content/core-p0/build.sh    \# ou: cristal base update
\prov{sinónimos: exemplos}
\prov{relações: contradiz contraexemplo; confirma definicao; relaciona word; relaciona while; relaciona virtualizacao; relaciona vontade\_de\_poder}
\prov{proveniência: README\_CRISTAL.md \textperiodcentered{} fato \textperiodcentered{} confiança 1.0 \textperiodcentered{} domínio manual \textperiodcentered{} história 68 versões no jornal}

%CRISTAL {"arestas":[{"alvo":"pow_espectral_investigacao_broca","confianca":1.0,"epistemico":"fato","origem":"POW_ESPECTRAL.md","rel":"parte_de"}],"confianca":1.0,"contraexemplos":[],"descricao":"- válidos (não degenerados): **2048** / 2048 - degenerados: **0** (0.00%) - hits calor ≤ 0.15: **2048** (taxa 1.0000) - calor médio: **0.0000** - entropia calor: **1.4612** / max 2.9957","epistemico":"fato","exemplos":[],"id":"experimento_1_distribuicao_espectral_dos_candidatos","memoria":"semantica","meta":{"dominio":"manual","nivel":6,"verificacao":"slug idêntico — enriquecer, não duplicar"},"origem":"POW_ESPECTRAL.md","palavras_chave":[],"sinonimos":[],"tipo":"conceito","titulo":"Experimento 1 — Distribuição espectral dos candidatos"}
\section{Experimento 1 — Distribuição espectral dos candidatos}
- válidos (não degenerados): **2048** / 2048 - degenerados: **0** (0.00\%) - hits calor \ensuremath{\le} 0.15: **2048** (taxa 1.0000) - calor médio: **0.0000** - entropia calor: **1.4612** / max 2.9957
\prov{relações: parte\_de pow\_espectral\_investigacao\_broca}
\prov{proveniência: POW\_ESPECTRAL.md \textperiodcentered{} fato \textperiodcentered{} confiança 1.0 \textperiodcentered{} domínio manual \textperiodcentered{} história 6 versões no jornal}

%CRISTAL {"arestas":[{"alvo":"metais","confianca":0.95,"epistemico":"fato","origem":"paper","rel":"parte_de"},{"alvo":"virtualizacao","confianca":0.5,"epistemico":"fato","origem":"wiki","rel":"relaciona"},{"alvo":"word","confianca":0.5,"epistemico":"fato","origem":"wiki","rel":"relaciona"},{"alvo":"universidade_curriculo_em_camadas","confianca":0.5,"epistemico":"fato","origem":"wiki","rel":"relaciona"},{"alvo":"gramatica","confianca":0.5,"epistemico":"fato","origem":"wiki","rel":"relaciona"},{"alvo":"proposition_gerador_de_escala_ipropderiv_com_tddt","confianca":0.5,"epistemico":"fato","origem":"wiki","rel":"relaciona"},{"alvo":"soma_velocidades","confianca":0.5,"epistemico":"fato","origem":"wiki","rel":"relaciona"}],"confianca":1.0,"contraexemplos":[],"descricao":"Órbita Aₙ·(1,0) para n = 1…8. A parábola **não depende** de n — depende só dos pares (Word,Word).","epistemico":"fato","exemplos":[],"id":"experimento_1_trocar_a_prata","memoria":"semantica","meta":{"dominio":"manual","nivel":6,"verificacao":"slug idêntico — enriquecer, não duplicar"},"origem":"METAIS.md","palavras_chave":[],"sinonimos":[],"tipo":"conceito","titulo":"Experimento 1 — Trocar a prata"}
\section{Experimento 1 — Trocar a prata}
Órbita A\ensuremath{_{n}}\textperiodcentered (1,0) para n = 1…8. A parábola **não depende** de n — depende só dos pares (Word,Word).
\prov{relações: parte\_de metais; relaciona virtualizacao; relaciona word; relaciona universidade\_curriculo\_em\_camadas; relaciona gramatica; relaciona proposition\_gerador\_de\_escala\_ipropderiv\_com\_tddt; relaciona soma\_velocidades}
\prov{proveniência: METAIS.md \textperiodcentered{} fato \textperiodcentered{} confiança 1.0 \textperiodcentered{} domínio manual \textperiodcentered{} história 10 versões no jornal}

%CRISTAL {"arestas":[{"alvo":"metais","confianca":0.95,"epistemico":"fato","origem":"paper","rel":"parte_de"},{"alvo":"pitagorismo","confianca":0.5,"epistemico":"fato","origem":"wiki","rel":"relaciona"},{"alvo":"teoria_do_valor","confianca":0.5,"epistemico":"fato","origem":"wiki","rel":"relaciona"},{"alvo":"utilitarismo","confianca":0.5,"epistemico":"fato","origem":"wiki","rel":"relaciona"}],"confianca":1.0,"contraexemplos":[],"descricao":"Fit de Aₙ sobre sequências `total` dos programas ISA (sem impor matriz):","epistemico":"fato","exemplos":[],"id":"experimento_2_a_maquina_escolhe","memoria":"semantica","meta":{"dominio":"manual","nivel":6,"verificacao":"slug idêntico — enriquecer, não duplicar"},"origem":"METAIS.md","palavras_chave":[],"sinonimos":[],"tipo":"conceito","titulo":"Experimento 2 — A máquina escolhe"}
\section{Experimento 2 — A máquina escolhe}
Fit de A\ensuremath{_{n}} sobre sequências `total` dos programas ISA (sem impor matriz):
\prov{relações: parte\_de metais; relaciona pitagorismo; relaciona teoria\_do\_valor; relaciona utilitarismo}
\prov{proveniência: METAIS.md \textperiodcentered{} fato \textperiodcentered{} confiança 1.0 \textperiodcentered{} domínio manual \textperiodcentered{} história 6 versões no jornal}

%CRISTAL {"arestas":[{"alvo":"colapso","confianca":0.95,"epistemico":"fato","origem":"paper","rel":"parte_de"},{"alvo":"fase_2_tecido_de_conhecimento","confianca":0.5,"epistemico":"fato","origem":"wiki","rel":"relaciona"},{"alvo":"utilitarismo","confianca":0.5,"epistemico":"fato","origem":"wiki","rel":"relaciona"},{"alvo":"ipc","confianca":0.5,"epistemico":"fato","origem":"wiki","rel":"relaciona"},{"alvo":"word","confianca":0.5,"epistemico":"fato","origem":"wiki","rel":"relaciona"}],"confianca":1.0,"contraexemplos":[],"descricao":"Cada passo: estado ideal vs perturbado → a, c², **Δ = 1−(a+c²)**.","epistemico":"fato","exemplos":[],"id":"experimento_2_mapa_do_erro","memoria":"semantica","meta":{"dominio":"manual","nivel":6,"verificacao":"slug idêntico — enriquecer, não duplicar"},"origem":"RUÍDO.md","palavras_chave":[],"sinonimos":[],"tipo":"conceito","titulo":"Experimento 2 — Mapa do erro"}
\section{Experimento 2 — Mapa do erro}
Cada passo: estado ideal vs perturbado \ensuremath{\to} a, c\ensuremath{^{2}}, **\ensuremath{\Delta} = 1\ensuremath{-}(a+c\ensuremath{^{2}})**.
\prov{relações: parte\_de colapso; relaciona fase\_2\_tecido\_de\_conhecimento; relaciona utilitarismo; relaciona ipc; relaciona word}
\prov{proveniência: RUÍDO.md \textperiodcentered{} fato \textperiodcentered{} confiança 1.0 \textperiodcentered{} domínio manual \textperiodcentered{} história 7 versões no jornal}

%CRISTAL {"arestas":[{"alvo":"pow_espectral_investigacao_broca","confianca":1.0,"epistemico":"fato","origem":"POW_ESPECTRAL.md","rel":"parte_de"},{"alvo":"word","confianca":0.5,"epistemico":"fato","origem":"wiki","rel":"relaciona"},{"alvo":"virtualizacao","confianca":0.5,"epistemico":"fato","origem":"wiki","rel":"relaciona"},{"alvo":"syscall","confianca":0.5,"epistemico":"fato","origem":"wiki","rel":"relaciona"},{"alvo":"sunyata_vazio","confianca":0.5,"epistemico":"fato","origem":"wiki","rel":"relaciona"}],"confianca":1.0,"contraexemplos":[],"descricao":"- melhor fit `A_n`: **n=1** (score 1.0000) - falhas Pell (n=2) na sequência de totais: **2046** (1.0000)","epistemico":"fato","exemplos":[],"id":"experimento_2_recorrencias_e_metais","memoria":"semantica","meta":{"dominio":"manual","nivel":6,"verificacao":"slug idêntico — enriquecer, não duplicar"},"origem":"POW_ESPECTRAL.md","palavras_chave":[],"sinonimos":[],"tipo":"conceito","titulo":"Experimento 2 — Recorrências e metais"}
\section{Experimento 2 — Recorrências e metais}
- melhor fit `A\_n`: **n=1** (score 1.0000) - falhas Pell (n=2) na sequência de totais: **2046** (1.0000)
\prov{relações: parte\_de pow\_espectral\_investigacao\_broca; relaciona word; relaciona virtualizacao; relaciona syscall; relaciona sunyata\_vazio}
\prov{proveniência: POW\_ESPECTRAL.md \textperiodcentered{} fato \textperiodcentered{} confiança 1.0 \textperiodcentered{} domínio manual \textperiodcentered{} história 10 versões no jornal}

%CRISTAL {"arestas":[{"alvo":"metais","confianca":0.95,"epistemico":"fato","origem":"paper","rel":"parte_de"},{"alvo":"transcricao","confianca":0.5,"epistemico":"fato","origem":"wiki","rel":"relaciona"},{"alvo":"utilitarismo","confianca":0.5,"epistemico":"fato","origem":"wiki","rel":"relaciona"},{"alvo":"ruido","confianca":0.5,"epistemico":"fato","origem":"wiki","rel":"relaciona"},{"alvo":"virtualizacao","confianca":0.5,"epistemico":"fato","origem":"wiki","rel":"relaciona"},{"alvo":"pedagogia","confianca":0.5,"epistemico":"fato","origem":"wiki","rel":"relaciona"},{"alvo":"vida-morte","confianca":0.5,"epistemico":"fato","origem":"wiki","rel":"relaciona"},{"alvo":"pow_metal_ressonancia_val","confianca":0.5,"epistemico":"fato","origem":"wiki","rel":"relaciona"},{"alvo":"mecanica","confianca":0.5,"epistemico":"fato","origem":"wiki","rel":"relaciona"},{"alvo":"scheduler","confianca":0.5,"epistemico":"fato","origem":"wiki","rel":"relaciona"},{"alvo":"resultado_anterior_identidade_nao_cobertura","confianca":0.5,"epistemico":"fato","origem":"wiki","rel":"relaciona"},{"alvo":"modo_espectral_spect","confianca":0.5,"epistemico":"fato","origem":"wiki","rel":"relaciona"},{"alvo":"paper_2_reta_ouro","confianca":0.5,"epistemico":"fato","origem":"wiki","rel":"relaciona"},{"alvo":"resultado_completo","confianca":0.5,"epistemico":"fato","origem":"wiki","rel":"relaciona"},{"alvo":"pitagorismo","confianca":0.5,"epistemico":"fato","origem":"wiki","rel":"relaciona"},{"alvo":"tipos_de_interpretante","confianca":0.5,"epistemico":"fato","origem":"wiki","rel":"relaciona"},{"alvo":"ficheiros","confianca":0.5,"epistemico":"fato","origem":"wiki","rel":"relaciona"},{"alvo":"veredito_experimental","confianca":0.5,"epistemico":"fato","origem":"wiki","rel":"relaciona"},{"alvo":"orbita_pell_hardware","confianca":0.5,"epistemico":"fato","origem":"wiki","rel":"relaciona"},{"alvo":"beta_numeracao_metalica","confianca":1.0,"epistemico":"fato","origem":"wiki","rel":"confirma"}],"confianca":1.0,"contraexemplos":[],"descricao":"Fit cego x(n+2)=a₁·x(n+1)+a₀·x(n) **sem** mencionar Pell:","epistemico":"fato","exemplos":[],"id":"experimento_3_inverter","memoria":"semantica","meta":{"dominio":"manual","nivel":6,"verificacao":"slug idêntico — enriquecer, não duplicar"},"origem":"METAIS.md","palavras_chave":[],"sinonimos":[],"tipo":"conceito","titulo":"Experimento 3 — Inverter"}
\section{Experimento 3 — Inverter}
Fit cego x(n+2)=a\ensuremath{_{1}}\textperiodcentered x(n+1)+a\ensuremath{_{0}}\textperiodcentered x(n) **sem** mencionar Pell:
\prov{relações: parte\_de metais; relaciona transcricao; relaciona utilitarismo; relaciona ruido; relaciona virtualizacao; relaciona pedagogia; relaciona vida-morte; relaciona pow\_metal\_ressonancia\_val; relaciona mecanica; relaciona scheduler; relaciona resultado\_anterior\_identidade\_nao\_cobertura; relaciona modo\_espectral\_spect; relaciona paper\_2\_reta\_ouro; relaciona resultado\_completo; relaciona pitagorismo; relaciona tipos\_de\_interpretante; relaciona ficheiros; relaciona veredito\_experimental; relaciona orbita\_pell\_hardware; confirma beta\_numeracao\_metalica}
\prov{proveniência: METAIS.md \textperiodcentered{} fato \textperiodcentered{} confiança 1.0 \textperiodcentered{} domínio manual \textperiodcentered{} história 22 versões no jornal}

%CRISTAL {"arestas":[{"alvo":"colapso","confianca":0.95,"epistemico":"fato","origem":"paper","rel":"parte_de"}],"confianca":1.0,"contraexemplos":[],"descricao":"- amostras Δ: **17** - Δ médio: **1.96e-17** - variância: **3.44e-33** - |Δ| máximo: **2.22e-16** - fit ordem 2: a₁=0.000, a₀=0.000, n≈0 - parece aleatório? **não** - compatível Pell? **não**","epistemico":"fato","exemplos":[],"id":"experimento_3_o_erro_possui_estrutura","memoria":"semantica","meta":{"dominio":"manual","nivel":6,"verificacao":"slug idêntico — enriquecer, não duplicar"},"origem":"RUÍDO.md","palavras_chave":[],"sinonimos":[],"tipo":"conceito","titulo":"Experimento 3 — O erro possui estrutura?"}
\section{Experimento 3 — O erro possui estrutura?}
- amostras \ensuremath{\Delta}: **17** - \ensuremath{\Delta} médio: **1.96e-17** - variância: **3.44e-33** - |\ensuremath{\Delta}| máximo: **2.22e-16** - fit ordem 2: a\ensuremath{_{1}}=0.000, a\ensuremath{_{0}}=0.000, n\ensuremath{\approx}0 - parece aleatório? **não** - compatível Pell? **não**
\prov{relações: parte\_de colapso}
\prov{proveniência: RUÍDO.md \textperiodcentered{} fato \textperiodcentered{} confiança 1.0 \textperiodcentered{} domínio manual \textperiodcentered{} história 3 versões no jornal}

%CRISTAL {"arestas":[{"alvo":"pow_espectral_investigacao_broca","confianca":1.0,"epistemico":"fato","origem":"POW_ESPECTRAL.md","rel":"parte_de"},{"alvo":"vetores","confianca":0.5,"epistemico":"fato","origem":"wiki","rel":"relaciona"},{"alvo":"virtualizacao","confianca":0.5,"epistemico":"fato","origem":"wiki","rel":"relaciona"},{"alvo":"word","confianca":0.5,"epistemico":"fato","origem":"wiki","rel":"relaciona"}],"confianca":1.0,"contraexemplos":[],"descricao":"Seed primeiro candidato: `[86, 38]` · horizonte metal estável: **n=1**","epistemico":"fato","exemplos":[],"id":"experimento_3_orbitas","memoria":"semantica","meta":{"dominio":"manual","nivel":6,"verificacao":"slug idêntico — enriquecer, não duplicar"},"origem":"POW_ESPECTRAL.md","palavras_chave":[],"sinonimos":[],"tipo":"conceito","titulo":"Experimento 3 — Órbitas"}
\section{Experimento 3 — Órbitas}
Seed primeiro candidato: `[86, 38]` \textperiodcentered  horizonte metal estável: **n=1**
\prov{relações: parte\_de pow\_espectral\_investigacao\_broca; relaciona vetores; relaciona virtualizacao; relaciona word}
\prov{proveniência: POW\_ESPECTRAL.md \textperiodcentered{} fato \textperiodcentered{} confiança 1.0 \textperiodcentered{} domínio manual \textperiodcentered{} história 9 versões no jornal}

%CRISTAL {"arestas":[{"alvo":"pow_espectral_investigacao_broca","confianca":1.0,"epistemico":"fato","origem":"POW_ESPECTRAL.md","rel":"parte_de"}],"confianca":1.0,"contraexemplos":[],"descricao":"- variância total: **2.75** · variância e: **1.25** - autocorr total (lag 1): **0.6436** - pares calor adjacentes próximos: **1.0000**","epistemico":"fato","exemplos":[],"id":"experimento_46_degeneracoes_invariantes_simetria","memoria":"semantica","meta":{"dominio":"manual","nivel":6,"verificacao":"slug idêntico — enriquecer, não duplicar"},"origem":"POW_ESPECTRAL.md","palavras_chave":[],"sinonimos":[],"tipo":"conceito","titulo":"Experimento 4–6 — Degenerações, invariantes, simetria"}
\section{Experimento 4–6 — Degenerações, invariantes, simetria}
- variância total: **2.75** \textperiodcentered  variância e: **1.25** - autocorr total (lag 1): **0.6436** - pares calor adjacentes próximos: **1.0000**
\prov{relações: parte\_de pow\_espectral\_investigacao\_broca}
\prov{proveniência: POW\_ESPECTRAL.md \textperiodcentered{} fato \textperiodcentered{} confiança 1.0 \textperiodcentered{} domínio manual \textperiodcentered{} história 6 versões no jornal}

%CRISTAL {"arestas":[{"alvo":"colapso","confianca":0.95,"epistemico":"fato","origem":"paper","rel":"parte_de"},{"alvo":"word","confianca":0.5,"epistemico":"fato","origem":"wiki","rel":"relaciona"},{"alvo":"vontade_de_poder","confianca":0.5,"epistemico":"fato","origem":"wiki","rel":"relaciona"},{"alvo":"while","confianca":0.5,"epistemico":"fato","origem":"wiki","rel":"relaciona"},{"alvo":"virtualizacao","confianca":0.5,"epistemico":"fato","origem":"wiki","rel":"relaciona"}],"confianca":1.0,"contraexemplos":[],"descricao":"Órbita prata A₂·(1,0) com ruído `offset` crescente. Cada passo regista o par ideal (sobre a+c²=1) e o perturbado.","epistemico":"fato","exemplos":[],"id":"experimento_5_trajetorias_no_plano_calor_fase","memoria":"semantica","meta":{"dominio":"manual","nivel":6,"verificacao":"slug idêntico — enriquecer, não duplicar"},"origem":"RUÍDO.md","palavras_chave":[],"sinonimos":[],"tipo":"conceito","titulo":"Experimento 5 — Trajetórias no plano (calor, fase)"}
\section{Experimento 5 — Trajetórias no plano (calor, fase)}
Órbita prata A\ensuremath{_{2}}\textperiodcentered (1,0) com ruído `offset` crescente. Cada passo regista o par ideal (sobre a+c\ensuremath{^{2}}=1) e o perturbado.
\prov{relações: parte\_de colapso; relaciona word; relaciona vontade\_de\_poder; relaciona while; relaciona virtualizacao}
\prov{proveniência: RUÍDO.md \textperiodcentered{} fato \textperiodcentered{} confiança 1.0 \textperiodcentered{} domínio manual \textperiodcentered{} história 7 versões no jornal}

%CRISTAL {"arestas":[{"alvo":"pow_espectral_investigacao_broca","confianca":1.0,"epistemico":"fato","origem":"POW_ESPECTRAL.md","rel":"parte_de"}],"confianca":1.0,"contraexemplos":[],"descricao":"| método | entropia calor | hits | taxa | |--------|----------------|------|------| | espectral (neurônio) | 1.4612 | 2048 | 1.0000 | | linear (nonce→word) | 2.6360 | 954 | 0.4658 |","epistemico":"fato","exemplos":[],"id":"experimento_7_comparacao_spectral_vs_linear","memoria":"semantica","meta":{"dominio":"manual","nivel":6,"verificacao":"slug idêntico — enriquecer, não duplicar"},"origem":"POW_ESPECTRAL.md","palavras_chave":[],"sinonimos":[],"tipo":"conceito","titulo":"Experimento 7 — Comparação spectral vs linear"}
\section{Experimento 7 — Comparação spectral vs linear}
| método | entropia calor | hits | taxa | |--------|----------------|------|------| | espectral (neurônio) | 1.4612 | 2048 | 1.0000 | | linear (nonce\ensuremath{\to}word) | 2.6360 | 954 | 0.4658 |
\prov{relações: parte\_de pow\_espectral\_investigacao\_broca}
\prov{proveniência: POW\_ESPECTRAL.md \textperiodcentered{} fato \textperiodcentered{} confiança 1.0 \textperiodcentered{} domínio manual \textperiodcentered{} história 6 versões no jornal}

%CRISTAL {"arestas":[{"alvo":"construcao_vs_crivo","confianca":1.0,"epistemico":"fato","origem":"wiki","rel":"usa"},{"alvo":"verdade_por_refutacao_gentil","confianca":1.0,"epistemico":"fato","origem":"wiki","rel":"usa"}],"confianca":1.0,"contraexemplos":[],"descricao":"Toda alegação vem acompanhada do experimento mínimo (numpy) que a decide ou refuta — o experimento que promove um item é sua 'certidão de nascimento'; distinguem-se sempre demonstrado / verificado / conjectural.","epistemico":"fato","exemplos":[],"id":"experimento_minimo_decide","memoria":"semantica","meta":{"autor":"Aarão/broca-so","dominio":"broca_repos","fonte":""},"origem":"wiki","palavras_chave":[],"sinonimos":[],"tipo":"conceito","titulo":"O Experimento Mínimo que Decide (certidão de nascimento)"}
\section{O Experimento Mínimo que Decide (certidão de nascimento)}
Toda alegação vem acompanhada do experimento mínimo (numpy) que a decide ou refuta — o experimento que promove um item é sua 'certidão de nascimento'; distinguem-se sempre demonstrado / verificado / conjectural.
\prov{relações: usa construcao\_vs\_crivo; usa verdade\_por\_refutacao\_gentil}
\prov{proveniência: wiki \textperiodcentered{} fato \textperiodcentered{} confiança 1.0 \textperiodcentered{} domínio broca\_repos \textperiodcentered{} história 3 versões no jornal}

%CRISTAL {"arestas":[{"alvo":"computacao","confianca":1.0,"epistemico":"fato","origem":"manual","rel":"referencia"},{"alvo":"goldchain_mercado_hub","confianca":0.5,"epistemico":"fato","origem":"wiki","rel":"relaciona"},{"alvo":"utilitarismo","confianca":0.5,"epistemico":"fato","origem":"wiki","rel":"relaciona"},{"alvo":"topologia","confianca":0.5,"epistemico":"fato","origem":"wiki","rel":"relaciona"},{"alvo":"nao_dualidade","confianca":0.5,"epistemico":"fato","origem":"wiki","rel":"relaciona"},{"alvo":"testes","confianca":0.5,"epistemico":"fato","origem":"wiki","rel":"relaciona"},{"alvo":"regua_associativa_e_o_risco_de_vazamento","confianca":0.5,"epistemico":"fato","origem":"wiki","rel":"relaciona"},{"alvo":"fase_2_tecido_de_conhecimento","confianca":0.5,"epistemico":"fato","origem":"wiki","rel":"relaciona"}],"confianca":1.0,"contraexemplos":[],"descricao":"**Pergunta:** qual é o poder computacional desta máquina?","epistemico":"fato","exemplos":[],"id":"expressividade","memoria":"semantica","meta":{"dominio":"manual","duplicata_sugerida":[],"nivel":6,"verificacao":"parte de conceito mais amplo 'expressividade'"},"origem":"EXPRESSIVIDADE.md","palavras_chave":[],"sinonimos":[],"tipo":"conceito","titulo":"EXPRESSIVIDADE"}
\section{EXPRESSIVIDADE}
**Pergunta:** qual é o poder computacional desta máquina?
\prov{relações: referencia computacao; relaciona goldchain\_mercado\_hub; relaciona utilitarismo; relaciona topologia; relaciona nao\_dualidade; relaciona testes; relaciona regua\_associativa\_e\_o\_risco\_de\_vazamento; relaciona fase\_2\_tecido\_de\_conhecimento}
\prov{proveniência: EXPRESSIVIDADE.md \textperiodcentered{} fato \textperiodcentered{} confiança 1.0 \textperiodcentered{} domínio manual \textperiodcentered{} história 13 versões no jornal}

%CRISTAL {"arestas":[{"alvo":"face_turno_fail_closed","confianca":1.0,"epistemico":"fato","origem":"wiki","rel":"parte_de"}],"confianca":1.0,"contraexemplos":[],"descricao":"Tags com :rNN (_face_only_tag) marcam a face isolada da Tiffany — não são turno 'usuário — resposta'; _turno_parts as trata como linha simples e score_hit as descarta com score −999, de modo que só a linha-turno completa pode casar com a fala do usuário. Complementa o fail-closed: só o turno duplo válido colapsa.","epistemico":"fato","exemplos":[],"id":"face_only_r_line","memoria":"semantica","meta":{"autor":"broca-so","dominio":"conversa","fonte":"conversa/colapso.py:_face_only_tag/_turno_parts/score_hit"},"origem":"wiki","palavras_chave":[],"sinonimos":[],"tipo":"conceito","titulo":"Tiffany — linhas face-só excluídas do casamento de turno"}
\section{Tiffany — linhas face-só excluídas do casamento de turno}
Tags com :rNN (\_face\_only\_tag) marcam a face isolada da Tiffany — não são turno 'usuário — resposta'; \_turno\_parts as trata como linha simples e score\_hit as descarta com score \ensuremath{-}999, de modo que só a linha-turno completa pode casar com a fala do usuário. Complementa o fail-closed: só o turno duplo válido colapsa.
\prov{relações: parte\_de face\_turno\_fail\_closed}
\prov{proveniência: wiki \textperiodcentered{} conversa/colapso.py:\_face\_only\_tag/\_turno\_parts/score\_hit \textperiodcentered{} fato \textperiodcentered{} confiança 1.0 \textperiodcentered{} domínio conversa \textperiodcentered{} história 10 versões no jornal}

%CRISTAL {"arestas":[{"alvo":"regua_de_gentil","confianca":1.0,"epistemico":"fato","origem":"paper","rel":"relaciona"},{"alvo":"invariancia_de_regua_arx","confianca":1.0,"epistemico":"fato","origem":"wiki","rel":"usa"},{"alvo":"cifra","confianca":1.0,"epistemico":"fato","origem":"wiki","rel":"relaciona"}],"confianca":1.0,"contraexemplos":[],"descricao":"Uma família dGb interpolando entre Hamming (b=1) e circular modular (b=n), onde b é o 'alcance de carry' que a régua enxerga; a isometria da adição cresce e a de XOR/rotação decresce monotonicamente, com compromisso máximo de apenas 0,19 em b=4 — nenhuma régua torna adição e XOR simultaneamente rígidas.","epistemico":"fato","exemplos":[],"id":"familia_metricas_gentil","memoria":"semantica","meta":{"autor":"Aarão/broca-so","dominio":"broca_repos","fonte":""},"origem":"wiki","palavras_chave":[],"sinonimos":[],"tipo":"conceito","titulo":"A Família Contínua de Métricas de Lopes"}
\section{A Família Contínua de Métricas de Lopes}
Uma família dGb interpolando entre Hamming (b=1) e circular modular (b=n), onde b é o 'alcance de carry' que a régua enxerga; a isometria da adição cresce e a de XOR/rotação decresce monotonicamente, com compromisso máximo de apenas 0,19 em b=4 — nenhuma régua torna adição e XOR simultaneamente rígidas.
\prov{relações: relaciona regua\_de\_gentil; usa invariancia\_de\_regua\_arx; relaciona cifra}
\prov{proveniência: wiki \textperiodcentered{} fato \textperiodcentered{} confiança 1.0 \textperiodcentered{} domínio broca\_repos \textperiodcentered{} história 8 versões no jornal}

%CRISTAL {"arestas":[{"alvo":"geometria","confianca":0.95,"epistemico":"fato","origem":"wiki","rel":"confirma"},{"alvo":"gato","confianca":0.95,"epistemico":"fato","origem":"wiki","rel":"confirma"},{"alvo":"mecanica_quantica","confianca":0.95,"epistemico":"fato","origem":"wiki","rel":"confirma"},{"alvo":"a_funcao_de_onda_a_esfera_de_bloch","confianca":0.95,"epistemico":"fato","origem":"wiki","rel":"confirma"},{"alvo":"a_dinamica_a_precessao_de_bloch","confianca":0.95,"epistemico":"fato","origem":"wiki","rel":"confirma"}],"confianca":0.95,"contraexemplos":[],"descricao":"Referência verificada: fase_geometrica_berry.md","epistemico":"fato","exemplos":[],"id":"fase_geometrica_berry","memoria":"semantica","meta":{"dominio":"referencia","fonte":"web","nivel":5},"origem":"wiki","palavras_chave":[],"sinonimos":[],"tipo":"conceito","titulo":"fase geometrica berry"}
\section{fase geometrica berry}
Referência verificada: fase\_geometrica\_berry.md
\prov{relações: confirma geometria; confirma gato; confirma mecanica\_quantica; confirma a\_funcao\_de\_onda\_a\_esfera\_de\_bloch; confirma a\_dinamica\_a\_precessao\_de\_bloch}
\prov{proveniência: wiki \textperiodcentered{} web \textperiodcentered{} fato \textperiodcentered{} confiança 0.95 \textperiodcentered{} domínio referencia \textperiodcentered{} história 35 versões no jornal}

%CRISTAL {"arestas":[{"alvo":"fase_geometrica_berry_fase_geometrica_de_berry_pancharatnamberry","confianca":1.0,"epistemico":"fato","origem":"fase_geometrica_berry.md","rel":"parte_de"},{"alvo":"virtualizacao","confianca":0.5,"epistemico":"fato","origem":"wiki","rel":"relaciona"},{"alvo":"word","confianca":0.5,"epistemico":"fato","origem":"wiki","rel":"relaciona"}],"confianca":1.0,"contraexemplos":[],"descricao":"Atribuir fase Berry a toros φ-Dougherty sem derivação — extrapolação.","epistemico":"fato","exemplos":[],"id":"fase_geometrica_berry_contraexemplos","memoria":"semantica","meta":{"dominio":"referencia","fonte":"web","nivel":5},"origem":"fase_geometrica_berry.md","palavras_chave":[],"sinonimos":[],"tipo":"conceito","titulo":"Contraexemplos"}
\section{Contraexemplos}
Atribuir fase Berry a toros \ensuremath{\varphi}-Dougherty sem derivação — extrapolação.
\prov{relações: parte\_de fase\_geometrica\_berry\_fase\_geometrica\_de\_berry\_pancharatnamberry; relaciona virtualizacao; relaciona word}
\prov{proveniência: fase\_geometrica\_berry.md \textperiodcentered{} web \textperiodcentered{} fato \textperiodcentered{} confiança 1.0 \textperiodcentered{} domínio referencia \textperiodcentered{} história 5 versões no jornal}

%CRISTAL {"fusao":[{"arestas":[{"alvo":"fase_geometrica_berry_fase_geometrica_de_berry_pancharatnamberry","confianca":1.0,"epistemico":"fato","origem":"fase_geometrica_berry.md","rel":"parte_de"},{"alvo":"observador_e_observado","confianca":0.95,"epistemico":"fato","origem":"wiki","rel":"referencia"}],"confianca":1.0,"contraexemplos":[],"descricao":"| Tipo | Depende de | Observador | |------|------------|------------| | Dinâmica | tempo, energia E_n | evolução temporal | | Geométrica | forma do ciclo C em R | geometria do processo |","epistemico":"fato","exemplos":[],"id":"fase_geometrica_berry_fase_dinamica_vs_geometrica","memoria":"semantica","meta":{"dominio":"referencia","fonte":"web","nivel":5},"origem":"fase_geometrica_berry.md","palavras_chave":[],"sinonimos":[],"tipo":"conceito","titulo":"Fase dinâmica vs geométrica"},{"arestas":[{"alvo":"fase_geometrica_de_berry_pancharatnamberry","confianca":1.0,"epistemico":"fato","origem":"fase_geometrica_berry.md","rel":"parte_de"},{"alvo":"observador_e_observado","confianca":0.95,"epistemico":"fato","origem":"wiki","rel":"referencia"}],"confianca":1.0,"contraexemplos":[],"descricao":"| Tipo | Depende de | Observador | |------|------------|------------| | Dinâmica | tempo, energia E_n | evolução temporal | | Geométrica | forma do ciclo C em R | geometria do processo |","epistemico":"fato","exemplos":[],"id":"fase_dinamica_vs_geometrica","memoria":"semantica","meta":{"dominio":"referencia","nivel":5},"origem":"fase_geometrica_berry.md","palavras_chave":[],"sinonimos":[],"tipo":"conceito","titulo":"Fase dinâmica vs geométrica"}],"id":"fase_geometrica_berry_fase_dinamica_vs_geometrica","tipo":"conceito"}
\section{Fase dinâmica vs geométrica}
| Tipo | Depende de | Observador | |------|------------|------------| | Dinâmica | tempo, energia E\_n | evolução temporal | | Geométrica | forma do ciclo C em R | geometria do processo |
\prov{relações: parte\_de fase\_geometrica\_berry\_fase\_geometrica\_de\_berry\_pancharatnamberry; referencia observador\_e\_observado}
\prov{proveniência: fase\_geometrica\_berry.md \textperiodcentered{} web \textperiodcentered{} fato \textperiodcentered{} confiança 1.0 \textperiodcentered{} domínio referencia \textperiodcentered{} história 4 versões no jornal}
\prov{fusão de curadoria (texto idêntico): absorve fase\_dinamica\_vs\_geometrica --- o mesmo conteúdo por dois esquemas de endereço}

%CRISTAL {"fusao":[{"arestas":[{"alvo":"fase_geometrica_berry","confianca":0.95,"epistemico":"fato","origem":"wiki","rel":"parte_de"}],"confianca":1.0,"contraexemplos":[],"descricao":"> **Epistemologia:** fato matemático-físico estabelecido (MQ + geometria). > **No grafo:** `epistemico=fato`, confiança 0.95. > **Papel:** formalizar fase adquirida em ciclos adiabáticos na esfera de Bloch; complementa Hopf.","epistemico":"fato","exemplos":[],"id":"fase_geometrica_berry_fase_geometrica_de_berry_pancharatnamberry","memoria":"semantica","meta":{"dominio":"referencia","fonte":"web","nivel":5},"origem":"fase_geometrica_berry.md","palavras_chave":[],"sinonimos":[],"tipo":"conceito","titulo":"Fase geométrica de Berry (Pancharatnam–Berry)"},{"arestas":[{"alvo":"fase_geometrica_berry","confianca":0.95,"epistemico":"fato","origem":"wiki","rel":"parte_de"}],"confianca":1.0,"contraexemplos":[],"descricao":"> **Epistemologia:** fato matemático-físico estabelecido (MQ + geometria). > **No grafo:** `epistemico=fato`, confiança 0.95. > **Papel:** formalizar fase adquirida em ciclos adiabáticos na esfera de Bloch; complementa Hopf.","epistemico":"fato","exemplos":[],"id":"fase_geometrica_de_berry_pancharatnamberry","memoria":"semantica","meta":{"dominio":"referencia","nivel":5},"origem":"fase_geometrica_berry.md","palavras_chave":[],"sinonimos":[],"tipo":"conceito","titulo":"Fase geométrica de Berry (Pancharatnam–Berry)"}],"id":"fase_geometrica_berry_fase_geometrica_de_berry_pancharatnamberry","tipo":"conceito"}
\section{Fase geométrica de Berry (Pancharatnam–Berry)}
> **Epistemologia:** fato matemático-físico estabelecido (MQ + geometria). > **No grafo:** `epistemico=fato`, confiança 0.95. > **Papel:** formalizar fase adquirida em ciclos adiabáticos na esfera de Bloch; complementa Hopf.
\prov{relações: parte\_de fase\_geometrica\_berry}
\prov{proveniência: fase\_geometrica\_berry.md \textperiodcentered{} web \textperiodcentered{} fato \textperiodcentered{} confiança 1.0 \textperiodcentered{} domínio referencia \textperiodcentered{} história 3 versões no jornal}
\prov{fusão de curadoria (texto idêntico): absorve fase\_geometrica\_de\_berry\_pancharatnamberry --- o mesmo conteúdo por dois esquemas de endereço}

%CRISTAL {"fusao":[{"arestas":[{"alvo":"fase_geometrica_berry_fase_geometrica_de_berry_pancharatnamberry","confianca":1.0,"epistemico":"fato","origem":"fase_geometrica_berry.md","rel":"parte_de"}],"confianca":1.0,"contraexemplos":[],"descricao":"Quando o Hamiltoniano H(R) evolui **cyclicamente e adiabaticamente** no espaço de parâmetros R, o estado |ψ⟩ adquire factor de fase composto:","epistemico":"fato","exemplos":[],"id":"fase_geometrica_berry_fase_pancharatnamberry","memoria":"semantica","meta":{"dominio":"referencia","fonte":"web","nivel":5},"origem":"fase_geometrica_berry.md","palavras_chave":[],"sinonimos":[],"tipo":"conceito","titulo":"Fase Pancharatnam–Berry"},{"arestas":[{"alvo":"fase_geometrica_de_berry_pancharatnamberry","confianca":1.0,"epistemico":"fato","origem":"fase_geometrica_berry.md","rel":"parte_de"}],"confianca":1.0,"contraexemplos":[],"descricao":"Quando o Hamiltoniano H(R) evolui **cyclicamente e adiabaticamente** no espaço de parâmetros R, o estado |ψ⟩ adquire factor de fase composto:","epistemico":"fato","exemplos":[],"id":"fase_pancharatnamberry","memoria":"semantica","meta":{"dominio":"referencia","nivel":5},"origem":"fase_geometrica_berry.md","palavras_chave":[],"sinonimos":[],"tipo":"conceito","titulo":"Fase Pancharatnam–Berry"}],"id":"fase_geometrica_berry_fase_pancharatnamberry","tipo":"conceito"}
\section{Fase Pancharatnam–Berry}
Quando o Hamiltoniano H(R) evolui **cyclicamente e adiabaticamente** no espaço de parâmetros R, o estado |\ensuremath{\psi}\ensuremath{\rangle} adquire factor de fase composto:
\prov{relações: parte\_de fase\_geometrica\_berry\_fase\_geometrica\_de\_berry\_pancharatnamberry}
\prov{proveniência: fase\_geometrica\_berry.md \textperiodcentered{} web \textperiodcentered{} fato \textperiodcentered{} confiança 1.0 \textperiodcentered{} domínio referencia \textperiodcentered{} história 3 versões no jornal}
\prov{fusão de curadoria (texto idêntico): absorve fase\_pancharatnamberry --- o mesmo conteúdo por dois esquemas de endereço}

%CRISTAL {"arestas":[{"alvo":"fase_geometrica_berry_fase_geometrica_de_berry_pancharatnamberry","confianca":1.0,"epistemico":"fato","origem":"fase_geometrica_berry.md","rel":"parte_de"}],"confianca":1.0,"contraexemplos":[],"descricao":"- https://en.wikipedia.org/wiki/Geometric_phase - https://doi.org/10.1098/rspa.1984.0023 - https://materia.fisica.unimi.it/manini/berryphase.html","epistemico":"fato","exemplos":[],"id":"fase_geometrica_berry_fontes","memoria":"semantica","meta":{"dominio":"referencia","fonte":"web","nivel":5},"origem":"fase_geometrica_berry.md","palavras_chave":[],"sinonimos":[],"tipo":"conceito","titulo":"Fontes"}
\section{Fontes}
- https://en.wikipedia.org/wiki/Geometric\_phase - https://doi.org/10.1098/rspa.1984.0023 - https://materia.fisica.unimi.it/manini/berryphase.html
\prov{relações: parte\_de fase\_geometrica\_berry\_fase\_geometrica\_de\_berry\_pancharatnamberry}
\prov{proveniência: fase\_geometrica\_berry.md \textperiodcentered{} web \textperiodcentered{} fato \textperiodcentered{} confiança 1.0 \textperiodcentered{} domínio referencia \textperiodcentered{} história 3 versões no jornal}

%CRISTAL {"arestas":[{"alvo":"fase_geometrica_berry_fase_geometrica_de_berry_pancharatnamberry","confianca":1.0,"epistemico":"fato","origem":"fase_geometrica_berry.md","rel":"parte_de"},{"alvo":"a_medida_o_calor_de_volta","confianca":0.95,"epistemico":"fato","origem":"wiki","rel":"confirma"},{"alvo":"a_dinamica_a_precessao_de_bloch","confianca":0.95,"epistemico":"fato","origem":"wiki","rel":"confirma"}],"confianca":1.0,"contraexemplos":[],"descricao":"| Berry (fato) | Nó Broca-SO | relação | |--------------|-------------|---------| | Bloch S² ciclo fechado | a_funcao_de_onda_a_esfera_de_bloch | confirma | | holonomia / curvatura | a_dinamica_a_precessao_de_bloch | confirma | | fibrado estados | a_medida_o_calor_de_volta | confirma | | qubit / fase | gato, mecanica_quantica | confirma | | S³→S² Hopf | fibracao_hopf_bloch | confirma | | ciclo adiabático | observador_e_observado | referencia (processo de medição) |","epistemico":"fato","exemplos":[],"id":"fase_geometrica_berry_ponte_broca-so","memoria":"semantica","meta":{"dominio":"referencia","fonte":"web","nivel":5},"origem":"fase_geometrica_berry.md","palavras_chave":[],"sinonimos":[],"tipo":"conceito","titulo":"Ponte Broca-SO"}
\section{Ponte Broca-SO}
| Berry (fato) | Nó Broca-SO | relação | |--------------|-------------|---------| | Bloch S\ensuremath{^{2}} ciclo fechado | a\_funcao\_de\_onda\_a\_esfera\_de\_bloch | confirma | | holonomia / curvatura | a\_dinamica\_a\_precessao\_de\_bloch | confirma | | fibrado estados | a\_medida\_o\_calor\_de\_volta | confirma | | qubit / fase | gato, mecanica\_quantica | confirma | | S\ensuremath{^{3}}\ensuremath{\to}S\ensuremath{^{2}} Hopf | fibracao\_hopf\_bloch | confirma | | ciclo adiabático | observador\_e\_observado | referencia (processo de medição) |
\prov{relações: parte\_de fase\_geometrica\_berry\_fase\_geometrica\_de\_berry\_pancharatnamberry; confirma a\_medida\_o\_calor\_de\_volta; confirma a\_dinamica\_a\_precessao\_de\_bloch}
\prov{proveniência: fase\_geometrica\_berry.md \textperiodcentered{} web \textperiodcentered{} fato \textperiodcentered{} confiança 1.0 \textperiodcentered{} domínio referencia \textperiodcentered{} história 5 versões no jornal}

%CRISTAL {"arestas":[{"alvo":"fase_geometrica_berry_fase_geometrica_de_berry_pancharatnamberry","confianca":1.0,"epistemico":"fato","origem":"fase_geometrica_berry.md","rel":"parte_de"}],"confianca":1.0,"contraexemplos":[],"descricao":"- **S. Pancharatnam (1956):** fase geométrica em óptica clássica - **Michael Berry (1984):** *Quantal phase factors accompanying adiabatic changes*, Proc. R. Soc. A 392, 45–57 — [doi:10.1098/rspa.1984.0023](https://doi.org/10.1098/rspa.1984.0023) - **Wikipedia:** [Geometric phase](https://en.wikipedia.org/wiki/Geometric_phase) - **Revisão:** [Berry's geometric phase](https://materia.fisica.unimi.it/manini/berryphase.html) (Manini)","epistemico":"fato","exemplos":[],"id":"fase_geometrica_berry_referencia","memoria":"semantica","meta":{"dominio":"referencia","fonte":"web","nivel":5},"origem":"fase_geometrica_berry.md","palavras_chave":[],"sinonimos":[],"tipo":"conceito","titulo":"Referência"}
\section{Referência}
- **S. Pancharatnam (1956):** fase geométrica em óptica clássica - **Michael Berry (1984):** *Quantal phase factors accompanying adiabatic changes*, Proc. R. Soc. A 392, 45–57 — [doi:10.1098/rspa.1984.0023](https://doi.org/10.1098/rspa.1984.0023) - **Wikipedia:** [Geometric phase](https://en.wikipedia.org/wiki/Geometric\_phase) - **Revisão:** [Berry's geometric phase](https://materia.fisica.unimi.it/manini/berryphase.html) (Manini)
\prov{relações: parte\_de fase\_geometrica\_berry\_fase\_geometrica\_de\_berry\_pancharatnamberry}
\prov{proveniência: fase\_geometrica\_berry.md \textperiodcentered{} web \textperiodcentered{} fato \textperiodcentered{} confiança 1.0 \textperiodcentered{} domínio referencia \textperiodcentered{} história 3 versões no jornal}

%CRISTAL {"fusao":[{"arestas":[{"alvo":"fase_geometrica_berry_fase_geometrica_de_berry_pancharatnamberry","confianca":1.0,"epistemico":"fato","origem":"fase_geometrica_berry.md","rel":"parte_de"},{"alvo":"ponte_quantica_dougherty","confianca":0.95,"epistemico":"fato","origem":"wiki","rel":"referencia"},{"alvo":"conjunto_dougherty","confianca":0.95,"epistemico":"fato","origem":"wiki","rel":"analogia"}],"confianca":1.0,"contraexemplos":[],"descricao":"MQ standard + Berry descrevem fases em ciclos adiabáticos **sem** coordenadas toroidais φ.","epistemico":"fato","exemplos":[],"id":"fase_geometrica_berry_relacao_com_hipotese_dougherty","memoria":"semantica","meta":{"dominio":"referencia","fonte":"web","nivel":5},"origem":"fase_geometrica_berry.md","palavras_chave":[],"sinonimos":[],"tipo":"conceito","titulo":"Relação com hipótese Dougherty"},{"arestas":[{"alvo":"microscopia_de_fotoionizacao_orbitais_de_hidrogenio_stodolna_2013","confianca":1.0,"epistemico":"fato","origem":"microscopia_fotoionizacao_hidrogenio.md","rel":"parte_de"},{"alvo":"fibracao_de_hopf_e_esfera_de_bloch","confianca":1.0,"epistemico":"fato","origem":"fibracao_hopf_bloch.md","rel":"parte_de"},{"alvo":"semiotica_triadica_de_charles_sanders_peirce","confianca":1.0,"epistemico":"fato","origem":"semiotica_peirce_triade.md","rel":"parte_de"},{"alvo":"fase_geometrica_de_berry_pancharatnamberry","confianca":1.0,"epistemico":"fato","origem":"fase_geometrica_berry.md","rel":"parte_de"},{"alvo":"observador_e_observado","confianca":0.95,"epistemico":"fato","origem":"wiki","rel":"analogia"},{"alvo":"ponte_quantica_dougherty","confianca":0.95,"epistemico":"fato","origem":"wiki","rel":"referencia"},{"alvo":"conjunto_dougherty","confianca":0.95,"epistemico":"fato","origem":"wiki","rel":"analogia"}],"confianca":1.0,"contraexemplos":[],"descricao":"MQ standard + Berry descrevem fases em ciclos adiabáticos **sem** coordenadas toroidais φ.","epistemico":"fato","exemplos":[],"id":"relacao_com_hipotese_dougherty","memoria":"semantica","meta":{"dominio":"referencia","nivel":5},"origem":"fase_geometrica_berry.md","palavras_chave":[],"sinonimos":[],"tipo":"conceito","titulo":"Relação com hipótese Dougherty"}],"id":"fase_geometrica_berry_relacao_com_hipotese_dougherty","tipo":"conceito"}
\section{Relação com hipótese Dougherty}
MQ standard + Berry descrevem fases em ciclos adiabáticos **sem** coordenadas toroidais \ensuremath{\varphi}.
\prov{relações: parte\_de fase\_geometrica\_berry\_fase\_geometrica\_de\_berry\_pancharatnamberry; referencia ponte\_quantica\_dougherty; analogia conjunto\_dougherty}
\prov{proveniência: fase\_geometrica\_berry.md \textperiodcentered{} web \textperiodcentered{} fato \textperiodcentered{} confiança 1.0 \textperiodcentered{} domínio referencia \textperiodcentered{} história 5 versões no jornal}
\prov{fusão de curadoria (texto idêntico): absorve relacao\_com\_hipotese\_dougherty --- o mesmo conteúdo por dois esquemas de endereço}

%CRISTAL {"fusao":[{"arestas":[{"alvo":"fase_geometrica_berry_fase_geometrica_de_berry_pancharatnamberry","confianca":1.0,"epistemico":"fato","origem":"fase_geometrica_berry.md","rel":"parte_de"},{"alvo":"a_funcao_de_onda_a_esfera_de_bloch","confianca":0.95,"epistemico":"fato","origem":"wiki","rel":"confirma"},{"alvo":"fibracao_hopf_bloch","confianca":0.95,"epistemico":"fato","origem":"wiki","rel":"confirma"}],"confianca":1.0,"contraexemplos":[],"descricao":"A fase geométrica é anholonomia de **transporte paralelo** no fibrado de estados sobre o espaço de parâmetros.","epistemico":"fato","exemplos":[],"id":"fase_geometrica_berry_transporte_paralelo_e_bloch_sphere","memoria":"semantica","meta":{"dominio":"referencia","fonte":"web","nivel":5},"origem":"fase_geometrica_berry.md","palavras_chave":[],"sinonimos":[],"tipo":"conceito","titulo":"Transporte paralelo e Bloch sphere"},{"arestas":[{"alvo":"fase_geometrica_de_berry_pancharatnamberry","confianca":1.0,"epistemico":"fato","origem":"fase_geometrica_berry.md","rel":"parte_de"},{"alvo":"a_funcao_de_onda_a_esfera_de_bloch","confianca":0.95,"epistemico":"fato","origem":"wiki","rel":"confirma"},{"alvo":"fibracao_hopf_bloch","confianca":0.95,"epistemico":"fato","origem":"wiki","rel":"confirma"}],"confianca":1.0,"contraexemplos":[],"descricao":"A fase geométrica é anholonomia de **transporte paralelo** no fibrado de estados sobre o espaço de parâmetros.","epistemico":"fato","exemplos":[],"id":"transporte_paralelo_e_bloch_sphere","memoria":"semantica","meta":{"dominio":"referencia","nivel":5},"origem":"fase_geometrica_berry.md","palavras_chave":[],"sinonimos":[],"tipo":"conceito","titulo":"Transporte paralelo e Bloch sphere"}],"id":"fase_geometrica_berry_transporte_paralelo_e_bloch_sphere","tipo":"conceito"}
\section{Transporte paralelo e Bloch sphere}
A fase geométrica é anholonomia de **transporte paralelo** no fibrado de estados sobre o espaço de parâmetros.
\prov{relações: parte\_de fase\_geometrica\_berry\_fase\_geometrica\_de\_berry\_pancharatnamberry; confirma a\_funcao\_de\_onda\_a\_esfera\_de\_bloch; confirma fibracao\_hopf\_bloch}
\prov{proveniência: fase\_geometrica\_berry.md \textperiodcentered{} web \textperiodcentered{} fato \textperiodcentered{} confiança 1.0 \textperiodcentered{} domínio referencia \textperiodcentered{} história 5 versões no jornal}
\prov{fusão de curadoria (texto idêntico): absorve transporte\_paralelo\_e\_bloch\_sphere --- o mesmo conteúdo por dois esquemas de endereço}

%CRISTAL {"arestas":[{"alvo":"gato","confianca":0.95,"epistemico":"fato","origem":"paper","rel":"parte_de"}],"confianca":1.0,"contraexemplos":[],"descricao":"1. Mudança de assinatura? **não** 2. Mudança de posto? **não** 3. Mudança de ramo Koch? **não** 4. Aproximação ao degenerado? **não** 5. Comportamento binário (liga/desliga)? **sim**","epistemico":"fato","exemplos":[],"id":"fecho_1572863","memoria":"semantica","meta":{"dominio":"manual","nivel":6,"verificacao":"slug idêntico — enriquecer, não duplicar"},"origem":"POW_METAL_SETOR.md","palavras_chave":[],"sinonimos":[],"tipo":"conceito","titulo":"Fecho `1572863`"}
\section{Fecho `1572863`}
1. Mudança de assinatura? **não** 2. Mudança de posto? **não** 3. Mudança de ramo Koch? **não** 4. Aproximação ao degenerado? **não** 5. Comportamento binário (liga/desliga)? **sim**
\prov{relações: parte\_de gato}
\prov{proveniência: POW\_METAL\_SETOR.md \textperiodcentered{} fato \textperiodcentered{} confiança 1.0 \textperiodcentered{} domínio manual \textperiodcentered{} história 3 versões no jornal}

%CRISTAL {"arestas":[{"alvo":"gato","confianca":0.95,"epistemico":"fato","origem":"paper","rel":"parte_de"}],"confianca":1.0,"contraexemplos":[],"descricao":"1. Mudança de assinatura? **não** 2. Mudança de posto? **não** 3. Mudança de ramo Koch? **não** 4. Aproximação ao degenerado? **não** 5. Comportamento binário (liga/desliga)? **sim**","epistemico":"fato","exemplos":[],"id":"fecho_1966079","memoria":"semantica","meta":{"dominio":"manual","nivel":6,"verificacao":"slug idêntico — enriquecer, não duplicar"},"origem":"POW_METAL_SETOR.md","palavras_chave":[],"sinonimos":[],"tipo":"conceito","titulo":"Fecho `1966079`"}
\section{Fecho `1966079`}
1. Mudança de assinatura? **não** 2. Mudança de posto? **não** 3. Mudança de ramo Koch? **não** 4. Aproximação ao degenerado? **não** 5. Comportamento binário (liga/desliga)? **sim**
\prov{relações: parte\_de gato}
\prov{proveniência: POW\_METAL\_SETOR.md \textperiodcentered{} fato \textperiodcentered{} confiança 1.0 \textperiodcentered{} domínio manual \textperiodcentered{} história 3 versões no jornal}

%CRISTAL {"arestas":[{"alvo":"habitantes_gentil_de_sitter","confianca":1.0,"epistemico":"fato","origem":"DS_UNIVERSO.md","rel":"parte_de"},{"alvo":"mecanica_quantica_hub","confianca":0.5,"epistemico":"fato","origem":"wiki","rel":"relaciona"},{"alvo":"geometria_hub_broca","confianca":0.5,"epistemico":"fato","origem":"wiki","rel":"relaciona"},{"alvo":"thread","confianca":0.5,"epistemico":"fato","origem":"wiki","rel":"relaciona"},{"alvo":"topologia","confianca":0.5,"epistemico":"fato","origem":"wiki","rel":"relaciona"},{"alvo":"resposta_a_pergunta_dificil","confianca":0.5,"epistemico":"fato","origem":"wiki","rel":"relaciona"},{"alvo":"projeto_cristaljson","confianca":0.5,"epistemico":"fato","origem":"wiki","rel":"relaciona"},{"alvo":"parabola_pa","confianca":0.5,"epistemico":"fato","origem":"wiki","rel":"relaciona"},{"alvo":"teste_ordem","confianca":0.5,"epistemico":"fato","origem":"wiki","rel":"relaciona"}],"confianca":1.0,"contraexemplos":[],"descricao":"Secção de DS_UNIVERSO.md.","epistemico":"fato","exemplos":[],"id":"fechos_centro_da_janela","memoria":"semantica","meta":{"dominio":"manual","nivel":6,"verificacao":"slug idêntico — enriquecer, não duplicar"},"origem":"DS_UNIVERSO.md","palavras_chave":[],"sinonimos":[],"tipo":"conceito","titulo":"Fechos (centro da janela)"}
\section{Fechos (centro da janela)}
Secção de DS\_UNIVERSO.md.
\prov{relações: parte\_de habitantes\_gentil\_de\_sitter; relaciona mecanica\_quantica\_hub; relaciona geometria\_hub\_broca; relaciona thread; relaciona topologia; relaciona resposta\_a\_pergunta\_dificil; relaciona projeto\_cristaljson; relaciona parabola\_pa; relaciona teste\_ordem}
\prov{proveniência: DS\_UNIVERSO.md \textperiodcentered{} fato \textperiodcentered{} confiança 1.0 \textperiodcentered{} domínio manual \textperiodcentered{} história 14 versões no jornal}

%CRISTAL {"arestas":[{"alvo":"mecanica_quantica","confianca":0.95,"epistemico":"fato","origem":"wiki","rel":"confirma"},{"alvo":"gato","confianca":0.95,"epistemico":"fato","origem":"wiki","rel":"confirma"},{"alvo":"a_funcao_de_onda_a_esfera_de_bloch","confianca":0.95,"epistemico":"fato","origem":"wiki","rel":"confirma"},{"alvo":"geometria","confianca":1.0,"epistemico":"fato","origem":"wiki","rel":"confirma"},{"alvo":"a_medida_o_calor_de_volta","confianca":1.0,"epistemico":"fato","origem":"wiki","rel":"confirma"},{"alvo":"conjunto_dougherty","confianca":1.0,"epistemico":"fato","origem":"wiki","rel":"analogia"}],"confianca":0.95,"contraexemplos":[],"descricao":"Referência verificada: fibracao_hopf_bloch.md","epistemico":"fato","exemplos":[],"id":"fibracao_hopf_bloch","memoria":"semantica","meta":{"dominio":"referencia","fonte":"web","nivel":5},"origem":"wiki","palavras_chave":[],"sinonimos":[],"tipo":"conceito","titulo":"fibracao hopf bloch"}
\section{fibracao hopf bloch}
Referência verificada: fibracao\_hopf\_bloch.md
\prov{relações: confirma mecanica\_quantica; confirma gato; confirma a\_funcao\_de\_onda\_a\_esfera\_de\_bloch; confirma geometria; confirma a\_medida\_o\_calor\_de\_volta; analogia conjunto\_dougherty}
\prov{proveniência: wiki \textperiodcentered{} web \textperiodcentered{} fato \textperiodcentered{} confiança 0.95 \textperiodcentered{} domínio referencia \textperiodcentered{} história 102 versões no jornal}

%CRISTAL {"arestas":[{"alvo":"fibracao_hopf_bloch_fibracao_de_hopf_e_esfera_de_bloch","confianca":1.0,"epistemico":"fato","origem":"fibracao_hopf_bloch.md","rel":"parte_de"},{"alvo":"word","confianca":0.5,"epistemico":"fato","origem":"wiki","rel":"relaciona"},{"alvo":"vontade_de_poder","confianca":0.5,"epistemico":"fato","origem":"wiki","rel":"relaciona"},{"alvo":"virtualizacao","confianca":0.5,"epistemico":"fato","origem":"wiki","rel":"relaciona"},{"alvo":"while","confianca":0.5,"epistemico":"fato","origem":"wiki","rel":"relaciona"}],"confianca":1.0,"contraexemplos":[],"descricao":"Confundir Hopf S³→S² com “toro em todo eixo” Dougherty — dimensões e grupos de homotopia diferentes.","epistemico":"fato","exemplos":[],"id":"fibracao_hopf_bloch_contraexemplos","memoria":"semantica","meta":{"dominio":"referencia","fonte":"web","nivel":5},"origem":"fibracao_hopf_bloch.md","palavras_chave":[],"sinonimos":[],"tipo":"conceito","titulo":"Contraexemplos"}
\section{Contraexemplos}
Confundir Hopf S\ensuremath{^{3}}\ensuremath{\to}S\ensuremath{^{2}} com “toro em todo eixo” Dougherty — dimensões e grupos de homotopia diferentes.
\prov{relações: parte\_de fibracao\_hopf\_bloch\_fibracao\_de\_hopf\_e\_esfera\_de\_bloch; relaciona word; relaciona vontade\_de\_poder; relaciona virtualizacao; relaciona while}
\prov{proveniência: fibracao\_hopf\_bloch.md \textperiodcentered{} web \textperiodcentered{} fato \textperiodcentered{} confiança 1.0 \textperiodcentered{} domínio referencia \textperiodcentered{} história 7 versões no jornal}

%CRISTAL {"fusao":[{"arestas":[{"alvo":"fibracao_hopf_bloch_fibracao_de_hopf_e_esfera_de_bloch","confianca":1.0,"epistemico":"fato","origem":"fibracao_hopf_bloch.md","rel":"parte_de"}],"confianca":1.0,"contraexemplos":[],"descricao":"Hilbert space de 2 qubits = S⁷. Fibrado de Hopf S⁷→S⁴ com fibra S³ — **sensível ao emaranhamento** (Mosseri-Dandoloff).","epistemico":"fato","exemplos":[],"id":"fibracao_hopf_bloch_dois_qubits_s7_s4","memoria":"semantica","meta":{"dominio":"referencia","fonte":"web","nivel":5},"origem":"fibracao_hopf_bloch.md","palavras_chave":[],"sinonimos":[],"tipo":"conceito","titulo":"Dois qubits — S⁷ → S⁴"},{"arestas":[{"alvo":"fibracao_de_hopf_e_esfera_de_bloch","confianca":1.0,"epistemico":"fato","origem":"fibracao_hopf_bloch.md","rel":"parte_de"},{"alvo":"word","confianca":0.5,"epistemico":"fato","origem":"wiki","rel":"relaciona"},{"alvo":"virtualizacao","confianca":0.5,"epistemico":"fato","origem":"wiki","rel":"relaciona"},{"alvo":"vontade_de_poder","confianca":0.5,"epistemico":"fato","origem":"wiki","rel":"relaciona"},{"alvo":"while","confianca":0.5,"epistemico":"fato","origem":"wiki","rel":"relaciona"},{"alvo":"esfera_de_bloch","confianca":0.5,"epistemico":"fato","origem":"wiki","rel":"relaciona"}],"confianca":1.0,"contraexemplos":[],"descricao":"Hilbert space de 2 qubits = S⁷. Fibrado de Hopf S⁷→S⁴ com fibra S³ — **sensível ao emaranhamento** (Mosseri-Dandoloff).","epistemico":"fato","exemplos":[],"id":"dois_qubits_s7_s4","memoria":"semantica","meta":{"dominio":"referencia","nivel":5},"origem":"fibracao_hopf_bloch.md","palavras_chave":[],"sinonimos":[],"tipo":"conceito","titulo":"Dois qubits — S⁷ → S⁴"}],"id":"fibracao_hopf_bloch_dois_qubits_s7_s4","tipo":"conceito"}
\section{Dois qubits — S\ensuremath{^{7}} \ensuremath{\to} S\ensuremath{^{4}}}
Hilbert space de 2 qubits = S\ensuremath{^{7}}. Fibrado de Hopf S\ensuremath{^{7}}\ensuremath{\to}S\ensuremath{^{4}} com fibra S\ensuremath{^{3}} — **sensível ao emaranhamento** (Mosseri-Dandoloff).
\prov{relações: parte\_de fibracao\_hopf\_bloch\_fibracao\_de\_hopf\_e\_esfera\_de\_bloch}
\prov{proveniência: fibracao\_hopf\_bloch.md \textperiodcentered{} web \textperiodcentered{} fato \textperiodcentered{} confiança 1.0 \textperiodcentered{} domínio referencia \textperiodcentered{} história 3 versões no jornal}
\prov{fusão de curadoria (texto idêntico): absorve dois\_qubits\_s7\_s4 --- o mesmo conteúdo por dois esquemas de endereço}

%CRISTAL {"fusao":[{"arestas":[{"alvo":"fibracao_hopf_bloch_fibracao_de_hopf_e_esfera_de_bloch","confianca":1.0,"epistemico":"fato","origem":"fibracao_hopf_bloch.md","rel":"parte_de"},{"alvo":"a_funcao_de_onda_a_esfera_de_bloch","confianca":0.95,"epistemico":"fato","origem":"wiki","rel":"confirma"}],"confianca":1.0,"contraexemplos":[],"descricao":"Estados puros de um qubit = rays em ℂ² = espaço projectivo CP¹ ≅ S² (Riemann sphere).","epistemico":"fato","exemplos":[],"id":"fibracao_hopf_bloch_esfera_de_bloch","memoria":"semantica","meta":{"dominio":"referencia","fonte":"web","nivel":5},"origem":"fibracao_hopf_bloch.md","palavras_chave":[],"sinonimos":[],"tipo":"conceito","titulo":"Esfera de Bloch"},{"arestas":[{"alvo":"fibracao_de_hopf_e_esfera_de_bloch","confianca":1.0,"epistemico":"fato","origem":"fibracao_hopf_bloch.md","rel":"parte_de"},{"alvo":"a_funcao_de_onda_a_esfera_de_bloch","confianca":1.0,"epistemico":"fato","origem":"wiki","rel":"confirma"},{"alvo":"word","confianca":0.5,"epistemico":"fato","origem":"wiki","rel":"relaciona"},{"alvo":"vontade_de_poder","confianca":0.5,"epistemico":"fato","origem":"wiki","rel":"relaciona"},{"alvo":"transcricao","confianca":0.5,"epistemico":"fato","origem":"wiki","rel":"relaciona"}],"confianca":1.0,"contraexemplos":[],"descricao":"Estados puros de um qubit = rays em ℂ² = espaço projectivo CP¹ ≅ S² (Riemann sphere).","epistemico":"fato","exemplos":[],"id":"esfera_de_bloch","memoria":"semantica","meta":{"dominio":"referencia","nivel":5},"origem":"fibracao_hopf_bloch.md","palavras_chave":[],"sinonimos":[],"tipo":"conceito","titulo":"Esfera de Bloch"}],"id":"fibracao_hopf_bloch_esfera_de_bloch","tipo":"conceito"}
\section{Esfera de Bloch}
Estados puros de um qubit = rays em \ensuremath{\mathbb{C}}\ensuremath{^{2}} = espaço projectivo CP\ensuremath{^{1}} \ensuremath{\cong} S\ensuremath{^{2}} (Riemann sphere).
\prov{relações: parte\_de fibracao\_hopf\_bloch\_fibracao\_de\_hopf\_e\_esfera\_de\_bloch; confirma a\_funcao\_de\_onda\_a\_esfera\_de\_bloch}
\prov{proveniência: fibracao\_hopf\_bloch.md \textperiodcentered{} web \textperiodcentered{} fato \textperiodcentered{} confiança 1.0 \textperiodcentered{} domínio referencia \textperiodcentered{} história 4 versões no jornal}
\prov{fusão de curadoria (texto idêntico): absorve esfera\_de\_bloch --- o mesmo conteúdo por dois esquemas de endereço}

%CRISTAL {"fusao":[{"arestas":[{"alvo":"fibracao_hopf_bloch_fibracao_de_hopf_e_esfera_de_bloch","confianca":1.0,"epistemico":"fato","origem":"fibracao_hopf_bloch.md","rel":"parte_de"},{"alvo":"observador_e_observado","confianca":0.95,"epistemico":"fato","origem":"wiki","rel":"confirma"}],"confianca":1.0,"contraexemplos":[],"descricao":"A fase global U(1) **não é observável** — só a base S² (|ψ|², operadores). Isto formaliza observador vs grau de liberdade eliminado.","epistemico":"fato","exemplos":[],"id":"fibracao_hopf_bloch_fase_global_e_observador","memoria":"semantica","meta":{"dominio":"referencia","fonte":"web","nivel":5},"origem":"fibracao_hopf_bloch.md","palavras_chave":[],"sinonimos":[],"tipo":"conceito","titulo":"Fase global e observador"},{"arestas":[{"alvo":"fibracao_de_hopf_e_esfera_de_bloch","confianca":1.0,"epistemico":"fato","origem":"fibracao_hopf_bloch.md","rel":"parte_de"},{"alvo":"observador_e_observado","confianca":0.95,"epistemico":"fato","origem":"wiki","rel":"confirma"}],"confianca":1.0,"contraexemplos":[],"descricao":"A fase global U(1) **não é observável** — só a base S² (|ψ|², operadores). Isto formaliza observador vs grau de liberdade eliminado.","epistemico":"fato","exemplos":[],"id":"fase_global_e_observador","memoria":"semantica","meta":{"dominio":"referencia","nivel":5},"origem":"fibracao_hopf_bloch.md","palavras_chave":[],"sinonimos":[],"tipo":"conceito","titulo":"Fase global e observador"}],"id":"fibracao_hopf_bloch_fase_global_e_observador","tipo":"conceito"}
\section{Fase global e observador}
A fase global U(1) **não é observável** — só a base S\ensuremath{^{2}} (|\ensuremath{\psi}|\ensuremath{^{2}}, operadores). Isto formaliza observador vs grau de liberdade eliminado.
\prov{relações: parte\_de fibracao\_hopf\_bloch\_fibracao\_de\_hopf\_e\_esfera\_de\_bloch; confirma observador\_e\_observado}
\prov{proveniência: fibracao\_hopf\_bloch.md \textperiodcentered{} web \textperiodcentered{} fato \textperiodcentered{} confiança 1.0 \textperiodcentered{} domínio referencia \textperiodcentered{} história 4 versões no jornal}
\prov{fusão de curadoria (texto idêntico): absorve fase\_global\_e\_observador --- o mesmo conteúdo por dois esquemas de endereço}

%CRISTAL {"fusao":[{"arestas":[{"alvo":"fibracao_hopf_bloch_fibracao_de_hopf_e_esfera_de_bloch","confianca":1.0,"epistemico":"fato","origem":"fibracao_hopf_bloch.md","rel":"parte_de"},{"alvo":"a_medida_o_calor_de_volta","confianca":0.95,"epistemico":"fato","origem":"wiki","rel":"confirma"}],"confianca":1.0,"contraexemplos":[],"descricao":"Mapa p: S³ → S². Cada ponto da base S² (estado observável) corresponde a uma fibra S¹ (fase global).","epistemico":"fato","exemplos":[],"id":"fibracao_hopf_bloch_fibracao_de_hopf","memoria":"semantica","meta":{"dominio":"referencia","fonte":"web","nivel":5},"origem":"fibracao_hopf_bloch.md","palavras_chave":[],"sinonimos":[],"tipo":"conceito","titulo":"Fibração de Hopf"},{"arestas":[{"alvo":"fibracao_de_hopf_e_esfera_de_bloch","confianca":1.0,"epistemico":"fato","origem":"fibracao_hopf_bloch.md","rel":"parte_de"},{"alvo":"a_medida_o_calor_de_volta","confianca":0.95,"epistemico":"fato","origem":"wiki","rel":"confirma"},{"alvo":"a_metrica_de_sitter_global","confianca":1.0,"epistemico":"fato","origem":"wiki","rel":"confirma"}],"confianca":1.0,"contraexemplos":[],"descricao":"Mapa p: S³ → S². Cada ponto da base S² (estado observável) corresponde a uma fibra S¹ (fase global).","epistemico":"fato","exemplos":[],"id":"fibracao_de_hopf","memoria":"semantica","meta":{"dominio":"referencia","nivel":5},"origem":"fibracao_hopf_bloch.md","palavras_chave":[],"sinonimos":[],"tipo":"conceito","titulo":"Fibração de Hopf"}],"id":"fibracao_hopf_bloch_fibracao_de_hopf","tipo":"conceito"}
\section{Fibração de Hopf}
Mapa p: S\ensuremath{^{3}} \ensuremath{\to} S\ensuremath{^{2}}. Cada ponto da base S\ensuremath{^{2}} (estado observável) corresponde a uma fibra S\ensuremath{^{1}} (fase global).
\prov{relações: parte\_de fibracao\_hopf\_bloch\_fibracao\_de\_hopf\_e\_esfera\_de\_bloch; confirma a\_medida\_o\_calor\_de\_volta}
\prov{proveniência: fibracao\_hopf\_bloch.md \textperiodcentered{} web \textperiodcentered{} fato \textperiodcentered{} confiança 1.0 \textperiodcentered{} domínio referencia \textperiodcentered{} história 4 versões no jornal}
\prov{fusão de curadoria (texto idêntico): absorve fibracao\_de\_hopf --- o mesmo conteúdo por dois esquemas de endereço}

%CRISTAL {"fusao":[{"arestas":[{"alvo":"fibracao_hopf_bloch","confianca":0.95,"epistemico":"fato","origem":"wiki","rel":"parte_de"}],"confianca":1.0,"contraexemplos":[],"descricao":"> **Epistemologia:** fato matemático-físico estabelecido (topologia + MQ). > **No grafo:** `epistemico=fato`, confiança 0.95. > **Papel:** confirma geometria S³→S² já presente nos papers Broca-SO; distingue de toros Dougherty.","epistemico":"fato","exemplos":[],"id":"fibracao_hopf_bloch_fibracao_de_hopf_e_esfera_de_bloch","memoria":"semantica","meta":{"dominio":"referencia","fonte":"web","nivel":5},"origem":"fibracao_hopf_bloch.md","palavras_chave":[],"sinonimos":[],"tipo":"conceito","titulo":"Fibração de Hopf e esfera de Bloch"},{"arestas":[{"alvo":"fibracao_hopf_bloch","confianca":0.95,"epistemico":"fato","origem":"wiki","rel":"parte_de"}],"confianca":1.0,"contraexemplos":[],"descricao":"> **Epistemologia:** fato matemático-físico estabelecido (topologia + MQ). > **No grafo:** `epistemico=fato`, confiança 0.95. > **Papel:** confirma geometria S³→S² já presente nos papers Broca-SO; distingue de toros Dougherty.","epistemico":"fato","exemplos":[],"id":"fibracao_de_hopf_e_esfera_de_bloch","memoria":"semantica","meta":{"dominio":"referencia","nivel":5},"origem":"fibracao_hopf_bloch.md","palavras_chave":[],"sinonimos":[],"tipo":"conceito","titulo":"Fibração de Hopf e esfera de Bloch"}],"id":"fibracao_hopf_bloch_fibracao_de_hopf_e_esfera_de_bloch","tipo":"conceito"}
\section{Fibração de Hopf e esfera de Bloch}
> **Epistemologia:** fato matemático-físico estabelecido (topologia + MQ). > **No grafo:** `epistemico=fato`, confiança 0.95. > **Papel:** confirma geometria S\ensuremath{^{3}}\ensuremath{\to}S\ensuremath{^{2}} já presente nos papers Broca-SO; distingue de toros Dougherty.
\prov{relações: parte\_de fibracao\_hopf\_bloch}
\prov{proveniência: fibracao\_hopf\_bloch.md \textperiodcentered{} web \textperiodcentered{} fato \textperiodcentered{} confiança 1.0 \textperiodcentered{} domínio referencia \textperiodcentered{} história 3 versões no jornal}
\prov{fusão de curadoria (texto idêntico): absorve fibracao\_de\_hopf\_e\_esfera\_de\_bloch --- o mesmo conteúdo por dois esquemas de endereço}

%CRISTAL {"arestas":[{"alvo":"fibracao_hopf_bloch_fibracao_de_hopf_e_esfera_de_bloch","confianca":1.0,"epistemico":"fato","origem":"fibracao_hopf_bloch.md","rel":"parte_de"}],"confianca":1.0,"contraexemplos":[],"descricao":"- https://en.wikipedia.org/wiki/Hopf_fibration - https://en.wikipedia.org/wiki/Bloch_sphere - https://iopscience.iop.org/article/10.1088/0305-4470/34/47/324","epistemico":"fato","exemplos":[],"id":"fibracao_hopf_bloch_fontes","memoria":"semantica","meta":{"dominio":"referencia","fonte":"web","nivel":5},"origem":"fibracao_hopf_bloch.md","palavras_chave":[],"sinonimos":[],"tipo":"conceito","titulo":"Fontes"}
\section{Fontes}
- https://en.wikipedia.org/wiki/Hopf\_fibration - https://en.wikipedia.org/wiki/Bloch\_sphere - https://iopscience.iop.org/article/10.1088/0305-4470/34/47/324
\prov{relações: parte\_de fibracao\_hopf\_bloch\_fibracao\_de\_hopf\_e\_esfera\_de\_bloch}
\prov{proveniência: fibracao\_hopf\_bloch.md \textperiodcentered{} web \textperiodcentered{} fato \textperiodcentered{} confiança 1.0 \textperiodcentered{} domínio referencia \textperiodcentered{} história 3 versões no jornal}

%CRISTAL {"fusao":[{"arestas":[{"alvo":"fibracao_hopf_bloch_fibracao_de_hopf_e_esfera_de_bloch","confianca":1.0,"epistemico":"fato","origem":"fibracao_hopf_bloch.md","rel":"parte_de"},{"alvo":"a_medida_o_calor_de_volta","confianca":0.95,"epistemico":"fato","origem":"wiki","rel":"confirma"},{"alvo":"a_dinamica_a_precessao_de_bloch","confianca":0.95,"epistemico":"fato","origem":"wiki","rel":"confirma"},{"alvo":"a_metrica_de_sitter_global","confianca":0.95,"epistemico":"fato","origem":"wiki","rel":"confirma"}],"confianca":1.0,"contraexemplos":[],"descricao":"| Conceito web | Nó Broca-SO | relação | |--------------|-------------|---------| | S³ versor | a_medida_o_calor_de_volta | confirma | | Bloch / CP¹ | a_funcao_de_onda_a_esfera_de_bloch | confirma | | precessão | a_dinamica_a_precessao_de_bloch | confirma | | qubit / gato A | gato, mecanica_quantica | confirma | | S³ métrica | a_metrica_de_sitter_global | confirma |","epistemico":"fato","exemplos":[],"id":"fibracao_hopf_bloch_ponte_broca-so_confirmacao","memoria":"semantica","meta":{"dominio":"referencia","fonte":"web","nivel":5},"origem":"fibracao_hopf_bloch.md","palavras_chave":[],"sinonimos":[],"tipo":"conceito","titulo":"Ponte Broca-SO (confirmação)"},{"arestas":[{"alvo":"fibracao_de_hopf_e_esfera_de_bloch","confianca":1.0,"epistemico":"fato","origem":"fibracao_hopf_bloch.md","rel":"parte_de"},{"alvo":"a_medida_o_calor_de_volta","confianca":0.95,"epistemico":"fato","origem":"wiki","rel":"confirma"},{"alvo":"a_dinamica_a_precessao_de_bloch","confianca":0.95,"epistemico":"fato","origem":"wiki","rel":"confirma"},{"alvo":"a_metrica_de_sitter_global","confianca":0.95,"epistemico":"fato","origem":"wiki","rel":"confirma"}],"confianca":1.0,"contraexemplos":[],"descricao":"| Conceito web | Nó Broca-SO | relação | |--------------|-------------|---------| | S³ versor | a_medida_o_calor_de_volta | confirma | | Bloch / CP¹ | a_funcao_de_onda_a_esfera_de_bloch | confirma | | precessão | a_dinamica_a_precessao_de_bloch | confirma | | qubit / gato A | gato, mecanica_quantica | confirma | | S³ métrica | a_metrica_de_sitter_global | confirma |","epistemico":"fato","exemplos":[],"id":"ponte_broca-so_confirmacao","memoria":"semantica","meta":{"dominio":"referencia","nivel":5},"origem":"fibracao_hopf_bloch.md","palavras_chave":[],"sinonimos":[],"tipo":"conceito","titulo":"Ponte Broca-SO (confirmação)"}],"id":"fibracao_hopf_bloch_ponte_broca-so_confirmacao","tipo":"conceito"}
\section{Ponte Broca-SO (confirmação)}
| Conceito web | Nó Broca-SO | relação | |--------------|-------------|---------| | S\ensuremath{^{3}} versor | a\_medida\_o\_calor\_de\_volta | confirma | | Bloch / CP\ensuremath{^{1}} | a\_funcao\_de\_onda\_a\_esfera\_de\_bloch | confirma | | precessão | a\_dinamica\_a\_precessao\_de\_bloch | confirma | | qubit / gato A | gato, mecanica\_quantica | confirma | | S\ensuremath{^{3}} métrica | a\_metrica\_de\_sitter\_global | confirma |
\prov{relações: parte\_de fibracao\_hopf\_bloch\_fibracao\_de\_hopf\_e\_esfera\_de\_bloch; confirma a\_medida\_o\_calor\_de\_volta; confirma a\_dinamica\_a\_precessao\_de\_bloch; confirma a\_metrica\_de\_sitter\_global}
\prov{proveniência: fibracao\_hopf\_bloch.md \textperiodcentered{} web \textperiodcentered{} fato \textperiodcentered{} confiança 1.0 \textperiodcentered{} domínio referencia \textperiodcentered{} história 6 versões no jornal}
\prov{fusão de curadoria (texto idêntico): absorve ponte\_broca-so\_confirmacao --- o mesmo conteúdo por dois esquemas de endereço}

%CRISTAL {"arestas":[{"alvo":"fibracao_hopf_bloch_fibracao_de_hopf_e_esfera_de_bloch","confianca":1.0,"epistemico":"fato","origem":"fibracao_hopf_bloch.md","rel":"parte_de"}],"confianca":1.0,"contraexemplos":[],"descricao":"- **Hopf (1931):** fibration S³ → S², fibra S¹ - **Bloch (1946):** representação esférica do spin / qubit - **Mosseri & Dandoloff (2001):** *Geometry of entangled states, Bloch spheres and Hopf fibrations*, J. Phys. A 34 10243 — [arXiv:quant-ph/0108137](https://arxiv.org/abs/quant-ph/0108137) - **Wikipedia:** [Hopf fibration](https://en.wikipedia.org/wiki/Hopf_fibration), [Bloch sphere](https://en.wikipedia.org/wiki/Bloch_sphere)","epistemico":"fato","exemplos":[],"id":"fibracao_hopf_bloch_referencia","memoria":"semantica","meta":{"dominio":"referencia","fonte":"web","nivel":5},"origem":"fibracao_hopf_bloch.md","palavras_chave":[],"sinonimos":[],"tipo":"conceito","titulo":"Referência"}
\section{Referência}
- **Hopf (1931):** fibration S\ensuremath{^{3}} \ensuremath{\to} S\ensuremath{^{2}}, fibra S\ensuremath{^{1}} - **Bloch (1946):** representação esférica do spin / qubit - **Mosseri \& Dandoloff (2001):** *Geometry of entangled states, Bloch spheres and Hopf fibrations*, J. Phys. A 34 10243 — [arXiv:quant-ph/0108137](https://arxiv.org/abs/quant-ph/0108137) - **Wikipedia:** [Hopf fibration](https://en.wikipedia.org/wiki/Hopf\_fibration), [Bloch sphere](https://en.wikipedia.org/wiki/Bloch\_sphere)
\prov{relações: parte\_de fibracao\_hopf\_bloch\_fibracao\_de\_hopf\_e\_esfera\_de\_bloch}
\prov{proveniência: fibracao\_hopf\_bloch.md \textperiodcentered{} web \textperiodcentered{} fato \textperiodcentered{} confiança 1.0 \textperiodcentered{} domínio referencia \textperiodcentered{} história 3 versões no jornal}

%CRISTAL {"arestas":[{"alvo":"fibracao_hopf_bloch_fibracao_de_hopf_e_esfera_de_bloch","confianca":1.0,"epistemico":"fato","origem":"fibracao_hopf_bloch.md","rel":"parte_de"},{"alvo":"conjunto_dougherty","confianca":0.95,"epistemico":"fato","origem":"wiki","rel":"analogia"},{"alvo":"ponte_quantica_dougherty","confianca":0.95,"epistemico":"fato","origem":"wiki","rel":"referencia"}],"confianca":1.0,"contraexemplos":[],"descricao":"Dougherty propõe **toros** e escalas φ; Hopf propõe **S³→S²** com fibra S¹ — ambos geométricos, **estruturas distintas**.","epistemico":"fato","exemplos":[],"id":"fibracao_hopf_bloch_relacao_com_hipotese_dougherty","memoria":"semantica","meta":{"dominio":"referencia","fonte":"web","nivel":5},"origem":"fibracao_hopf_bloch.md","palavras_chave":[],"sinonimos":[],"tipo":"conceito","titulo":"Relação com hipótese Dougherty"}
\section{Relação com hipótese Dougherty}
Dougherty propõe **toros** e escalas \ensuremath{\varphi}; Hopf propõe **S\ensuremath{^{3}}\ensuremath{\to}S\ensuremath{^{2}}** com fibra S\ensuremath{^{1}} — ambos geométricos, **estruturas distintas**.
\prov{relações: parte\_de fibracao\_hopf\_bloch\_fibracao\_de\_hopf\_e\_esfera\_de\_bloch; analogia conjunto\_dougherty; referencia ponte\_quantica\_dougherty}
\prov{proveniência: fibracao\_hopf\_bloch.md \textperiodcentered{} web \textperiodcentered{} fato \textperiodcentered{} confiança 1.0 \textperiodcentered{} domínio referencia \textperiodcentered{} história 5 versões no jornal}

%CRISTAL {"arestas":[{"alvo":"fase_2_tecido_de_conhecimento","confianca":1.0,"epistemico":"fato","origem":"CONHECIMENTO.md","rel":"parte_de"},{"alvo":"transcricao","confianca":0.5,"epistemico":"fato","origem":"wiki","rel":"relaciona"},{"alvo":"utilitarismo","confianca":0.5,"epistemico":"fato","origem":"wiki","rel":"relaciona"},{"alvo":"virtualizacao","confianca":0.5,"epistemico":"fato","origem":"wiki","rel":"relaciona"},{"alvo":"scheduler","confianca":0.5,"epistemico":"fato","origem":"wiki","rel":"relaciona"},{"alvo":"resultado_anterior_identidade_nao_cobertura","confianca":0.5,"epistemico":"fato","origem":"wiki","rel":"relaciona"},{"alvo":"vida-morte","confianca":0.5,"epistemico":"fato","origem":"wiki","rel":"relaciona"},{"alvo":"ruido","confianca":0.5,"epistemico":"fato","origem":"wiki","rel":"relaciona"},{"alvo":"modo_espectral_spect","confianca":0.5,"epistemico":"fato","origem":"wiki","rel":"relaciona"},{"alvo":"pow_metal_ressonancia_val","confianca":0.5,"epistemico":"fato","origem":"wiki","rel":"relaciona"},{"alvo":"teoria_do_valor","confianca":0.5,"epistemico":"fato","origem":"wiki","rel":"relaciona"},{"alvo":"tipos_de_interpretante","confianca":0.5,"epistemico":"fato","origem":"wiki","rel":"relaciona"},{"alvo":"resultado_completo","confianca":0.5,"epistemico":"fato","origem":"wiki","rel":"relaciona"},{"alvo":"pedagogia","confianca":0.5,"epistemico":"fato","origem":"wiki","rel":"relaciona"},{"alvo":"veredito_experimental","confianca":0.5,"epistemico":"fato","origem":"wiki","rel":"relaciona"},{"alvo":"filesystem","confianca":0.5,"epistemico":"fato","origem":"wiki","rel":"relaciona"},{"alvo":"setor_algebrico_descontinuidades_no_fecho","confianca":0.5,"epistemico":"fato","origem":"wiki","rel":"relaciona"},{"alvo":"pitagorismo","confianca":0.5,"epistemico":"fato","origem":"wiki","rel":"relaciona"},{"alvo":"orbita_pell_hardware","confianca":0.5,"epistemico":"fato","origem":"wiki","rel":"relaciona"},{"alvo":"mecanica","confianca":0.5,"epistemico":"fato","origem":"wiki","rel":"relaciona"},{"alvo":"transicao","confianca":0.5,"epistemico":"fato","origem":"wiki","rel":"relaciona"},{"alvo":"pell","confianca":0.5,"epistemico":"fato","origem":"wiki","rel":"relaciona"},{"alvo":"gentil_fundamentos_22_relacoes_especiais","confianca":0.5,"epistemico":"fato","origem":"wiki","rel":"relaciona"},{"alvo":"paper_2_reta_ouro","confianca":0.5,"epistemico":"fato","origem":"wiki","rel":"relaciona"}],"confianca":1.0,"contraexemplos":[],"descricao":"**Problema:** após `ADD`, `R` não persiste em `B` no próximo `LOAD`.","epistemico":"fato","exemplos":[],"id":"ficheiros","memoria":"semantica","meta":{"dominio":"manual","duplicata_sugerida":[],"nivel":6,"verificacao":"parte de conceito mais amplo 'ficheiros'"},"origem":"PROVA_ISA.md","palavras_chave":[],"sinonimos":[],"tipo":"conceito","titulo":"Ficheiros"}
\section{Ficheiros}
**Problema:** após `ADD`, `R` não persiste em `B` no próximo `LOAD`.
\prov{relações: parte\_de fase\_2\_tecido\_de\_conhecimento; relaciona transcricao; relaciona utilitarismo; relaciona virtualizacao; relaciona scheduler; relaciona resultado\_anterior\_identidade\_nao\_cobertura; relaciona vida-morte; relaciona ruido; relaciona modo\_espectral\_spect; relaciona pow\_metal\_ressonancia\_val; relaciona teoria\_do\_valor; relaciona tipos\_de\_interpretante; relaciona resultado\_completo; relaciona pedagogia; relaciona veredito\_experimental; relaciona filesystem; relaciona setor\_algebrico\_descontinuidades\_no\_fecho; relaciona pitagorismo; relaciona orbita\_pell\_hardware; relaciona mecanica; relaciona transicao; relaciona pell; relaciona gentil\_fundamentos\_22\_relacoes\_especiais; relaciona paper\_2\_reta\_ouro}
\prov{proveniência: PROVA\_ISA.md \textperiodcentered{} fato \textperiodcentered{} confiança 1.0 \textperiodcentered{} domínio manual \textperiodcentered{} história 27 versões no jornal}

%CRISTAL {"arestas":[{"alvo":"2_semantica_formal","confianca":1.0,"epistemico":"fato","origem":"PROVA_ISA.md","rel":"parte_de"},{"alvo":"word","confianca":0.5,"epistemico":"fato","origem":"wiki","rel":"relaciona"},{"alvo":"virtualizacao","confianca":0.5,"epistemico":"fato","origem":"wiki","rel":"relaciona"},{"alvo":"semiotica_peirce_triade","confianca":0.5,"epistemico":"fato","origem":"wiki","rel":"relaciona"},{"alvo":"proposition_gerador_de_escala_ipropderiv_com_tddt","confianca":0.5,"epistemico":"fato","origem":"wiki","rel":"relaciona"},{"alvo":"goldchain_mercado_hub_fontes","confianca":0.5,"epistemico":"fato","origem":"wiki","rel":"relaciona"},{"alvo":"gato_operador_hub_contraexemplos","confianca":0.5,"epistemico":"fato","origem":"wiki","rel":"relaciona"},{"alvo":"vida-morte","confianca":0.5,"epistemico":"fato","origem":"wiki","rel":"relaciona"},{"alvo":"panpsiquismo","confianca":0.5,"epistemico":"fato","origem":"wiki","rel":"relaciona"}],"confianca":1.0,"contraexemplos":[],"descricao":"| Bit | Condição | |-----|----------| | `ZERO` | `A=0 ∧ B=0` | | `EQ` | `A=B` | | `LT` | `A.total < B.total` |","epistemico":"fato","exemplos":[],"id":"flags","memoria":"semantica","meta":{"dominio":"manual","nivel":6,"verificacao":"slug idêntico — enriquecer, não duplicar"},"origem":"PROVA_ISA.md","palavras_chave":[],"sinonimos":[],"tipo":"conceito","titulo":"FLAGS"}
\section{FLAGS}
| Bit | Condição | |-----|----------| | `ZERO` | `A=0 \ensuremath{\wedge} B=0` | | `EQ` | `A=B` | | `LT` | `A.total < B.total` |
\prov{relações: parte\_de 2\_semantica\_formal; relaciona word; relaciona virtualizacao; relaciona semiotica\_peirce\_triade; relaciona proposition\_gerador\_de\_escala\_ipropderiv\_com\_tddt; relaciona goldchain\_mercado\_hub\_fontes; relaciona gato\_operador\_hub\_contraexemplos; relaciona vida-morte; relaciona panpsiquismo}
\prov{proveniência: PROVA\_ISA.md \textperiodcentered{} fato \textperiodcentered{} confiança 1.0 \textperiodcentered{} domínio manual \textperiodcentered{} história 12 versões no jornal}

%CRISTAL {"arestas":[{"alvo":"o_substrato_a_maquina_e_funcao_de_onda","confianca":1.0,"epistemico":"fato","origem":"wiki","rel":"explica"},{"alvo":"corpo_dual","confianca":1.0,"epistemico":"fato","origem":"wiki","rel":"usa"},{"alvo":"conjugados_galois_parabola","confianca":1.0,"epistemico":"fato","origem":"wiki","rel":"relaciona"}],"confianca":1.0,"contraexemplos":[],"descricao":"O flip de Wick ε não é postulado: é ε=sign(−disc) da recorrência do gap. disc<0→ε=+1 (ℂ, divisão); disc>0→ε=−1 (split, divisores de zero); disc=0→ε=0 (duais, nilpotente). O gap de ouro tem disc=5>0, logo nasce split. A origem algébrica do flip usado em toda a máquina.","epistemico":"fato","exemplos":[],"id":"flip_epsilon_discriminante","memoria":"semantica","meta":{"autor":"Aarão/broca-so","dominio":"broca_repos","fonte":""},"origem":"wiki","palavras_chave":[],"sinonimos":[],"tipo":"conceito","titulo":"O Flip é o Sinal do Discriminante"}
\section{O Flip é o Sinal do Discriminante}
O flip de Wick \ensuremath{\varepsilon} não é postulado: é \ensuremath{\varepsilon}=sign(\ensuremath{-}disc) da recorrência do gap. disc<0\ensuremath{\to}\ensuremath{\varepsilon}=+1 (\ensuremath{\mathbb{C}}, divisão); disc>0\ensuremath{\to}\ensuremath{\varepsilon}=\ensuremath{-}1 (split, divisores de zero); disc=0\ensuremath{\to}\ensuremath{\varepsilon}=0 (duais, nilpotente). O gap de ouro tem disc=5>0, logo nasce split. A origem algébrica do flip usado em toda a máquina.
\prov{relações: explica o\_substrato\_a\_maquina\_e\_funcao\_de\_onda; usa corpo\_dual; relaciona conjugados\_galois\_parabola}
\prov{proveniência: wiki \textperiodcentered{} fato \textperiodcentered{} confiança 1.0 \textperiodcentered{} domínio broca\_repos \textperiodcentered{} história 4 versões no jornal}

%CRISTAL {"arestas":[{"alvo":"gato","confianca":1.0,"epistemico":"fato","origem":"paper","rel":"referencia"},{"alvo":"computacao_fractal","confianca":1.0,"epistemico":"fato","origem":"wiki","rel":"parte_de"},{"alvo":"koch_pell","confianca":1.0,"epistemico":"fato","origem":"wiki","rel":"relaciona"},{"alvo":"koch_parabola","confianca":1.0,"epistemico":"fato","origem":"wiki","rel":"relaciona"},{"alvo":"dissipacao_na_passagem","confianca":1.0,"epistemico":"fato","origem":"wiki","rel":"relaciona"}],"confianca":1.0,"contraexemplos":[],"descricao":"O pipeline canónico w ──gold──► ouro (Fibonacci) ──silver──► prata (Pell) ──parábola──► s, Cantor/Fib (as RETAS ASSOCIATIVAS) ──z²──► floco Koch (o RECIPIENTE não-associativo, mod p=1009 — NÃO é Pell). A leitura: cada estágio à esquerda do z² é reta associativa (codifica, reconhece por overlap); o floco à direita é o resto ortogonal — o que a reta perde na passagem (ver dissipacao_na_passagem, floco_phi). Área finita / perímetro infinito é só a metáfora que nomeia; o floco real é Φ(w).","epistemico":"fato","exemplos":[],"id":"floco_de_neve","memoria":"semantica","meta":{"dominio":"manual","duplicata_sugerida":[],"nivel":6,"verificacao":"parte de conceito mais amplo 'floco_de_neve'"},"origem":"KOCH_PELL.md","palavras_chave":[],"sinonimos":[],"tipo":"conceito","titulo":"floco_de_neve"}
\section{floco\_de\_neve}
O pipeline canónico w --gold--$\triangleright$ ouro (Fibonacci) --silver--$\triangleright$ prata (Pell) --parábola--$\triangleright$ s, Cantor/Fib (as RETAS ASSOCIATIVAS) --z\ensuremath{^{2}}--$\triangleright$ floco Koch (o RECIPIENTE não-associativo, mod p=1009 — NÃO é Pell). A leitura: cada estágio à esquerda do z\ensuremath{^{2}} é reta associativa (codifica, reconhece por overlap); o floco à direita é o resto ortogonal — o que a reta perde na passagem (ver dissipacao\_na\_passagem, floco\_phi). Área finita / perímetro infinito é só a metáfora que nomeia; o floco real é \ensuremath{\Phi}(w).
\prov{relações: referencia gato; parte\_de computacao\_fractal; relaciona koch\_pell; relaciona koch\_parabola; relaciona dissipacao\_na\_passagem}
\prov{proveniência: KOCH\_PELL.md \textperiodcentered{} fato \textperiodcentered{} confiança 1.0 \textperiodcentered{} domínio manual \textperiodcentered{} história 38 versões no jornal}

%CRISTAL {"arestas":[{"alvo":"rede_neural_algebra_aprendivel","confianca":1.0,"epistemico":"fato","origem":"wiki","rel":"usa"},{"alvo":"hopfield","confianca":1.0,"epistemico":"fato","origem":"wiki","rel":"relaciona"}],"confianca":1.0,"contraexemplos":[],"descricao":"Com parâmetros por neurônio, o conjunto (p,r,t,s,u)ᵢ forma um fluido no espaço de álgebras governado por Langevin m·s̈=2m·s−k(s−salvo)+ruído — três osciladores invertidos desacoplados de frequência √2. Primo dinâmico da relaxação de Hopfield.","epistemico":"fato","exemplos":[],"id":"fluido_algebrico","memoria":"semantica","meta":{"autor":"Lopes/Aarão","dominio":"hiper","fonte":""},"origem":"wiki","palavras_chave":[],"sinonimos":[],"tipo":"conceito","titulo":"Fluido Algébrico"}
\section{Fluido Algébrico}
Com parâmetros por neurônio, o conjunto (p,r,t,s,u)\ensuremath{_{i}} forma um fluido no espaço de álgebras governado por Langevin m\textperiodcentered s\ensuremath{}=2m\textperiodcentered s\ensuremath{-}k(s\ensuremath{-}salvo)+ruído — três osciladores invertidos desacoplados de frequência \ensuremath{\sqrt{\,}}2. Primo dinâmico da relaxação de Hopfield.
\prov{relações: usa rede\_neural\_algebra\_aprendivel; relaciona hopfield}
\prov{proveniência: wiki \textperiodcentered{} fato \textperiodcentered{} confiança 1.0 \textperiodcentered{} domínio hiper \textperiodcentered{} história 5 versões no jornal}

%CRISTAL {"arestas":[{"alvo":"fibracao_de_hopf","confianca":1.0,"epistemico":"fato","origem":"wiki","rel":"usa"},{"alvo":"razao_aurea","confianca":1.0,"epistemico":"fato","origem":"wiki","rel":"usa"},{"alvo":"emaranhador","confianca":1.0,"epistemico":"fato","origem":"wiki","rel":"relaciona"}],"confianca":1.0,"contraexemplos":[],"descricao":"Teorema: cada campo local ℂ_â contém um atrator Cantor-φ² (IFS com contração 1/φ², dim Hausdorff log2/logφ²≈0,72); o fibrado é uma família SO(3)-equivariante desses fractais sobre S². A razão áurea emerge como autovalor universal de contração (R²=R+1). A distribuição-limite é mista (contínua+picos).","epistemico":"fato","exemplos":[],"id":"fractal_phi_fibrado","memoria":"semantica","meta":{"autor":"Lopes/Aarão","dominio":"hiper","fonte":""},"origem":"wiki","palavras_chave":[],"sinonimos":[],"tipo":"conceito","titulo":"Estrutura Fractal-φ² no Fibrado"}
\section{Estrutura Fractal-\ensuremath{\varphi}\ensuremath{^{2}} no Fibrado}
Teorema: cada campo local \ensuremath{\mathbb{C}}\_â contém um atrator Cantor-\ensuremath{\varphi}\ensuremath{^{2}} (IFS com contração 1/\ensuremath{\varphi}\ensuremath{^{2}}, dim Hausdorff log2/log\ensuremath{\varphi}\ensuremath{^{2}}\ensuremath{\approx}0,72); o fibrado é uma família SO(3)-equivariante desses fractais sobre S\ensuremath{^{2}}. A razão áurea emerge como autovalor universal de contração (R\ensuremath{^{2}}=R+1). A distribuição-limite é mista (contínua+picos).
\prov{relações: usa fibracao\_de\_hopf; usa razao\_aurea; relaciona emaranhador}
\prov{proveniência: wiki \textperiodcentered{} fato \textperiodcentered{} confiança 1.0 \textperiodcentered{} domínio hiper \textperiodcentered{} história 5 versões no jornal}

%CRISTAL {"arestas":[{"alvo":"fibracao_de_hopf","confianca":1.0,"epistemico":"fato","origem":"wiki","rel":"usa"},{"alvo":"mecanica_quantica","confianca":1.0,"epistemico":"fato","origem":"wiki","rel":"explica"},{"alvo":"o_substrato_a_maquina_e_funcao_de_onda","confianca":1.0,"epistemico":"fato","origem":"wiki","rel":"relaciona"}],"confianca":1.0,"contraexemplos":[],"descricao":"Teorema: o espaço de estados do qubit ℂ² é canonicamente isomorfo às seções tautológicas do fibrado em retas E⁺→S² dos autoespaços de R_â:v↦vâ em ℍ; E⁺≅O(−1) (Chern −1). As amplitudes são valores da seção; |⟨+â|ψ⟩|²=‖P⁺_â(ψ)‖² é identidade algébrica em ℍ, não postulado.","epistemico":"fato","exemplos":[],"id":"funcao_onda_secao_hopf","memoria":"semantica","meta":{"autor":"Lopes/Aarão","dominio":"hiper","fonte":""},"origem":"wiki","palavras_chave":[],"sinonimos":[],"tipo":"conceito","titulo":"A Função de Onda é Seção do Fibrado de Hopf O(−1)"}
\section{A Função de Onda é Seção do Fibrado de Hopf O(\ensuremath{-}1)}
Teorema: o espaço de estados do qubit \ensuremath{\mathbb{C}}\ensuremath{^{2}} é canonicamente isomorfo às seções tautológicas do fibrado em retas E\ensuremath{^{+}}\ensuremath{\to}S\ensuremath{^{2}} dos autoespaços de R\_â:v\ensuremath{\mapsto}vâ em \ensuremath{\mathbb{H}}; E\ensuremath{^{+}}\ensuremath{\cong}O(\ensuremath{-}1) (Chern \ensuremath{-}1). As amplitudes são valores da seção; |\ensuremath{\langle}+â|\ensuremath{\psi}\ensuremath{\rangle}|\ensuremath{^{2}}=\ensuremath{\|}P\ensuremath{^{+}}\_â(\ensuremath{\psi})\ensuremath{\|}\ensuremath{^{2}} é identidade algébrica em \ensuremath{\mathbb{H}}, não postulado.
\prov{relações: usa fibracao\_de\_hopf; explica mecanica\_quantica; relaciona o\_substrato\_a\_maquina\_e\_funcao\_de\_onda}
\prov{proveniência: wiki \textperiodcentered{} fato \textperiodcentered{} confiança 1.0 \textperiodcentered{} domínio hiper \textperiodcentered{} história 5 versões no jornal}

%CRISTAL {"arestas":[{"alvo":"gato","confianca":0.95,"epistemico":"fato","origem":"paper","rel":"parte_de"},{"alvo":"memoria_virtual","confianca":0.5,"epistemico":"fato","origem":"wiki","rel":"relaciona"},{"alvo":"mecanica_quantica_hub","confianca":0.5,"epistemico":"fato","origem":"wiki","rel":"relaciona"},{"alvo":"o_metodo_em_escala_conhecida_o_sistema_solar","confianca":0.5,"epistemico":"fato","origem":"wiki","rel":"relaciona"},{"alvo":"geometria_hub_broca","confianca":0.5,"epistemico":"fato","origem":"wiki","rel":"relaciona"},{"alvo":"utilitarismo","confianca":0.5,"epistemico":"fato","origem":"wiki","rel":"relaciona"}],"confianca":1.0,"contraexemplos":[],"descricao":"```text Header............    2,000,000 (100%)","epistemico":"fato","exemplos":[],"id":"funil_cumulativo","memoria":"semantica","meta":{"dominio":"manual","nivel":6,"verificacao":"slug idêntico — enriquecer, não duplicar"},"origem":"POW_METAL_FUNIL.md","palavras_chave":[],"sinonimos":[],"tipo":"conceito","titulo":"Funil cumulativo"}
\section{Funil cumulativo}
```text Header............    2,000,000 (100\%)
\prov{relações: parte\_de gato; relaciona memoria\_virtual; relaciona mecanica\_quantica\_hub; relaciona o\_metodo\_em\_escala\_conhecida\_o\_sistema\_solar; relaciona geometria\_hub\_broca; relaciona utilitarismo}
\prov{proveniência: POW\_METAL\_FUNIL.md \textperiodcentered{} fato \textperiodcentered{} confiança 1.0 \textperiodcentered{} domínio manual \textperiodcentered{} história 8 versões no jornal}

%CRISTAL {"arestas":[{"alvo":"gato","confianca":0.95,"epistemico":"fato","origem":"wiki","rel":"confirma"},{"alvo":"hopfield_rede_hub","confianca":0.5,"epistemico":"fato","origem":"wiki","rel":"relaciona"},{"alvo":"goldcoin","confianca":0.5,"epistemico":"fato","origem":"wiki","rel":"relaciona"},{"alvo":"rust","confianca":0.5,"epistemico":"fato","origem":"wiki","rel":"relaciona"},{"alvo":"gpu","confianca":0.5,"epistemico":"fato","origem":"wiki","rel":"relaciona"},{"alvo":"vida-morte","confianca":0.5,"epistemico":"fato","origem":"wiki","rel":"relaciona"},{"alvo":"koch_atlas","confianca":0.5,"epistemico":"fato","origem":"wiki","rel":"relaciona"},{"alvo":"main","confianca":0.5,"epistemico":"fato","origem":"wiki","rel":"relaciona"},{"alvo":"motivos","confianca":0.5,"epistemico":"fato","origem":"wiki","rel":"relaciona"},{"alvo":"koch_memoria","confianca":0.5,"epistemico":"fato","origem":"wiki","rel":"relaciona"},{"alvo":"por_sessao_conversadadostelemetriasessaojsonl","confianca":0.5,"epistemico":"fato","origem":"wiki","rel":"relaciona"},{"alvo":"pow_metal_pre_fecho","confianca":0.5,"epistemico":"fato","origem":"wiki","rel":"relaciona"},{"alvo":"geometria_hub_broca","confianca":0.5,"epistemico":"fato","origem":"wiki","rel":"relaciona"},{"alvo":"java","confianca":0.5,"epistemico":"fato","origem":"wiki","rel":"relaciona"},{"alvo":"recursao","confianca":0.5,"epistemico":"fato","origem":"wiki","rel":"relaciona"}],"confianca":0.95,"contraexemplos":[],"descricao":"Referência verificada: gato_operador_hub.md","epistemico":"fato","exemplos":[],"id":"gato_operador_hub","memoria":"semantica","meta":{"dominio":"referencia","fonte":"web","nivel":5},"origem":"wiki","palavras_chave":[],"sinonimos":[],"tipo":"conceito","titulo":"gato operador hub"}
\section{gato operador hub}
Referência verificada: gato\_operador\_hub.md
\prov{relações: confirma gato; relaciona hopfield\_rede\_hub; relaciona goldcoin; relaciona rust; relaciona gpu; relaciona vida-morte; relaciona koch\_atlas; relaciona main; relaciona motivos; relaciona koch\_memoria; relaciona por\_sessao\_conversadadostelemetriasessaojsonl; relaciona pow\_metal\_pre\_fecho; relaciona geometria\_hub\_broca; relaciona java; relaciona recursao}
\prov{proveniência: wiki \textperiodcentered{} web \textperiodcentered{} fato \textperiodcentered{} confiança 0.95 \textperiodcentered{} domínio referencia \textperiodcentered{} história 17 versões no jornal}

%CRISTAL {"arestas":[{"alvo":"gato_operador_hub_gato_a_hub_do_operador_algebrico_pow","confianca":1.0,"epistemico":"fato","origem":"gato_operador_hub.md","rel":"parte_de"},{"alvo":"virtualizacao","confianca":0.5,"epistemico":"fato","origem":"wiki","rel":"relaciona"},{"alvo":"word","confianca":0.5,"epistemico":"fato","origem":"wiki","rel":"relaciona"},{"alvo":"universidade_curriculo_em_camadas","confianca":0.5,"epistemico":"fato","origem":"wiki","rel":"relaciona"},{"alvo":"semiotica_peirce_triade","confianca":0.5,"epistemico":"fato","origem":"wiki","rel":"relaciona"},{"alvo":"goldchain_mercado_hub_fontes","confianca":0.5,"epistemico":"fato","origem":"wiki","rel":"relaciona"}],"confianca":1.0,"contraexemplos":[],"descricao":"**gato A ≠ gato felino** — homónimo; semilla mamífero ≠ operador PoW.","epistemico":"fato","exemplos":[],"id":"gato_operador_hub_contraexemplos","memoria":"semantica","meta":{"dominio":"referencia","fonte":"web","nivel":5},"origem":"gato_operador_hub.md","palavras_chave":[],"sinonimos":[],"tipo":"conceito","titulo":"Contraexemplos"}
\section{Contraexemplos}
**gato A \ensuremath{\ne} gato felino** — homónimo; semilla mamífero \ensuremath{\ne} operador PoW.
\prov{relações: parte\_de gato\_operador\_hub\_gato\_a\_hub\_do\_operador\_algebrico\_pow; relaciona virtualizacao; relaciona word; relaciona universidade\_curriculo\_em\_camadas; relaciona semiotica\_peirce\_triade; relaciona goldchain\_mercado\_hub\_fontes}
\prov{proveniência: gato\_operador\_hub.md \textperiodcentered{} web \textperiodcentered{} fato \textperiodcentered{} confiança 1.0 \textperiodcentered{} domínio referencia \textperiodcentered{} história 8 versões no jornal}

%CRISTAL {"arestas":[{"alvo":"gato_operador_hub_gato_a_hub_do_operador_algebrico_pow","confianca":1.0,"epistemico":"fato","origem":"gato_operador_hub.md","rel":"parte_de"},{"alvo":"ranking_global_por_acoplamento","confianca":0.5,"epistemico":"fato","origem":"wiki","rel":"relaciona"},{"alvo":"readme","confianca":0.5,"epistemico":"fato","origem":"wiki","rel":"relaciona"},{"alvo":"vontade_de_poder","confianca":0.5,"epistemico":"fato","origem":"wiki","rel":"relaciona"}],"confianca":1.0,"contraexemplos":[],"descricao":"- `neuronio/neuronio.c` - `papers/paper_bridge_areas.tex` - `ula/PAINEL_ESTELAR.md`","epistemico":"fato","exemplos":[],"id":"gato_operador_hub_fontes","memoria":"semantica","meta":{"dominio":"referencia","fonte":"web","nivel":5},"origem":"gato_operador_hub.md","palavras_chave":[],"sinonimos":[],"tipo":"conceito","titulo":"Fontes"}
\section{Fontes}
- `neuronio/neuronio.c` - `papers/paper\_bridge\_areas.tex` - `ula/PAINEL\_ESTELAR.md`
\prov{relações: parte\_de gato\_operador\_hub\_gato\_a\_hub\_do\_operador\_algebrico\_pow; relaciona ranking\_global\_por\_acoplamento; relaciona readme; relaciona vontade\_de\_poder}
\prov{proveniência: gato\_operador\_hub.md \textperiodcentered{} web \textperiodcentered{} fato \textperiodcentered{} confiança 1.0 \textperiodcentered{} domínio referencia \textperiodcentered{} história 6 versões no jornal}

%CRISTAL {"arestas":[{"alvo":"gato_operador_hub","confianca":0.95,"epistemico":"fato","origem":"wiki","rel":"parte_de"}],"confianca":1.0,"contraexemplos":[],"descricao":"> **Epistemologia:** fato de implementação (papers + código + currículo). > **No grafo:** `epistemico=fato`, confiança 0.95. > **Papel:** fechar o operador **A** do Broca-SO — distinto do mamífero homónimo.","epistemico":"fato","exemplos":[],"id":"gato_operador_hub_gato_a_hub_do_operador_algebrico_pow","memoria":"semantica","meta":{"dominio":"referencia","fonte":"web","nivel":5},"origem":"gato_operador_hub.md","palavras_chave":[],"sinonimos":[],"tipo":"conceito","titulo":"Gato A — hub do operador algébrico PoW"}
\section{Gato A — hub do operador algébrico PoW}
> **Epistemologia:** fato de implementação (papers + código + currículo). > **No grafo:** `epistemico=fato`, confiança 0.95. > **Papel:** fechar o operador **A** do Broca-SO — distinto do mamífero homónimo.
\prov{relações: parte\_de gato\_operador\_hub}
\prov{proveniência: gato\_operador\_hub.md \textperiodcentered{} web \textperiodcentered{} fato \textperiodcentered{} confiança 1.0 \textperiodcentered{} domínio referencia \textperiodcentered{} história 3 versões no jornal}

%CRISTAL {"arestas":[{"alvo":"gato_operador_hub_gato_a_hub_do_operador_algebrico_pow","confianca":1.0,"epistemico":"fato","origem":"gato_operador_hub.md","rel":"parte_de"},{"alvo":"a_decomposicao_de_peirce_celulas_e_canais","confianca":0.95,"epistemico":"fato","origem":"wiki","rel":"depende"},{"alvo":"word","confianca":0.5,"epistemico":"fato","origem":"wiki","rel":"relaciona"}],"confianca":1.0,"contraexemplos":[],"descricao":"**Operador 2×2** A = [[1,1],[1,0]] sobre bytes — núcleo da medição Broca.","epistemico":"fato","exemplos":[],"id":"gato_operador_hub_o_que_e","memoria":"semantica","meta":{"dominio":"referencia","fonte":"web","nivel":5},"origem":"gato_operador_hub.md","palavras_chave":[],"sinonimos":[],"tipo":"conceito","titulo":"O que é"}
\section{O que é}
**Operador 2$\times$2** A = [[1,1],[1,0]] sobre bytes — núcleo da medição Broca.
\prov{relações: parte\_de gato\_operador\_hub\_gato\_a\_hub\_do\_operador\_algebrico\_pow; depende a\_decomposicao\_de\_peirce\_celulas\_e\_canais; relaciona word}
\prov{proveniência: gato\_operador\_hub.md \textperiodcentered{} web \textperiodcentered{} fato \textperiodcentered{} confiança 1.0 \textperiodcentered{} domínio referencia \textperiodcentered{} história 5 versões no jornal}

%CRISTAL {"arestas":[{"alvo":"gato_operador_hub_gato_a_hub_do_operador_algebrico_pow","confianca":1.0,"epistemico":"fato","origem":"gato_operador_hub.md","rel":"parte_de"},{"alvo":"ponte_simbolica_gentil","confianca":0.95,"epistemico":"fato","origem":"wiki","rel":"confirma"},{"alvo":"corpo_tropical","confianca":0.95,"epistemico":"fato","origem":"wiki","rel":"confirma"},{"alvo":"corpo_glacial","confianca":0.95,"epistemico":"fato","origem":"wiki","rel":"confirma"},{"alvo":"broca_os","confianca":0.95,"epistemico":"fato","origem":"wiki","rel":"referencia"},{"alvo":"colapso","confianca":0.95,"epistemico":"fato","origem":"wiki","rel":"usa"},{"alvo":"hopfield","confianca":0.95,"epistemico":"fato","origem":"wiki","rel":"referencia"},{"alvo":"painel_estelar","confianca":0.95,"epistemico":"fato","origem":"wiki","rel":"explica"},{"alvo":"mecanica_quantica","confianca":0.95,"epistemico":"fato","origem":"wiki","rel":"analogia"}],"confianca":1.0,"contraexemplos":[],"descricao":"| Gato A | Nó | relação | |--------|-----|---------| | operador | gato | confirma | | decomposição | a_decomposicao_de_peirce_celulas_e_canais | depende | | ponte Lopes | ponte_simbolica_gentil | confirma | | corpo algébrico | corpo_tropical | confirma | | corpo algébrico | corpo_glacial | confirma | | medição SO | broca_os | referencia | | colapso overlap | colapso | usa | | rede associativa | hopfield | referencia |","epistemico":"fato","exemplos":[],"id":"gato_operador_hub_ponte_broca-so","memoria":"semantica","meta":{"dominio":"referencia","fonte":"web","nivel":5},"origem":"gato_operador_hub.md","palavras_chave":[],"sinonimos":[],"tipo":"conceito","titulo":"Ponte Broca-SO"}
\section{Ponte Broca-SO}
| Gato A | Nó | relação | |--------|-----|---------| | operador | gato | confirma | | decomposição | a\_decomposicao\_de\_peirce\_celulas\_e\_canais | depende | | ponte Lopes | ponte\_simbolica\_gentil | confirma | | corpo algébrico | corpo\_tropical | confirma | | corpo algébrico | corpo\_glacial | confirma | | medição SO | broca\_os | referencia | | colapso overlap | colapso | usa | | rede associativa | hopfield | referencia |
\prov{relações: parte\_de gato\_operador\_hub\_gato\_a\_hub\_do\_operador\_algebrico\_pow; confirma ponte\_simbolica\_gentil; confirma corpo\_tropical; confirma corpo\_glacial; referencia broca\_os; usa colapso; referencia hopfield; explica painel\_estelar; analogia mecanica\_quantica}
\prov{proveniência: gato\_operador\_hub.md \textperiodcentered{} web \textperiodcentered{} fato \textperiodcentered{} confiança 1.0 \textperiodcentered{} domínio referencia \textperiodcentered{} história 12 versões no jornal}

%CRISTAL {"arestas":[{"alvo":"gato_operador_hub_gato_a_hub_do_operador_algebrico_pow","confianca":1.0,"epistemico":"fato","origem":"gato_operador_hub.md","rel":"parte_de"}],"confianca":1.0,"contraexemplos":[],"descricao":"- **Paper:** `papers/paper_bridge_areas.tex`, `papers/computacao_fundamentos_broca.tex` - **Código:** `neuronio/neuronio.c` — streaming byte a byte; `neuronio/primitivas_broca.py` - **Manual:** `ula/PAINEL_ESTELAR.md` — três braços do gato - **Currículo:** fase 4½ computação → álgebra gato → PoW","epistemico":"fato","exemplos":[],"id":"gato_operador_hub_referencia","memoria":"semantica","meta":{"dominio":"referencia","fonte":"web","nivel":5},"origem":"gato_operador_hub.md","palavras_chave":[],"sinonimos":[],"tipo":"conceito","titulo":"Referência"}
\section{Referência}
- **Paper:** `papers/paper\_bridge\_areas.tex`, `papers/computacao\_fundamentos\_broca.tex` - **Código:** `neuronio/neuronio.c` — streaming byte a byte; `neuronio/primitivas\_broca.py` - **Manual:** `ula/PAINEL\_ESTELAR.md` — três braços do gato - **Currículo:** fase 4\ensuremath{\tfrac12} computação \ensuremath{\to} álgebra gato \ensuremath{\to} PoW
\prov{relações: parte\_de gato\_operador\_hub\_gato\_a\_hub\_do\_operador\_algebrico\_pow}
\prov{proveniência: gato\_operador\_hub.md \textperiodcentered{} web \textperiodcentered{} fato \textperiodcentered{} confiança 1.0 \textperiodcentered{} domínio referencia \textperiodcentered{} história 3 versões no jornal}

%CRISTAL {"arestas":[{"alvo":"razao_aurea","confianca":1.0,"epistemico":"fato","origem":"wiki","rel":"usa"},{"alvo":"corpo_tropical","confianca":1.0,"epistemico":"fato","origem":"wiki","rel":"relaciona"},{"alvo":"matematica_estelar","confianca":1.0,"epistemico":"fato","origem":"wiki","rel":"relaciona"},{"alvo":"transformada_metalica","confianca":1.0,"epistemico":"fato","origem":"wiki","rel":"relaciona"}],"confianca":1.0,"contraexemplos":[],"descricao":"Quasicristal de Penrose como 'teto da escala': fase Φ₅ pentagonal (gap 0, só gira) + plano φ² (Lucas pares 3,7,18,47…) unidos em ℤ[ζ₅]. Nova mecânica de pressão Π=1−s², imposto V=Π·‖a×b‖², 'evaporação=despressurização s→±1' = o flip de Wick como transição de fase. Três superfícies esgotam o repertório espectral: reta / espiral / pura fase.","epistemico":"inferencia","exemplos":[],"id":"genoma_penrose_zeta5","memoria":"semantica","meta":{"autor":"Lopes/Aarão","dominio":"hiper","fonte":""},"origem":"wiki","palavras_chave":[],"sinonimos":[],"tipo":"conceito","titulo":"O Genoma de Penrose e a Pressão Algébrica"}
\section{O Genoma de Penrose e a Pressão Algébrica}
Quasicristal de Penrose como 'teto da escala': fase \ensuremath{\Phi}\ensuremath{_{5}} pentagonal (gap 0, só gira) + plano \ensuremath{\varphi}\ensuremath{^{2}} (Lucas pares 3,7,18,47…) unidos em \ensuremath{\mathbb{Z}}[\ensuremath{\zeta}\ensuremath{_{5}}]. Nova mecânica de pressão \ensuremath{\Pi}=1\ensuremath{-}s\ensuremath{^{2}}, imposto V=\ensuremath{\Pi}\textperiodcentered \ensuremath{\|}a$\times$b\ensuremath{\|}\ensuremath{^{2}}, 'evaporação=despressurização s\ensuremath{\to}$\pm$1' = o flip de Wick como transição de fase. Três superfícies esgotam o repertório espectral: reta / espiral / pura fase.
\prov{relações: usa razao\_aurea; relaciona corpo\_tropical; relaciona matematica\_estelar; relaciona transformada\_metalica}
\prov{proveniência: wiki \textperiodcentered{} inferencia \textperiodcentered{} confiança 1.0 \textperiodcentered{} domínio hiper \textperiodcentered{} história 6 versões no jornal}

%CRISTAL {"arestas":[{"alvo":"gentil_fundamentos_dos_numeros","confianca":1.0,"epistemico":"fato","origem":"gentil/Gentil_Lopes_FUNDAMENTOS_DOS_NUMEROS_pdf.pdf","rel":"parte_de"},{"alvo":"gentil_fundamentos_capitulo_10_n_umeros_reais_por_cantor","confianca":1.0,"epistemico":"fato","origem":"gentil/Gentil_Lopes_FUNDAMENTOS_DOS_NUMEROS_pdf.pdf","rel":"parte_de"}],"confianca":1.0,"contraexemplos":[],"descricao":"- Linux x86_64, Python ≥ 3.10 - `make`, `gcc` (para `colapso` / `libtiffany.so` — opcional mas recomendado) - `zstd` (opcional, para tarball `.tar.zst`)","epistemico":"fato","exemplos":[],"id":"gentil_fundamentos_101_pre-requisitos","memoria":"semantica","meta":{"dominio":"manual","duplicata_sugerida":[],"nivel":6,"verificacao":"parte de conceito mais amplo 'gentil_fundamentos_101_pre-requisitos'"},"origem":"INSTALL.md","palavras_chave":[],"sinonimos":[],"tipo":"conceito","titulo":"gentil fundamentos 10.1 — pré-requisitos . . . . . . . . . . . . . . . . . . . . . . . . . . ."}
\section{gentil fundamentos 10.1 — pré-requisitos . . . . . . . . . . . . . . . . . . . . . . . . . . .}
- Linux x86\_64, Python \ensuremath{\ge} 3.10 - `make`, `gcc` (para `colapso` / `libtiffany.so` — opcional mas recomendado) - `zstd` (opcional, para tarball `.tar.zst`)
\prov{relações: parte\_de gentil\_fundamentos\_dos\_numeros; parte\_de gentil\_fundamentos\_capitulo\_10\_n\_umeros\_reais\_por\_cantor}
\prov{proveniência: INSTALL.md \textperiodcentered{} fato \textperiodcentered{} confiança 1.0 \textperiodcentered{} domínio manual \textperiodcentered{} história 6 versões no jornal}

%CRISTAL {"arestas":[{"alvo":"matematica_estelar","confianca":1.0,"epistemico":"fato","origem":"wiki","rel":"parte_de"},{"alvo":"nao_associatividade_amputacao","confianca":1.0,"epistemico":"fato","origem":"wiki","rel":"usa"},{"alvo":"regua_de_gentil","confianca":1.0,"epistemico":"fato","origem":"wiki","rel":"relaciona"}],"confianca":1.0,"contraexemplos":[],"descricao":"Métrica ‖z‖²=1−s·c² (euclidiana em s=0, tipo-Minkowski em s>0); a curvatura é gerada pelo gap K∝s(1−s)·Associador². O espaço plano (s=0) é estéril; o estelar (s>0) gera estrutura. Em dim 8, estrutura G₂ e geometria 'de Sitter estelar'.","epistemico":"fato","exemplos":[],"id":"geometria_estelar","memoria":"semantica","meta":{"autor":"Lopes/Aarão","dominio":"hiper","fonte":""},"origem":"wiki","palavras_chave":[],"sinonimos":[],"tipo":"conceito","titulo":"Geometria Estelar (curvatura pelo gap)"}
\section{Geometria Estelar (curvatura pelo gap)}
Métrica \ensuremath{\|}z\ensuremath{\|}\ensuremath{^{2}}=1\ensuremath{-}s\textperiodcentered c\ensuremath{^{2}} (euclidiana em s=0, tipo-Minkowski em s>0); a curvatura é gerada pelo gap K\ensuremath{\propto}s(1\ensuremath{-}s)\textperiodcentered Associador\ensuremath{^{2}}. O espaço plano (s=0) é estéril; o estelar (s>0) gera estrutura. Em dim 8, estrutura G\ensuremath{_{2}} e geometria 'de Sitter estelar'.
\prov{relações: parte\_de matematica\_estelar; usa nao\_associatividade\_amputacao; relaciona regua\_de\_gentil}
\prov{proveniência: wiki \textperiodcentered{} fato \textperiodcentered{} confiança 1.0 \textperiodcentered{} domínio hiper \textperiodcentered{} história 5 versões no jornal}

%CRISTAL {"arestas":[{"alvo":"geometria","confianca":0.95,"epistemico":"fato","origem":"wiki","rel":"confirma"},{"alvo":"motivos","confianca":0.5,"epistemico":"fato","origem":"wiki","rel":"relaciona"},{"alvo":"por_sessao_conversadadostelemetriasessaojsonl","confianca":0.5,"epistemico":"fato","origem":"wiki","rel":"relaciona"},{"alvo":"topologia","confianca":0.5,"epistemico":"fato","origem":"wiki","rel":"relaciona"},{"alvo":"transcricao","confianca":0.5,"epistemico":"fato","origem":"wiki","rel":"relaciona"},{"alvo":"projeto_cristaljson","confianca":0.5,"epistemico":"fato","origem":"wiki","rel":"relaciona"},{"alvo":"pow_geometria","confianca":0.5,"epistemico":"fato","origem":"wiki","rel":"relaciona"},{"alvo":"parabola_pa","confianca":0.5,"epistemico":"fato","origem":"wiki","rel":"relaciona"},{"alvo":"resposta_a_pergunta_dificil","confianca":0.5,"epistemico":"fato","origem":"wiki","rel":"relaciona"},{"alvo":"mecanica_quantica_hub","confianca":0.5,"epistemico":"fato","origem":"wiki","rel":"relaciona"},{"alvo":"thread","confianca":0.5,"epistemico":"fato","origem":"wiki","rel":"relaciona"},{"alvo":"ponteiro","confianca":0.5,"epistemico":"fato","origem":"wiki","rel":"relaciona"},{"alvo":"recursao","confianca":0.5,"epistemico":"fato","origem":"wiki","rel":"relaciona"},{"alvo":"pow_metal_pre_fecho","confianca":0.5,"epistemico":"fato","origem":"wiki","rel":"relaciona"},{"alvo":"log_legado_conversadadostelemetriajsonl","confianca":0.5,"epistemico":"fato","origem":"wiki","rel":"relaciona"},{"alvo":"metais","confianca":0.5,"epistemico":"fato","origem":"wiki","rel":"relaciona"},{"alvo":"testes","confianca":0.5,"epistemico":"fato","origem":"wiki","rel":"relaciona"},{"alvo":"rust","confianca":0.5,"epistemico":"fato","origem":"wiki","rel":"relaciona"},{"alvo":"pow_metal_ressonancia","confianca":0.5,"epistemico":"fato","origem":"wiki","rel":"relaciona"},{"alvo":"vida-morte","confianca":0.5,"epistemico":"fato","origem":"wiki","rel":"relaciona"},{"alvo":"pilha","confianca":0.5,"epistemico":"fato","origem":"wiki","rel":"relaciona"},{"alvo":"pow_metal_setor","confianca":0.5,"epistemico":"fato","origem":"wiki","rel":"relaciona"},{"alvo":"pow_metal_funil","confianca":0.5,"epistemico":"fato","origem":"wiki","rel":"relaciona"},{"alvo":"tiffany_cabeca_broca","confianca":0.5,"epistemico":"fato","origem":"wiki","rel":"relaciona"},{"alvo":"teste_ordem","confianca":0.5,"epistemico":"fato","origem":"wiki","rel":"relaciona"},{"alvo":"numa","confianca":0.5,"epistemico":"fato","origem":"wiki","rel":"relaciona"},{"alvo":"hardware_sequencias_de_cauchy","confianca":0.5,"epistemico":"fato","origem":"wiki","rel":"relaciona"},{"alvo":"lobo_frontal","confianca":0.5,"epistemico":"fato","origem":"wiki","rel":"relaciona"}],"confianca":0.95,"contraexemplos":[],"descricao":"Referência verificada: geometria_hub_broca.md","epistemico":"fato","exemplos":[],"id":"geometria_hub_broca","memoria":"semantica","meta":{"dominio":"referencia","fonte":"web","nivel":5},"origem":"wiki","palavras_chave":[],"sinonimos":[],"tipo":"conceito","titulo":"geometria hub broca"}
\section{geometria hub broca}
Referência verificada: geometria\_hub\_broca.md
\prov{relações: confirma geometria; relaciona motivos; relaciona por\_sessao\_conversadadostelemetriasessaojsonl; relaciona topologia; relaciona transcricao; relaciona projeto\_cristaljson; relaciona pow\_geometria; relaciona parabola\_pa; relaciona resposta\_a\_pergunta\_dificil; relaciona mecanica\_quantica\_hub; relaciona thread; relaciona ponteiro; relaciona recursao; relaciona pow\_metal\_pre\_fecho; relaciona log\_legado\_conversadadostelemetriajsonl; relaciona metais; relaciona testes; relaciona rust; relaciona pow\_metal\_ressonancia; relaciona vida-morte; relaciona pilha; relaciona pow\_metal\_setor; relaciona pow\_metal\_funil; relaciona tiffany\_cabeca\_broca; relaciona teste\_ordem; relaciona numa; relaciona hardware\_sequencias\_de\_cauchy; relaciona lobo\_frontal}
\prov{proveniência: wiki \textperiodcentered{} web \textperiodcentered{} fato \textperiodcentered{} confiança 0.95 \textperiodcentered{} domínio referencia \textperiodcentered{} história 36 versões no jornal}

%CRISTAL {"arestas":[{"alvo":"geometria_hub_broca_geometria_hub_broca-so","confianca":1.0,"epistemico":"fato","origem":"geometria_hub_broca.md","rel":"parte_de"},{"alvo":"word","confianca":0.5,"epistemico":"fato","origem":"wiki","rel":"relaciona"},{"alvo":"vontade_de_poder","confianca":0.5,"epistemico":"fato","origem":"wiki","rel":"relaciona"},{"alvo":"virtualizacao","confianca":0.5,"epistemico":"fato","origem":"wiki","rel":"relaciona"}],"confianca":1.0,"contraexemplos":[],"descricao":"Identificar **toro φ-Dougherty** com **fibrado Hopf S³→S²** — topologias distintas (confirmado em `fibracao_hopf_bloch`).","epistemico":"fato","exemplos":[],"id":"geometria_hub_broca_contraexemplos","memoria":"semantica","meta":{"dominio":"referencia","fonte":"web","nivel":5},"origem":"geometria_hub_broca.md","palavras_chave":[],"sinonimos":[],"tipo":"conceito","titulo":"Contraexemplos"}
\section{Contraexemplos}
Identificar **toro \ensuremath{\varphi}-Dougherty** com **fibrado Hopf S\ensuremath{^{3}}\ensuremath{\to}S\ensuremath{^{2}}** — topologias distintas (confirmado em `fibracao\_hopf\_bloch`).
\prov{relações: parte\_de geometria\_hub\_broca\_geometria\_hub\_broca-so; relaciona word; relaciona vontade\_de\_poder; relaciona virtualizacao}
\prov{proveniência: geometria\_hub\_broca.md \textperiodcentered{} web \textperiodcentered{} fato \textperiodcentered{} confiança 1.0 \textperiodcentered{} domínio referencia \textperiodcentered{} história 6 versões no jornal}

%CRISTAL {"arestas":[{"alvo":"geometria_hub_broca_geometria_hub_broca-so","confianca":1.0,"epistemico":"fato","origem":"geometria_hub_broca.md","rel":"parte_de"},{"alvo":"word","confianca":0.5,"epistemico":"fato","origem":"wiki","rel":"relaciona"},{"alvo":"virtualizacao","confianca":0.5,"epistemico":"fato","origem":"wiki","rel":"relaciona"}],"confianca":1.0,"contraexemplos":[],"descricao":"- `papers/geometria.tex` - `ula/POW_GEOMETRIA.md` - `papers/fibracao_hopf_bloch.md` - `papers/fase_geometrica_berry.md`","epistemico":"fato","exemplos":[],"id":"geometria_hub_broca_fontes","memoria":"semantica","meta":{"dominio":"referencia","fonte":"web","nivel":5},"origem":"geometria_hub_broca.md","palavras_chave":[],"sinonimos":[],"tipo":"conceito","titulo":"Fontes"}
\section{Fontes}
- `papers/geometria.tex` - `ula/POW\_GEOMETRIA.md` - `papers/fibracao\_hopf\_bloch.md` - `papers/fase\_geometrica\_berry.md`
\prov{relações: parte\_de geometria\_hub\_broca\_geometria\_hub\_broca-so; relaciona word; relaciona virtualizacao}
\prov{proveniência: geometria\_hub\_broca.md \textperiodcentered{} web \textperiodcentered{} fato \textperiodcentered{} confiança 1.0 \textperiodcentered{} domínio referencia \textperiodcentered{} história 5 versões no jornal}

%CRISTAL {"fusao":[{"arestas":[{"alvo":"geometria_hub_broca","confianca":0.95,"epistemico":"fato","origem":"wiki","rel":"parte_de"}],"confianca":1.0,"contraexemplos":[],"descricao":"> **Epistemologia:** fato matemático-físico + leituras Broca (papers verificados). > **No grafo:** `epistemico=fato`, confiança 0.95. > **Papel:** fechar contraexemplos e consolidar pontes externas (Hopf, Dougherty, PoW).","epistemico":"fato","exemplos":[],"id":"geometria_hub_broca_geometria_hub_broca-so","memoria":"semantica","meta":{"dominio":"referencia","fonte":"web","nivel":5},"origem":"geometria_hub_broca.md","palavras_chave":[],"sinonimos":[],"tipo":"conceito","titulo":"Geometria — hub Broca-SO"},{"arestas":[{"alvo":"geometria_hub_broca","confianca":0.95,"epistemico":"fato","origem":"wiki","rel":"parte_de"}],"confianca":1.0,"contraexemplos":[],"descricao":"> **Epistemologia:** fato matemático-físico + leituras Broca (papers verificados). > **No grafo:** `epistemico=fato`, confiança 0.95. > **Papel:** fechar contraexemplos e consolidar pontes externas (Hopf, Dougherty, PoW).","epistemico":"fato","exemplos":[],"id":"geometria_hub_broca-so","memoria":"semantica","meta":{"dominio":"referencia","nivel":5},"origem":"geometria_hub_broca.md","palavras_chave":[],"sinonimos":[],"tipo":"conceito","titulo":"Geometria — hub Broca-SO"}],"id":"geometria_hub_broca_geometria_hub_broca-so","tipo":"conceito"}
\section{Geometria — hub Broca-SO}
> **Epistemologia:** fato matemático-físico + leituras Broca (papers verificados). > **No grafo:** `epistemico=fato`, confiança 0.95. > **Papel:** fechar contraexemplos e consolidar pontes externas (Hopf, Dougherty, PoW).
\prov{relações: parte\_de geometria\_hub\_broca}
\prov{proveniência: geometria\_hub\_broca.md \textperiodcentered{} web \textperiodcentered{} fato \textperiodcentered{} confiança 1.0 \textperiodcentered{} domínio referencia \textperiodcentered{} história 3 versões no jornal}
\prov{fusão de curadoria (texto idêntico): absorve geometria\_hub\_broca-so --- o mesmo conteúdo por dois esquemas de endereço}

%CRISTAL {"fusao":[{"arestas":[{"alvo":"geometria_hub_broca_geometria_hub_broca-so","confianca":1.0,"epistemico":"fato","origem":"geometria_hub_broca.md","rel":"parte_de"},{"alvo":"o_que_e_no_projecto","confianca":0.5,"epistemico":"fato","origem":"wiki","rel":"relaciona"},{"alvo":"habitantes_gentil_de_sitter","confianca":0.5,"epistemico":"fato","origem":"wiki","rel":"relaciona"},{"alvo":"word","confianca":0.5,"epistemico":"fato","origem":"wiki","rel":"relaciona"}],"confianca":1.0,"contraexemplos":[],"descricao":"Geometria = forma, distância, curvatura, transformações.","epistemico":"fato","exemplos":[],"id":"geometria_hub_broca_o_que_e_no_projecto","memoria":"semantica","meta":{"dominio":"referencia","fonte":"web","nivel":5},"origem":"geometria_hub_broca.md","palavras_chave":[],"sinonimos":[],"tipo":"conceito","titulo":"O que é (no projecto)"},{"arestas":[{"alvo":"geometria_hub_broca-so","confianca":1.0,"epistemico":"fato","origem":"geometria_hub_broca.md","rel":"parte_de"},{"alvo":"utilitarismo","confianca":0.5,"epistemico":"fato","origem":"wiki","rel":"relaciona"},{"alvo":"teoria_do_valor","confianca":0.5,"epistemico":"fato","origem":"wiki","rel":"relaciona"},{"alvo":"virtualizacao","confianca":0.5,"epistemico":"fato","origem":"wiki","rel":"relaciona"},{"alvo":"word","confianca":0.5,"epistemico":"fato","origem":"wiki","rel":"relaciona"},{"alvo":"vontade_de_poder","confianca":0.5,"epistemico":"fato","origem":"wiki","rel":"relaciona"}],"confianca":1.0,"contraexemplos":[],"descricao":"Geometria = forma, distância, curvatura, transformações.","epistemico":"fato","exemplos":[],"id":"o_que_e_no_projecto","memoria":"semantica","meta":{"dominio":"referencia","nivel":5},"origem":"geometria_hub_broca.md","palavras_chave":[],"sinonimos":[],"tipo":"conceito","titulo":"O que é (no projecto)"}],"id":"geometria_hub_broca_o_que_e_no_projecto","tipo":"conceito"}
\section{O que é (no projecto)}
Geometria = forma, distância, curvatura, transformações.
\prov{relações: parte\_de geometria\_hub\_broca\_geometria\_hub\_broca-so; relaciona o\_que\_e\_no\_projecto; relaciona habitantes\_gentil\_de\_sitter; relaciona word}
\prov{proveniência: geometria\_hub\_broca.md \textperiodcentered{} web \textperiodcentered{} fato \textperiodcentered{} confiança 1.0 \textperiodcentered{} domínio referencia \textperiodcentered{} história 6 versões no jornal}
\prov{fusão de curadoria (texto idêntico): absorve o\_que\_e\_no\_projecto --- o mesmo conteúdo por dois esquemas de endereço}

%CRISTAL {"fusao":[{"arestas":[{"alvo":"geometria_hub_broca_geometria_hub_broca-so","confianca":1.0,"epistemico":"fato","origem":"geometria_hub_broca.md","rel":"parte_de"},{"alvo":"ponte_quantica_dougherty","confianca":0.95,"epistemico":"fato","origem":"wiki","rel":"analogia"},{"alvo":"ginzburg_toro_helice","confianca":0.95,"epistemico":"fato","origem":"wiki","rel":"analogia"},{"alvo":"conjunto_dougherty","confianca":0.95,"epistemico":"fato","origem":"wiki","rel":"analogia"},{"alvo":"birkeland_correntes_plasma","confianca":0.95,"epistemico":"fato","origem":"wiki","rel":"analogia"}],"confianca":1.0,"contraexemplos":[],"descricao":"| Hipótese | Geometria | epistemologia | |----------|-----------|---------------| | Dougherty toros φ | geometria | analogia | | Ginzburg 3D-SST | geometria | analogia | | Birkeland filamentos | geometria | analogia visual |","epistemico":"fato","exemplos":[],"id":"geometria_hub_broca_pontes_conjecturais_hipotese_analogia","memoria":"semantica","meta":{"dominio":"referencia","fonte":"web","nivel":5},"origem":"geometria_hub_broca.md","palavras_chave":[],"sinonimos":[],"tipo":"conceito","titulo":"Pontes conjecturais (hipótese / analogia)"},{"arestas":[{"alvo":"geometria_hub_broca-so","confianca":1.0,"epistemico":"fato","origem":"geometria_hub_broca.md","rel":"parte_de"},{"alvo":"ponte_quantica_dougherty","confianca":0.95,"epistemico":"fato","origem":"wiki","rel":"analogia"},{"alvo":"ginzburg_toro_helice","confianca":0.95,"epistemico":"fato","origem":"wiki","rel":"analogia"},{"alvo":"conjunto_dougherty","confianca":0.95,"epistemico":"fato","origem":"wiki","rel":"analogia"},{"alvo":"birkeland_correntes_plasma","confianca":0.95,"epistemico":"fato","origem":"wiki","rel":"analogia"}],"confianca":1.0,"contraexemplos":[],"descricao":"| Hipótese | Geometria | epistemologia | |----------|-----------|---------------| | Dougherty toros φ | geometria | analogia | | Ginzburg 3D-SST | geometria | analogia | | Birkeland filamentos | geometria | analogia visual |","epistemico":"fato","exemplos":[],"id":"pontes_conjecturais_hipotese_analogia","memoria":"semantica","meta":{"dominio":"referencia","nivel":5},"origem":"geometria_hub_broca.md","palavras_chave":[],"sinonimos":[],"tipo":"conceito","titulo":"Pontes conjecturais (hipótese / analogia)"}],"id":"geometria_hub_broca_pontes_conjecturais_hipotese_analogia","tipo":"conceito"}
\section{Pontes conjecturais (hipótese / analogia)}
| Hipótese | Geometria | epistemologia | |----------|-----------|---------------| | Dougherty toros \ensuremath{\varphi} | geometria | analogia | | Ginzburg 3D-SST | geometria | analogia | | Birkeland filamentos | geometria | analogia visual |
\prov{relações: parte\_de geometria\_hub\_broca\_geometria\_hub\_broca-so; analogia ponte\_quantica\_dougherty; analogia ginzburg\_toro\_helice; analogia conjunto\_dougherty; analogia birkeland\_correntes\_plasma}
\prov{proveniência: geometria\_hub\_broca.md \textperiodcentered{} web \textperiodcentered{} fato \textperiodcentered{} confiança 1.0 \textperiodcentered{} domínio referencia \textperiodcentered{} história 7 versões no jornal}
\prov{fusão de curadoria (texto idêntico): absorve pontes\_conjecturais\_hipotese\_analogia --- o mesmo conteúdo por dois esquemas de endereço}

%CRISTAL {"fusao":[{"arestas":[{"alvo":"geometria_hub_broca_geometria_hub_broca-so","confianca":1.0,"epistemico":"fato","origem":"geometria_hub_broca.md","rel":"parte_de"},{"alvo":"fibracao_hopf_bloch","confianca":0.95,"epistemico":"fato","origem":"wiki","rel":"confirma"},{"alvo":"fase_geometrica_berry","confianca":0.95,"epistemico":"fato","origem":"wiki","rel":"confirma"},{"alvo":"pow_geometria","confianca":0.95,"epistemico":"fato","origem":"wiki","rel":"referencia"},{"alvo":"a_metrica_de_sitter_global","confianca":0.95,"epistemico":"fato","origem":"wiki","rel":"confirma"}],"confianca":1.0,"contraexemplos":[],"descricao":"| Fonte | Geometria Broca | relação | |-------|-----------------|---------| | Hopf S³→S² | a_medida_o_calor_de_volta | confirma | | Bloch CP¹ | a_funcao_de_onda_a_esfera_de_bloch | confirma | | Berry γ | a_dinamica_a_precessao_de_bloch | confirma | | de Sitter | a_metrica_de_sitter_global | confirma | | PoW atlas | pow_geometria | explica |","epistemico":"fato","exemplos":[],"id":"geometria_hub_broca_pontes_verificadas_fato","memoria":"semantica","meta":{"dominio":"referencia","fonte":"web","nivel":5},"origem":"geometria_hub_broca.md","palavras_chave":[],"sinonimos":[],"tipo":"conceito","titulo":"Pontes verificadas (fato)"},{"arestas":[{"alvo":"geometria_hub_broca-so","confianca":1.0,"epistemico":"fato","origem":"geometria_hub_broca.md","rel":"parte_de"},{"alvo":"fibracao_hopf_bloch","confianca":0.95,"epistemico":"fato","origem":"wiki","rel":"confirma"},{"alvo":"fase_geometrica_berry","confianca":0.95,"epistemico":"fato","origem":"wiki","rel":"confirma"},{"alvo":"pow_geometria","confianca":0.95,"epistemico":"fato","origem":"wiki","rel":"referencia"},{"alvo":"a_metrica_de_sitter_global","confianca":0.95,"epistemico":"fato","origem":"wiki","rel":"confirma"}],"confianca":1.0,"contraexemplos":[],"descricao":"| Fonte | Geometria Broca | relação | |-------|-----------------|---------| | Hopf S³→S² | a_medida_o_calor_de_volta | confirma | | Bloch CP¹ | a_funcao_de_onda_a_esfera_de_bloch | confirma | | Berry γ | a_dinamica_a_precessao_de_bloch | confirma | | de Sitter | a_metrica_de_sitter_global | confirma | | PoW atlas | pow_geometria | explica |","epistemico":"fato","exemplos":[],"id":"pontes_verificadas_fato","memoria":"semantica","meta":{"dominio":"referencia","nivel":5},"origem":"geometria_hub_broca.md","palavras_chave":[],"sinonimos":[],"tipo":"conceito","titulo":"Pontes verificadas (fato)"}],"id":"geometria_hub_broca_pontes_verificadas_fato","tipo":"conceito"}
\section{Pontes verificadas (fato)}
| Fonte | Geometria Broca | relação | |-------|-----------------|---------| | Hopf S\ensuremath{^{3}}\ensuremath{\to}S\ensuremath{^{2}} | a\_medida\_o\_calor\_de\_volta | confirma | | Bloch CP\ensuremath{^{1}} | a\_funcao\_de\_onda\_a\_esfera\_de\_bloch | confirma | | Berry \ensuremath{\gamma} | a\_dinamica\_a\_precessao\_de\_bloch | confirma | | de Sitter | a\_metrica\_de\_sitter\_global | confirma | | PoW atlas | pow\_geometria | explica |
\prov{relações: parte\_de geometria\_hub\_broca\_geometria\_hub\_broca-so; confirma fibracao\_hopf\_bloch; confirma fase\_geometrica\_berry; referencia pow\_geometria; confirma a\_metrica\_de\_sitter\_global}
\prov{proveniência: geometria\_hub\_broca.md \textperiodcentered{} web \textperiodcentered{} fato \textperiodcentered{} confiança 1.0 \textperiodcentered{} domínio referencia \textperiodcentered{} história 7 versões no jornal}
\prov{fusão de curadoria (texto idêntico): absorve pontes\_verificadas\_fato --- o mesmo conteúdo por dois esquemas de endereço}

%CRISTAL {"arestas":[{"alvo":"geometria_hub_broca_geometria_hub_broca-so","confianca":1.0,"epistemico":"fato","origem":"geometria_hub_broca.md","rel":"parte_de"}],"confianca":1.0,"contraexemplos":[],"descricao":"- **Paper interno:** `papers/geometria.tex`, `papers/quatro_ciencias.tex` - **Manual:** `ula/POW_GEOMETRIA.md` — coordenadas espectrais PoW - **Confirmado web:** `fibracao_hopf_bloch` (S³→S²), `fase_geometrica_berry` - **Hipótese externa:** `ponte_quantica_dougherty`, `ginzburg_toro_helice` — **analogia**","epistemico":"fato","exemplos":[],"id":"geometria_hub_broca_referencia","memoria":"semantica","meta":{"dominio":"referencia","fonte":"web","nivel":5},"origem":"geometria_hub_broca.md","palavras_chave":[],"sinonimos":[],"tipo":"conceito","titulo":"Referência"}
\section{Referência}
- **Paper interno:** `papers/geometria.tex`, `papers/quatro\_ciencias.tex` - **Manual:** `ula/POW\_GEOMETRIA.md` — coordenadas espectrais PoW - **Confirmado web:** `fibracao\_hopf\_bloch` (S\ensuremath{^{3}}\ensuremath{\to}S\ensuremath{^{2}}), `fase\_geometrica\_berry` - **Hipótese externa:** `ponte\_quantica\_dougherty`, `ginzburg\_toro\_helice` — **analogia**
\prov{relações: parte\_de geometria\_hub\_broca\_geometria\_hub\_broca-so}
\prov{proveniência: geometria\_hub\_broca.md \textperiodcentered{} web \textperiodcentered{} fato \textperiodcentered{} confiança 1.0 \textperiodcentered{} domínio referencia \textperiodcentered{} história 3 versões no jornal}

%CRISTAL {"arestas":[{"alvo":"corpo_tropical","confianca":1.0,"epistemico":"fato","origem":"wiki","rel":"usa"},{"alvo":"o_substrato_a_maquina_e_funcao_de_onda","confianca":1.0,"epistemico":"fato","origem":"wiki","rel":"relaciona"}],"confianca":1.0,"contraexemplos":[],"descricao":"A realização riemanniana derivada (não importada) da dinâmica de ouro: da topologia da torre e da álgebra da balança nasce a métrica, via a cadeia S³ dos versores → de Sitter global → flip de Wick espacial → palco hiperbólico ℍ³/Γ de Bianchi.","epistemico":"fato","exemplos":[],"id":"geometria_riemanniana_de_ouro","memoria":"semantica","meta":{"autor":"Aarão/broca-so","dominio":"broca_repos","fonte":""},"origem":"wiki","palavras_chave":[],"sinonimos":[],"tipo":"conceito","titulo":"Geometria Riemanniana de Ouro"}
\section{Geometria Riemanniana de Ouro}
A realização riemanniana derivada (não importada) da dinâmica de ouro: da topologia da torre e da álgebra da balança nasce a métrica, via a cadeia S\ensuremath{^{3}} dos versores \ensuremath{\to} de Sitter global \ensuremath{\to} flip de Wick espacial \ensuremath{\to} palco hiperbólico \ensuremath{\mathbb{H}}\ensuremath{^{3}}/\ensuremath{\Gamma} de Bianchi.
\prov{relações: usa corpo\_tropical; relaciona o\_substrato\_a\_maquina\_e\_funcao\_de\_onda}
\prov{proveniência: wiki \textperiodcentered{} fato \textperiodcentered{} confiança 1.0 \textperiodcentered{} domínio broca\_repos \textperiodcentered{} história 3 versões no jornal}

%CRISTAL {"arestas":[{"alvo":"goldchain","confianca":0.95,"epistemico":"fato","origem":"wiki","rel":"confirma"},{"alvo":"mecanica_quantica_hub","confianca":0.5,"epistemico":"fato","origem":"wiki","rel":"relaciona"},{"alvo":"historia","confianca":0.5,"epistemico":"fato","origem":"wiki","rel":"relaciona"},{"alvo":"kernel","confianca":0.5,"epistemico":"fato","origem":"wiki","rel":"relaciona"},{"alvo":"topologia","confianca":0.5,"epistemico":"fato","origem":"wiki","rel":"relaciona"},{"alvo":"word","confianca":0.5,"epistemico":"fato","origem":"wiki","rel":"relaciona"}],"confianca":0.95,"contraexemplos":[],"descricao":"Referência verificada: goldchain_mercado_hub.md","epistemico":"fato","exemplos":[],"id":"goldchain_mercado_hub","memoria":"semantica","meta":{"dominio":"referencia","fonte":"web","nivel":5},"origem":"wiki","palavras_chave":[],"sinonimos":[],"tipo":"conceito","titulo":"goldchain mercado hub"}
\section{goldchain mercado hub}
Referência verificada: goldchain\_mercado\_hub.md
\prov{relações: confirma goldchain; relaciona mecanica\_quantica\_hub; relaciona historia; relaciona kernel; relaciona topologia; relaciona word}
\prov{proveniência: wiki \textperiodcentered{} web \textperiodcentered{} fato \textperiodcentered{} confiança 0.95 \textperiodcentered{} domínio referencia \textperiodcentered{} história 8 versões no jornal}

%CRISTAL {"arestas":[{"alvo":"goldchain_mercado_hub_goldchain_hub_do_mercado_broca","confianca":1.0,"epistemico":"fato","origem":"goldchain_mercado_hub.md","rel":"parte_de"},{"alvo":"word","confianca":0.5,"epistemico":"fato","origem":"wiki","rel":"relaciona"},{"alvo":"virtualizacao","confianca":0.5,"epistemico":"fato","origem":"wiki","rel":"relaciona"},{"alvo":"universidade_curriculo_em_camadas","confianca":0.5,"epistemico":"fato","origem":"wiki","rel":"relaciona"}],"confianca":1.0,"contraexemplos":[],"descricao":"GoldChain **≠** DeFi genérico com gas ETH — modelo lastreado PoW→Pell→tesouraria.","epistemico":"fato","exemplos":[],"id":"goldchain_mercado_hub_contraexemplos","memoria":"semantica","meta":{"dominio":"referencia","fonte":"web","nivel":5},"origem":"goldchain_mercado_hub.md","palavras_chave":[],"sinonimos":[],"tipo":"conceito","titulo":"Contraexemplos"}
\section{Contraexemplos}
GoldChain **\ensuremath{\ne}** DeFi genérico com gas ETH — modelo lastreado PoW\ensuremath{\to}Pell\ensuremath{\to}tesouraria.
\prov{relações: parte\_de goldchain\_mercado\_hub\_goldchain\_hub\_do\_mercado\_broca; relaciona word; relaciona virtualizacao; relaciona universidade\_curriculo\_em\_camadas}
\prov{proveniência: goldchain\_mercado\_hub.md \textperiodcentered{} web \textperiodcentered{} fato \textperiodcentered{} confiança 1.0 \textperiodcentered{} domínio referencia \textperiodcentered{} história 6 versões no jornal}

%CRISTAL {"arestas":[{"alvo":"goldchain_mercado_hub_goldchain_hub_do_mercado_broca","confianca":1.0,"epistemico":"fato","origem":"goldchain_mercado_hub.md","rel":"parte_de"},{"alvo":"word","confianca":0.5,"epistemico":"fato","origem":"wiki","rel":"relaciona"},{"alvo":"virtualizacao","confianca":0.5,"epistemico":"fato","origem":"wiki","rel":"relaciona"},{"alvo":"universidade_curriculo_em_camadas","confianca":0.5,"epistemico":"fato","origem":"wiki","rel":"relaciona"},{"alvo":"vida-morte","confianca":0.5,"epistemico":"fato","origem":"wiki","rel":"relaciona"},{"alvo":"semiotica_peirce_triade","confianca":0.5,"epistemico":"fato","origem":"wiki","rel":"relaciona"},{"alvo":"proposition_gerador_de_escala_ipropderiv_com_tddt","confianca":0.5,"epistemico":"fato","origem":"wiki","rel":"relaciona"},{"alvo":"inconsciente","confianca":0.5,"epistemico":"fato","origem":"wiki","rel":"relaciona"}],"confianca":1.0,"contraexemplos":[],"descricao":"- `goldchain/README.md` - `kernel/algebra_dos_contratos.tex` - `goldchain/fluxo.py`","epistemico":"fato","exemplos":[],"id":"goldchain_mercado_hub_fontes","memoria":"semantica","meta":{"dominio":"referencia","fonte":"web","nivel":5},"origem":"goldchain_mercado_hub.md","palavras_chave":[],"sinonimos":[],"tipo":"conceito","titulo":"Fontes"}
\section{Fontes}
- `goldchain/README.md` - `kernel/algebra\_dos\_contratos.tex` - `goldchain/fluxo.py`
\prov{relações: parte\_de goldchain\_mercado\_hub\_goldchain\_hub\_do\_mercado\_broca; relaciona word; relaciona virtualizacao; relaciona universidade\_curriculo\_em\_camadas; relaciona vida-morte; relaciona semiotica\_peirce\_triade; relaciona proposition\_gerador\_de\_escala\_ipropderiv\_com\_tddt; relaciona inconsciente}
\prov{proveniência: goldchain\_mercado\_hub.md \textperiodcentered{} web \textperiodcentered{} fato \textperiodcentered{} confiança 1.0 \textperiodcentered{} domínio referencia \textperiodcentered{} história 10 versões no jornal}

%CRISTAL {"arestas":[{"alvo":"goldchain_mercado_hub","confianca":0.95,"epistemico":"fato","origem":"wiki","rel":"parte_de"}],"confianca":1.0,"contraexemplos":[],"descricao":"> **Epistemologia:** fato de implementação (código + papers + currículo). > **No grafo:** `epistemico=fato`, confiança 0.95. > **Papel:** fechar álgebra dos contratos, generalizações e contraexemplos.","epistemico":"fato","exemplos":[],"id":"goldchain_mercado_hub_goldchain_hub_do_mercado_broca","memoria":"semantica","meta":{"dominio":"referencia","fonte":"web","nivel":5},"origem":"goldchain_mercado_hub.md","palavras_chave":[],"sinonimos":[],"tipo":"conceito","titulo":"GoldChain — hub do mercado Broca"}
\section{GoldChain — hub do mercado Broca}
> **Epistemologia:** fato de implementação (código + papers + currículo). > **No grafo:** `epistemico=fato`, confiança 0.95. > **Papel:** fechar álgebra dos contratos, generalizações e contraexemplos.
\prov{relações: parte\_de goldchain\_mercado\_hub}
\prov{proveniência: goldchain\_mercado\_hub.md \textperiodcentered{} web \textperiodcentered{} fato \textperiodcentered{} confiança 1.0 \textperiodcentered{} domínio referencia \textperiodcentered{} história 3 versões no jornal}

%CRISTAL {"arestas":[{"alvo":"goldchain_mercado_hub_goldchain_hub_do_mercado_broca","confianca":1.0,"epistemico":"fato","origem":"goldchain_mercado_hub.md","rel":"parte_de"},{"alvo":"valor","confianca":0.95,"epistemico":"fato","origem":"wiki","rel":"depende"},{"alvo":"dinheiro","confianca":0.95,"epistemico":"fato","origem":"wiki","rel":"depende"},{"alvo":"sha256","confianca":0.95,"epistemico":"fato","origem":"wiki","rel":"usa"}],"confianca":1.0,"contraexemplos":[],"descricao":"**Álgebra dos contratos Broca:** C = (P, φ, λ) — partes, condição, liquidação.","epistemico":"fato","exemplos":[],"id":"goldchain_mercado_hub_o_que_e","memoria":"semantica","meta":{"dominio":"referencia","fonte":"web","nivel":5},"origem":"goldchain_mercado_hub.md","palavras_chave":[],"sinonimos":[],"tipo":"conceito","titulo":"O que é"}
\section{O que é}
**Álgebra dos contratos Broca:** C = (P, \ensuremath{\varphi}, \ensuremath{\lambda}) — partes, condição, liquidação.
\prov{relações: parte\_de goldchain\_mercado\_hub\_goldchain\_hub\_do\_mercado\_broca; depende valor; depende dinheiro; usa sha256}
\prov{proveniência: goldchain\_mercado\_hub.md \textperiodcentered{} web \textperiodcentered{} fato \textperiodcentered{} confiança 1.0 \textperiodcentered{} domínio referencia \textperiodcentered{} história 6 versões no jornal}

%CRISTAL {"arestas":[{"alvo":"goldchain_mercado_hub_goldchain_hub_do_mercado_broca","confianca":1.0,"epistemico":"fato","origem":"goldchain_mercado_hub.md","rel":"parte_de"},{"alvo":"blockchain","confianca":0.95,"epistemico":"fato","origem":"wiki","rel":"especializa"},{"alvo":"cifra_de_ouro_branco","confianca":0.95,"epistemico":"fato","origem":"wiki","rel":"referencia"},{"alvo":"contabilidade","confianca":0.95,"epistemico":"fato","origem":"wiki","rel":"explica"},{"alvo":"broca_os","confianca":0.95,"epistemico":"fato","origem":"wiki","rel":"referencia"},{"alvo":"codigo_goldchain_fluxo","confianca":0.95,"epistemico":"fato","origem":"wiki","rel":"confirma"},{"alvo":"gato","confianca":0.95,"epistemico":"fato","origem":"wiki","rel":"usa"}],"confianca":1.0,"contraexemplos":[],"descricao":"| GoldChain | Nó | relação | |-----------|-----|---------| | núcleo | goldchain | confirma | | cadeia | blockchain | especializa | | valor | valor | depende | | moeda | dinheiro | depende | | hash | sha256 | usa | | cifra | cifra_de_ouro_branco | referencia | | contabilidade | contabilidade | explica | | SO | broca_os | referencia | | código | codigo_goldchain_fluxo | confirma | | PoW | gato | usa |","epistemico":"fato","exemplos":[],"id":"goldchain_mercado_hub_ponte_broca-so","memoria":"semantica","meta":{"dominio":"referencia","fonte":"web","nivel":5},"origem":"goldchain_mercado_hub.md","palavras_chave":[],"sinonimos":[],"tipo":"conceito","titulo":"Ponte Broca-SO"}
\section{Ponte Broca-SO}
| GoldChain | Nó | relação | |-----------|-----|---------| | núcleo | goldchain | confirma | | cadeia | blockchain | especializa | | valor | valor | depende | | moeda | dinheiro | depende | | hash | sha256 | usa | | cifra | cifra\_de\_ouro\_branco | referencia | | contabilidade | contabilidade | explica | | SO | broca\_os | referencia | | código | codigo\_goldchain\_fluxo | confirma | | PoW | gato | usa |
\prov{relações: parte\_de goldchain\_mercado\_hub\_goldchain\_hub\_do\_mercado\_broca; especializa blockchain; referencia cifra\_de\_ouro\_branco; explica contabilidade; referencia broca\_os; confirma codigo\_goldchain\_fluxo; usa gato}
\prov{proveniência: goldchain\_mercado\_hub.md \textperiodcentered{} web \textperiodcentered{} fato \textperiodcentered{} confiança 1.0 \textperiodcentered{} domínio referencia \textperiodcentered{} história 9 versões no jornal}

%CRISTAL {"arestas":[{"alvo":"goldchain_mercado_hub_goldchain_hub_do_mercado_broca","confianca":1.0,"epistemico":"fato","origem":"goldchain_mercado_hub.md","rel":"parte_de"}],"confianca":1.0,"contraexemplos":[],"descricao":"- **Código:** `goldchain/fluxo.py`, `goldchain/producao.py`, `goldchain/contrato.py` - **Paper:** `kernel/algebra_dos_contratos.tex` - **Manual:** `goldchain/README.md`, `goldchain/PAPER_CONTABILIDADE.md` - **Currículo:** fase mercado — valor → sha256 → GoldChain → Broca OS","epistemico":"fato","exemplos":[],"id":"goldchain_mercado_hub_referencia","memoria":"semantica","meta":{"dominio":"referencia","fonte":"web","nivel":5},"origem":"goldchain_mercado_hub.md","palavras_chave":[],"sinonimos":[],"tipo":"conceito","titulo":"Referência"}
\section{Referência}
- **Código:** `goldchain/fluxo.py`, `goldchain/producao.py`, `goldchain/contrato.py` - **Paper:** `kernel/algebra\_dos\_contratos.tex` - **Manual:** `goldchain/README.md`, `goldchain/PAPER\_CONTABILIDADE.md` - **Currículo:** fase mercado — valor \ensuremath{\to} sha256 \ensuremath{\to} GoldChain \ensuremath{\to} Broca OS
\prov{relações: parte\_de goldchain\_mercado\_hub\_goldchain\_hub\_do\_mercado\_broca}
\prov{proveniência: goldchain\_mercado\_hub.md \textperiodcentered{} web \textperiodcentered{} fato \textperiodcentered{} confiança 1.0 \textperiodcentered{} domínio referencia \textperiodcentered{} história 3 versões no jornal}

%CRISTAL {"arestas":[{"alvo":"o_conselho","confianca":1.0,"epistemico":"fato","origem":"wiki","rel":"parte_de"},{"alvo":"tiffany","confianca":1.0,"epistemico":"fato","origem":"wiki","rel":"explica"}],"confianca":1.0,"contraexemplos":[],"descricao":"As descobertas não são embeddings soltos, mas um grafo de arestas de proveniência — 'Aarão descobriu X; GPT formalizou; Grok refutou; o Conselho registrou X' — que preserva quem propôs, corrigiu e refutou cada resultado. É a própria estrutura da Universidade Tiffany.","epistemico":"inferencia","exemplos":[],"id":"grafo_de_proveniencia","memoria":"semantica","meta":{"autor":"Aarão/broca-so","dominio":"broca_repos","fonte":""},"origem":"wiki","palavras_chave":[],"sinonimos":[],"tipo":"conceito","titulo":"O Grafo de Proveniência (agente ≠ tecido)"}
\section{O Grafo de Proveniência (agente \ensuremath{\ne} tecido)}
As descobertas não são embeddings soltos, mas um grafo de arestas de proveniência — 'Aarão descobriu X; GPT formalizou; Grok refutou; o Conselho registrou X' — que preserva quem propôs, corrigiu e refutou cada resultado. É a própria estrutura da Universidade Tiffany.
\prov{relações: parte\_de o\_conselho; explica tiffany}
\prov{proveniência: wiki \textperiodcentered{} inferencia \textperiodcentered{} confiança 1.0 \textperiodcentered{} domínio broca\_repos \textperiodcentered{} história 3 versões no jornal}

%CRISTAL {"arestas":[{"alvo":"mult_prtsu","confianca":1.0,"epistemico":"fato","origem":"wiki","rel":"usa"},{"alvo":"torre_hurwitz","confianca":1.0,"epistemico":"fato","origem":"wiki","rel":"relaciona"}],"confianca":1.0,"contraexemplos":[],"descricao":"As trocas de sinal (p,r,t,s)→(±p,±r,±t,±s) com ε_p·ε_r=+1 formam ℤ₂³ (8 elementos) que preserva a norma; com o 'boost algébrico' contínuo (transformação de (p,r) que preserva o produto), dá o grupo de Lie ℤ₂³⋊ℝ. A órbita gera 8 representantes de divisão normada por classe de Hurwitz.","epistemico":"fato","exemplos":[],"id":"grupo_z2cubo_boost","memoria":"semantica","meta":{"autor":"Lopes/Aarão","dominio":"hiper","fonte":""},"origem":"wiki","palavras_chave":[],"sinonimos":[],"tipo":"conceito","titulo":"Grupo ℤ₂³⋊ℝ e o Boost Algébrico"}
\section{Grupo \ensuremath{\mathbb{Z}}\ensuremath{_{2}}\ensuremath{^{3}}\ensuremath{\rtimes}\ensuremath{\mathbb{R}} e o Boost Algébrico}
As trocas de sinal (p,r,t,s)\ensuremath{\to}($\pm$p,$\pm$r,$\pm$t,$\pm$s) com \ensuremath{\varepsilon}\_p\textperiodcentered \ensuremath{\varepsilon}\_r=+1 formam \ensuremath{\mathbb{Z}}\ensuremath{_{2}}\ensuremath{^{3}} (8 elementos) que preserva a norma; com o 'boost algébrico' contínuo (transformação de (p,r) que preserva o produto), dá o grupo de Lie \ensuremath{\mathbb{Z}}\ensuremath{_{2}}\ensuremath{^{3}}\ensuremath{\rtimes}\ensuremath{\mathbb{R}}. A órbita gera 8 representantes de divisão normada por classe de Hurwitz.
\prov{relações: usa mult\_prtsu; relaciona torre\_hurwitz}
\prov{proveniência: wiki \textperiodcentered{} fato \textperiodcentered{} confiança 1.0 \textperiodcentered{} domínio hiper \textperiodcentered{} história 4 versões no jornal}

%CRISTAL {"arestas":[{"alvo":"ds_universo","confianca":1.0,"epistemico":"fato","origem":"manual","rel":"parte_de"},{"alvo":"geometria","confianca":1.0,"epistemico":"fato","origem":"wiki","rel":"usa"},{"alvo":"complexos","confianca":1.0,"epistemico":"fato","origem":"wiki","rel":"usa"},{"alvo":"vetores","confianca":1.0,"epistemico":"fato","origem":"wiki","rel":"usa"},{"alvo":"ontologia","confianca":0.5,"epistemico":"fato","origem":"wiki","rel":"relaciona"},{"alvo":"imagem_do_floco_---_validacao_no","confianca":0.5,"epistemico":"fato","origem":"wiki","rel":"relaciona"},{"alvo":"word","confianca":0.5,"epistemico":"fato","origem":"wiki","rel":"relaciona"},{"alvo":"virtualizacao","confianca":0.5,"epistemico":"fato","origem":"wiki","rel":"relaciona"}],"confianca":1.0,"contraexemplos":[],"descricao":"Secção de DS_UNIVERSO.md.","epistemico":"fato","exemplos":[],"id":"habitantes_gentil_de_sitter","memoria":"semantica","meta":{"dominio":"manual","nivel":6,"verificacao":"slug idêntico — enriquecer, não duplicar"},"origem":"DS_UNIVERSO.md","palavras_chave":[],"sinonimos":[],"tipo":"conceito","titulo":"Habitantes → Lopes → de Sitter"}
\section{Habitantes \ensuremath{\to} Lopes \ensuremath{\to} de Sitter}
Secção de DS\_UNIVERSO.md.
\prov{relações: parte\_de ds\_universo; usa geometria; usa complexos; usa vetores; relaciona ontologia; relaciona imagem\_do\_floco\_---\_validacao\_no; relaciona word; relaciona virtualizacao}
\prov{proveniência: DS\_UNIVERSO.md \textperiodcentered{} fato \textperiodcentered{} confiança 1.0 \textperiodcentered{} domínio manual \textperiodcentered{} história 160 versões no jornal}

%CRISTAL {"arestas":[{"alvo":"koch_codifica_cantorfibonacci_prova_e_evidencia","confianca":1.0,"epistemico":"fato","origem":"KOCH_CANTOR_FIB_PROVA.md","rel":"parte_de"},{"alvo":"word","confianca":0.5,"epistemico":"fato","origem":"wiki","rel":"relaciona"},{"alvo":"virtualizacao","confianca":0.5,"epistemico":"fato","origem":"wiki","rel":"relaciona"},{"alvo":"while","confianca":0.5,"epistemico":"fato","origem":"wiki","rel":"relaciona"},{"alvo":"vontade_de_poder","confianca":0.5,"epistemico":"fato","origem":"wiki","rel":"relaciona"}],"confianca":1.0,"contraexemplos":[],"descricao":"O atlas Koch codifica o par **Cantor (real) + Fibonacci (ouro)** do parâmetro (s) sobre a parábola desdistorcida","epistemico":"fato","exemplos":[],"id":"hipotese_central","memoria":"semantica","meta":{"dominio":"manual","nivel":6,"verificacao":"slug idêntico — enriquecer, não duplicar"},"origem":"KOCH_CANTOR_FIB_PROVA.md","palavras_chave":[],"sinonimos":[],"tipo":"conceito","titulo":"Hipótese central"}
\section{Hipótese central}
O atlas Koch codifica o par **Cantor (real) + Fibonacci (ouro)** do parâmetro (s) sobre a parábola desdistorcida
\prov{relações: parte\_de koch\_codifica\_cantorfibonacci\_prova\_e\_evidencia; relaciona word; relaciona virtualizacao; relaciona while; relaciona vontade\_de\_poder}
\prov{proveniência: KOCH\_CANTOR\_FIB\_PROVA.md \textperiodcentered{} fato \textperiodcentered{} confiança 1.0 \textperiodcentered{} domínio manual \textperiodcentered{} história 9 versões no jornal}

%CRISTAL {"arestas":[{"alvo":"setor_algebrico_descontinuidades_no_fecho","confianca":1.0,"epistemico":"fato","origem":"POW_METAL_SETOR.md","rel":"parte_de"}],"confianca":1.0,"contraexemplos":[],"descricao":"O fecho ocorre quando a trajetória entra num setor algébrico diferente (divisores de zero / degenerado / complemento da representação principal), sem aproximação contínua no espaço observado.","epistemico":"fato","exemplos":[],"id":"hipotese_nao_assumida","memoria":"semantica","meta":{"dominio":"manual","nivel":6,"verificacao":"slug idêntico — enriquecer, não duplicar"},"origem":"POW_METAL_SETOR.md","palavras_chave":[],"sinonimos":[],"tipo":"conceito","titulo":"Hipótese (não assumida)"}
\section{Hipótese (não assumida)}
O fecho ocorre quando a trajetória entra num setor algébrico diferente (divisores de zero / degenerado / complemento da representação principal), sem aproximação contínua no espaço observado.
\prov{relações: parte\_de setor\_algebrico\_descontinuidades\_no\_fecho}
\prov{proveniência: POW\_METAL\_SETOR.md \textperiodcentered{} fato \textperiodcentered{} confiança 1.0 \textperiodcentered{} domínio manual \textperiodcentered{} história 4 versões no jornal}

%CRISTAL {"arestas":[{"alvo":"gato","confianca":0.95,"epistemico":"fato","origem":"paper","rel":"parte_de"}],"confianca":1.0,"contraexemplos":[],"descricao":"| Motivo | Contagem | |--------|----------| | `degenerado_real` | 0 | | `degenerado_ouro` | 0 | | `fora_corpo_gauss` | 0 | | `fora_lemniscata` | 0 | | `orbita_ouro` | 0 | | `fora_gauss` | 0 | | `fora_interface` | 0 | | `ref_fora_lemniscata` | 0 | | `fecho_nao` | 1,999,998 | | `hash_invalido` | 2 | | `outros` | 0 |","epistemico":"fato","exemplos":[],"id":"histograma_de_eliminacao","memoria":"semantica","meta":{"dominio":"manual","nivel":6,"verificacao":"slug idêntico — enriquecer, não duplicar"},"origem":"POW_METAL_FUNIL.md","palavras_chave":[],"sinonimos":[],"tipo":"conceito","titulo":"Histograma de eliminacao"}
\section{Histograma de eliminacao}
| Motivo | Contagem | |--------|----------| | `degenerado\_real` | 0 | | `degenerado\_ouro` | 0 | | `fora\_corpo\_gauss` | 0 | | `fora\_lemniscata` | 0 | | `orbita\_ouro` | 0 | | `fora\_gauss` | 0 | | `fora\_interface` | 0 | | `ref\_fora\_lemniscata` | 0 | | `fecho\_nao` | 1,999,998 | | `hash\_invalido` | 2 | | `outros` | 0 |
\prov{relações: parte\_de gato}
\prov{proveniência: POW\_METAL\_FUNIL.md \textperiodcentered{} fato \textperiodcentered{} confiança 1.0 \textperiodcentered{} domínio manual \textperiodcentered{} história 3 versões no jornal}

%CRISTAL {"arestas":[{"alvo":"continua_ultimo_face","confianca":1.0,"epistemico":"fato","origem":"wiki","rel":"especializa"}],"confianca":1.0,"contraexemplos":[],"descricao":"Além do ultimo.face, um historico.faces mantém as últimas MAX_HIST=4 respostas; _historico_bonus reaplica _continua (com peso pela metade) sobre as 3 faces mais recentes — continuidade além do turno imediatamente anterior.","epistemico":"fato","exemplos":[],"id":"historico_faces_janela","memoria":"semantica","meta":{"autor":"broca-so","dominio":"conversa","fonte":"conversa/colapso.py:load_historico/append_historico/_historico_bonus"},"origem":"wiki","palavras_chave":[],"sinonimos":[],"tipo":"conceito","titulo":"Tiffany — janela de histórico das últimas faces"}
\section{Tiffany — janela de histórico das últimas faces}
Além do ultimo.face, um historico.faces mantém as últimas MAX\_HIST=4 respostas; \_historico\_bonus reaplica \_continua (com peso pela metade) sobre as 3 faces mais recentes — continuidade além do turno imediatamente anterior.
\prov{relações: especializa continua\_ultimo\_face}
\prov{proveniência: wiki \textperiodcentered{} conversa/colapso.py:load\_historico/append\_historico/\_historico\_bonus \textperiodcentered{} fato \textperiodcentered{} confiança 1.0 \textperiodcentered{} domínio conversa \textperiodcentered{} história 10 versões no jornal}

%CRISTAL {"arestas":[{"alvo":"hopfield","confianca":0.95,"epistemico":"fato","origem":"wiki","rel":"confirma"},{"alvo":"rust","confianca":0.5,"epistemico":"fato","origem":"wiki","rel":"relaciona"},{"alvo":"por_sessao_conversadadostelemetriasessaojsonl","confianca":0.5,"epistemico":"fato","origem":"wiki","rel":"relaciona"},{"alvo":"misticismo_nao_dual","confianca":0.5,"epistemico":"fato","origem":"wiki","rel":"relaciona"},{"alvo":"virtualizacao","confianca":0.5,"epistemico":"fato","origem":"wiki","rel":"relaciona"},{"alvo":"word","confianca":0.5,"epistemico":"fato","origem":"wiki","rel":"relaciona"}],"confianca":0.95,"contraexemplos":[],"descricao":"Referência verificada: hopfield_rede_hub.md","epistemico":"fato","exemplos":[],"id":"hopfield_rede_hub","memoria":"semantica","meta":{"dominio":"referencia","fonte":"web","nivel":5},"origem":"wiki","palavras_chave":[],"sinonimos":[],"tipo":"conceito","titulo":"hopfield rede hub"}
\section{hopfield rede hub}
Referência verificada: hopfield\_rede\_hub.md
\prov{relações: confirma hopfield; relaciona rust; relaciona por\_sessao\_conversadadostelemetriasessaojsonl; relaciona misticismo\_nao\_dual; relaciona virtualizacao; relaciona word}
\prov{proveniência: wiki \textperiodcentered{} web \textperiodcentered{} fato \textperiodcentered{} confiança 0.95 \textperiodcentered{} domínio referencia \textperiodcentered{} história 8 versões no jornal}

%CRISTAL {"arestas":[{"alvo":"hopfield_rede_hub_hopfield_hub_da_rede_associativa_broca","confianca":1.0,"epistemico":"fato","origem":"hopfield_rede_hub.md","rel":"parte_de"},{"alvo":"virtualizacao","confianca":0.5,"epistemico":"fato","origem":"wiki","rel":"relaciona"},{"alvo":"word","confianca":0.5,"epistemico":"fato","origem":"wiki","rel":"relaciona"},{"alvo":"misticismo_nao_dual","confianca":0.5,"epistemico":"fato","origem":"wiki","rel":"relaciona"}],"confianca":1.0,"contraexemplos":[],"descricao":"**Hopfield ≠ backprop / deep learning** — memória associativa clássica, não SGD.","epistemico":"fato","exemplos":[],"id":"hopfield_rede_hub_contraexemplos","memoria":"semantica","meta":{"dominio":"referencia","fonte":"web","nivel":5},"origem":"hopfield_rede_hub.md","palavras_chave":[],"sinonimos":[],"tipo":"conceito","titulo":"Contraexemplos"}
\section{Contraexemplos}
**Hopfield \ensuremath{\ne} backprop / deep learning** — memória associativa clássica, não SGD.
\prov{relações: parte\_de hopfield\_rede\_hub\_hopfield\_hub\_da\_rede\_associativa\_broca; relaciona virtualizacao; relaciona word; relaciona misticismo\_nao\_dual}
\prov{proveniência: hopfield\_rede\_hub.md \textperiodcentered{} web \textperiodcentered{} fato \textperiodcentered{} confiança 1.0 \textperiodcentered{} domínio referencia \textperiodcentered{} história 6 versões no jornal}

%CRISTAL {"arestas":[{"alvo":"hopfield_rede_hub_hopfield_hub_da_rede_associativa_broca","confianca":1.0,"epistemico":"fato","origem":"hopfield_rede_hub.md","rel":"parte_de"},{"alvo":"virtualizacao","confianca":0.5,"epistemico":"fato","origem":"wiki","rel":"relaciona"},{"alvo":"word","confianca":0.5,"epistemico":"fato","origem":"wiki","rel":"relaciona"}],"confianca":1.0,"contraexemplos":[],"descricao":"- `papers/tiffany_cabeca_broca.md` - `neuronio/experimentos_tiffany.py` - currículo fase IA","epistemico":"fato","exemplos":[],"id":"hopfield_rede_hub_fontes","memoria":"semantica","meta":{"dominio":"referencia","fonte":"web","nivel":5},"origem":"hopfield_rede_hub.md","palavras_chave":[],"sinonimos":[],"tipo":"conceito","titulo":"Fontes"}
\section{Fontes}
- `papers/tiffany\_cabeca\_broca.md` - `neuronio/experimentos\_tiffany.py` - currículo fase IA
\prov{relações: parte\_de hopfield\_rede\_hub\_hopfield\_hub\_da\_rede\_associativa\_broca; relaciona virtualizacao; relaciona word}
\prov{proveniência: hopfield\_rede\_hub.md \textperiodcentered{} web \textperiodcentered{} fato \textperiodcentered{} confiança 1.0 \textperiodcentered{} domínio referencia \textperiodcentered{} história 5 versões no jornal}

%CRISTAL {"arestas":[{"alvo":"hopfield_rede_hub","confianca":0.95,"epistemico":"fato","origem":"wiki","rel":"parte_de"}],"confianca":1.0,"contraexemplos":[],"descricao":"> **Epistemologia:** fato histórico + implementação Broca (McCulloch-Pitts → Hopfield → Tiffany). > **No grafo:** `epistemico=fato`, confiança 0.95. > **Papel:** fechar cadeia IA antes da cabeça Tiffany.","epistemico":"fato","exemplos":[],"id":"hopfield_rede_hub_hopfield_hub_da_rede_associativa_broca","memoria":"semantica","meta":{"dominio":"referencia","fonte":"web","nivel":5},"origem":"hopfield_rede_hub.md","palavras_chave":[],"sinonimos":[],"tipo":"conceito","titulo":"Hopfield — hub da rede associativa Broca"}
\section{Hopfield — hub da rede associativa Broca}
> **Epistemologia:** fato histórico + implementação Broca (McCulloch-Pitts \ensuremath{\to} Hopfield \ensuremath{\to} Tiffany). > **No grafo:** `epistemico=fato`, confiança 0.95. > **Papel:** fechar cadeia IA antes da cabeça Tiffany.
\prov{relações: parte\_de hopfield\_rede\_hub}
\prov{proveniência: hopfield\_rede\_hub.md \textperiodcentered{} web \textperiodcentered{} fato \textperiodcentered{} confiança 1.0 \textperiodcentered{} domínio referencia \textperiodcentered{} história 3 versões no jornal}

%CRISTAL {"arestas":[{"alvo":"hopfield_rede_hub_hopfield_hub_da_rede_associativa_broca","confianca":1.0,"epistemico":"fato","origem":"hopfield_rede_hub.md","rel":"parte_de"},{"alvo":"mcculloch_pitts","confianca":0.95,"epistemico":"fato","origem":"wiki","rel":"depende"},{"alvo":"gato","confianca":0.95,"epistemico":"fato","origem":"wiki","rel":"depende"},{"alvo":"word","confianca":0.5,"epistemico":"fato","origem":"wiki","rel":"relaciona"}],"confianca":1.0,"contraexemplos":[],"descricao":"**Rede associativa** (Hopfield 1982): memória content-addressable por overlap/padrões.","epistemico":"fato","exemplos":[],"id":"hopfield_rede_hub_o_que_e","memoria":"semantica","meta":{"dominio":"referencia","fonte":"web","nivel":5},"origem":"hopfield_rede_hub.md","palavras_chave":[],"sinonimos":[],"tipo":"conceito","titulo":"O que é"}
\section{O que é}
**Rede associativa** (Hopfield 1982): memória content-addressable por overlap/padrões.
\prov{relações: parte\_de hopfield\_rede\_hub\_hopfield\_hub\_da\_rede\_associativa\_broca; depende mcculloch\_pitts; depende gato; relaciona word}
\prov{proveniência: hopfield\_rede\_hub.md \textperiodcentered{} web \textperiodcentered{} fato \textperiodcentered{} confiança 1.0 \textperiodcentered{} domínio referencia \textperiodcentered{} história 6 versões no jornal}

%CRISTAL {"arestas":[{"alvo":"hopfield_rede_hub_hopfield_hub_da_rede_associativa_broca","confianca":1.0,"epistemico":"fato","origem":"hopfield_rede_hub.md","rel":"parte_de"},{"alvo":"tiffany","confianca":0.95,"epistemico":"fato","origem":"wiki","rel":"referencia"},{"alvo":"broca_os","confianca":0.95,"epistemico":"fato","origem":"wiki","rel":"explica"},{"alvo":"inteligencia_artificial","confianca":0.95,"epistemico":"fato","origem":"wiki","rel":"parte_de"},{"alvo":"codigo_neuronio_experimentos_tiffany","confianca":0.95,"epistemico":"fato","origem":"wiki","rel":"referencia"}],"confianca":1.0,"contraexemplos":[],"descricao":"| Hopfield | Nó | relação | |----------|-----|---------| | rede | hopfield | confirma | | predecessor | mcculloch_pitts | depende | | operador base | gato | depende | | sucessor | tiffany | especializa | | SO fractal | broca_os | explica | | domínio | inteligencia_artificial | parte_de | | experimentos | codigo_neuronio_experimentos_tiffany | referencia |","epistemico":"fato","exemplos":[],"id":"hopfield_rede_hub_ponte_broca-so","memoria":"semantica","meta":{"dominio":"referencia","fonte":"web","nivel":5},"origem":"hopfield_rede_hub.md","palavras_chave":[],"sinonimos":[],"tipo":"conceito","titulo":"Ponte Broca-SO"}
\section{Ponte Broca-SO}
| Hopfield | Nó | relação | |----------|-----|---------| | rede | hopfield | confirma | | predecessor | mcculloch\_pitts | depende | | operador base | gato | depende | | sucessor | tiffany | especializa | | SO fractal | broca\_os | explica | | domínio | inteligencia\_artificial | parte\_de | | experimentos | codigo\_neuronio\_experimentos\_tiffany | referencia |
\prov{relações: parte\_de hopfield\_rede\_hub\_hopfield\_hub\_da\_rede\_associativa\_broca; referencia tiffany; explica broca\_os; parte\_de inteligencia\_artificial; referencia codigo\_neuronio\_experimentos\_tiffany}
\prov{proveniência: hopfield\_rede\_hub.md \textperiodcentered{} web \textperiodcentered{} fato \textperiodcentered{} confiança 1.0 \textperiodcentered{} domínio referencia \textperiodcentered{} história 7 versões no jornal}

%CRISTAL {"arestas":[{"alvo":"hopfield_rede_hub_hopfield_hub_da_rede_associativa_broca","confianca":1.0,"epistemico":"fato","origem":"hopfield_rede_hub.md","rel":"parte_de"}],"confianca":1.0,"contraexemplos":[],"descricao":"- **Currículo:** fase IA — `mcculloch_pitts` → **hopfield** → `tiffany` - **Paper:** `papers/Tiffany.tex`, `papers/computacao_fundamentos_broca.tex` - **Código:** `neuronio/experimentos_tiffany.py` — memória associativa sobre tecido - **Hub Tiffany:** `papers/tiffany_cabeca_broca.md`","epistemico":"fato","exemplos":[],"id":"hopfield_rede_hub_referencia","memoria":"semantica","meta":{"dominio":"referencia","fonte":"web","nivel":5},"origem":"hopfield_rede_hub.md","palavras_chave":[],"sinonimos":[],"tipo":"conceito","titulo":"Referência"}
\section{Referência}
- **Currículo:** fase IA — `mcculloch\_pitts` \ensuremath{\to} **hopfield** \ensuremath{\to} `tiffany` - **Paper:** `papers/Tiffany.tex`, `papers/computacao\_fundamentos\_broca.tex` - **Código:** `neuronio/experimentos\_tiffany.py` — memória associativa sobre tecido - **Hub Tiffany:** `papers/tiffany\_cabeca\_broca.md`
\prov{relações: parte\_de hopfield\_rede\_hub\_hopfield\_hub\_da\_rede\_associativa\_broca}
\prov{proveniência: hopfield\_rede\_hub.md \textperiodcentered{} web \textperiodcentered{} fato \textperiodcentered{} confiança 1.0 \textperiodcentered{} domínio referencia \textperiodcentered{} história 3 versões no jornal}

%CRISTAL {"arestas":[{"alvo":"morte","confianca":1.0,"epistemico":"fato","origem":"paper","rel":"parte_de"},{"alvo":"vida","confianca":1.0,"epistemico":"fato","origem":"paper","rel":"parte_de"},{"alvo":"maquina_de_turing","confianca":1.0,"epistemico":"fato","origem":"paper","rel":"analogia"}],"confianca":1.0,"contraexemplos":[],"descricao":"| metal | n | ε ruptura (1º deg>2 ou Δ>1e-6) | |-------|---|--------------------------------| | ouro | 1 | 200 | | prata | 2 | 2000 | | bronze | 3 | 0 | | cobre | 4 | 0 | | estanho | 5 | 0 | | metal | 6 | 0 | | metal | 7 | 0 | | metal | 8 | 0 |","epistemico":"fato","exemplos":[],"id":"horizonte","memoria":"semantica","meta":{"dominio":"manual","duplicata_sugerida":[],"nivel":6,"verificacao":"parte de conceito mais amplo 'horizonte'"},"origem":"RUÍDO.md","palavras_chave":[],"sinonimos":[],"tipo":"conceito","titulo":"Horizonte"}
\section{Horizonte}
| metal | n | \ensuremath{\varepsilon} ruptura (1º deg>2 ou \ensuremath{\Delta}>1e-6) | |-------|---|--------------------------------| | ouro | 1 | 200 | | prata | 2 | 2000 | | bronze | 3 | 0 | | cobre | 4 | 0 | | estanho | 5 | 0 | | metal | 6 | 0 | | metal | 7 | 0 | | metal | 8 | 0 |
\prov{relações: parte\_de morte; parte\_de vida; analogia maquina\_de\_turing}
\prov{proveniência: RUÍDO.md \textperiodcentered{} fato \textperiodcentered{} confiança 1.0 \textperiodcentered{} domínio manual \textperiodcentered{} história 42 versões no jornal}

%CRISTAL {"arestas":[{"alvo":"mult_prtsu","confianca":1.0,"epistemico":"fato","origem":"wiki","rel":"usa"},{"alvo":"nco_hipercomplexo_maxcut","confianca":1.0,"epistemico":"fato","origem":"wiki","rel":"relaciona"},{"alvo":"assinatura_algebrica_financeira","confianca":1.0,"epistemico":"fato","origem":"wiki","rel":"usa"}],"confianca":1.0,"contraexemplos":[],"descricao":"Um MLP com camadas HyperLinear (botões PRTSUV por neurônio, regime fermiônico) supera a RealMLP em 6/7 ativos usando 1/3 dos parâmetros; os botões core convergem para ~(1,1,1) com exóticos ~0 — a rede 'escolhe' a álgebra base, resgatando os ativos não-lineares (BTC, USD/BRL, ouro) onde a regressão linear falha.","epistemico":"fato","exemplos":[],"id":"hypermlp_escolhe_algebra","memoria":"semantica","meta":{"autor":"Aarão/broca-so","dominio":"broca_repos","fonte":""},"origem":"wiki","palavras_chave":[],"sinonimos":[],"tipo":"conceito","titulo":"HyperMLP Converge para a Álgebra Base"}
\section{HyperMLP Converge para a Álgebra Base}
Um MLP com camadas HyperLinear (botões PRTSUV por neurônio, regime fermiônico) supera a RealMLP em 6/7 ativos usando 1/3 dos parâmetros; os botões core convergem para \~{}(1,1,1) com exóticos \~{}0 — a rede 'escolhe' a álgebra base, resgatando os ativos não-lineares (BTC, USD/BRL, ouro) onde a regressão linear falha.
\prov{relações: usa mult\_prtsu; relaciona nco\_hipercomplexo\_maxcut; usa assinatura\_algebrica\_financeira}
\prov{proveniência: wiki \textperiodcentered{} fato \textperiodcentered{} confiança 1.0 \textperiodcentered{} domínio broca\_repos \textperiodcentered{} história 4 versões no jornal}

%CRISTAL {"arestas":[{"alvo":"funcao_onda_secao_hopf","confianca":1.0,"epistemico":"fato","origem":"wiki","rel":"usa"},{"alvo":"colapso_observador","confianca":1.0,"epistemico":"fato","origem":"wiki","rel":"explica"},{"alvo":"corpo_dual","confianca":1.0,"epistemico":"fato","origem":"wiki","rel":"relaciona"}],"confianca":1.0,"contraexemplos":[],"descricao":"Cada versor imaginário â∈S² traz sua própria unidade imaginária (ℂâ=span1,â≅ℂ); a 'i' universal da QM é só a escolha de referência global. Evolução = transporte paralelo; fase de Berry = holonomia no O(−1); 'colapso na base â' = ler o estado na carta local, não um processo físico distinto.","epistemico":"inferencia","exemplos":[],"id":"i_de_schrodinger_carta","memoria":"semantica","meta":{"autor":"Lopes/Aarão","dominio":"hiper","fonte":""},"origem":"wiki","palavras_chave":[],"sinonimos":[],"tipo":"conceito","titulo":"A 'i' de Schrödinger é uma Carta Local"}
\section{A 'i' de Schrödinger é uma Carta Local}
Cada versor imaginário â\ensuremath{\in}S\ensuremath{^{2}} traz sua própria unidade imaginária (\ensuremath{\mathbb{C}}â=span1,â\ensuremath{\cong}\ensuremath{\mathbb{C}}); a 'i' universal da QM é só a escolha de referência global. Evolução = transporte paralelo; fase de Berry = holonomia no O(\ensuremath{-}1); 'colapso na base â' = ler o estado na carta local, não um processo físico distinto.
\prov{relações: usa funcao\_onda\_secao\_hopf; explica colapso\_observador; relaciona corpo\_dual}
\prov{proveniência: wiki \textperiodcentered{} inferencia \textperiodcentered{} confiança 1.0 \textperiodcentered{} domínio hiper \textperiodcentered{} história 6 versões no jornal}

%CRISTAL {"arestas":[{"alvo":"koch_parabola_v2_sobre_a_curva_desdistorcida","confianca":1.0,"epistemico":"fato","origem":"KOCH_PARABOLA.md","rel":"parte_de"}],"confianca":1.0,"contraexemplos":[],"descricao":"1. **Marcar pontos na parábola** — cada Word dá `(c², calor)` com `calor = 1 − c²`. 2. **Desdistorcer parábola → reta** — coordenada `s = arco(c) / L` (comprimento de arco normalizado). 3. **Lado real** — Cantor ternário em `s` (sem dígito 1, base 3). 4. **Lado ouro** — Cantor φ (expansão β sem `11`) + convergentes Fibonacci `Fₙ+₁/Fₙ`. 5. **AGM** — acelera o comprimento; ganho quadrático, **nunca exacto** (atrator projetado).","epistemico":"fato","exemplos":[],"id":"ideia","memoria":"semantica","meta":{"dominio":"manual","nivel":6,"verificacao":"slug idêntico — enriquecer, não duplicar"},"origem":"KOCH_PARABOLA.md","palavras_chave":[],"sinonimos":[],"tipo":"conceito","titulo":"Ideia"}
\section{Ideia}
1. **Marcar pontos na parábola** — cada Word dá `(c\ensuremath{^{2}}, calor)` com `calor = 1 \ensuremath{-} c\ensuremath{^{2}}`. 2. **Desdistorcer parábola \ensuremath{\to} reta** — coordenada `s = arco(c) / L` (comprimento de arco normalizado). 3. **Lado real** — Cantor ternário em `s` (sem dígito 1, base 3). 4. **Lado ouro** — Cantor \ensuremath{\varphi} (expansão \ensuremath{\beta} sem `11`) + convergentes Fibonacci `F\ensuremath{_{n}}+\ensuremath{_{1}}/F\ensuremath{_{n}}`. 5. **AGM** — acelera o comprimento; ganho quadrático, **nunca exacto** (atrator projetado).
\prov{relações: parte\_de koch\_parabola\_v2\_sobre\_a\_curva\_desdistorcida}
\prov{proveniência: KOCH\_PARABOLA.md \textperiodcentered{} fato \textperiodcentered{} confiança 1.0 \textperiodcentered{} domínio manual \textperiodcentered{} história 5 versões no jornal}

%CRISTAL {"arestas":[{"alvo":"cifra","confianca":1.0,"epistemico":"fato","origem":"wiki","rel":"relaciona"},{"alvo":"identidade_obra_holografica","confianca":1.0,"epistemico":"fato","origem":"wiki","rel":"relaciona"}],"confianca":1.0,"contraexemplos":[],"descricao":"A identidade intrínseca de um tecido não é o ponto v na reta de ouro (que carrega o frame de quem traçou a régua), mas w(v)=(v−φ)/(v−ψ), a coordenada de Möbius nos autovetores φ,ψ: invariante afim, no círculo projetivo ℙ¹. Identificar por v é forjável (1000/1000 personificações aceitas por v; 0/1000 por w) — é integridade de endereçamento, não confidencialidade.","epistemico":"fato","exemplos":[],"id":"identidade_invariante_w","memoria":"semantica","meta":{"autor":"Aarão/broca-so","dominio":"broca_repos","fonte":""},"origem":"wiki","palavras_chave":[],"sinonimos":[],"tipo":"conceito","titulo":"A Identidade Invariante w (o observador-régua morre)"}
\section{A Identidade Invariante w (o observador-régua morre)}
A identidade intrínseca de um tecido não é o ponto v na reta de ouro (que carrega o frame de quem traçou a régua), mas w(v)=(v\ensuremath{-}\ensuremath{\varphi})/(v\ensuremath{-}\ensuremath{\psi}), a coordenada de Möbius nos autovetores \ensuremath{\varphi},\ensuremath{\psi}: invariante afim, no círculo projetivo \ensuremath{\mathbb{P}}\ensuremath{^{1}}. Identificar por v é forjável (1000/1000 personificações aceitas por v; 0/1000 por w) — é integridade de endereçamento, não confidencialidade.
\prov{relações: relaciona cifra; relaciona identidade\_obra\_holografica}
\prov{proveniência: wiki \textperiodcentered{} fato \textperiodcentered{} confiança 1.0 \textperiodcentered{} domínio broca\_repos \textperiodcentered{} história 3 versões no jornal}

%CRISTAL {"arestas":[{"alvo":"olho_de_koch","confianca":1.0,"epistemico":"fato","origem":"wiki","rel":"usa"},{"alvo":"colapso_koch_tiffany","confianca":1.0,"epistemico":"fato","origem":"wiki","rel":"relaciona"}],"confianca":1.0,"contraexemplos":[],"descricao":"A noção certa é identidade estrutural (Id(A)≈Id(B), 'perdi um braço, continuo eu'), não igualdade bit-a-bit. O tecido é holográfico (cada fragmento carrega o todo: 25% da obra → 0,999) e o invariante real é topológico — a classe em GL(2,ℤ) na homologia H₁(T²); lente/buraco-negro são conjugações que não movem a classe. Habilita detectar pirataria pela alma.","epistemico":"fato","exemplos":[],"id":"identidade_obra_holografica","memoria":"semantica","meta":{"autor":"Aarão/broca-so","dominio":"broca_repos","fonte":""},"origem":"wiki","palavras_chave":[],"sinonimos":[],"tipo":"conceito","titulo":"A Identidade da Obra é Holográfica e Topológica"}
\section{A Identidade da Obra é Holográfica e Topológica}
A noção certa é identidade estrutural (Id(A)\ensuremath{\approx}Id(B), 'perdi um braço, continuo eu'), não igualdade bit-a-bit. O tecido é holográfico (cada fragmento carrega o todo: 25\% da obra \ensuremath{\to} 0,999) e o invariante real é topológico — a classe em GL(2,\ensuremath{\mathbb{Z}}) na homologia H\ensuremath{_{1}}(T\ensuremath{^{2}}); lente/buraco-negro são conjugações que não movem a classe. Habilita detectar pirataria pela alma.
\prov{relações: usa olho\_de\_koch; relaciona colapso\_koch\_tiffany}
\prov{proveniência: wiki \textperiodcentered{} fato \textperiodcentered{} confiança 1.0 \textperiodcentered{} domínio broca\_repos \textperiodcentered{} história 3 versões no jornal}

%CRISTAL {"arestas":[{"alvo":"maquina_algebrica","confianca":1.0,"epistemico":"fato","origem":"wiki","rel":"usa"},{"alvo":"cordas_ilha_hurwitz","confianca":1.0,"epistemico":"fato","origem":"wiki","rel":"relaciona"}],"confianca":1.0,"contraexemplos":[],"descricao":"A auto-contração da 4-forma de Cayley canônica em ℝ⁸ (estabilizador Spin(7)) fecha exatamente em duas peças irredutíveis Spin(7)-invariantes, com resíduo zero; a constante calibrada b=4 coincide com a de ℝ⁷, a métrica a=6 depende da dimensão. O maquinário algébrico estendido a ℝ⁸/𝕆.","epistemico":"fato","exemplos":[],"id":"identidade_omega_r8","memoria":"semantica","meta":{"autor":"Aarão/broca-so","dominio":"broca_repos","fonte":""},"origem":"wiki","palavras_chave":[],"sinonimos":[],"tipo":"conceito","titulo":"A Identidade Ω·Ω = 6K + 4Ω (ℝ⁸)"}
\section{A Identidade \ensuremath{\Omega}\textperiodcentered \ensuremath{\Omega} = 6K + 4\ensuremath{\Omega} (\ensuremath{\mathbb{R}}\ensuremath{^{8}})}
A auto-contração da 4-forma de Cayley canônica em \ensuremath{\mathbb{R}}\ensuremath{^{8}} (estabilizador Spin(7)) fecha exatamente em duas peças irredutíveis Spin(7)-invariantes, com resíduo zero; a constante calibrada b=4 coincide com a de \ensuremath{\mathbb{R}}\ensuremath{^{7}}, a métrica a=6 depende da dimensão. O maquinário algébrico estendido a \ensuremath{\mathbb{R}}\ensuremath{^{8}}/\ensuremath{\mathbb{O}}.
\prov{relações: usa maquina\_algebrica; relaciona cordas\_ilha\_hurwitz}
\prov{proveniência: wiki \textperiodcentered{} fato \textperiodcentered{} confiança 1.0 \textperiodcentered{} domínio broca\_repos \textperiodcentered{} história 3 versões no jornal}

%CRISTAL {"arestas":[{"alvo":"ontologia_da_multiplicacao","confianca":1.0,"epistemico":"fato","origem":"wiki","rel":"usa"},{"alvo":"cordas_ilha_hurwitz","confianca":1.0,"epistemico":"fato","origem":"wiki","rel":"relaciona"}],"confianca":1.0,"contraexemplos":[],"descricao":"No espaço de bilineares (dim 27 sob SO(3)), cada ponto físico (positividade+unitariedade+causalidade) é um 'mundo possível', e o domínio físico tem medida <3% (teorema da pequenez). A novidade central: o setor spin-3, 7-dimensional, ausente da classificação de Hurwitz, com aplicações em oscilações quirais e modificação de Navier-Stokes.","epistemico":"inferencia","exemplos":[],"id":"ilha_fractal_spin3","memoria":"semantica","meta":{"autor":"Aarão/broca-so","dominio":"broca_repos","fonte":""},"origem":"wiki","palavras_chave":[],"sinonimos":[],"tipo":"conceito","titulo":"A Ilha Fractal e o Setor Spin-3"}
\section{A Ilha Fractal e o Setor Spin-3}
No espaço de bilineares (dim 27 sob SO(3)), cada ponto físico (positividade+unitariedade+causalidade) é um 'mundo possível', e o domínio físico tem medida <3\% (teorema da pequenez). A novidade central: o setor spin-3, 7-dimensional, ausente da classificação de Hurwitz, com aplicações em oscilações quirais e modificação de Navier-Stokes.
\prov{relações: usa ontologia\_da\_multiplicacao; relaciona cordas\_ilha\_hurwitz}
\prov{proveniência: wiki \textperiodcentered{} inferencia \textperiodcentered{} confiança 1.0 \textperiodcentered{} domínio broca\_repos \textperiodcentered{} história 3 versões no jornal}

%CRISTAL {"arestas":[{"alvo":"floco_de_neve","confianca":1.0,"epistemico":"fato","origem":"paper","rel":"parte_de"},{"alvo":"mecanica_quantica_hub","confianca":0.5,"epistemico":"fato","origem":"wiki","rel":"relaciona"},{"alvo":"memoria_virtual","confianca":0.5,"epistemico":"fato","origem":"wiki","rel":"relaciona"},{"alvo":"teoria_do_valor","confianca":0.5,"epistemico":"fato","origem":"wiki","rel":"relaciona"},{"alvo":"utilitarismo","confianca":0.5,"epistemico":"fato","origem":"wiki","rel":"relaciona"}],"confianca":1.0,"contraexemplos":[],"descricao":"assert not tiffany.validate(isolation=True) ```","epistemico":"fato","exemplos":[],"id":"imagem_do_floco_---_validacao_no","memoria":"semantica","meta":{"dominio":"manual","duplicata_sugerida":[],"nivel":6,"verificacao":"parte de conceito mais amplo 'imagem_do_floco_---_validacao_no'"},"origem":"INSTALL.md","palavras_chave":[],"sinonimos":[],"tipo":"conceito","titulo":"Imagem do floco --- validação no"}
\section{Imagem do floco --- validação no}
assert not tiffany.validate(isolation=True) ```
\prov{relações: parte\_de floco\_de\_neve; relaciona mecanica\_quantica\_hub; relaciona memoria\_virtual; relaciona teoria\_do\_valor; relaciona utilitarismo}
\prov{proveniência: INSTALL.md \textperiodcentered{} fato \textperiodcentered{} confiança 1.0 \textperiodcentered{} domínio manual \textperiodcentered{} história 8 versões no jornal}

%CRISTAL {"arestas":[{"alvo":"mult_prtsu","confianca":1.0,"epistemico":"fato","origem":"wiki","rel":"usa"},{"alvo":"invariante_estelar","confianca":1.0,"epistemico":"fato","origem":"wiki","rel":"relaciona"},{"alvo":"prova_de_trabalho","confianca":1.0,"epistemico":"fato","origem":"wiki","rel":"relaciona"}],"confianca":1.0,"contraexemplos":[],"descricao":"‖z·w‖²=‖z‖²‖w‖²−(1−p²)(a₀b₀)²−(1−r²)D²−(1−t²)Q−(1−s²)S+2(t²−pr)a₀b₀D; o défice é o 'imposto' que a álgebra cobra por desviar-se das álgebras de divisão. Análogo estrutural do custo da prova de trabalho.","epistemico":"fato","exemplos":[],"id":"imposto_algebrico","memoria":"semantica","meta":{"autor":"Lopes/Aarão","dominio":"hiper","fonte":""},"origem":"wiki","palavras_chave":[],"sinonimos":[],"tipo":"conceito","titulo":"O Imposto Algébrico (custo de norma)"}
\section{O Imposto Algébrico (custo de norma)}
\ensuremath{\|}z\textperiodcentered w\ensuremath{\|}\ensuremath{^{2}}=\ensuremath{\|}z\ensuremath{\|}\ensuremath{^{2}}\ensuremath{\|}w\ensuremath{\|}\ensuremath{^{2}}\ensuremath{-}(1\ensuremath{-}p\ensuremath{^{2}})(a\ensuremath{_{0}}b\ensuremath{_{0}})\ensuremath{^{2}}\ensuremath{-}(1\ensuremath{-}r\ensuremath{^{2}})D\ensuremath{^{2}}\ensuremath{-}(1\ensuremath{-}t\ensuremath{^{2}})Q\ensuremath{-}(1\ensuremath{-}s\ensuremath{^{2}})S+2(t\ensuremath{^{2}}\ensuremath{-}pr)a\ensuremath{_{0}}b\ensuremath{_{0}}D; o défice é o 'imposto' que a álgebra cobra por desviar-se das álgebras de divisão. Análogo estrutural do custo da prova de trabalho.
\prov{relações: usa mult\_prtsu; relaciona invariante\_estelar; relaciona prova\_de\_trabalho}
\prov{proveniência: wiki \textperiodcentered{} fato \textperiodcentered{} confiança 1.0 \textperiodcentered{} domínio hiper \textperiodcentered{} história 5 versões no jornal}

%CRISTAL {"arestas":[{"alvo":"validacao_hipotese_h_ressonancia","confianca":1.0,"epistemico":"fato","origem":"POW_METAL_RESSONANCIA_VAL.md","rel":"parte_de"},{"alvo":"respiracao_dos_programas","confianca":0.5,"epistemico":"fato","origem":"wiki","rel":"relaciona"},{"alvo":"pell","confianca":0.5,"epistemico":"fato","origem":"wiki","rel":"relaciona"},{"alvo":"word","confianca":0.5,"epistemico":"fato","origem":"wiki","rel":"relaciona"},{"alvo":"virtualizacao","confianca":0.5,"epistemico":"fato","origem":"wiki","rel":"relaciona"},{"alvo":"vontade_de_poder","confianca":0.5,"epistemico":"fato","origem":"wiki","rel":"relaciona"}],"confianca":1.0,"contraexemplos":[],"descricao":"Secção de POW_METAL_RESSONANCIA_VAL.md.","epistemico":"fato","exemplos":[],"id":"inspecao_local_32","memoria":"semantica","meta":{"dominio":"manual","nivel":6,"verificacao":"slug idêntico — enriquecer, não duplicar"},"origem":"POW_METAL_RESSONANCIA_VAL.md","palavras_chave":[],"sinonimos":[],"tipo":"conceito","titulo":"Inspeção local (±32)"}
\section{Inspeção local ($\pm$32)}
Secção de POW\_METAL\_RESSONANCIA\_VAL.md.
\prov{relações: parte\_de validacao\_hipotese\_h\_ressonancia; relaciona respiracao\_dos\_programas; relaciona pell; relaciona word; relaciona virtualizacao; relaciona vontade\_de\_poder}
\prov{proveniência: POW\_METAL\_RESSONANCIA\_VAL.md \textperiodcentered{} fato \textperiodcentered{} confiança 1.0 \textperiodcentered{} domínio manual \textperiodcentered{} história 9 versões no jornal}

%CRISTAL {"arestas":[{"alvo":"broca_os","confianca":1.0,"epistemico":"fato","origem":"manual","rel":"referencia"},{"alvo":"transcricao","confianca":0.5,"epistemico":"fato","origem":"wiki","rel":"relaciona"},{"alvo":"prova_isa","confianca":0.5,"epistemico":"fato","origem":"wiki","rel":"relaciona"},{"alvo":"pow_gato","confianca":0.5,"epistemico":"fato","origem":"wiki","rel":"relaciona"},{"alvo":"montagem","confianca":0.5,"epistemico":"fato","origem":"wiki","rel":"relaciona"},{"alvo":"teste_ordem","confianca":0.5,"epistemico":"fato","origem":"wiki","rel":"relaciona"},{"alvo":"readme","confianca":0.5,"epistemico":"fato","origem":"wiki","rel":"relaciona"},{"alvo":"sessao_interativa","confianca":0.5,"epistemico":"fato","origem":"wiki","rel":"relaciona"},{"alvo":"main","confianca":0.5,"epistemico":"fato","origem":"wiki","rel":"relaciona"},{"alvo":"koch_pell","confianca":0.5,"epistemico":"fato","origem":"wiki","rel":"relaciona"},{"alvo":"mamifero","confianca":0.5,"epistemico":"fato","origem":"wiki","rel":"relaciona"},{"alvo":"pow_espectral","confianca":0.5,"epistemico":"fato","origem":"wiki","rel":"relaciona"},{"alvo":"kernel","confianca":0.5,"epistemico":"fato","origem":"wiki","rel":"relaciona"}],"confianca":1.0,"contraexemplos":[],"descricao":"Manual: INSTALL.md","epistemico":"fato","exemplos":[],"id":"install","memoria":"semantica","meta":{"arquivo":"INSTALL.md","dominio":"manual","nivel":6},"origem":"manual","palavras_chave":[],"sinonimos":[],"tipo":"conceito","titulo":"INSTALL"}
\section{INSTALL}
Manual: INSTALL.md
\prov{relações: referencia broca\_os; relaciona transcricao; relaciona prova\_isa; relaciona pow\_gato; relaciona montagem; relaciona teste\_ordem; relaciona readme; relaciona sessao\_interativa; relaciona main; relaciona koch\_pell; relaciona mamifero; relaciona pow\_espectral; relaciona kernel}
\prov{proveniência: manual \textperiodcentered{} fato \textperiodcentered{} confiança 1.0 \textperiodcentered{} domínio manual \textperiodcentered{} história 16 versões no jornal}

%CRISTAL {"arestas":[{"alvo":"gato","confianca":0.95,"epistemico":"fato","origem":"paper","rel":"parte_de"}],"confianca":1.0,"contraexemplos":[],"descricao":"| fecho | pre-nonce | acopl | fase | atrator | Δatr | koch | lemn_erro | |-------|-----------|-------|------|---------|------|------|-----------| | 1572863 | 1572862 | 0.6022 | 0.1293 | 520 | 152 | 0.0858 | 0.00000 | | 1966079 | 1966078 | 0.6027 | 0.1293 | 520 | 152 | 0.0858 | 0.00000 |","epistemico":"fato","exemplos":[],"id":"instante_anterior_ao_fecho_nonce1","memoria":"semantica","meta":{"dominio":"manual","nivel":6,"verificacao":"slug idêntico — enriquecer, não duplicar"},"origem":"POW_METAL_PRE_FECHO.md","palavras_chave":[],"sinonimos":[],"tipo":"conceito","titulo":"Instante anterior ao fecho (nonce−1)"}
\section{Instante anterior ao fecho (nonce\ensuremath{-}1)}
| fecho | pre-nonce | acopl | fase | atrator | \ensuremath{\Delta}atr | koch | lemn\_erro | |-------|-----------|-------|------|---------|------|------|-----------| | 1572863 | 1572862 | 0.6022 | 0.1293 | 520 | 152 | 0.0858 | 0.00000 | | 1966079 | 1966078 | 0.6027 | 0.1293 | 520 | 152 | 0.0858 | 0.00000 |
\prov{relações: parte\_de gato}
\prov{proveniência: POW\_METAL\_PRE\_FECHO.md \textperiodcentered{} fato \textperiodcentered{} confiança 1.0 \textperiodcentered{} domínio manual \textperiodcentered{} história 3 versões no jornal}

%CRISTAL {"arestas":[{"alvo":"cifra_ouro_branco_nucleo_espectral","confianca":1.0,"epistemico":"fato","origem":"CIFRA.md","rel":"parte_de"}],"confianca":1.0,"contraexemplos":[],"descricao":"| Opcode | Operando | Efeito | |--------|----------|--------| | `SILVER` | A | `A ← A₂·A`, `R ← A` | | `GOLD` | A | `A ← A₁·A`, `R ← A` | | `BRONZE` | A | `A ← A₃·A`, `R ← A` | | `UNFOLD` | A | `A ← gold(A)`, `R ← A` | | `PROJECT` | A, B | `A ← project(A,B)`, `R ← A` | | `FOLD` | A, B | `A ← fold(A,B)`, `R ← A` | | `LIFT` | A, B | `A ← lift(A,B)`, `R ← A` | | `LOADS` | slot | neurônio → `A₂·word` → A (cifra implícita) |","epistemico":"fato","exemplos":[],"id":"instrucoes_espectrais","memoria":"semantica","meta":{"dominio":"manual","nivel":6,"verificacao":"slug idêntico — enriquecer, não duplicar"},"origem":"CIFRA.md","palavras_chave":[],"sinonimos":[],"tipo":"conceito","titulo":"Instruções espectrais"}
\section{Instruções espectrais}
| Opcode | Operando | Efeito | |--------|----------|--------| | `SILVER` | A | `A \ensuremath{\leftarrow} A\ensuremath{_{2}}\textperiodcentered A`, `R \ensuremath{\leftarrow} A` | | `GOLD` | A | `A \ensuremath{\leftarrow} A\ensuremath{_{1}}\textperiodcentered A`, `R \ensuremath{\leftarrow} A` | | `BRONZE` | A | `A \ensuremath{\leftarrow} A\ensuremath{_{3}}\textperiodcentered A`, `R \ensuremath{\leftarrow} A` | | `UNFOLD` | A | `A \ensuremath{\leftarrow} gold(A)`, `R \ensuremath{\leftarrow} A` | | `PROJECT` | A, B | `A \ensuremath{\leftarrow} project(A,B)`, `R \ensuremath{\leftarrow} A` | | `FOLD` | A, B | `A \ensuremath{\leftarrow} fold(A,B)`, `R \ensuremath{\leftarrow} A` | | `LIFT` | A, B | `A \ensuremath{\leftarrow} lift(A,B)`, `R \ensuremath{\leftarrow} A` | | `LOADS` | slot | neurônio \ensuremath{\to} `A\ensuremath{_{2}}\textperiodcentered word` \ensuremath{\to} A (cifra implícita) |
\prov{relações: parte\_de cifra\_ouro\_branco\_nucleo\_espectral}
\prov{proveniência: CIFRA.md \textperiodcentered{} fato \textperiodcentered{} confiança 1.0 \textperiodcentered{} domínio manual \textperiodcentered{} história 6 versões no jornal}

%CRISTAL {"arestas":[{"alvo":"casl-propagation","confianca":0.094,"epistemico":"fato","origem":"manual","rel":"referencia"}],"confianca":1.0,"contraexemplos":[],"descricao":"- A parábola **não é** o fractal Koch D≈1.26 — é uma curva 1D com parâmetro intrínseco `s`. - O experimento v1 errava ao fazer box-counting em `(calor, c²)`; v2 mede **sobre** a curva. - Koch mora na órbita `z²` mod P **e** parcialmente nos endereços Cantor/Fibonacci de `s`. - AGM não fecha o comprimento de arco em forma elementar — só aproxima com duplicação de dígitos.","epistemico":"fato","exemplos":[],"id":"interpretacao_fisica_opcional","memoria":"semantica","meta":{"dominio":"manual","duplicata_sugerida":[],"nivel":6,"verificacao":"parte de conceito mais amplo 'interpretacao_fisica_opcional'"},"origem":"KOCH_PARABOLA.md","palavras_chave":[],"sinonimos":[],"tipo":"conceito","titulo":"Interpretação física (opcional)"}
\section{Interpretação física (opcional)}
- A parábola **não é** o fractal Koch D\ensuremath{\approx}1.26 — é uma curva 1D com parâmetro intrínseco `s`. - O experimento v1 errava ao fazer box-counting em `(calor, c\ensuremath{^{2}})`; v2 mede **sobre** a curva. - Koch mora na órbita `z\ensuremath{^{2}}` mod P **e** parcialmente nos endereços Cantor/Fibonacci de `s`. - AGM não fecha o comprimento de arco em forma elementar — só aproxima com duplicação de dígitos.
\prov{relações: referencia casl-propagation}
\prov{proveniência: KOCH\_PARABOLA.md \textperiodcentered{} fato \textperiodcentered{} confiança 1.0 \textperiodcentered{} domínio manual \textperiodcentered{} história 7 versões no jornal}

%CRISTAL {"arestas":[{"alvo":"habitantes_gentil_de_sitter","confianca":1.0,"epistemico":"fato","origem":"DS_UNIVERSO.md","rel":"parte_de"}],"confianca":1.0,"contraexemplos":[],"descricao":"Cada candidato PoW é um **habitante** com espectro binário; o passo `word→ouro` fixa coordenadas **Lopes** na lemniscata; o embedding Hopf coloca o habitante numa **fatia S³** do universo **de Sitter**, com escala `cosh(τ)` dada pela massa espectral `total`.","epistemico":"fato","exemplos":[],"id":"interpretacao_hipotese_nao_prova","memoria":"semantica","meta":{"dominio":"manual","nivel":6,"verificacao":"slug idêntico — enriquecer, não duplicar"},"origem":"DS_UNIVERSO.md","palavras_chave":[],"sinonimos":[],"tipo":"conceito","titulo":"Interpretação (hipótese, não prova)"}
\section{Interpretação (hipótese, não prova)}
Cada candidato PoW é um **habitante** com espectro binário; o passo `word\ensuremath{\to}ouro` fixa coordenadas **Lopes** na lemniscata; o embedding Hopf coloca o habitante numa **fatia S\ensuremath{^{3}}** do universo **de Sitter**, com escala `cosh(\ensuremath{\tau})` dada pela massa espectral `total`.
\prov{relações: parte\_de habitantes\_gentil\_de\_sitter}
\prov{proveniência: DS\_UNIVERSO.md \textperiodcentered{} fato \textperiodcentered{} confiança 1.0 \textperiodcentered{} domínio manual \textperiodcentered{} história 7 versões no jornal}

%CRISTAL {"arestas":[{"alvo":"distinguisher_sha_hipercomplexo","confianca":1.0,"epistemico":"fato","origem":"wiki","rel":"usa"},{"alvo":"cifra","confianca":1.0,"epistemico":"fato","origem":"wiki","rel":"relaciona"}],"confianca":1.0,"contraexemplos":[],"descricao":"A fronteira de difusão r≈6 do SHA é invariante à escolha do anel base (ℝ, GF(2), ℤ/2ⁿ): o ARX alterna operações lineares em réguas mutuamente incompatíveis (adição rígida na métrica circular U(1); XOR/rotação rígidas em Hamming), e nenhuma trivialização única lineariza tudo — a robustez do SHA É essa impossibilidade.","epistemico":"fato","exemplos":[],"id":"invariancia_de_regua_arx","memoria":"semantica","meta":{"autor":"Aarão/broca-so","dominio":"broca_repos","fonte":""},"origem":"wiki","palavras_chave":[],"sinonimos":[],"tipo":"conceito","titulo":"Invariância de Régua (a robustez do ARX)"}
\section{Invariância de Régua (a robustez do ARX)}
A fronteira de difusão r\ensuremath{\approx}6 do SHA é invariante à escolha do anel base (\ensuremath{\mathbb{R}}, GF(2), \ensuremath{\mathbb{Z}}/2\ensuremath{^{n}}): o ARX alterna operações lineares em réguas mutuamente incompatíveis (adição rígida na métrica circular U(1); XOR/rotação rígidas em Hamming), e nenhuma trivialização única lineariza tudo — a robustez do SHA É essa impossibilidade.
\prov{relações: usa distinguisher\_sha\_hipercomplexo; relaciona cifra}
\prov{proveniência: wiki \textperiodcentered{} fato \textperiodcentered{} confiança 1.0 \textperiodcentered{} domínio broca\_repos \textperiodcentered{} história 3 versões no jornal}

%CRISTAL {"arestas":[{"alvo":"matematica_estelar","confianca":1.0,"epistemico":"fato","origem":"wiki","rel":"parte_de"},{"alvo":"reta_de_ouro","confianca":1.0,"epistemico":"fato","origem":"wiki","rel":"relaciona"},{"alvo":"conjugados_galois_parabola","confianca":1.0,"epistemico":"fato","origem":"wiki","rel":"relaciona"}],"confianca":1.0,"contraexemplos":[],"descricao":"A lei que costura o lado euclidiano e o relativístico da geometria de ouro — funciona como conservação (análogo de energia-momento) preservada pela multiplicação e pela transformada. O custo em comutatividade é quadrático: os primeiros incrementos de produto vetorial são 'baratos'.","epistemico":"fato","exemplos":[],"id":"invariante_estelar","memoria":"semantica","meta":{"autor":"Lopes/Aarão","dominio":"hiper","fonte":""},"origem":"wiki","palavras_chave":[],"sinonimos":[],"tipo":"conceito","titulo":"O Invariante a+c²=1"}
\section{O Invariante a+c\ensuremath{^{2}}=1}
A lei que costura o lado euclidiano e o relativístico da geometria de ouro — funciona como conservação (análogo de energia-momento) preservada pela multiplicação e pela transformada. O custo em comutatividade é quadrático: os primeiros incrementos de produto vetorial são 'baratos'.
\prov{relações: parte\_de matematica\_estelar; relaciona reta\_de\_ouro; relaciona conjugados\_galois\_parabola}
\prov{proveniência: wiki \textperiodcentered{} fato \textperiodcentered{} confiança 1.0 \textperiodcentered{} domínio hiper \textperiodcentered{} história 5 versões no jornal}

%CRISTAL {"arestas":[{"alvo":"borda_critica_vazio","confianca":1.0,"epistemico":"fato","origem":"wiki","rel":"usa"},{"alvo":"regua_de_gentil","confianca":1.0,"epistemico":"fato","origem":"wiki","rel":"relaciona"},{"alvo":"godel_incompletude","confianca":1.0,"epistemico":"fato","origem":"wiki","rel":"usa"}],"confianca":1.0,"contraexemplos":[],"descricao":"Duas provas de que o atrator AGM (o ponto onde real e relativístico coincidem) existe se e somente se é indecidível: (a) s*=AGM(a,b) é transcendente salvo degeneração (aproximável, nunca decidido); (b) os mensuráveis são contáveis (medida 0), o imensurável tem medida 1 — se 'medir tudo' fosse possível, 1=0. Chaitin Ω o habita.","epistemico":"fato","exemplos":[],"id":"invariante_indecidivel","memoria":"semantica","meta":{"autor":"Lopes/Aarão","dominio":"hiper","fonte":""},"origem":"wiki","palavras_chave":[],"sinonimos":[],"tipo":"conceito","titulo":"Existência ⟺ Indecidibilidade (o invariante imensurável)"}
\section{Existência \ensuremath{\Longleftrightarrow} Indecidibilidade (o invariante imensurável)}
Duas provas de que o atrator AGM (o ponto onde real e relativístico coincidem) existe se e somente se é indecidível: (a) s*=AGM(a,b) é transcendente salvo degeneração (aproximável, nunca decidido); (b) os mensuráveis são contáveis (medida 0), o imensurável tem medida 1 — se 'medir tudo' fosse possível, 1=0. Chaitin \ensuremath{\Omega} o habita.
\prov{relações: usa borda\_critica\_vazio; relaciona regua\_de\_gentil; usa godel\_incompletude}
\prov{proveniência: wiki \textperiodcentered{} fato \textperiodcentered{} confiança 1.0 \textperiodcentered{} domínio hiper \textperiodcentered{} história 6 versões no jornal}

%CRISTAL {"arestas":[{"alvo":"fechos_centro_da_janela","confianca":1.0,"epistemico":"fato","origem":"DS_UNIVERSO.md","rel":"parte_de"},{"alvo":"observacao","confianca":1.0,"epistemico":"fato","origem":"POW_METAL_SETOR.md","rel":"parte_de"}],"confianca":1.0,"contraexemplos":[],"descricao":"| nonce | acopl | atrator | fase | koch | lemn_erro | dist | fecho | |-------|-------|---------|------|------|-----------|------|-------| | 1572799 | 0.6022 | 520 | 0.1293 | 0.0858 | 0.00000 | 64 |  | | 1572800 | 0.5987 | 238 | 0.1300 | 0.0857 | 0.00000 | 63 |  | | 1572801 | 0.6155 | 773 | 0.1305 | 0.0857 | 0.00000 | 62 |  | | 1572802 | 0.5970 | 353 | 0.1294 | 0.0858 | 0.00000 | 61 |  | | 1572803 | 0.8652 | 685 | 0.1299 | 0.0857 | 0.00000 | 60 |  |","epistemico":"fato","exemplos":[],"id":"janela_1572863csv","memoria":"semantica","meta":{"dominio":"manual","duplicata_sugerida":[],"nivel":6,"verificacao":"parte de conceito mais amplo 'janela_1572863csv'"},"origem":"POW_METAL_PRE_FECHO.md","palavras_chave":[],"sinonimos":[],"tipo":"conceito","titulo":"`janela_1572863.csv`"}
\section{`janela\_1572863.csv`}
| nonce | acopl | atrator | fase | koch | lemn\_erro | dist | fecho | |-------|-------|---------|------|------|-----------|------|-------| | 1572799 | 0.6022 | 520 | 0.1293 | 0.0858 | 0.00000 | 64 |  | | 1572800 | 0.5987 | 238 | 0.1300 | 0.0857 | 0.00000 | 63 |  | | 1572801 | 0.6155 | 773 | 0.1305 | 0.0857 | 0.00000 | 62 |  | | 1572802 | 0.5970 | 353 | 0.1294 | 0.0858 | 0.00000 | 61 |  | | 1572803 | 0.8652 | 685 | 0.1299 | 0.0857 | 0.00000 | 60 |  |
\prov{relações: parte\_de fechos\_centro\_da\_janela; parte\_de observacao}
\prov{proveniência: POW\_METAL\_PRE\_FECHO.md \textperiodcentered{} fato \textperiodcentered{} confiança 1.0 \textperiodcentered{} domínio manual \textperiodcentered{} história 14 versões no jornal}

%CRISTAL {"arestas":[{"alvo":"fechos_centro_da_janela","confianca":1.0,"epistemico":"fato","origem":"DS_UNIVERSO.md","rel":"parte_de"},{"alvo":"observacao","confianca":1.0,"epistemico":"fato","origem":"POW_METAL_SETOR.md","rel":"parte_de"}],"confianca":1.0,"contraexemplos":[],"descricao":"| nonce | acopl | atrator | fase | koch | lemn_erro | dist | fecho | |-------|-------|---------|------|------|-----------|------|-------| | 1966015 | 0.6027 | 520 | 0.1293 | 0.0858 | 0.00000 | 64 |  | | 1966016 | 0.5992 | 238 | 0.1300 | 0.0857 | 0.00000 | 63 |  | | 1966017 | 0.6163 | 773 | 0.1305 | 0.0857 | 0.00000 | 62 |  | | 1966018 | 0.5976 | 353 | 0.1294 | 0.0858 | 0.00000 | 61 |  | | 1966019 | 0.8657 | 685 | 0.1299 | 0.0857 | 0.00000 | 60 |  |","epistemico":"fato","exemplos":[],"id":"janela_1966079csv","memoria":"semantica","meta":{"dominio":"manual","duplicata_sugerida":[],"nivel":6,"verificacao":"parte de conceito mais amplo 'janela_1966079csv'"},"origem":"POW_METAL_PRE_FECHO.md","palavras_chave":[],"sinonimos":[],"tipo":"conceito","titulo":"`janela_1966079.csv`"}
\section{`janela\_1966079.csv`}
| nonce | acopl | atrator | fase | koch | lemn\_erro | dist | fecho | |-------|-------|---------|------|------|-----------|------|-------| | 1966015 | 0.6027 | 520 | 0.1293 | 0.0858 | 0.00000 | 64 |  | | 1966016 | 0.5992 | 238 | 0.1300 | 0.0857 | 0.00000 | 63 |  | | 1966017 | 0.6163 | 773 | 0.1305 | 0.0857 | 0.00000 | 62 |  | | 1966018 | 0.5976 | 353 | 0.1294 | 0.0858 | 0.00000 | 61 |  | | 1966019 | 0.8657 | 685 | 0.1299 | 0.0857 | 0.00000 | 60 |  |
\prov{relações: parte\_de fechos\_centro\_da\_janela; parte\_de observacao}
\prov{proveniência: POW\_METAL\_PRE\_FECHO.md \textperiodcentered{} fato \textperiodcentered{} confiança 1.0 \textperiodcentered{} domínio manual \textperiodcentered{} história 14 versões no jornal}

%CRISTAL {"arestas":[{"alvo":"observacao","confianca":1.0,"epistemico":"fato","origem":"POW_METAL_SETOR.md","rel":"parte_de"}],"confianca":1.0,"contraexemplos":[],"descricao":"| nonce | assinatura | posto | gl | dist_deg | prox_dz | atrator | koch | lemn | fecho | |-------|------------|-------|----|---------|---------|---------|------|------|-------| | 2855 | `S1:R0:A730:P2:D0` | 2 | 24 | 0.3286 | 0.0029 | 730 | 0.0858 | 0.00000 |  | | 2856 | `S1:R0:A709:P2:D0` | 2 | 24 | 0.3301 | 0.0027 | 709 | 0.0859 | 0.00000 |  | | 2857 | `S1:R0:A403:P2:D0` | 2 | 24 | 0.3294 | 0.0028 | 403 | 0.0858 | 0.00000 |  |","epistemico":"fato","exemplos":[],"id":"janela_2863csv","memoria":"semantica","meta":{"dominio":"manual","nivel":6,"verificacao":"slug idêntico — enriquecer, não duplicar"},"origem":"POW_METAL_SETOR.md","palavras_chave":[],"sinonimos":[],"tipo":"conceito","titulo":"`janela_2863.csv`"}
\section{`janela\_2863.csv`}
| nonce | assinatura | posto | gl | dist\_deg | prox\_dz | atrator | koch | lemn | fecho | |-------|------------|-------|----|---------|---------|---------|------|------|-------| | 2855 | `S1:R0:A730:P2:D0` | 2 | 24 | 0.3286 | 0.0029 | 730 | 0.0858 | 0.00000 |  | | 2856 | `S1:R0:A709:P2:D0` | 2 | 24 | 0.3301 | 0.0027 | 709 | 0.0859 | 0.00000 |  | | 2857 | `S1:R0:A403:P2:D0` | 2 | 24 | 0.3294 | 0.0028 | 403 | 0.0858 | 0.00000 |  |
\prov{relações: parte\_de observacao}
\prov{proveniência: POW\_METAL\_SETOR.md \textperiodcentered{} fato \textperiodcentered{} confiança 1.0 \textperiodcentered{} domínio manual \textperiodcentered{} história 4 versões no jornal}

%CRISTAL {"arestas":[{"alvo":"observacao","confianca":1.0,"epistemico":"fato","origem":"POW_METAL_SETOR.md","rel":"parte_de"}],"confianca":1.0,"contraexemplos":[],"descricao":"| nonce | assinatura | posto | gl | dist_deg | prox_dz | atrator | koch | lemn | fecho | |-------|------------|-------|----|---------|---------|---------|------|------|-------| | 0 | `S1:R0:A577:P2:D0` | 2 | 24 | 0.3276 | 0.0030 | 577 | 0.0858 | 0.00000 |  | | 1 | `S1:R0:A43:P2:D0` | 2 | 23 | 0.3268 | 0.0031 | 43 | 0.0857 | 0.00000 |  | | 2 | `S1:R0:A696:P2:D0` | 2 | 23 | 0.3285 | 0.0029 | 696 | 0.0858 | 0.00000 |  |","epistemico":"fato","exemplos":[],"id":"janela_3csv","memoria":"semantica","meta":{"dominio":"manual","nivel":6,"verificacao":"slug idêntico — enriquecer, não duplicar"},"origem":"POW_METAL_SETOR.md","palavras_chave":[],"sinonimos":[],"tipo":"conceito","titulo":"`janela_3.csv`"}
\section{`janela\_3.csv`}
| nonce | assinatura | posto | gl | dist\_deg | prox\_dz | atrator | koch | lemn | fecho | |-------|------------|-------|----|---------|---------|---------|------|------|-------| | 0 | `S1:R0:A577:P2:D0` | 2 | 24 | 0.3276 | 0.0030 | 577 | 0.0858 | 0.00000 |  | | 1 | `S1:R0:A43:P2:D0` | 2 | 23 | 0.3268 | 0.0031 | 43 | 0.0857 | 0.00000 |  | | 2 | `S1:R0:A696:P2:D0` | 2 | 23 | 0.3285 | 0.0029 | 696 | 0.0858 | 0.00000 |  |
\prov{relações: parte\_de observacao}
\prov{proveniência: POW\_METAL\_SETOR.md \textperiodcentered{} fato \textperiodcentered{} confiança 1.0 \textperiodcentered{} domínio manual \textperiodcentered{} história 4 versões no jornal}

%CRISTAL {"arestas":[{"alvo":"observacao","confianca":1.0,"epistemico":"fato","origem":"POW_METAL_SETOR.md","rel":"parte_de"}],"confianca":1.0,"contraexemplos":[],"descricao":"| nonce | assinatura | posto | gl | dist_deg | prox_dz | atrator | koch | lemn | fecho | |-------|------------|-------|----|---------|---------|---------|------|------|-------| | 572855 | `S1:R0:A149:P2:D0` | 2 | 23 | 0.3287 | 0.0029 | 149 | 0.0858 | 0.00000 |  | | 572856 | `S1:R0:A253:P2:D0` | 2 | 23 | 0.3302 | 0.0027 | 253 | 0.0859 | 0.00000 |  | | 572857 | `S1:R0:A296:P2:D0` | 2 | 23 | 0.3294 | 0.0028 | 296 | 0.0858 | 0.00000 |  |","epistemico":"fato","exemplos":[],"id":"janela_572863csv","memoria":"semantica","meta":{"dominio":"manual","nivel":6,"verificacao":"slug idêntico — enriquecer, não duplicar"},"origem":"POW_METAL_SETOR.md","palavras_chave":[],"sinonimos":[],"tipo":"conceito","titulo":"`janela_572863.csv`"}
\section{`janela\_572863.csv`}
| nonce | assinatura | posto | gl | dist\_deg | prox\_dz | atrator | koch | lemn | fecho | |-------|------------|-------|----|---------|---------|---------|------|------|-------| | 572855 | `S1:R0:A149:P2:D0` | 2 | 23 | 0.3287 | 0.0029 | 149 | 0.0858 | 0.00000 |  | | 572856 | `S1:R0:A253:P2:D0` | 2 | 23 | 0.3302 | 0.0027 | 253 | 0.0859 | 0.00000 |  | | 572857 | `S1:R0:A296:P2:D0` | 2 | 23 | 0.3294 | 0.0028 | 296 | 0.0858 | 0.00000 |  |
\prov{relações: parte\_de observacao}
\prov{proveniência: POW\_METAL\_SETOR.md \textperiodcentered{} fato \textperiodcentered{} confiança 1.0 \textperiodcentered{} domínio manual \textperiodcentered{} história 4 versões no jornal}

%CRISTAL {"arestas":[{"alvo":"observacao","confianca":1.0,"epistemico":"fato","origem":"POW_METAL_SETOR.md","rel":"parte_de"}],"confianca":1.0,"contraexemplos":[],"descricao":"| nonce | assinatura | posto | gl | dist_deg | prox_dz | atrator | koch | lemn | fecho | |-------|------------|-------|----|---------|---------|---------|------|------|-------| | 6071 | `S1:R0:A829:P2:D0` | 2 | 23 | 0.3262 | 0.0031 | 829 | 0.0857 | 0.00000 |  | | 6072 | `S1:R0:A863:P2:D0` | 2 | 23 | 0.3278 | 0.0030 | 863 | 0.0858 | 0.00000 |  | | 6073 | `S1:R0:A396:P2:D0` | 2 | 23 | 0.3270 | 0.0031 | 396 | 0.0857 | 0.00000 |  |","epistemico":"fato","exemplos":[],"id":"janela_6079csv","memoria":"semantica","meta":{"dominio":"manual","nivel":6,"verificacao":"slug idêntico — enriquecer, não duplicar"},"origem":"POW_METAL_SETOR.md","palavras_chave":[],"sinonimos":[],"tipo":"conceito","titulo":"`janela_6079.csv`"}
\section{`janela\_6079.csv`}
| nonce | assinatura | posto | gl | dist\_deg | prox\_dz | atrator | koch | lemn | fecho | |-------|------------|-------|----|---------|---------|---------|------|------|-------| | 6071 | `S1:R0:A829:P2:D0` | 2 | 23 | 0.3262 | 0.0031 | 829 | 0.0857 | 0.00000 |  | | 6072 | `S1:R0:A863:P2:D0` | 2 | 23 | 0.3278 | 0.0030 | 863 | 0.0858 | 0.00000 |  | | 6073 | `S1:R0:A396:P2:D0` | 2 | 23 | 0.3270 | 0.0031 | 396 | 0.0857 | 0.00000 |  |
\prov{relações: parte\_de observacao}
\prov{proveniência: POW\_METAL\_SETOR.md \textperiodcentered{} fato \textperiodcentered{} confiança 1.0 \textperiodcentered{} domínio manual \textperiodcentered{} história 4 versões no jornal}

%CRISTAL {"arestas":[{"alvo":"observacao","confianca":1.0,"epistemico":"fato","origem":"POW_METAL_SETOR.md","rel":"parte_de"}],"confianca":1.0,"contraexemplos":[],"descricao":"| nonce | assinatura | posto | gl | dist_deg | prox_dz | atrator | koch | lemn | fecho | |-------|------------|-------|----|---------|---------|---------|------|------|-------| | 55 | `S1:R0:A262:P2:D0` | 2 | 23 | 0.3269 | 0.0031 | 262 | 0.0857 | 0.00000 |  | | 56 | `S1:R0:A269:P2:D0` | 2 | 23 | 0.3285 | 0.0029 | 269 | 0.0858 | 0.00000 |  | | 57 | `S1:R0:A134:P2:D0` | 2 | 23 | 0.3277 | 0.0030 | 134 | 0.0858 | 0.00000 |  |","epistemico":"fato","exemplos":[],"id":"janela_63csv","memoria":"semantica","meta":{"dominio":"manual","nivel":6,"verificacao":"slug idêntico — enriquecer, não duplicar"},"origem":"POW_METAL_SETOR.md","palavras_chave":[],"sinonimos":[],"tipo":"conceito","titulo":"`janela_63.csv`"}
\section{`janela\_63.csv`}
| nonce | assinatura | posto | gl | dist\_deg | prox\_dz | atrator | koch | lemn | fecho | |-------|------------|-------|----|---------|---------|---------|------|------|-------| | 55 | `S1:R0:A262:P2:D0` | 2 | 23 | 0.3269 | 0.0031 | 262 | 0.0857 | 0.00000 |  | | 56 | `S1:R0:A269:P2:D0` | 2 | 23 | 0.3285 | 0.0029 | 269 | 0.0858 | 0.00000 |  | | 57 | `S1:R0:A134:P2:D0` | 2 | 23 | 0.3277 | 0.0030 | 134 | 0.0858 | 0.00000 |  |
\prov{relações: parte\_de observacao}
\prov{proveniência: POW\_METAL\_SETOR.md \textperiodcentered{} fato \textperiodcentered{} confiança 1.0 \textperiodcentered{} domínio manual \textperiodcentered{} história 4 versões no jornal}

%CRISTAL {"arestas":[{"alvo":"observacao","confianca":1.0,"epistemico":"fato","origem":"POW_METAL_SETOR.md","rel":"parte_de"}],"confianca":1.0,"contraexemplos":[],"descricao":"| nonce | assinatura | posto | gl | dist_deg | prox_dz | atrator | koch | lemn | fecho | |-------|------------|-------|----|---------|---------|---------|------|------|-------| | 66071 | `S1:R0:A844:P2:D0` | 2 | 23 | 0.3261 | 0.0032 | 844 | 0.0857 | 0.00000 |  | | 66072 | `S1:R0:A134:P2:D0` | 2 | 23 | 0.3277 | 0.0030 | 134 | 0.0858 | 0.00000 |  | | 66073 | `S1:R0:A262:P2:D0` | 2 | 22 | 0.3269 | 0.0031 | 262 | 0.0857 | 0.00000 |  |","epistemico":"fato","exemplos":[],"id":"janela_66079csv","memoria":"semantica","meta":{"dominio":"manual","nivel":6,"verificacao":"slug idêntico — enriquecer, não duplicar"},"origem":"POW_METAL_SETOR.md","palavras_chave":[],"sinonimos":[],"tipo":"conceito","titulo":"`janela_66079.csv`"}
\section{`janela\_66079.csv`}
| nonce | assinatura | posto | gl | dist\_deg | prox\_dz | atrator | koch | lemn | fecho | |-------|------------|-------|----|---------|---------|---------|------|------|-------| | 66071 | `S1:R0:A844:P2:D0` | 2 | 23 | 0.3261 | 0.0032 | 844 | 0.0857 | 0.00000 |  | | 66072 | `S1:R0:A134:P2:D0` | 2 | 23 | 0.3277 | 0.0030 | 134 | 0.0858 | 0.00000 |  | | 66073 | `S1:R0:A262:P2:D0` | 2 | 22 | 0.3269 | 0.0031 | 262 | 0.0857 | 0.00000 |  |
\prov{relações: parte\_de observacao}
\prov{proveniência: POW\_METAL\_SETOR.md \textperiodcentered{} fato \textperiodcentered{} confiança 1.0 \textperiodcentered{} domínio manual \textperiodcentered{} história 4 versões no jornal}

%CRISTAL {"arestas":[{"alvo":"observacao","confianca":1.0,"epistemico":"fato","origem":"POW_METAL_SETOR.md","rel":"parte_de"}],"confianca":1.0,"contraexemplos":[],"descricao":"| nonce | assinatura | posto | gl | dist_deg | prox_dz | atrator | koch | lemn | fecho | |-------|------------|-------|----|---------|---------|---------|------|------|-------| | 72855 | `S1:R0:A267:P2:D0` | 2 | 23 | 0.3254 | 0.0032 | 267 | 0.0857 | 0.00000 |  | | 72856 | `S1:R0:A734:P2:D0` | 2 | 23 | 0.3270 | 0.0031 | 734 | 0.0857 | 0.00000 |  | | 72857 | `S1:R0:A956:P2:D0` | 2 | 23 | 0.3262 | 0.0032 | 956 | 0.0857 | 0.00000 |  |","epistemico":"fato","exemplos":[],"id":"janela_72863csv","memoria":"semantica","meta":{"dominio":"manual","nivel":6,"verificacao":"slug idêntico — enriquecer, não duplicar"},"origem":"POW_METAL_SETOR.md","palavras_chave":[],"sinonimos":[],"tipo":"conceito","titulo":"`janela_72863.csv`"}
\section{`janela\_72863.csv`}
| nonce | assinatura | posto | gl | dist\_deg | prox\_dz | atrator | koch | lemn | fecho | |-------|------------|-------|----|---------|---------|---------|------|------|-------| | 72855 | `S1:R0:A267:P2:D0` | 2 | 23 | 0.3254 | 0.0032 | 267 | 0.0857 | 0.00000 |  | | 72856 | `S1:R0:A734:P2:D0` | 2 | 23 | 0.3270 | 0.0031 | 734 | 0.0857 | 0.00000 |  | | 72857 | `S1:R0:A956:P2:D0` | 2 | 23 | 0.3262 | 0.0032 | 956 | 0.0857 | 0.00000 |  |
\prov{relações: parte\_de observacao}
\prov{proveniência: POW\_METAL\_SETOR.md \textperiodcentered{} fato \textperiodcentered{} confiança 1.0 \textperiodcentered{} domínio manual \textperiodcentered{} história 4 versões no jornal}

%CRISTAL {"arestas":[{"alvo":"observacao","confianca":1.0,"epistemico":"fato","origem":"POW_METAL_SETOR.md","rel":"parte_de"}],"confianca":1.0,"contraexemplos":[],"descricao":"| nonce | assinatura | posto | gl | dist_deg | prox_dz | atrator | koch | lemn | fecho | |-------|------------|-------|----|---------|---------|---------|------|------|-------| | 71 | `S1:R0:A192:P2:D0` | 2 | 23 | 0.3261 | 0.0032 | 192 | 0.0857 | 0.00000 |  | | 72 | `S1:R0:A392:P2:D0` | 2 | 23 | 0.3277 | 0.0030 | 392 | 0.0858 | 0.00000 |  | | 73 | `S1:R0:A858:P2:D0` | 2 | 23 | 0.3269 | 0.0031 | 858 | 0.0857 | 0.00000 |  |","epistemico":"fato","exemplos":[],"id":"janela_79csv","memoria":"semantica","meta":{"dominio":"manual","nivel":6,"verificacao":"slug idêntico — enriquecer, não duplicar"},"origem":"POW_METAL_SETOR.md","palavras_chave":[],"sinonimos":[],"tipo":"conceito","titulo":"`janela_79.csv`"}
\section{`janela\_79.csv`}
| nonce | assinatura | posto | gl | dist\_deg | prox\_dz | atrator | koch | lemn | fecho | |-------|------------|-------|----|---------|---------|---------|------|------|-------| | 71 | `S1:R0:A192:P2:D0` | 2 | 23 | 0.3261 | 0.0032 | 192 | 0.0857 | 0.00000 |  | | 72 | `S1:R0:A392:P2:D0` | 2 | 23 | 0.3277 | 0.0030 | 392 | 0.0858 | 0.00000 |  | | 73 | `S1:R0:A858:P2:D0` | 2 | 23 | 0.3269 | 0.0031 | 858 | 0.0857 | 0.00000 |  |
\prov{relações: parte\_de observacao}
\prov{proveniência: POW\_METAL\_SETOR.md \textperiodcentered{} fato \textperiodcentered{} confiança 1.0 \textperiodcentered{} domínio manual \textperiodcentered{} história 4 versões no jornal}

%CRISTAL {"arestas":[{"alvo":"observacao","confianca":1.0,"epistemico":"fato","origem":"POW_METAL_SETOR.md","rel":"parte_de"}],"confianca":1.0,"contraexemplos":[],"descricao":"| nonce | assinatura | posto | gl | dist_deg | prox_dz | atrator | koch | lemn | fecho | |-------|------------|-------|----|---------|---------|---------|------|------|-------| | 855 | `S1:R0:A326:P2:D0` | 2 | 22 | 0.3254 | 0.0032 | 326 | 0.0857 | 0.00000 |  | | 856 | `S1:R0:A262:P2:D0` | 2 | 23 | 0.3269 | 0.0031 | 262 | 0.0857 | 0.00000 |  | | 857 | `S1:R0:A844:P2:D0` | 2 | 22 | 0.3261 | 0.0032 | 844 | 0.0857 | 0.00000 |  |","epistemico":"fato","exemplos":[],"id":"janela_863csv","memoria":"semantica","meta":{"dominio":"manual","nivel":6,"verificacao":"slug idêntico — enriquecer, não duplicar"},"origem":"POW_METAL_SETOR.md","palavras_chave":[],"sinonimos":[],"tipo":"conceito","titulo":"`janela_863.csv`"}
\section{`janela\_863.csv`}
| nonce | assinatura | posto | gl | dist\_deg | prox\_dz | atrator | koch | lemn | fecho | |-------|------------|-------|----|---------|---------|---------|------|------|-------| | 855 | `S1:R0:A326:P2:D0` | 2 | 22 | 0.3254 | 0.0032 | 326 | 0.0857 | 0.00000 |  | | 856 | `S1:R0:A262:P2:D0` | 2 | 23 | 0.3269 | 0.0031 | 262 | 0.0857 | 0.00000 |  | | 857 | `S1:R0:A844:P2:D0` | 2 | 22 | 0.3261 | 0.0032 | 844 | 0.0857 | 0.00000 |  |
\prov{relações: parte\_de observacao}
\prov{proveniência: POW\_METAL\_SETOR.md \textperiodcentered{} fato \textperiodcentered{} confiança 1.0 \textperiodcentered{} domínio manual \textperiodcentered{} história 4 versões no jornal}

%CRISTAL {"arestas":[{"alvo":"observacao","confianca":1.0,"epistemico":"fato","origem":"POW_METAL_SETOR.md","rel":"parte_de"}],"confianca":1.0,"contraexemplos":[],"descricao":"| nonce | assinatura | posto | gl | dist_deg | prox_dz | atrator | koch | lemn | fecho | |-------|------------|-------|----|---------|---------|---------|------|------|-------| | 966071 | `S1:R0:A982:P2:D0` | 2 | 24 | 0.3287 | 0.0029 | 982 | 0.0858 | 0.00000 |  | | 966072 | `S1:R0:A312:P2:D0` | 2 | 24 | 0.3302 | 0.0027 | 312 | 0.0859 | 0.00000 |  | | 966073 | `S1:R0:A276:P2:D0` | 2 | 24 | 0.3295 | 0.0028 | 276 | 0.0858 | 0.00000 |  |","epistemico":"fato","exemplos":[],"id":"janela_966079csv","memoria":"semantica","meta":{"dominio":"manual","nivel":6,"verificacao":"slug idêntico — enriquecer, não duplicar"},"origem":"POW_METAL_SETOR.md","palavras_chave":[],"sinonimos":[],"tipo":"conceito","titulo":"`janela_966079.csv`"}
\section{`janela\_966079.csv`}
| nonce | assinatura | posto | gl | dist\_deg | prox\_dz | atrator | koch | lemn | fecho | |-------|------------|-------|----|---------|---------|---------|------|------|-------| | 966071 | `S1:R0:A982:P2:D0` | 2 | 24 | 0.3287 | 0.0029 | 982 | 0.0858 | 0.00000 |  | | 966072 | `S1:R0:A312:P2:D0` | 2 | 24 | 0.3302 | 0.0027 | 312 | 0.0859 | 0.00000 |  | | 966073 | `S1:R0:A276:P2:D0` | 2 | 24 | 0.3295 | 0.0028 | 276 | 0.0858 | 0.00000 |  |
\prov{relações: parte\_de observacao}
\prov{proveniência: POW\_METAL\_SETOR.md \textperiodcentered{} fato \textperiodcentered{} confiança 1.0 \textperiodcentered{} domínio manual \textperiodcentered{} história 4 versões no jornal}

%CRISTAL {"arestas":[{"alvo":"observacao","confianca":1.0,"epistemico":"fato","origem":"POW_METAL_SETOR.md","rel":"parte_de"}],"confianca":1.0,"contraexemplos":[],"descricao":"| nonce | assinatura | posto | gl | dist_deg | prox_dz | atrator | koch | lemn | fecho | |-------|------------|-------|----|---------|---------|---------|------|------|-------| | 1 | `S1:R0:A43:P2:D0` | 2 | 23 | 0.3268 | 0.0031 | 43 | 0.0857 | 0.00000 |  | | 2 | `S1:R0:A696:P2:D0` | 2 | 23 | 0.3285 | 0.0029 | 696 | 0.0858 | 0.00000 |  | | 3 | `S1:R0:A392:P2:D0` | 2 | 23 | 0.3277 | 0.0030 | 392 | 0.0858 | 0.00000 |  |","epistemico":"fato","exemplos":[],"id":"janela_9csv","memoria":"semantica","meta":{"dominio":"manual","nivel":6,"verificacao":"slug idêntico — enriquecer, não duplicar"},"origem":"POW_METAL_SETOR.md","palavras_chave":[],"sinonimos":[],"tipo":"conceito","titulo":"`janela_9.csv`"}
\section{`janela\_9.csv`}
| nonce | assinatura | posto | gl | dist\_deg | prox\_dz | atrator | koch | lemn | fecho | |-------|------------|-------|----|---------|---------|---------|------|------|-------| | 1 | `S1:R0:A43:P2:D0` | 2 | 23 | 0.3268 | 0.0031 | 43 | 0.0857 | 0.00000 |  | | 2 | `S1:R0:A696:P2:D0` | 2 | 23 | 0.3285 | 0.0029 | 696 | 0.0858 | 0.00000 |  | | 3 | `S1:R0:A392:P2:D0` | 2 | 23 | 0.3277 | 0.0030 | 392 | 0.0858 | 0.00000 |  |
\prov{relações: parte\_de observacao}
\prov{proveniência: POW\_METAL\_SETOR.md \textperiodcentered{} fato \textperiodcentered{} confiança 1.0 \textperiodcentered{} domínio manual \textperiodcentered{} história 4 versões no jornal}

%CRISTAL {"arestas":[{"alvo":"colapso","confianca":1.0,"epistemico":"fato","origem":"manual","rel":"referencia"},{"alvo":"pow_gato","confianca":0.5,"epistemico":"fato","origem":"wiki","rel":"relaciona"},{"alvo":"sessao_interativa","confianca":0.5,"epistemico":"fato","origem":"wiki","rel":"relaciona"},{"alvo":"prova_isa","confianca":0.5,"epistemico":"fato","origem":"wiki","rel":"relaciona"},{"alvo":"koch_memoria","confianca":0.5,"epistemico":"fato","origem":"wiki","rel":"relaciona"},{"alvo":"transcricao","confianca":0.5,"epistemico":"fato","origem":"wiki","rel":"relaciona"},{"alvo":"montagem","confianca":0.5,"epistemico":"fato","origem":"wiki","rel":"relaciona"},{"alvo":"python","confianca":0.5,"epistemico":"fato","origem":"wiki","rel":"relaciona"},{"alvo":"koch_parabola","confianca":0.5,"epistemico":"fato","origem":"wiki","rel":"relaciona"},{"alvo":"koch_pell","confianca":0.5,"epistemico":"fato","origem":"wiki","rel":"relaciona"},{"alvo":"pow_bit_taxa","confianca":0.5,"epistemico":"fato","origem":"wiki","rel":"relaciona"},{"alvo":"parabola_isa","confianca":0.5,"epistemico":"fato","origem":"wiki","rel":"relaciona"}],"confianca":1.0,"contraexemplos":[],"descricao":"Manual: ula/KOCH_ATLAS.md","epistemico":"fato","exemplos":[],"id":"koch_atlas","memoria":"semantica","meta":{"arquivo":"ula/KOCH_ATLAS.md","dominio":"manual","nivel":6},"origem":"manual","palavras_chave":[],"sinonimos":[],"tipo":"conceito","titulo":"KOCH_ATLAS"}
\section{KOCH\_ATLAS}
Manual: ula/KOCH\_ATLAS.md
\prov{relações: referencia colapso; relaciona pow\_gato; relaciona sessao\_interativa; relaciona prova\_isa; relaciona koch\_memoria; relaciona transcricao; relaciona montagem; relaciona python; relaciona koch\_parabola; relaciona koch\_pell; relaciona pow\_bit\_taxa; relaciona parabola\_isa}
\prov{proveniência: manual \textperiodcentered{} fato \textperiodcentered{} confiança 1.0 \textperiodcentered{} domínio manual \textperiodcentered{} história 15 versões no jornal}

%CRISTAL {"arestas":[{"alvo":"koch_atlas","confianca":1.0,"epistemico":"fato","origem":"manual","rel":"parte_de"},{"alvo":"colapso","confianca":1.0,"epistemico":"fato","origem":"manual","rel":"explica"}],"confianca":1.0,"contraexemplos":[],"descricao":"> Koch como **sistema de coordenadas** sobre Z[i] na interface real×ouro. > **Sem poda.** Medição rigorosa de estrutura geométrica.","epistemico":"fato","exemplos":[],"id":"koch_atlas_experimento_de_parametrizacao","memoria":"semantica","meta":{"dominio":"manual","nivel":6,"verificacao":"slug idêntico — enriquecer, não duplicar"},"origem":"KOCH_ATLAS.md","palavras_chave":[],"sinonimos":[],"tipo":"conceito","titulo":"Koch Atlas — Experimento de Parametrização"}
\section{Koch Atlas — Experimento de Parametrização}
> Koch como **sistema de coordenadas** sobre Z[i] na interface real$\times$ouro. > **Sem poda.** Medição rigorosa de estrutura geométrica.
\prov{relações: parte\_de koch\_atlas; explica colapso}
\prov{proveniência: KOCH\_ATLAS.md \textperiodcentered{} fato \textperiodcentered{} confiança 1.0 \textperiodcentered{} domínio manual \textperiodcentered{} história 56 versões no jornal}

%CRISTAL {"arestas":[{"alvo":"gato","confianca":1.0,"epistemico":"fato","origem":"manual","rel":"prova"},{"alvo":"transcricao","confianca":0.5,"epistemico":"fato","origem":"wiki","rel":"relaciona"},{"alvo":"vida-morte","confianca":0.5,"epistemico":"fato","origem":"wiki","rel":"relaciona"},{"alvo":"parabola_destruir","confianca":0.5,"epistemico":"fato","origem":"wiki","rel":"relaciona"},{"alvo":"teste_ordem","confianca":0.5,"epistemico":"fato","origem":"wiki","rel":"relaciona"},{"alvo":"pow_metal_pre_fecho","confianca":0.5,"epistemico":"fato","origem":"wiki","rel":"relaciona"},{"alvo":"por_sessao_conversadadostelemetriasessaojsonl","confianca":0.5,"epistemico":"fato","origem":"wiki","rel":"relaciona"},{"alvo":"pow_metal_funil","confianca":0.5,"epistemico":"fato","origem":"wiki","rel":"relaciona"},{"alvo":"pow_metal_ressonancia","confianca":0.5,"epistemico":"fato","origem":"wiki","rel":"relaciona"},{"alvo":"resposta_a_pergunta_dificil","confianca":0.5,"epistemico":"fato","origem":"wiki","rel":"relaciona"},{"alvo":"pow_geometria","confianca":0.5,"epistemico":"fato","origem":"wiki","rel":"relaciona"},{"alvo":"parabola_pa","confianca":0.5,"epistemico":"fato","origem":"wiki","rel":"relaciona"},{"alvo":"topologia","confianca":0.5,"epistemico":"fato","origem":"wiki","rel":"relaciona"},{"alvo":"pilha","confianca":0.5,"epistemico":"fato","origem":"wiki","rel":"relaciona"},{"alvo":"log_legado_conversadadostelemetriajsonl","confianca":0.5,"epistemico":"fato","origem":"wiki","rel":"relaciona"},{"alvo":"testes","confianca":0.5,"epistemico":"fato","origem":"wiki","rel":"relaciona"},{"alvo":"metais","confianca":0.5,"epistemico":"fato","origem":"wiki","rel":"relaciona"},{"alvo":"pow_metal_setor","confianca":0.5,"epistemico":"fato","origem":"wiki","rel":"relaciona"},{"alvo":"lobo_frontal","confianca":0.5,"epistemico":"fato","origem":"wiki","rel":"relaciona"}],"confianca":1.0,"contraexemplos":[],"descricao":"Manual: ula/KOCH_CANTOR_FIB_PROVA.md","epistemico":"fato","exemplos":[],"id":"koch_cantor_fib_prova","memoria":"semantica","meta":{"arquivo":"ula/KOCH_CANTOR_FIB_PROVA.md","dominio":"manual","nivel":6},"origem":"manual","palavras_chave":[],"sinonimos":[],"tipo":"conceito","titulo":"KOCH_CANTOR_FIB_PROVA"}
\section{KOCH\_CANTOR\_FIB\_PROVA}
Manual: ula/KOCH\_CANTOR\_FIB\_PROVA.md
\prov{relações: prova gato; relaciona transcricao; relaciona vida-morte; relaciona parabola\_destruir; relaciona teste\_ordem; relaciona pow\_metal\_pre\_fecho; relaciona por\_sessao\_conversadadostelemetriasessaojsonl; relaciona pow\_metal\_funil; relaciona pow\_metal\_ressonancia; relaciona resposta\_a\_pergunta\_dificil; relaciona pow\_geometria; relaciona parabola\_pa; relaciona topologia; relaciona pilha; relaciona log\_legado\_conversadadostelemetriajsonl; relaciona testes; relaciona metais; relaciona pow\_metal\_setor; relaciona lobo\_frontal}
\prov{proveniência: manual \textperiodcentered{} fato \textperiodcentered{} confiança 1.0 \textperiodcentered{} domínio manual \textperiodcentered{} história 22 versões no jornal}

%CRISTAL {"arestas":[{"alvo":"koch_cantor_fib_prova","confianca":1.0,"epistemico":"fato","origem":"manual","rel":"parte_de"},{"alvo":"word","confianca":0.5,"epistemico":"fato","origem":"wiki","rel":"relaciona"},{"alvo":"virtualizacao","confianca":0.5,"epistemico":"fato","origem":"wiki","rel":"relaciona"}],"confianca":1.0,"contraexemplos":[],"descricao":"> `make koch-cantor-fib` colhe dados; este documento formaliza a hipótese e o esboço de prova.","epistemico":"fato","exemplos":[],"id":"koch_codifica_cantorfibonacci_prova_e_evidencia","memoria":"semantica","meta":{"dominio":"manual","nivel":6,"verificacao":"slug idêntico — enriquecer, não duplicar"},"origem":"KOCH_CANTOR_FIB_PROVA.md","palavras_chave":[],"sinonimos":[],"tipo":"conceito","titulo":"Koch codifica Cantor–Fibonacci — prova e evidência"}
\section{Koch codifica Cantor–Fibonacci — prova e evidência}
> `make koch-cantor-fib` colhe dados; este documento formaliza a hipótese e o esboço de prova.
\prov{relações: parte\_de koch\_cantor\_fib\_prova; relaciona word; relaciona virtualizacao}
\prov{proveniência: KOCH\_CANTOR\_FIB\_PROVA.md \textperiodcentered{} fato \textperiodcentered{} confiança 1.0 \textperiodcentered{} domínio manual \textperiodcentered{} história 6 versões no jornal}

%CRISTAL {"arestas":[{"alvo":"koch_memoria","confianca":1.0,"epistemico":"fato","origem":"manual","rel":"parte_de"},{"alvo":"colapso","confianca":1.0,"epistemico":"fato","origem":"manual","rel":"usa"},{"alvo":"semantica","confianca":1.0,"epistemico":"fato","origem":"manual","rel":"referencia"}],"confianca":1.0,"contraexemplos":[],"descricao":"> `make koch-memoria` — verifica se Koch **guarda** informação (recipiente) em vez de **codificar** como Cantor/Fibonacci.","epistemico":"fato","exemplos":[],"id":"koch_como_recipiente_memoria_nao-associativa","memoria":"semantica","meta":{"dominio":"manual","nivel":6,"verificacao":"slug idêntico — enriquecer, não duplicar"},"origem":"KOCH_MEMORIA.md","palavras_chave":[],"sinonimos":[],"tipo":"conceito","titulo":"Koch como recipiente — memória não-associativa"}
\section{Koch como recipiente — memória não-associativa}
> `make koch-memoria` — verifica se Koch **guarda** informação (recipiente) em vez de **codificar** como Cantor/Fibonacci.
\prov{relações: parte\_de koch\_memoria; usa colapso; referencia semantica}
\prov{proveniência: KOCH\_MEMORIA.md \textperiodcentered{} fato \textperiodcentered{} confiança 1.0 \textperiodcentered{} domínio manual \textperiodcentered{} história 104 versões no jornal}

%CRISTAL {"arestas":[{"alvo":"hopfield","confianca":1.0,"epistemico":"fato","origem":"manual","rel":"analogia"},{"alvo":"pow_espectral","confianca":0.5,"epistemico":"fato","origem":"wiki","rel":"relaciona"},{"alvo":"transcricao","confianca":0.5,"epistemico":"fato","origem":"wiki","rel":"relaciona"},{"alvo":"vida-morte","confianca":0.5,"epistemico":"fato","origem":"wiki","rel":"relaciona"},{"alvo":"prova_isa","confianca":0.5,"epistemico":"fato","origem":"wiki","rel":"relaciona"},{"alvo":"pow_bit_taxa","confianca":0.5,"epistemico":"fato","origem":"wiki","rel":"relaciona"},{"alvo":"parabola_isa","confianca":0.5,"epistemico":"fato","origem":"wiki","rel":"relaciona"},{"alvo":"word","confianca":0.5,"epistemico":"fato","origem":"wiki","rel":"relaciona"}],"confianca":1.0,"contraexemplos":[],"descricao":"Manual: ula/KOCH_MEMORIA.md","epistemico":"fato","exemplos":[],"id":"koch_memoria","memoria":"semantica","meta":{"arquivo":"ula/KOCH_MEMORIA.md","dominio":"manual","nivel":6},"origem":"manual","palavras_chave":[],"sinonimos":[],"tipo":"conceito","titulo":"KOCH_MEMORIA"}
\section{KOCH\_MEMORIA}
Manual: ula/KOCH\_MEMORIA.md
\prov{relações: analogia hopfield; relaciona pow\_espectral; relaciona transcricao; relaciona vida-morte; relaciona prova\_isa; relaciona pow\_bit\_taxa; relaciona parabola\_isa; relaciona word}
\prov{proveniência: manual \textperiodcentered{} fato \textperiodcentered{} confiança 1.0 \textperiodcentered{} domínio manual \textperiodcentered{} história 11 versões no jornal}

%CRISTAL {"arestas":[{"alvo":"vida","confianca":1.0,"epistemico":"fato","origem":"manual","rel":"referencia"},{"alvo":"pow_espectral","confianca":0.5,"epistemico":"fato","origem":"wiki","rel":"relaciona"},{"alvo":"transcricao","confianca":0.5,"epistemico":"fato","origem":"wiki","rel":"relaciona"},{"alvo":"prova_isa","confianca":0.5,"epistemico":"fato","origem":"wiki","rel":"relaciona"},{"alvo":"teste_ordem","confianca":0.5,"epistemico":"fato","origem":"wiki","rel":"relaciona"},{"alvo":"python","confianca":0.5,"epistemico":"fato","origem":"wiki","rel":"relaciona"},{"alvo":"sessao_interativa","confianca":0.5,"epistemico":"fato","origem":"wiki","rel":"relaciona"},{"alvo":"pow_gato","confianca":0.5,"epistemico":"fato","origem":"wiki","rel":"relaciona"}],"confianca":1.0,"contraexemplos":[],"descricao":"Manual: ula/KOCH_PARABOLA.md","epistemico":"fato","exemplos":[],"id":"koch_parabola","memoria":"semantica","meta":{"arquivo":"ula/KOCH_PARABOLA.md","dominio":"manual","nivel":6},"origem":"manual","palavras_chave":[],"sinonimos":[],"tipo":"conceito","titulo":"KOCH_PARABOLA"}
\section{KOCH\_PARABOLA}
Manual: ula/KOCH\_PARABOLA.md
\prov{relações: referencia vida; relaciona pow\_espectral; relaciona transcricao; relaciona prova\_isa; relaciona teste\_ordem; relaciona python; relaciona sessao\_interativa; relaciona pow\_gato}
\prov{proveniência: manual \textperiodcentered{} fato \textperiodcentered{} confiança 1.0 \textperiodcentered{} domínio manual \textperiodcentered{} história 11 versões no jornal}

%CRISTAL {"arestas":[{"alvo":"koch_parabola","confianca":1.0,"epistemico":"fato","origem":"manual","rel":"parte_de"},{"alvo":"word","confianca":0.5,"epistemico":"fato","origem":"wiki","rel":"relaciona"},{"alvo":"syscall","confianca":0.5,"epistemico":"fato","origem":"wiki","rel":"relaciona"},{"alvo":"virtualizacao","confianca":0.5,"epistemico":"fato","origem":"wiki","rel":"relaciona"},{"alvo":"vontade_de_poder","confianca":0.5,"epistemico":"fato","origem":"wiki","rel":"relaciona"}],"confianca":1.0,"contraexemplos":[],"descricao":"> **Nota:** Koch **não codifica** Cantor/Fibonacci — é **recipiente** (ver `KOCH_MEMORIA.md`, `make koch-memoria`).","epistemico":"fato","exemplos":[],"id":"koch_parabola_v2_sobre_a_curva_desdistorcida","memoria":"semantica","meta":{"dominio":"manual","nivel":6,"verificacao":"slug idêntico — enriquecer, não duplicar"},"origem":"KOCH_PARABOLA.md","palavras_chave":[],"sinonimos":[],"tipo":"conceito","titulo":"Koch × Parábola v2 — sobre a curva, desdistorcida"}
\section{Koch $\times$ Parábola v2 — sobre a curva, desdistorcida}
> **Nota:** Koch **não codifica** Cantor/Fibonacci — é **recipiente** (ver `KOCH\_MEMORIA.md`, `make koch-memoria`).
\prov{relações: parte\_de koch\_parabola; relaciona word; relaciona syscall; relaciona virtualizacao; relaciona vontade\_de\_poder}
\prov{proveniência: KOCH\_PARABOLA.md \textperiodcentered{} fato \textperiodcentered{} confiança 1.0 \textperiodcentered{} domínio manual \textperiodcentered{} história 8 versões no jornal}

%CRISTAL {"arestas":[{"alvo":"goldchain","confianca":1.0,"epistemico":"fato","origem":"manual","rel":"referencia"},{"alvo":"prova_isa","confianca":0.5,"epistemico":"fato","origem":"wiki","rel":"relaciona"},{"alvo":"python","confianca":0.5,"epistemico":"fato","origem":"wiki","rel":"relaciona"},{"alvo":"pow_gato","confianca":0.5,"epistemico":"fato","origem":"wiki","rel":"relaciona"},{"alvo":"pow_espectral","confianca":0.5,"epistemico":"fato","origem":"wiki","rel":"relaciona"},{"alvo":"virtualizacao","confianca":0.5,"epistemico":"fato","origem":"wiki","rel":"relaciona"},{"alvo":"word","confianca":0.5,"epistemico":"fato","origem":"wiki","rel":"relaciona"},{"alvo":"vida-morte","confianca":0.5,"epistemico":"fato","origem":"wiki","rel":"relaciona"}],"confianca":1.0,"contraexemplos":[],"descricao":"Manual: ula/KOCH_PELL.md","epistemico":"fato","exemplos":[],"id":"koch_pell","memoria":"semantica","meta":{"arquivo":"ula/KOCH_PELL.md","dominio":"manual","nivel":6},"origem":"manual","palavras_chave":[],"sinonimos":[],"tipo":"conceito","titulo":"KOCH_PELL"}
\section{KOCH\_PELL}
Manual: ula/KOCH\_PELL.md
\prov{relações: referencia goldchain; relaciona prova\_isa; relaciona python; relaciona pow\_gato; relaciona pow\_espectral; relaciona virtualizacao; relaciona word; relaciona vida-morte}
\prov{proveniência: manual \textperiodcentered{} fato \textperiodcentered{} confiança 1.0 \textperiodcentered{} domínio manual \textperiodcentered{} história 11 versões no jornal}

%CRISTAL {"arestas":[{"alvo":"koch_pell","confianca":1.0,"epistemico":"fato","origem":"manual","rel":"parte_de"},{"alvo":"vida-morte","confianca":0.5,"epistemico":"fato","origem":"wiki","rel":"relaciona"},{"alvo":"word","confianca":0.5,"epistemico":"fato","origem":"wiki","rel":"relaciona"},{"alvo":"virtualizacao","confianca":0.5,"epistemico":"fato","origem":"wiki","rel":"relaciona"},{"alvo":"while","confianca":0.5,"epistemico":"fato","origem":"wiki","rel":"relaciona"}],"confianca":1.0,"contraexemplos":[],"descricao":"> `make koch-pell` — testa se o floco Koch guarda **apenas** números de Pell (prata `A₂`).","epistemico":"fato","exemplos":[],"id":"koch_pell_verificacao","memoria":"semantica","meta":{"dominio":"manual","nivel":6,"verificacao":"slug idêntico — enriquecer, não duplicar"},"origem":"KOCH_PELL.md","palavras_chave":[],"sinonimos":[],"tipo":"conceito","titulo":"Koch × Pell — verificação"}
\section{Koch $\times$ Pell — verificação}
> `make koch-pell` — testa se o floco Koch guarda **apenas** números de Pell (prata `A\ensuremath{_{2}}`).
\prov{relações: parte\_de koch\_pell; relaciona vida-morte; relaciona word; relaciona virtualizacao; relaciona while}
\prov{proveniência: KOCH\_PELL.md \textperiodcentered{} fato \textperiodcentered{} confiança 1.0 \textperiodcentered{} domínio manual \textperiodcentered{} história 8 versões no jornal}

%CRISTAL {"arestas":[{"alvo":"soberania_da_chave","confianca":1.0,"epistemico":"fato","origem":"wiki","rel":"usa"},{"alvo":"sistema_e_lei_nao_dono","confianca":1.0,"epistemico":"fato","origem":"wiki","rel":"relaciona"},{"alvo":"goldchain","confianca":1.0,"epistemico":"fato","origem":"wiki","rel":"relaciona"}],"confianca":1.0,"contraexemplos":[],"descricao":"O lastro não é escassez (roubável) mas a impossibilidade de arrombamento por teorema (A=0): o cofre vale porque não abre, não porque guarda metal. O valor mora em três camadas — Pell/ouro-branco cifrado (interno), liquidação ETH (externo) e a cadeia de certificados sha256 (a história). Única perda irreversível: a chave.","epistemico":"inferencia","exemplos":[],"id":"lastro_do_cofre","memoria":"semantica","meta":{"autor":"Aarão/broca-so","dominio":"broca_repos","fonte":""},"origem":"wiki","palavras_chave":[],"sinonimos":[],"tipo":"conceito","titulo":"O Lastro é a Irreversibilidade, não a Escassez"}
\section{O Lastro é a Irreversibilidade, não a Escassez}
O lastro não é escassez (roubável) mas a impossibilidade de arrombamento por teorema (A=0): o cofre vale porque não abre, não porque guarda metal. O valor mora em três camadas — Pell/ouro-branco cifrado (interno), liquidação ETH (externo) e a cadeia de certificados sha256 (a história). Única perda irreversível: a chave.
\prov{relações: usa soberania\_da\_chave; relaciona sistema\_e\_lei\_nao\_dono; relaciona goldchain}
\prov{proveniência: wiki \textperiodcentered{} inferencia \textperiodcentered{} confiança 1.0 \textperiodcentered{} domínio broca\_repos \textperiodcentered{} história 4 versões no jornal}

%CRISTAL {"arestas":[{"alvo":"memoria_viva","confianca":1.0,"epistemico":"fato","origem":"paper","rel":"parte_de"},{"alvo":"ula_da_broca_arquitetura_minima","confianca":1.0,"epistemico":"fato","origem":"README.md","rel":"parte_de"}],"confianca":1.0,"contraexemplos":[],"descricao":"``` ula/ neuron_io.c/h     espelho do átomo (mesmo loop que neuronio.c) word.h            tipo Word registradores.c/h ula.c/h cifra.c/h         núcleo Ouro Branco — silver/gold/bronze memoria.c/h       mem/NNNN.bin instrucoes.c/h cpu.c/h assembler.c testes/ mem/ ```","epistemico":"fato","exemplos":[],"id":"layout","memoria":"semantica","meta":{"dominio":"manual","nivel":6,"verificacao":"slug idêntico — enriquecer, não duplicar"},"origem":"README.md","palavras_chave":[],"sinonimos":[],"tipo":"conceito","titulo":"Layout"}
\section{Layout}
``` ula/ neuron\_io.c/h     espelho do átomo (mesmo loop que neuronio.c) word.h            tipo Word registradores.c/h ula.c/h cifra.c/h         núcleo Ouro Branco — silver/gold/bronze memoria.c/h       mem/NNNN.bin instrucoes.c/h cpu.c/h assembler.c testes/ mem/ ```
\prov{relações: parte\_de memoria\_viva; parte\_de ula\_da\_broca\_arquitetura\_minima}
\prov{proveniência: README.md \textperiodcentered{} fato \textperiodcentered{} confiança 1.0 \textperiodcentered{} domínio manual \textperiodcentered{} história 6 versões no jornal}

%CRISTAL {"arestas":[{"alvo":"pow_lemniscata_koch_corpo_de_gauss","confianca":1.0,"epistemico":"fato","origem":"POW_METAL.md","rel":"parte_de"}],"confianca":1.0,"contraexemplos":[],"descricao":"| Motivo | Significado | |--------|-------------| | `degenerado_real/ouro` | sem distribuição | | `fora_lemniscata` | calor + c² ≠ 1 | | `fora_corpo_gauss` | norma nula ou interface vazia | | `orbita_ouro` | passo A₁ fora da lemniscata |","epistemico":"fato","exemplos":[],"id":"leis_de_impossibilidade","memoria":"semantica","meta":{"dominio":"manual","nivel":6,"verificacao":"título igual (slug diferente)"},"origem":"POW_METAL.md","palavras_chave":[],"sinonimos":[],"tipo":"conceito","titulo":"Leis de impossibilidade"}
\section{Leis de impossibilidade}
| Motivo | Significado | |--------|-------------| | `degenerado\_real/ouro` | sem distribuição | | `fora\_lemniscata` | calor + c\ensuremath{^{2}} \ensuremath{\ne} 1 | | `fora\_corpo\_gauss` | norma nula ou interface vazia | | `orbita\_ouro` | passo A\ensuremath{_{1}} fora da lemniscata |
\prov{relações: parte\_de pow\_lemniscata\_koch\_corpo\_de\_gauss}
\prov{proveniência: POW\_METAL.md \textperiodcentered{} fato \textperiodcentered{} confiança 1.0 \textperiodcentered{} domínio manual \textperiodcentered{} história 4 versões no jornal}

%CRISTAL {"arestas":[{"alvo":"lemniscata_ponto_i","confianca":1.0,"epistemico":"fato","origem":"paper","rel":"especializa"},{"alvo":"maquina_evolutiva","confianca":1.0,"epistemico":"fato","origem":"wiki","rel":"explica"},{"alvo":"transformada_dourada","confianca":1.0,"epistemico":"fato","origem":"wiki","rel":"relaciona"}],"confianca":1.0,"contraexemplos":[],"descricao":"Os dois laços da lemniscata (∞) da máquina evolutiva são duais de Pontryagin: Ẑ=S¹ e Ŝ¹=ℤ — o laço discreto (bits, contagem, Perron) e o contínuo (fase, ℂâ) determinam-se por Fourier (verificado a 10⁻¹⁶). A cúspide central é o ponto auto-dual: a gaussiana e−πx², único ponto fixo da transformada.","epistemico":"fato","exemplos":[],"id":"lemniscata_pontryagin","memoria":"semantica","meta":{"autor":"Lopes/Aarão","dominio":"hiper","fonte":""},"origem":"wiki","palavras_chave":[],"sinonimos":[],"tipo":"conceito","titulo":"A Lemniscata como Dualidade de Pontryagin"}
\section{A Lemniscata como Dualidade de Pontryagin}
Os dois laços da lemniscata (\ensuremath{\infty}) da máquina evolutiva são duais de Pontryagin: \ensuremath{\hat{Z}}=S\ensuremath{^{1}} e \ensuremath{\hat{S}}\ensuremath{^{1}}=\ensuremath{\mathbb{Z}} — o laço discreto (bits, contagem, Perron) e o contínuo (fase, \ensuremath{\mathbb{C}}â) determinam-se por Fourier (verificado a 10\ensuremath{^{-16}}). A cúspide central é o ponto auto-dual: a gaussiana e\ensuremath{-}\ensuremath{\pi}x\ensuremath{^{2}}, único ponto fixo da transformada.
\prov{relações: especializa lemniscata\_ponto\_i; explica maquina\_evolutiva; relaciona transformada\_dourada}
\prov{proveniência: wiki \textperiodcentered{} fato \textperiodcentered{} confiança 1.0 \textperiodcentered{} domínio hiper \textperiodcentered{} história 7 versões no jornal}

%CRISTAL {"arestas":[{"alvo":"transformada_dourada","confianca":1.0,"epistemico":"fato","origem":"wiki","rel":"usa"},{"alvo":"hopfield","confianca":1.0,"epistemico":"fato","origem":"wiki","rel":"relaciona"},{"alvo":"mente_computacao","confianca":1.0,"epistemico":"fato","origem":"wiki","rel":"relaciona"}],"confianca":1.0,"contraexemplos":[],"descricao":"Teorema: uma rede de profundidade L que acessa um cristal W por produtos matriz-vetor extrai o modo dominante com erro angular εL∼ρL>0 para todo L finito (power iteration) — quebra o cristal mas nunca o evapora. A evaporação exata é a diagonalização = flip 𝓖. Aplica-se à atenção dos transformers (QKᵀ); o RoPE já usa escala geométrica ≈Ω.","epistemico":"fato","exemplos":[],"id":"limite_atencao_cristalina","memoria":"semantica","meta":{"autor":"Lopes/Aarão","dominio":"hiper","fonte":""},"origem":"wiki","palavras_chave":[],"sinonimos":[],"tipo":"conceito","titulo":"O Limite de Evaporação da Atenção Cristalina"}
\section{O Limite de Evaporação da Atenção Cristalina}
Teorema: uma rede de profundidade L que acessa um cristal W por produtos matriz-vetor extrai o modo dominante com erro angular \ensuremath{\varepsilon}L\ensuremath{\sim}\ensuremath{\rho}L>0 para todo L finito (power iteration) — quebra o cristal mas nunca o evapora. A evaporação exata é a diagonalização = flip \ensuremath{\mathcal{G}}. Aplica-se à atenção dos transformers (QK\ensuremath{^{T}}); o RoPE já usa escala geométrica \ensuremath{\approx}\ensuremath{\Omega}.
\prov{relações: usa transformada\_dourada; relaciona hopfield; relaciona mente\_computacao}
\prov{proveniência: wiki \textperiodcentered{} fato \textperiodcentered{} confiança 1.0 \textperiodcentered{} domínio hiper \textperiodcentered{} história 6 versões no jornal}

%CRISTAL {"arestas":[{"alvo":"destruir_a_parabola","confianca":1.0,"epistemico":"fato","origem":"PARABOLA_DESTRUIR.md","rel":"parte_de"}],"confianca":1.0,"contraexemplos":[],"descricao":"``` === DESTRUIR PARÁBOLA === passos totais ........ 1716582 validos .............. 699682 degenerados .......... 1016900  (sem distribuição — constante 0.5 seria necessária) fallback_05 .......... 1016900  (estados zero — hipótese não testável sem convencao) violacoes ............ 0  (|a+c²−1| > 1e-09) max erro ............. 0.00000000000000000e+00 programas ............ 1013","epistemico":"fato","exemplos":[],"id":"log_bruto","memoria":"semantica","meta":{"dominio":"manual","nivel":6,"verificacao":"slug idêntico — enriquecer, não duplicar"},"origem":"PARABOLA_DESTRUIR.md","palavras_chave":[],"sinonimos":[],"tipo":"conceito","titulo":"Log bruto"}
\section{Log bruto}
``` === DESTRUIR PARÁBOLA === passos totais ........ 1716582 validos .............. 699682 degenerados .......... 1016900  (sem distribuição — constante 0.5 seria necessária) fallback\_05 .......... 1016900  (estados zero — hipótese não testável sem convencao) violacoes ............ 0  (|a+c\ensuremath{^{2}}\ensuremath{-}1| > 1e-09) max erro ............. 0.00000000000000000e+00 programas ............ 1013
\prov{relações: parte\_de destruir\_a\_parabola}
\prov{proveniência: PARABOLA\_DESTRUIR.md \textperiodcentered{} fato \textperiodcentered{} confiança 1.0 \textperiodcentered{} domínio manual \textperiodcentered{} história 4 versões no jornal}

%CRISTAL {"arestas":[{"alvo":"a_cacada_da_parabola_mapa_vivo_da_isa","confianca":1.0,"epistemico":"fato","origem":"PARABOLA_ISA.md","rel":"parte_de"},{"alvo":"metafisica","confianca":0.5,"epistemico":"fato","origem":"wiki","rel":"relaciona"},{"alvo":"utilitarismo","confianca":0.5,"epistemico":"fato","origem":"wiki","rel":"relaciona"},{"alvo":"teoria_do_valor","confianca":0.5,"epistemico":"fato","origem":"wiki","rel":"relaciona"},{"alvo":"regua_associativa_e_o_risco_de_vazamento","confianca":0.5,"epistemico":"fato","origem":"wiki","rel":"relaciona"}],"confianca":1.0,"contraexemplos":[],"descricao":"Cada passo: parábola entre estado antes/depois (`parabola_entre`).","epistemico":"fato","exemplos":[],"id":"mapa_por_opcode","memoria":"semantica","meta":{"dominio":"manual","nivel":6,"verificacao":"slug idêntico — enriquecer, não duplicar"},"origem":"PARABOLA_ISA.md","palavras_chave":[],"sinonimos":[],"tipo":"conceito","titulo":"Mapa por opcode"}
\section{Mapa por opcode}
Cada passo: parábola entre estado antes/depois (`parabola\_entre`).
\prov{relações: parte\_de a\_cacada\_da\_parabola\_mapa\_vivo\_da\_isa; relaciona metafisica; relaciona utilitarismo; relaciona teoria\_do\_valor; relaciona regua\_associativa\_e\_o\_risco\_de\_vazamento}
\prov{proveniência: PARABOLA\_ISA.md \textperiodcentered{} fato \textperiodcentered{} confiança 1.0 \textperiodcentered{} domínio manual \textperiodcentered{} história 10 versões no jornal}

%CRISTAL {"arestas":[{"alvo":"maquina_tropical","confianca":1.0,"epistemico":"fato","origem":"wiki","rel":"especializa"},{"alvo":"corpo_f2k","confianca":1.0,"epistemico":"fato","origem":"wiki","rel":"usa"},{"alvo":"cayley_dickson","confianca":1.0,"epistemico":"fato","origem":"wiki","rel":"usa"},{"alvo":"associador","confianca":1.0,"epistemico":"fato","origem":"wiki","rel":"relaciona"},{"alvo":"matematica_estelar","confianca":1.0,"epistemico":"fato","origem":"wiki","rel":"relaciona"}],"confianca":1.0,"contraexemplos":[],"descricao":"Organismo computacional definido só por ⊕ (XOR=gap), ⊙ (XNOR=correlação), pop (peso de Hamming=medida), shift e NOT. Gap=distância de Hamming; reconhecimento=colapso winner-take-all (Born booleana); o flip Φₛ é involutivo (⊕ é sua própria inversa); o 'erro perpendicular' gera a torre de Cayley-Dickson em bits. Substrato-independente (CPU/GPU/FPGA computam o mesmo organismo).","epistemico":"fato","exemplos":[],"id":"maquina_booleana","memoria":"semantica","meta":{"autor":"Lopes/Aarão","dominio":"hiper","fonte":""},"origem":"wiki","palavras_chave":[],"sinonimos":[],"tipo":"conceito","titulo":"A Máquina Booleana (organismo sobre 𝔽₂)"}
\section{A Máquina Booleana (organismo sobre \ensuremath{\mathbb{F}}\ensuremath{_{2}})}
Organismo computacional definido só por \ensuremath{\oplus} (XOR=gap), \ensuremath{\odot} (XNOR=correlação), pop (peso de Hamming=medida), shift e NOT. Gap=distância de Hamming; reconhecimento=colapso winner-take-all (Born booleana); o flip \ensuremath{\Phi}\ensuremath{_{s}} é involutivo (\ensuremath{\oplus} é sua própria inversa); o 'erro perpendicular' gera a torre de Cayley-Dickson em bits. Substrato-independente (CPU/GPU/FPGA computam o mesmo organismo).
\prov{relações: especializa maquina\_tropical; usa corpo\_f2k; usa cayley\_dickson; relaciona associador; relaciona matematica\_estelar}
\prov{proveniência: wiki \textperiodcentered{} fato \textperiodcentered{} confiança 1.0 \textperiodcentered{} domínio hiper \textperiodcentered{} história 7 versões no jornal}

%CRISTAL {"arestas":[{"alvo":"algebras_hipercomplexas","confianca":1.0,"epistemico":"fato","origem":"wiki","rel":"usa"},{"alvo":"reta_de_ouro","confianca":1.0,"epistemico":"fato","origem":"wiki","rel":"usa"},{"alvo":"torre_hurwitz","confianca":1.0,"epistemico":"fato","origem":"wiki","rel":"relaciona"}],"confianca":1.0,"contraexemplos":[],"descricao":"O framework 'estelar': a geometria de ouro em regime relativístico — um produto s-deformado (⋆ₛ) que interpola do plano/associativo (s=0, split-complex) ao curvo/octoniônico (s=1), sob o invariante a+c²=1.","epistemico":"fato","exemplos":[],"id":"matematica_estelar","memoria":"semantica","meta":{"autor":"Lopes/Aarão","dominio":"hiper","fonte":""},"origem":"wiki","palavras_chave":[],"sinonimos":[],"tipo":"conceito","titulo":"Matemática Estelar"}
\section{Matemática Estelar}
O framework 'estelar': a geometria de ouro em regime relativístico — um produto s-deformado (\ensuremath{\star}\ensuremath{_{s}}) que interpola do plano/associativo (s=0, split-complex) ao curvo/octoniônico (s=1), sob o invariante a+c\ensuremath{^{2}}=1.
\prov{relações: usa algebras\_hipercomplexas; usa reta\_de\_ouro; relaciona torre\_hurwitz}
\prov{proveniência: wiki \textperiodcentered{} fato \textperiodcentered{} confiança 1.0 \textperiodcentered{} domínio hiper \textperiodcentered{} história 5 versões no jornal}

%CRISTAL {"arestas":[{"alvo":"pow_taxa_de_sucesso_do_bit_broca","confianca":1.0,"epistemico":"fato","origem":"POW_BIT_TAXA.md","rel":"parte_de"},{"alvo":"testes","confianca":0.5,"epistemico":"fato","origem":"wiki","rel":"relaciona"},{"alvo":"pow_metal_pre_fecho","confianca":0.5,"epistemico":"fato","origem":"wiki","rel":"relaciona"},{"alvo":"projeto_cristaljson","confianca":0.5,"epistemico":"fato","origem":"wiki","rel":"relaciona"},{"alvo":"mecanica_quantica_hub","confianca":0.5,"epistemico":"fato","origem":"wiki","rel":"relaciona"},{"alvo":"topologia","confianca":0.5,"epistemico":"fato","origem":"wiki","rel":"relaciona"},{"alvo":"utilitarismo","confianca":0.5,"epistemico":"fato","origem":"wiki","rel":"relaciona"}],"confianca":1.0,"contraexemplos":[],"descricao":"Linhas = SHA real · Colunas = bit","epistemico":"fato","exemplos":[],"id":"matriz_de_confusao","memoria":"semantica","meta":{"dominio":"manual","nivel":6,"verificacao":"slug idêntico — enriquecer, não duplicar"},"origem":"POW_BIT_TAXA.md","palavras_chave":[],"sinonimos":[],"tipo":"conceito","titulo":"Matriz de confusão"}
\section{Matriz de confusão}
Linhas = SHA real \textperiodcentered  Colunas = bit
\prov{relações: parte\_de pow\_taxa\_de\_sucesso\_do\_bit\_broca; relaciona testes; relaciona pow\_metal\_pre\_fecho; relaciona projeto\_cristaljson; relaciona mecanica\_quantica\_hub; relaciona topologia; relaciona utilitarismo}
\prov{proveniência: POW\_BIT\_TAXA.md \textperiodcentered{} fato \textperiodcentered{} confiança 1.0 \textperiodcentered{} domínio manual \textperiodcentered{} história 10 versões no jornal}

%CRISTAL {"arestas":[{"alvo":"matematica_estelar","confianca":1.0,"epistemico":"fato","origem":"wiki","rel":"parte_de"},{"alvo":"transformada_estelar","confianca":1.0,"epistemico":"fato","origem":"wiki","rel":"usa"},{"alvo":"colapso_observador","confianca":1.0,"epistemico":"fato","origem":"wiki","rel":"relaciona"}],"confianca":1.0,"contraexemplos":[],"descricao":"O sistema dinâmico da transformada: estados z com a+c²=1, evolução i∂ψ/∂t=Hₛψ com Hₛ da multiplicação em 𝓔ₛ⁸; o associador (gap) é o 'propulsor' que nunca fecha, e o colapso quântico ocorre quando o associador satura. Inclui o oscilador invertido s̈=+s (sombra mecânica da compacidade de G₂).","epistemico":"fato","exemplos":[],"id":"mecanica_estelar","memoria":"semantica","meta":{"autor":"Lopes/Aarão","dominio":"hiper","fonte":""},"origem":"wiki","palavras_chave":[],"sinonimos":[],"tipo":"conceito","titulo":"Mecânica Estelar"}
\section{Mecânica Estelar}
O sistema dinâmico da transformada: estados z com a+c\ensuremath{^{2}}=1, evolução i\ensuremath{\partial}\ensuremath{\psi}/\ensuremath{\partial}t=H\ensuremath{_{s}}\ensuremath{\psi} com H\ensuremath{_{s}} da multiplicação em \ensuremath{\mathcal{E}}\ensuremath{_{s}}\ensuremath{^{8}}; o associador (gap) é o 'propulsor' que nunca fecha, e o colapso quântico ocorre quando o associador satura. Inclui o oscilador invertido s\ensuremath{}=+s (sombra mecânica da compacidade de G\ensuremath{_{2}}).
\prov{relações: parte\_de matematica\_estelar; usa transformada\_estelar; relaciona colapso\_observador}
\prov{proveniência: wiki \textperiodcentered{} fato \textperiodcentered{} confiança 1.0 \textperiodcentered{} domínio hiper \textperiodcentered{} história 5 versões no jornal}

%CRISTAL {"arestas":[{"alvo":"corpo_dual","confianca":1.0,"epistemico":"fato","origem":"wiki","rel":"usa"},{"alvo":"o_substrato_a_maquina_e_funcao_de_onda","confianca":1.0,"epistemico":"fato","origem":"wiki","rel":"relaciona"}],"confianca":1.0,"contraexemplos":[],"descricao":"O regime reversível (lado c², não-comutativo) da balança, onde o calor a se anula: a forma de Robertson-Schrödinger É a lei a+c²=1 lida na não-comutatividade, e a representação de Hilbert É a própria álgebra dos quaternions (su(2), spin), com Bloch como dinâmica e CHSH=2√2 genuíno (erro 10⁻¹³).","epistemico":"fato","exemplos":[],"id":"mecanica_quantica_de_ouro","memoria":"semantica","meta":{"autor":"Aarão/broca-so","dominio":"broca_repos","fonte":""},"origem":"wiki","palavras_chave":[],"sinonimos":[],"tipo":"conceito","titulo":"Mecânica Quântica de Ouro"}
\section{Mecânica Quântica de Ouro}
O regime reversível (lado c\ensuremath{^{2}}, não-comutativo) da balança, onde o calor a se anula: a forma de Robertson-Schrödinger É a lei a+c\ensuremath{^{2}}=1 lida na não-comutatividade, e a representação de Hilbert É a própria álgebra dos quaternions (su(2), spin), com Bloch como dinâmica e CHSH=2\ensuremath{\sqrt{\,}}2 genuíno (erro 10\ensuremath{^{-13}}).
\prov{relações: usa corpo\_dual; relaciona o\_substrato\_a\_maquina\_e\_funcao\_de\_onda}
\prov{proveniência: wiki \textperiodcentered{} fato \textperiodcentered{} confiança 1.0 \textperiodcentered{} domínio broca\_repos \textperiodcentered{} história 3 versões no jornal}

%CRISTAL {"arestas":[{"alvo":"mecanica_quantica","confianca":0.95,"epistemico":"fato","origem":"wiki","rel":"confirma"},{"alvo":"topologia","confianca":0.5,"epistemico":"fato","origem":"wiki","rel":"relaciona"},{"alvo":"testes","confianca":0.5,"epistemico":"fato","origem":"wiki","rel":"relaciona"},{"alvo":"processos","confianca":0.5,"epistemico":"fato","origem":"wiki","rel":"relaciona"},{"alvo":"recursao","confianca":0.5,"epistemico":"fato","origem":"wiki","rel":"relaciona"},{"alvo":"opcao_a_tarball_release","confianca":0.5,"epistemico":"fato","origem":"wiki","rel":"relaciona"},{"alvo":"ponteiro","confianca":0.5,"epistemico":"fato","origem":"wiki","rel":"relaciona"},{"alvo":"regua_associativa_e_o_risco_de_vazamento","confianca":0.5,"epistemico":"fato","origem":"wiki","rel":"relaciona"},{"alvo":"projeto_cristaljson","confianca":0.5,"epistemico":"fato","origem":"wiki","rel":"relaciona"},{"alvo":"por_sessao_conversadadostelemetriasessaojsonl","confianca":0.5,"epistemico":"fato","origem":"wiki","rel":"relaciona"},{"alvo":"nucleo_broca_vivo","confianca":0.5,"epistemico":"fato","origem":"wiki","rel":"relaciona"},{"alvo":"rust","confianca":0.5,"epistemico":"fato","origem":"wiki","rel":"relaciona"},{"alvo":"syscall","confianca":0.5,"epistemico":"fato","origem":"wiki","rel":"relaciona"},{"alvo":"pow_metal_pre_fecho","confianca":0.5,"epistemico":"fato","origem":"wiki","rel":"relaciona"},{"alvo":"numa","confianca":0.5,"epistemico":"fato","origem":"wiki","rel":"relaciona"},{"alvo":"pow_geometria","confianca":0.5,"epistemico":"fato","origem":"wiki","rel":"relaciona"},{"alvo":"teste_ordem","confianca":0.5,"epistemico":"fato","origem":"wiki","rel":"relaciona"},{"alvo":"pilha","confianca":0.5,"epistemico":"fato","origem":"wiki","rel":"relaciona"},{"alvo":"pipeline_de_ingestao","confianca":0.5,"epistemico":"fato","origem":"wiki","rel":"relaciona"},{"alvo":"motivos","confianca":0.5,"epistemico":"fato","origem":"wiki","rel":"relaciona"},{"alvo":"metais","confianca":0.5,"epistemico":"fato","origem":"wiki","rel":"relaciona"}],"confianca":0.95,"contraexemplos":[],"descricao":"Referência verificada: mecanica_quantica_hub.md","epistemico":"fato","exemplos":[],"id":"mecanica_quantica_hub","memoria":"semantica","meta":{"dominio":"referencia","fonte":"web","nivel":5},"origem":"wiki","palavras_chave":[],"sinonimos":[],"tipo":"conceito","titulo":"mecanica quantica hub"}
\section{mecanica quantica hub}
Referência verificada: mecanica\_quantica\_hub.md
\prov{relações: confirma mecanica\_quantica; relaciona topologia; relaciona testes; relaciona processos; relaciona recursao; relaciona opcao\_a\_tarball\_release; relaciona ponteiro; relaciona regua\_associativa\_e\_o\_risco\_de\_vazamento; relaciona projeto\_cristaljson; relaciona por\_sessao\_conversadadostelemetriasessaojsonl; relaciona nucleo\_broca\_vivo; relaciona rust; relaciona syscall; relaciona pow\_metal\_pre\_fecho; relaciona numa; relaciona pow\_geometria; relaciona teste\_ordem; relaciona pilha; relaciona pipeline\_de\_ingestao; relaciona motivos; relaciona metais}
\prov{proveniência: wiki \textperiodcentered{} web \textperiodcentered{} fato \textperiodcentered{} confiança 0.95 \textperiodcentered{} domínio referencia \textperiodcentered{} história 23 versões no jornal}

%CRISTAL {"arestas":[{"alvo":"mecanica_quantica_hub_mecanica_quantica_hub_fisico_broca","confianca":1.0,"epistemico":"fato","origem":"mecanica_quantica_hub.md","rel":"parte_de"},{"alvo":"word","confianca":0.5,"epistemico":"fato","origem":"wiki","rel":"relaciona"},{"alvo":"vontade_de_poder","confianca":0.5,"epistemico":"fato","origem":"wiki","rel":"relaciona"},{"alvo":"while","confianca":0.5,"epistemico":"fato","origem":"wiki","rel":"relaciona"}],"confianca":1.0,"contraexemplos":[],"descricao":"**Hipótese Dougherty ≠ mecânica quântica standard** — toro φ é conjectura, ψ_nlm é fato.","epistemico":"fato","exemplos":[],"id":"mecanica_quantica_hub_contraexemplos","memoria":"semantica","meta":{"dominio":"referencia","fonte":"web","nivel":5},"origem":"mecanica_quantica_hub.md","palavras_chave":[],"sinonimos":[],"tipo":"conceito","titulo":"Contraexemplos"}
\section{Contraexemplos}
**Hipótese Dougherty \ensuremath{\ne} mecânica quântica standard** — toro \ensuremath{\varphi} é conjectura, \ensuremath{\psi}\_nlm é fato.
\prov{relações: parte\_de mecanica\_quantica\_hub\_mecanica\_quantica\_hub\_fisico\_broca; relaciona word; relaciona vontade\_de\_poder; relaciona while}
\prov{proveniência: mecanica\_quantica\_hub.md \textperiodcentered{} web \textperiodcentered{} fato \textperiodcentered{} confiança 1.0 \textperiodcentered{} domínio referencia \textperiodcentered{} história 6 versões no jornal}

%CRISTAL {"arestas":[{"alvo":"mecanica_quantica_hub_mecanica_quantica_hub_fisico_broca","confianca":1.0,"epistemico":"fato","origem":"mecanica_quantica_hub.md","rel":"parte_de"}],"confianca":1.0,"contraexemplos":[],"descricao":"- `papers/fibracao_hopf_bloch.md` - `papers/microscopia_fotoionizacao_hidrogenio.md` - `papers/fase_geometrica_berry.md` - `papers/ponte_quantica_dougherty.md` (hipótese)","epistemico":"fato","exemplos":[],"id":"mecanica_quantica_hub_fontes","memoria":"semantica","meta":{"dominio":"referencia","fonte":"web","nivel":5},"origem":"mecanica_quantica_hub.md","palavras_chave":[],"sinonimos":[],"tipo":"conceito","titulo":"Fontes"}
\section{Fontes}
- `papers/fibracao\_hopf\_bloch.md` - `papers/microscopia\_fotoionizacao\_hidrogenio.md` - `papers/fase\_geometrica\_berry.md` - `papers/ponte\_quantica\_dougherty.md` (hipótese)
\prov{relações: parte\_de mecanica\_quantica\_hub\_mecanica\_quantica\_hub\_fisico\_broca}
\prov{proveniência: mecanica\_quantica\_hub.md \textperiodcentered{} web \textperiodcentered{} fato \textperiodcentered{} confiança 1.0 \textperiodcentered{} domínio referencia \textperiodcentered{} história 3 versões no jornal}

%CRISTAL {"arestas":[{"alvo":"mecanica_quantica_hub","confianca":0.95,"epistemico":"fato","origem":"wiki","rel":"parte_de"}],"confianca":1.0,"contraexemplos":[],"descricao":"> **Epistemologia:** fato físico-matemático (papers verificados + currículo); hipóteses φ-Dougherty = analogia. > **No grafo:** `epistemico=fato`, confiança 0.95. > **Papel:** fechar dependências, confirmações web e contraexemplos epistemológicos.","epistemico":"fato","exemplos":[],"id":"mecanica_quantica_hub_mecanica_quantica_hub_fisico_broca","memoria":"semantica","meta":{"dominio":"referencia","fonte":"web","nivel":5},"origem":"mecanica_quantica_hub.md","palavras_chave":[],"sinonimos":[],"tipo":"conceito","titulo":"Mecânica quântica — hub físico Broca"}
\section{Mecânica quântica — hub físico Broca}
> **Epistemologia:** fato físico-matemático (papers verificados + currículo); hipóteses \ensuremath{\varphi}-Dougherty = analogia. > **No grafo:** `epistemico=fato`, confiança 0.95. > **Papel:** fechar dependências, confirmações web e contraexemplos epistemológicos.
\prov{relações: parte\_de mecanica\_quantica\_hub}
\prov{proveniência: mecanica\_quantica\_hub.md \textperiodcentered{} web \textperiodcentered{} fato \textperiodcentered{} confiança 1.0 \textperiodcentered{} domínio referencia \textperiodcentered{} história 3 versões no jornal}

%CRISTAL {"arestas":[{"alvo":"mecanica_quantica_hub_mecanica_quantica_hub_fisico_broca","confianca":1.0,"epistemico":"fato","origem":"mecanica_quantica_hub.md","rel":"parte_de"},{"alvo":"matematica","confianca":0.95,"epistemico":"fato","origem":"wiki","rel":"depende"},{"alvo":"mecanica","confianca":0.95,"epistemico":"fato","origem":"wiki","rel":"especializa"},{"alvo":"fisica","confianca":0.95,"epistemico":"fato","origem":"wiki","rel":"parte_de"}],"confianca":1.0,"contraexemplos":[],"descricao":"**Estados, operadores, |ψ|²** — domínio físico-matemático standard.","epistemico":"fato","exemplos":[],"id":"mecanica_quantica_hub_o_que_e","memoria":"semantica","meta":{"dominio":"referencia","fonte":"web","nivel":5},"origem":"mecanica_quantica_hub.md","palavras_chave":[],"sinonimos":[],"tipo":"conceito","titulo":"O que é"}
\section{O que é}
**Estados, operadores, |\ensuremath{\psi}|\ensuremath{^{2}}** — domínio físico-matemático standard.
\prov{relações: parte\_de mecanica\_quantica\_hub\_mecanica\_quantica\_hub\_fisico\_broca; depende matematica; especializa mecanica; parte\_de fisica}
\prov{proveniência: mecanica\_quantica\_hub.md \textperiodcentered{} web \textperiodcentered{} fato \textperiodcentered{} confiança 1.0 \textperiodcentered{} domínio referencia \textperiodcentered{} história 6 versões no jornal}

%CRISTAL {"arestas":[{"alvo":"mecanica_quantica_hub_mecanica_quantica_hub_fisico_broca","confianca":1.0,"epistemico":"fato","origem":"mecanica_quantica_hub.md","rel":"parte_de"},{"alvo":"geometria","confianca":0.95,"epistemico":"fato","origem":"wiki","rel":"referencia"},{"alvo":"quatro_ciencias","confianca":0.95,"epistemico":"fato","origem":"wiki","rel":"parte_de"}],"confianca":1.0,"contraexemplos":[],"descricao":"| MQ | Nó | relação | |----|-----|---------| | núcleo | mecanica_quantica | confirma | | clássica | mecanica | especializa | | matemática | matematica | depende | | física | fisica | parte_de | | geometria | geometria | referencia | | quatro ciências | quatro_ciencias | parte_de |","epistemico":"fato","exemplos":[],"id":"mecanica_quantica_hub_ponte_broca-so","memoria":"semantica","meta":{"dominio":"referencia","fonte":"web","nivel":5},"origem":"mecanica_quantica_hub.md","palavras_chave":[],"sinonimos":[],"tipo":"conceito","titulo":"Ponte Broca-SO"}
\section{Ponte Broca-SO}
| MQ | Nó | relação | |----|-----|---------| | núcleo | mecanica\_quantica | confirma | | clássica | mecanica | especializa | | matemática | matematica | depende | | física | fisica | parte\_de | | geometria | geometria | referencia | | quatro ciências | quatro\_ciencias | parte\_de |
\prov{relações: parte\_de mecanica\_quantica\_hub\_mecanica\_quantica\_hub\_fisico\_broca; referencia geometria; parte\_de quatro\_ciencias}
\prov{proveniência: mecanica\_quantica\_hub.md \textperiodcentered{} web \textperiodcentered{} fato \textperiodcentered{} confiança 1.0 \textperiodcentered{} domínio referencia \textperiodcentered{} história 5 versões no jornal}

%CRISTAL {"arestas":[{"alvo":"mecanica_quantica_hub_mecanica_quantica_hub_fisico_broca","confianca":1.0,"epistemico":"fato","origem":"mecanica_quantica_hub.md","rel":"parte_de"},{"alvo":"ponte_quantica_dougherty","confianca":0.95,"epistemico":"fato","origem":"wiki","rel":"analogia"},{"alvo":"ginzburg_toro_helice","confianca":0.95,"epistemico":"fato","origem":"wiki","rel":"analogia"}],"confianca":1.0,"contraexemplos":[],"descricao":"| Hipótese | Nó | relação | |----------|-----|---------| | Dougherty φ | ponte_quantica_dougherty | analogia | | Ginzburg 3D-SST | ginzburg_toro_helice | analogia |","epistemico":"fato","exemplos":[],"id":"mecanica_quantica_hub_ponte_conjectural_hipotese","memoria":"semantica","meta":{"dominio":"referencia","fonte":"web","nivel":5},"origem":"mecanica_quantica_hub.md","palavras_chave":[],"sinonimos":[],"tipo":"conceito","titulo":"Ponte conjectural (hipótese)"}
\section{Ponte conjectural (hipótese)}
| Hipótese | Nó | relação | |----------|-----|---------| | Dougherty \ensuremath{\varphi} | ponte\_quantica\_dougherty | analogia | | Ginzburg 3D-SST | ginzburg\_toro\_helice | analogia |
\prov{relações: parte\_de mecanica\_quantica\_hub\_mecanica\_quantica\_hub\_fisico\_broca; analogia ponte\_quantica\_dougherty; analogia ginzburg\_toro\_helice}
\prov{proveniência: mecanica\_quantica\_hub.md \textperiodcentered{} web \textperiodcentered{} fato \textperiodcentered{} confiança 1.0 \textperiodcentered{} domínio referencia \textperiodcentered{} história 5 versões no jornal}

%CRISTAL {"arestas":[{"alvo":"mecanica_quantica_hub_mecanica_quantica_hub_fisico_broca","confianca":1.0,"epistemico":"fato","origem":"mecanica_quantica_hub.md","rel":"parte_de"},{"alvo":"fibracao_hopf_bloch","confianca":0.95,"epistemico":"fato","origem":"wiki","rel":"confirma"},{"alvo":"microscopia_fotoionizacao_hidrogenio","confianca":0.95,"epistemico":"fato","origem":"wiki","rel":"confirma"},{"alvo":"fase_geometrica_berry","confianca":0.95,"epistemico":"fato","origem":"wiki","rel":"confirma"},{"alvo":"gato","confianca":0.95,"epistemico":"fato","origem":"wiki","rel":"analogia"},{"alvo":"cosmologia_quantica","confianca":0.95,"epistemico":"fato","origem":"wiki","rel":"referencia"}],"confianca":1.0,"contraexemplos":[],"descricao":"| Referência | Nó | relação | |------------|-----|---------| | Hopf/Bloch | fibracao_hopf_bloch | confirma | | H fotoionização | microscopia_fotoionizacao_hidrogenio | confirma | | Berry | fase_geometrica_berry | confirma | | qubit | gato | analogia | | cosmologia | cosmologia_quantica | referencia |","epistemico":"fato","exemplos":[],"id":"mecanica_quantica_hub_pontes_verificadas_fato","memoria":"semantica","meta":{"dominio":"referencia","fonte":"web","nivel":5},"origem":"mecanica_quantica_hub.md","palavras_chave":[],"sinonimos":[],"tipo":"conceito","titulo":"Pontes verificadas (fato)"}
\section{Pontes verificadas (fato)}
| Referência | Nó | relação | |------------|-----|---------| | Hopf/Bloch | fibracao\_hopf\_bloch | confirma | | H fotoionização | microscopia\_fotoionizacao\_hidrogenio | confirma | | Berry | fase\_geometrica\_berry | confirma | | qubit | gato | analogia | | cosmologia | cosmologia\_quantica | referencia |
\prov{relações: parte\_de mecanica\_quantica\_hub\_mecanica\_quantica\_hub\_fisico\_broca; confirma fibracao\_hopf\_bloch; confirma microscopia\_fotoionizacao\_hidrogenio; confirma fase\_geometrica\_berry; analogia gato; referencia cosmologia\_quantica}
\prov{proveniência: mecanica\_quantica\_hub.md \textperiodcentered{} web \textperiodcentered{} fato \textperiodcentered{} confiança 1.0 \textperiodcentered{} domínio referencia \textperiodcentered{} história 8 versões no jornal}

%CRISTAL {"arestas":[{"alvo":"mecanica_quantica_hub_mecanica_quantica_hub_fisico_broca","confianca":1.0,"epistemico":"fato","origem":"mecanica_quantica_hub.md","rel":"parte_de"}],"confianca":1.0,"contraexemplos":[],"descricao":"- **Factos web:** `papers/fibracao_hopf_bloch.md`, `papers/microscopia_fotoionizacao_hidrogenio.md`, `papers/fase_geometrica_berry.md` - **Hipótese (≤0.55):** `papers/ponte_quantica_dougherty.md` — analogia, não axioma - **Currículo:** fase física — mecânica → MQ → cosmologia quântica - **Broca:** operador gato A, Painel Estelar — **analogia** qubit/Bloch","epistemico":"fato","exemplos":[],"id":"mecanica_quantica_hub_referencia","memoria":"semantica","meta":{"dominio":"referencia","fonte":"web","nivel":5},"origem":"mecanica_quantica_hub.md","palavras_chave":[],"sinonimos":[],"tipo":"conceito","titulo":"Referência"}
\section{Referência}
- **Factos web:** `papers/fibracao\_hopf\_bloch.md`, `papers/microscopia\_fotoionizacao\_hidrogenio.md`, `papers/fase\_geometrica\_berry.md` - **Hipótese (\ensuremath{\le}0.55):** `papers/ponte\_quantica\_dougherty.md` — analogia, não axioma - **Currículo:** fase física — mecânica \ensuremath{\to} MQ \ensuremath{\to} cosmologia quântica - **Broca:** operador gato A, Painel Estelar — **analogia** qubit/Bloch
\prov{relações: parte\_de mecanica\_quantica\_hub\_mecanica\_quantica\_hub\_fisico\_broca}
\prov{proveniência: mecanica\_quantica\_hub.md \textperiodcentered{} web \textperiodcentered{} fato \textperiodcentered{} confiança 1.0 \textperiodcentered{} domínio referencia \textperiodcentered{} história 3 versões no jornal}

%CRISTAL {"arestas":[{"alvo":"gato","confianca":0.95,"epistemico":"fato","origem":"paper","rel":"parte_de"},{"alvo":"o_que_e_no_projecto","confianca":0.5,"epistemico":"fato","origem":"wiki","rel":"relaciona"},{"alvo":"testes","confianca":0.5,"epistemico":"fato","origem":"wiki","rel":"relaciona"},{"alvo":"utilitarismo","confianca":0.5,"epistemico":"fato","origem":"wiki","rel":"relaciona"},{"alvo":"recursao","confianca":0.5,"epistemico":"fato","origem":"wiki","rel":"relaciona"},{"alvo":"nucleo_broca_vivo","confianca":0.5,"epistemico":"fato","origem":"wiki","rel":"relaciona"}],"confianca":1.0,"contraexemplos":[],"descricao":"- nonce 1,966,079 · acoplamento 0.999095 · atrator 672 · fecho=True · sha=False","epistemico":"fato","exemplos":[],"id":"melhor_candidato","memoria":"semantica","meta":{"dominio":"manual","nivel":6,"verificacao":"slug idêntico — enriquecer, não duplicar"},"origem":"POW_METAL_RESSONANCIA.md","palavras_chave":[],"sinonimos":[],"tipo":"conceito","titulo":"Melhor candidato"}
\section{Melhor candidato}
- nonce 1,966,079 \textperiodcentered  acoplamento 0.999095 \textperiodcentered  atrator 672 \textperiodcentered  fecho=True \textperiodcentered  sha=False
\prov{relações: parte\_de gato; relaciona o\_que\_e\_no\_projecto; relaciona testes; relaciona utilitarismo; relaciona recursao; relaciona nucleo\_broca\_vivo}
\prov{proveniência: POW\_METAL\_RESSONANCIA.md \textperiodcentered{} fato \textperiodcentered{} confiança 1.0 \textperiodcentered{} domínio manual \textperiodcentered{} história 8 versões no jornal}

%CRISTAL {"arestas":[{"alvo":"koch_como_recipiente_memoria_nao-associativa","confianca":1.0,"epistemico":"fato","origem":"KOCH_MEMORIA.md","rel":"parte_de"}],"confianca":1.0,"contraexemplos":[],"descricao":"``` Célula branca e₀  →  Word [total,e]     (conteúdo vivo, associativo via gold) Célula branca e₁  →  gold(w)            (Fibonacci, associativo) Canal negro E₀₁   →  parábola a+c²=1   (ressonância — parâmetro s, Cantor/Fib) Recipiente Koch   →  órbita z² mod P    (memória do que aconteceu no canal) ```","epistemico":"fato","exemplos":[],"id":"metafora_peirce","memoria":"semantica","meta":{"dominio":"manual","nivel":6,"verificacao":"slug idêntico — enriquecer, não duplicar"},"origem":"KOCH_MEMORIA.md","palavras_chave":[],"sinonimos":[],"tipo":"conceito","titulo":"Metáfora Peirce"}
\section{Metáfora Peirce}
``` Célula branca e\ensuremath{_{0}}  \ensuremath{\to}  Word [total,e]     (conteúdo vivo, associativo via gold) Célula branca e\ensuremath{_{1}}  \ensuremath{\to}  gold(w)            (Fibonacci, associativo) Canal negro E\ensuremath{_{01}}   \ensuremath{\to}  parábola a+c\ensuremath{^{2}}=1   (ressonância — parâmetro s, Cantor/Fib) Recipiente Koch   \ensuremath{\to}  órbita z\ensuremath{^{2}} mod P    (memória do que aconteceu no canal) ```
\prov{relações: parte\_de koch\_como\_recipiente\_memoria\_nao-associativa}
\prov{proveniência: KOCH\_MEMORIA.md \textperiodcentered{} fato \textperiodcentered{} confiança 1.0 \textperiodcentered{} domínio manual \textperiodcentered{} história 6 versões no jornal}

%CRISTAL {"arestas":[{"alvo":"corpo_dual","confianca":1.0,"epistemico":"fato","origem":"manual","rel":"referencia"},{"alvo":"gato","confianca":1.0,"epistemico":"fato","origem":"paper","rel":"usa"},{"alvo":"readme","confianca":0.5,"epistemico":"fato","origem":"wiki","rel":"relaciona"},{"alvo":"vida-morte","confianca":0.5,"epistemico":"fato","origem":"wiki","rel":"relaciona"},{"alvo":"teste_ordem","confianca":0.5,"epistemico":"fato","origem":"wiki","rel":"relaciona"},{"alvo":"testes","confianca":0.5,"epistemico":"fato","origem":"wiki","rel":"relaciona"},{"alvo":"virtualizacao","confianca":0.5,"epistemico":"fato","origem":"wiki","rel":"relaciona"},{"alvo":"transcricao","confianca":0.5,"epistemico":"fato","origem":"wiki","rel":"relaciona"},{"alvo":"projeto_cristaljson","confianca":0.5,"epistemico":"fato","origem":"wiki","rel":"relaciona"},{"alvo":"motivos","confianca":0.5,"epistemico":"fato","origem":"wiki","rel":"relaciona"},{"alvo":"pow_metal_pre_fecho","confianca":0.5,"epistemico":"fato","origem":"wiki","rel":"relaciona"},{"alvo":"thread","confianca":0.5,"epistemico":"fato","origem":"wiki","rel":"relaciona"},{"alvo":"pilha","confianca":0.5,"epistemico":"fato","origem":"wiki","rel":"relaciona"},{"alvo":"opcao_a_tarball_release","confianca":0.5,"epistemico":"fato","origem":"wiki","rel":"relaciona"},{"alvo":"rust","confianca":0.5,"epistemico":"fato","origem":"wiki","rel":"relaciona"},{"alvo":"parabola_pa","confianca":0.5,"epistemico":"fato","origem":"wiki","rel":"relaciona"},{"alvo":"topologia","confianca":0.5,"epistemico":"fato","origem":"wiki","rel":"relaciona"},{"alvo":"parabola_destruir","confianca":0.5,"epistemico":"fato","origem":"wiki","rel":"relaciona"},{"alvo":"ponteiro","confianca":0.5,"epistemico":"fato","origem":"wiki","rel":"relaciona"},{"alvo":"area_metalurgia","confianca":1.0,"epistemico":"fato","origem":"wiki","rel":"parte_de"}],"confianca":1.0,"contraexemplos":[],"descricao":"Manual: ula/METAIS.md","epistemico":"fato","exemplos":[],"id":"metais","memoria":"semantica","meta":{"arquivo":"ula/METAIS.md","dominio":"manual","nivel":6},"origem":"manual","palavras_chave":[],"sinonimos":[],"tipo":"conceito","titulo":"METAIS"}
\section{METAIS}
Manual: ula/METAIS.md
\prov{relações: referencia corpo\_dual; usa gato; relaciona readme; relaciona vida-morte; relaciona teste\_ordem; relaciona testes; relaciona virtualizacao; relaciona transcricao; relaciona projeto\_cristaljson; relaciona motivos; relaciona pow\_metal\_pre\_fecho; relaciona thread; relaciona pilha; relaciona opcao\_a\_tarball\_release; relaciona rust; relaciona parabola\_pa; relaciona topologia; relaciona parabola\_destruir; relaciona ponteiro; parte\_de area\_metalurgia}
\prov{proveniência: manual \textperiodcentered{} fato \textperiodcentered{} confiança 1.0 \textperiodcentered{} domínio manual \textperiodcentered{} história 30 versões no jornal}

%CRISTAL {"arestas":[{"alvo":"falsificabilidade","confianca":0.95,"epistemico":"fato","origem":"wiki","rel":"confirma"},{"alvo":"mecanica_quantica","confianca":1.0,"epistemico":"fato","origem":"wiki","rel":"confirma"},{"alvo":"orbitais_de_hidrogenio_evidencia_empirica","confianca":1.0,"epistemico":"fato","origem":"wiki","rel":"confirma"},{"alvo":"ponte_quantica_dougherty","confianca":1.0,"epistemico":"fato","origem":"wiki","rel":"referencia"}],"confianca":0.95,"contraexemplos":[],"descricao":"Referência verificada: microscopia_fotoionizacao_hidrogenio.md","epistemico":"fato","exemplos":[],"id":"microscopia_fotoionizacao_hidrogenio","memoria":"semantica","meta":{"dominio":"referencia","fonte":"web","nivel":5},"origem":"wiki","palavras_chave":[],"sinonimos":[],"tipo":"conceito","titulo":"microscopia fotoionizacao hidrogenio"}
\section{microscopia fotoionizacao hidrogenio}
Referência verificada: microscopia\_fotoionizacao\_hidrogenio.md
\prov{relações: confirma falsificabilidade; confirma mecanica\_quantica; confirma orbitais\_de\_hidrogenio\_evidencia\_empirica; referencia ponte\_quantica\_dougherty}
\prov{proveniência: wiki \textperiodcentered{} web \textperiodcentered{} fato \textperiodcentered{} confiança 0.95 \textperiodcentered{} domínio referencia \textperiodcentered{} história 100 versões no jornal}

%CRISTAL {"arestas":[{"alvo":"microscopia_fotoionizacao_hidrogenio_microscopia_de_fotoionizacao_orbitais_de_hi","confianca":1.0,"epistemico":"fato","origem":"microscopia_fotoionizacao_hidrogenio.md","rel":"parte_de"}],"confianca":1.0,"contraexemplos":[],"descricao":"Interpretações “electric universe” ou toros empíricos **sem** predizer nodos Stark medidos são metáfora, não física atómica verificada.","epistemico":"fato","exemplos":[],"id":"microscopia_fotoionizacao_hidrogenio_contraexemplos","memoria":"semantica","meta":{"dominio":"referencia","fonte":"web","nivel":5},"origem":"microscopia_fotoionizacao_hidrogenio.md","palavras_chave":[],"sinonimos":[],"tipo":"conceito","titulo":"Contraexemplos"}
\section{Contraexemplos}
Interpretações “electric universe” ou toros empíricos **sem** predizer nodos Stark medidos são metáfora, não física atómica verificada.
\prov{relações: parte\_de microscopia\_fotoionizacao\_hidrogenio\_microscopia\_de\_fotoionizacao\_orbitais\_de\_hi}
\prov{proveniência: microscopia\_fotoionizacao\_hidrogenio.md \textperiodcentered{} web \textperiodcentered{} fato \textperiodcentered{} confiança 1.0 \textperiodcentered{} domínio referencia \textperiodcentered{} história 3 versões no jornal}

%CRISTAL {"fusao":[{"arestas":[{"alvo":"microscopia_fotoionizacao_hidrogenio_microscopia_de_fotoionizacao_orbitais_de_hi","confianca":1.0,"epistemico":"fato","origem":"microscopia_fotoionizacao_hidrogenio.md","rel":"parte_de"},{"alvo":"mecanica_quantica","confianca":0.95,"epistemico":"fato","origem":"wiki","rel":"prova"}],"confianca":1.0,"contraexemplos":[],"descricao":"ψₙlm(r,θ,φ) = Rₙl(r) Ylm(θ,φ) — solução analítica do átomo de hidrogénio; Schrödinger + Coulomb. Não requer geometria toroidal.","epistemico":"fato","exemplos":[],"id":"microscopia_fotoionizacao_hidrogenio_equacao_de_referencia_mq_standard","memoria":"semantica","meta":{"dominio":"referencia","fonte":"web","nivel":5},"origem":"microscopia_fotoionizacao_hidrogenio.md","palavras_chave":[],"sinonimos":[],"tipo":"conceito","titulo":"Equação de referência (MQ standard)"},{"arestas":[{"alvo":"microscopia_de_fotoionizacao_orbitais_de_hidrogenio_stodolna_2013","confianca":1.0,"epistemico":"fato","origem":"microscopia_fotoionizacao_hidrogenio.md","rel":"parte_de"},{"alvo":"mecanica_quantica","confianca":0.95,"epistemico":"fato","origem":"wiki","rel":"prova"}],"confianca":1.0,"contraexemplos":[],"descricao":"ψₙlm(r,θ,φ) = Rₙl(r) Ylm(θ,φ) — solução analítica do átomo de hidrogénio; Schrödinger + Coulomb. Não requer geometria toroidal.","epistemico":"fato","exemplos":[],"id":"equacao_de_referencia_mq_standard","memoria":"semantica","meta":{"dominio":"referencia","nivel":5},"origem":"microscopia_fotoionizacao_hidrogenio.md","palavras_chave":[],"sinonimos":[],"tipo":"conceito","titulo":"Equação de referência (MQ standard)"}],"id":"microscopia_fotoionizacao_hidrogenio_equacao_de_referencia_mq_standard","tipo":"conceito"}
\section{Equação de referência (MQ standard)}
\ensuremath{\psi}\ensuremath{_{n}}lm(r,\ensuremath{\theta},\ensuremath{\varphi}) = R\ensuremath{_{n}}l(r) Ylm(\ensuremath{\theta},\ensuremath{\varphi}) — solução analítica do átomo de hidrogénio; Schrödinger + Coulomb. Não requer geometria toroidal.
\prov{relações: parte\_de microscopia\_fotoionizacao\_hidrogenio\_microscopia\_de\_fotoionizacao\_orbitais\_de\_hi; prova mecanica\_quantica}
\prov{proveniência: microscopia\_fotoionizacao\_hidrogenio.md \textperiodcentered{} web \textperiodcentered{} fato \textperiodcentered{} confiança 1.0 \textperiodcentered{} domínio referencia \textperiodcentered{} história 5 versões no jornal}
\prov{fusão de curadoria (texto idêntico): absorve equacao\_de\_referencia\_mq\_standard --- o mesmo conteúdo por dois esquemas de endereço}

%CRISTAL {"fusao":[{"arestas":[{"alvo":"microscopia_fotoionizacao_hidrogenio_microscopia_de_fotoionizacao_orbitais_de_hi","confianca":1.0,"epistemico":"fato","origem":"microscopia_fotoionizacao_hidrogenio.md","rel":"parte_de"}],"confianca":1.0,"contraexemplos":[],"descricao":"- https://physics.aps.org/articles/v6/58 - https://journals.aps.org/prl/abstract/10.1103/PhysRevLett.110.213001 - https://mbi-berlin.de/research/highlights/details/hydrogen-atoms-under-the-microscope","epistemico":"fato","exemplos":[],"id":"microscopia_fotoionizacao_hidrogenio_fontes_web","memoria":"semantica","meta":{"dominio":"referencia","fonte":"web","nivel":5},"origem":"microscopia_fotoionizacao_hidrogenio.md","palavras_chave":[],"sinonimos":[],"tipo":"conceito","titulo":"Fontes web"},{"arestas":[{"alvo":"microscopia_de_fotoionizacao_orbitais_de_hidrogenio_stodolna_2013","confianca":1.0,"epistemico":"fato","origem":"microscopia_fotoionizacao_hidrogenio.md","rel":"parte_de"}],"confianca":1.0,"contraexemplos":[],"descricao":"- https://physics.aps.org/articles/v6/58 - https://journals.aps.org/prl/abstract/10.1103/PhysRevLett.110.213001 - https://mbi-berlin.de/research/highlights/details/hydrogen-atoms-under-the-microscope","epistemico":"fato","exemplos":[],"id":"fontes_web","memoria":"semantica","meta":{"dominio":"referencia","nivel":5},"origem":"microscopia_fotoionizacao_hidrogenio.md","palavras_chave":[],"sinonimos":[],"tipo":"conceito","titulo":"Fontes web"}],"id":"microscopia_fotoionizacao_hidrogenio_fontes_web","tipo":"conceito"}
\section{Fontes web}
- https://physics.aps.org/articles/v6/58 - https://journals.aps.org/prl/abstract/10.1103/PhysRevLett.110.213001 - https://mbi-berlin.de/research/highlights/details/hydrogen-atoms-under-the-microscope
\prov{relações: parte\_de microscopia\_fotoionizacao\_hidrogenio\_microscopia\_de\_fotoionizacao\_orbitais\_de\_hi}
\prov{proveniência: microscopia\_fotoionizacao\_hidrogenio.md \textperiodcentered{} web \textperiodcentered{} fato \textperiodcentered{} confiança 1.0 \textperiodcentered{} domínio referencia \textperiodcentered{} história 3 versões no jornal}
\prov{fusão de curadoria (texto idêntico): absorve fontes\_web --- o mesmo conteúdo por dois esquemas de endereço}

%CRISTAL {"fusao":[{"arestas":[{"alvo":"microscopia_fotoionizacao_hidrogenio_microscopia_de_fotoionizacao_orbitais_de_hi","confianca":1.0,"epistemico":"fato","origem":"microscopia_fotoionizacao_hidrogenio.md","rel":"parte_de"},{"alvo":"mecanica_quantica","confianca":0.95,"epistemico":"fato","origem":"wiki","rel":"usa"}],"confianca":1.0,"contraexemplos":[],"descricao":"1. Feixe de átomos H preparados por dissociação laser (H₂S). 2. Excitação two-photon: 2s/2p → estado de Rydberg (n≈30, nodos controláveis). 3. Campo eléctrico DC: Hamiltoniano de Stark separável em coordenadas parabólicas (η, ξ, φ). 4. Fotoionização + **lente electrostática** (~20 000×) preservando coerência quântica. 5. Detector 2D (MCP): padrão de interferência = estrutura nodal de |ψ|².","epistemico":"fato","exemplos":[],"id":"microscopia_fotoionizacao_hidrogenio_metodo","memoria":"semantica","meta":{"dominio":"referencia","fonte":"web","nivel":5},"origem":"microscopia_fotoionizacao_hidrogenio.md","palavras_chave":[],"sinonimos":[],"tipo":"conceito","titulo":"Método"},{"arestas":[{"alvo":"microscopia_de_fotoionizacao_orbitais_de_hidrogenio_stodolna_2013","confianca":1.0,"epistemico":"fato","origem":"microscopia_fotoionizacao_hidrogenio.md","rel":"parte_de"},{"alvo":"mecanica_quantica","confianca":0.95,"epistemico":"fato","origem":"wiki","rel":"usa"}],"confianca":1.0,"contraexemplos":[],"descricao":"1. Feixe de átomos H preparados por dissociação laser (H₂S). 2. Excitação two-photon: 2s/2p → estado de Rydberg (n≈30, nodos controláveis). 3. Campo eléctrico DC: Hamiltoniano de Stark separável em coordenadas parabólicas (η, ξ, φ). 4. Fotoionização + **lente electrostática** (~20 000×) preservando coerência quântica. 5. Detector 2D (MCP): padrão de interferência = estrutura nodal de |ψ|².","epistemico":"fato","exemplos":[],"id":"metodo","memoria":"semantica","meta":{"dominio":"referencia","nivel":5},"origem":"microscopia_fotoionizacao_hidrogenio.md","palavras_chave":[],"sinonimos":[],"tipo":"conceito","titulo":"Método"}],"id":"microscopia_fotoionizacao_hidrogenio_metodo","tipo":"conceito"}
\section{Método}
1. Feixe de átomos H preparados por dissociação laser (H\ensuremath{_{2}}S). 2. Excitação two-photon: 2s/2p \ensuremath{\to} estado de Rydberg (n\ensuremath{\approx}30, nodos controláveis). 3. Campo eléctrico DC: Hamiltoniano de Stark separável em coordenadas parabólicas (\ensuremath{\eta}, \ensuremath{\xi}, \ensuremath{\varphi}). 4. Fotoionização + **lente electrostática** (\~{}20 000$\times$) preservando coerência quântica. 5. Detector 2D (MCP): padrão de interferência = estrutura nodal de |\ensuremath{\psi}|\ensuremath{^{2}}.
\prov{relações: parte\_de microscopia\_fotoionizacao\_hidrogenio\_microscopia\_de\_fotoionizacao\_orbitais\_de\_hi; usa mecanica\_quantica}
\prov{proveniência: microscopia\_fotoionizacao\_hidrogenio.md \textperiodcentered{} web \textperiodcentered{} fato \textperiodcentered{} confiança 1.0 \textperiodcentered{} domínio referencia \textperiodcentered{} história 4 versões no jornal}
\prov{fusão de curadoria (texto idêntico): absorve metodo --- o mesmo conteúdo por dois esquemas de endereço}

%CRISTAL {"fusao":[{"arestas":[{"alvo":"microscopia_fotoionizacao_hidrogenio","confianca":0.95,"epistemico":"fato","origem":"wiki","rel":"parte_de"}],"confianca":1.0,"contraexemplos":[],"descricao":"> **Epistemologia:** fato experimental publicado (PRL, replicável). > **No grafo:** `epistemico=fato`, confiança 0.95. > **Papel:** ancora empírica para MQ standard; teste de falsificação da hipótese Dougherty.","epistemico":"fato","exemplos":[],"id":"microscopia_fotoionizacao_hidrogenio_microscopia_de_fotoionizacao_orbitais_de_hi","memoria":"semantica","meta":{"dominio":"referencia","fonte":"web","nivel":5},"origem":"microscopia_fotoionizacao_hidrogenio.md","palavras_chave":[],"sinonimos":[],"tipo":"conceito","titulo":"Microscopia de fotoionização — orbitais de hidrogénio (Stodolna 2013)"},{"arestas":[{"alvo":"microscopia_fotoionizacao_hidrogenio","confianca":0.95,"epistemico":"fato","origem":"wiki","rel":"parte_de"}],"confianca":1.0,"contraexemplos":[],"descricao":"> **Epistemologia:** fato experimental publicado (PRL, replicável). > **No grafo:** `epistemico=fato`, confiança 0.95. > **Papel:** ancora empírica para MQ standard; teste de falsificação da hipótese Dougherty.","epistemico":"fato","exemplos":[],"id":"microscopia_de_fotoionizacao_orbitais_de_hidrogenio_stodolna_2013","memoria":"semantica","meta":{"dominio":"referencia","nivel":5},"origem":"microscopia_fotoionizacao_hidrogenio.md","palavras_chave":[],"sinonimos":[],"tipo":"conceito","titulo":"Microscopia de fotoionização — orbitais de hidrogénio (Stodolna 2013)"}],"id":"microscopia_fotoionizacao_hidrogenio_microscopia_de_fotoionizacao_orbitais_de_hi","tipo":"conceito"}
\section{Microscopia de fotoionização — orbitais de hidrogénio (Stodolna 2013)}
> **Epistemologia:** fato experimental publicado (PRL, replicável). > **No grafo:** `epistemico=fato`, confiança 0.95. > **Papel:** ancora empírica para MQ standard; teste de falsificação da hipótese Dougherty.
\prov{relações: parte\_de microscopia\_fotoionizacao\_hidrogenio}
\prov{proveniência: microscopia\_fotoionizacao\_hidrogenio.md \textperiodcentered{} web \textperiodcentered{} fato \textperiodcentered{} confiança 1.0 \textperiodcentered{} domínio referencia \textperiodcentered{} história 3 versões no jornal}
\prov{fusão de curadoria (texto idêntico): absorve microscopia\_de\_fotoionizacao\_orbitais\_de\_hidrogenio\_stodolna\_2013 --- o mesmo conteúdo por dois esquemas de endereço}

%CRISTAL {"arestas":[{"alvo":"microscopia_fotoionizacao_hidrogenio_microscopia_de_fotoionizacao_orbitais_de_hi","confianca":1.0,"epistemico":"fato","origem":"microscopia_fotoionizacao_hidrogenio.md","rel":"parte_de"},{"alvo":"colapso","confianca":0.95,"epistemico":"fato","origem":"wiki","rel":"analogia"},{"alvo":"geometria","confianca":0.95,"epistemico":"fato","origem":"wiki","rel":"referencia"}],"confianca":1.0,"contraexemplos":[],"descricao":"| Aspecto | Ligação no grafo | |---------|------------------| | |ψ|² observável | mecânica quântica — confirma | | nodos / lóbulos | orbitais hidrogénio (hipótese Dougherty) — confirma base empírica | | microscópio quântico | colapso / medição — analogia (leitura projectiva) | | coordenadas parabólicas | geometria — referencia (não toro) |","epistemico":"fato","exemplos":[],"id":"microscopia_fotoionizacao_hidrogenio_ponte_broca-so","memoria":"semantica","meta":{"dominio":"referencia","fonte":"web","nivel":5},"origem":"microscopia_fotoionizacao_hidrogenio.md","palavras_chave":[],"sinonimos":[],"tipo":"conceito","titulo":"Ponte Broca-SO"}
\section{Ponte Broca-SO}
| Aspecto | Ligação no grafo | |---------|------------------| | |\ensuremath{\psi}|\ensuremath{^{2}} observável | mecânica quântica — confirma | | nodos / lóbulos | orbitais hidrogénio (hipótese Dougherty) — confirma base empírica | | microscópio quântico | colapso / medição — analogia (leitura projectiva) | | coordenadas parabólicas | geometria — referencia (não toro) |
\prov{relações: parte\_de microscopia\_fotoionizacao\_hidrogenio\_microscopia\_de\_fotoionizacao\_orbitais\_de\_hi; analogia colapso; referencia geometria}
\prov{proveniência: microscopia\_fotoionizacao\_hidrogenio.md \textperiodcentered{} web \textperiodcentered{} fato \textperiodcentered{} confiança 1.0 \textperiodcentered{} domínio referencia \textperiodcentered{} história 5 versões no jornal}

%CRISTAL {"arestas":[{"alvo":"microscopia_fotoionizacao_hidrogenio_microscopia_de_fotoionizacao_orbitais_de_hi","confianca":1.0,"epistemico":"fato","origem":"microscopia_fotoionizacao_hidrogenio.md","rel":"parte_de"}],"confianca":1.0,"contraexemplos":[],"descricao":"- **Título:** Hydrogen Atoms under Magnification: Direct Observation of the Nodal Structure of Stark States - **Autores:** Aneta Stodolna, Marc Vrakking et al. (AMOLF, Max-Born-Institut) - **Revista:** Physical Review Letters **110**, 213001 (2013) - **DOI:** [10.1103/PhysRevLett.110.213001](https://doi.org/10.1103/PhysRevLett.110.213001) - **Viewpoint APS:** [Physics 6, 58 — A New Look at the Hydrogen Wave Function](https://physics.aps.org/articles/v6/58)","epistemico":"fato","exemplos":[],"id":"microscopia_fotoionizacao_hidrogenio_referencia","memoria":"semantica","meta":{"dominio":"referencia","fonte":"web","nivel":5},"origem":"microscopia_fotoionizacao_hidrogenio.md","palavras_chave":[],"sinonimos":[],"tipo":"conceito","titulo":"Referência"}
\section{Referência}
- **Título:** Hydrogen Atoms under Magnification: Direct Observation of the Nodal Structure of Stark States - **Autores:** Aneta Stodolna, Marc Vrakking et al. (AMOLF, Max-Born-Institut) - **Revista:** Physical Review Letters **110**, 213001 (2013) - **DOI:** [10.1103/PhysRevLett.110.213001](https://doi.org/10.1103/PhysRevLett.110.213001) - **Viewpoint APS:** [Physics 6, 58 — A New Look at the Hydrogen Wave Function](https://physics.aps.org/articles/v6/58)
\prov{relações: parte\_de microscopia\_fotoionizacao\_hidrogenio\_microscopia\_de\_fotoionizacao\_orbitais\_de\_hi}
\prov{proveniência: microscopia\_fotoionizacao\_hidrogenio.md \textperiodcentered{} web \textperiodcentered{} fato \textperiodcentered{} confiança 1.0 \textperiodcentered{} domínio referencia \textperiodcentered{} história 3 versões no jornal}

%CRISTAL {"arestas":[{"alvo":"microscopia_fotoionizacao_hidrogenio_microscopia_de_fotoionizacao_orbitais_de_hi","confianca":1.0,"epistemico":"fato","origem":"microscopia_fotoionizacao_hidrogenio.md","rel":"parte_de"},{"alvo":"ponte_quantica_dougherty","confianca":0.95,"epistemico":"fato","origem":"wiki","rel":"referencia"},{"alvo":"conjunto_dougherty","confianca":0.55,"epistemico":"inferencia","origem":"wiki","rel":"contradiz"}],"confianca":1.0,"contraexemplos":[],"descricao":"A microscopia **confirma** a topologia nodal de ψ que qualquer ponte toroidal deve reproduzir. **Falsificação:** se equação toroidal Dougherty não reproduzir contagem nodal e padrões Betti das imagens PRL, a ponte cai na escala atómica.","epistemico":"fato","exemplos":[],"id":"microscopia_fotoionizacao_hidrogenio_relacao_com_hipotese_dougherty","memoria":"semantica","meta":{"dominio":"referencia","fonte":"web","nivel":5},"origem":"microscopia_fotoionizacao_hidrogenio.md","palavras_chave":[],"sinonimos":[],"tipo":"conceito","titulo":"Relação com hipótese Dougherty"}
\section{Relação com hipótese Dougherty}
A microscopia **confirma** a topologia nodal de \ensuremath{\psi} que qualquer ponte toroidal deve reproduzir. **Falsificação:** se equação toroidal Dougherty não reproduzir contagem nodal e padrões Betti das imagens PRL, a ponte cai na escala atómica.
\prov{relações: parte\_de microscopia\_fotoionizacao\_hidrogenio\_microscopia\_de\_fotoionizacao\_orbitais\_de\_hi; referencia ponte\_quantica\_dougherty; contradiz conjunto\_dougherty}
\prov{proveniência: microscopia\_fotoionizacao\_hidrogenio.md \textperiodcentered{} web \textperiodcentered{} fato \textperiodcentered{} confiança 1.0 \textperiodcentered{} domínio referencia \textperiodcentered{} história 5 versões no jornal}

%CRISTAL {"fusao":[{"arestas":[{"alvo":"microscopia_fotoionizacao_hidrogenio_microscopia_de_fotoionizacao_orbitais_de_hi","confianca":1.0,"epistemico":"fato","origem":"microscopia_fotoionizacao_hidrogenio.md","rel":"parte_de"},{"alvo":"orbitais_de_hidrogenio_evidencia_empirica","confianca":0.95,"epistemico":"fato","origem":"wiki","rel":"confirma"}],"confianca":1.0,"contraexemplos":[],"descricao":"A função de onda microscópica ao longo de ξ e a projeção do continuo medida a distância macroscópica **partilham a mesma estrutura nodal**. Estados com 0, 1, 2, 3 nodos observados directamente na câmara.","epistemico":"fato","exemplos":[],"id":"microscopia_fotoionizacao_hidrogenio_resultado-chave","memoria":"semantica","meta":{"dominio":"referencia","fonte":"web","nivel":5},"origem":"microscopia_fotoionizacao_hidrogenio.md","palavras_chave":[],"sinonimos":[],"tipo":"conceito","titulo":"Resultado-chave"},{"arestas":[{"alvo":"microscopia_de_fotoionizacao_orbitais_de_hidrogenio_stodolna_2013","confianca":1.0,"epistemico":"fato","origem":"microscopia_fotoionizacao_hidrogenio.md","rel":"parte_de"},{"alvo":"orbitais_de_hidrogenio_evidencia_empirica","confianca":0.95,"epistemico":"fato","origem":"wiki","rel":"confirma"}],"confianca":1.0,"contraexemplos":[],"descricao":"A função de onda microscópica ao longo de ξ e a projeção do continuo medida a distância macroscópica **partilham a mesma estrutura nodal**. Estados com 0, 1, 2, 3 nodos observados directamente na câmara.","epistemico":"fato","exemplos":[],"id":"resultado-chave","memoria":"semantica","meta":{"dominio":"referencia","nivel":5},"origem":"microscopia_fotoionizacao_hidrogenio.md","palavras_chave":[],"sinonimos":[],"tipo":"conceito","titulo":"Resultado-chave"}],"id":"microscopia_fotoionizacao_hidrogenio_resultado-chave","tipo":"conceito"}
\section{Resultado-chave}
A função de onda microscópica ao longo de \ensuremath{\xi} e a projeção do continuo medida a distância macroscópica **partilham a mesma estrutura nodal**. Estados com 0, 1, 2, 3 nodos observados directamente na câmara.
\prov{relações: parte\_de microscopia\_fotoionizacao\_hidrogenio\_microscopia\_de\_fotoionizacao\_orbitais\_de\_hi; confirma orbitais\_de\_hidrogenio\_evidencia\_empirica}
\prov{proveniência: microscopia\_fotoionizacao\_hidrogenio.md \textperiodcentered{} web \textperiodcentered{} fato \textperiodcentered{} confiança 1.0 \textperiodcentered{} domínio referencia \textperiodcentered{} história 4 versões no jornal}
\prov{fusão de curadoria (texto idêntico): absorve resultado-chave --- o mesmo conteúdo por dois esquemas de endereço}

%CRISTAL {"fusao":[{"arestas":[{"alvo":"microscopia_fotoionizacao_hidrogenio_microscopia_de_fotoionizacao_orbitais_de_hi","confianca":1.0,"epistemico":"fato","origem":"microscopia_fotoionizacao_hidrogenio.md","rel":"parte_de"}],"confianca":1.0,"contraexemplos":[],"descricao":"Primeira observação **directa** da estrutura nodal da função de onda electrónica do hidrogénio excitado. Valida predições teóricas de três décadas sobre visualização de orbitais quânticos.","epistemico":"fato","exemplos":[],"id":"microscopia_fotoionizacao_hidrogenio_resumo","memoria":"semantica","meta":{"dominio":"referencia","fonte":"web","nivel":5},"origem":"microscopia_fotoionizacao_hidrogenio.md","palavras_chave":[],"sinonimos":[],"tipo":"conceito","titulo":"Resumo"},{"arestas":[{"alvo":"microscopia_de_fotoionizacao_orbitais_de_hidrogenio_stodolna_2013","confianca":1.0,"epistemico":"fato","origem":"microscopia_fotoionizacao_hidrogenio.md","rel":"parte_de"}],"confianca":1.0,"contraexemplos":[],"descricao":"Primeira observação **directa** da estrutura nodal da função de onda electrónica do hidrogénio excitado. Valida predições teóricas de três décadas sobre visualização de orbitais quânticos.","epistemico":"fato","exemplos":[],"id":"resumo","memoria":"semantica","meta":{"dominio":"referencia","nivel":5},"origem":"microscopia_fotoionizacao_hidrogenio.md","palavras_chave":[],"sinonimos":[],"tipo":"conceito","titulo":"Resumo"}],"id":"microscopia_fotoionizacao_hidrogenio_resumo","tipo":"conceito"}
\section{Resumo}
Primeira observação **directa** da estrutura nodal da função de onda electrónica do hidrogénio excitado. Valida predições teóricas de três décadas sobre visualização de orbitais quânticos.
\prov{relações: parte\_de microscopia\_fotoionizacao\_hidrogenio\_microscopia\_de\_fotoionizacao\_orbitais\_de\_hi}
\prov{proveniência: microscopia\_fotoionizacao\_hidrogenio.md \textperiodcentered{} web \textperiodcentered{} fato \textperiodcentered{} confiança 1.0 \textperiodcentered{} domínio referencia \textperiodcentered{} história 3 versões no jornal}
\prov{fusão de curadoria (texto idêntico): absorve resumo --- o mesmo conteúdo por dois esquemas de endereço}

%CRISTAL {"arestas":[{"alvo":"prova_de_trabalho","confianca":1.0,"epistemico":"fato","origem":"wiki","rel":"relaciona"},{"alvo":"ouro_branco","confianca":1.0,"epistemico":"fato","origem":"wiki","rel":"usa"}],"confianca":1.0,"contraexemplos":[],"descricao":"A prata minera-se de dentro pelos 6 modos dominantes (hexagrama, estrela de Salomão): a emissão exige exatamente rank-6; rank-5 (pentagrama = o ouro, φ) é lastro e não minera; rank-7/ruído também não. Oferta com teto (~13×10⁶).","epistemico":"fato","exemplos":[],"id":"mineracao_hexagrama","memoria":"semantica","meta":{"autor":"Aarão/broca-so","dominio":"broca_repos","fonte":""},"origem":"wiki","palavras_chave":[],"sinonimos":[],"tipo":"conceito","titulo":"Mineração pela Estrela de Prata (rank-6)"}
\section{Mineração pela Estrela de Prata (rank-6)}
A prata minera-se de dentro pelos 6 modos dominantes (hexagrama, estrela de Salomão): a emissão exige exatamente rank-6; rank-5 (pentagrama = o ouro, \ensuremath{\varphi}) é lastro e não minera; rank-7/ruído também não. Oferta com teto (\~{}13$\times$10\ensuremath{^{6}}).
\prov{relações: relaciona prova\_de\_trabalho; usa ouro\_branco}
\prov{proveniência: wiki \textperiodcentered{} fato \textperiodcentered{} confiança 1.0 \textperiodcentered{} domínio broca\_repos \textperiodcentered{} história 3 versões no jornal}

%CRISTAL {"arestas":[{"alvo":"prova_da_maquina_espectral","confianca":1.0,"epistemico":"fato","origem":"CIFRA_PROVA.md","rel":"parte_de"},{"alvo":"utilitarismo","confianca":0.5,"epistemico":"fato","origem":"wiki","rel":"relaciona"},{"alvo":"virtualizacao","confianca":0.5,"epistemico":"fato","origem":"wiki","rel":"relaciona"},{"alvo":"resultado_anterior_identidade_nao_cobertura","confianca":0.5,"epistemico":"fato","origem":"wiki","rel":"relaciona"},{"alvo":"vida-morte","confianca":0.5,"epistemico":"fato","origem":"wiki","rel":"relaciona"},{"alvo":"transcricao","confianca":0.5,"epistemico":"fato","origem":"wiki","rel":"relaciona"},{"alvo":"scheduler","confianca":0.5,"epistemico":"fato","origem":"wiki","rel":"relaciona"},{"alvo":"pedagogia","confianca":0.5,"epistemico":"fato","origem":"wiki","rel":"relaciona"},{"alvo":"ruido","confianca":0.5,"epistemico":"fato","origem":"wiki","rel":"relaciona"},{"alvo":"teoria_do_valor","confianca":0.5,"epistemico":"fato","origem":"wiki","rel":"relaciona"},{"alvo":"resultado_completo","confianca":0.5,"epistemico":"fato","origem":"wiki","rel":"relaciona"},{"alvo":"pow_metal_ressonancia_val","confianca":0.5,"epistemico":"fato","origem":"wiki","rel":"relaciona"},{"alvo":"toryces_e_helyces","confianca":0.5,"epistemico":"fato","origem":"wiki","rel":"relaciona"},{"alvo":"pitagorismo","confianca":0.5,"epistemico":"fato","origem":"wiki","rel":"relaciona"},{"alvo":"paper_2_reta_ouro","confianca":0.5,"epistemico":"fato","origem":"wiki","rel":"relaciona"},{"alvo":"tipos_de_interpretante","confianca":0.5,"epistemico":"fato","origem":"wiki","rel":"relaciona"},{"alvo":"orbita_pell_hardware","confianca":0.5,"epistemico":"fato","origem":"wiki","rel":"relaciona"},{"alvo":"veredito_experimental","confianca":0.5,"epistemico":"fato","origem":"wiki","rel":"relaciona"}],"confianca":1.0,"contraexemplos":[],"descricao":"`SPECT` liga: todo `LOAD` aplica `A₂` automaticamente.","epistemico":"fato","exemplos":[],"id":"modo_espectral_spect","memoria":"semantica","meta":{"dominio":"manual","nivel":6,"verificacao":"slug idêntico — enriquecer, não duplicar"},"origem":"CIFRA_PROVA.md","palavras_chave":[],"sinonimos":[],"tipo":"conceito","titulo":"Modo espectral (`SPECT`)"}
\section{Modo espectral (`SPECT`)}
`SPECT` liga: todo `LOAD` aplica `A\ensuremath{_{2}}` automaticamente.
\prov{relações: parte\_de prova\_da\_maquina\_espectral; relaciona utilitarismo; relaciona virtualizacao; relaciona resultado\_anterior\_identidade\_nao\_cobertura; relaciona vida-morte; relaciona transcricao; relaciona scheduler; relaciona pedagogia; relaciona ruido; relaciona teoria\_do\_valor; relaciona resultado\_completo; relaciona pow\_metal\_ressonancia\_val; relaciona toryces\_e\_helyces; relaciona pitagorismo; relaciona paper\_2\_reta\_ouro; relaciona tipos\_de\_interpretante; relaciona orbita\_pell\_hardware; relaciona veredito\_experimental}
\prov{proveniência: CIFRA\_PROVA.md \textperiodcentered{} fato \textperiodcentered{} confiança 1.0 \textperiodcentered{} domínio manual \textperiodcentered{} história 21 versões no jornal}

%CRISTAL {"arestas":[{"alvo":"gato","confianca":0.95,"epistemico":"fato","origem":"paper","rel":"parte_de"}],"confianca":1.0,"contraexemplos":[],"descricao":"| Modo | Significado | |------|-------------| | par/ímpar | decimação neurônio (Cooley–Tukey) | | grade Koch | 64 modos FFT-like (`koch_scalar` mod P) | | calor, c² | eixo lemniscata real↔ouro e ref↔word | | órbita A₁ | fração de passos com a+c²=1 | | atrator | nó espectral Koch mod 1009 |","epistemico":"fato","exemplos":[],"id":"modos_medidos_por_nonce","memoria":"semantica","meta":{"dominio":"manual","nivel":6,"verificacao":"slug idêntico — enriquecer, não duplicar"},"origem":"POW_METAL_RESSONANCIA.md","palavras_chave":[],"sinonimos":[],"tipo":"conceito","titulo":"Modos medidos por nonce"}
\section{Modos medidos por nonce}
| Modo | Significado | |------|-------------| | par/ímpar | decimação neurônio (Cooley–Tukey) | | grade Koch | 64 modos FFT-like (`koch\_scalar` mod P) | | calor, c\ensuremath{^{2}} | eixo lemniscata real\ensuremath{\leftrightarrow}ouro e ref\ensuremath{\leftrightarrow}word | | órbita A\ensuremath{_{1}} | fração de passos com a+c\ensuremath{^{2}}=1 | | atrator | nó espectral Koch mod 1009 |
\prov{relações: parte\_de gato}
\prov{proveniência: POW\_METAL\_RESSONANCIA.md \textperiodcentered{} fato \textperiodcentered{} confiança 1.0 \textperiodcentered{} domínio manual \textperiodcentered{} história 3 versões no jornal}

%CRISTAL {"arestas":[{"alvo":"colapso_camadas_ordem_parada","confianca":1.0,"epistemico":"fato","origem":"wiki","rel":"parte_de"}],"confianca":1.0,"contraexemplos":[],"descricao":"montar compõe o tecido em blocos priorizados: (1) curados (saudações + turnos ping-pong à mão) + dialogos.tsv manual, (2) Tatoeba conversacional filtrado por _RE_CONV, (3) resto do Tatoeba. Grava camadas separadas (dialogos = bloco 1; corpus = 2+3) e indexa cada uma; a ordem no disco define prioridade em empates Koch. A fonte física das camadas do colapso.","epistemico":"fato","exemplos":[],"id":"montar_conversa_tecido","memoria":"semantica","meta":{"autor":"broca-so","dominio":"conversa","fonte":"conversa/montar_conversa.py:montar/_curados"},"origem":"wiki","palavras_chave":[],"sinonimos":[],"tipo":"conceito","titulo":"Tiffany — diálogos viram tecido (curados + Tatoeba)"}
\section{Tiffany — diálogos viram tecido (curados + Tatoeba)}
montar compõe o tecido em blocos priorizados: (1) curados (saudações + turnos ping-pong à mão) + dialogos.tsv manual, (2) Tatoeba conversacional filtrado por \_RE\_CONV, (3) resto do Tatoeba. Grava camadas separadas (dialogos = bloco 1; corpus = 2+3) e indexa cada uma; a ordem no disco define prioridade em empates Koch. A fonte física das camadas do colapso.
\prov{relações: parte\_de colapso\_camadas\_ordem\_parada}
\prov{proveniência: wiki \textperiodcentered{} conversa/montar\_conversa.py:montar/\_curados \textperiodcentered{} fato \textperiodcentered{} confiança 1.0 \textperiodcentered{} domínio conversa \textperiodcentered{} história 10 versões no jornal}

%CRISTAL {"arestas":[{"alvo":"colapso_koch_tiffany","confianca":1.0,"epistemico":"fato","origem":"wiki","rel":"relaciona"},{"alvo":"corpo_dual","confianca":1.0,"epistemico":"fato","origem":"wiki","rel":"usa"}],"confianca":1.0,"contraexemplos":[],"descricao":"Critério formal e decidível de má-formulação: um problema-par (Geo,Ari) sofre morte sse exige a identidade estática Geo=Ari (cristalização) sob três propriedades restritivas — assimetria de base, gap de continuação analítica (Δ>0), dissipação volumétrica (z↦z², 2-para-1). Sob as três, Geo=Ari é mal posta por construção. Custo O(n)/O(2N).","epistemico":"fato","exemplos":[],"id":"morte_matematica","memoria":"semantica","meta":{"autor":"Aarão/broca-so","dominio":"broca_repos","fonte":""},"origem":"wiki","palavras_chave":[],"sinonimos":[],"tipo":"conceito","titulo":"A Morte Matemática"}
\section{A Morte Matemática}
Critério formal e decidível de má-formulação: um problema-par (Geo,Ari) sofre morte sse exige a identidade estática Geo=Ari (cristalização) sob três propriedades restritivas — assimetria de base, gap de continuação analítica (\ensuremath{\Delta}>0), dissipação volumétrica (z\ensuremath{\mapsto}z\ensuremath{^{2}}, 2-para-1). Sob as três, Geo=Ari é mal posta por construção. Custo O(n)/O(2N).
\prov{relações: relaciona colapso\_koch\_tiffany; usa corpo\_dual}
\prov{proveniência: wiki \textperiodcentered{} fato \textperiodcentered{} confiança 1.0 \textperiodcentered{} domínio broca\_repos \textperiodcentered{} história 4 versões no jornal}

%CRISTAL {"arestas":[{"alvo":"corrida","confianca":1.0,"epistemico":"fato","origem":"POW_METAL.md","rel":"parte_de"},{"alvo":"transcricao","confianca":0.5,"epistemico":"fato","origem":"wiki","rel":"relaciona"},{"alvo":"parabola_pa","confianca":0.5,"epistemico":"fato","origem":"wiki","rel":"relaciona"},{"alvo":"topologia","confianca":0.5,"epistemico":"fato","origem":"wiki","rel":"relaciona"},{"alvo":"resposta_a_pergunta_dificil","confianca":0.5,"epistemico":"fato","origem":"wiki","rel":"relaciona"},{"alvo":"por_sessao_conversadadostelemetriasessaojsonl","confianca":0.5,"epistemico":"fato","origem":"wiki","rel":"relaciona"},{"alvo":"teste_ordem","confianca":0.5,"epistemico":"fato","origem":"wiki","rel":"relaciona"},{"alvo":"thread","confianca":0.5,"epistemico":"fato","origem":"wiki","rel":"relaciona"},{"alvo":"parabola_destruir","confianca":0.5,"epistemico":"fato","origem":"wiki","rel":"relaciona"},{"alvo":"vida-morte","confianca":0.5,"epistemico":"fato","origem":"wiki","rel":"relaciona"},{"alvo":"pow_metal_pre_fecho","confianca":0.5,"epistemico":"fato","origem":"wiki","rel":"relaciona"},{"alvo":"rust","confianca":0.5,"epistemico":"fato","origem":"wiki","rel":"relaciona"},{"alvo":"pilha","confianca":0.5,"epistemico":"fato","origem":"wiki","rel":"relaciona"},{"alvo":"projeto_cristaljson","confianca":0.5,"epistemico":"fato","origem":"wiki","rel":"relaciona"},{"alvo":"ponteiro","confianca":0.5,"epistemico":"fato","origem":"wiki","rel":"relaciona"},{"alvo":"testes","confianca":0.5,"epistemico":"fato","origem":"wiki","rel":"relaciona"},{"alvo":"numa","confianca":0.5,"epistemico":"fato","origem":"wiki","rel":"relaciona"},{"alvo":"tiffany_cabeca_broca","confianca":0.5,"epistemico":"fato","origem":"wiki","rel":"relaciona"},{"alvo":"pow_metal_setor","confianca":0.5,"epistemico":"fato","origem":"wiki","rel":"relaciona"},{"alvo":"pow_geometria","confianca":0.5,"epistemico":"fato","origem":"wiki","rel":"relaciona"},{"alvo":"processos","confianca":0.5,"epistemico":"fato","origem":"wiki","rel":"relaciona"},{"alvo":"recursao","confianca":0.5,"epistemico":"fato","origem":"wiki","rel":"relaciona"}],"confianca":1.0,"contraexemplos":[],"descricao":"- **ambos**: 1,005,083 - **antes**: 5,929 - **depois**: 5,883","epistemico":"fato","exemplos":[],"id":"motivos","memoria":"semantica","meta":{"dominio":"manual","duplicata_sugerida":[],"nivel":6,"verificacao":"parte de conceito mais amplo 'motivos'"},"origem":"DEGENERADOS.md","palavras_chave":[],"sinonimos":[],"tipo":"conceito","titulo":"Motivos"}
\section{Motivos}
- **ambos**: 1,005,083 - **antes**: 5,929 - **depois**: 5,883
\prov{relações: parte\_de corrida; relaciona transcricao; relaciona parabola\_pa; relaciona topologia; relaciona resposta\_a\_pergunta\_dificil; relaciona por\_sessao\_conversadadostelemetriasessaojsonl; relaciona teste\_ordem; relaciona thread; relaciona parabola\_destruir; relaciona vida-morte; relaciona pow\_metal\_pre\_fecho; relaciona rust; relaciona pilha; relaciona projeto\_cristaljson; relaciona ponteiro; relaciona testes; relaciona numa; relaciona tiffany\_cabeca\_broca; relaciona pow\_metal\_setor; relaciona pow\_geometria; relaciona processos; relaciona recursao}
\prov{proveniência: DEGENERADOS.md \textperiodcentered{} fato \textperiodcentered{} confiança 1.0 \textperiodcentered{} domínio manual \textperiodcentered{} história 27 versões no jornal}

%CRISTAL {"arestas":[{"alvo":"mult_prtsu","confianca":1.0,"epistemico":"fato","origem":"wiki","rel":"generaliza"},{"alvo":"mecanica_quantica","confianca":1.0,"epistemico":"fato","origem":"wiki","rel":"relaciona"},{"alvo":"cordas_ilha_hurwitz","confianca":1.0,"epistemico":"fato","origem":"wiki","rel":"relaciona"}],"confianca":1.0,"contraexemplos":[],"descricao":"Extensão A⋆r,t,s,u,vB sobre matrizes Hermitianas (comutador substitui cross product); o ponto canônico (−1,1,1,0,1) recupera a MQ padrão (Bell, Tsirelson 2√2). Varreduras mostram que só ~0,3–2,5% do espaço de parâmetros é físico-admissível: a MQ é uma ilha estreita no espaço de álgebras.","epistemico":"inferencia","exemplos":[],"id":"mq_ilha_estreita","memoria":"semantica","meta":{"autor":"Lopes/Aarão","dominio":"hiper","fonte":""},"origem":"wiki","palavras_chave":[],"sinonimos":[],"tipo":"conceito","titulo":"A Mecânica Quântica é uma Ilha Estreita"}
\section{A Mecânica Quântica é uma Ilha Estreita}
Extensão A\ensuremath{\star}r,t,s,u,vB sobre matrizes Hermitianas (comutador substitui cross product); o ponto canônico (\ensuremath{-}1,1,1,0,1) recupera a MQ padrão (Bell, Tsirelson 2\ensuremath{\sqrt{\,}}2). Varreduras mostram que só \~{}0,3–2,5\% do espaço de parâmetros é físico-admissível: a MQ é uma ilha estreita no espaço de álgebras.
\prov{relações: generaliza mult\_prtsu; relaciona mecanica\_quantica; relaciona cordas\_ilha\_hurwitz}
\prov{proveniência: wiki \textperiodcentered{} inferencia \textperiodcentered{} confiança 1.0 \textperiodcentered{} domínio hiper \textperiodcentered{} história 6 versões no jornal}

%CRISTAL {"arestas":[{"alvo":"nucleo_g2_reabsorcao","confianca":1.0,"epistemico":"fato","origem":"wiki","rel":"usa"},{"alvo":"borda_critica_vazio","confianca":1.0,"epistemico":"fato","origem":"wiki","rel":"contradiz"},{"alvo":"associador","confianca":1.0,"epistemico":"fato","origem":"wiki","rel":"explica"}],"confianca":1.0,"contraexemplos":[],"descricao":"φ₇ preserva a norma mas descarta 2/3 da informação geométrica (a orientação dentro de 𝔤₂) — essa informação amputada É a fonte da não-associatividade octoniônica. O 'resto que nunca zera' é o mass gap algébrico, dim 𝔤₂=14, com estado fundamental a ≥ a_min > 0.","epistemico":"fato","exemplos":[],"id":"nao_associatividade_amputacao","memoria":"semantica","meta":{"autor":"Lopes/Aarão","dominio":"hiper","fonte":""},"origem":"wiki","palavras_chave":[],"sinonimos":[],"tipo":"conceito","titulo":"Não-Associatividade = Informação Amputada (mass gap algébrico)"}
\section{Não-Associatividade = Informação Amputada (mass gap algébrico)}
\ensuremath{\varphi}\ensuremath{_{7}} preserva a norma mas descarta 2/3 da informação geométrica (a orientação dentro de \ensuremath{\mathfrak{g}}\ensuremath{_{2}}) — essa informação amputada É a fonte da não-associatividade octoniônica. O 'resto que nunca zera' é o mass gap algébrico, dim \ensuremath{\mathfrak{g}}\ensuremath{_{2}}=14, com estado fundamental a \ensuremath{\ge} a\_min > 0.
\prov{relações: usa nucleo\_g2\_reabsorcao; contradiz borda\_critica\_vazio; explica associador}
\prov{proveniência: wiki \textperiodcentered{} fato \textperiodcentered{} confiança 1.0 \textperiodcentered{} domínio hiper \textperiodcentered{} história 5 versões no jornal}

%CRISTAL {"arestas":[{"alvo":"maquina_algebrica","confianca":1.0,"epistemico":"fato","origem":"wiki","rel":"usa"},{"alvo":"prova_de_trabalho","confianca":1.0,"epistemico":"fato","origem":"wiki","rel":"relaciona"},{"alvo":"criptoeconomia","confianca":1.0,"epistemico":"fato","origem":"wiki","rel":"relaciona"}],"confianca":1.0,"contraexemplos":[],"descricao":"Rede hipercomplexa vence a GNN clássica idêntica em MAX-CUT com vantagem crescente Δ(n)≈0,04·n³/² (expoente de Edwards), até +57% em n=10000 — aplicada a varejo (mix de SKUs) e smart-grid (islanding, N-1/N-2 como k-cut). O regime MAX-CUT é o mesmo que fundamenta a prova de trabalho.","epistemico":"fato","exemplos":[],"id":"nco_hipercomplexo_maxcut","memoria":"semantica","meta":{"autor":"Lopes/Aarão","dominio":"hiper","fonte":""},"origem":"wiki","palavras_chave":[],"sinonimos":[],"tipo":"conceito","titulo":"Otimização Combinatória Neural Hipercomplexa (MAX-CUT)"}
\section{Otimização Combinatória Neural Hipercomplexa (MAX-CUT)}
Rede hipercomplexa vence a GNN clássica idêntica em MAX-CUT com vantagem crescente \ensuremath{\Delta}(n)\ensuremath{\approx}0,04\textperiodcentered n\ensuremath{^{3}}/\ensuremath{^{2}} (expoente de Edwards), até +57\% em n=10000 — aplicada a varejo (mix de SKUs) e smart-grid (islanding, N-1/N-2 como k-cut). O regime MAX-CUT é o mesmo que fundamenta a prova de trabalho.
\prov{relações: usa maquina\_algebrica; relaciona prova\_de\_trabalho; relaciona criptoeconomia}
\prov{proveniência: wiki \textperiodcentered{} fato \textperiodcentered{} confiança 1.0 \textperiodcentered{} domínio hiper \textperiodcentered{} história 6 versões no jornal}

%CRISTAL {"arestas":[{"alvo":"painel_estelar","confianca":1.0,"epistemico":"fato","origem":"PAINEL_ESTELAR.md","rel":"parte_de"},{"alvo":"fechos_centro_da_janela","confianca":1.0,"epistemico":"fato","origem":"DS_UNIVERSO.md","rel":"parte_de"}],"confianca":1.0,"contraexemplos":[],"descricao":"- habitante `[295, 143]` · ouro `[438, 295]` - Lopes: r=0.078946, c=0.996879, θ=0.129854, Q=1.000000 - dS: τ=0.548285, a(τ)=1.154112 - S³ = (0.5033, 0.5350, -0.4943, 0.4650) - Hopf z=0.0789, φ=0.1299","epistemico":"fato","exemplos":[],"id":"nonce_1572863","memoria":"semantica","meta":{"dominio":"manual","nivel":6,"verificacao":"slug idêntico — enriquecer, não duplicar"},"origem":"DS_UNIVERSO.md","palavras_chave":[],"sinonimos":[],"tipo":"conceito","titulo":"nonce `1572863`"}
\section{nonce `1572863`}
- habitante `[295, 143]` \textperiodcentered  ouro `[438, 295]` - Lopes: r=0.078946, c=0.996879, \ensuremath{\theta}=0.129854, Q=1.000000 - dS: \ensuremath{\tau}=0.548285, a(\ensuremath{\tau})=1.154112 - S\ensuremath{^{3}} = (0.5033, 0.5350, -0.4943, 0.4650) - Hopf z=0.0789, \ensuremath{\varphi}=0.1299
\prov{relações: parte\_de painel\_estelar; parte\_de fechos\_centro\_da\_janela}
\prov{proveniência: DS\_UNIVERSO.md \textperiodcentered{} fato \textperiodcentered{} confiança 1.0 \textperiodcentered{} domínio manual \textperiodcentered{} história 11 versões no jornal}

%CRISTAL {"arestas":[{"alvo":"painel_estelar","confianca":1.0,"epistemico":"fato","origem":"PAINEL_ESTELAR.md","rel":"parte_de"},{"alvo":"fechos_centro_da_janela","confianca":1.0,"epistemico":"fato","origem":"DS_UNIVERSO.md","rel":"parte_de"}],"confianca":1.0,"contraexemplos":[],"descricao":"- habitante `[295, 143]` · ouro `[438, 295]` - Lopes: r=0.078946, c=0.996879, θ=0.129854, Q=1.000000 - dS: τ=0.548285, a(τ)=1.154112 - S³ = (0.5033, 0.5350, -0.4943, 0.4650) - Hopf z=0.0789, φ=0.1299","epistemico":"fato","exemplos":[],"id":"nonce_1966079","memoria":"semantica","meta":{"dominio":"manual","nivel":6,"verificacao":"slug idêntico — enriquecer, não duplicar"},"origem":"DS_UNIVERSO.md","palavras_chave":[],"sinonimos":[],"tipo":"conceito","titulo":"nonce `1966079`"}
\section{nonce `1966079`}
- habitante `[295, 143]` \textperiodcentered  ouro `[438, 295]` - Lopes: r=0.078946, c=0.996879, \ensuremath{\theta}=0.129854, Q=1.000000 - dS: \ensuremath{\tau}=0.548285, a(\ensuremath{\tau})=1.154112 - S\ensuremath{^{3}} = (0.5033, 0.5350, -0.4943, 0.4650) - Hopf z=0.0789, \ensuremath{\varphi}=0.1299
\prov{relações: parte\_de painel\_estelar; parte\_de fechos\_centro\_da\_janela}
\prov{proveniência: DS\_UNIVERSO.md \textperiodcentered{} fato \textperiodcentered{} confiança 1.0 \textperiodcentered{} domínio manual \textperiodcentered{} história 11 versões no jornal}

%CRISTAL {"arestas":[{"alvo":"broca-so_cristalizado","confianca":1.0,"epistemico":"fato","origem":"BROCA_CRISTAL.md","rel":"parte_de"},{"alvo":"recursao","confianca":0.5,"epistemico":"fato","origem":"wiki","rel":"relaciona"},{"alvo":"pilar_i_---_ponte_simbolica_gentil_chart_a_chart_b","confianca":0.5,"epistemico":"fato","origem":"wiki","rel":"relaciona"},{"alvo":"regua_associativa_e_o_risco_de_vazamento","confianca":0.5,"epistemico":"fato","origem":"wiki","rel":"relaciona"},{"alvo":"testes","confianca":0.5,"epistemico":"fato","origem":"wiki","rel":"relaciona"},{"alvo":"processos","confianca":0.5,"epistemico":"fato","origem":"wiki","rel":"relaciona"},{"alvo":"parabola_pa","confianca":0.5,"epistemico":"fato","origem":"wiki","rel":"relaciona"}],"confianca":1.0,"contraexemplos":[],"descricao":"Origem: `neuronio/neuronio.c` — **não modificar**.","epistemico":"fato","exemplos":[],"id":"nucleo_broca_vivo","memoria":"semantica","meta":{"dominio":"manual","nivel":6,"verificacao":"slug idêntico — enriquecer, não duplicar"},"origem":"BROCA_CRISTAL.md","palavras_chave":[],"sinonimos":[],"tipo":"conceito","titulo":"Núcleo (broca vivo)"}
\section{Núcleo (broca vivo)}
Origem: `neuronio/neuronio.c` — **não modificar**.
\prov{relações: parte\_de broca-so\_cristalizado; relaciona recursao; relaciona pilar\_i\_---\_ponte\_simbolica\_gentil\_chart\_a\_chart\_b; relaciona regua\_associativa\_e\_o\_risco\_de\_vazamento; relaciona testes; relaciona processos; relaciona parabola\_pa}
\prov{proveniência: BROCA\_CRISTAL.md \textperiodcentered{} fato \textperiodcentered{} confiança 1.0 \textperiodcentered{} domínio manual \textperiodcentered{} história 10 versões no jornal}

%CRISTAL {"arestas":[{"alvo":"estrela_clifford_projetado","confianca":1.0,"epistemico":"fato","origem":"wiki","rel":"usa"},{"alvo":"torre_hurwitz","confianca":1.0,"epistemico":"fato","origem":"wiki","rel":"usa"},{"alvo":"botao_h_mtheory_g2","confianca":1.0,"epistemico":"fato","origem":"wiki","rel":"relaciona"}],"confianca":1.0,"contraexemplos":[],"descricao":"Teorema: o mapa φₘ:Λ²ℝᵐ→ℝᵐ é isometria em m=3 (quatérnios associativos, núcleo 0), mas em m=7 tem núcleo de dimensão 14 = 𝔤₂ = Der(𝕆); a decomposição Λ²ℝ⁷ = V₇ ⊕ 𝔤₂ (7+14=21) é a dos irredutíveis de G₂.","epistemico":"fato","exemplos":[],"id":"nucleo_g2_reabsorcao","memoria":"semantica","meta":{"autor":"Lopes/Aarão","dominio":"hiper","fonte":""},"origem":"wiki","palavras_chave":[],"sinonimos":[],"tipo":"conceito","titulo":"Núcleo 𝔤₂ e o Mapa de Reabsorção φₘ"}
\section{Núcleo \ensuremath{\mathfrak{g}}\ensuremath{_{2}} e o Mapa de Reabsorção \ensuremath{\varphi}\ensuremath{_{m}}}
Teorema: o mapa \ensuremath{\varphi}\ensuremath{_{m}}:\ensuremath{\Lambda}\ensuremath{^{2}}\ensuremath{\mathbb{R}}\ensuremath{^{m}}\ensuremath{\to}\ensuremath{\mathbb{R}}\ensuremath{^{m}} é isometria em m=3 (quatérnios associativos, núcleo 0), mas em m=7 tem núcleo de dimensão 14 = \ensuremath{\mathfrak{g}}\ensuremath{_{2}} = Der(\ensuremath{\mathbb{O}}); a decomposição \ensuremath{\Lambda}\ensuremath{^{2}}\ensuremath{\mathbb{R}}\ensuremath{^{7}} = V\ensuremath{_{7}} \ensuremath{\oplus} \ensuremath{\mathfrak{g}}\ensuremath{_{2}} (7+14=21) é a dos irredutíveis de G\ensuremath{_{2}}.
\prov{relações: usa estrela\_clifford\_projetado; usa torre\_hurwitz; relaciona botao\_h\_mtheory\_g2}
\prov{proveniência: wiki \textperiodcentered{} fato \textperiodcentered{} confiança 1.0 \textperiodcentered{} domínio hiper \textperiodcentered{} história 5 versões no jornal}

%CRISTAL {"arestas":[{"alvo":"cifra_ouro_branco_nucleo_espectral","confianca":1.0,"epistemico":"fato","origem":"CIFRA.md","rel":"parte_de"},{"alvo":"word","confianca":0.5,"epistemico":"fato","origem":"wiki","rel":"relaciona"}],"confianca":1.0,"contraexemplos":[],"descricao":"```c Word silver(Word x);  /* A₂ — prata / Pell */ Word gold(Word x);    /* A₁ — ouro / Fibonacci */ Word bronze(Word x);  /* A₃ — bronze */ ```","epistemico":"fato","exemplos":[],"id":"nucleo_metalico_cifrac","memoria":"semantica","meta":{"dominio":"manual","nivel":6,"verificacao":"slug idêntico — enriquecer, não duplicar"},"origem":"CIFRA.md","palavras_chave":[],"sinonimos":[],"tipo":"conceito","titulo":"Núcleo metálico (`cifra.c`)"}
\section{Núcleo metálico (`cifra.c`)}
```c Word silver(Word x);  /* A\ensuremath{_{2}} — prata / Pell */ Word gold(Word x);    /* A\ensuremath{_{1}} — ouro / Fibonacci */ Word bronze(Word x);  /* A\ensuremath{_{3}} — bronze */ ```
\prov{relações: parte\_de cifra\_ouro\_branco\_nucleo\_espectral; relaciona word}
\prov{proveniência: CIFRA.md \textperiodcentered{} fato \textperiodcentered{} confiança 1.0 \textperiodcentered{} domínio manual \textperiodcentered{} história 7 versões no jornal}

%CRISTAL {"arestas":[{"alvo":"cifra_ouro_branco_nucleo_espectral","confianca":1.0,"epistemico":"fato","origem":"CIFRA.md","rel":"parte_de"}],"confianca":1.0,"contraexemplos":[],"descricao":"```c Word project(Word antes, Word depois);  /* (calor,fase²) → word */ Word fold(Word x, Word ref);            /* project(silver(x), ref) */ Word unfold(Word x);                    /* gold(x) */ Word lift(Word x, Word ref);            /* project(bronze(x), ref) */ Word cifra_load(Word w);                /* silver(w) — caminho LOADS */ ```","epistemico":"fato","exemplos":[],"id":"nucleo_parabolico","memoria":"semantica","meta":{"dominio":"manual","nivel":6,"verificacao":"slug idêntico — enriquecer, não duplicar"},"origem":"CIFRA.md","palavras_chave":[],"sinonimos":[],"tipo":"conceito","titulo":"Núcleo parabólico"}
\section{Núcleo parabólico}
```c Word project(Word antes, Word depois);  /* (calor,fase\ensuremath{^{2}}) \ensuremath{\to} word */ Word fold(Word x, Word ref);            /* project(silver(x), ref) */ Word unfold(Word x);                    /* gold(x) */ Word lift(Word x, Word ref);            /* project(bronze(x), ref) */ Word cifra\_load(Word w);                /* silver(w) — caminho LOADS */ ```
\prov{relações: parte\_de cifra\_ouro\_branco\_nucleo\_espectral}
\prov{proveniência: CIFRA.md \textperiodcentered{} fato \textperiodcentered{} confiança 1.0 \textperiodcentered{} domínio manual \textperiodcentered{} história 6 versões no jornal}

%CRISTAL {"arestas":[{"alvo":"honestidade_intelectual","confianca":1.0,"epistemico":"fato","origem":"wiki","rel":"usa"},{"alvo":"verdade_por_refutacao_gentil","confianca":1.0,"epistemico":"fato","origem":"wiki","rel":"usa"},{"alvo":"tiffany","confianca":1.0,"epistemico":"fato","origem":"wiki","rel":"relaciona"},{"alvo":"miopia_coletiva","confianca":1.0,"epistemico":"fato","origem":"wiki","rel":"contradiz"}],"confianca":1.0,"contraexemplos":[],"descricao":"O método de construção do broca-so: uma consulta técnica vai em paralelo a várias IAs independentes (GPT, Grok, Gemini, Claude), cada uma com voto honesto, para verificar, refutar e costurar matemática/design antes de qualquer resultado ser aceito. É a epistemologia desta Universidade encarnada.","epistemico":"fato","exemplos":[],"id":"o_conselho","memoria":"semantica","meta":{"autor":"Aarão/broca-so","dominio":"broca_repos","fonte":""},"origem":"wiki","palavras_chave":[],"sinonimos":[],"tipo":"conceito","titulo":"O Conselho (deliberação multi-IA)"}
\section{O Conselho (deliberação multi-IA)}
O método de construção do broca-so: uma consulta técnica vai em paralelo a várias IAs independentes (GPT, Grok, Gemini, Claude), cada uma com voto honesto, para verificar, refutar e costurar matemática/design antes de qualquer resultado ser aceito. É a epistemologia desta Universidade encarnada.
\prov{relações: usa honestidade\_intelectual; usa verdade\_por\_refutacao\_gentil; relaciona tiffany; contradiz miopia\_coletiva}
\prov{proveniência: wiki \textperiodcentered{} fato \textperiodcentered{} confiança 1.0 \textperiodcentered{} domínio broca\_repos \textperiodcentered{} história 5 versões no jornal}

%CRISTAL {"arestas":[{"alvo":"resposta_a_pergunta_dificil","confianca":1.0,"epistemico":"fato","origem":"METAIS.md","rel":"parte_de"},{"alvo":"metais","confianca":1.0,"epistemico":"fato","origem":"manual","rel":"parte_de"}],"confianca":1.0,"contraexemplos":[],"descricao":"A família completa existe: ouro (Fibonacci), prata (Pell), bronze, … O que falta descobrir é **qual dinâmica da Broca seleciona qual n** — provavelmente ligada ao neurônio (fronteira calor↔fase) e não à ISA crua.","epistemico":"fato","exemplos":[],"id":"o_principio_de_selecao","memoria":"semantica","meta":{"dominio":"manual","nivel":6,"verificacao":"slug idêntico — enriquecer, não duplicar"},"origem":"METAIS.md","palavras_chave":[],"sinonimos":[],"tipo":"conceito","titulo":"O princípio de seleção"}
\section{O princípio de seleção}
A família completa existe: ouro (Fibonacci), prata (Pell), bronze, … O que falta descobrir é **qual dinâmica da Broca seleciona qual n** — provavelmente ligada ao neurônio (fronteira calor\ensuremath{\leftrightarrow}fase) e não à ISA crua.
\prov{relações: parte\_de resposta\_a\_pergunta\_dificil; parte\_de metais}
\prov{proveniência: METAIS.md \textperiodcentered{} fato \textperiodcentered{} confiança 1.0 \textperiodcentered{} domínio manual \textperiodcentered{} história 10 versões no jornal}

%CRISTAL {"arestas":[{"alvo":"resposta_a_pergunta_dificil","confianca":1.0,"epistemico":"fato","origem":"METAIS.md","rel":"parte_de"},{"alvo":"metais","confianca":1.0,"epistemico":"fato","origem":"manual","rel":"parte_de"}],"confianca":1.0,"contraexemplos":[],"descricao":"1. **A parábola** é anterior a A₂ — é lei dos pares válidos, independente do metal. 2. **A₂** é o membro n=2 da família Aₙ, cuja recorrência é Pell. 3. A **máquina ISA** não converge espontaneamente para A₂ — os programas actuais são composições ad-hoc, não órbitas puras de metal. 4. **A₂ foi escolhida** na exploração Pell como seed (1,0) — não emergiu da ISA sozinha.","epistemico":"fato","exemplos":[],"id":"o_que_determina_a2","memoria":"semantica","meta":{"dominio":"manual","nivel":6,"verificacao":"slug idêntico — enriquecer, não duplicar"},"origem":"METAIS.md","palavras_chave":[],"sinonimos":[],"tipo":"conceito","titulo":"O que determina A₂?"}
\section{O que determina A\ensuremath{_{2}}?}
1. **A parábola** é anterior a A\ensuremath{_{2}} — é lei dos pares válidos, independente do metal. 2. **A\ensuremath{_{2}}** é o membro n=2 da família A\ensuremath{_{n}}, cuja recorrência é Pell. 3. A **máquina ISA** não converge espontaneamente para A\ensuremath{_{2}} — os programas actuais são composições ad-hoc, não órbitas puras de metal. 4. **A\ensuremath{_{2}} foi escolhida** na exploração Pell como seed (1,0) — não emergiu da ISA sozinha.
\prov{relações: parte\_de resposta\_a\_pergunta\_dificil; parte\_de metais}
\prov{proveniência: METAIS.md \textperiodcentered{} fato \textperiodcentered{} confiança 1.0 \textperiodcentered{} domínio manual \textperiodcentered{} história 10 versões no jornal}

%CRISTAL {"arestas":[{"alvo":"destruir_a_parabola","confianca":1.0,"epistemico":"fato","origem":"PARABOLA_DESTRUIR.md","rel":"parte_de"}],"confianca":1.0,"contraexemplos":[],"descricao":"| Ataque | Descrição | |--------|-----------| | Combinatório | todos pares de valores extremos (±2³¹, zero, etc.) | | Aleatório | milhões de Words aleatórios | | 5000× LOAD | milhares de entradas do neurônio | | 5000× STORE | milhares de gravações | | 10000× ADD | acumulação consecutiva | | 10000× CMP | comparações consecutivas | | Loop CMP | 50k iterações de CMP+JMP | | 1000 programas aleatórios | bytecode randômico |","epistemico":"fato","exemplos":[],"id":"o_que_foi_atacado","memoria":"semantica","meta":{"dominio":"manual","nivel":6,"verificacao":"slug idêntico — enriquecer, não duplicar"},"origem":"PARABOLA_DESTRUIR.md","palavras_chave":[],"sinonimos":[],"tipo":"conceito","titulo":"O que foi atacado"}
\section{O que foi atacado}
| Ataque | Descrição | |--------|-----------| | Combinatório | todos pares de valores extremos ($\pm$2\ensuremath{^{31}}, zero, etc.) | | Aleatório | milhões de Words aleatórios | | 5000$\times$ LOAD | milhares de entradas do neurônio | | 5000$\times$ STORE | milhares de gravações | | 10000$\times$ ADD | acumulação consecutiva | | 10000$\times$ CMP | comparações consecutivas | | Loop CMP | 50k iterações de CMP+JMP | | 1000 programas aleatórios | bytecode randômico |
\prov{relações: parte\_de destruir\_a\_parabola}
\prov{proveniência: PARABOLA\_DESTRUIR.md \textperiodcentered{} fato \textperiodcentered{} confiança 1.0 \textperiodcentered{} domínio manual \textperiodcentered{} história 4 versões no jornal}

%CRISTAL {"arestas":[{"alvo":"delta","confianca":-0.031,"epistemico":"fato","origem":"manual","rel":"referencia"},{"alvo":"bai","confianca":0.023,"epistemico":"fato","origem":"manual","rel":"referencia"},{"alvo":"koch_como_recipiente_memoria_nao-associativa","confianca":1.0,"epistemico":"fato","origem":"KOCH_MEMORIA.md","rel":"parte_de"}],"confianca":1.0,"contraexemplos":[],"descricao":"1. **Associatividade Cantor/Fib** — (∆ s) entre vizinhos do grid; correlação word↔(s). 2. **Não-associatividade Koch** — Hamming `endereco` entre vizinhos (~80 bits); baixa correlação word↔`endereco`. 3. **Colisões recipiente** — mesmo `endereco`/`atrator`, (c²) distinto (várias memórias no mesmo compartimento). 4. **Memória slot** — `memstoreword` + snapshot Koch determinístico; **sem** recall do Word só pelo atrator.","epistemico":"fato","exemplos":[],"id":"o_que_o_experimento_mede","memoria":"semantica","meta":{"dominio":"manual","nivel":6,"verificacao":"slug idêntico — enriquecer, não duplicar"},"origem":"KOCH_MEMORIA.md","palavras_chave":[],"sinonimos":[],"tipo":"conceito","titulo":"O que o experimento mede"}
\section{O que o experimento mede}
1. **Associatividade Cantor/Fib** — (\ensuremath{\Delta} s) entre vizinhos do grid; correlação word\ensuremath{\leftrightarrow}(s). 2. **Não-associatividade Koch** — Hamming `endereco` entre vizinhos (\~{}80 bits); baixa correlação word\ensuremath{\leftrightarrow}`endereco`. 3. **Colisões recipiente** — mesmo `endereco`/`atrator`, (c\ensuremath{^{2}}) distinto (várias memórias no mesmo compartimento). 4. **Memória slot** — `memstoreword` + snapshot Koch determinístico; **sem** recall do Word só pelo atrator.
\prov{relações: referencia delta; referencia bai; parte\_de koch\_como\_recipiente\_memoria\_nao-associativa}
\prov{proveniência: KOCH\_MEMORIA.md \textperiodcentered{} fato \textperiodcentered{} confiança 1.0 \textperiodcentered{} domínio manual \textperiodcentered{} história 13 versões no jornal}

%CRISTAL {"arestas":[{"alvo":"a_cacada_da_parabola_mapa_vivo_da_isa","confianca":1.0,"epistemico":"fato","origem":"PARABOLA_ISA.md","rel":"parte_de"},{"alvo":"broca-so_camada_de_trabalho_isomorfica","confianca":1.0,"epistemico":"fato","origem":"codigo","rel":"parte_de"},{"alvo":"word","confianca":0.5,"epistemico":"fato","origem":"wiki","rel":"relaciona"}],"confianca":1.0,"contraexemplos":[],"descricao":"A Broca não imita Turing. Ela **mediu espectros** e **fechou parábolas**.","epistemico":"fato","exemplos":[],"id":"o_universo_que_a_broca_ja_carrega","memoria":"semantica","meta":{"dominio":"manual","nivel":6,"verificacao":"slug idêntico — enriquecer, não duplicar"},"origem":"PARABOLA_ISA.md","palavras_chave":[],"sinonimos":[],"tipo":"conceito","titulo":"O universo que a Broca já carrega"}
\section{O universo que a Broca já carrega}
A Broca não imita Turing. Ela **mediu espectros** e **fechou parábolas**.
\prov{relações: parte\_de a\_cacada\_da\_parabola\_mapa\_vivo\_da\_isa; parte\_de broca-so\_camada\_de\_trabalho\_isomorfica; relaciona word}
\prov{proveniência: PARABOLA\_ISA.md \textperiodcentered{} fato \textperiodcentered{} confiança 1.0 \textperiodcentered{} domínio manual \textperiodcentered{} história 60 versões no jornal}

%CRISTAL {"arestas":[{"alvo":"setor_algebrico_descontinuidades_no_fecho","confianca":1.0,"epistemico":"fato","origem":"POW_METAL_SETOR.md","rel":"parte_de"},{"alvo":"word","confianca":0.5,"epistemico":"fato","origem":"wiki","rel":"relaciona"}],"confianca":1.0,"contraexemplos":[],"descricao":"> Medir se o acoplamento de modos **corta** o espaço PoW antes do fecho. > Apenas observação — sem alterar mineração, SHA ou fecho.","epistemico":"fato","exemplos":[],"id":"observacao","memoria":"semantica","meta":{"dominio":"manual","duplicata_sugerida":[],"nivel":6,"verificacao":"parte de conceito mais amplo 'observacao'"},"origem":"POW_METAL_PRE_FECHO.md","palavras_chave":[],"sinonimos":[],"tipo":"conceito","titulo":"Observação"}
\section{Observação}
> Medir se o acoplamento de modos **corta** o espaço PoW antes do fecho. > Apenas observação — sem alterar mineração, SHA ou fecho.
\prov{relações: parte\_de setor\_algebrico\_descontinuidades\_no\_fecho; relaciona word}
\prov{proveniência: POW\_METAL\_PRE\_FECHO.md \textperiodcentered{} fato \textperiodcentered{} confiança 1.0 \textperiodcentered{} domínio manual \textperiodcentered{} história 9 versões no jornal}

%CRISTAL {"arestas":[{"alvo":"conversa_recupera_nao_gera","confianca":1.0,"epistemico":"fato","origem":"wiki","rel":"explica"},{"alvo":"telemetria_shadow","confianca":1.0,"epistemico":"fato","origem":"wiki","rel":"relaciona"},{"alvo":"mede_flocos_phi","confianca":1.0,"epistemico":"fato","origem":"wiki","rel":"relaciona"}],"confianca":1.0,"contraexemplos":[],"descricao":"observar.py expõe a conversa sem gravar/fabricar: passo mostra a sequência Koch → contexto → blend → varredura por camada → colapso; fluxo roda uma sequência de turnos (ou uma cadeia) mostrando Δ Koch↔contexto e acumulando o recipiente Φ; ranking lista o Ψ com score/ov/δ/cadeia. Modo --seco não grava contexto. A lente humana sobre o colapso.","epistemico":"fato","exemplos":[],"id":"observar_passo_fluxo","memoria":"semantica","meta":{"autor":"broca-so","dominio":"conversa","fonte":"conversa/observar.py:cmd_passo/cmd_fluxo/cmd_ranking"},"origem":"wiki","palavras_chave":[],"sinonimos":[],"tipo":"conceito","titulo":"Tiffany — observação passo-a-passo do colapso"}
\section{Tiffany — observação passo-a-passo do colapso}
observar.py expõe a conversa sem gravar/fabricar: passo mostra a sequência Koch \ensuremath{\to} contexto \ensuremath{\to} blend \ensuremath{\to} varredura por camada \ensuremath{\to} colapso; fluxo roda uma sequência de turnos (ou uma cadeia) mostrando \ensuremath{\Delta} Koch\ensuremath{\leftrightarrow}contexto e acumulando o recipiente \ensuremath{\Phi}; ranking lista o \ensuremath{\Psi} com score/ov/\ensuremath{\delta}/cadeia. Modo --seco não grava contexto. A lente humana sobre o colapso.
\prov{relações: explica conversa\_recupera\_nao\_gera; relaciona telemetria\_shadow; relaciona mede\_flocos\_phi}
\prov{proveniência: wiki \textperiodcentered{} conversa/observar.py:cmd\_passo/cmd\_fluxo/cmd\_ranking \textperiodcentered{} fato \textperiodcentered{} confiança 1.0 \textperiodcentered{} domínio conversa \textperiodcentered{} história 20 versões no jornal}

%CRISTAL {"arestas":[{"alvo":"koch_pell_verificacao","confianca":1.0,"epistemico":"fato","origem":"KOCH_PELL.md","rel":"parte_de"}],"confianca":1.0,"contraexemplos":[],"descricao":"| Camada | Metal | Objecto | Pell? | |--------|-------|---------|-------| | Cifra | `A₂` | `silver(w) = (2·total+e, total)` | **sim** — órbita `(1,0)` | | Cifra | `A₁` | `gold(w)` | Fibonacci (ouro) | | Canal Koch | — | `z_canal` ← parábola `(√c², √a)` | **não** — ressonância | | Floco Koch | — | `z² mod 1009`, `atrator` XOR | **não** — geometria finita |","epistemico":"fato","exemplos":[],"id":"onde_mora_a_prata","memoria":"semantica","meta":{"dominio":"manual","nivel":6,"verificacao":"slug idêntico — enriquecer, não duplicar"},"origem":"KOCH_PELL.md","palavras_chave":[],"sinonimos":[],"tipo":"conceito","titulo":"Onde mora a prata"}
\section{Onde mora a prata}
| Camada | Metal | Objecto | Pell? | |--------|-------|---------|-------| | Cifra | `A\ensuremath{_{2}}` | `silver(w) = (2\textperiodcentered total+e, total)` | **sim** — órbita `(1,0)` | | Cifra | `A\ensuremath{_{1}}` | `gold(w)` | Fibonacci (ouro) | | Canal Koch | — | `z\_canal` \ensuremath{\leftarrow} parábola `(\ensuremath{\sqrt{\,}}c\ensuremath{^{2}}, \ensuremath{\sqrt{\,}}a)` | **não** — ressonância | | Floco Koch | — | `z\ensuremath{^{2}} mod 1009`, `atrator` XOR | **não** — geometria finita |
\prov{relações: parte\_de koch\_pell\_verificacao}
\prov{proveniência: KOCH\_PELL.md \textperiodcentered{} fato \textperiodcentered{} confiança 1.0 \textperiodcentered{} domínio manual \textperiodcentered{} história 4 versões no jornal}

%CRISTAL {"arestas":[{"alvo":"a_cacada_da_parabola_mapa_vivo_da_isa","confianca":1.0,"epistemico":"fato","origem":"PARABOLA_ISA.md","rel":"parte_de"}],"confianca":1.0,"contraexemplos":[],"descricao":"| Padrão | Onde | |--------|------| | **Neurônio entra** | `LOAD` — fronteira viva, traz fase do arquivo | | **Ressonância** | `ADD`, `CMP` — alta fase entre operandos | | **Cicatriz** | `STORE` — calor grava scratch | | **Dissipação** | `SUB`, `XOR` — abre calor | | **Danza pura** | `JMP/JZ/JNZ` — fluxo sem alterar espectro |","epistemico":"fato","exemplos":[],"id":"onde_respira_onde_cristaliza","memoria":"semantica","meta":{"dominio":"manual","nivel":6,"verificacao":"slug idêntico — enriquecer, não duplicar"},"origem":"PARABOLA_ISA.md","palavras_chave":[],"sinonimos":[],"tipo":"conceito","titulo":"Onde respira / onde cristaliza"}
\section{Onde respira / onde cristaliza}
| Padrão | Onde | |--------|------| | **Neurônio entra** | `LOAD` — fronteira viva, traz fase do arquivo | | **Ressonância** | `ADD`, `CMP` — alta fase entre operandos | | **Cicatriz** | `STORE` — calor grava scratch | | **Dissipação** | `SUB`, `XOR` — abre calor | | **Danza pura** | `JMP/JZ/JNZ` — fluxo sem alterar espectro |
\prov{relações: parte\_de a\_cacada\_da\_parabola\_mapa\_vivo\_da\_isa}
\prov{proveniência: PARABOLA\_ISA.md \textperiodcentered{} fato \textperiodcentered{} confiança 1.0 \textperiodcentered{} domínio manual \textperiodcentered{} história 6 versões no jornal}

%CRISTAL {"arestas":[{"alvo":"mult_prtsu","confianca":1.0,"epistemico":"fato","origem":"wiki","rel":"generaliza"},{"alvo":"matematica_estelar","confianca":1.0,"epistemico":"fato","origem":"wiki","rel":"relaciona"},{"alvo":"cordas_ilha_hurwitz","confianca":1.0,"epistemico":"fato","origem":"wiki","rel":"relaciona"}],"confianca":1.0,"contraexemplos":[],"descricao":"A tese do livro: o objeto fundamental não são partículas, campos ou geometrias, mas o tensor de multiplicação μ∈V^⊗3; cada teoria física (Galileu r=0, Einstein r=1, Maxwell, Newton, RG) emerge como ponto ou trajetória da família μs — não é acrescentada ao formalismo, é lida dele.","epistemico":"hipotese","exemplos":[],"id":"ontologia_da_multiplicacao","memoria":"semantica","meta":{"autor":"Aarão/broca-so","dominio":"broca_repos","fonte":""},"origem":"wiki","palavras_chave":[],"sinonimos":[],"tipo":"conceito","titulo":"A Ontologia da Multiplicação"}
\section{A Ontologia da Multiplicação}
A tese do livro: o objeto fundamental não são partículas, campos ou geometrias, mas o tensor de multiplicação \ensuremath{\mu}\ensuremath{\in}V\^{}\ensuremath{\otimes}3; cada teoria física (Galileu r=0, Einstein r=1, Maxwell, Newton, RG) emerge como ponto ou trajetória da família \ensuremath{\mu}s — não é acrescentada ao formalismo, é lida dele.
\prov{relações: generaliza mult\_prtsu; relaciona matematica\_estelar; relaciona cordas\_ilha\_hurwitz}
\prov{proveniência: wiki \textperiodcentered{} hipotese \textperiodcentered{} confiança 1.0 \textperiodcentered{} domínio broca\_repos \textperiodcentered{} história 5 versões no jornal}

%CRISTAL {"arestas":[{"alvo":"broca_os","confianca":0.95,"epistemico":"fato","origem":"paper","rel":"parte_de"},{"alvo":"testes","confianca":0.5,"epistemico":"fato","origem":"wiki","rel":"relaciona"},{"alvo":"rust","confianca":0.5,"epistemico":"fato","origem":"wiki","rel":"relaciona"},{"alvo":"recursao","confianca":0.5,"epistemico":"fato","origem":"wiki","rel":"relaciona"},{"alvo":"tiffany_cabeca_broca","confianca":0.5,"epistemico":"fato","origem":"wiki","rel":"relaciona"},{"alvo":"respiracao_dos_programas","confianca":0.5,"epistemico":"fato","origem":"wiki","rel":"relaciona"},{"alvo":"vida-morte","confianca":0.5,"epistemico":"fato","origem":"wiki","rel":"relaciona"},{"alvo":"regua_associativa_e_o_risco_de_vazamento","confianca":0.5,"epistemico":"fato","origem":"wiki","rel":"relaciona"}],"confianca":1.0,"contraexemplos":[],"descricao":"tar -xf dist/tiffany-core-1.0.0.tar.* -C $TIFFANY_HOME","epistemico":"fato","exemplos":[],"id":"opcao_a_tarball_release","memoria":"semantica","meta":{"dominio":"manual","nivel":6,"verificacao":"slug idêntico — enriquecer, não duplicar"},"origem":"INSTALL.md","palavras_chave":[],"sinonimos":[],"tipo":"conceito","titulo":"opção A: tarball release"}
\section{opção A: tarball release}
tar -xf dist/tiffany-core-1.0.0.tar.* -C \$TIFFANY\_HOME
\prov{relações: parte\_de broca\_os; relaciona testes; relaciona rust; relaciona recursao; relaciona tiffany\_cabeca\_broca; relaciona respiracao\_dos\_programas; relaciona vida-morte; relaciona regua\_associativa\_e\_o\_risco\_de\_vazamento}
\prov{proveniência: INSTALL.md \textperiodcentered{} fato \textperiodcentered{} confiança 1.0 \textperiodcentered{} domínio manual \textperiodcentered{} história 10 versões no jornal}

%CRISTAL {"arestas":[{"alvo":"broca_os","confianca":1.0,"epistemico":"fato","origem":"manual","rel":"parte_de"},{"alvo":"utilitarismo","confianca":0.5,"epistemico":"fato","origem":"wiki","rel":"relaciona"},{"alvo":"virtualizacao","confianca":0.5,"epistemico":"fato","origem":"wiki","rel":"relaciona"},{"alvo":"word","confianca":0.5,"epistemico":"fato","origem":"wiki","rel":"relaciona"},{"alvo":"vontade_de_poder","confianca":0.5,"epistemico":"fato","origem":"wiki","rel":"relaciona"}],"confianca":1.0,"contraexemplos":[],"descricao":"cd code && ./tiffany-init --home $TIFFANY_HOME sh ../content/core-p0/build.sh ```","epistemico":"fato","exemplos":[],"id":"opcao_b_desenvolvimento_no_repo","memoria":"semantica","meta":{"dominio":"manual","nivel":6,"verificacao":"slug idêntico — enriquecer, não duplicar"},"origem":"INSTALL.md","palavras_chave":[],"sinonimos":[],"tipo":"conceito","titulo":"opção B: desenvolvimento no repo"}
\section{opção B: desenvolvimento no repo}
cd code \&\& ./tiffany-init --home \$TIFFANY\_HOME sh ../content/core-p0/build.sh ```
\prov{relações: parte\_de broca\_os; relaciona utilitarismo; relaciona virtualizacao; relaciona word; relaciona vontade\_de\_poder}
\prov{proveniência: INSTALL.md \textperiodcentered{} fato \textperiodcentered{} confiança 1.0 \textperiodcentered{} domínio manual \textperiodcentered{} história 12 versões no jornal}

%CRISTAL {"arestas":[{"alvo":"colapso","confianca":0.055,"epistemico":"fato","origem":"manual","rel":"referencia"},{"alvo":"spec_trabalho","confianca":1.0,"epistemico":"fato","origem":"BROCA_SO.md","rel":"parte_de"}],"confianca":1.0,"contraexemplos":[],"descricao":"```text GATO_STREAM → TORO → REGISTRO_N → COLAPSO_6 → BANDA → NONCE └────── BROCA_OP_POW_SEED (composite) ──────┘ ```","epistemico":"fato","exemplos":[],"id":"ops_sequencia_pow_nativo","memoria":"semantica","meta":{"dominio":"manual","nivel":6,"verificacao":"slug idêntico — enriquecer, não duplicar"},"origem":"BROCA_SO.md","palavras_chave":[],"sinonimos":[],"tipo":"conceito","titulo":"Ops (sequência PoW nativo)"}
\section{Ops (sequência PoW nativo)}
```text GATO\_STREAM \ensuremath{\to} TORO \ensuremath{\to} REGISTRO\_N \ensuremath{\to} COLAPSO\_6 \ensuremath{\to} BANDA \ensuremath{\to} NONCE ------ BROCA\_OP\_POW\_SEED (composite) ------ ```
\prov{relações: referencia colapso; parte\_de spec\_trabalho}
\prov{proveniência: BROCA\_SO.md \textperiodcentered{} fato \textperiodcentered{} confiança 1.0 \textperiodcentered{} domínio manual \textperiodcentered{} história 9 versões no jornal}

%CRISTAL {"arestas":[{"alvo":"prova_da_maquina_espectral","confianca":1.0,"epistemico":"fato","origem":"CIFRA_PROVA.md","rel":"parte_de"},{"alvo":"utilitarismo","confianca":0.5,"epistemico":"fato","origem":"wiki","rel":"relaciona"},{"alvo":"transcricao","confianca":0.5,"epistemico":"fato","origem":"wiki","rel":"relaciona"},{"alvo":"pedagogia","confianca":0.5,"epistemico":"fato","origem":"wiki","rel":"relaciona"},{"alvo":"resultado_completo","confianca":0.5,"epistemico":"fato","origem":"wiki","rel":"relaciona"},{"alvo":"resultado_anterior_identidade_nao_cobertura","confianca":0.5,"epistemico":"fato","origem":"wiki","rel":"relaciona"},{"alvo":"scheduler","confianca":0.5,"epistemico":"fato","origem":"wiki","rel":"relaciona"},{"alvo":"vida-morte","confianca":0.5,"epistemico":"fato","origem":"wiki","rel":"relaciona"}],"confianca":1.0,"contraexemplos":[],"descricao":"Seed `(1,0)` + 4× `SILVER` → `(29,12)` via `09_pell.asm`.","epistemico":"fato","exemplos":[],"id":"orbita_pell_hardware","memoria":"semantica","meta":{"dominio":"manual","nivel":6,"verificacao":"slug idêntico — enriquecer, não duplicar"},"origem":"CIFRA_PROVA.md","palavras_chave":[],"sinonimos":[],"tipo":"conceito","titulo":"Órbita Pell — hardware"}
\section{Órbita Pell — hardware}
Seed `(1,0)` + 4$\times$ `SILVER` \ensuremath{\to} `(29,12)` via `09\_pell.asm`.
\prov{relações: parte\_de prova\_da\_maquina\_espectral; relaciona utilitarismo; relaciona transcricao; relaciona pedagogia; relaciona resultado\_completo; relaciona resultado\_anterior\_identidade\_nao\_cobertura; relaciona scheduler; relaciona vida-morte}
\prov{proveniência: CIFRA\_PROVA.md \textperiodcentered{} fato \textperiodcentered{} confiança 1.0 \textperiodcentered{} domínio manual \textperiodcentered{} história 11 versões no jornal}

%CRISTAL {"arestas":[{"alvo":"metais","confianca":0.95,"epistemico":"fato","origem":"paper","rel":"parte_de"}],"confianca":1.0,"contraexemplos":[],"descricao":"| k | Pell | estado (total,e) | degenerado | parábola (calor, fase) | |---|------|------------------|------------|------------------------| | 1 | 1 | `1,0` | não | (0.333, 0.816) | | 2 | 2 | `2,1` | não | (0.003, 0.999) | | 3 | 5 | `5,2` | não | (0.000, 1.000) | | 4 | 12 | `12,5` | não | (0.000, 1.000) | | 5 | 29 | `29,12` | não | (0.000, 1.000) | | 6 | 70 | `70,29` | não | (0.000, 1.000) | | 7 | 169 | `169,70` | não | (0.000, 1.000) |","epistemico":"fato","exemplos":[],"id":"orbita_prata_a210","memoria":"semantica","meta":{"dominio":"manual","nivel":6,"verificacao":"slug idêntico — enriquecer, não duplicar"},"origem":"PELL.md","palavras_chave":[],"sinonimos":[],"tipo":"conceito","titulo":"Órbita prata A₂·(1,0)"}
\section{Órbita prata A\ensuremath{_{2}}\textperiodcentered (1,0)}
| k | Pell | estado (total,e) | degenerado | parábola (calor, fase) | |---|------|------------------|------------|------------------------| | 1 | 1 | `1,0` | não | (0.333, 0.816) | | 2 | 2 | `2,1` | não | (0.003, 0.999) | | 3 | 5 | `5,2` | não | (0.000, 1.000) | | 4 | 12 | `12,5` | não | (0.000, 1.000) | | 5 | 29 | `29,12` | não | (0.000, 1.000) | | 6 | 70 | `70,29` | não | (0.000, 1.000) | | 7 | 169 | `169,70` | não | (0.000, 1.000) |
\prov{relações: parte\_de metais}
\prov{proveniência: PELL.md \textperiodcentered{} fato \textperiodcentered{} confiança 1.0 \textperiodcentered{} domínio manual \textperiodcentered{} história 3 versões no jornal}

%CRISTAL {"arestas":[{"alvo":"broca_os","confianca":1.0,"epistemico":"fato","origem":"wiki","rel":"usa"},{"alvo":"mult_prtsu","confianca":1.0,"epistemico":"fato","origem":"wiki","rel":"usa"},{"alvo":"a_bolha_infraestrutura","confianca":1.0,"epistemico":"fato","origem":"wiki","rel":"relaciona"}],"confianca":1.0,"contraexemplos":[],"descricao":"A arquitetura de implementação no servidor central: boy (traz o dado bruto sem picotar), Central Estelar (única que toca o bruto, selada por Ed25519), universos U2–U9 (digerem em células PRTSU, nomeados por pensadores), bus (coração + meta-veia PI + diagnose de 113 células), sync Patria a cada 5 min. DNA em transição dos genes PRTSU para a máquina de torção (o gap).","epistemico":"inferencia","exemplos":[],"id":"organismo_gex44","memoria":"semantica","meta":{"autor":"Aarão/broca-so","dominio":"broca_repos","fonte":""},"origem":"wiki","palavras_chave":[],"sinonimos":[],"tipo":"conceito","titulo":"O Organismo GEX44 (grafo de fluxo, não multiverso físico)"}
\section{O Organismo GEX44 (grafo de fluxo, não multiverso físico)}
A arquitetura de implementação no servidor central: boy (traz o dado bruto sem picotar), Central Estelar (única que toca o bruto, selada por Ed25519), universos U2–U9 (digerem em células PRTSU, nomeados por pensadores), bus (coração + meta-veia PI + diagnose de 113 células), sync Patria a cada 5 min. DNA em transição dos genes PRTSU para a máquina de torção (o gap).
\prov{relações: usa broca\_os; usa mult\_prtsu; relaciona a\_bolha\_infraestrutura}
\prov{proveniência: wiki \textperiodcentered{} inferencia \textperiodcentered{} confiança 1.0 \textperiodcentered{} domínio broca\_repos \textperiodcentered{} história 4 versões no jornal}

%CRISTAL {"arestas":[{"alvo":"o_substrato_a_maquina_e_funcao_de_onda","confianca":1.0,"epistemico":"fato","origem":"wiki","rel":"relaciona"},{"alvo":"topologia_vazia_do_suporte","confianca":1.0,"epistemico":"fato","origem":"wiki","rel":"usa"}],"confianca":1.0,"contraexemplos":[],"descricao":"Decomposição do mesmo objeto em três camadas separáveis: osso = dim_H (geometria do suporte, ε=+1), carne = ρ=|ψ|² (densidade), sangue = j=ρ∇S (fluxo do gradiente de fase, ε=−1, via Madelung/Cole-Hopf). A identificação osso/carne/sangue é conjectura; a matemática de Madelung é teorema.","epistemico":"hipotese","exemplos":[],"id":"osso_carne_sangue","memoria":"semantica","meta":{"autor":"Aarão/broca-so","dominio":"broca_repos","fonte":""},"origem":"wiki","palavras_chave":[],"sinonimos":[],"tipo":"conceito","titulo":"Osso, Carne e Sangue (as três camadas do suporte)"}
\section{Osso, Carne e Sangue (as três camadas do suporte)}
Decomposição do mesmo objeto em três camadas separáveis: osso = dim\_H (geometria do suporte, \ensuremath{\varepsilon}=+1), carne = \ensuremath{\rho}=|\ensuremath{\psi}|\ensuremath{^{2}} (densidade), sangue = j=\ensuremath{\rho}\ensuremath{\nabla}S (fluxo do gradiente de fase, \ensuremath{\varepsilon}=\ensuremath{-}1, via Madelung/Cole-Hopf). A identificação osso/carne/sangue é conjectura; a matemática de Madelung é teorema.
\prov{relações: relaciona o\_substrato\_a\_maquina\_e\_funcao\_de\_onda; usa topologia\_vazia\_do\_suporte}
\prov{proveniência: wiki \textperiodcentered{} hipotese \textperiodcentered{} confiança 1.0 \textperiodcentered{} domínio broca\_repos \textperiodcentered{} história 3 versões no jornal}

%CRISTAL {"arestas":[{"alvo":"goldchain","confianca":1.0,"epistemico":"fato","origem":"wiki","rel":"usa"},{"alvo":"cifra","confianca":1.0,"epistemico":"fato","origem":"wiki","rel":"usa"},{"alvo":"dimensao_de_prata_pell","confianca":1.0,"epistemico":"fato","origem":"wiki","rel":"usa"}],"confianca":1.0,"contraexemplos":[],"descricao":"A moeda interna da Engenharia de Ouro — a liga do ouro (lastro externo, os swaps cross-chain) com a face prateada da dimensão de prata σ₂. Nasce dentro da máquina (não entra de fora), minerada por estrela de prata e cifrada por mudança de base em 𝔽₂ᵏ, tudo em bits.","epistemico":"inferencia","exemplos":[],"id":"ouro_branco","memoria":"semantica","meta":{"autor":"Aarão/broca-so","dominio":"broca_repos","fonte":""},"origem":"wiki","palavras_chave":[],"sinonimos":[],"tipo":"conceito","titulo":"Ouro Branco (a moeda interna)"}
\section{Ouro Branco (a moeda interna)}
A moeda interna da Engenharia de Ouro — a liga do ouro (lastro externo, os swaps cross-chain) com a face prateada da dimensão de prata \ensuremath{\sigma}\ensuremath{_{2}}. Nasce dentro da máquina (não entra de fora), minerada por estrela de prata e cifrada por mudança de base em \ensuremath{\mathbb{F}}\ensuremath{_{2}}\ensuremath{^{k}}, tudo em bits.
\prov{relações: usa goldchain; usa cifra; usa dimensao\_de\_prata\_pell}
\prov{proveniência: wiki \textperiodcentered{} inferencia \textperiodcentered{} confiança 1.0 \textperiodcentered{} domínio broca\_repos \textperiodcentered{} história 4 versões no jornal}

%CRISTAL {"arestas":[{"alvo":"painel_estelar","confianca":1.0,"epistemico":"fato","origem":"manual","rel":"parte_de"},{"alvo":"gato","confianca":1.0,"epistemico":"fato","origem":"manual","rel":"especializa"},{"alvo":"colapso","confianca":1.0,"epistemico":"fato","origem":"manual","rel":"usa"},{"alvo":"vida-morte","confianca":0.5,"epistemico":"fato","origem":"wiki","rel":"relaciona"},{"alvo":"word","confianca":0.5,"epistemico":"fato","origem":"wiki","rel":"relaciona"},{"alvo":"pell","confianca":0.5,"epistemico":"fato","origem":"wiki","rel":"relaciona"},{"alvo":"virtualizacao","confianca":0.5,"epistemico":"fato","origem":"wiki","rel":"relaciona"},{"alvo":"telemetria_shadow_producao","confianca":0.5,"epistemico":"fato","origem":"wiki","rel":"relaciona"},{"alvo":"pipeline_de_ingestao","confianca":0.5,"epistemico":"fato","origem":"wiki","rel":"relaciona"}],"confianca":1.0,"contraexemplos":[],"descricao":"Observatório das **três consequências** a partir do invariante do gato.","epistemico":"fato","exemplos":[],"id":"painel_estelar","memoria":"semantica","meta":{"dominio":"manual","nivel":6,"verificacao":"slug idêntico — enriquecer, não duplicar"},"origem":"PAINEL_ESTELAR.md","palavras_chave":[],"sinonimos":[],"tipo":"conceito","titulo":"PAINEL_ESTELAR"}
\section{PAINEL\_ESTELAR}
Observatório das **três consequências** a partir do invariante do gato.
\prov{relações: parte\_de painel\_estelar; especializa gato; usa colapso; relaciona vida-morte; relaciona word; relaciona pell; relaciona virtualizacao; relaciona telemetria\_shadow\_producao; relaciona pipeline\_de\_ingestao}
\prov{proveniência: PAINEL\_ESTELAR.md \textperiodcentered{} fato \textperiodcentered{} confiança 1.0 \textperiodcentered{} domínio manual \textperiodcentered{} história 88 versões no jornal}

%CRISTAL {"arestas":[{"alvo":"koch_como_recipiente_memoria_nao-associativa","confianca":1.0,"epistemico":"fato","origem":"KOCH_MEMORIA.md","rel":"parte_de"}],"confianca":1.0,"contraexemplos":[],"descricao":"| Peça | Papel | Associativo? | |------|-------|--------------| | **Cantor** (em (s)) | Endereço na reta real desdistorcida | **Sim** — vizinhos têm (s) próximo | | **Fibonacci** (em (s)) | Marca convergentes (Fₙ/Fₙ+₁) | **Sim** — estrutura da reta ouro | | **Koch** (`endereco`, `atrator`) | **Recipiente** no canal negro Peirce | **Não** — salto grande, colisões, sem recall por similaridade |","epistemico":"fato","exemplos":[],"id":"papeis_distintos","memoria":"semantica","meta":{"dominio":"manual","nivel":6,"verificacao":"slug idêntico — enriquecer, não duplicar"},"origem":"KOCH_MEMORIA.md","palavras_chave":[],"sinonimos":[],"tipo":"conceito","titulo":"Papéis distintos"}
\section{Papéis distintos}
| Peça | Papel | Associativo? | |------|-------|--------------| | **Cantor** (em (s)) | Endereço na reta real desdistorcida | **Sim** — vizinhos têm (s) próximo | | **Fibonacci** (em (s)) | Marca convergentes (F\ensuremath{_{n}}/F\ensuremath{_{n}}+\ensuremath{_{1}}) | **Sim** — estrutura da reta ouro | | **Koch** (`endereco`, `atrator`) | **Recipiente** no canal negro Peirce | **Não** — salto grande, colisões, sem recall por similaridade |
\prov{relações: parte\_de koch\_como\_recipiente\_memoria\_nao-associativa}
\prov{proveniência: KOCH\_MEMORIA.md \textperiodcentered{} fato \textperiodcentered{} confiança 1.0 \textperiodcentered{} domínio manual \textperiodcentered{} história 7 versões no jornal}

%CRISTAL {"arestas":[{"alvo":"fisica","confianca":1.0,"epistemico":"fato","origem":"paper","rel":"parte_de"},{"alvo":"vontade_de_poder","confianca":0.5,"epistemico":"fato","origem":"wiki","rel":"relaciona"},{"alvo":"virtualizacao","confianca":0.5,"epistemico":"fato","origem":"wiki","rel":"relaciona"},{"alvo":"word","confianca":0.5,"epistemico":"fato","origem":"wiki","rel":"relaciona"},{"alvo":"vida-morte","confianca":0.5,"epistemico":"fato","origem":"wiki","rel":"relaciona"}],"confianca":1.0,"contraexemplos":[],"descricao":"Ver `papers/painel_estelar.tex` — derivação das três ciências a partir do gato (`quatro_ciencias.tex`).","epistemico":"fato","exemplos":[],"id":"paper_0_sintese","memoria":"semantica","meta":{"dominio":"manual","duplicata_sugerida":[],"nivel":6,"verificacao":"parte de conceito mais amplo 'paper_0_sintese'"},"origem":"PAINEL_ESTELAR.md","palavras_chave":[],"sinonimos":[],"tipo":"conceito","titulo":"paper_0_sintese"}
\section{paper\_0\_sintese}
Ver `papers/painel\_estelar.tex` — derivação das três ciências a partir do gato (`quatro\_ciencias.tex`).
\prov{relações: parte\_de fisica; relaciona vontade\_de\_poder; relaciona virtualizacao; relaciona word; relaciona vida-morte}
\prov{proveniência: PAINEL\_ESTELAR.md \textperiodcentered{} fato \textperiodcentered{} confiança 1.0 \textperiodcentered{} domínio manual \textperiodcentered{} história 37 versões no jornal}

%CRISTAL {"arestas":[{"alvo":"morte","confianca":1.0,"epistemico":"fato","origem":"manual","rel":"referencia"},{"alvo":"transcricao","confianca":0.5,"epistemico":"fato","origem":"wiki","rel":"relaciona"},{"alvo":"vida-morte","confianca":0.5,"epistemico":"fato","origem":"wiki","rel":"relaciona"},{"alvo":"resposta_a_pergunta_dificil","confianca":0.5,"epistemico":"fato","origem":"wiki","rel":"relaciona"},{"alvo":"pow_metal_pre_fecho","confianca":0.5,"epistemico":"fato","origem":"wiki","rel":"relaciona"},{"alvo":"teste_ordem","confianca":0.5,"epistemico":"fato","origem":"wiki","rel":"relaciona"},{"alvo":"por_sessao_conversadadostelemetriasessaojsonl","confianca":0.5,"epistemico":"fato","origem":"wiki","rel":"relaciona"},{"alvo":"parabola_pa","confianca":0.5,"epistemico":"fato","origem":"wiki","rel":"relaciona"},{"alvo":"pow_metal_funil","confianca":0.5,"epistemico":"fato","origem":"wiki","rel":"relaciona"},{"alvo":"pilha","confianca":0.5,"epistemico":"fato","origem":"wiki","rel":"relaciona"},{"alvo":"pow_metal_setor","confianca":0.5,"epistemico":"fato","origem":"wiki","rel":"relaciona"},{"alvo":"tiffany_cabeca_broca","confianca":0.5,"epistemico":"fato","origem":"wiki","rel":"relaciona"},{"alvo":"pow_geometria","confianca":0.5,"epistemico":"fato","origem":"wiki","rel":"relaciona"},{"alvo":"testes","confianca":0.5,"epistemico":"fato","origem":"wiki","rel":"relaciona"},{"alvo":"topologia","confianca":0.5,"epistemico":"fato","origem":"wiki","rel":"relaciona"},{"alvo":"virtualizacao","confianca":0.5,"epistemico":"fato","origem":"wiki","rel":"relaciona"}],"confianca":1.0,"contraexemplos":[],"descricao":"Manual: ula/PARABOLA_DESTRUIR.md","epistemico":"fato","exemplos":[],"id":"parabola_destruir","memoria":"semantica","meta":{"arquivo":"ula/PARABOLA_DESTRUIR.md","dominio":"manual","nivel":6},"origem":"manual","palavras_chave":[],"sinonimos":[],"tipo":"conceito","titulo":"PARABOLA_DESTRUIR"}
\section{PARABOLA\_DESTRUIR}
Manual: ula/PARABOLA\_DESTRUIR.md
\prov{relações: referencia morte; relaciona transcricao; relaciona vida-morte; relaciona resposta\_a\_pergunta\_dificil; relaciona pow\_metal\_pre\_fecho; relaciona teste\_ordem; relaciona por\_sessao\_conversadadostelemetriasessaojsonl; relaciona parabola\_pa; relaciona pow\_metal\_funil; relaciona pilha; relaciona pow\_metal\_setor; relaciona tiffany\_cabeca\_broca; relaciona pow\_geometria; relaciona testes; relaciona topologia; relaciona virtualizacao}
\prov{proveniência: manual \textperiodcentered{} fato \textperiodcentered{} confiança 1.0 \textperiodcentered{} domínio manual \textperiodcentered{} história 19 versões no jornal}

%CRISTAL {"arestas":[{"alvo":"broca_os","confianca":1.0,"epistemico":"fato","origem":"manual","rel":"explica"},{"alvo":"pow_bit_taxa","confianca":0.5,"epistemico":"fato","origem":"wiki","rel":"relaciona"},{"alvo":"transcricao","confianca":0.5,"epistemico":"fato","origem":"wiki","rel":"relaciona"},{"alvo":"prova_isa","confianca":0.5,"epistemico":"fato","origem":"wiki","rel":"relaciona"},{"alvo":"utilitarismo","confianca":0.5,"epistemico":"fato","origem":"wiki","rel":"relaciona"},{"alvo":"word","confianca":0.5,"epistemico":"fato","origem":"wiki","rel":"relaciona"},{"alvo":"virtualizacao","confianca":0.5,"epistemico":"fato","origem":"wiki","rel":"relaciona"},{"alvo":"vida-morte","confianca":0.5,"epistemico":"fato","origem":"wiki","rel":"relaciona"}],"confianca":1.0,"contraexemplos":[],"descricao":"Manual: ula/PARABOLA_ISA.md","epistemico":"fato","exemplos":[],"id":"parabola_isa","memoria":"semantica","meta":{"arquivo":"ula/PARABOLA_ISA.md","dominio":"manual","nivel":6},"origem":"manual","palavras_chave":[],"sinonimos":[],"tipo":"conceito","titulo":"PARABOLA_ISA"}
\section{PARABOLA\_ISA}
Manual: ula/PARABOLA\_ISA.md
\prov{relações: explica broca\_os; relaciona pow\_bit\_taxa; relaciona transcricao; relaciona prova\_isa; relaciona utilitarismo; relaciona word; relaciona virtualizacao; relaciona vida-morte}
\prov{proveniência: manual \textperiodcentered{} fato \textperiodcentered{} confiança 1.0 \textperiodcentered{} domínio manual \textperiodcentered{} história 38 versões no jornal}

%CRISTAL {"arestas":[{"alvo":"koch_codifica_cantorfibonacci_prova_e_evidencia","confianca":1.0,"epistemico":"fato","origem":"KOCH_CANTOR_FIB_PROVA.md","rel":"parte_de"}],"confianca":1.0,"contraexemplos":[],"descricao":"1. **Lema gold–Fibonacci:** provar que `gold(w)` minimiza distância a (Fₙ/Fₙ+₁) no parâmetro (c²). 2. **Lema poda–órbita:** provar que LSB da órbita (z²) gera só palavras em (Σ). 3. **Lema parábola–reta:** (Φ) conjugada monótona entre ((c²,a)) e (s). 4. **Teorema produto:** (π(w)) isomorfismo semigrupal para palavras admissíveis × Cantor ternário.","epistemico":"fato","exemplos":[],"id":"passos_para_prova_fechada","memoria":"semantica","meta":{"dominio":"manual","nivel":6,"verificacao":"slug idêntico — enriquecer, não duplicar"},"origem":"KOCH_CANTOR_FIB_PROVA.md","palavras_chave":[],"sinonimos":[],"tipo":"conceito","titulo":"Passos para prova fechada"}
\section{Passos para prova fechada}
1. **Lema gold–Fibonacci:** provar que `gold(w)` minimiza distância a (F\ensuremath{_{n}}/F\ensuremath{_{n}}+\ensuremath{_{1}}) no parâmetro (c\ensuremath{^{2}}). 2. **Lema poda–órbita:** provar que LSB da órbita (z\ensuremath{^{2}}) gera só palavras em (\ensuremath{\Sigma}). 3. **Lema parábola–reta:** (\ensuremath{\Phi}) conjugada monótona entre ((c\ensuremath{^{2}},a)) e (s). 4. **Teorema produto:** (\ensuremath{\pi}(w)) isomorfismo semigrupal para palavras admissíveis $\times$ Cantor ternário.
\prov{relações: parte\_de koch\_codifica\_cantorfibonacci\_prova\_e\_evidencia}
\prov{proveniência: KOCH\_CANTOR\_FIB\_PROVA.md \textperiodcentered{} fato \textperiodcentered{} confiança 1.0 \textperiodcentered{} domínio manual \textperiodcentered{} história 5 versões no jornal}

%CRISTAL {"arestas":[{"alvo":"decomposicao_peirce","confianca":1.0,"epistemico":"fato","origem":"wiki","rel":"especializa"},{"alvo":"broca_os","confianca":1.0,"epistemico":"fato","origem":"wiki","rel":"relaciona"},{"alvo":"chave_do_coracao","confianca":1.0,"epistemico":"fato","origem":"wiki","rel":"relaciona"}],"confianca":1.0,"contraexemplos":[],"descricao":"Na decomposição de Peirce (1=Σe_i): a diagonal e_i𝒜e_i são células (idempotentes, o estado); a off-diagonal e_i𝒜e_j são canais (cada um square-zero/nilpotente, o fluxo). O mux-demux isola o canal i→j com contaminação cruzada 0. A dinâmica É o acoplamento — isolar uma face quebra.","epistemico":"fato","exemplos":[],"id":"peirce_celulas_canais","memoria":"semantica","meta":{"autor":"Aarão/broca-so","dominio":"broca_repos","fonte":""},"origem":"wiki","palavras_chave":[],"sinonimos":[],"tipo":"conceito","titulo":"Peirce: Células (brancas) e Canais (negros)"}
\section{Peirce: Células (brancas) e Canais (negros)}
Na decomposição de Peirce (1=\ensuremath{\Sigma}e\_i): a diagonal e\_i\ensuremath{\mathcal{A}}e\_i são células (idempotentes, o estado); a off-diagonal e\_i\ensuremath{\mathcal{A}}e\_j são canais (cada um square-zero/nilpotente, o fluxo). O mux-demux isola o canal i\ensuremath{\to}j com contaminação cruzada 0. A dinâmica É o acoplamento — isolar uma face quebra.
\prov{relações: especializa decomposicao\_peirce; relaciona broca\_os; relaciona chave\_do\_coracao}
\prov{proveniência: wiki \textperiodcentered{} fato \textperiodcentered{} confiança 1.0 \textperiodcentered{} domínio broca\_repos \textperiodcentered{} história 4 versões no jornal}

%CRISTAL {"arestas":[{"alvo":"cifra_de_ouro_branco","confianca":1.0,"epistemico":"fato","origem":"manual","rel":"define"},{"alvo":"utilitarismo","confianca":0.5,"epistemico":"fato","origem":"wiki","rel":"relaciona"},{"alvo":"virtualizacao","confianca":0.5,"epistemico":"fato","origem":"wiki","rel":"relaciona"},{"alvo":"vida-morte","confianca":0.5,"epistemico":"fato","origem":"wiki","rel":"relaciona"},{"alvo":"word","confianca":0.5,"epistemico":"fato","origem":"wiki","rel":"relaciona"},{"alvo":"readme","confianca":0.5,"epistemico":"fato","origem":"wiki","rel":"relaciona"},{"alvo":"telemetria_shadow_producao","confianca":0.5,"epistemico":"fato","origem":"wiki","rel":"relaciona"},{"alvo":"sintaxe","confianca":0.5,"epistemico":"fato","origem":"wiki","rel":"relaciona"},{"alvo":"sequencia_pow_nativa","confianca":0.5,"epistemico":"fato","origem":"wiki","rel":"relaciona"},{"alvo":"pipeline_de_ingestao","confianca":0.5,"epistemico":"fato","origem":"wiki","rel":"relaciona"},{"alvo":"ula_da_broca_arquitetura_minima","confianca":0.5,"epistemico":"fato","origem":"wiki","rel":"relaciona"},{"alvo":"poder_disciplinar_foucault","confianca":0.5,"epistemico":"fato","origem":"wiki","rel":"relaciona"},{"alvo":"vontade_de_poder","confianca":0.5,"epistemico":"fato","origem":"wiki","rel":"relaciona"},{"alvo":"pow_taxa_de_sucesso_do_bit_broca","confianca":0.5,"epistemico":"fato","origem":"wiki","rel":"relaciona"},{"alvo":"semiotica_peirce_triade","confianca":0.5,"epistemico":"fato","origem":"wiki","rel":"relaciona"},{"alvo":"religiao","confianca":0.5,"epistemico":"fato","origem":"wiki","rel":"relaciona"}],"confianca":1.0,"contraexemplos":[],"descricao":"> Hipótese: estados degenerados de interesse têm estrutura compatível com A₂.","epistemico":"fato","exemplos":[],"id":"pell","memoria":"semantica","meta":{"dominio":"manual","duplicata_sugerida":[],"nivel":6,"verificacao":"parte de conceito mais amplo 'pell'"},"origem":"PELL.md","palavras_chave":[],"sinonimos":[],"tipo":"conceito","titulo":"PELL"}
\section{PELL}
> Hipótese: estados degenerados de interesse têm estrutura compatível com A\ensuremath{_{2}}.
\prov{relações: define cifra\_de\_ouro\_branco; relaciona utilitarismo; relaciona virtualizacao; relaciona vida-morte; relaciona word; relaciona readme; relaciona telemetria\_shadow\_producao; relaciona sintaxe; relaciona sequencia\_pow\_nativa; relaciona pipeline\_de\_ingestao; relaciona ula\_da\_broca\_arquitetura\_minima; relaciona poder\_disciplinar\_foucault; relaciona vontade\_de\_poder; relaciona pow\_taxa\_de\_sucesso\_do\_bit\_broca; relaciona semiotica\_peirce\_triade; relaciona religiao}
\prov{proveniência: PELL.md \textperiodcentered{} fato \textperiodcentered{} confiança 1.0 \textperiodcentered{} domínio manual \textperiodcentered{} história 21 versões no jornal}

%CRISTAL {"arestas":[{"alvo":"koch_parabola_v2_sobre_a_curva_desdistorcida","confianca":1.0,"epistemico":"fato","origem":"KOCH_PARABOLA.md","rel":"parte_de"},{"alvo":"validacao_hipotese_h_ressonancia","confianca":1.0,"epistemico":"fato","origem":"POW_METAL_RESSONANCIA_VAL.md","rel":"parte_de"},{"alvo":"utilitarismo","confianca":0.5,"epistemico":"fato","origem":"wiki","rel":"relaciona"},{"alvo":"virtualizacao","confianca":0.5,"epistemico":"fato","origem":"wiki","rel":"relaciona"},{"alvo":"word","confianca":0.5,"epistemico":"fato","origem":"wiki","rel":"relaciona"},{"alvo":"vontade_de_poder","confianca":0.5,"epistemico":"fato","origem":"wiki","rel":"relaciona"}],"confianca":1.0,"contraexemplos":[],"descricao":"Secção de POW_METAL_PRE_FECHO.md.","epistemico":"fato","exemplos":[],"id":"perguntas","memoria":"semantica","meta":{"dominio":"manual","nivel":6,"verificacao":"slug idêntico — enriquecer, não duplicar"},"origem":"POW_METAL_PRE_FECHO.md","palavras_chave":[],"sinonimos":[],"tipo":"conceito","titulo":"Perguntas"}
\section{Perguntas}
Secção de POW\_METAL\_PRE\_FECHO.md.
\prov{relações: parte\_de koch\_parabola\_v2\_sobre\_a\_curva\_desdistorcida; parte\_de validacao\_hipotese\_h\_ressonancia; relaciona utilitarismo; relaciona virtualizacao; relaciona word; relaciona vontade\_de\_poder}
\prov{proveniência: POW\_METAL\_PRE\_FECHO.md \textperiodcentered{} fato \textperiodcentered{} confiança 1.0 \textperiodcentered{} domínio manual \textperiodcentered{} história 16 versões no jornal}

%CRISTAL {"arestas":[{"alvo":"broca-so_cristalizado","confianca":1.0,"epistemico":"fato","origem":"BROCA_CRISTAL.md","rel":"parte_de"}],"confianca":1.0,"contraexemplos":[],"descricao":"1. Koch atlas → SVG/CSV → monitor clássico. 2. JSON, papers, Python de visualização sobre **saída projectada**. 3. Goldchain UI, tesouraria, relatórios (Pell + verify nativo no núcleo; display na borda). 4. Qualquer CPU clássica a **renderizar** o que Koch projectou.","epistemico":"fato","exemplos":[],"id":"permitido_encorajado_na_borda","memoria":"semantica","meta":{"dominio":"manual","nivel":6,"verificacao":"slug idêntico — enriquecer, não duplicar"},"origem":"BROCA_CRISTAL.md","palavras_chave":[],"sinonimos":[],"tipo":"conceito","titulo":"Permitido / encorajado na borda"}
\section{Permitido / encorajado na borda}
1. Koch atlas \ensuremath{\to} SVG/CSV \ensuremath{\to} monitor clássico. 2. JSON, papers, Python de visualização sobre **saída projectada**. 3. Goldchain UI, tesouraria, relatórios (Pell + verify nativo no núcleo; display na borda). 4. Qualquer CPU clássica a **renderizar** o que Koch projectou.
\prov{relações: parte\_de broca-so\_cristalizado}
\prov{proveniência: BROCA\_CRISTAL.md \textperiodcentered{} fato \textperiodcentered{} confiança 1.0 \textperiodcentered{} domínio manual \textperiodcentered{} história 4 versões no jornal}

%CRISTAL {"arestas":[{"alvo":"respiracao_dos_programas","confianca":1.0,"epistemico":"fato","origem":"PARABOLA_ISA.md","rel":"parte_de"},{"alvo":"cifra_ouro_branco_nucleo_espectral","confianca":1.0,"epistemico":"fato","origem":"CIFRA.md","rel":"parte_de"},{"alvo":"ula_da_broca_arquitetura_minima","confianca":1.0,"epistemico":"fato","origem":"README.md","rel":"parte_de"}],"confianca":1.0,"contraexemplos":[],"descricao":"``` arquivo → neurônio → word → cifra (Aₙ) → registradores → ULA → resultado ```","epistemico":"fato","exemplos":[],"id":"pipeline","memoria":"semantica","meta":{"dominio":"manual","nivel":6,"verificacao":"slug idêntico — enriquecer, não duplicar"},"origem":"README.md","palavras_chave":[],"sinonimos":[],"tipo":"conceito","titulo":"🌱 pipeline"}
\section{[semente] pipeline}
``` arquivo \ensuremath{\to} neurônio \ensuremath{\to} word \ensuremath{\to} cifra (A\ensuremath{_{n}}) \ensuremath{\to} registradores \ensuremath{\to} ULA \ensuremath{\to} resultado ```
\prov{relações: parte\_de respiracao\_dos\_programas; parte\_de cifra\_ouro\_branco\_nucleo\_espectral; parte\_de ula\_da\_broca\_arquitetura\_minima}
\prov{proveniência: README.md \textperiodcentered{} fato \textperiodcentered{} confiança 1.0 \textperiodcentered{} domínio manual \textperiodcentered{} história 16 versões no jornal}

%CRISTAL {"arestas":[{"alvo":"broca_os","confianca":0.95,"epistemico":"fato","origem":"paper","rel":"parte_de"},{"alvo":"virtualizacao","confianca":0.5,"epistemico":"fato","origem":"wiki","rel":"relaciona"},{"alvo":"word","confianca":0.5,"epistemico":"fato","origem":"wiki","rel":"relaciona"}],"confianca":1.0,"contraexemplos":[],"descricao":"``` abre projeto → memória → tecidos (.idx) → agente software → gera patch → testes → bench → diff → (paper) ```","epistemico":"fato","exemplos":[],"id":"pipeline_code_refactor_fase_03","memoria":"semantica","meta":{"dominio":"manual","nivel":6,"verificacao":"slug idêntico — enriquecer, não duplicar"},"origem":"CRISTAL_DEV.md","palavras_chave":[],"sinonimos":[],"tipo":"conceito","titulo":"Pipeline `code refactor` (Fase 0.3)"}
\section{Pipeline `code refactor` (Fase 0.3)}
``` abre projeto \ensuremath{\to} memória \ensuremath{\to} tecidos (.idx) \ensuremath{\to} agente software \ensuremath{\to} gera patch \ensuremath{\to} testes \ensuremath{\to} bench \ensuremath{\to} diff \ensuremath{\to} (paper) ```
\prov{relações: parte\_de broca\_os; relaciona virtualizacao; relaciona word}
\prov{proveniência: CRISTAL\_DEV.md \textperiodcentered{} fato \textperiodcentered{} confiança 1.0 \textperiodcentered{} domínio manual \textperiodcentered{} história 5 versões no jornal}

%CRISTAL {"arestas":[{"alvo":"fase3","confianca":1.0,"epistemico":"fato","origem":"paper","rel":"parte_de"},{"alvo":"recursao","confianca":0.5,"epistemico":"fato","origem":"wiki","rel":"relaciona"},{"alvo":"rust","confianca":0.5,"epistemico":"fato","origem":"wiki","rel":"relaciona"},{"alvo":"vida-morte","confianca":0.5,"epistemico":"fato","origem":"wiki","rel":"relaciona"},{"alvo":"virtualizacao","confianca":0.5,"epistemico":"fato","origem":"wiki","rel":"relaciona"},{"alvo":"word","confianca":0.5,"epistemico":"fato","origem":"wiki","rel":"relaciona"},{"alvo":"sequencia_pow_nativa","confianca":0.5,"epistemico":"fato","origem":"wiki","rel":"relaciona"},{"alvo":"vontade_de_poder","confianca":0.5,"epistemico":"fato","origem":"wiki","rel":"relaciona"},{"alvo":"ula_da_broca_arquitetura_minima","confianca":0.5,"epistemico":"fato","origem":"wiki","rel":"relaciona"},{"alvo":"telemetria_shadow_producao","confianca":0.5,"epistemico":"fato","origem":"wiki","rel":"relaciona"},{"alvo":"poder_disciplinar_foucault","confianca":0.5,"epistemico":"fato","origem":"wiki","rel":"relaciona"},{"alvo":"readme","confianca":0.5,"epistemico":"fato","origem":"wiki","rel":"relaciona"},{"alvo":"vetores","confianca":0.5,"epistemico":"fato","origem":"wiki","rel":"relaciona"}],"confianca":1.0,"contraexemplos":[],"descricao":"tiffany.ingest(\"documentacao/markdown\", [\"README.md\"], root=\"base\")","epistemico":"fato","exemplos":[],"id":"pipeline_de_ingestao","memoria":"semantica","meta":{"dominio":"manual","duplicata_sugerida":[],"nivel":6,"verificacao":"parte de conceito mais amplo 'pipeline_de_ingestao'"},"origem":"INSTALL.md","palavras_chave":[],"sinonimos":[],"tipo":"conceito","titulo":"Pipeline de ingestão"}
\section{Pipeline de ingestão}
tiffany.ingest("documentacao/markdown", ["README.md"], root="base")
\prov{relações: parte\_de fase3; relaciona recursao; relaciona rust; relaciona vida-morte; relaciona virtualizacao; relaciona word; relaciona sequencia\_pow\_nativa; relaciona vontade\_de\_poder; relaciona ula\_da\_broca\_arquitetura\_minima; relaciona telemetria\_shadow\_producao; relaciona poder\_disciplinar\_foucault; relaciona readme; relaciona vetores}
\prov{proveniência: INSTALL.md \textperiodcentered{} fato \textperiodcentered{} confiança 1.0 \textperiodcentered{} domínio manual \textperiodcentered{} história 16 versões no jornal}

%CRISTAL {"arestas":[{"alvo":"koch_codifica_cantorfibonacci_prova_e_evidencia","confianca":1.0,"epistemico":"fato","origem":"KOCH_CANTOR_FIB_PROVA.md","rel":"parte_de"}],"confianca":1.0,"contraexemplos":[],"descricao":"``` Word w  →  gold(w)  →  parabola (c²,a)  →  s = parabola_para_reta(c²) →  z_canal  →  z² orbita  →  endereco (σ₀σ₁…σₙ) ```","epistemico":"fato","exemplos":[],"id":"pipeline_implementado","memoria":"semantica","meta":{"dominio":"manual","nivel":6,"verificacao":"slug idêntico — enriquecer, não duplicar"},"origem":"KOCH_CANTOR_FIB_PROVA.md","palavras_chave":[],"sinonimos":[],"tipo":"conceito","titulo":"Pipeline (implementado)"}
\section{Pipeline (implementado)}
``` Word w  \ensuremath{\to}  gold(w)  \ensuremath{\to}  parabola (c\ensuremath{^{2}},a)  \ensuremath{\to}  s = parabola\_para\_reta(c\ensuremath{^{2}}) \ensuremath{\to}  z\_canal  \ensuremath{\to}  z\ensuremath{^{2}} orbita  \ensuremath{\to}  endereco (\ensuremath{\sigma}\ensuremath{_{0}}\ensuremath{\sigma}\ensuremath{_{1}}…\ensuremath{\sigma}\ensuremath{_{n}}) ```
\prov{relações: parte\_de koch\_codifica\_cantorfibonacci\_prova\_e\_evidencia}
\prov{proveniência: KOCH\_CANTOR\_FIB\_PROVA.md \textperiodcentered{} fato \textperiodcentered{} confiança 1.0 \textperiodcentered{} domínio manual \textperiodcentered{} história 4 versões no jornal}

%CRISTAL {"arestas":[{"alvo":"koch_atlas_experimento_de_parametrizacao","confianca":1.0,"epistemico":"fato","origem":"KOCH_ATLAS.md","rel":"parte_de"}],"confianca":1.0,"contraexemplos":[],"descricao":"```text Header ↓ Neurônio ↓ Word [total,e]          ← célula branca e₀ ↓ gold(word)              ← célula branca e₁ ↓ parabola_auditar(w,ouro) ← canal negro E₀₁ ↓ canal → z² em Z[i] ↓ Koch Atlas  ← koch_coord(w, ouro) ↓ (endereço fractal, ramo, profundidade, atrator) ↓ Lemniscata → Fecho → SHA  (não medido neste experimento) ```","epistemico":"fato","exemplos":[],"id":"pipeline_medido","memoria":"semantica","meta":{"dominio":"manual","nivel":6,"verificacao":"slug idêntico — enriquecer, não duplicar"},"origem":"KOCH_ATLAS.md","palavras_chave":[],"sinonimos":[],"tipo":"conceito","titulo":"Pipeline medido"}
\section{Pipeline medido}
```text Header \ensuremath{\downarrow} Neurônio \ensuremath{\downarrow} Word [total,e]          \ensuremath{\leftarrow} célula branca e\ensuremath{_{0}} \ensuremath{\downarrow} gold(word)              \ensuremath{\leftarrow} célula branca e\ensuremath{_{1}} \ensuremath{\downarrow} parabola\_auditar(w,ouro) \ensuremath{\leftarrow} canal negro E\ensuremath{_{01}} \ensuremath{\downarrow} canal \ensuremath{\to} z\ensuremath{^{2}} em Z[i] \ensuremath{\downarrow} Koch Atlas  \ensuremath{\leftarrow} koch\_coord(w, ouro) \ensuremath{\downarrow} (endereço fractal, ramo, profundidade, atrator) \ensuremath{\downarrow} Lemniscata \ensuremath{\to} Fecho \ensuremath{\to} SHA  (não medido neste experimento) ```
\prov{relações: parte\_de koch\_atlas\_experimento\_de\_parametrizacao}
\prov{proveniência: KOCH\_ATLAS.md \textperiodcentered{} fato \textperiodcentered{} confiança 1.0 \textperiodcentered{} domínio manual \textperiodcentered{} história 6 versões no jornal}

%CRISTAL {"arestas":[{"alvo":"incerteza_algebrica","confianca":1.0,"epistemico":"fato","origem":"paper","rel":"parte_de"},{"alvo":"mecanica_quantica_de_ouro","confianca":1.0,"epistemico":"fato","origem":"wiki","rel":"usa"},{"alvo":"corpo_dual","confianca":1.0,"epistemico":"fato","origem":"wiki","rel":"relaciona"}],"confianca":1.0,"contraexemplos":[],"descricao":"A pressão algébrica Π=1−x² muda de sinal em |x|=1: o interior (|x|<1) é dissipativo/termodinâmico (norma contrai); a fronteira (|x|=1, s=±1, a=0) é o ponto hermitiano onde a norma é exatamente conservada e a mecânica quântica vive. O eixo dissipativo↔hermitiano é a parábola.","epistemico":"fato","exemplos":[],"id":"ponto_hermitiano","memoria":"semantica","meta":{"autor":"Aarão/broca-so","dominio":"broca_repos","fonte":""},"origem":"wiki","palavras_chave":[],"sinonimos":[],"tipo":"conceito","titulo":"O Ponto Hermitiano"}
\section{O Ponto Hermitiano}
A pressão algébrica \ensuremath{\Pi}=1\ensuremath{-}x\ensuremath{^{2}} muda de sinal em |x|=1: o interior (|x|<1) é dissipativo/termodinâmico (norma contrai); a fronteira (|x|=1, s=$\pm$1, a=0) é o ponto hermitiano onde a norma é exatamente conservada e a mecânica quântica vive. O eixo dissipativo\ensuremath{\leftrightarrow}hermitiano é a parábola.
\prov{relações: parte\_de incerteza\_algebrica; usa mecanica\_quantica\_de\_ouro; relaciona corpo\_dual}
\prov{proveniência: wiki \textperiodcentered{} fato \textperiodcentered{} confiança 1.0 \textperiodcentered{} domínio broca\_repos \textperiodcentered{} história 5 versões no jornal}

%CRISTAL {"arestas":[{"alvo":"readme_cristal","confianca":1.0,"epistemico":"fato","origem":"README_CRISTAL.md","rel":"parte_de"}],"confianca":1.0,"contraexemplos":[],"descricao":"| Ferramenta | Função | |------------|--------| | libtiffany | Engine (overlap, índice, mmap) — **congelada** | | **cristal** | Administra conhecimento | | tiffany | Conversa com o usuário | | broca | Orquestra agentes e sistema |","epistemico":"fato","exemplos":[],"id":"por_que_existe","memoria":"semantica","meta":{"dominio":"manual","nivel":6,"verificacao":"slug idêntico — enriquecer, não duplicar"},"origem":"README_CRISTAL.md","palavras_chave":[],"sinonimos":[],"tipo":"conceito","titulo":"Por que existe"}
\section{Por que existe}
| Ferramenta | Função | |------------|--------| | libtiffany | Engine (overlap, índice, mmap) — **congelada** | | **cristal** | Administra conhecimento | | tiffany | Conversa com o usuário | | broca | Orquestra agentes e sistema |
\prov{relações: parte\_de readme\_cristal}
\prov{proveniência: README\_CRISTAL.md \textperiodcentered{} fato \textperiodcentered{} confiança 1.0 \textperiodcentered{} domínio manual \textperiodcentered{} história 4 versões no jornal}

%CRISTAL {"arestas":[{"alvo":"bit","confianca":1.0,"epistemico":"fato","origem":"manual","rel":"explica"},{"alvo":"transcricao","confianca":0.5,"epistemico":"fato","origem":"wiki","rel":"relaciona"},{"alvo":"pow_gato","confianca":0.5,"epistemico":"fato","origem":"wiki","rel":"relaciona"},{"alvo":"regua_associativa_e_o_risco_de_vazamento","confianca":0.5,"epistemico":"fato","origem":"wiki","rel":"relaciona"},{"alvo":"virtualizacao","confianca":0.5,"epistemico":"fato","origem":"wiki","rel":"relaciona"},{"alvo":"word","confianca":0.5,"epistemico":"fato","origem":"wiki","rel":"relaciona"},{"alvo":"vida-morte","confianca":0.5,"epistemico":"fato","origem":"wiki","rel":"relaciona"}],"confianca":1.0,"contraexemplos":[],"descricao":"Manual: ula/POW_BIT_TAXA.md","epistemico":"fato","exemplos":[],"id":"pow_bit_taxa","memoria":"semantica","meta":{"arquivo":"ula/POW_BIT_TAXA.md","dominio":"manual","nivel":6},"origem":"manual","palavras_chave":[],"sinonimos":[],"tipo":"conceito","titulo":"POW_BIT_TAXA"}
\section{POW\_BIT\_TAXA}
Manual: ula/POW\_BIT\_TAXA.md
\prov{relações: explica bit; relaciona transcricao; relaciona pow\_gato; relaciona regua\_associativa\_e\_o\_risco\_de\_vazamento; relaciona virtualizacao; relaciona word; relaciona vida-morte}
\prov{proveniência: manual \textperiodcentered{} fato \textperiodcentered{} confiança 1.0 \textperiodcentered{} domínio manual \textperiodcentered{} história 37 versões no jornal}

%CRISTAL {"arestas":[{"alvo":"pow_geometria","confianca":1.0,"epistemico":"fato","origem":"manual","rel":"parte_de"},{"alvo":"virtualizacao","confianca":0.5,"epistemico":"fato","origem":"wiki","rel":"relaciona"},{"alvo":"word","confianca":0.5,"epistemico":"fato","origem":"wiki","rel":"relaciona"},{"alvo":"vida-morte","confianca":0.5,"epistemico":"fato","origem":"wiki","rel":"relaciona"},{"alvo":"semiotica_peirce_triade","confianca":0.5,"epistemico":"fato","origem":"wiki","rel":"relaciona"}],"confianca":1.0,"contraexemplos":[],"descricao":"> **Pergunta:** existe mudança de coordenadas que organize melhor o espaço de busca?","epistemico":"fato","exemplos":[],"id":"pow_coordenadas_espectrais_broca","memoria":"semantica","meta":{"dominio":"manual","nivel":6,"verificacao":"slug idêntico — enriquecer, não duplicar"},"origem":"POW_GEOMETRIA.md","palavras_chave":[],"sinonimos":[],"tipo":"conceito","titulo":"PoW — Coordenadas Espectrais (Broca)"}
\section{PoW — Coordenadas Espectrais (Broca)}
> **Pergunta:** existe mudança de coordenadas que organize melhor o espaço de busca?
\prov{relações: parte\_de pow\_geometria; relaciona virtualizacao; relaciona word; relaciona vida-morte; relaciona semiotica\_peirce\_triade}
\prov{proveniência: POW\_GEOMETRIA.md \textperiodcentered{} fato \textperiodcentered{} confiança 1.0 \textperiodcentered{} domínio manual \textperiodcentered{} história 9 versões no jornal}

%CRISTAL {"arestas":[{"alvo":"gato","confianca":1.0,"epistemico":"fato","origem":"manual","rel":"confirma"},{"alvo":"pow_geometria","confianca":0.5,"epistemico":"fato","origem":"wiki","rel":"relaciona"},{"alvo":"python","confianca":0.5,"epistemico":"fato","origem":"wiki","rel":"relaciona"},{"alvo":"teste_ordem","confianca":0.5,"epistemico":"fato","origem":"wiki","rel":"relaciona"},{"alvo":"prova_isa","confianca":0.5,"epistemico":"fato","origem":"wiki","rel":"relaciona"},{"alvo":"sessao_interativa","confianca":0.5,"epistemico":"fato","origem":"wiki","rel":"relaciona"},{"alvo":"word","confianca":0.5,"epistemico":"fato","origem":"wiki","rel":"relaciona"},{"alvo":"virtualizacao","confianca":0.5,"epistemico":"fato","origem":"wiki","rel":"relaciona"}],"confianca":1.0,"contraexemplos":[],"descricao":"Manual: ula/POW_ESPECTRAL.md","epistemico":"fato","exemplos":[],"id":"pow_espectral","memoria":"semantica","meta":{"arquivo":"ula/POW_ESPECTRAL.md","dominio":"manual","nivel":6},"origem":"manual","palavras_chave":[],"sinonimos":[],"tipo":"conceito","titulo":"POW_ESPECTRAL"}
\section{POW\_ESPECTRAL}
Manual: ula/POW\_ESPECTRAL.md
\prov{relações: confirma gato; relaciona pow\_geometria; relaciona python; relaciona teste\_ordem; relaciona prova\_isa; relaciona sessao\_interativa; relaciona word; relaciona virtualizacao}
\prov{proveniência: manual \textperiodcentered{} fato \textperiodcentered{} confiança 1.0 \textperiodcentered{} domínio manual \textperiodcentered{} história 11 versões no jornal}

%CRISTAL {"arestas":[{"alvo":"geometria","confianca":0.289,"epistemico":"fato","origem":"manual","rel":"referencia"},{"alvo":"pow_espectral","confianca":1.0,"epistemico":"fato","origem":"manual","rel":"parte_de"},{"alvo":"transformacao_linear","confianca":1.0,"epistemico":"fato","origem":"manual","rel":"usa"},{"alvo":"word","confianca":0.5,"epistemico":"fato","origem":"wiki","rel":"relaciona"},{"alvo":"vontade_de_poder","confianca":0.5,"epistemico":"fato","origem":"wiki","rel":"relaciona"}],"confianca":1.0,"contraexemplos":[],"descricao":"> **Pergunta:** como o Proof-of-Work se apresenta na geometria espectral?","epistemico":"fato","exemplos":[],"id":"pow_espectral_investigacao_broca","memoria":"semantica","meta":{"dominio":"manual","nivel":6,"verificacao":"slug idêntico — enriquecer, não duplicar"},"origem":"POW_ESPECTRAL.md","palavras_chave":[],"sinonimos":[],"tipo":"conceito","titulo":"PoW Espectral — Investigação Broca"}
\section{PoW Espectral — Investigação Broca}
> **Pergunta:** como o Proof-of-Work se apresenta na geometria espectral?
\prov{relações: referencia geometria; parte\_de pow\_espectral; usa transformacao\_linear; relaciona word; relaciona vontade\_de\_poder}
\prov{proveniência: POW\_ESPECTRAL.md \textperiodcentered{} fato \textperiodcentered{} confiança 1.0 \textperiodcentered{} domínio manual \textperiodcentered{} história 61 versões no jornal}

%CRISTAL {"arestas":[{"alvo":"geometria","confianca":1.0,"epistemico":"fato","origem":"manual","rel":"explica"},{"alvo":"python","confianca":0.5,"epistemico":"fato","origem":"wiki","rel":"relaciona"},{"alvo":"transcricao","confianca":0.5,"epistemico":"fato","origem":"wiki","rel":"relaciona"},{"alvo":"teste_ordem","confianca":0.5,"epistemico":"fato","origem":"wiki","rel":"relaciona"},{"alvo":"pow_metal_setor","confianca":0.5,"epistemico":"fato","origem":"wiki","rel":"relaciona"},{"alvo":"projeto_cristaljson","confianca":0.5,"epistemico":"fato","origem":"wiki","rel":"relaciona"},{"alvo":"topologia","confianca":0.5,"epistemico":"fato","origem":"wiki","rel":"relaciona"},{"alvo":"recursao","confianca":0.5,"epistemico":"fato","origem":"wiki","rel":"relaciona"},{"alvo":"pow_metal_funil","confianca":0.5,"epistemico":"fato","origem":"wiki","rel":"relaciona"},{"alvo":"rust","confianca":0.5,"epistemico":"fato","origem":"wiki","rel":"relaciona"},{"alvo":"pow_metal_pre_fecho","confianca":0.5,"epistemico":"fato","origem":"wiki","rel":"relaciona"},{"alvo":"tiffany_cabeca_broca","confianca":0.5,"epistemico":"fato","origem":"wiki","rel":"relaciona"}],"confianca":1.0,"contraexemplos":[],"descricao":"Manual: ula/POW_GEOMETRIA.md","epistemico":"fato","exemplos":[],"id":"pow_geometria","memoria":"semantica","meta":{"arquivo":"ula/POW_GEOMETRIA.md","dominio":"manual","nivel":6},"origem":"manual","palavras_chave":[],"sinonimos":[],"tipo":"conceito","titulo":"POW_GEOMETRIA"}
\section{POW\_GEOMETRIA}
Manual: ula/POW\_GEOMETRIA.md
\prov{relações: explica geometria; relaciona python; relaciona transcricao; relaciona teste\_ordem; relaciona pow\_metal\_setor; relaciona projeto\_cristaljson; relaciona topologia; relaciona recursao; relaciona pow\_metal\_funil; relaciona rust; relaciona pow\_metal\_pre\_fecho; relaciona tiffany\_cabeca\_broca}
\prov{proveniência: manual \textperiodcentered{} fato \textperiodcentered{} confiança 1.0 \textperiodcentered{} domínio manual \textperiodcentered{} história 42 versões no jornal}

%CRISTAL {"arestas":[{"alvo":"pow_metal","confianca":1.0,"epistemico":"fato","origem":"manual","rel":"parte_de"},{"alvo":"word","confianca":0.5,"epistemico":"fato","origem":"wiki","rel":"relaciona"}],"confianca":1.0,"contraexemplos":[],"descricao":"> PoW no **corpo Z[i] resolvido** entre reta real e reta de ouro. > Koch codifica a **interface** entre os fractais das duas retas (z↦z², 2-para-1).","epistemico":"fato","exemplos":[],"id":"pow_lemniscata_koch_corpo_de_gauss","memoria":"semantica","meta":{"dominio":"manual","nivel":6,"verificacao":"slug idêntico — enriquecer, não duplicar"},"origem":"POW_METAL.md","palavras_chave":[],"sinonimos":[],"tipo":"conceito","titulo":"PoW — Lemniscata + Koch (corpo de Gauss)"}
\section{PoW — Lemniscata + Koch (corpo de Gauss)}
> PoW no **corpo Z[i] resolvido** entre reta real e reta de ouro. > Koch codifica a **interface** entre os fractais das duas retas (z\ensuremath{\mapsto}z\ensuremath{^{2}}, 2-para-1).
\prov{relações: parte\_de pow\_metal; relaciona word}
\prov{proveniência: POW\_METAL.md \textperiodcentered{} fato \textperiodcentered{} confiança 1.0 \textperiodcentered{} domínio manual \textperiodcentered{} história 5 versões no jornal}

%CRISTAL {"arestas":[{"alvo":"termodinamica","confianca":1.0,"epistemico":"fato","origem":"manual","rel":"referencia"},{"alvo":"vida-morte","confianca":0.5,"epistemico":"fato","origem":"wiki","rel":"relaciona"},{"alvo":"virtualizacao","confianca":0.5,"epistemico":"fato","origem":"wiki","rel":"relaciona"}],"confianca":1.0,"contraexemplos":[],"descricao":"Busca **vertical**: modos vibrando no nó a+c²=1 (vida.tex). > Mede se acoplamento prediz fecho/SHA **antes** da prova irreversível.","epistemico":"fato","exemplos":[],"id":"pow_metal","memoria":"semantica","meta":{"dominio":"manual","duplicata_sugerida":[],"nivel":6,"verificacao":"parte de conceito mais amplo 'pow_metal'"},"origem":"POW_METAL_RESSONANCIA.md","palavras_chave":[],"sinonimos":[],"tipo":"conceito","titulo":"POW_METAL"}
\section{POW\_METAL}
Busca **vertical**: modos vibrando no nó a+c\ensuremath{^{2}}=1 (vida.tex). > Mede se acoplamento prediz fecho/SHA **antes** da prova irreversível.
\prov{relações: referencia termodinamica; relaciona vida-morte; relaciona virtualizacao}
\prov{proveniência: POW\_METAL\_RESSONANCIA.md \textperiodcentered{} fato \textperiodcentered{} confiança 1.0 \textperiodcentered{} domínio manual \textperiodcentered{} história 10 versões no jornal}

%CRISTAL {"arestas":[{"alvo":"goldchain","confianca":1.0,"epistemico":"fato","origem":"manual","rel":"referencia"},{"alvo":"transcricao","confianca":0.5,"epistemico":"fato","origem":"wiki","rel":"relaciona"},{"alvo":"pow_metal_setor","confianca":0.5,"epistemico":"fato","origem":"wiki","rel":"relaciona"},{"alvo":"vida-morte","confianca":0.5,"epistemico":"fato","origem":"wiki","rel":"relaciona"},{"alvo":"teste_ordem","confianca":0.5,"epistemico":"fato","origem":"wiki","rel":"relaciona"},{"alvo":"resposta_a_pergunta_dificil","confianca":0.5,"epistemico":"fato","origem":"wiki","rel":"relaciona"},{"alvo":"pow_metal_pre_fecho","confianca":0.5,"epistemico":"fato","origem":"wiki","rel":"relaciona"},{"alvo":"testes","confianca":0.5,"epistemico":"fato","origem":"wiki","rel":"relaciona"},{"alvo":"pow_metal_ressonancia","confianca":0.5,"epistemico":"fato","origem":"wiki","rel":"relaciona"},{"alvo":"thread","confianca":0.5,"epistemico":"fato","origem":"wiki","rel":"relaciona"}],"confianca":1.0,"contraexemplos":[],"descricao":"Manual: ula/POW_METAL_FUNIL.md","epistemico":"fato","exemplos":[],"id":"pow_metal_funil","memoria":"semantica","meta":{"arquivo":"ula/POW_METAL_FUNIL.md","dominio":"manual","nivel":6},"origem":"manual","palavras_chave":[],"sinonimos":[],"tipo":"conceito","titulo":"POW_METAL_FUNIL"}
\section{POW\_METAL\_FUNIL}
Manual: ula/POW\_METAL\_FUNIL.md
\prov{relações: referencia goldchain; relaciona transcricao; relaciona pow\_metal\_setor; relaciona vida-morte; relaciona teste\_ordem; relaciona resposta\_a\_pergunta\_dificil; relaciona pow\_metal\_pre\_fecho; relaciona testes; relaciona pow\_metal\_ressonancia; relaciona thread}
\prov{proveniência: manual \textperiodcentered{} fato \textperiodcentered{} confiança 1.0 \textperiodcentered{} domínio manual \textperiodcentered{} história 13 versões no jornal}

%CRISTAL {"arestas":[{"alvo":"morte","confianca":1.0,"epistemico":"fato","origem":"manual","rel":"referencia"},{"alvo":"transcricao","confianca":0.5,"epistemico":"fato","origem":"wiki","rel":"relaciona"},{"alvo":"vida-morte","confianca":0.5,"epistemico":"fato","origem":"wiki","rel":"relaciona"},{"alvo":"resposta_a_pergunta_dificil","confianca":0.5,"epistemico":"fato","origem":"wiki","rel":"relaciona"},{"alvo":"pow_metal_setor","confianca":0.5,"epistemico":"fato","origem":"wiki","rel":"relaciona"},{"alvo":"teste_ordem","confianca":0.5,"epistemico":"fato","origem":"wiki","rel":"relaciona"},{"alvo":"topologia","confianca":0.5,"epistemico":"fato","origem":"wiki","rel":"relaciona"},{"alvo":"projeto_cristaljson","confianca":0.5,"epistemico":"fato","origem":"wiki","rel":"relaciona"},{"alvo":"testes","confianca":0.5,"epistemico":"fato","origem":"wiki","rel":"relaciona"},{"alvo":"processos","confianca":0.5,"epistemico":"fato","origem":"wiki","rel":"relaciona"},{"alvo":"pow_metal_ressonancia","confianca":0.5,"epistemico":"fato","origem":"wiki","rel":"relaciona"},{"alvo":"rust","confianca":0.5,"epistemico":"fato","origem":"wiki","rel":"relaciona"},{"alvo":"pow_metal_ressonancia_val","confianca":0.5,"epistemico":"fato","origem":"wiki","rel":"relaciona"},{"alvo":"virtualizacao","confianca":0.5,"epistemico":"fato","origem":"wiki","rel":"relaciona"},{"alvo":"regua_associativa_e_o_risco_de_vazamento","confianca":0.5,"epistemico":"fato","origem":"wiki","rel":"relaciona"}],"confianca":1.0,"contraexemplos":[],"descricao":"Manual: ula/POW_METAL_PRE_FECHO.md","epistemico":"fato","exemplos":[],"id":"pow_metal_pre_fecho","memoria":"semantica","meta":{"arquivo":"ula/POW_METAL_PRE_FECHO.md","dominio":"manual","nivel":6},"origem":"manual","palavras_chave":[],"sinonimos":[],"tipo":"conceito","titulo":"POW_METAL_PRE_FECHO"}
\section{POW\_METAL\_PRE\_FECHO}
Manual: ula/POW\_METAL\_PRE\_FECHO.md
\prov{relações: referencia morte; relaciona transcricao; relaciona vida-morte; relaciona resposta\_a\_pergunta\_dificil; relaciona pow\_metal\_setor; relaciona teste\_ordem; relaciona topologia; relaciona projeto\_cristaljson; relaciona testes; relaciona processos; relaciona pow\_metal\_ressonancia; relaciona rust; relaciona pow\_metal\_ressonancia\_val; relaciona virtualizacao; relaciona regua\_associativa\_e\_o\_risco\_de\_vazamento}
\prov{proveniência: manual \textperiodcentered{} fato \textperiodcentered{} confiança 1.0 \textperiodcentered{} domínio manual \textperiodcentered{} história 18 versões no jornal}

%CRISTAL {"arestas":[{"alvo":"vida","confianca":1.0,"epistemico":"fato","origem":"manual","rel":"referencia"},{"alvo":"teste_ordem","confianca":0.5,"epistemico":"fato","origem":"wiki","rel":"relaciona"},{"alvo":"transcricao","confianca":0.5,"epistemico":"fato","origem":"wiki","rel":"relaciona"},{"alvo":"resposta_a_pergunta_dificil","confianca":0.5,"epistemico":"fato","origem":"wiki","rel":"relaciona"},{"alvo":"pow_metal_setor","confianca":0.5,"epistemico":"fato","origem":"wiki","rel":"relaciona"},{"alvo":"vida-morte","confianca":0.5,"epistemico":"fato","origem":"wiki","rel":"relaciona"}],"confianca":1.0,"contraexemplos":[],"descricao":"Manual: ula/POW_METAL_RESSONANCIA.md","epistemico":"fato","exemplos":[],"id":"pow_metal_ressonancia","memoria":"semantica","meta":{"arquivo":"ula/POW_METAL_RESSONANCIA.md","dominio":"manual","nivel":6},"origem":"manual","palavras_chave":[],"sinonimos":[],"tipo":"conceito","titulo":"POW_METAL_RESSONANCIA"}
\section{POW\_METAL\_RESSONANCIA}
Manual: ula/POW\_METAL\_RESSONANCIA.md
\prov{relações: referencia vida; relaciona teste\_ordem; relaciona transcricao; relaciona resposta\_a\_pergunta\_dificil; relaciona pow\_metal\_setor; relaciona vida-morte}
\prov{proveniência: manual \textperiodcentered{} fato \textperiodcentered{} confiança 1.0 \textperiodcentered{} domínio manual \textperiodcentered{} história 9 versões no jornal}

%CRISTAL {"arestas":[{"alvo":"valor","confianca":1.0,"epistemico":"fato","origem":"manual","rel":"referencia"},{"alvo":"virtualizacao","confianca":0.5,"epistemico":"fato","origem":"wiki","rel":"relaciona"},{"alvo":"transcricao","confianca":0.5,"epistemico":"fato","origem":"wiki","rel":"relaciona"},{"alvo":"vida-morte","confianca":0.5,"epistemico":"fato","origem":"wiki","rel":"relaciona"},{"alvo":"ruido","confianca":0.5,"epistemico":"fato","origem":"wiki","rel":"relaciona"},{"alvo":"scheduler","confianca":0.5,"epistemico":"fato","origem":"wiki","rel":"relaciona"},{"alvo":"resultado_anterior_identidade_nao_cobertura","confianca":0.5,"epistemico":"fato","origem":"wiki","rel":"relaciona"},{"alvo":"utilitarismo","confianca":0.5,"epistemico":"fato","origem":"wiki","rel":"relaciona"},{"alvo":"veredito_experimental","confianca":0.5,"epistemico":"fato","origem":"wiki","rel":"relaciona"},{"alvo":"prova_da_maquina_espectral","confianca":0.5,"epistemico":"fato","origem":"wiki","rel":"relaciona"},{"alvo":"tipos_de_interpretante","confianca":0.5,"epistemico":"fato","origem":"wiki","rel":"relaciona"}],"confianca":1.0,"contraexemplos":[],"descricao":"Manual: ula/POW_METAL_RESSONANCIA_VAL.md","epistemico":"fato","exemplos":[],"id":"pow_metal_ressonancia_val","memoria":"semantica","meta":{"arquivo":"ula/POW_METAL_RESSONANCIA_VAL.md","dominio":"manual","nivel":6},"origem":"manual","palavras_chave":[],"sinonimos":[],"tipo":"conceito","titulo":"POW_METAL_RESSONANCIA_VAL"}
\section{POW\_METAL\_RESSONANCIA\_VAL}
Manual: ula/POW\_METAL\_RESSONANCIA\_VAL.md
\prov{relações: referencia valor; relaciona virtualizacao; relaciona transcricao; relaciona vida-morte; relaciona ruido; relaciona scheduler; relaciona resultado\_anterior\_identidade\_nao\_cobertura; relaciona utilitarismo; relaciona veredito\_experimental; relaciona prova\_da\_maquina\_espectral; relaciona tipos\_de\_interpretante}
\prov{proveniência: manual \textperiodcentered{} fato \textperiodcentered{} confiança 1.0 \textperiodcentered{} domínio manual \textperiodcentered{} história 14 versões no jornal}

%CRISTAL {"arestas":[{"alvo":"mecanica_quantica","confianca":1.0,"epistemico":"fato","origem":"manual","rel":"referencia"},{"alvo":"transcricao","confianca":0.5,"epistemico":"fato","origem":"wiki","rel":"relaciona"},{"alvo":"vida-morte","confianca":0.5,"epistemico":"fato","origem":"wiki","rel":"relaciona"},{"alvo":"resposta_a_pergunta_dificil","confianca":0.5,"epistemico":"fato","origem":"wiki","rel":"relaciona"},{"alvo":"testes","confianca":0.5,"epistemico":"fato","origem":"wiki","rel":"relaciona"},{"alvo":"teste_ordem","confianca":0.5,"epistemico":"fato","origem":"wiki","rel":"relaciona"},{"alvo":"topologia","confianca":0.5,"epistemico":"fato","origem":"wiki","rel":"relaciona"}],"confianca":1.0,"contraexemplos":[],"descricao":"Manual: ula/POW_METAL_SETOR.md","epistemico":"fato","exemplos":[],"id":"pow_metal_setor","memoria":"semantica","meta":{"arquivo":"ula/POW_METAL_SETOR.md","dominio":"manual","nivel":6},"origem":"manual","palavras_chave":[],"sinonimos":[],"tipo":"conceito","titulo":"POW_METAL_SETOR"}
\section{POW\_METAL\_SETOR}
Manual: ula/POW\_METAL\_SETOR.md
\prov{relações: referencia mecanica\_quantica; relaciona transcricao; relaciona vida-morte; relaciona resposta\_a\_pergunta\_dificil; relaciona testes; relaciona teste\_ordem; relaciona topologia}
\prov{proveniência: manual \textperiodcentered{} fato \textperiodcentered{} confiança 1.0 \textperiodcentered{} domínio manual \textperiodcentered{} história 10 versões no jornal}

%CRISTAL {"arestas":[{"alvo":"pow_bit_taxa","confianca":1.0,"epistemico":"fato","origem":"manual","rel":"parte_de"},{"alvo":"word","confianca":0.5,"epistemico":"fato","origem":"wiki","rel":"relaciona"},{"alvo":"virtualizacao","confianca":0.5,"epistemico":"fato","origem":"wiki","rel":"relaciona"},{"alvo":"vida-morte","confianca":0.5,"epistemico":"fato","origem":"wiki","rel":"relaciona"},{"alvo":"readme","confianca":0.5,"epistemico":"fato","origem":"wiki","rel":"relaciona"}],"confianca":1.0,"contraexemplos":[],"descricao":"> **Pergunta:** o neurônio + parábola funciona como filtro barato antes do SHA?","epistemico":"fato","exemplos":[],"id":"pow_taxa_de_sucesso_do_bit_broca","memoria":"semantica","meta":{"dominio":"manual","nivel":6,"verificacao":"slug idêntico — enriquecer, não duplicar"},"origem":"POW_BIT_TAXA.md","palavras_chave":[],"sinonimos":[],"tipo":"conceito","titulo":"PoW — Taxa de Sucesso do Bit (Broca)"}
\section{PoW — Taxa de Sucesso do Bit (Broca)}
> **Pergunta:** o neurônio + parábola funciona como filtro barato antes do SHA?
\prov{relações: parte\_de pow\_bit\_taxa; relaciona word; relaciona virtualizacao; relaciona vida-morte; relaciona readme}
\prov{proveniência: POW\_BIT\_TAXA.md \textperiodcentered{} fato \textperiodcentered{} confiança 1.0 \textperiodcentered{} domínio manual \textperiodcentered{} história 8 versões no jornal}

%CRISTAL {"arestas":[{"alvo":"cifra_ouro_branco_nucleo_espectral","confianca":1.0,"epistemico":"fato","origem":"CIFRA.md","rel":"parte_de"}],"confianca":1.0,"contraexemplos":[],"descricao":"| Programa | Física | Resultado | |----------|--------|-----------| | `01_soma.asm` | LOAD+ADD | combina `[2,0]+[2,1]=[4,1]` | | `09_pell.asm` | SILVER×4 | evolui `(1,0)→(29,12)` | | `11_acum_pell.asm` | loop SILVER+JNZ | idêntico a 09, N escalável | | `10_spect.asm` | SPECT+LOAD | `'A'`→`[4,2]` automático |","epistemico":"fato","exemplos":[],"id":"programas_espectrais","memoria":"semantica","meta":{"dominio":"manual","nivel":6,"verificacao":"slug idêntico — enriquecer, não duplicar"},"origem":"CIFRA.md","palavras_chave":[],"sinonimos":[],"tipo":"conceito","titulo":"Programas espectrais"}
\section{Programas espectrais}
| Programa | Física | Resultado | |----------|--------|-----------| | `01\_soma.asm` | LOAD+ADD | combina `[2,0]+[2,1]=[4,1]` | | `09\_pell.asm` | SILVER$\times$4 | evolui `(1,0)\ensuremath{\to}(29,12)` | | `11\_acum\_pell.asm` | loop SILVER+JNZ | idêntico a 09, N escalável | | `10\_spect.asm` | SPECT+LOAD | `'A'`\ensuremath{\to}`[4,2]` automático |
\prov{relações: parte\_de cifra\_ouro\_branco\_nucleo\_espectral}
\prov{proveniência: CIFRA.md \textperiodcentered{} fato \textperiodcentered{} confiança 1.0 \textperiodcentered{} domínio manual \textperiodcentered{} história 6 versões no jornal}

%CRISTAL {"arestas":[{"alvo":"broca-so_cristalizado","confianca":1.0,"epistemico":"fato","origem":"BROCA_CRISTAL.md","rel":"parte_de"}],"confianca":1.0,"contraexemplos":[],"descricao":"1. Duplicar `Morto n ` ou VM cadáver em runtime de produção. 2. FFT/Fourier no espectro gato (mata `e`). 3. SHA imposto no PoW espectral (`POWESPECTRAL.md`). 4. Amputar `e` **antes** do Olho de Koch (perder fase no meio do pipeline).","epistemico":"fato","exemplos":[],"id":"proibido_no_nucleo","memoria":"semantica","meta":{"dominio":"manual","nivel":6,"verificacao":"slug idêntico — enriquecer, não duplicar"},"origem":"BROCA_CRISTAL.md","palavras_chave":[],"sinonimos":[],"tipo":"conceito","titulo":"Proibido no núcleo"}
\section{Proibido no núcleo}
1. Duplicar `Morto n ` ou VM cadáver em runtime de produção. 2. FFT/Fourier no espectro gato (mata `e`). 3. SHA imposto no PoW espectral (`POWESPECTRAL.md`). 4. Amputar `e` **antes** do Olho de Koch (perder fase no meio do pipeline).
\prov{relações: parte\_de broca-so\_cristalizado}
\prov{proveniência: BROCA\_CRISTAL.md \textperiodcentered{} fato \textperiodcentered{} confiança 1.0 \textperiodcentered{} domínio manual \textperiodcentered{} história 5 versões no jornal}

%CRISTAL {"arestas":[{"alvo":"broca_os","confianca":0.95,"epistemico":"fato","origem":"paper","rel":"parte_de"},{"alvo":"topologia","confianca":0.5,"epistemico":"fato","origem":"wiki","rel":"relaciona"},{"alvo":"teste_ordem","confianca":0.5,"epistemico":"fato","origem":"wiki","rel":"relaciona"},{"alvo":"transcricao","confianca":0.5,"epistemico":"fato","origem":"wiki","rel":"relaciona"},{"alvo":"rust","confianca":0.5,"epistemico":"fato","origem":"wiki","rel":"relaciona"},{"alvo":"syscall","confianca":0.5,"epistemico":"fato","origem":"wiki","rel":"relaciona"},{"alvo":"testes","confianca":0.5,"epistemico":"fato","origem":"wiki","rel":"relaciona"},{"alvo":"regua_associativa_e_o_risco_de_vazamento","confianca":0.5,"epistemico":"fato","origem":"wiki","rel":"relaciona"},{"alvo":"thread","confianca":0.5,"epistemico":"fato","origem":"wiki","rel":"relaciona"},{"alvo":"vida-morte","confianca":0.5,"epistemico":"fato","origem":"wiki","rel":"relaciona"},{"alvo":"resposta_a_pergunta_dificil","confianca":0.5,"epistemico":"fato","origem":"wiki","rel":"relaciona"},{"alvo":"tiffany_cabeca_broca","confianca":0.5,"epistemico":"fato","origem":"wiki","rel":"relaciona"}],"confianca":1.0,"contraexemplos":[],"descricao":"Todo repo Tiffany/Broca carrega contexto:","epistemico":"fato","exemplos":[],"id":"projeto_cristaljson","memoria":"semantica","meta":{"dominio":"manual","nivel":6,"verificacao":"slug idêntico — enriquecer, não duplicar"},"origem":"CRISTAL_DEV.md","palavras_chave":[],"sinonimos":[],"tipo":"conceito","titulo":"Projeto (`cristal.json`)"}
\section{Projeto (`cristal.json`)}
Todo repo Tiffany/Broca carrega contexto:
\prov{relações: parte\_de broca\_os; relaciona topologia; relaciona teste\_ordem; relaciona transcricao; relaciona rust; relaciona syscall; relaciona testes; relaciona regua\_associativa\_e\_o\_risco\_de\_vazamento; relaciona thread; relaciona vida-morte; relaciona resposta\_a\_pergunta\_dificil; relaciona tiffany\_cabeca\_broca}
\prov{proveniência: CRISTAL\_DEV.md \textperiodcentered{} fato \textperiodcentered{} confiança 1.0 \textperiodcentered{} domínio manual \textperiodcentered{} história 14 versões no jornal}

%CRISTAL {"arestas":[{"alvo":"o_conselho","confianca":1.0,"epistemico":"fato","origem":"wiki","rel":"parte_de"},{"alvo":"regua_dnci","confianca":1.0,"epistemico":"fato","origem":"wiki","rel":"usa"}],"confianca":1.0,"contraexemplos":[],"descricao":"Nada é cristalizado automaticamente: um achado sobe por Hipótese → Validação → Conselho → Promoção → Cristal, e o cristal (texto canônico) nunca é reescrito sozinho pelo modelo. Só vira fato canônico após ratificação — 'passou no meu teste' ≠ 'o Conselho ratificou'.","epistemico":"fato","exemplos":[],"id":"promocao_hipotese_cristal","memoria":"semantica","meta":{"autor":"Aarão/broca-so","dominio":"broca_repos","fonte":""},"origem":"wiki","palavras_chave":[],"sinonimos":[],"tipo":"conceito","titulo":"Promoção: Hipótese → Cristal"}
\section{Promoção: Hipótese \ensuremath{\to} Cristal}
Nada é cristalizado automaticamente: um achado sobe por Hipótese \ensuremath{\to} Validação \ensuremath{\to} Conselho \ensuremath{\to} Promoção \ensuremath{\to} Cristal, e o cristal (texto canônico) nunca é reescrito sozinho pelo modelo. Só vira fato canônico após ratificação — 'passou no meu teste' \ensuremath{\ne} 'o Conselho ratificou'.
\prov{relações: parte\_de o\_conselho; usa regua\_dnci}
\prov{proveniência: wiki \textperiodcentered{} fato \textperiodcentered{} confiança 1.0 \textperiodcentered{} domínio broca\_repos \textperiodcentered{} história 3 versões no jornal}

%CRISTAL {"arestas":[{"alvo":"campo_psicologico_cuspide","confianca":1.0,"epistemico":"fato","origem":"wiki","rel":"usa"},{"alvo":"cuspide_pisot","confianca":1.0,"epistemico":"fato","origem":"wiki","rel":"relaciona"}],"confianca":1.0,"contraexemplos":[],"descricao":"Oscilador s̈+2γṡ+ω₀²s=0: calma=super-amortecido (γ>1), agitação=sub-amortecido (γ<1, overshoot), prontidão=crítico (γ=1, raiz dupla λ₁=λ₂=a cúspide, retorno mais rápido sem overshoot). O tempo de assentamento é mínimo no crítico e cresce dos dois lados — a lei de Yerkes-Dodson (U invertido).","epistemico":"inferencia","exemplos":[],"id":"prontidao_amortecimento_critico","memoria":"semantica","meta":{"autor":"Lopes/Aarão","dominio":"hiper","fonte":""},"origem":"wiki","palavras_chave":[],"sinonimos":[],"tipo":"conceito","titulo":"A Prontidão como Amortecimento Crítico (Yerkes-Dodson)"}
\section{A Prontidão como Amortecimento Crítico (Yerkes-Dodson)}
Oscilador s\ensuremath{}+2\ensuremath{\gamma}\ensuremath{\dot{s}}+\ensuremath{\omega}\ensuremath{_{0}}\ensuremath{^{2}}s=0: calma=super-amortecido (\ensuremath{\gamma}>1), agitação=sub-amortecido (\ensuremath{\gamma}<1, overshoot), prontidão=crítico (\ensuremath{\gamma}=1, raiz dupla \ensuremath{\lambda}\ensuremath{_{1}}=\ensuremath{\lambda}\ensuremath{_{2}}=a cúspide, retorno mais rápido sem overshoot). O tempo de assentamento é mínimo no crítico e cresce dos dois lados — a lei de Yerkes-Dodson (U invertido).
\prov{relações: usa campo\_psicologico\_cuspide; relaciona cuspide\_pisot}
\prov{proveniência: wiki \textperiodcentered{} inferencia \textperiodcentered{} confiança 1.0 \textperiodcentered{} domínio hiper \textperiodcentered{} história 4 versões no jornal}

%CRISTAL {"arestas":[{"alvo":"koch_codifica_cantorfibonacci_prova_e_evidencia","confianca":1.0,"epistemico":"fato","origem":"KOCH_CANTOR_FIB_PROVA.md","rel":"parte_de"}],"confianca":1.0,"contraexemplos":[],"descricao":"**Enunciado.** (s=Φ(c²)) herda discretude de Word: (c²ınℚ∩[0,1]⇒ sınℚ). A fração em (C₃) (Cantor ternário, profundidade 8) é (≈ 10%) — consistente com medida zero (μ(C₃)=0).","epistemico":"fato","exemplos":[],"id":"proposicao_3_cantor_ternario_em_s_directo","memoria":"semantica","meta":{"dominio":"manual","nivel":6,"verificacao":"slug idêntico — enriquecer, não duplicar"},"origem":"KOCH_CANTOR_FIB_PROVA.md","palavras_chave":[],"sinonimos":[],"tipo":"conceito","titulo":"Proposição 3 (Cantor ternário em s) — **directo**"}
\section{Proposição 3 (Cantor ternário em s) — **directo**}
**Enunciado.** (s=\ensuremath{\Phi}(c\ensuremath{^{2}})) herda discretude de Word: (c\ensuremath{^{2}}\ensuremath{\imath}n\ensuremath{\mathbb{Q}}\ensuremath{\cap}[0,1]\ensuremath{\Rightarrow} s\ensuremath{\imath}n\ensuremath{\mathbb{Q}}). A fração em (C\ensuremath{_{3}}) (Cantor ternário, profundidade 8) é (\ensuremath{\approx} 10\%) — consistente com medida zero (\ensuremath{\mu}(C\ensuremath{_{3}})=0).
\prov{relações: parte\_de koch\_codifica\_cantorfibonacci\_prova\_e\_evidencia}
\prov{proveniência: KOCH\_CANTOR\_FIB\_PROVA.md \textperiodcentered{} fato \textperiodcentered{} confiança 1.0 \textperiodcentered{} domínio manual \textperiodcentered{} história 5 versões no jornal}

%CRISTAL {"arestas":[{"alvo":"koch_como_recipiente_memoria_nao-associativa","confianca":1.0,"epistemico":"fato","origem":"KOCH_MEMORIA.md","rel":"parte_de"},{"alvo":"word","confianca":0.5,"epistemico":"fato","origem":"wiki","rel":"relaciona"}],"confianca":1.0,"contraexemplos":[],"descricao":"**Enunciado.** Seja (R(w)=(endereco,atrator,zₖoch)) o recipiente Koch. Então","epistemico":"fato","exemplos":[],"id":"proposicao_recipiente","memoria":"semantica","meta":{"dominio":"manual","nivel":6,"verificacao":"slug idêntico — enriquecer, não duplicar"},"origem":"KOCH_MEMORIA.md","palavras_chave":[],"sinonimos":[],"tipo":"conceito","titulo":"Proposição (recipiente)"}
\section{Proposição (recipiente)}
**Enunciado.** Seja (R(w)=(endereco,atrator,z\ensuremath{_{k}}och)) o recipiente Koch. Então
\prov{relações: parte\_de koch\_como\_recipiente\_memoria\_nao-associativa; relaciona word}
\prov{proveniência: KOCH\_MEMORIA.md \textperiodcentered{} fato \textperiodcentered{} confiança 1.0 \textperiodcentered{} domínio manual \textperiodcentered{} história 8 versões no jornal}

%CRISTAL {"arestas":[{"alvo":"gato","confianca":0.95,"epistemico":"fato","origem":"paper","rel":"parte_de"}],"confianca":1.0,"contraexemplos":[],"descricao":"| Parâmetro | Valor | |-----------|-------| | Candidatos | 2,000,000 | | Alvo SHA (bits zero MSB) | 16 | | Fração alvo aceitação bit | 10.00% | | Limiar calor calibrado | 0.000000175 | | Ref Word | `[295, 143]` | | Sucessos SHA | 31 | | Cifra emitida | True |","epistemico":"fato","exemplos":[],"id":"proposition_agulha_de_dirac_o_amostradorpropdirac_seja_","memoria":"semantica","meta":{"dominio":"manual","duplicata_sugerida":[],"nivel":6,"verificacao":"parte de conceito mais amplo 'proposition_agulha_de_dirac_o_amostradorpropdirac_seja_'"},"origem":"POW_BIT_TAXA.md","palavras_chave":[],"sinonimos":[],"tipo":"conceito","titulo":"Agulha de Dirac: o amostrador"}
\section{Agulha de Dirac: o amostrador}
| Parâmetro | Valor | |-----------|-------| | Candidatos | 2,000,000 | | Alvo SHA (bits zero MSB) | 16 | | Fração alvo aceitação bit | 10.00\% | | Limiar calor calibrado | 0.000000175 | | Ref Word | `[295, 143]` | | Sucessos SHA | 31 | | Cifra emitida | True |
\prov{relações: parte\_de gato}
\prov{proveniência: POW\_BIT\_TAXA.md \textperiodcentered{} fato \textperiodcentered{} confiança 1.0 \textperiodcentered{} domínio manual \textperiodcentered{} história 5 versões no jornal}

%CRISTAL {"arestas":[{"alvo":"htlc","confianca":1.0,"epistemico":"fato","origem":"wiki","rel":"relaciona"},{"alvo":"confianca_minimizada","confianca":1.0,"epistemico":"fato","origem":"wiki","rel":"usa"}],"confianca":1.0,"contraexemplos":[],"descricao":"O complemento atemporal do atomic swap: em vez de atomicidade no tempo, registra-se o certificado imutável de cada evento e resolve-se de forma assíncrona sobre fundos pré-bloqueados. Não elimina liveness — defere (troca liveness temporal por eventual). Pipeline evento→cert→mux→demux→reordena→escrow.","epistemico":"fato","exemplos":[],"id":"protocolo_assincrono","memoria":"semantica","meta":{"autor":"Aarão/broca-so","dominio":"broca_repos","fonte":""},"origem":"wiki","palavras_chave":[],"sinonimos":[],"tipo":"conceito","titulo":"Protocolo Assíncrono (evidência + ordenação)"}
\section{Protocolo Assíncrono (evidência + ordenação)}
O complemento atemporal do atomic swap: em vez de atomicidade no tempo, registra-se o certificado imutável de cada evento e resolve-se de forma assíncrona sobre fundos pré-bloqueados. Não elimina liveness — defere (troca liveness temporal por eventual). Pipeline evento\ensuremath{\to}cert\ensuremath{\to}mux\ensuremath{\to}demux\ensuremath{\to}reordena\ensuremath{\to}escrow.
\prov{relações: relaciona htlc; usa confianca\_minimizada}
\prov{proveniência: wiki \textperiodcentered{} fato \textperiodcentered{} confiança 1.0 \textperiodcentered{} domínio broca\_repos \textperiodcentered{} história 3 versões no jornal}

%CRISTAL {"arestas":[{"alvo":"hypermlp_escolhe_algebra","confianca":1.0,"epistemico":"fato","origem":"wiki","rel":"usa"},{"alvo":"honestidade_intelectual","confianca":1.0,"epistemico":"fato","origem":"wiki","rel":"relaciona"}],"confianca":1.0,"contraexemplos":[],"descricao":"Continual learning (prevê, espera h dias, compara, atualiza com replay) + warmup por embedding de ticker é o único protocolo sem vazamento que gera R² out-of-sample positivo; produz Sharpe OOS 1,51–1,78 em IXIC (2017-2024) vs 0,81 do buy-and-hold. A honestidade metodológica (anti-leak) é o resultado.","epistemico":"fato","exemplos":[],"id":"protocolo_continual_anti_leak","memoria":"semantica","meta":{"autor":"Aarão/broca-so","dominio":"broca_repos","fonte":""},"origem":"wiki","palavras_chave":[],"sinonimos":[],"tipo":"conceito","titulo":"O Protocolo Honesto predict→wait→compare→train"}
\section{O Protocolo Honesto predict\ensuremath{\to}wait\ensuremath{\to}compare\ensuremath{\to}train}
Continual learning (prevê, espera h dias, compara, atualiza com replay) + warmup por embedding de ticker é o único protocolo sem vazamento que gera R\ensuremath{^{2}} out-of-sample positivo; produz Sharpe OOS 1,51–1,78 em IXIC (2017-2024) vs 0,81 do buy-and-hold. A honestidade metodológica (anti-leak) é o resultado.
\prov{relações: usa hypermlp\_escolhe\_algebra; relaciona honestidade\_intelectual}
\prov{proveniência: wiki \textperiodcentered{} fato \textperiodcentered{} confiança 1.0 \textperiodcentered{} domínio broca\_repos \textperiodcentered{} história 4 versões no jornal}

%CRISTAL {"arestas":[{"alvo":"ula_da_broca_arquitetura_minima","confianca":1.0,"epistemico":"fato","origem":"README.md","rel":"parte_de"}],"confianca":1.0,"contraexemplos":[],"descricao":"| # | Programa | Métricas (exemplo) | |---|----------|-------------------| | 1 | Soma | 5 instr, 2 LOAD, 1 STORE | | 2 | Igualdade | 12 instr, 1 salto | | 3 | Loop + sentinela | 19 instr, CMP/JZ | | 4 | Acumulador dir. | 14 instr, scratch 8000 | | 5 | Pipeline | 4 instr | | 6 | Tesouraria | 16 instr, contrato real | | 7 | Sistema completo | 27 instr, sem Python no runtime |","epistemico":"fato","exemplos":[],"id":"prova_da_isa_make_isa-prova","memoria":"semantica","meta":{"dominio":"manual","nivel":6,"verificacao":"slug idêntico — enriquecer, não duplicar"},"origem":"README.md","palavras_chave":[],"sinonimos":[],"tipo":"conceito","titulo":"Prova da ISA (`make isa-prova`)"}
\section{Prova da ISA (`make isa-prova`)}
| \# | Programa | Métricas (exemplo) | |---|----------|-------------------| | 1 | Soma | 5 instr, 2 LOAD, 1 STORE | | 2 | Igualdade | 12 instr, 1 salto | | 3 | Loop + sentinela | 19 instr, CMP/JZ | | 4 | Acumulador dir. | 14 instr, scratch 8000 | | 5 | Pipeline | 4 instr | | 6 | Tesouraria | 16 instr, contrato real | | 7 | Sistema completo | 27 instr, sem Python no runtime |
\prov{relações: parte\_de ula\_da\_broca\_arquitetura\_minima}
\prov{proveniência: README.md \textperiodcentered{} fato \textperiodcentered{} confiança 1.0 \textperiodcentered{} domínio manual \textperiodcentered{} história 4 versões no jornal}

%CRISTAL {"arestas":[{"alvo":"cifra_prova","confianca":1.0,"epistemico":"fato","origem":"manual","rel":"parte_de"},{"alvo":"utilitarismo","confianca":0.5,"epistemico":"fato","origem":"wiki","rel":"relaciona"},{"alvo":"ruido","confianca":0.5,"epistemico":"fato","origem":"wiki","rel":"relaciona"},{"alvo":"vida-morte","confianca":0.5,"epistemico":"fato","origem":"wiki","rel":"relaciona"},{"alvo":"tipos_de_interpretante","confianca":0.5,"epistemico":"fato","origem":"wiki","rel":"relaciona"},{"alvo":"transcricao","confianca":0.5,"epistemico":"fato","origem":"wiki","rel":"relaciona"},{"alvo":"resultado_anterior_identidade_nao_cobertura","confianca":0.5,"epistemico":"fato","origem":"wiki","rel":"relaciona"},{"alvo":"virtualizacao","confianca":0.5,"epistemico":"fato","origem":"wiki","rel":"relaciona"},{"alvo":"veredito_experimental","confianca":0.5,"epistemico":"fato","origem":"wiki","rel":"relaciona"},{"alvo":"scheduler","confianca":0.5,"epistemico":"fato","origem":"wiki","rel":"relaciona"}],"confianca":1.0,"contraexemplos":[],"descricao":"> ADD combina. SILVER **evolui**. São físicas diferentes.","epistemico":"fato","exemplos":[],"id":"prova_da_maquina_espectral","memoria":"semantica","meta":{"dominio":"manual","nivel":6,"verificacao":"slug idêntico — enriquecer, não duplicar"},"origem":"CIFRA_PROVA.md","palavras_chave":[],"sinonimos":[],"tipo":"conceito","titulo":"Prova da Máquina Espectral"}
\section{Prova da Máquina Espectral}
> ADD combina. SILVER **evolui**. São físicas diferentes.
\prov{relações: parte\_de cifra\_prova; relaciona utilitarismo; relaciona ruido; relaciona vida-morte; relaciona tipos\_de\_interpretante; relaciona transcricao; relaciona resultado\_anterior\_identidade\_nao\_cobertura; relaciona virtualizacao; relaciona veredito\_experimental; relaciona scheduler}
\prov{proveniência: CIFRA\_PROVA.md \textperiodcentered{} fato \textperiodcentered{} confiança 1.0 \textperiodcentered{} domínio manual \textperiodcentered{} história 13 versões no jornal}

%CRISTAL {"arestas":[{"alvo":"broca_os","confianca":1.0,"epistemico":"fato","origem":"manual","rel":"prova"},{"alvo":"transcricao","confianca":0.5,"epistemico":"fato","origem":"wiki","rel":"relaciona"},{"alvo":"teste_ordem","confianca":0.5,"epistemico":"fato","origem":"wiki","rel":"relaciona"},{"alvo":"python","confianca":0.5,"epistemico":"fato","origem":"wiki","rel":"relaciona"},{"alvo":"readme","confianca":0.5,"epistemico":"fato","origem":"wiki","rel":"relaciona"},{"alvo":"recursao","confianca":0.5,"epistemico":"fato","origem":"wiki","rel":"relaciona"},{"alvo":"sessao_interativa","confianca":0.5,"epistemico":"fato","origem":"wiki","rel":"relaciona"},{"alvo":"word","confianca":0.5,"epistemico":"fato","origem":"wiki","rel":"relaciona"},{"alvo":"syscall","confianca":0.5,"epistemico":"fato","origem":"wiki","rel":"relaciona"}],"confianca":1.0,"contraexemplos":[],"descricao":"Manual: ula/PROVA_ISA.md","epistemico":"fato","exemplos":[],"id":"prova_isa","memoria":"semantica","meta":{"arquivo":"ula/PROVA_ISA.md","dominio":"manual","nivel":6},"origem":"manual","palavras_chave":[],"sinonimos":[],"tipo":"conceito","titulo":"PROVA_ISA"}
\section{PROVA\_ISA}
Manual: ula/PROVA\_ISA.md
\prov{relações: prova broca\_os; relaciona transcricao; relaciona teste\_ordem; relaciona python; relaciona readme; relaciona recursao; relaciona sessao\_interativa; relaciona word; relaciona syscall}
\prov{proveniência: manual \textperiodcentered{} fato \textperiodcentered{} confiança 1.0 \textperiodcentered{} domínio manual \textperiodcentered{} história 12 versões no jornal}

%CRISTAL {"arestas":[{"alvo":"grava_fato_resposta_normal_em_seguida","confianca":1.0,"epistemico":"fato","origem":"CONHECIMENTO.md","rel":"parte_de"},{"alvo":"broca-so_camada_de_trabalho_isomorfica","confianca":1.0,"epistemico":"fato","origem":"BROCA_SO.md","rel":"parte_de"}],"confianca":1.0,"contraexemplos":[],"descricao":"1. Port `TORO…POW_SEED` para C no backend CPU (ou prova cruzada Python↔C). 2. Backend GPU PTX — mesmas ops, batch gato 2×2. 3. Refactor stratum: borda só I/O + `broca_so_exec(POW_SEED)`. 4. Android NDK — terceiro backend, mesma prova.","epistemico":"fato","exemplos":[],"id":"proximos_passos","memoria":"semantica","meta":{"dominio":"manual","nivel":6,"verificacao":"slug idêntico — enriquecer, não duplicar"},"origem":"BROCA_SO.md","palavras_chave":[],"sinonimos":[],"tipo":"conceito","titulo":"Próximos passos"}
\section{Próximos passos}
1. Port `TORO…POW\_SEED` para C no backend CPU (ou prova cruzada Python\ensuremath{\leftrightarrow}C). 2. Backend GPU PTX — mesmas ops, batch gato 2$\times$2. 3. Refactor stratum: borda só I/O + `broca\_so\_exec(POW\_SEED)`. 4. Android NDK — terceiro backend, mesma prova.
\prov{relações: parte\_de grava\_fato\_resposta\_normal\_em\_seguida; parte\_de broca-so\_camada\_de\_trabalho\_isomorfica}
\prov{proveniência: BROCA\_SO.md \textperiodcentered{} fato \textperiodcentered{} confiança 1.0 \textperiodcentered{} domínio manual \textperiodcentered{} história 8 versões no jornal}

%CRISTAL {"arestas":[{"alvo":"broca_os","confianca":0.95,"epistemico":"fato","origem":"paper","rel":"parte_de"},{"alvo":"virtualizacao","confianca":0.5,"epistemico":"fato","origem":"wiki","rel":"relaciona"},{"alvo":"word","confianca":0.5,"epistemico":"fato","origem":"wiki","rel":"relaciona"},{"alvo":"while","confianca":0.5,"epistemico":"fato","origem":"wiki","rel":"relaciona"}],"confianca":1.0,"contraexemplos":[],"descricao":"```bash cristal qa                    # test + bench + validate cristal qa regression cristal release full          # test_cristal + fase3 + … ```","epistemico":"fato","exemplos":[],"id":"qa_e_release","memoria":"semantica","meta":{"dominio":"manual","nivel":6,"verificacao":"slug idêntico — enriquecer, não duplicar"},"origem":"CRISTAL_DEV.md","palavras_chave":[],"sinonimos":[],"tipo":"conceito","titulo":"QA e Release"}
\section{QA e Release}
```bash cristal qa                    \# test + bench + validate cristal qa regression cristal release full          \# test\_cristal + fase3 + … ```
\prov{relações: parte\_de broca\_os; relaciona virtualizacao; relaciona word; relaciona while}
\prov{proveniência: CRISTAL\_DEV.md \textperiodcentered{} fato \textperiodcentered{} confiança 1.0 \textperiodcentered{} domínio manual \textperiodcentered{} história 6 versões no jornal}

%CRISTAL {"arestas":[{"alvo":"gato","confianca":0.95,"epistemico":"fato","origem":"paper","rel":"parte_de"},{"alvo":"virtualizacao","confianca":0.5,"epistemico":"fato","origem":"wiki","rel":"relaciona"},{"alvo":"word","confianca":0.5,"epistemico":"fato","origem":"wiki","rel":"relaciona"},{"alvo":"vontade_de_poder","confianca":0.5,"epistemico":"fato","origem":"wiki","rel":"relaciona"},{"alvo":"while","confianca":0.5,"epistemico":"fato","origem":"wiki","rel":"relaciona"}],"confianca":1.0,"contraexemplos":[],"descricao":"- nonce **1572863** → rank **#2** (acopl=0.9986) - nonce **1966079** → rank **#1** (acopl=0.9991)","epistemico":"fato","exemplos":[],"id":"ranking_global_por_acoplamento","memoria":"semantica","meta":{"dominio":"manual","nivel":6,"verificacao":"slug idêntico — enriquecer, não duplicar"},"origem":"POW_METAL_PRE_FECHO.md","palavras_chave":[],"sinonimos":[],"tipo":"conceito","titulo":"Ranking global por acoplamento"}
\section{Ranking global por acoplamento}
- nonce **1572863** \ensuremath{\to} rank **\#2** (acopl=0.9986) - nonce **1966079** \ensuremath{\to} rank **\#1** (acopl=0.9991)
\prov{relações: parte\_de gato; relaciona virtualizacao; relaciona word; relaciona vontade\_de\_poder; relaciona while}
\prov{proveniência: POW\_METAL\_PRE\_FECHO.md \textperiodcentered{} fato \textperiodcentered{} confiança 1.0 \textperiodcentered{} domínio manual \textperiodcentered{} história 7 versões no jornal}

%CRISTAL {"arestas":[{"alvo":"ciencias","confianca":1.0,"epistemico":"fato","origem":"manual","rel":"referencia"},{"alvo":"transcricao","confianca":0.5,"epistemico":"fato","origem":"wiki","rel":"relaciona"},{"alvo":"sessao_interativa","confianca":0.5,"epistemico":"fato","origem":"wiki","rel":"relaciona"},{"alvo":"utilitarismo","confianca":0.5,"epistemico":"fato","origem":"wiki","rel":"relaciona"},{"alvo":"vida-morte","confianca":0.5,"epistemico":"fato","origem":"wiki","rel":"relaciona"},{"alvo":"word","confianca":0.5,"epistemico":"fato","origem":"wiki","rel":"relaciona"},{"alvo":"virtualizacao","confianca":0.5,"epistemico":"fato","origem":"wiki","rel":"relaciona"},{"alvo":"telemetria_shadow_producao","confianca":0.5,"epistemico":"fato","origem":"wiki","rel":"relaciona"},{"alvo":"vetores","confianca":0.5,"epistemico":"fato","origem":"wiki","rel":"relaciona"},{"alvo":"sequencia_pow_nativa","confianca":0.5,"epistemico":"fato","origem":"wiki","rel":"relaciona"},{"alvo":"semiotica_peirce_triade","confianca":0.5,"epistemico":"fato","origem":"wiki","rel":"relaciona"}],"confianca":1.0,"contraexemplos":[],"descricao":"Manual: ula/README.md","epistemico":"fato","exemplos":[],"id":"readme","memoria":"semantica","meta":{"arquivo":"ula/README.md","dominio":"manual","nivel":6},"origem":"manual","palavras_chave":[],"sinonimos":[],"tipo":"conceito","titulo":"README"}
\section{README}
Manual: ula/README.md
\prov{relações: referencia ciencias; relaciona transcricao; relaciona sessao\_interativa; relaciona utilitarismo; relaciona vida-morte; relaciona word; relaciona virtualizacao; relaciona telemetria\_shadow\_producao; relaciona vetores; relaciona sequencia\_pow\_nativa; relaciona semiotica\_peirce\_triade}
\prov{proveniência: manual \textperiodcentered{} fato \textperiodcentered{} confiança 1.0 \textperiodcentered{} domínio manual \textperiodcentered{} história 14 versões no jornal}

%CRISTAL {"arestas":[{"alvo":"readme_cristal","confianca":1.0,"epistemico":"fato","origem":"manual","rel":"parte_de"},{"alvo":"tiffany","confianca":1.0,"epistemico":"fato","origem":"manual","rel":"explica"},{"alvo":"utilitarismo","confianca":0.5,"epistemico":"fato","origem":"wiki","rel":"relaciona"},{"alvo":"teoria_do_valor","confianca":0.5,"epistemico":"fato","origem":"wiki","rel":"relaciona"},{"alvo":"ruido","confianca":0.5,"epistemico":"fato","origem":"wiki","rel":"relaciona"},{"alvo":"word","confianca":0.5,"epistemico":"fato","origem":"wiki","rel":"relaciona"},{"alvo":"virtualizacao","confianca":0.5,"epistemico":"fato","origem":"wiki","rel":"relaciona"}],"confianca":1.0,"contraexemplos":[],"descricao":"**Cristal** é a ferramenta de infraestrutura da plataforma Broca/Tiffany.","epistemico":"fato","exemplos":[],"id":"readme_cristal","memoria":"semantica","meta":{"dominio":"manual","nivel":6,"verificacao":"slug idêntico — enriquecer, não duplicar"},"origem":"README_CRISTAL.md","palavras_chave":[],"sinonimos":[],"tipo":"conceito","titulo":"README_CRISTAL"}
\section{README\_CRISTAL}
**Cristal** é a ferramenta de infraestrutura da plataforma Broca/Tiffany.
\prov{relações: parte\_de readme\_cristal; explica tiffany; relaciona utilitarismo; relaciona teoria\_do\_valor; relaciona ruido; relaciona word; relaciona virtualizacao}
\prov{proveniência: README\_CRISTAL.md \textperiodcentered{} fato \textperiodcentered{} confiança 1.0 \textperiodcentered{} domínio manual \textperiodcentered{} história 40 versões no jornal}

%CRISTAL {"arestas":[{"alvo":"peirce_celulas_canais","confianca":1.0,"epistemico":"fato","origem":"wiki","rel":"usa"},{"alvo":"broca_os","confianca":1.0,"epistemico":"fato","origem":"wiki","rel":"relaciona"}],"confianca":1.0,"contraexemplos":[],"descricao":"A rede vive na parte negra: nós = células (idempotentes); links = canais (radical, nilpotente); roteamento store-and-forward = multiplicação (E_ij·E_jk=E_ik). Invariante: o roteamento termina (sem deadlock) ⟺ rad(kQ) é nilpotente ⟺ o quiver Q é acíclico. Fecha por blocos (a recursão broca-de-brocas) só se o quociente permanece acíclico.","epistemico":"fato","exemplos":[],"id":"rede_algebra_de_caminhos","memoria":"semantica","meta":{"autor":"Aarão/broca-so","dominio":"broca_repos","fonte":""},"origem":"wiki","palavras_chave":[],"sinonimos":[],"tipo":"conceito","titulo":"Rede: Roteamento Termina ⟺ Radical Nilpotente"}
\section{Rede: Roteamento Termina \ensuremath{\Longleftrightarrow} Radical Nilpotente}
A rede vive na parte negra: nós = células (idempotentes); links = canais (radical, nilpotente); roteamento store-and-forward = multiplicação (E\_ij\textperiodcentered E\_jk=E\_ik). Invariante: o roteamento termina (sem deadlock) \ensuremath{\Longleftrightarrow} rad(kQ) é nilpotente \ensuremath{\Longleftrightarrow} o quiver Q é acíclico. Fecha por blocos (a recursão broca-de-brocas) só se o quociente permanece acíclico.
\prov{relações: usa peirce\_celulas\_canais; relaciona broca\_os}
\prov{proveniência: wiki \textperiodcentered{} fato \textperiodcentered{} confiança 1.0 \textperiodcentered{} domínio broca\_repos \textperiodcentered{} história 3 versões no jornal}

%CRISTAL {"arestas":[{"alvo":"peirce_celulas_canais","confianca":1.0,"epistemico":"fato","origem":"wiki","rel":"usa"},{"alvo":"cifra","confianca":1.0,"epistemico":"fato","origem":"wiki","rel":"usa"},{"alvo":"a_bolha_infraestrutura","confianca":1.0,"epistemico":"fato","origem":"wiki","rel":"relaciona"}],"confianca":1.0,"contraexemplos":[],"descricao":"O canal limpo é o bump = msg ⊕ keystream(banda), banda=sha256(tecido), cru no datagrama UDP (sem handshake). TCP/SSH é 'sujo': joga uma segunda assinatura por cima, destruindo a função de onda. O canal é sempre P2P; o central só faz descoberta (rendezvous/hole-punching) e pode ser cofre (backup cifrado), nunca correio. Banda por ECDH X25519.","epistemico":"fato","exemplos":[],"id":"rede_limpa_bump","memoria":"semantica","meta":{"autor":"Aarão/broca-so","dominio":"broca_repos","fonte":""},"origem":"wiki","palavras_chave":[],"sinonimos":[],"tipo":"conceito","titulo":"A Rede Limpa (o bump P2P)"}
\section{A Rede Limpa (o bump P2P)}
O canal limpo é o bump = msg \ensuremath{\oplus} keystream(banda), banda=sha256(tecido), cru no datagrama UDP (sem handshake). TCP/SSH é 'sujo': joga uma segunda assinatura por cima, destruindo a função de onda. O canal é sempre P2P; o central só faz descoberta (rendezvous/hole-punching) e pode ser cofre (backup cifrado), nunca correio. Banda por ECDH X25519.
\prov{relações: usa peirce\_celulas\_canais; usa cifra; relaciona a\_bolha\_infraestrutura}
\prov{proveniência: wiki \textperiodcentered{} fato \textperiodcentered{} confiança 1.0 \textperiodcentered{} domínio broca\_repos \textperiodcentered{} história 4 versões no jornal}

%CRISTAL {"arestas":[{"alvo":"mult_prtsu","confianca":1.0,"epistemico":"fato","origem":"wiki","rel":"usa"},{"alvo":"hopfield","confianca":1.0,"epistemico":"fato","origem":"wiki","rel":"relaciona"},{"alvo":"maquina_algebrica","confianca":1.0,"epistemico":"fato","origem":"wiki","rel":"relaciona"}],"confianca":1.0,"contraexemplos":[],"descricao":"Arquitetura cuja operação de propagação é a multiplicação de 5 parâmetros (p,r,t,s,u) ajustados por gradiente junto aos pesos — a rede 'escolhe sua própria álgebra' (global, por camada ou por neurônio). Variante aprendível da álgebra dos contratos.","epistemico":"fato","exemplos":[],"id":"rede_neural_algebra_aprendivel","memoria":"semantica","meta":{"autor":"Lopes/Aarão","dominio":"hiper","fonte":""},"origem":"wiki","palavras_chave":[],"sinonimos":[],"tipo":"conceito","titulo":"Rede Neural com Álgebra Aprendível"}
\section{Rede Neural com Álgebra Aprendível}
Arquitetura cuja operação de propagação é a multiplicação de 5 parâmetros (p,r,t,s,u) ajustados por gradiente junto aos pesos — a rede 'escolhe sua própria álgebra' (global, por camada ou por neurônio). Variante aprendível da álgebra dos contratos.
\prov{relações: usa mult\_prtsu; relaciona hopfield; relaciona maquina\_algebrica}
\prov{proveniência: wiki \textperiodcentered{} fato \textperiodcentered{} confiança 1.0 \textperiodcentered{} domínio hiper \textperiodcentered{} história 5 versões no jornal}

%CRISTAL {"arestas":[{"alvo":"respiracao_dos_programas","confianca":1.0,"epistemico":"fato","origem":"PARABOLA_ISA.md","rel":"parte_de"},{"alvo":"word","confianca":0.5,"epistemico":"fato","origem":"wiki","rel":"relaciona"},{"alvo":"virtualizacao","confianca":0.5,"epistemico":"fato","origem":"wiki","rel":"relaciona"}],"confianca":1.0,"contraexemplos":[],"descricao":"- calor acumulado: **0.651** - fase acumulada: **15.629** - perda: **0.000** - potência (calor/passo): **0.041** - passos: **16**","epistemico":"fato","exemplos":[],"id":"reduce","memoria":"semantica","meta":{"dominio":"manual","nivel":6,"verificacao":"slug idêntico — enriquecer, não duplicar"},"origem":"PARABOLA_ISA.md","palavras_chave":[],"sinonimos":[],"tipo":"conceito","titulo":"🌱 reduce"}
\section{[semente] reduce}
- calor acumulado: **0.651** - fase acumulada: **15.629** - perda: **0.000** - potência (calor/passo): **0.041** - passos: **16**
\prov{relações: parte\_de respiracao\_dos\_programas; relaciona word; relaciona virtualizacao}
\prov{proveniência: PARABOLA\_ISA.md \textperiodcentered{} fato \textperiodcentered{} confiança 1.0 \textperiodcentered{} domínio manual \textperiodcentered{} história 8 versões no jornal}

%CRISTAL {"arestas":[{"alvo":"main","confianca":1.0,"epistemico":"fato","origem":"paper","rel":"parte_de"},{"alvo":"broca-so_cristalizado","confianca":1.0,"epistemico":"fato","origem":"BROCA_CRISTAL.md","rel":"parte_de"},{"alvo":"koch_como_recipiente_memoria_nao-associativa","confianca":1.0,"epistemico":"fato","origem":"KOCH_MEMORIA.md","rel":"parte_de"},{"alvo":"koch_codifica_cantorfibonacci_prova_e_evidencia","confianca":1.0,"epistemico":"fato","origem":"KOCH_CANTOR_FIB_PROVA.md","rel":"parte_de"}],"confianca":1.0,"contraexemplos":[],"descricao":"`papers/paper₂retaouro.tex` — Teo. Cantor, (D₂) vs (D_varphi) - `papers/olhodeₖoch.tex` — endereço fractal, canal Peirce - `papers/transformadadourada.tex` — AGM, atrator - `ula/kochatlas.c`, `ula/kochparabola.c`","epistemico":"fato","exemplos":[],"id":"referencias","memoria":"semantica","meta":{"dominio":"manual","nivel":6,"verificacao":"slug idêntico — enriquecer, não duplicar"},"origem":"KOCH_CANTOR_FIB_PROVA.md","palavras_chave":[],"sinonimos":[],"tipo":"conceito","titulo":"Referências"}
\section{Referências}
`papers/paper\ensuremath{_{2}}retaouro.tex` — Teo. Cantor, (D\ensuremath{_{2}}) vs (D\_varphi) - `papers/olhode\ensuremath{_{k}}och.tex` — endereço fractal, canal Peirce - `papers/transformadadourada.tex` — AGM, atrator - `ula/kochatlas.c`, `ula/kochparabola.c`
\prov{relações: parte\_de main; parte\_de broca-so\_cristalizado; parte\_de koch\_como\_recipiente\_memoria\_nao-associativa; parte\_de koch\_codifica\_cantorfibonacci\_prova\_e\_evidencia}
\prov{proveniência: KOCH\_CANTOR\_FIB\_PROVA.md \textperiodcentered{} fato \textperiodcentered{} confiança 1.0 \textperiodcentered{} domínio manual \textperiodcentered{} história 21 versões no jornal}

%CRISTAL {"arestas":[{"alvo":"2_semantica_formal","confianca":1.0,"epistemico":"fato","origem":"PROVA_ISA.md","rel":"parte_de"},{"alvo":"word","confianca":0.5,"epistemico":"fato","origem":"wiki","rel":"relaciona"},{"alvo":"vontade_de_poder","confianca":0.5,"epistemico":"fato","origem":"wiki","rel":"relaciona"}],"confianca":1.0,"contraexemplos":[],"descricao":"| Reg | Função | |-----|--------| | A | operando / entrada LOAD | | B | operando | | R | resultado | | PC | program counter | | FLAGS | ZERO, EQ, LT |","epistemico":"fato","exemplos":[],"id":"registradores","memoria":"semantica","meta":{"dominio":"manual","duplicata_sugerida":[],"nivel":6,"verificacao":"parte de conceito mais amplo 'registradores'"},"origem":"README.md","palavras_chave":[],"sinonimos":[],"tipo":"conceito","titulo":"Registradores"}
\section{Registradores}
| Reg | Função | |-----|--------| | A | operando / entrada LOAD | | B | operando | | R | resultado | | PC | program counter | | FLAGS | ZERO, EQ, LT |
\prov{relações: parte\_de 2\_semantica\_formal; relaciona word; relaciona vontade\_de\_poder}
\prov{proveniência: README.md \textperiodcentered{} fato \textperiodcentered{} confiança 1.0 \textperiodcentered{} domínio manual \textperiodcentered{} história 8 versões no jornal}

%CRISTAL {"arestas":[{"alvo":"broca_os","confianca":0.95,"epistemico":"fato","origem":"paper","rel":"parte_de"},{"alvo":"word","confianca":0.5,"epistemico":"fato","origem":"wiki","rel":"relaciona"}],"confianca":1.0,"contraexemplos":[],"descricao":"- **libtiffany:** não alterar. - **APIs existentes:** tiffany_platform, prosa, orquestrador, tiffany pip. - **Usuário:** nunca edita `.tsv`/`.idx`/manifest manualmente.","epistemico":"fato","exemplos":[],"id":"regra_de_ouro","memoria":"semantica","meta":{"dominio":"manual","nivel":6,"verificacao":"slug idêntico — enriquecer, não duplicar"},"origem":"CRISTAL_DEV.md","palavras_chave":[],"sinonimos":[],"tipo":"conceito","titulo":"Regra de ouro"}
\section{Regra de ouro}
- **libtiffany:** não alterar. - **APIs existentes:** tiffany\_platform, prosa, orquestrador, tiffany pip. - **Usuário:** nunca edita `.tsv`/`.idx`/manifest manualmente.
\prov{relações: parte\_de broca\_os; relaciona word}
\prov{proveniência: CRISTAL\_DEV.md \textperiodcentered{} fato \textperiodcentered{} confiança 1.0 \textperiodcentered{} domínio manual \textperiodcentered{} história 4 versões no jornal}

%CRISTAL {"arestas":[{"alvo":"o_conselho","confianca":1.0,"epistemico":"fato","origem":"wiki","rel":"parte_de"},{"alvo":"honestidade_intelectual","confianca":1.0,"epistemico":"fato","origem":"wiki","rel":"usa"},{"alvo":"vida_matematica","confianca":1.0,"epistemico":"fato","origem":"wiki","rel":"relaciona"}],"confianca":1.0,"contraexemplos":[],"descricao":"Cada afirmação é etiquetada como Derivado / Numérico / proJeto (posição-de-design) / Interpretação — separando o que é demonstrado do que é escolha ou leitura. A codificação explícita de 'posição≠fato' que rege todo o grafo.","epistemico":"fato","exemplos":[],"id":"regua_dnci","memoria":"semantica","meta":{"autor":"Aarão/broca-so","dominio":"broca_repos","fonte":""},"origem":"wiki","palavras_chave":[],"sinonimos":[],"tipo":"conceito","titulo":"A Régua Epistêmica [D]/[N]/[C]/[I]"}
\section{A Régua Epistêmica [D]/[N]/[C]/[I]}
Cada afirmação é etiquetada como Derivado / Numérico / proJeto (posição-de-design) / Interpretação — separando o que é demonstrado do que é escolha ou leitura. A codificação explícita de 'posição\ensuremath{\ne}fato' que rege todo o grafo.
\prov{relações: parte\_de o\_conselho; usa honestidade\_intelectual; relaciona vida\_matematica}
\prov{proveniência: wiki \textperiodcentered{} fato \textperiodcentered{} confiança 1.0 \textperiodcentered{} domínio broca\_repos \textperiodcentered{} história 4 versões no jornal}

%CRISTAL {"arestas":[{"alvo":"cifra_ouro_branco_nucleo_espectral","confianca":1.0,"epistemico":"fato","origem":"CIFRA.md","rel":"parte_de"},{"alvo":"utilitarismo","confianca":0.5,"epistemico":"fato","origem":"wiki","rel":"relaciona"},{"alvo":"teoria_do_valor","confianca":0.5,"epistemico":"fato","origem":"wiki","rel":"relaciona"},{"alvo":"vetores","confianca":0.5,"epistemico":"fato","origem":"wiki","rel":"relaciona"},{"alvo":"word","confianca":0.5,"epistemico":"fato","origem":"wiki","rel":"relaciona"},{"alvo":"setor_algebrico_descontinuidades_no_fecho","confianca":0.5,"epistemico":"fato","origem":"wiki","rel":"relaciona"}],"confianca":1.0,"contraexemplos":[],"descricao":"```bash make cifra ./broca_asm programas/08_cifra.asm programas/08_cifra.bro ```","epistemico":"fato","exemplos":[],"id":"reproduzir","memoria":"semantica","meta":{"dominio":"manual","nivel":6,"verificacao":"slug idêntico — enriquecer, não duplicar"},"origem":"CIFRA.md","palavras_chave":[],"sinonimos":[],"tipo":"conceito","titulo":"Reproduzir"}
\section{Reproduzir}
```bash make cifra ./broca\_asm programas/08\_cifra.asm programas/08\_cifra.bro ```
\prov{relações: parte\_de cifra\_ouro\_branco\_nucleo\_espectral; relaciona utilitarismo; relaciona teoria\_do\_valor; relaciona vetores; relaciona word; relaciona setor\_algebrico\_descontinuidades\_no\_fecho}
\prov{proveniência: CIFRA.md \textperiodcentered{} fato \textperiodcentered{} confiança 1.0 \textperiodcentered{} domínio manual \textperiodcentered{} história 11 versões no jornal}

%CRISTAL {"arestas":[{"alvo":"a_cacada_da_parabola_mapa_vivo_da_isa","confianca":1.0,"epistemico":"fato","origem":"PARABOLA_ISA.md","rel":"parte_de"},{"alvo":"testes","confianca":0.5,"epistemico":"fato","origem":"wiki","rel":"relaciona"},{"alvo":"syscall","confianca":0.5,"epistemico":"fato","origem":"wiki","rel":"relaciona"},{"alvo":"rust","confianca":0.5,"epistemico":"fato","origem":"wiki","rel":"relaciona"},{"alvo":"thread","confianca":0.5,"epistemico":"fato","origem":"wiki","rel":"relaciona"},{"alvo":"utilitarismo","confianca":0.5,"epistemico":"fato","origem":"wiki","rel":"relaciona"},{"alvo":"word","confianca":0.5,"epistemico":"fato","origem":"wiki","rel":"relaciona"}],"confianca":1.0,"contraexemplos":[],"descricao":"Secção de PARABOLA_ISA.md.","epistemico":"fato","exemplos":[],"id":"respiracao_dos_programas","memoria":"semantica","meta":{"dominio":"manual","nivel":6,"verificacao":"slug idêntico — enriquecer, não duplicar"},"origem":"PARABOLA_ISA.md","palavras_chave":[],"sinonimos":[],"tipo":"conceito","titulo":"Respiração dos programas"}
\section{Respiração dos programas}
Secção de PARABOLA\_ISA.md.
\prov{relações: parte\_de a\_cacada\_da\_parabola\_mapa\_vivo\_da\_isa; relaciona testes; relaciona syscall; relaciona rust; relaciona thread; relaciona utilitarismo; relaciona word}
\prov{proveniência: PARABOLA\_ISA.md \textperiodcentered{} fato \textperiodcentered{} confiança 1.0 \textperiodcentered{} domínio manual \textperiodcentered{} história 12 versões no jornal}

%CRISTAL {"arestas":[{"alvo":"metais","confianca":1.0,"epistemico":"fato","origem":"manual","rel":"parte_de"},{"alvo":"transcricao","confianca":0.5,"epistemico":"fato","origem":"wiki","rel":"relaciona"},{"alvo":"vida-morte","confianca":0.5,"epistemico":"fato","origem":"wiki","rel":"relaciona"},{"alvo":"teste_ordem","confianca":0.5,"epistemico":"fato","origem":"wiki","rel":"relaciona"},{"alvo":"testes","confianca":0.5,"epistemico":"fato","origem":"wiki","rel":"relaciona"},{"alvo":"thread","confianca":0.5,"epistemico":"fato","origem":"wiki","rel":"relaciona"},{"alvo":"tiffany_cabeca_broca","confianca":0.5,"epistemico":"fato","origem":"wiki","rel":"relaciona"},{"alvo":"topologia","confianca":0.5,"epistemico":"fato","origem":"wiki","rel":"relaciona"},{"alvo":"virtualizacao","confianca":0.5,"epistemico":"fato","origem":"wiki","rel":"relaciona"},{"alvo":"syscall","confianca":0.5,"epistemico":"fato","origem":"wiki","rel":"relaciona"},{"alvo":"rust","confianca":0.5,"epistemico":"fato","origem":"wiki","rel":"relaciona"}],"confianca":1.0,"contraexemplos":[],"descricao":"Secção de METAIS.md.","epistemico":"fato","exemplos":[],"id":"resposta_a_pergunta_dificil","memoria":"semantica","meta":{"dominio":"manual","nivel":6,"verificacao":"slug idêntico — enriquecer, não duplicar"},"origem":"METAIS.md","palavras_chave":[],"sinonimos":[],"tipo":"conceito","titulo":"Resposta à pergunta difícil"}
\section{Resposta à pergunta difícil}
Secção de METAIS.md.
\prov{relações: parte\_de metais; relaciona transcricao; relaciona vida-morte; relaciona teste\_ordem; relaciona testes; relaciona thread; relaciona tiffany\_cabeca\_broca; relaciona topologia; relaciona virtualizacao; relaciona syscall; relaciona rust}
\prov{proveniência: METAIS.md \textperiodcentered{} fato \textperiodcentered{} confiança 1.0 \textperiodcentered{} domínio manual \textperiodcentered{} história 18 versões no jornal}

%CRISTAL {"arestas":[{"alvo":"1_modelo_formal","confianca":1.0,"epistemico":"fato","origem":"EXPRESSIVIDADE.md","rel":"parte_de"},{"alvo":"broca_os","confianca":1.0,"epistemico":"fato","origem":"manual","rel":"parte_de"},{"alvo":"word","confianca":0.5,"epistemico":"fato","origem":"wiki","rel":"relaciona"}],"confianca":1.0,"contraexemplos":[],"descricao":"``` LOAD n  ⟹  n ∈ ℕ literal no bytecode (3 bytes) STORE n ⟹  n literal JMP/JZ/JNZ o ⟹  o literal relativo (assemble-time) ```","epistemico":"fato","exemplos":[],"id":"restricao_estrutural_critica","memoria":"semantica","meta":{"dominio":"manual","nivel":6,"verificacao":"slug idêntico — enriquecer, não duplicar"},"origem":"EXPRESSIVIDADE.md","palavras_chave":[],"sinonimos":[],"tipo":"conceito","titulo":"Restrição estrutural (crítica)"}
\section{Restrição estrutural (crítica)}
``` LOAD n  \ensuremath{\Longrightarrow}  n \ensuremath{\in} \ensuremath{\mathbb{N}} literal no bytecode (3 bytes) STORE n \ensuremath{\Longrightarrow}  n literal JMP/JZ/JNZ o \ensuremath{\Longrightarrow}  o literal relativo (assemble-time) ```
\prov{relações: parte\_de 1\_modelo\_formal; parte\_de broca\_os; relaciona word}
\prov{proveniência: EXPRESSIVIDADE.md \textperiodcentered{} fato \textperiodcentered{} confiança 1.0 \textperiodcentered{} domínio manual \textperiodcentered{} história 11 versões no jornal}

%CRISTAL {"arestas":[{"alvo":"koch_pell_verificacao","confianca":1.0,"epistemico":"fato","origem":"KOCH_PELL.md","rel":"parte_de"},{"alvo":"utilitarismo","confianca":0.5,"epistemico":"fato","origem":"wiki","rel":"relaciona"},{"alvo":"ruido","confianca":0.5,"epistemico":"fato","origem":"wiki","rel":"relaciona"},{"alvo":"transcricao","confianca":0.5,"epistemico":"fato","origem":"wiki","rel":"relaciona"},{"alvo":"vida-morte","confianca":0.5,"epistemico":"fato","origem":"wiki","rel":"relaciona"},{"alvo":"virtualizacao","confianca":0.5,"epistemico":"fato","origem":"wiki","rel":"relaciona"},{"alvo":"teoria_do_valor","confianca":0.5,"epistemico":"fato","origem":"wiki","rel":"relaciona"},{"alvo":"veredito_experimental","confianca":0.5,"epistemico":"fato","origem":"wiki","rel":"relaciona"},{"alvo":"scheduler","confianca":0.5,"epistemico":"fato","origem":"wiki","rel":"relaciona"},{"alvo":"resultado_completo","confianca":0.5,"epistemico":"fato","origem":"wiki","rel":"relaciona"},{"alvo":"tipos_de_interpretante","confianca":0.5,"epistemico":"fato","origem":"wiki","rel":"relaciona"},{"alvo":"setor_algebrico_descontinuidades_no_fecho","confianca":0.5,"epistemico":"fato","origem":"wiki","rel":"relaciona"}],"confianca":1.0,"contraexemplos":[],"descricao":"Secção de PARABOLA_DESTRUIR.md.","epistemico":"fato","exemplos":[],"id":"resultado_anterior_identidade_nao_cobertura","memoria":"semantica","meta":{"dominio":"manual","duplicata_sugerida":[],"nivel":6,"verificacao":"parte de conceito mais amplo 'resultado_anterior_identidade_nao_cobertura'"},"origem":"PARABOLA_DESTRUIR.md","palavras_chave":[],"sinonimos":[],"tipo":"conceito","titulo":"Resultado anterior (identidade, não cobertura)"}
\section{Resultado anterior (identidade, não cobertura)}
Secção de PARABOLA\_DESTRUIR.md.
\prov{relações: parte\_de koch\_pell\_verificacao; relaciona utilitarismo; relaciona ruido; relaciona transcricao; relaciona vida-morte; relaciona virtualizacao; relaciona teoria\_do\_valor; relaciona veredito\_experimental; relaciona scheduler; relaciona resultado\_completo; relaciona tipos\_de\_interpretante; relaciona setor\_algebrico\_descontinuidades\_no\_fecho}
\prov{proveniência: PARABOLA\_DESTRUIR.md \textperiodcentered{} fato \textperiodcentered{} confiança 1.0 \textperiodcentered{} domínio manual \textperiodcentered{} história 19 versões no jornal}

%CRISTAL {"arestas":[{"alvo":"broca_os","confianca":0.95,"epistemico":"fato","origem":"paper","rel":"parte_de"},{"alvo":"utilitarismo","confianca":0.5,"epistemico":"fato","origem":"wiki","rel":"relaciona"},{"alvo":"transcricao","confianca":0.5,"epistemico":"fato","origem":"wiki","rel":"relaciona"},{"alvo":"teoria_do_valor","confianca":0.5,"epistemico":"fato","origem":"wiki","rel":"relaciona"},{"alvo":"vida-morte","confianca":0.5,"epistemico":"fato","origem":"wiki","rel":"relaciona"},{"alvo":"ruido","confianca":0.5,"epistemico":"fato","origem":"wiki","rel":"relaciona"},{"alvo":"setor_algebrico_descontinuidades_no_fecho","confianca":0.5,"epistemico":"fato","origem":"wiki","rel":"relaciona"},{"alvo":"tipos_de_interpretante","confianca":0.5,"epistemico":"fato","origem":"wiki","rel":"relaciona"},{"alvo":"scheduler","confianca":0.5,"epistemico":"fato","origem":"wiki","rel":"relaciona"},{"alvo":"toryces_e_helyces","confianca":0.5,"epistemico":"fato","origem":"wiki","rel":"relaciona"},{"alvo":"vetores","confianca":0.5,"epistemico":"fato","origem":"wiki","rel":"relaciona"},{"alvo":"virtualizacao","confianca":0.5,"epistemico":"fato","origem":"wiki","rel":"relaciona"}],"confianca":1.0,"contraexemplos":[],"descricao":"hit = tiffany.consulta(\"TAG_GIT_001 zqwm01 clone\") print(hit[\"domain\"], hit[\"text\"])","epistemico":"fato","exemplos":[],"id":"resultado_completo","memoria":"semantica","meta":{"dominio":"manual","nivel":6,"verificacao":"slug idêntico — enriquecer, não duplicar"},"origem":"INSTALL.md","palavras_chave":[],"sinonimos":[],"tipo":"conceito","titulo":"resultado completo"}
\section{resultado completo}
hit = tiffany.consulta("TAG\_GIT\_001 zqwm01 clone") print(hit["domain"], hit["text"])
\prov{relações: parte\_de broca\_os; relaciona utilitarismo; relaciona transcricao; relaciona teoria\_do\_valor; relaciona vida-morte; relaciona ruido; relaciona setor\_algebrico\_descontinuidades\_no\_fecho; relaciona tipos\_de\_interpretante; relaciona scheduler; relaciona toryces\_e\_helyces; relaciona vetores; relaciona virtualizacao}
\prov{proveniência: INSTALL.md \textperiodcentered{} fato \textperiodcentered{} confiança 1.0 \textperiodcentered{} domínio manual \textperiodcentered{} história 14 versões no jornal}

%CRISTAL {"arestas":[{"alvo":"bai","confianca":0.078,"epistemico":"fato","origem":"manual","rel":"referencia"},{"alvo":"koch_como_recipiente_memoria_nao-associativa","confianca":1.0,"epistemico":"fato","origem":"KOCH_MEMORIA.md","rel":"parte_de"}],"confianca":1.0,"contraexemplos":[],"descricao":"| Métrica | Cantor/Fib | Koch | |---------|------------|------| | Δ vizinho | pequeno ((s) contínuo) | ~80 bits no endereco | | corr(word, ·) | alta com (s) | baixa com endereco | | Colisões | — | muitas (conteúdo distinto) | | Recall por similaridade | sim ((s)) | não (atrator) |","epistemico":"fato","exemplos":[],"id":"resultados_esperados","memoria":"semantica","meta":{"dominio":"manual","nivel":6,"verificacao":"slug idêntico — enriquecer, não duplicar"},"origem":"KOCH_MEMORIA.md","palavras_chave":[],"sinonimos":[],"tipo":"conceito","titulo":"Resultados esperados"}
\section{Resultados esperados}
| Métrica | Cantor/Fib | Koch | |---------|------------|------| | \ensuremath{\Delta} vizinho | pequeno ((s) contínuo) | \~{}80 bits no endereco | | corr(word, \textperiodcentered ) | alta com (s) | baixa com endereco | | Colisões | — | muitas (conteúdo distinto) | | Recall por similaridade | sim ((s)) | não (atrator) |
\prov{relações: referencia bai; parte\_de koch\_como\_recipiente\_memoria\_nao-associativa}
\prov{proveniência: KOCH\_MEMORIA.md \textperiodcentered{} fato \textperiodcentered{} confiança 1.0 \textperiodcentered{} domínio manual \textperiodcentered{} história 10 versões no jornal}

%CRISTAL {"arestas":[{"alvo":"koch_codifica_cantorfibonacci_prova_e_evidencia","confianca":1.0,"epistemico":"fato","origem":"wiki","rel":"usa"},{"alvo":"razao_aurea","confianca":1.0,"epistemico":"fato","origem":"wiki","rel":"relaciona"},{"alvo":"decomposicao_peirce","confianca":1.0,"epistemico":"fato","origem":"wiki","rel":"relaciona"}],"confianca":1.0,"contraexemplos":[],"descricao":"A substituição fractal de Koch satisfaz as 7 propriedades do ribossomo (tradutor, processivo, universal, código fixo, ribozima=auto-similar) e traduz o código (3 cúspides algébricas) em carne (curva fractal; floco 1,30; dragão 1,55). Linearizado: N(t)=Σλᵢt=Σet·log λᵢ estende a substituição a todo t real. No coração: maturidade=floco (área converge, perímetro diverge).","epistemico":"fato","exemplos":[],"id":"ribossomo_koch","memoria":"semantica","meta":{"autor":"Lopes/Aarão","dominio":"hiper","fonte":""},"origem":"wiki","palavras_chave":[],"sinonimos":[],"tipo":"conceito","titulo":"O Ribossomo = Substituição de Koch (código→carne)"}
\section{O Ribossomo = Substituição de Koch (código\ensuremath{\to}carne)}
A substituição fractal de Koch satisfaz as 7 propriedades do ribossomo (tradutor, processivo, universal, código fixo, ribozima=auto-similar) e traduz o código (3 cúspides algébricas) em carne (curva fractal; floco 1,30; dragão 1,55). Linearizado: N(t)=\ensuremath{\Sigma}\ensuremath{\lambda}\ensuremath{_{i}}t=\ensuremath{\Sigma}et\textperiodcentered log \ensuremath{\lambda}\ensuremath{_{i}} estende a substituição a todo t real. No coração: maturidade=floco (área converge, perímetro diverge).
\prov{relações: usa koch\_codifica\_cantorfibonacci\_prova\_e\_evidencia; relaciona razao\_aurea; relaciona decomposicao\_peirce}
\prov{proveniência: wiki \textperiodcentered{} fato \textperiodcentered{} confiança 1.0 \textperiodcentered{} domínio hiper \textperiodcentered{} história 6 versões no jornal}

%CRISTAL {"arestas":[{"alvo":"assistente","confianca":1.0,"epistemico":"fato","origem":"paper","rel":"parte_de"},{"alvo":"readme_cristal","confianca":1.0,"epistemico":"fato","origem":"README_CRISTAL.md","rel":"parte_de"}],"confianca":1.0,"contraexemplos":[],"descricao":"| Fase | Entrega | |------|---------| | **0.1** | Tecidos: ingest, index, validate ✓ | | **0.2** | Projeto, agentes, REPL, code/qa/paper/release ✓ | | **0.3** | Pipeline code (refactor→patch→diff), task/work automático | | **0.4** | deploy, plugins hooks, paper publish | | **0.5** | Integração editor, sessão persistente JSONL |","epistemico":"fato","exemplos":[],"id":"roadmap","memoria":"semantica","meta":{"dominio":"manual","nivel":6,"verificacao":"slug idêntico — enriquecer, não duplicar"},"origem":"CRISTAL_DEV.md","palavras_chave":[],"sinonimos":[],"tipo":"conceito","titulo":"Roadmap"}
\section{Roadmap}
| Fase | Entrega | |------|---------| | **0.1** | Tecidos: ingest, index, validate \checkmark  | | **0.2** | Projeto, agentes, REPL, code/qa/paper/release \checkmark  | | **0.3** | Pipeline code (refactor\ensuremath{\to}patch\ensuremath{\to}diff), task/work automático | | **0.4** | deploy, plugins hooks, paper publish | | **0.5** | Integração editor, sessão persistente JSONL |
\prov{relações: parte\_de assistente; parte\_de readme\_cristal}
\prov{proveniência: CRISTAL\_DEV.md \textperiodcentered{} fato \textperiodcentered{} confiança 1.0 \textperiodcentered{} domínio manual \textperiodcentered{} história 9 versões no jornal}

%CRISTAL {"arestas":[{"alvo":"termodinamica","confianca":1.0,"epistemico":"fato","origem":"manual","rel":"referencia"},{"alvo":"virtualizacao","confianca":0.5,"epistemico":"fato","origem":"wiki","rel":"relaciona"},{"alvo":"transcricao","confianca":0.5,"epistemico":"fato","origem":"wiki","rel":"relaciona"},{"alvo":"vida-morte","confianca":0.5,"epistemico":"fato","origem":"wiki","rel":"relaciona"},{"alvo":"utilitarismo","confianca":0.5,"epistemico":"fato","origem":"wiki","rel":"relaciona"},{"alvo":"scheduler","confianca":0.5,"epistemico":"fato","origem":"wiki","rel":"relaciona"},{"alvo":"veredito_experimental","confianca":0.5,"epistemico":"fato","origem":"wiki","rel":"relaciona"},{"alvo":"tipos_de_interpretante","confianca":0.5,"epistemico":"fato","origem":"wiki","rel":"relaciona"},{"alvo":"teoria_do_valor","confianca":0.5,"epistemico":"fato","origem":"wiki","rel":"relaciona"},{"alvo":"setor_algebrico_descontinuidades_no_fecho","confianca":0.5,"epistemico":"fato","origem":"wiki","rel":"relaciona"}],"confianca":1.0,"contraexemplos":[],"descricao":"Perturbação `offset` e `bit_flip` sobre órbita prata A₂·(1,0).","epistemico":"fato","exemplos":[],"id":"ruido","memoria":"semantica","meta":{"dominio":"manual","duplicata_sugerida":[],"nivel":6,"verificacao":"parte de conceito mais amplo 'ruido'"},"origem":"RUÍDO.md","palavras_chave":[],"sinonimos":[],"tipo":"conceito","titulo":"RUÍDO"}
\section{RUÍDO}
Perturbação `offset` e `bit\_flip` sobre órbita prata A\ensuremath{_{2}}\textperiodcentered (1,0).
\prov{relações: referencia termodinamica; relaciona virtualizacao; relaciona transcricao; relaciona vida-morte; relaciona utilitarismo; relaciona scheduler; relaciona veredito\_experimental; relaciona tipos\_de\_interpretante; relaciona teoria\_do\_valor; relaciona setor\_algebrico\_descontinuidades\_no\_fecho}
\prov{proveniência: RUÍDO.md \textperiodcentered{} fato \textperiodcentered{} confiança 1.0 \textperiodcentered{} domínio manual \textperiodcentered{} história 17 versões no jornal}

%CRISTAL {"arestas":[{"alvo":"koch_como_recipiente_memoria_nao-associativa","confianca":1.0,"epistemico":"fato","origem":"KOCH_MEMORIA.md","rel":"parte_de"},{"alvo":"koch_parabola_v2_sobre_a_curva_desdistorcida","confianca":1.0,"epistemico":"fato","origem":"KOCH_PARABOLA.md","rel":"parte_de"}],"confianca":1.0,"contraexemplos":[],"descricao":"- `dados/koch_parabola/report.json` - `atlas/koch_parabola/marcacao_reta.csv` — marcas uniformes em `s`, mapeadas de volta à parábola - `atlas/koch_parabola/oscilacao.csv` - `atlas/koch_parabola/discreta_reta.csv` - `atlas/koch_parabola/cantor_fibonacci.csv` - `atlas/koch_parabola/agm.csv`","epistemico":"fato","exemplos":[],"id":"saidas","memoria":"semantica","meta":{"dominio":"manual","nivel":6,"verificacao":"slug idêntico — enriquecer, não duplicar"},"origem":"KOCH_PARABOLA.md","palavras_chave":[],"sinonimos":[],"tipo":"conceito","titulo":"Saídas"}
\section{Saídas}
- `dados/koch\_parabola/report.json` - `atlas/koch\_parabola/marcacao\_reta.csv` — marcas uniformes em `s`, mapeadas de volta à parábola - `atlas/koch\_parabola/oscilacao.csv` - `atlas/koch\_parabola/discreta\_reta.csv` - `atlas/koch\_parabola/cantor\_fibonacci.csv` - `atlas/koch\_parabola/agm.csv`
\prov{relações: parte\_de koch\_como\_recipiente\_memoria\_nao-associativa; parte\_de koch\_parabola\_v2\_sobre\_a\_curva\_desdistorcida}
\prov{proveniência: KOCH\_PARABOLA.md \textperiodcentered{} fato \textperiodcentered{} confiança 1.0 \textperiodcentered{} domínio manual \textperiodcentered{} história 10 versões no jornal}

%CRISTAL {"arestas":[{"alvo":"peirce_celulas_canais","confianca":1.0,"epistemico":"fato","origem":"wiki","rel":"usa"},{"alvo":"agulha_admissao","confianca":1.0,"epistemico":"fato","origem":"wiki","rel":"relaciona"},{"alvo":"cifra","confianca":1.0,"epistemico":"fato","origem":"wiki","rel":"relaciona"}],"confianca":1.0,"contraexemplos":[],"descricao":"A segurança é dimensional: dentro de uma banda (≤ k_cut) a Máquina é segura por design (não-contaminação e_ie_j=0 dá vazamento interno zero); o side-channel é a banda que excede k_cut e vive na extensão Â∖A (o degrau de Cayley-Dickson), e a criptografia é o band-limiting na fronteira. Zero ganho criptográfico sobre as primitivas.","epistemico":"inferencia","exemplos":[],"id":"seguranca_dimensional","memoria":"semantica","meta":{"autor":"Aarão/broca-so","dominio":"broca_repos","fonte":""},"origem":"wiki","palavras_chave":[],"sinonimos":[],"tipo":"conceito","titulo":"Segurança Dimensional"}
\section{Segurança Dimensional}
A segurança é dimensional: dentro de uma banda (\ensuremath{\le} k\_cut) a Máquina é segura por design (não-contaminação e\_ie\_j=0 dá vazamento interno zero); o side-channel é a banda que excede k\_cut e vive na extensão Â\ensuremath{\setminus}A (o degrau de Cayley-Dickson), e a criptografia é o band-limiting na fronteira. Zero ganho criptográfico sobre as primitivas.
\prov{relações: usa peirce\_celulas\_canais; relaciona agulha\_admissao; relaciona cifra}
\prov{proveniência: wiki \textperiodcentered{} inferencia \textperiodcentered{} confiança 1.0 \textperiodcentered{} domínio broca\_repos \textperiodcentered{} história 4 versões no jornal}

%CRISTAL {"arestas":[{"alvo":"observador_e_observado","confianca":0.95,"epistemico":"fato","origem":"wiki","rel":"confirma"},{"alvo":"semantica","confianca":0.95,"epistemico":"fato","origem":"wiki","rel":"confirma"},{"alvo":"semiotica","confianca":1.0,"epistemico":"fato","origem":"wiki","rel":"parte_de"},{"alvo":"virtualizacao","confianca":0.5,"epistemico":"fato","origem":"wiki","rel":"relaciona"},{"alvo":"word","confianca":0.5,"epistemico":"fato","origem":"wiki","rel":"relaciona"},{"alvo":"vida-morte","confianca":0.5,"epistemico":"fato","origem":"wiki","rel":"relaciona"},{"alvo":"sessao_interativa","confianca":0.5,"epistemico":"fato","origem":"wiki","rel":"relaciona"},{"alvo":"sociologia","confianca":0.5,"epistemico":"fato","origem":"wiki","rel":"relaciona"}],"confianca":0.95,"contraexemplos":[],"descricao":"Referência verificada: semiotica_peirce_triade.md","epistemico":"fato","exemplos":[],"id":"semiotica_peirce_triade","memoria":"semantica","meta":{"dominio":"referencia","fonte":"web","nivel":5},"origem":"wiki","palavras_chave":[],"sinonimos":[],"tipo":"conceito","titulo":"semiotica peirce triade"}
\section{semiotica peirce triade}
Referência verificada: semiotica\_peirce\_triade.md
\prov{relações: confirma observador\_e\_observado; confirma semantica; parte\_de semiotica; relaciona virtualizacao; relaciona word; relaciona vida-morte; relaciona sessao\_interativa; relaciona sociologia}
\prov{proveniência: wiki \textperiodcentered{} web \textperiodcentered{} fato \textperiodcentered{} confiança 0.95 \textperiodcentered{} domínio referencia \textperiodcentered{} história 33 versões no jornal}

%CRISTAL {"arestas":[{"alvo":"semiotica_peirce_triade_semiotica_triadica_de_charles_sanders_peirce","confianca":1.0,"epistemico":"fato","origem":"semiotica_peirce_triade.md","rel":"parte_de"}],"confianca":1.0,"contraexemplos":[],"descricao":"Identificar semiótica triádica Peirce com decomposição idempotente 1=∑ eᵢ — categorias distintas (filosofia vs álgebra).","epistemico":"fato","exemplos":[],"id":"semiotica_peirce_triade_contraexemplos","memoria":"semantica","meta":{"dominio":"referencia","fonte":"web","nivel":5},"origem":"semiotica_peirce_triade.md","palavras_chave":[],"sinonimos":[],"tipo":"conceito","titulo":"Contraexemplos"}
\section{Contraexemplos}
Identificar semiótica triádica Peirce com decomposição idempotente 1=\ensuremath{\sum} e\ensuremath{_{i}} — categorias distintas (filosofia vs álgebra).
\prov{relações: parte\_de semiotica\_peirce\_triade\_semiotica\_triadica\_de\_charles\_sanders\_peirce}
\prov{proveniência: semiotica\_peirce\_triade.md \textperiodcentered{} web \textperiodcentered{} fato \textperiodcentered{} confiança 1.0 \textperiodcentered{} domínio referencia \textperiodcentered{} história 4 versões no jornal}

%CRISTAL {"fusao":[{"arestas":[{"alvo":"semiotica_peirce_triade_semiotica_triadica_de_charles_sanders_peirce","confianca":1.0,"epistemico":"fato","origem":"semiotica_peirce_triade.md","rel":"parte_de"},{"alvo":"a_decomposicao_de_peirce_celulas_e_canais","confianca":0.95,"epistemico":"fato","origem":"wiki","rel":"referencia"}],"confianca":1.0,"contraexemplos":[],"descricao":"**Homónimo crítico** — dois “Peirce” no projecto:","epistemico":"fato","exemplos":[],"id":"semiotica_peirce_triade_distincao_peirce_filosofo_vs_algebra_broca-so","memoria":"semantica","meta":{"dominio":"referencia","fonte":"web","nivel":5},"origem":"semiotica_peirce_triade.md","palavras_chave":[],"sinonimos":[],"tipo":"conceito","titulo":"Distinção Peirce filósofo vs álgebra Broca-SO"},{"arestas":[{"alvo":"semiotica_triadica_de_charles_sanders_peirce","confianca":1.0,"epistemico":"fato","origem":"semiotica_peirce_triade.md","rel":"parte_de"},{"alvo":"gato","confianca":0.95,"epistemico":"fato","origem":"wiki","rel":"referencia"},{"alvo":"a_decomposicao_de_peirce_celulas_e_canais","confianca":0.95,"epistemico":"fato","origem":"wiki","rel":"referencia"}],"confianca":1.0,"contraexemplos":[],"descricao":"**Homónimo crítico** — dois “Peirce” no projecto:","epistemico":"fato","exemplos":[],"id":"distincao_peirce_filosofo_vs_algebra_broca-so","memoria":"semantica","meta":{"dominio":"referencia","nivel":5},"origem":"semiotica_peirce_triade.md","palavras_chave":[],"sinonimos":[],"tipo":"conceito","titulo":"Distinção Peirce filósofo vs álgebra Broca-SO"}],"id":"semiotica_peirce_triade_distincao_peirce_filosofo_vs_algebra_broca-so","tipo":"conceito"}
\section{Distinção Peirce filósofo vs álgebra Broca-SO}
**Homónimo crítico** — dois “Peirce” no projecto:
\prov{relações: parte\_de semiotica\_peirce\_triade\_semiotica\_triadica\_de\_charles\_sanders\_peirce; referencia a\_decomposicao\_de\_peirce\_celulas\_e\_canais}
\prov{proveniência: semiotica\_peirce\_triade.md \textperiodcentered{} web \textperiodcentered{} fato \textperiodcentered{} confiança 1.0 \textperiodcentered{} domínio referencia \textperiodcentered{} história 4 versões no jornal}
\prov{fusão de curadoria (texto idêntico): absorve distincao\_peirce\_filosofo\_vs\_algebra\_broca-so --- o mesmo conteúdo por dois esquemas de endereço}

%CRISTAL {"arestas":[{"alvo":"semiotica_peirce_triade_semiotica_triadica_de_charles_sanders_peirce","confianca":1.0,"epistemico":"fato","origem":"semiotica_peirce_triade.md","rel":"parte_de"}],"confianca":1.0,"contraexemplos":[],"descricao":"- https://plato.stanford.edu/entries/peirce-semiotics/ - https://en.wikipedia.org/wiki/Semiotic_theory_of_Charles_Sanders_Peirce - https://en.wikipedia.org/wiki/Interpretant","epistemico":"fato","exemplos":[],"id":"semiotica_peirce_triade_fontes","memoria":"semantica","meta":{"dominio":"referencia","fonte":"web","nivel":5},"origem":"semiotica_peirce_triade.md","palavras_chave":[],"sinonimos":[],"tipo":"conceito","titulo":"Fontes"}
\section{Fontes}
- https://plato.stanford.edu/entries/peirce-semiotics/ - https://en.wikipedia.org/wiki/Semiotic\_theory\_of\_Charles\_Sanders\_Peirce - https://en.wikipedia.org/wiki/Interpretant
\prov{relações: parte\_de semiotica\_peirce\_triade\_semiotica\_triadica\_de\_charles\_sanders\_peirce}
\prov{proveniência: semiotica\_peirce\_triade.md \textperiodcentered{} web \textperiodcentered{} fato \textperiodcentered{} confiança 1.0 \textperiodcentered{} domínio referencia \textperiodcentered{} história 3 versões no jornal}

%CRISTAL {"arestas":[{"alvo":"semiotica_peirce_triade_semiotica_triadica_de_charles_sanders_peirce","confianca":1.0,"epistemico":"fato","origem":"semiotica_peirce_triade.md","rel":"parte_de"},{"alvo":"semantica","confianca":0.95,"epistemico":"fato","origem":"wiki","rel":"confirma"},{"alvo":"observador_e_observado","confianca":0.95,"epistemico":"fato","origem":"wiki","rel":"confirma"},{"alvo":"colapso","confianca":0.95,"epistemico":"fato","origem":"wiki","rel":"analogia"},{"alvo":"gato","confianca":0.95,"epistemico":"fato","origem":"wiki","rel":"analogia"}],"confianca":1.0,"contraexemplos":[],"descricao":"| Peirce (SEP) | Nó Broca-SO | relação | |--------------|-------------|---------| | interpretante / leitor | observadoreobservado | confirma | | significação triádica | semantica | confirma | | regressão de signos | colapso | analogia | | representamen | analogia | referencia | | decomposição eᵢ (nome) | adecomposicaodepeircecelulasecanais | referencia (homónimo) | | pesos φ/ψ do habitante | gato | analogia (espectro binário, não tríade) |","epistemico":"fato","exemplos":[],"id":"semiotica_peirce_triade_ponte_broca-so","memoria":"semantica","meta":{"dominio":"referencia","fonte":"web","nivel":5},"origem":"semiotica_peirce_triade.md","palavras_chave":[],"sinonimos":[],"tipo":"conceito","titulo":"Ponte Broca-SO"}
\section{Ponte Broca-SO}
| Peirce (SEP) | Nó Broca-SO | relação | |--------------|-------------|---------| | interpretante / leitor | observadoreobservado | confirma | | significação triádica | semantica | confirma | | regressão de signos | colapso | analogia | | representamen | analogia | referencia | | decomposição e\ensuremath{_{i}} (nome) | adecomposicaodepeircecelulasecanais | referencia (homónimo) | | pesos \ensuremath{\varphi}/\ensuremath{\psi} do habitante | gato | analogia (espectro binário, não tríade) |
\prov{relações: parte\_de semiotica\_peirce\_triade\_semiotica\_triadica\_de\_charles\_sanders\_peirce; confirma semantica; confirma observador\_e\_observado; analogia colapso; analogia gato}
\prov{proveniência: semiotica\_peirce\_triade.md \textperiodcentered{} web \textperiodcentered{} fato \textperiodcentered{} confiança 1.0 \textperiodcentered{} domínio referencia \textperiodcentered{} história 8 versões no jornal}

%CRISTAL {"arestas":[{"alvo":"semiotica_peirce_triade_semiotica_triadica_de_charles_sanders_peirce","confianca":1.0,"epistemico":"fato","origem":"semiotica_peirce_triade.md","rel":"parte_de"}],"confianca":1.0,"contraexemplos":[],"descricao":"- **Charles Sanders Peirce (1839–1914):** lógica como semeiótica formal - **Stanford Encyclopedia:** [Peirce's Theory of Signs](https://plato.stanford.edu/entries/peirce-semiotics/) (Atkin, rev. 2025) - **Wikipedia:** [Semiotic theory of Charles Sanders Peirce](https://en.wikipedia.org/wiki/Semiotic_theory_of_Charles_Sanders_Peirce) - **Termo técnico:** *semeiosis* — acção triádica irredutível entre signo, objecto e interpretante","epistemico":"fato","exemplos":[],"id":"semiotica_peirce_triade_referencia","memoria":"semantica","meta":{"dominio":"referencia","fonte":"web","nivel":5},"origem":"semiotica_peirce_triade.md","palavras_chave":[],"sinonimos":[],"tipo":"conceito","titulo":"Referência"}
\section{Referência}
- **Charles Sanders Peirce (1839–1914):** lógica como semeiótica formal - **Stanford Encyclopedia:** [Peirce's Theory of Signs](https://plato.stanford.edu/entries/peirce-semiotics/) (Atkin, rev. 2025) - **Wikipedia:** [Semiotic theory of Charles Sanders Peirce](https://en.wikipedia.org/wiki/Semiotic\_theory\_of\_Charles\_Sanders\_Peirce) - **Termo técnico:** *semeiosis* — acção triádica irredutível entre signo, objecto e interpretante
\prov{relações: parte\_de semiotica\_peirce\_triade\_semiotica\_triadica\_de\_charles\_sanders\_peirce}
\prov{proveniência: semiotica\_peirce\_triade.md \textperiodcentered{} web \textperiodcentered{} fato \textperiodcentered{} confiança 1.0 \textperiodcentered{} domínio referencia \textperiodcentered{} história 3 versões no jornal}

%CRISTAL {"arestas":[{"alvo":"semiotica_peirce_triade_semiotica_triadica_de_charles_sanders_peirce","confianca":1.0,"epistemico":"fato","origem":"semiotica_peirce_triade.md","rel":"parte_de"},{"alvo":"ponte_quantica_dougherty","confianca":0.95,"epistemico":"fato","origem":"wiki","rel":"referencia"},{"alvo":"observador_e_observado","confianca":0.95,"epistemico":"fato","origem":"wiki","rel":"analogia"}],"confianca":1.0,"contraexemplos":[],"descricao":"Dougherty menciona leitura projectiva (quem mapeia, instrumento) nos orbitais. Peirce oferece **quadro filosófico** para observador↔signo↔objecto — **analogia** apenas.","epistemico":"fato","exemplos":[],"id":"semiotica_peirce_triade_relacao_com_hipotese_dougherty","memoria":"semantica","meta":{"dominio":"referencia","fonte":"web","nivel":5},"origem":"semiotica_peirce_triade.md","palavras_chave":[],"sinonimos":[],"tipo":"conceito","titulo":"Relação com hipótese Dougherty"}
\section{Relação com hipótese Dougherty}
Dougherty menciona leitura projectiva (quem mapeia, instrumento) nos orbitais. Peirce oferece **quadro filosófico** para observador\ensuremath{\leftrightarrow}signo\ensuremath{\leftrightarrow}objecto — **analogia** apenas.
\prov{relações: parte\_de semiotica\_peirce\_triade\_semiotica\_triadica\_de\_charles\_sanders\_peirce; referencia ponte\_quantica\_dougherty; analogia observador\_e\_observado}
\prov{proveniência: semiotica\_peirce\_triade.md \textperiodcentered{} web \textperiodcentered{} fato \textperiodcentered{} confiança 1.0 \textperiodcentered{} domínio referencia \textperiodcentered{} história 5 versões no jornal}

%CRISTAL {"fusao":[{"arestas":[{"alvo":"semiotica_peirce_triade_semiotica_triadica_de_charles_sanders_peirce","confianca":1.0,"epistemico":"fato","origem":"semiotica_peirce_triade.md","rel":"parte_de"},{"alvo":"observador_e_observado","confianca":0.95,"epistemico":"fato","origem":"wiki","rel":"confirma"}],"confianca":1.0,"contraexemplos":[],"descricao":"Peirce: onde há signo há sequência infinita de representamens do mesmo objecto. A mente (interpretante) não é opcional — é constitutiva do signo.","epistemico":"fato","exemplos":[],"id":"semiotica_peirce_triade_semeiosis_e_observador","memoria":"semantica","meta":{"dominio":"referencia","fonte":"web","nivel":5},"origem":"semiotica_peirce_triade.md","palavras_chave":[],"sinonimos":[],"tipo":"conceito","titulo":"Semeiosis e observador"},{"arestas":[{"alvo":"semiotica_triadica_de_charles_sanders_peirce","confianca":1.0,"epistemico":"fato","origem":"semiotica_peirce_triade.md","rel":"parte_de"},{"alvo":"colapso","confianca":0.95,"epistemico":"fato","origem":"wiki","rel":"analogia"},{"alvo":"observador_e_observado","confianca":0.95,"epistemico":"fato","origem":"wiki","rel":"confirma"}],"confianca":1.0,"contraexemplos":[],"descricao":"Peirce: onde há signo há sequência infinita de representamens do mesmo objecto. A mente (interpretante) não é opcional — é constitutiva do signo.","epistemico":"fato","exemplos":[],"id":"semeiosis_e_observador","memoria":"semantica","meta":{"dominio":"referencia","nivel":5},"origem":"semiotica_peirce_triade.md","palavras_chave":[],"sinonimos":[],"tipo":"conceito","titulo":"Semeiosis e observador"}],"id":"semiotica_peirce_triade_semeiosis_e_observador","tipo":"conceito"}
\section{Semeiosis e observador}
Peirce: onde há signo há sequência infinita de representamens do mesmo objecto. A mente (interpretante) não é opcional — é constitutiva do signo.
\prov{relações: parte\_de semiotica\_peirce\_triade\_semiotica\_triadica\_de\_charles\_sanders\_peirce; confirma observador\_e\_observado}
\prov{proveniência: semiotica\_peirce\_triade.md \textperiodcentered{} web \textperiodcentered{} fato \textperiodcentered{} confiança 1.0 \textperiodcentered{} domínio referencia \textperiodcentered{} história 4 versões no jornal}
\prov{fusão de curadoria (texto idêntico): absorve semeiosis\_e\_observador --- o mesmo conteúdo por dois esquemas de endereço}

%CRISTAL {"fusao":[{"arestas":[{"alvo":"semiotica_peirce_triade","confianca":0.95,"epistemico":"fato","origem":"wiki","rel":"parte_de"}],"confianca":1.0,"contraexemplos":[],"descricao":"> **Epistemologia:** fato filosófico estabelecido (Stanford SEP, obra Peirce). > **No grafo:** `epistemico=fato`, confiança 0.95. > **Papel:** fundamentar observador/semântica no Broca-SO; distinguir de decomposição algébrica Peirce no código.","epistemico":"fato","exemplos":[],"id":"semiotica_peirce_triade_semiotica_triadica_de_charles_sanders_peirce","memoria":"semantica","meta":{"dominio":"referencia","fonte":"web","nivel":5},"origem":"semiotica_peirce_triade.md","palavras_chave":[],"sinonimos":[],"tipo":"conceito","titulo":"Semiótica triádica de Charles Sanders Peirce"},{"arestas":[{"alvo":"semiotica_peirce_triade","confianca":0.95,"epistemico":"fato","origem":"wiki","rel":"parte_de"}],"confianca":1.0,"contraexemplos":[],"descricao":"> **Epistemologia:** fato filosófico estabelecido (Stanford SEP, obra Peirce). > **No grafo:** `epistemico=fato`, confiança 0.95. > **Papel:** fundamentar observador/semântica no Broca-SO; distinguir de decomposição algébrica Peirce no código.","epistemico":"fato","exemplos":[],"id":"semiotica_triadica_de_charles_sanders_peirce","memoria":"semantica","meta":{"dominio":"referencia","nivel":5},"origem":"semiotica_peirce_triade.md","palavras_chave":[],"sinonimos":[],"tipo":"conceito","titulo":"Semiótica triádica de Charles Sanders Peirce"}],"id":"semiotica_peirce_triade_semiotica_triadica_de_charles_sanders_peirce","tipo":"conceito"}
\section{Semiótica triádica de Charles Sanders Peirce}
> **Epistemologia:** fato filosófico estabelecido (Stanford SEP, obra Peirce). > **No grafo:** `epistemico=fato`, confiança 0.95. > **Papel:** fundamentar observador/semântica no Broca-SO; distinguir de decomposição algébrica Peirce no código.
\prov{relações: parte\_de semiotica\_peirce\_triade}
\prov{proveniência: semiotica\_peirce\_triade.md \textperiodcentered{} web \textperiodcentered{} fato \textperiodcentered{} confiança 1.0 \textperiodcentered{} domínio referencia \textperiodcentered{} história 3 versões no jornal}
\prov{fusão de curadoria (texto idêntico): absorve semiotica\_triadica\_de\_charles\_sanders\_peirce --- o mesmo conteúdo por dois esquemas de endereço}

%CRISTAL {"fusao":[{"arestas":[{"alvo":"semiotica_peirce_triade_semiotica_triadica_de_charles_sanders_peirce","confianca":1.0,"epistemico":"fato","origem":"semiotica_peirce_triade.md","rel":"parte_de"},{"alvo":"tipos_de_interpretante","confianca":0.5,"epistemico":"fato","origem":"wiki","rel":"relaciona"},{"alvo":"word","confianca":0.5,"epistemico":"fato","origem":"wiki","rel":"relaciona"},{"alvo":"vontade_de_poder","confianca":0.5,"epistemico":"fato","origem":"wiki","rel":"relaciona"},{"alvo":"transcricao","confianca":0.5,"epistemico":"fato","origem":"wiki","rel":"relaciona"}],"confianca":1.0,"contraexemplos":[],"descricao":"Peirce distingue (entre outros):","epistemico":"fato","exemplos":[],"id":"semiotica_peirce_triade_tipos_de_interpretante","memoria":"semantica","meta":{"dominio":"referencia","fonte":"web","nivel":5},"origem":"semiotica_peirce_triade.md","palavras_chave":[],"sinonimos":[],"tipo":"conceito","titulo":"Tipos de interpretante"},{"arestas":[{"alvo":"semiotica_triadica_de_charles_sanders_peirce","confianca":1.0,"epistemico":"fato","origem":"semiotica_peirce_triade.md","rel":"parte_de"},{"alvo":"utilitarismo","confianca":0.5,"epistemico":"fato","origem":"wiki","rel":"relaciona"},{"alvo":"vida-morte","confianca":0.5,"epistemico":"fato","origem":"wiki","rel":"relaciona"},{"alvo":"transcricao","confianca":0.5,"epistemico":"fato","origem":"wiki","rel":"relaciona"},{"alvo":"virtualizacao","confianca":0.5,"epistemico":"fato","origem":"wiki","rel":"relaciona"},{"alvo":"veredito_experimental","confianca":0.5,"epistemico":"fato","origem":"wiki","rel":"relaciona"}],"confianca":1.0,"contraexemplos":[],"descricao":"Peirce distingue (entre outros):","epistemico":"fato","exemplos":[],"id":"tipos_de_interpretante","memoria":"semantica","meta":{"dominio":"referencia","nivel":5},"origem":"semiotica_peirce_triade.md","palavras_chave":[],"sinonimos":[],"tipo":"conceito","titulo":"Tipos de interpretante"}],"id":"semiotica_peirce_triade_tipos_de_interpretante","tipo":"conceito"}
\section{Tipos de interpretante}
Peirce distingue (entre outros):
\prov{relações: parte\_de semiotica\_peirce\_triade\_semiotica\_triadica\_de\_charles\_sanders\_peirce; relaciona tipos\_de\_interpretante; relaciona word; relaciona vontade\_de\_poder; relaciona transcricao}
\prov{proveniência: semiotica\_peirce\_triade.md \textperiodcentered{} web \textperiodcentered{} fato \textperiodcentered{} confiança 1.0 \textperiodcentered{} domínio referencia \textperiodcentered{} história 7 versões no jornal}
\prov{fusão de curadoria (texto idêntico): absorve tipos\_de\_interpretante --- o mesmo conteúdo por dois esquemas de endereço}

%CRISTAL {"fusao":[{"arestas":[{"alvo":"semiotica_peirce_triade_semiotica_triadica_de_charles_sanders_peirce","confianca":1.0,"epistemico":"fato","origem":"semiotica_peirce_triade.md","rel":"parte_de"},{"alvo":"observador_e_observado","confianca":0.95,"epistemico":"fato","origem":"wiki","rel":"confirma"},{"alvo":"word","confianca":0.5,"epistemico":"fato","origem":"wiki","rel":"relaciona"},{"alvo":"virtualizacao","confianca":0.5,"epistemico":"fato","origem":"wiki","rel":"relaciona"}],"confianca":1.0,"contraexemplos":[],"descricao":"Relação **triádica irredutível** (não decomponível em pares):","epistemico":"fato","exemplos":[],"id":"semiotica_peirce_triade_triade_signo-objeto-interpretante","memoria":"semantica","meta":{"dominio":"referencia","fonte":"web","nivel":5},"origem":"semiotica_peirce_triade.md","palavras_chave":[],"sinonimos":[],"tipo":"conceito","titulo":"Tríade signo-objeto-interpretante"},{"arestas":[{"alvo":"semiotica_triadica_de_charles_sanders_peirce","confianca":1.0,"epistemico":"fato","origem":"semiotica_peirce_triade.md","rel":"parte_de"},{"alvo":"observador_e_observado","confianca":0.95,"epistemico":"fato","origem":"wiki","rel":"confirma"},{"alvo":"word","confianca":0.5,"epistemico":"fato","origem":"wiki","rel":"relaciona"}],"confianca":1.0,"contraexemplos":[],"descricao":"Relação **triádica irredutível** (não decomponível em pares):","epistemico":"fato","exemplos":[],"id":"triade_signo-objeto-interpretante","memoria":"semantica","meta":{"dominio":"referencia","nivel":5},"origem":"semiotica_peirce_triade.md","palavras_chave":[],"sinonimos":[],"tipo":"conceito","titulo":"Tríade signo-objeto-interpretante"}],"id":"semiotica_peirce_triade_triade_signo-objeto-interpretante","tipo":"conceito"}
\section{Tríade signo-objeto-interpretante}
Relação **triádica irredutível** (não decomponível em pares):
\prov{relações: parte\_de semiotica\_peirce\_triade\_semiotica\_triadica\_de\_charles\_sanders\_peirce; confirma observador\_e\_observado; relaciona word; relaciona virtualizacao}
\prov{proveniência: semiotica\_peirce\_triade.md \textperiodcentered{} web \textperiodcentered{} fato \textperiodcentered{} confiança 1.0 \textperiodcentered{} domínio referencia \textperiodcentered{} história 6 versões no jornal}
\prov{fusão de curadoria (texto idêntico): absorve triade\_signo-objeto-interpretante --- o mesmo conteúdo por dois esquemas de endereço}

%CRISTAL {"arestas":[{"alvo":"broca_os","confianca":0.95,"epistemico":"fato","origem":"paper","rel":"parte_de"}],"confianca":1.0,"contraexemplos":[],"descricao":"| Componente | Papel | Analogia | |------------|-------|----------| | **libtiffany** | Motor de recuperação (overlap, idx, mmap) | LLVM / runtime | | **Tiffany** | Assistente conversacional | Copilot chat | | **Cristal** | **Ambiente de desenvolvimento** da plataforma | git + cargo + npm | | **Broca** | Sistema operacional / orquestração | OS |","epistemico":"fato","exemplos":[],"id":"separacao_de_papeis","memoria":"semantica","meta":{"dominio":"manual","nivel":6,"verificacao":"slug idêntico — enriquecer, não duplicar"},"origem":"CRISTAL_DEV.md","palavras_chave":[],"sinonimos":[],"tipo":"conceito","titulo":"Separação de papéis"}
\section{Separação de papéis}
| Componente | Papel | Analogia | |------------|-------|----------| | **libtiffany** | Motor de recuperação (overlap, idx, mmap) | LLVM / runtime | | **Tiffany** | Assistente conversacional | Copilot chat | | **Cristal** | **Ambiente de desenvolvimento** da plataforma | git + cargo + npm | | **Broca** | Sistema operacional / orquestração | OS |
\prov{relações: parte\_de broca\_os}
\prov{proveniência: CRISTAL\_DEV.md \textperiodcentered{} fato \textperiodcentered{} confiança 1.0 \textperiodcentered{} domínio manual \textperiodcentered{} história 3 versões no jornal}

%CRISTAL {"arestas":[{"alvo":"broca_os","confianca":0.95,"epistemico":"fato","origem":"paper","rel":"parte_de"},{"alvo":"transcricao","confianca":0.5,"epistemico":"fato","origem":"wiki","rel":"relaciona"},{"alvo":"teste_ordem","confianca":0.5,"epistemico":"fato","origem":"wiki","rel":"relaciona"},{"alvo":"utilitarismo","confianca":0.5,"epistemico":"fato","origem":"wiki","rel":"relaciona"},{"alvo":"topologia","confianca":0.5,"epistemico":"fato","origem":"wiki","rel":"relaciona"},{"alvo":"virtualizacao","confianca":0.5,"epistemico":"fato","origem":"wiki","rel":"relaciona"},{"alvo":"word","confianca":0.5,"epistemico":"fato","origem":"wiki","rel":"relaciona"},{"alvo":"while","confianca":0.5,"epistemico":"fato","origem":"wiki","rel":"relaciona"}],"confianca":1.0,"contraexemplos":[],"descricao":"```bash cristal ```","epistemico":"fato","exemplos":[],"id":"sessao_interativa","memoria":"semantica","meta":{"dominio":"manual","nivel":6,"verificacao":"slug idêntico — enriquecer, não duplicar"},"origem":"CRISTAL_DEV.md","palavras_chave":[],"sinonimos":[],"tipo":"conceito","titulo":"Sessão interativa"}
\section{Sessão interativa}
```bash cristal ```
\prov{relações: parte\_de broca\_os; relaciona transcricao; relaciona teste\_ordem; relaciona utilitarismo; relaciona topologia; relaciona virtualizacao; relaciona word; relaciona while}
\prov{proveniência: CRISTAL\_DEV.md \textperiodcentered{} fato \textperiodcentered{} confiança 1.0 \textperiodcentered{} domínio manual \textperiodcentered{} história 10 versões no jornal}

%CRISTAL {"arestas":[{"alvo":"pow_metal_setor","confianca":1.0,"epistemico":"fato","origem":"manual","rel":"parte_de"},{"alvo":"utilitarismo","confianca":0.5,"epistemico":"fato","origem":"wiki","rel":"relaciona"},{"alvo":"vetores","confianca":0.5,"epistemico":"fato","origem":"wiki","rel":"relaciona"},{"alvo":"vida-morte","confianca":0.5,"epistemico":"fato","origem":"wiki","rel":"relaciona"},{"alvo":"teoria_do_valor","confianca":0.5,"epistemico":"fato","origem":"wiki","rel":"relaciona"},{"alvo":"transcricao","confianca":0.5,"epistemico":"fato","origem":"wiki","rel":"relaciona"}],"confianca":1.0,"contraexemplos":[],"descricao":"Secção de POW_METAL_SETOR.md.","epistemico":"fato","exemplos":[],"id":"setor_algebrico_descontinuidades_no_fecho","memoria":"semantica","meta":{"dominio":"manual","nivel":6,"verificacao":"slug idêntico — enriquecer, não duplicar"},"origem":"POW_METAL_SETOR.md","palavras_chave":[],"sinonimos":[],"tipo":"conceito","titulo":"Setor Algébrico — Descontinuidades no Fecho"}
\section{Setor Algébrico — Descontinuidades no Fecho}
Secção de POW\_METAL\_SETOR.md.
\prov{relações: parte\_de pow\_metal\_setor; relaciona utilitarismo; relaciona vetores; relaciona vida-morte; relaciona teoria\_do\_valor; relaciona transcricao}
\prov{proveniência: POW\_METAL\_SETOR.md \textperiodcentered{} fato \textperiodcentered{} confiança 1.0 \textperiodcentered{} domínio manual \textperiodcentered{} história 9 versões no jornal}

%CRISTAL {"arestas":[{"alvo":"torre_corpos_sinh","confianca":1.0,"epistemico":"fato","origem":"wiki","rel":"usa"},{"alvo":"cifra","confianca":1.0,"epistemico":"fato","origem":"wiki","rel":"explica"},{"alvo":"prova_de_trabalho","confianca":1.0,"epistemico":"fato","origem":"wiki","rel":"relaciona"}],"confianca":1.0,"contraexemplos":[],"descricao":"Leitura algébrica da cripto: o SHA implementa o 'osso' de um número de grau infinito (topo da torre; forma do acaso, Marchenko-Pastur). Verificado por PSLQ: a cúspide θ₂ fecha em grau 4, mas o SHA não fecha com altura <10¹². Inverter exigiria reconstruir a órbita de Galois infinita — nada a agarrar. (Distinção honesta: o SHA real é intratável 2²⁵⁶, não indecidível.)","epistemico":"fato","exemplos":[],"id":"sha_osso_do_topo","memoria":"semantica","meta":{"autor":"Lopes/Aarão","dominio":"hiper","fonte":""},"origem":"wiki","palavras_chave":[],"sinonimos":[],"tipo":"conceito","titulo":"O SHA como Osso de um Número do Topo"}
\section{O SHA como Osso de um Número do Topo}
Leitura algébrica da cripto: o SHA implementa o 'osso' de um número de grau infinito (topo da torre; forma do acaso, Marchenko-Pastur). Verificado por PSLQ: a cúspide \ensuremath{\theta}\ensuremath{_{2}} fecha em grau 4, mas o SHA não fecha com altura <10\ensuremath{^{12}}. Inverter exigiria reconstruir a órbita de Galois infinita — nada a agarrar. (Distinção honesta: o SHA real é intratável 2\ensuremath{^{256}}, não indecidível.)
\prov{relações: usa torre\_corpos\_sinh; explica cifra; relaciona prova\_de\_trabalho}
\prov{proveniência: wiki \textperiodcentered{} fato \textperiodcentered{} confiança 1.0 \textperiodcentered{} domínio hiper \textperiodcentered{} história 5 versões no jornal}

%CRISTAL {"arestas":[{"alvo":"spec_trabalho","confianca":1.0,"epistemico":"fato","origem":"BROCA_SO.md","rel":"parte_de"}],"confianca":1.0,"contraexemplos":[],"descricao":"| Slot | Conteúdo | |------|----------| | `TRABALHO` | bytes do job | | `GATO` | `Word total, e ` | | `ESPECTRO` | `double[12⁴]` | | `BANDA` | SHA256 hex do tecido | | `NONCE` | derivado da banda |","epistemico":"fato","exemplos":[],"id":"slots_logicos","memoria":"semantica","meta":{"dominio":"manual","nivel":6,"verificacao":"slug idêntico — enriquecer, não duplicar"},"origem":"BROCA_SO.md","palavras_chave":[],"sinonimos":[],"tipo":"conceito","titulo":"Slots lógicos"}
\section{Slots lógicos}
| Slot | Conteúdo | |------|----------| | `TRABALHO` | bytes do job | | `GATO` | `Word total, e ` | | `ESPECTRO` | `double[12\ensuremath{^{4}}]` | | `BANDA` | SHA256 hex do tecido | | `NONCE` | derivado da banda |
\prov{relações: parte\_de spec\_trabalho}
\prov{proveniência: BROCA\_SO.md \textperiodcentered{} fato \textperiodcentered{} confiança 1.0 \textperiodcentered{} domínio manual \textperiodcentered{} história 7 versões no jornal}

%CRISTAL {"arestas":[{"alvo":"broca_os","confianca":0.95,"epistemico":"fato","origem":"paper","rel":"parte_de"}],"confianca":1.0,"contraexemplos":[],"descricao":"- meta de passos: **1,000,000** - passos totais: **1,716,582** - passos válidos (com distribuição nos dados): **699,682** - violações |a+c²−1| > 1e-9: **0** - erro máximo: **0.0** - programas executados: **1013**","epistemico":"fato","exemplos":[],"id":"sobreviveu","memoria":"semantica","meta":{"dominio":"manual","nivel":6,"verificacao":"slug idêntico — enriquecer, não duplicar"},"origem":"PARABOLA_DESTRUIR.md","palavras_chave":[],"sinonimos":[],"tipo":"conceito","titulo":"✅ Sobreviveu"}
\section{[ok] Sobreviveu}
- meta de passos: **1,000,000** - passos totais: **1,716,582** - passos válidos (com distribuição nos dados): **699,682** - violações |a+c\ensuremath{^{2}}\ensuremath{-}1| > 1e-9: **0** - erro máximo: **0.0** - programas executados: **1013**
\prov{relações: parte\_de broca\_os}
\prov{proveniência: PARABOLA\_DESTRUIR.md \textperiodcentered{} fato \textperiodcentered{} confiança 1.0 \textperiodcentered{} domínio manual \textperiodcentered{} história 3 versões no jornal}

%CRISTAL {"arestas":[{"alvo":"respiracao_dos_programas","confianca":1.0,"epistemico":"fato","origem":"PARABOLA_ISA.md","rel":"parte_de"},{"alvo":"word","confianca":0.5,"epistemico":"fato","origem":"wiki","rel":"relaciona"},{"alvo":"virtualizacao","confianca":0.5,"epistemico":"fato","origem":"wiki","rel":"relaciona"}],"confianca":1.0,"contraexemplos":[],"descricao":"- calor acumulado: **0.700** - fase acumulada: **4.602** - perda: **0.000** - potência (calor/passo): **0.140** - passos: **5**","epistemico":"fato","exemplos":[],"id":"soma","memoria":"semantica","meta":{"dominio":"manual","nivel":6,"verificacao":"slug idêntico — enriquecer, não duplicar"},"origem":"PARABOLA_ISA.md","palavras_chave":[],"sinonimos":[],"tipo":"conceito","titulo":"🌱 soma"}
\section{[semente] soma}
- calor acumulado: **0.700** - fase acumulada: **4.602** - perda: **0.000** - potência (calor/passo): **0.140** - passos: **5**
\prov{relações: parte\_de respiracao\_dos\_programas; relaciona word; relaciona virtualizacao}
\prov{proveniência: PARABOLA\_ISA.md \textperiodcentered{} fato \textperiodcentered{} confiança 1.0 \textperiodcentered{} domínio manual \textperiodcentered{} história 8 versões no jornal}

%CRISTAL {"arestas":[{"alvo":"broca-so_camada_de_trabalho_isomorfica","confianca":1.0,"epistemico":"fato","origem":"BROCA_SO.md","rel":"parte_de"}],"confianca":1.0,"contraexemplos":[],"descricao":"| Ficheiro | Papel | |----------|-------| | `broca_so.h` | `BrocaMaquina`, slots, `BrocaOp`, `broca_so_exec()` | | `broca_backend.h` | `BrocaBackend` plugável | | `broca_borda.h` | Ops de contemplação (fora do núcleo) |","epistemico":"fato","exemplos":[],"id":"spec_trabalho","memoria":"semantica","meta":{"dominio":"manual","nivel":6,"verificacao":"slug idêntico — enriquecer, não duplicar"},"origem":"BROCA_SO.md","palavras_chave":[],"sinonimos":[],"tipo":"conceito","titulo":"Spec (trabalho)"}
\section{Spec (trabalho)}
| Ficheiro | Papel | |----------|-------| | `broca\_so.h` | `BrocaMaquina`, slots, `BrocaOp`, `broca\_so\_exec()` | | `broca\_backend.h` | `BrocaBackend` plugável | | `broca\_borda.h` | Ops de contemplação (fora do núcleo) |
\prov{relações: parte\_de broca-so\_camada\_de\_trabalho\_isomorfica}
\prov{proveniência: BROCA\_SO.md \textperiodcentered{} fato \textperiodcentered{} confiança 1.0 \textperiodcentered{} domínio manual \textperiodcentered{} história 6 versões no jornal}

%CRISTAL {"arestas":[{"alvo":"funcao_onda_secao_hopf","confianca":1.0,"epistemico":"fato","origem":"wiki","rel":"usa"},{"alvo":"fibracao_de_hopf","confianca":1.0,"epistemico":"fato","origem":"wiki","rel":"usa"}],"confianca":1.0,"contraexemplos":[],"descricao":"Construção explícita: Vs=Sym²s(ℍ) (dim 2s+1), geradores Jâ^(s)=−(iħ/2)Σ Râ^(k) satisfazem su(2), Casimir ħ²s(s+1), espectro −s,…,s; via Borel-Weil Sym²s(ℍ)≅H⁰(ℂP¹,O(2s)). Verificado até s=5/2 a ε de máquina.","epistemico":"fato","exemplos":[],"id":"spin_s_sym_quaternion","memoria":"semantica","meta":{"autor":"Lopes/Aarão","dominio":"hiper","fonte":""},"origem":"wiki","palavras_chave":[],"sinonimos":[],"tipo":"conceito","titulo":"Spin-s via Sym²s(ℍ)"}
\section{Spin-s via Sym\ensuremath{^{2}}s(\ensuremath{\mathbb{H}})}
Construção explícita: Vs=Sym\ensuremath{^{2}}s(\ensuremath{\mathbb{H}}) (dim 2s+1), geradores Jâ\^{}(s)=\ensuremath{-}(i\ensuremath{\hbar}/2)\ensuremath{\Sigma} Râ\^{}(k) satisfazem su(2), Casimir \ensuremath{\hbar}\ensuremath{^{2}}s(s+1), espectro \ensuremath{-}s,…,s; via Borel-Weil Sym\ensuremath{^{2}}s(\ensuremath{\mathbb{H}})\ensuremath{\cong}H\ensuremath{^{0}}(\ensuremath{\mathbb{C}}P\ensuremath{^{1}},O(2s)). Verificado até s=5/2 a \ensuremath{\varepsilon} de máquina.
\prov{relações: usa funcao\_onda\_secao\_hopf; usa fibracao\_de\_hopf}
\prov{proveniência: wiki \textperiodcentered{} fato \textperiodcentered{} confiança 1.0 \textperiodcentered{} domínio hiper \textperiodcentered{} história 6 versões no jornal}

%CRISTAL {"arestas":[{"alvo":"mult_prtsu","confianca":1.0,"epistemico":"fato","origem":"wiki","rel":"generaliza"},{"alvo":"torre_hurwitz","confianca":1.0,"epistemico":"fato","origem":"wiki","rel":"usa"},{"alvo":"algebras_hipercomplexas","confianca":1.0,"epistemico":"fato","origem":"wiki","rel":"usa"},{"alvo":"matematica_estelar","confianca":1.0,"epistemico":"fato","origem":"wiki","rel":"relaciona"}],"confianca":1.0,"contraexemplos":[],"descricao":"Multiplicação hipercomplexa em ℝⁿ parametrizada por (p,r,t,u,s,v,w_p,h)∈ℝ⁸ que recupera ℝ,ℂ,ℍ,𝕆 em pontos privilegiados; o oitavo parâmetro h acopla à 4-forma de Cayley e só existe em ℝ⁸. Extensão da família ×ₛ / do produto mult_prtsu.","epistemico":"fato","exemplos":[],"id":"substrato_prtsuv_h","memoria":"semantica","meta":{"autor":"Lopes/Aarão","dominio":"hiper","fonte":""},"origem":"wiki","palavras_chave":[],"sinonimos":[],"tipo":"conceito","titulo":"Substrato PRTSUV+h (8 parâmetros)"}
\section{Substrato PRTSUV+h (8 parâmetros)}
Multiplicação hipercomplexa em \ensuremath{\mathbb{R}}\ensuremath{^{n}} parametrizada por (p,r,t,u,s,v,w\_p,h)\ensuremath{\in}\ensuremath{\mathbb{R}}\ensuremath{^{8}} que recupera \ensuremath{\mathbb{R}},\ensuremath{\mathbb{C}},\ensuremath{\mathbb{H}},\ensuremath{\mathbb{O}} em pontos privilegiados; o oitavo parâmetro h acopla à 4-forma de Cayley e só existe em \ensuremath{\mathbb{R}}\ensuremath{^{8}}. Extensão da família $\times$\ensuremath{_{s}} / do produto mult\_prtsu.
\prov{relações: generaliza mult\_prtsu; usa torre\_hurwitz; usa algebras\_hipercomplexas; relaciona matematica\_estelar}
\prov{proveniência: wiki \textperiodcentered{} fato \textperiodcentered{} confiança 1.0 \textperiodcentered{} domínio hiper \textperiodcentered{} história 6 versões no jornal}

%CRISTAL {"arestas":[{"alvo":"koch_codifica_cantorfibonacci_prova_e_evidencia","confianca":1.0,"epistemico":"fato","origem":"KOCH_CANTOR_FIB_PROVA.md","rel":"parte_de"}],"confianca":1.0,"contraexemplos":[],"descricao":"| Métrica | Esperado | Interpretação | |---------|----------|---------------| | `fracaofibonacciproximo₀05` | (≫ 50%) | ouro marca (s) | | `taxaviolacao₁1lsb` | (≈ 26%) | máquina finita; limite contínuo em aberto | | `fracaoscantorreal` | (∼ 5–10%) | medida zero | | `fracaoscantorouro` | (∼ 10–20%) | medida zero | | `concordanciasectorcantor` | — | **não** usar como prova directa |","epistemico":"fato","exemplos":[],"id":"tabela_de_evidencia_make_koch-cantor-fib","memoria":"semantica","meta":{"dominio":"manual","nivel":6,"verificacao":"slug idêntico — enriquecer, não duplicar"},"origem":"KOCH_CANTOR_FIB_PROVA.md","palavras_chave":[],"sinonimos":[],"tipo":"conceito","titulo":"Tabela de evidência (`make koch-cantor-fib`)"}
\section{Tabela de evidência (`make koch-cantor-fib`)}
| Métrica | Esperado | Interpretação | |---------|----------|---------------| | `fracaofibonacciproximo\ensuremath{_{0}}05` | (\ensuremath{\gg} 50\%) | ouro marca (s) | | `taxaviolacao\ensuremath{_{1}}1lsb` | (\ensuremath{\approx} 26\%) | máquina finita; limite contínuo em aberto | | `fracaoscantorreal` | (\ensuremath{\sim} 5–10\%) | medida zero | | `fracaoscantorouro` | (\ensuremath{\sim} 10–20\%) | medida zero | | `concordanciasectorcantor` | — | **não** usar como prova directa |
\prov{relações: parte\_de koch\_codifica\_cantorfibonacci\_prova\_e\_evidencia}
\prov{proveniência: KOCH\_CANTOR\_FIB\_PROVA.md \textperiodcentered{} fato \textperiodcentered{} confiança 1.0 \textperiodcentered{} domínio manual \textperiodcentered{} história 5 versões no jornal}

%CRISTAL {"arestas":[{"alvo":"pow_taxa_de_sucesso_do_bit_broca","confianca":1.0,"epistemico":"fato","origem":"POW_BIT_TAXA.md","rel":"parte_de"}],"confianca":1.0,"contraexemplos":[],"descricao":"| Métrica | Valor | |---------|-------| | Taxa base P(SHA ok) | 0.00001550 (0.0016%) | | Taxa aceitação bit | 0.1392 (13.92%) | | P(SHA ok | bit aceito) | 0.00000718 (0.0007%) | | P(SHA ok | bit rejeitado) | 0.00001684 | | Enriquecimento (pós/base) | 0.4632× | | Precisão bit (TP/(TP+FP)) | 0.00000718 | | Recall bit (TP/(TP+FN)) | 0.0645 |","epistemico":"fato","exemplos":[],"id":"taxa_base_vs_pos-bit","memoria":"semantica","meta":{"dominio":"manual","nivel":6,"verificacao":"slug idêntico — enriquecer, não duplicar"},"origem":"POW_BIT_TAXA.md","palavras_chave":[],"sinonimos":[],"tipo":"conceito","titulo":"Taxa base vs pós-bit"}
\section{Taxa base vs pós-bit}
| Métrica | Valor | |---------|-------| | Taxa base P(SHA ok) | 0.00001550 (0.0016\%) | | Taxa aceitação bit | 0.1392 (13.92\%) | | P(SHA ok | bit aceito) | 0.00000718 (0.0007\%) | | P(SHA ok | bit rejeitado) | 0.00001684 | | Enriquecimento (pós/base) | 0.4632$\times$ | | Precisão bit (TP/(TP+FP)) | 0.00000718 | | Recall bit (TP/(TP+FN)) | 0.0645 |
\prov{relações: parte\_de pow\_taxa\_de\_sucesso\_do\_bit\_broca}
\prov{proveniência: POW\_BIT\_TAXA.md \textperiodcentered{} fato \textperiodcentered{} confiança 1.0 \textperiodcentered{} domínio manual \textperiodcentered{} história 5 versões no jornal}

%CRISTAL {"arestas":[{"alvo":"matematica_estelar","confianca":1.0,"epistemico":"fato","origem":"wiki","rel":"parte_de"},{"alvo":"nucleo_g2_reabsorcao","confianca":1.0,"epistemico":"fato","origem":"wiki","rel":"usa"},{"alvo":"decomposicao_peirce","confianca":1.0,"epistemico":"fato","origem":"wiki","rel":"relaciona"}],"confianca":1.0,"contraexemplos":[],"descricao":"Toda a estrutura estelar é o tensor de ordem 3 μlᵢj (μlᵢj=s·φlᵢj, φ = 3-forma de G₂); a operação de n corpos é a contração de n−1 cópias de μ numa árvore binária (Catalan(n−1) parêntesamentos). Cada corpo = uma cópia a mais de μ, sobe uma ordem tensorial.","epistemico":"fato","exemplos":[],"id":"tensor_mu_torre_corpos","memoria":"semantica","meta":{"autor":"Lopes/Aarão","dominio":"hiper","fonte":""},"origem":"wiki","palavras_chave":[],"sinonimos":[],"tipo":"conceito","titulo":"O Tensor μ e a Torre de n-Corpos"}
\section{O Tensor \ensuremath{\mu} e a Torre de n-Corpos}
Toda a estrutura estelar é o tensor de ordem 3 \ensuremath{\mu}l\ensuremath{_{i}}j (\ensuremath{\mu}l\ensuremath{_{i}}j=s\textperiodcentered \ensuremath{\varphi}l\ensuremath{_{i}}j, \ensuremath{\varphi} = 3-forma de G\ensuremath{_{2}}); a operação de n corpos é a contração de n\ensuremath{-}1 cópias de \ensuremath{\mu} numa árvore binária (Catalan(n\ensuremath{-}1) parêntesamentos). Cada corpo = uma cópia a mais de \ensuremath{\mu}, sobe uma ordem tensorial.
\prov{relações: parte\_de matematica\_estelar; usa nucleo\_g2\_reabsorcao; relaciona decomposicao\_peirce}
\prov{proveniência: wiki \textperiodcentered{} fato \textperiodcentered{} confiança 1.0 \textperiodcentered{} domínio hiper \textperiodcentered{} história 6 versões no jornal}

%CRISTAL {"arestas":[{"alvo":"4_obstrucao_formal_nao_turing-completa","confianca":1.0,"epistemico":"fato","origem":"EXPRESSIVIDADE.md","rel":"parte_de"},{"alvo":"broca_os","confianca":1.0,"epistemico":"fato","origem":"manual","rel":"parte_de"}],"confianca":1.0,"contraexemplos":[],"descricao":"**Afirmação:** não existe composição de opcodes que execute `μ[f(A,B,R)]` para função arbitrária `f` dos registradores.","epistemico":"fato","exemplos":[],"id":"teorema_1_enderecamento_estatico","memoria":"semantica","meta":{"dominio":"manual","nivel":6,"verificacao":"slug idêntico — enriquecer, não duplicar"},"origem":"EXPRESSIVIDADE.md","palavras_chave":[],"sinonimos":[],"tipo":"conceito","titulo":"Teorema 1 — Endereçamento estático"}
\section{Teorema 1 — Endereçamento estático}
**Afirmação:** não existe composição de opcodes que execute `\ensuremath{\mu}[f(A,B,R)]` para função arbitrária `f` dos registradores.
\prov{relações: parte\_de 4\_obstrucao\_formal\_nao\_turing-completa; parte\_de broca\_os}
\prov{proveniência: EXPRESSIVIDADE.md \textperiodcentered{} fato \textperiodcentered{} confiança 1.0 \textperiodcentered{} domínio manual \textperiodcentered{} história 10 versões no jornal}

%CRISTAL {"arestas":[{"alvo":"koch_codifica_cantorfibonacci_prova_e_evidencia","confianca":1.0,"epistemico":"fato","origem":"KOCH_CANTOR_FIB_PROVA.md","rel":"parte_de"},{"alvo":"word","confianca":0.5,"epistemico":"fato","origem":"wiki","rel":"relaciona"},{"alvo":"vontade_de_poder","confianca":0.5,"epistemico":"fato","origem":"wiki","rel":"relaciona"}],"confianca":1.0,"contraexemplos":[],"descricao":"**Enunciado.** Para Word inteiro (w), seja (s=Φ(c²(w,gold(w)))). Então","epistemico":"fato","exemplos":[],"id":"teorema_1_fibonacci_no_parametro_s_empirico_forte","memoria":"semantica","meta":{"dominio":"manual","nivel":6,"verificacao":"slug idêntico — enriquecer, não duplicar"},"origem":"KOCH_CANTOR_FIB_PROVA.md","palavras_chave":[],"sinonimos":[],"tipo":"conceito","titulo":"Teorema 1 (Fibonacci no parâmetro s) — **empírico forte**"}
\section{Teorema 1 (Fibonacci no parâmetro s) — **empírico forte**}
**Enunciado.** Para Word inteiro (w), seja (s=\ensuremath{\Phi}(c\ensuremath{^{2}}(w,gold(w)))). Então
\prov{relações: parte\_de koch\_codifica\_cantorfibonacci\_prova\_e\_evidencia; relaciona word; relaciona vontade\_de\_poder}
\prov{proveniência: KOCH\_CANTOR\_FIB\_PROVA.md \textperiodcentered{} fato \textperiodcentered{} confiança 1.0 \textperiodcentered{} domínio manual \textperiodcentered{} história 7 versões no jornal}

%CRISTAL {"arestas":[{"alvo":"koch_codifica_cantorfibonacci_prova_e_evidencia","confianca":1.0,"epistemico":"fato","origem":"KOCH_CANTOR_FIB_PROVA.md","rel":"parte_de"}],"confianca":1.0,"contraexemplos":[],"descricao":"**Enunciado.** Seja (σₖın0,ldots,15) o sector da órbita (zₖ+₁=zₖ² ±od p). Os bits LSB (bₖ=σₖ bmod 2) satisfazem (bₖ bₖ+₁ eq 11) com taxa de violação (≪ 1) (medida empírica: ver `taxaviolacao₁1lsb` em `report.json`).","epistemico":"fato","exemplos":[],"id":"teorema_2_alfabeto_sem_11_na_orbita_koch_estrutural_parcial","memoria":"semantica","meta":{"dominio":"manual","nivel":6,"verificacao":"slug idêntico — enriquecer, não duplicar"},"origem":"KOCH_CANTOR_FIB_PROVA.md","palavras_chave":[],"sinonimos":[],"tipo":"conceito","titulo":"Teorema 2 (alfabeto sem `11` na órbita Koch) — **estrutural parcial**"}
\section{Teorema 2 (alfabeto sem `11` na órbita Koch) — **estrutural parcial**}
**Enunciado.** Seja (\ensuremath{\sigma}\ensuremath{_{k}}\ensuremath{\imath}n0,ldots,15) o sector da órbita (z\ensuremath{_{k}}+\ensuremath{_{1}}=z\ensuremath{_{k}}\ensuremath{^{2}} $\pm$od p). Os bits LSB (b\ensuremath{_{k}}=\ensuremath{\sigma}\ensuremath{_{k}} bmod 2) satisfazem (b\ensuremath{_{k}} b\ensuremath{_{k}}+\ensuremath{_{1}} eq 11) com taxa de violação (\ensuremath{\ll} 1) (medida empírica: ver `taxaviolacao\ensuremath{_{1}}1lsb` em `report.json`).
\prov{relações: parte\_de koch\_codifica\_cantorfibonacci\_prova\_e\_evidencia}
\prov{proveniência: KOCH\_CANTOR\_FIB\_PROVA.md \textperiodcentered{} fato \textperiodcentered{} confiança 1.0 \textperiodcentered{} domínio manual \textperiodcentered{} história 5 versões no jornal}

%CRISTAL {"arestas":[{"alvo":"4_obstrucao_formal_nao_turing-completa","confianca":1.0,"epistemico":"fato","origem":"EXPRESSIVIDADE.md","rel":"parte_de"},{"alvo":"broca_os","confianca":1.0,"epistemico":"fato","origem":"manual","rel":"parte_de"},{"alvo":"word","confianca":0.5,"epistemico":"fato","origem":"wiki","rel":"relaciona"}],"confianca":1.0,"contraexemplos":[],"descricao":"**Afirmação:** não existe composição que execute `PC ← g(A,B,R)` para `g` arbitrária.","epistemico":"fato","exemplos":[],"id":"teorema_2_saltos_estaticos","memoria":"semantica","meta":{"dominio":"manual","nivel":6,"verificacao":"slug idêntico — enriquecer, não duplicar"},"origem":"EXPRESSIVIDADE.md","palavras_chave":[],"sinonimos":[],"tipo":"conceito","titulo":"Teorema 2 — Saltos estáticos"}
\section{Teorema 2 — Saltos estáticos}
**Afirmação:** não existe composição que execute `PC \ensuremath{\leftarrow} g(A,B,R)` para `g` arbitrária.
\prov{relações: parte\_de 4\_obstrucao\_formal\_nao\_turing-completa; parte\_de broca\_os; relaciona word}
\prov{proveniência: EXPRESSIVIDADE.md \textperiodcentered{} fato \textperiodcentered{} confiança 1.0 \textperiodcentered{} domínio manual \textperiodcentered{} história 11 versões no jornal}

%CRISTAL {"arestas":[{"alvo":"4_obstrucao_formal_nao_turing-completa","confianca":1.0,"epistemico":"fato","origem":"EXPRESSIVIDADE.md","rel":"parte_de"},{"alvo":"broca_os","confianca":1.0,"epistemico":"fato","origem":"manual","rel":"parte_de"},{"alvo":"word","confianca":0.5,"epistemico":"fato","origem":"wiki","rel":"relaciona"}],"confianca":1.0,"contraexemplos":[],"descricao":"**Afirmação:** um programa `π` acede a conjunto finito `S ⊂ Slot` fixo em assemble-time.","epistemico":"fato","exemplos":[],"id":"teorema_3_memoria_finita_nomeada","memoria":"semantica","meta":{"dominio":"manual","nivel":6,"verificacao":"slug idêntico — enriquecer, não duplicar"},"origem":"EXPRESSIVIDADE.md","palavras_chave":[],"sinonimos":[],"tipo":"conceito","titulo":"Teorema 3 — Memória finita nomeada"}
\section{Teorema 3 — Memória finita nomeada}
**Afirmação:** um programa `\ensuremath{\pi}` acede a conjunto finito `S \ensuremath{\subset} Slot` fixo em assemble-time.
\prov{relações: parte\_de 4\_obstrucao\_formal\_nao\_turing-completa; parte\_de broca\_os; relaciona word}
\prov{proveniência: EXPRESSIVIDADE.md \textperiodcentered{} fato \textperiodcentered{} confiança 1.0 \textperiodcentered{} domínio manual \textperiodcentered{} história 11 versões no jornal}

%CRISTAL {"arestas":[{"alvo":"4_obstrucao_formal_nao_turing-completa","confianca":1.0,"epistemico":"fato","origem":"EXPRESSIVIDADE.md","rel":"parte_de"},{"alvo":"broca_os","confianca":1.0,"epistemico":"fato","origem":"manual","rel":"parte_de"},{"alvo":"word","confianca":0.5,"epistemico":"fato","origem":"wiki","rel":"relaciona"},{"alvo":"virtualizacao","confianca":0.5,"epistemico":"fato","origem":"wiki","rel":"relaciona"}],"confianca":1.0,"contraexemplos":[],"descricao":"**Afirmação:** cada componente de `Word` é `long` finito; contadores não crescem além de `2⁶⁴`.","epistemico":"fato","exemplos":[],"id":"teorema_4_palavras_limitadas","memoria":"semantica","meta":{"dominio":"manual","nivel":6,"verificacao":"slug idêntico — enriquecer, não duplicar"},"origem":"EXPRESSIVIDADE.md","palavras_chave":[],"sinonimos":[],"tipo":"conceito","titulo":"Teorema 4 — Palavras limitadas"}
\section{Teorema 4 — Palavras limitadas}
**Afirmação:** cada componente de `Word` é `long` finito; contadores não crescem além de `2\ensuremath{^{64}}`.
\prov{relações: parte\_de 4\_obstrucao\_formal\_nao\_turing-completa; parte\_de broca\_os; relaciona word; relaciona virtualizacao}
\prov{proveniência: EXPRESSIVIDADE.md \textperiodcentered{} fato \textperiodcentered{} confiança 1.0 \textperiodcentered{} domínio manual \textperiodcentered{} história 12 versões no jornal}

%CRISTAL {"arestas":[{"alvo":"koch_codifica_cantorfibonacci_prova_e_evidencia","confianca":1.0,"epistemico":"fato","origem":"KOCH_CANTOR_FIB_PROVA.md","rel":"parte_de"},{"alvo":"word","confianca":0.5,"epistemico":"fato","origem":"wiki","rel":"relaciona"},{"alvo":"virtualizacao","confianca":0.5,"epistemico":"fato","origem":"wiki","rel":"relaciona"}],"confianca":1.0,"contraexemplos":[],"descricao":"**Conjectura.** No limite contínuo (p→ınfty), (N→ınfty)","epistemico":"fato","exemplos":[],"id":"teorema_4_produto_cantorfibonacci_conjectura","memoria":"semantica","meta":{"dominio":"manual","nivel":6,"verificacao":"slug idêntico — enriquecer, não duplicar"},"origem":"KOCH_CANTOR_FIB_PROVA.md","palavras_chave":[],"sinonimos":[],"tipo":"conceito","titulo":"Teorema 4 (produto Cantor–Fibonacci) — **conjectura**"}
\section{Teorema 4 (produto Cantor–Fibonacci) — **conjectura**}
**Conjectura.** No limite contínuo (p\ensuremath{\to}\ensuremath{\imath}nfty), (N\ensuremath{\to}\ensuremath{\imath}nfty)
\prov{relações: parte\_de koch\_codifica\_cantorfibonacci\_prova\_e\_evidencia; relaciona word; relaciona virtualizacao}
\prov{proveniência: KOCH\_CANTOR\_FIB\_PROVA.md \textperiodcentered{} fato \textperiodcentered{} confiança 1.0 \textperiodcentered{} domínio manual \textperiodcentered{} história 7 versões no jornal}

%CRISTAL {"arestas":[{"alvo":"respiracao_dos_programas","confianca":1.0,"epistemico":"fato","origem":"PARABOLA_ISA.md","rel":"parte_de"},{"alvo":"word","confianca":0.5,"epistemico":"fato","origem":"wiki","rel":"relaciona"},{"alvo":"virtualizacao","confianca":0.5,"epistemico":"fato","origem":"wiki","rel":"relaciona"},{"alvo":"joalheria_hub","confianca":1.0,"epistemico":"fato","origem":"wiki","rel":"relaciona"}],"confianca":1.0,"contraexemplos":[],"descricao":"- calor acumulado: **0.500** - fase acumulada: **15.707** - perda: **0.000** - potência (calor/passo): **0.031** - passos: **16**","epistemico":"fato","exemplos":[],"id":"tesouraria","memoria":"semantica","meta":{"dominio":"manual","nivel":6,"verificacao":"slug idêntico — enriquecer, não duplicar"},"origem":"PARABOLA_ISA.md","palavras_chave":[],"sinonimos":[],"tipo":"conceito","titulo":"🌱 tesouraria"}
\section{[semente] tesouraria}
- calor acumulado: **0.500** - fase acumulada: **15.707** - perda: **0.000** - potência (calor/passo): **0.031** - passos: **16**
\prov{relações: parte\_de respiracao\_dos\_programas; relaciona word; relaciona virtualizacao; relaciona joalheria\_hub}
\prov{proveniência: PARABOLA\_ISA.md \textperiodcentered{} fato \textperiodcentered{} confiança 1.0 \textperiodcentered{} domínio manual \textperiodcentered{} história 10 versões no jornal}

%CRISTAL {"arestas":[{"alvo":"readme_cristal","confianca":1.0,"epistemico":"fato","origem":"README_CRISTAL.md","rel":"parte_de"},{"alvo":"transcricao","confianca":0.5,"epistemico":"fato","origem":"wiki","rel":"relaciona"},{"alvo":"vida-morte","confianca":0.5,"epistemico":"fato","origem":"wiki","rel":"relaciona"},{"alvo":"topologia","confianca":0.5,"epistemico":"fato","origem":"wiki","rel":"relaciona"},{"alvo":"thread","confianca":0.5,"epistemico":"fato","origem":"wiki","rel":"relaciona"},{"alvo":"utilitarismo","confianca":0.5,"epistemico":"fato","origem":"wiki","rel":"relaciona"},{"alvo":"virtualizacao","confianca":0.5,"epistemico":"fato","origem":"wiki","rel":"relaciona"}],"confianca":1.0,"contraexemplos":[],"descricao":"```bash chmod +x test/test_cristal.sh sh test/test_cristal.sh ```","epistemico":"fato","exemplos":[],"id":"testes","memoria":"semantica","meta":{"dominio":"manual","nivel":6,"verificacao":"slug idêntico — enriquecer, não duplicar"},"origem":"README_CRISTAL.md","palavras_chave":[],"sinonimos":[],"tipo":"conceito","titulo":"Testes"}
\section{Testes}
```bash chmod +x test/test\_cristal.sh sh test/test\_cristal.sh ```
\prov{relações: parte\_de readme\_cristal; relaciona transcricao; relaciona vida-morte; relaciona topologia; relaciona thread; relaciona utilitarismo; relaciona virtualizacao}
\prov{proveniência: README\_CRISTAL.md \textperiodcentered{} fato \textperiodcentered{} confiança 1.0 \textperiodcentered{} domínio manual \textperiodcentered{} história 10 versões no jornal}

%CRISTAL {"arestas":[{"alvo":"gato","confianca":0.95,"epistemico":"fato","origem":"paper","rel":"parte_de"}],"confianca":1.0,"contraexemplos":[],"descricao":"| Etapa | cand/s | |-------|--------| | Neuronio | 2,482,188 | | Gauss | 40,230,041 | | Gold | 32,163,021 | | Interface | 25,020,476 | | Koch | 24,809,990 | | Lemniscata | 3,181,318 | | Fecho | 41,838,287 | | Hash valido | 79,224 |","epistemico":"fato","exemplos":[],"id":"throughput_por_etapa","memoria":"semantica","meta":{"dominio":"manual","nivel":6,"verificacao":"slug idêntico — enriquecer, não duplicar"},"origem":"POW_METAL_FUNIL.md","palavras_chave":[],"sinonimos":[],"tipo":"conceito","titulo":"Throughput por etapa"}
\section{Throughput por etapa}
| Etapa | cand/s | |-------|--------| | Neuronio | 2,482,188 | | Gauss | 40,230,041 | | Gold | 32,163,021 | | Interface | 25,020,476 | | Koch | 24,809,990 | | Lemniscata | 3,181,318 | | Fecho | 41,838,287 | | Hash valido | 79,224 |
\prov{relações: parte\_de gato}
\prov{proveniência: POW\_METAL\_FUNIL.md \textperiodcentered{} fato \textperiodcentered{} confiança 1.0 \textperiodcentered{} domínio manual \textperiodcentered{} história 3 versões no jornal}

%CRISTAL {"arestas":[{"alvo":"tiffany","confianca":0.95,"epistemico":"fato","origem":"wiki","rel":"confirma"},{"alvo":"transcricao","confianca":0.5,"epistemico":"fato","origem":"wiki","rel":"relaciona"},{"alvo":"topologia","confianca":0.5,"epistemico":"fato","origem":"wiki","rel":"relaciona"},{"alvo":"vida-morte","confianca":0.5,"epistemico":"fato","origem":"wiki","rel":"relaciona"},{"alvo":"virtualizacao","confianca":0.5,"epistemico":"fato","origem":"wiki","rel":"relaciona"},{"alvo":"word","confianca":0.5,"epistemico":"fato","origem":"wiki","rel":"relaciona"},{"alvo":"vontade_de_poder","confianca":0.5,"epistemico":"fato","origem":"wiki","rel":"relaciona"}],"confianca":0.95,"contraexemplos":[],"descricao":"Referência verificada: tiffany_cabeca_broca.md","epistemico":"fato","exemplos":[],"id":"tiffany_cabeca_broca","memoria":"semantica","meta":{"dominio":"referencia","fonte":"web","nivel":5},"origem":"wiki","palavras_chave":[],"sinonimos":[],"tipo":"conceito","titulo":"tiffany cabeca broca"}
\section{tiffany cabeca broca}
Referência verificada: tiffany\_cabeca\_broca.md
\prov{relações: confirma tiffany; relaciona transcricao; relaciona topologia; relaciona vida-morte; relaciona virtualizacao; relaciona word; relaciona vontade\_de\_poder}
\prov{proveniência: wiki \textperiodcentered{} web \textperiodcentered{} fato \textperiodcentered{} confiança 0.95 \textperiodcentered{} domínio referencia \textperiodcentered{} história 15 versões no jornal}

%CRISTAL {"fusao":[{"arestas":[{"alvo":"tiffany_cabeca_broca_tiffany_cabeca_conversacional_broca-so","confianca":1.0,"epistemico":"fato","origem":"tiffany_cabeca_broca.md","rel":"parte_de"},{"alvo":"hopfield","confianca":0.95,"epistemico":"fato","origem":"wiki","rel":"especializa"},{"alvo":"inteligencia_artificial","confianca":0.95,"epistemico":"fato","origem":"wiki","rel":"especializa"},{"alvo":"mcculloch_pitts","confianca":0.95,"epistemico":"fato","origem":"wiki","rel":"referencia"}],"confianca":1.0,"contraexemplos":[],"descricao":"``` McCulloch-Pitts (1943) → Hopfield (1982) → Tiffany (cabeça Broca-SO) ```","epistemico":"fato","exemplos":[],"id":"tiffany_cabeca_broca_cadeia_ia_no_projecto","memoria":"semantica","meta":{"dominio":"referencia","fonte":"web","nivel":5},"origem":"tiffany_cabeca_broca.md","palavras_chave":[],"sinonimos":[],"tipo":"conceito","titulo":"Cadeia IA no projecto"},{"arestas":[{"alvo":"tiffany_cabeca_conversacional_broca-so","confianca":1.0,"epistemico":"fato","origem":"tiffany_cabeca_broca.md","rel":"parte_de"},{"alvo":"hopfield","confianca":0.95,"epistemico":"fato","origem":"wiki","rel":"especializa"},{"alvo":"inteligencia_artificial","confianca":0.95,"epistemico":"fato","origem":"wiki","rel":"especializa"},{"alvo":"mcculloch_pitts","confianca":0.95,"epistemico":"fato","origem":"wiki","rel":"referencia"}],"confianca":1.0,"contraexemplos":[],"descricao":"``` McCulloch-Pitts (1943) → Hopfield (1982) → Tiffany (cabeça Broca-SO) ```","epistemico":"fato","exemplos":[],"id":"cadeia_ia_no_projecto","memoria":"semantica","meta":{"dominio":"referencia","nivel":5},"origem":"tiffany_cabeca_broca.md","palavras_chave":[],"sinonimos":[],"tipo":"conceito","titulo":"Cadeia IA no projecto"}],"id":"tiffany_cabeca_broca_cadeia_ia_no_projecto","tipo":"conceito"}
\section{Cadeia IA no projecto}
``` McCulloch-Pitts (1943) \ensuremath{\to} Hopfield (1982) \ensuremath{\to} Tiffany (cabeça Broca-SO) ```
\prov{relações: parte\_de tiffany\_cabeca\_broca\_tiffany\_cabeca\_conversacional\_broca-so; especializa hopfield; especializa inteligencia\_artificial; referencia mcculloch\_pitts}
\prov{proveniência: tiffany\_cabeca\_broca.md \textperiodcentered{} web \textperiodcentered{} fato \textperiodcentered{} confiança 1.0 \textperiodcentered{} domínio referencia \textperiodcentered{} história 6 versões no jornal}
\prov{fusão de curadoria (texto idêntico): absorve cadeia\_ia\_no\_projecto --- o mesmo conteúdo por dois esquemas de endereço}

%CRISTAL {"arestas":[{"alvo":"tiffany_cabeca_broca_tiffany_cabeca_conversacional_broca-so","confianca":1.0,"epistemico":"fato","origem":"tiffany_cabeca_broca.md","rel":"parte_de"},{"alvo":"word","confianca":0.5,"epistemico":"fato","origem":"wiki","rel":"relaciona"},{"alvo":"virtualizacao","confianca":0.5,"epistemico":"fato","origem":"wiki","rel":"relaciona"},{"alvo":"vida-morte","confianca":0.5,"epistemico":"fato","origem":"wiki","rel":"relaciona"}],"confianca":1.0,"contraexemplos":[],"descricao":"Tiffany **≠** modelo de linguagem estocástico (GPT) — sem sampling neural autoregressivo.","epistemico":"fato","exemplos":[],"id":"tiffany_cabeca_broca_contraexemplos","memoria":"semantica","meta":{"dominio":"referencia","fonte":"web","nivel":5},"origem":"tiffany_cabeca_broca.md","palavras_chave":[],"sinonimos":[],"tipo":"conceito","titulo":"Contraexemplos"}
\section{Contraexemplos}
Tiffany **\ensuremath{\ne}** modelo de linguagem estocástico (GPT) — sem sampling neural autoregressivo.
\prov{relações: parte\_de tiffany\_cabeca\_broca\_tiffany\_cabeca\_conversacional\_broca-so; relaciona word; relaciona virtualizacao; relaciona vida-morte}
\prov{proveniência: tiffany\_cabeca\_broca.md \textperiodcentered{} web \textperiodcentered{} fato \textperiodcentered{} confiança 1.0 \textperiodcentered{} domínio referencia \textperiodcentered{} história 6 versões no jornal}

%CRISTAL {"fusao":[{"arestas":[{"alvo":"tiffany_cabeca_broca_tiffany_cabeca_conversacional_broca-so","confianca":1.0,"epistemico":"fato","origem":"tiffany_cabeca_broca.md","rel":"parte_de"},{"alvo":"codigo_neuronio_experimentos_tiffany","confianca":0.95,"epistemico":"fato","origem":"wiki","rel":"confirma"},{"alvo":"word","confianca":0.5,"epistemico":"fato","origem":"wiki","rel":"relaciona"}],"confianca":1.0,"contraexemplos":[],"descricao":"`neuronio/experimentos_tiffany.py` — suites:","epistemico":"fato","exemplos":[],"id":"tiffany_cabeca_broca_experimentos_reprodutiveis","memoria":"semantica","meta":{"dominio":"referencia","fonte":"web","nivel":5},"origem":"tiffany_cabeca_broca.md","palavras_chave":[],"sinonimos":[],"tipo":"conceito","titulo":"Experimentos reprodutíveis"},{"arestas":[{"alvo":"tiffany_cabeca_conversacional_broca-so","confianca":1.0,"epistemico":"fato","origem":"tiffany_cabeca_broca.md","rel":"parte_de"},{"alvo":"codigo_neuronio_experimentos_tiffany","confianca":0.95,"epistemico":"fato","origem":"wiki","rel":"confirma"},{"alvo":"word","confianca":0.5,"epistemico":"fato","origem":"wiki","rel":"relaciona"},{"alvo":"virtualizacao","confianca":0.5,"epistemico":"fato","origem":"wiki","rel":"relaciona"},{"alvo":"vida-morte","confianca":0.5,"epistemico":"fato","origem":"wiki","rel":"relaciona"}],"confianca":1.0,"contraexemplos":[],"descricao":"`neuronio/experimentos_tiffany.py` — suites:","epistemico":"fato","exemplos":[],"id":"experimentos_reprodutiveis","memoria":"semantica","meta":{"dominio":"referencia","nivel":5},"origem":"tiffany_cabeca_broca.md","palavras_chave":[],"sinonimos":[],"tipo":"conceito","titulo":"Experimentos reprodutíveis"}],"id":"tiffany_cabeca_broca_experimentos_reprodutiveis","tipo":"conceito"}
\section{Experimentos reprodutíveis}
`neuronio/experimentos\_tiffany.py` — suites:
\prov{relações: parte\_de tiffany\_cabeca\_broca\_tiffany\_cabeca\_conversacional\_broca-so; confirma codigo\_neuronio\_experimentos\_tiffany; relaciona word}
\prov{proveniência: tiffany\_cabeca\_broca.md \textperiodcentered{} web \textperiodcentered{} fato \textperiodcentered{} confiança 1.0 \textperiodcentered{} domínio referencia \textperiodcentered{} história 5 versões no jornal}
\prov{fusão de curadoria (texto idêntico): absorve experimentos\_reprodutiveis --- o mesmo conteúdo por dois esquemas de endereço}

%CRISTAL {"arestas":[{"alvo":"tiffany_cabeca_broca_tiffany_cabeca_conversacional_broca-so","confianca":1.0,"epistemico":"fato","origem":"tiffany_cabeca_broca.md","rel":"parte_de"},{"alvo":"word","confianca":0.5,"epistemico":"fato","origem":"wiki","rel":"relaciona"},{"alvo":"vontade_de_poder","confianca":0.5,"epistemico":"fato","origem":"wiki","rel":"relaciona"}],"confianca":1.0,"contraexemplos":[],"descricao":"- `papers/Tiffany.tex` - `neuronio/experimentos_tiffany.py` - `conversa/colapso.py` - `README_CRISTAL.md`","epistemico":"fato","exemplos":[],"id":"tiffany_cabeca_broca_fontes","memoria":"semantica","meta":{"dominio":"referencia","fonte":"web","nivel":5},"origem":"tiffany_cabeca_broca.md","palavras_chave":[],"sinonimos":[],"tipo":"conceito","titulo":"Fontes"}
\section{Fontes}
- `papers/Tiffany.tex` - `neuronio/experimentos\_tiffany.py` - `conversa/colapso.py` - `README\_CRISTAL.md`
\prov{relações: parte\_de tiffany\_cabeca\_broca\_tiffany\_cabeca\_conversacional\_broca-so; relaciona word; relaciona vontade\_de\_poder}
\prov{proveniência: tiffany\_cabeca\_broca.md \textperiodcentered{} web \textperiodcentered{} fato \textperiodcentered{} confiança 1.0 \textperiodcentered{} domínio referencia \textperiodcentered{} história 5 versões no jornal}

%CRISTAL {"arestas":[{"alvo":"tiffany_cabeca_broca_tiffany_cabeca_conversacional_broca-so","confianca":1.0,"epistemico":"fato","origem":"tiffany_cabeca_broca.md","rel":"parte_de"},{"alvo":"colapso","confianca":0.95,"epistemico":"fato","origem":"wiki","rel":"usa"},{"alvo":"broca_os","confianca":0.95,"epistemico":"fato","origem":"wiki","rel":"parte_de"},{"alvo":"word","confianca":0.5,"epistemico":"fato","origem":"wiki","rel":"relaciona"}],"confianca":1.0,"contraexemplos":[],"descricao":"**Cabeça do sistema** Broca-SO: modelo **M = U·P** (unidade × percepção).","epistemico":"fato","exemplos":[],"id":"tiffany_cabeca_broca_o_que_e","memoria":"semantica","meta":{"dominio":"referencia","fonte":"web","nivel":5},"origem":"tiffany_cabeca_broca.md","palavras_chave":[],"sinonimos":[],"tipo":"conceito","titulo":"O que é"}
\section{O que é}
**Cabeça do sistema** Broca-SO: modelo **M = U\textperiodcentered P** (unidade $\times$ percepção).
\prov{relações: parte\_de tiffany\_cabeca\_broca\_tiffany\_cabeca\_conversacional\_broca-so; usa colapso; parte\_de broca\_os; relaciona word}
\prov{proveniência: tiffany\_cabeca\_broca.md \textperiodcentered{} web \textperiodcentered{} fato \textperiodcentered{} confiança 1.0 \textperiodcentered{} domínio referencia \textperiodcentered{} história 6 versões no jornal}

%CRISTAL {"fusao":[{"arestas":[{"alvo":"tiffany_cabeca_broca_tiffany_cabeca_conversacional_broca-so","confianca":1.0,"epistemico":"fato","origem":"tiffany_cabeca_broca.md","rel":"parte_de"},{"alvo":"colapso","confianca":0.95,"epistemico":"fato","origem":"wiki","rel":"usa"}],"confianca":1.0,"contraexemplos":[],"descricao":"1. Utilizador fala → `prosa.impressao` 2. `colapso()` varre camadas (dialogos → conhecimento → corpus) 3. Face vencedora + tag → resposta legível 4. `lab.py porque <fala>` explica ov/score/camada","epistemico":"fato","exemplos":[],"id":"tiffany_cabeca_broca_pipeline_conversacional","memoria":"semantica","meta":{"dominio":"referencia","fonte":"web","nivel":5},"origem":"tiffany_cabeca_broca.md","palavras_chave":[],"sinonimos":[],"tipo":"conceito","titulo":"Pipeline conversacional"},{"arestas":[{"alvo":"tiffany_cabeca_conversacional_broca-so","confianca":1.0,"epistemico":"fato","origem":"tiffany_cabeca_broca.md","rel":"parte_de"},{"alvo":"colapso","confianca":0.95,"epistemico":"fato","origem":"wiki","rel":"usa"}],"confianca":1.0,"contraexemplos":[],"descricao":"1. Utilizador fala → `prosa.impressao` 2. `colapso()` varre camadas (dialogos → conhecimento → corpus) 3. Face vencedora + tag → resposta legível 4. `lab.py porque <fala>` explica ov/score/camada","epistemico":"fato","exemplos":[],"id":"pipeline_conversacional","memoria":"semantica","meta":{"dominio":"referencia","nivel":5},"origem":"tiffany_cabeca_broca.md","palavras_chave":[],"sinonimos":[],"tipo":"conceito","titulo":"Pipeline conversacional"}],"id":"tiffany_cabeca_broca_pipeline_conversacional","tipo":"conceito"}
\section{Pipeline conversacional}
1. Utilizador fala \ensuremath{\to} `prosa.impressao` 2. `colapso()` varre camadas (dialogos \ensuremath{\to} conhecimento \ensuremath{\to} corpus) 3. Face vencedora + tag \ensuremath{\to} resposta legível 4. `lab.py porque <fala>` explica ov/score/camada
\prov{relações: parte\_de tiffany\_cabeca\_broca\_tiffany\_cabeca\_conversacional\_broca-so; usa colapso}
\prov{proveniência: tiffany\_cabeca\_broca.md \textperiodcentered{} web \textperiodcentered{} fato \textperiodcentered{} confiança 1.0 \textperiodcentered{} domínio referencia \textperiodcentered{} história 4 versões no jornal}
\prov{fusão de curadoria (texto idêntico): absorve pipeline\_conversacional --- o mesmo conteúdo por dois esquemas de endereço}

%CRISTAL {"arestas":[{"alvo":"tiffany_cabeca_broca_tiffany_cabeca_conversacional_broca-so","confianca":1.0,"epistemico":"fato","origem":"tiffany_cabeca_broca.md","rel":"parte_de"},{"alvo":"tiffany","confianca":0.95,"epistemico":"fato","origem":"wiki","rel":"confirma"},{"alvo":"painel_estelar","confianca":0.95,"epistemico":"fato","origem":"wiki","rel":"explica"},{"alvo":"lobo_frontal","confianca":0.95,"epistemico":"fato","origem":"wiki","rel":"explica"},{"alvo":"assistente","confianca":0.95,"epistemico":"fato","origem":"wiki","rel":"explica"},{"alvo":"colapso","confianca":0.95,"epistemico":"fato","origem":"wiki","rel":"usa"}],"confianca":1.0,"contraexemplos":[],"descricao":"| Tiffany | Nó | relação | |---------|-----|---------| | cabeça | tiffany | confirma | | motor Ψ | colapso | usa | | SO fractal | broca_os | parte_de | | rede associativa | hopfield | especializa | | domínio IA | inteligencia_artificial | especializa | | validação | codigo_neuronio_experimentos_tiffany | confirma | | três ciências | painel_estelar | explica | | turno A—B | lobo_frontal | explica | | interface | assistente | explica |","epistemico":"fato","exemplos":[],"id":"tiffany_cabeca_broca_ponte_broca-so","memoria":"semantica","meta":{"dominio":"referencia","fonte":"web","nivel":5},"origem":"tiffany_cabeca_broca.md","palavras_chave":[],"sinonimos":[],"tipo":"conceito","titulo":"Ponte Broca-SO"}
\section{Ponte Broca-SO}
| Tiffany | Nó | relação | |---------|-----|---------| | cabeça | tiffany | confirma | | motor \ensuremath{\Psi} | colapso | usa | | SO fractal | broca\_os | parte\_de | | rede associativa | hopfield | especializa | | domínio IA | inteligencia\_artificial | especializa | | validação | codigo\_neuronio\_experimentos\_tiffany | confirma | | três ciências | painel\_estelar | explica | | turno A—B | lobo\_frontal | explica | | interface | assistente | explica |
\prov{relações: parte\_de tiffany\_cabeca\_broca\_tiffany\_cabeca\_conversacional\_broca-so; confirma tiffany; explica painel\_estelar; explica lobo\_frontal; explica assistente; usa colapso}
\prov{proveniência: tiffany\_cabeca\_broca.md \textperiodcentered{} web \textperiodcentered{} fato \textperiodcentered{} confiança 1.0 \textperiodcentered{} domínio referencia \textperiodcentered{} história 8 versões no jornal}

%CRISTAL {"arestas":[{"alvo":"tiffany_cabeca_broca_tiffany_cabeca_conversacional_broca-so","confianca":1.0,"epistemico":"fato","origem":"tiffany_cabeca_broca.md","rel":"parte_de"}],"confianca":1.0,"contraexemplos":[],"descricao":"- **Paper:** `papers/Tiffany.tex`, `papers/computacao_fundamentos_broca.tex` - **Currículo:** fase IA — McCulloch-Pitts → Hopfield → **Tiffany** - **Código:** `neuronio/experimentos_tiffany.py`, `conversa/colapso.py`, `conversa/conhecimento/lab.py` - **Arquitectura:** `README_CRISTAL.md` — Tiffany conversa; Cristal administra tecido","epistemico":"fato","exemplos":[],"id":"tiffany_cabeca_broca_referencia","memoria":"semantica","meta":{"dominio":"referencia","fonte":"web","nivel":5},"origem":"tiffany_cabeca_broca.md","palavras_chave":[],"sinonimos":[],"tipo":"conceito","titulo":"Referência"}
\section{Referência}
- **Paper:** `papers/Tiffany.tex`, `papers/computacao\_fundamentos\_broca.tex` - **Currículo:** fase IA — McCulloch-Pitts \ensuremath{\to} Hopfield \ensuremath{\to} **Tiffany** - **Código:** `neuronio/experimentos\_tiffany.py`, `conversa/colapso.py`, `conversa/conhecimento/lab.py` - **Arquitectura:** `README\_CRISTAL.md` — Tiffany conversa; Cristal administra tecido
\prov{relações: parte\_de tiffany\_cabeca\_broca\_tiffany\_cabeca\_conversacional\_broca-so}
\prov{proveniência: tiffany\_cabeca\_broca.md \textperiodcentered{} web \textperiodcentered{} fato \textperiodcentered{} confiança 1.0 \textperiodcentered{} domínio referencia \textperiodcentered{} história 3 versões no jornal}

%CRISTAL {"fusao":[{"arestas":[{"alvo":"tiffany_cabeca_broca","confianca":0.95,"epistemico":"fato","origem":"wiki","rel":"parte_de"}],"confianca":1.0,"contraexemplos":[],"descricao":"> **Epistemologia:** fato de implementação (papers + código + currículo). > **No grafo:** `epistemico=fato`, confiança 0.95. > **Papel:** fechar hub Tiffany — exemplos, explicações, contraexemplos.","epistemico":"fato","exemplos":[],"id":"tiffany_cabeca_broca_tiffany_cabeca_conversacional_broca-so","memoria":"semantica","meta":{"dominio":"referencia","fonte":"web","nivel":5},"origem":"tiffany_cabeca_broca.md","palavras_chave":[],"sinonimos":[],"tipo":"conceito","titulo":"Tiffany — cabeça conversacional Broca-SO"},{"arestas":[{"alvo":"tiffany_cabeca_broca","confianca":0.95,"epistemico":"fato","origem":"wiki","rel":"parte_de"}],"confianca":1.0,"contraexemplos":[],"descricao":"> **Epistemologia:** fato de implementação (papers + código + currículo). > **No grafo:** `epistemico=fato`, confiança 0.95. > **Papel:** fechar hub Tiffany — exemplos, explicações, contraexemplos.","epistemico":"fato","exemplos":[],"id":"tiffany_cabeca_conversacional_broca-so","memoria":"semantica","meta":{"dominio":"referencia","nivel":5},"origem":"tiffany_cabeca_broca.md","palavras_chave":[],"sinonimos":[],"tipo":"conceito","titulo":"Tiffany — cabeça conversacional Broca-SO"}],"id":"tiffany_cabeca_broca_tiffany_cabeca_conversacional_broca-so","tipo":"conceito"}
\section{Tiffany — cabeça conversacional Broca-SO}
> **Epistemologia:** fato de implementação (papers + código + currículo). > **No grafo:** `epistemico=fato`, confiança 0.95. > **Papel:** fechar hub Tiffany — exemplos, explicações, contraexemplos.
\prov{relações: parte\_de tiffany\_cabeca\_broca}
\prov{proveniência: tiffany\_cabeca\_broca.md \textperiodcentered{} web \textperiodcentered{} fato \textperiodcentered{} confiança 1.0 \textperiodcentered{} domínio referencia \textperiodcentered{} história 3 versões no jornal}
\prov{fusão de curadoria (texto idêntico): absorve tiffany\_cabeca\_conversacional\_broca-so --- o mesmo conteúdo por dois esquemas de endereço}

%CRISTAL {"arestas":[{"alvo":"estatisticas","confianca":1.0,"epistemico":"fato","origem":"DEGENERADOS.md","rel":"parte_de"},{"alvo":"gato","confianca":1.0,"epistemico":"fato","origem":"manual","rel":"parte_de"},{"alvo":"word","confianca":0.5,"epistemico":"fato","origem":"wiki","rel":"relaciona"},{"alvo":"virtualizacao","confianca":0.5,"epistemico":"fato","origem":"wiki","rel":"relaciona"},{"alvo":"vita_contemplativa","confianca":0.5,"epistemico":"fato","origem":"wiki","rel":"relaciona"}],"confianca":1.0,"contraexemplos":[],"descricao":"``` 31250  0,0|0,0 22  0,0|2,1 15  0,0|2,0 5  0,0|5,2 3  0,0|1,1 3  2,2|0,0 3  0,0|4,2 2  0,0|1,0 2  0,0|2,2 2  2,1|0,0 ```","epistemico":"fato","exemplos":[],"id":"top_assinaturas","memoria":"semantica","meta":{"dominio":"manual","nivel":6,"verificacao":"slug idêntico — enriquecer, não duplicar"},"origem":"DEGENERADOS.md","palavras_chave":[],"sinonimos":[],"tipo":"conceito","titulo":"Top assinaturas"}
\section{Top assinaturas}
``` 31250  0,0|0,0 22  0,0|2,1 15  0,0|2,0 5  0,0|5,2 3  0,0|1,1 3  2,2|0,0 3  0,0|4,2 2  0,0|1,0 2  0,0|2,2 2  2,1|0,0 ```
\prov{relações: parte\_de estatisticas; parte\_de gato; relaciona word; relaciona virtualizacao; relaciona vita\_contemplativa}
\prov{proveniência: DEGENERADOS.md \textperiodcentered{} fato \textperiodcentered{} confiança 1.0 \textperiodcentered{} domínio manual \textperiodcentered{} história 13 versões no jornal}

%CRISTAL {"arestas":[{"alvo":"estatisticas","confianca":1.0,"epistemico":"fato","origem":"DEGENERADOS.md","rel":"parte_de"},{"alvo":"broca_os","confianca":1.0,"epistemico":"fato","origem":"manual","rel":"parte_de"},{"alvo":"word","confianca":0.5,"epistemico":"fato","origem":"wiki","rel":"relaciona"}],"confianca":1.0,"contraexemplos":[],"descricao":"- `—`: 985,600 - `JNZ`: 9,909 - `JMP`: 3,836 - `OR`: 2,749 - `JZ`: 2,068 - `SUB`: 2,055 - `AND`: 1,994 - `LOAD`: 1,948 - `ADD`: 1,888 - `XOR`: 1,660","epistemico":"fato","exemplos":[],"id":"top_opcodes","memoria":"semantica","meta":{"dominio":"manual","nivel":6,"verificacao":"slug idêntico — enriquecer, não duplicar"},"origem":"DEGENERADOS.md","palavras_chave":[],"sinonimos":[],"tipo":"conceito","titulo":"Top opcodes"}
\section{Top opcodes}
- `—`: 985,600 - `JNZ`: 9,909 - `JMP`: 3,836 - `OR`: 2,749 - `JZ`: 2,068 - `SUB`: 2,055 - `AND`: 1,994 - `LOAD`: 1,948 - `ADD`: 1,888 - `XOR`: 1,660
\prov{relações: parte\_de estatisticas; parte\_de broca\_os; relaciona word}
\prov{proveniência: DEGENERADOS.md \textperiodcentered{} fato \textperiodcentered{} confiança 1.0 \textperiodcentered{} domínio manual \textperiodcentered{} história 11 versões no jornal}

%CRISTAL {"arestas":[{"alvo":"colapso_koch_tiffany","confianca":1.0,"epistemico":"fato","origem":"wiki","rel":"relaciona"},{"alvo":"maquina_algebrica","confianca":1.0,"epistemico":"fato","origem":"wiki","rel":"relaciona"}],"confianca":1.0,"contraexemplos":[],"descricao":"A Máquina mede um suporte (atrator, um Cantor: nowhere dense, Lebesgue 0, totalmente desconexo) cuja topologia é trivial mas cuja medida (dim_H) é rica: o que a topologia não vê, a dimensão vê. A dimensão vem do repouso admitido — admitir a ramificação no vazio ('00') produz λ>1 e dim_H>0.","epistemico":"fato","exemplos":[],"id":"topologia_vazia_do_suporte","memoria":"semantica","meta":{"autor":"Aarão/broca-so","dominio":"broca_repos","fonte":""},"origem":"wiki","palavras_chave":[],"sinonimos":[],"tipo":"conceito","titulo":"Topologia Vazia, Medida Rica"}
\section{Topologia Vazia, Medida Rica}
A Máquina mede um suporte (atrator, um Cantor: nowhere dense, Lebesgue 0, totalmente desconexo) cuja topologia é trivial mas cuja medida (dim\_H) é rica: o que a topologia não vê, a dimensão vê. A dimensão vem do repouso admitido — admitir a ramificação no vazio ('00') produz \ensuremath{\lambda}>1 e dim\_H>0.
\prov{relações: relaciona colapso\_koch\_tiffany; relaciona maquina\_algebrica}
\prov{proveniência: wiki \textperiodcentered{} fato \textperiodcentered{} confiança 1.0 \textperiodcentered{} domínio broca\_repos \textperiodcentered{} história 3 versões no jornal}

%CRISTAL {"arestas":[{"alvo":"cuspide_pisot","confianca":1.0,"epistemico":"fato","origem":"wiki","rel":"usa"},{"alvo":"razao_aurea","confianca":1.0,"epistemico":"fato","origem":"wiki","rel":"usa"},{"alvo":"sequencia_pell","confianca":1.0,"epistemico":"fato","origem":"wiki","rel":"relaciona"},{"alvo":"conjugados_galois_parabola","confianca":1.0,"epistemico":"fato","origem":"wiki","rel":"usa"}],"confianca":1.0,"contraexemplos":[],"descricao":"Teorema (verificado em 6 andares): torre ℚ=K₀⊂K₁⊂… com θₙ a cúspide metálica de θₙ−₁ (θ₁=φ); o de cima abre o de baixo por θₙ−₁=θₙ−1/θₙ=2sinh(log θₙ), com [Kₙ:ℚ]=2ⁿ. 'O irracional de um andar é o racional do seguinte'; a descida x↦x−1/x só estabiliza com a órbita de Galois completa (o osso exige a carne).","epistemico":"fato","exemplos":[],"id":"torre_corpos_sinh","memoria":"semantica","meta":{"autor":"Lopes/Aarão","dominio":"hiper","fonte":""},"origem":"wiki","palavras_chave":[],"sinonimos":[],"tipo":"conceito","titulo":"A Torre de Corpos e a Abertura Hiperbólica"}
\section{A Torre de Corpos e a Abertura Hiperbólica}
Teorema (verificado em 6 andares): torre \ensuremath{\mathbb{Q}}=K\ensuremath{_{0}}\ensuremath{\subset}K\ensuremath{_{1}}\ensuremath{\subset}… com \ensuremath{\theta}\ensuremath{_{n}} a cúspide metálica de \ensuremath{\theta}\ensuremath{_{n}}\ensuremath{-}\ensuremath{_{1}} (\ensuremath{\theta}\ensuremath{_{1}}=\ensuremath{\varphi}); o de cima abre o de baixo por \ensuremath{\theta}\ensuremath{_{n}}\ensuremath{-}\ensuremath{_{1}}=\ensuremath{\theta}\ensuremath{_{n}}\ensuremath{-}1/\ensuremath{\theta}\ensuremath{_{n}}=2sinh(log \ensuremath{\theta}\ensuremath{_{n}}), com [K\ensuremath{_{n}}:\ensuremath{\mathbb{Q}}]=2\ensuremath{^{n}}. 'O irracional de um andar é o racional do seguinte'; a descida x\ensuremath{\mapsto}x\ensuremath{-}1/x só estabiliza com a órbita de Galois completa (o osso exige a carne).
\prov{relações: usa cuspide\_pisot; usa razao\_aurea; relaciona sequencia\_pell; usa conjugados\_galois\_parabola}
\prov{proveniência: wiki \textperiodcentered{} fato \textperiodcentered{} confiança 1.0 \textperiodcentered{} domínio hiper \textperiodcentered{} história 7 versões no jornal}

%CRISTAL {"arestas":[{"alvo":"analise","confianca":1.0,"epistemico":"fato","origem":"paper","rel":"parte_de"},{"alvo":"a_engenharia_de_ouro_pa_pgs_agm_fc_reta","confianca":1.0,"epistemico":"fato","origem":"paper","rel":"confirma"},{"alvo":"reta_de_ouro","confianca":1.0,"epistemico":"fato","origem":"wiki","rel":"usa"},{"alvo":"regua_de_gentil","confianca":1.0,"epistemico":"fato","origem":"wiki","rel":"relaciona"},{"alvo":"transformada_estelar","confianca":1.0,"epistemico":"fato","origem":"wiki","rel":"relaciona"}],"confianca":1.0,"contraexemplos":[],"descricao":"Transformada de Mellin restrita à linha crítica (s=iω), = Fourier no logaritmo: 𝓖[x](ω)=∫x(t)t−iωdt/t. Isometria unitária invertível (Parseval dourado) que leva os dados ao 'lado relativístico' (multiplicativo/logarítmico), onde a reta de ouro Ω=φ¹/φ é reta de fato, devolvendo amplitude+fase.","epistemico":"fato","exemplos":[],"id":"transformada_dourada","memoria":"semantica","meta":{"autor":"Lopes/Aarão","dominio":"hiper","fonte":""},"origem":"wiki","palavras_chave":[],"sinonimos":[],"tipo":"conceito","titulo":"A Transformada Dourada 𝓖=𝓕∘Λ"}
\section{A Transformada Dourada \ensuremath{\mathcal{G}}=\ensuremath{\mathcal{F}}\ensuremath{\circ}\ensuremath{\Lambda}}
Transformada de Mellin restrita à linha crítica (s=i\ensuremath{\omega}), = Fourier no logaritmo: \ensuremath{\mathcal{G}}[x](\ensuremath{\omega})=\ensuremath{\int}x(t)t\ensuremath{-}i\ensuremath{\omega}dt/t. Isometria unitária invertível (Parseval dourado) que leva os dados ao 'lado relativístico' (multiplicativo/logarítmico), onde a reta de ouro \ensuremath{\Omega}=\ensuremath{\varphi}\ensuremath{^{1}}/\ensuremath{\varphi} é reta de fato, devolvendo amplitude+fase.
\prov{relações: parte\_de analise; confirma a\_engenharia\_de\_ouro\_pa\_pgs\_agm\_fc\_reta; usa reta\_de\_ouro; relaciona regua\_de\_gentil; relaciona transformada\_estelar}
\prov{proveniência: wiki \textperiodcentered{} fato \textperiodcentered{} confiança 1.0 \textperiodcentered{} domínio hiper \textperiodcentered{} história 90 versões no jornal}

%CRISTAL {"arestas":[{"alvo":"transformada_dourada","confianca":1.0,"epistemico":"fato","origem":"wiki","rel":"generaliza"},{"alvo":"torre_hurwitz","confianca":1.0,"epistemico":"fato","origem":"wiki","rel":"usa"},{"alvo":"fibracao_de_hopf","confianca":1.0,"epistemico":"fato","origem":"wiki","rel":"relaciona"},{"alvo":"transformada_metalica","confianca":1.0,"epistemico":"fato","origem":"wiki","rel":"especializa"}],"confianca":1.0,"contraexemplos":[],"descricao":"Teorema: a elevação 𝓖A a A∈ℝ,ℂ,ℍ,𝕆 (dourada 1D no módulo ⊗ harmônicos esféricos na direção) — a fatoração ortogonal módulo/direção e a linearização ln ρ existem se e somente se a norma é multiplicativa, i.e. exatamente nas 4 álgebras de divisão normadas.","epistemico":"fato","exemplos":[],"id":"transformada_dourada_hurwitz","memoria":"semantica","meta":{"autor":"Lopes/Aarão","dominio":"hiper","fonte":""},"origem":"wiki","palavras_chave":[],"sinonimos":[],"tipo":"conceito","titulo":"A Transformada Dourada Multidimensional (só sobre Hurwitz)"}
\section{A Transformada Dourada Multidimensional (só sobre Hurwitz)}
Teorema: a elevação \ensuremath{\mathcal{G}}A a A\ensuremath{\in}\ensuremath{\mathbb{R}},\ensuremath{\mathbb{C}},\ensuremath{\mathbb{H}},\ensuremath{\mathbb{O}} (dourada 1D no módulo \ensuremath{\otimes} harmônicos esféricos na direção) — a fatoração ortogonal módulo/direção e a linearização ln \ensuremath{\rho} existem se e somente se a norma é multiplicativa, i.e. exatamente nas 4 álgebras de divisão normadas.
\prov{relações: generaliza transformada\_dourada; usa torre\_hurwitz; relaciona fibracao\_de\_hopf; especializa transformada\_metalica}
\prov{proveniência: wiki \textperiodcentered{} fato \textperiodcentered{} confiança 1.0 \textperiodcentered{} domínio hiper \textperiodcentered{} história 7 versões no jornal}

%CRISTAL {"arestas":[{"alvo":"matematica_estelar","confianca":1.0,"epistemico":"fato","origem":"wiki","rel":"parte_de"},{"alvo":"invariante_estelar","confianca":1.0,"epistemico":"fato","origem":"wiki","rel":"usa"},{"alvo":"maquina_evolutiva","confianca":1.0,"epistemico":"fato","origem":"wiki","rel":"relaciona"}],"confianca":1.0,"contraexemplos":[],"descricao":"Fluxo discreto cₖ+₁=cₖ¹−s·Ωs·Δₖ, aₖ+₁=1−c², com 'costura relativística' que injeta o associador — realiza a passagem do regime aditivo-plano ao multiplicativo-relativístico (ativa curvatura quando s>0). O lado relativístico vivo da geometria de ouro.","epistemico":"fato","exemplos":[],"id":"transformada_estelar","memoria":"semantica","meta":{"autor":"Lopes/Aarão","dominio":"hiper","fonte":""},"origem":"wiki","palavras_chave":[],"sinonimos":[],"tipo":"conceito","titulo":"Transformada Estelar"}
\section{Transformada Estelar}
Fluxo discreto c\ensuremath{_{k}}+\ensuremath{_{1}}=c\ensuremath{_{k}}\ensuremath{^{1}}\ensuremath{-}s\textperiodcentered \ensuremath{\Omega}s\textperiodcentered \ensuremath{\Delta}\ensuremath{_{k}}, a\ensuremath{_{k}}+\ensuremath{_{1}}=1\ensuremath{-}c\ensuremath{^{2}}, com 'costura relativística' que injeta o associador — realiza a passagem do regime aditivo-plano ao multiplicativo-relativístico (ativa curvatura quando s>0). O lado relativístico vivo da geometria de ouro.
\prov{relações: parte\_de matematica\_estelar; usa invariante\_estelar; relaciona maquina\_evolutiva}
\prov{proveniência: wiki \textperiodcentered{} fato \textperiodcentered{} confiança 1.0 \textperiodcentered{} domínio hiper \textperiodcentered{} história 6 versões no jornal}

%CRISTAL {"arestas":[{"alvo":"1_modelo_formal","confianca":1.0,"epistemico":"fato","origem":"EXPRESSIVIDADE.md","rel":"parte_de"},{"alvo":"broca_os","confianca":1.0,"epistemico":"fato","origem":"manual","rel":"parte_de"},{"alvo":"utilitarismo","confianca":0.5,"epistemico":"fato","origem":"wiki","rel":"relaciona"},{"alvo":"virtualizacao","confianca":0.5,"epistemico":"fato","origem":"wiki","rel":"relaciona"},{"alvo":"word","confianca":0.5,"epistemico":"fato","origem":"wiki","rel":"relaciona"},{"alvo":"universidade_curriculo_em_camadas","confianca":0.5,"epistemico":"fato","origem":"wiki","rel":"relaciona"}],"confianca":1.0,"contraexemplos":[],"descricao":"Regime: **abrupta** · salto no fecho: 0.3982 · gradiente médio: 0.0492.","epistemico":"fato","exemplos":[],"id":"transicao","memoria":"semantica","meta":{"dominio":"manual","duplicata_sugerida":[],"nivel":6,"verificacao":"parte de conceito mais amplo 'transicao'"},"origem":"POW_METAL_RESSONANCIA_VAL.md","palavras_chave":[],"sinonimos":[],"tipo":"conceito","titulo":"Transição"}
\section{Transição}
Regime: **abrupta** \textperiodcentered  salto no fecho: 0.3982 \textperiodcentered  gradiente médio: 0.0492.
\prov{relações: parte\_de 1\_modelo\_formal; parte\_de broca\_os; relaciona utilitarismo; relaciona virtualizacao; relaciona word; relaciona universidade\_curriculo\_em\_camadas}
\prov{proveniência: POW\_METAL\_RESSONANCIA\_VAL.md \textperiodcentered{} fato \textperiodcentered{} confiança 1.0 \textperiodcentered{} domínio manual \textperiodcentered{} história 16 versões no jornal}

%CRISTAL {"arestas":[{"alvo":"alternatividade_dimensoes_criticas","confianca":1.0,"epistemico":"fato","origem":"wiki","rel":"usa"},{"alvo":"associador","confianca":1.0,"epistemico":"fato","origem":"wiki","rel":"especializa"}],"confianca":1.0,"contraexemplos":[],"descricao":"A_s(ψ,χ)=[ψ̄,ψ,χ]ₛ, obstrução à independência de ordem de avaliação do produto; interpola como (1−s²)A₀+s²R, cuja magnitude |1−s²| funciona como 'massa de quebra de SUSY' algébrica.","epistemico":"fato","exemplos":[],"id":"tri_associador_espinorial","memoria":"semantica","meta":{"autor":"Lopes/Aarão","dominio":"hiper","fonte":""},"origem":"wiki","palavras_chave":[],"sinonimos":[],"tipo":"conceito","titulo":"Tri-Associador Espinorial"}
\section{Tri-Associador Espinorial}
A\_s(\ensuremath{\psi},\ensuremath{\chi})=[\ensuremath{\psi}\ensuremath{},\ensuremath{\psi},\ensuremath{\chi}]\ensuremath{_{s}}, obstrução à independência de ordem de avaliação do produto; interpola como (1\ensuremath{-}s\ensuremath{^{2}})A\ensuremath{_{0}}+s\ensuremath{^{2}}R, cuja magnitude |1\ensuremath{-}s\ensuremath{^{2}}| funciona como 'massa de quebra de SUSY' algébrica.
\prov{relações: usa alternatividade\_dimensoes\_criticas; especializa associador}
\prov{proveniência: wiki \textperiodcentered{} fato \textperiodcentered{} confiança 1.0 \textperiodcentered{} domínio hiper \textperiodcentered{} história 4 versões no jornal}

%CRISTAL {"arestas":[{"alvo":"cristalchain","confianca":1.0,"epistemico":"fato","origem":"wiki","rel":"eh_um"},{"alvo":"espectro_irredutivel_floco","confianca":1.0,"epistemico":"fato","origem":"wiki","rel":"usa"},{"alvo":"corpo_estelar","confianca":1.0,"epistemico":"fato","origem":"wiki","rel":"usa"},{"alvo":"numeros_de_pisot","confianca":1.0,"epistemico":"fato","origem":"wiki","rel":"relaciona"},{"alvo":"area_gemologia","confianca":1.0,"epistemico":"fato","origem":"wiki","rel":"parte_de"}],"confianca":1.0,"contraexemplos":[],"descricao":"A pedra-âncora da CristalChain: a Turmalina Paraíba (elbaíta cuprífera). Descoberta em 1989 por Heitor Dimas Barbosa na mina da Batalha (Paraíba), tem o azul-neon dado pelo COBRE (Cu²⁺) — a primeira turmalina colorida por cobre, o metal que a máquina lê. Elbaíta Na(Li,Al)₃Al₆(BO₃)₃Si₆O₁₈(OH)₄; trigonal, R3m, simetria 3-FOLD. A ESTRELA da pedra é o azul Paraíba (~490–509 nm), a janela de transmissão entre o violeta (Mn) e o vermelho (Cu); a simetria de 3 eixos dá 3 modos. Uma pedra brasileira, de um metal, de três eixos — a âncora concreta da cadeia. Verif. em turmalina.py (6/6).","epistemico":"fato","exemplos":[],"id":"turmalina_paraiba_ancora","memoria":"semantica","meta":{"dominio":"referencia","fonte":"turmalina paraíba"},"origem":"wiki","palavras_chave":[],"sinonimos":[],"tipo":"conceito","titulo":"A Turmalina Paraíba (a âncora da CristalChain)"}
\section{A Turmalina Paraíba (a âncora da CristalChain)}
A pedra-âncora da CristalChain: a Turmalina Paraíba (elbaíta cuprífera). Descoberta em 1989 por Heitor Dimas Barbosa na mina da Batalha (Paraíba), tem o azul-neon dado pelo COBRE (Cu\ensuremath{^{2+}}) — a primeira turmalina colorida por cobre, o metal que a máquina lê. Elbaíta Na(Li,Al)\ensuremath{_{3}}Al\ensuremath{_{6}}(BO\ensuremath{_{3}})\ensuremath{_{3}}Si\ensuremath{_{6}}O\ensuremath{_{18}}(OH)\ensuremath{_{4}}; trigonal, R3m, simetria 3-FOLD. A ESTRELA da pedra é o azul Paraíba (\~{}490–509 nm), a janela de transmissão entre o violeta (Mn) e o vermelho (Cu); a simetria de 3 eixos dá 3 modos. Uma pedra brasileira, de um metal, de três eixos — a âncora concreta da cadeia. Verif. em turmalina.py (6/6).
\prov{relações: eh\_um cristalchain; usa espectro\_irredutivel\_floco; usa corpo\_estelar; relaciona numeros\_de\_pisot; parte\_de area\_gemologia}
\prov{proveniência: wiki \textperiodcentered{} turmalina paraíba \textperiodcentered{} fato \textperiodcentered{} confiança 1.0 \textperiodcentered{} domínio referencia \textperiodcentered{} história 13 versões no jornal}

%CRISTAL {"arestas":[{"alvo":"readme","confianca":1.0,"epistemico":"fato","origem":"manual","rel":"parte_de"},{"alvo":"vida-morte","confianca":0.5,"epistemico":"fato","origem":"wiki","rel":"relaciona"},{"alvo":"virtualizacao","confianca":0.5,"epistemico":"fato","origem":"wiki","rel":"relaciona"},{"alvo":"vetores","confianca":0.5,"epistemico":"fato","origem":"wiki","rel":"relaciona"}],"confianca":1.0,"contraexemplos":[],"descricao":"Não reproduzimos CPU x86 nem RISC-V. Descobrimos a **menor ULA** cuja primitiva é o operador do neurônio.","epistemico":"fato","exemplos":[],"id":"ula_da_broca_arquitetura_minima","memoria":"semantica","meta":{"dominio":"manual","nivel":6,"verificacao":"slug idêntico — enriquecer, não duplicar"},"origem":"README.md","palavras_chave":[],"sinonimos":[],"tipo":"conceito","titulo":"ULA da Broca — Arquitetura Mínima"}
\section{ULA da Broca — Arquitetura Mínima}
Não reproduzimos CPU x86 nem RISC-V. Descobrimos a **menor ULA** cuja primitiva é o operador do neurônio.
\prov{relações: parte\_de readme; relaciona vida-morte; relaciona virtualizacao; relaciona vetores}
\prov{proveniência: README.md \textperiodcentered{} fato \textperiodcentered{} confiança 1.0 \textperiodcentered{} domínio manual \textperiodcentered{} história 7 versões no jornal}

%CRISTAL {"arestas":[{"alvo":"pow_lemniscata_koch_corpo_de_gauss","confianca":1.0,"epistemico":"fato","origem":"POW_METAL.md","rel":"parte_de"}],"confianca":1.0,"contraexemplos":[],"descricao":"```bash make pow-metal ./pow_metal_miner -n 5000000 -F koch      # atrator Koch == ref ./pow_metal_miner -n 5000000 -F ref       # word == ref ./pow_metal_miner -n 500000 -F ouro-ref ./pow_metal_miner -n 500000 -F project ```","epistemico":"fato","exemplos":[],"id":"uso","memoria":"semantica","meta":{"dominio":"manual","nivel":6,"verificacao":"slug idêntico — enriquecer, não duplicar"},"origem":"POW_METAL.md","palavras_chave":[],"sinonimos":[],"tipo":"conceito","titulo":"Uso"}
\section{Uso}
```bash make pow-metal ./pow\_metal\_miner -n 5000000 -F koch      \# atrator Koch == ref ./pow\_metal\_miner -n 5000000 -F ref       \# word == ref ./pow\_metal\_miner -n 500000 -F ouro-ref ./pow\_metal\_miner -n 500000 -F project ```
\prov{relações: parte\_de pow\_lemniscata\_koch\_corpo\_de\_gauss}
\prov{proveniência: POW\_METAL.md \textperiodcentered{} fato \textperiodcentered{} confiança 1.0 \textperiodcentered{} domínio manual \textperiodcentered{} história 4 versões no jornal}

%CRISTAL {"arestas":[{"alvo":"pow_metal_ressonancia_val","confianca":1.0,"epistemico":"fato","origem":"manual","rel":"parte_de"}],"confianca":1.0,"contraexemplos":[],"descricao":"> **H:** o fecho ocorre quando os modos entram em ressonância. > Ressonância **mede** proximidade ao estado coerente — não causa o fecho.","epistemico":"fato","exemplos":[],"id":"validacao_hipotese_h_ressonancia","memoria":"semantica","meta":{"dominio":"manual","nivel":6,"verificacao":"slug idêntico — enriquecer, não duplicar"},"origem":"POW_METAL_RESSONANCIA_VAL.md","palavras_chave":[],"sinonimos":[],"tipo":"conceito","titulo":"Validação — Hipótese H (Ressonância)"}
\section{Validação — Hipótese H (Ressonância)}
> **H:** o fecho ocorre quando os modos entram em ressonância. > Ressonância **mede** proximidade ao estado coerente — não causa o fecho.
\prov{relações: parte\_de pow\_metal\_ressonancia\_val}
\prov{proveniência: POW\_METAL\_RESSONANCIA\_VAL.md \textperiodcentered{} fato \textperiodcentered{} confiança 1.0 \textperiodcentered{} domínio manual \textperiodcentered{} história 4 versões no jornal}

%CRISTAL {"arestas":[{"alvo":"assistente","confianca":1.0,"epistemico":"fato","origem":"paper","rel":"parte_de"},{"alvo":"word","confianca":0.5,"epistemico":"fato","origem":"wiki","rel":"relaciona"}],"confianca":1.0,"contraexemplos":[],"descricao":"Geometria dos degenerados = **estrutura de zeros e clips negativos**, dominada por assinaturas simétricas `0,0|0,0` e variações de componente não-positiva.","epistemico":"fato","exemplos":[],"id":"veredito_e_honestidade","memoria":"semantica","meta":{"dominio":"manual","duplicata_sugerida":[],"nivel":6,"verificacao":"parte de conceito mais amplo 'veredito_e_honestidade'"},"origem":"DEGENERADOS.md","palavras_chave":[],"sinonimos":[],"tipo":"conceito","titulo":"Veredito e honestidade"}
\section{Veredito e honestidade}
Geometria dos degenerados = **estrutura de zeros e clips negativos**, dominada por assinaturas simétricas `0,0|0,0` e variações de componente não-positiva.
\prov{relações: parte\_de assistente; relaciona word}
\prov{proveniência: DEGENERADOS.md \textperiodcentered{} fato \textperiodcentered{} confiança 1.0 \textperiodcentered{} domínio manual \textperiodcentered{} história 10 versões no jornal}

%CRISTAL {"arestas":[{"alvo":"setor_algebrico_descontinuidades_no_fecho","confianca":1.0,"epistemico":"fato","origem":"POW_METAL_SETOR.md","rel":"parte_de"},{"alvo":"vida-morte","confianca":0.5,"epistemico":"fato","origem":"wiki","rel":"relaciona"},{"alvo":"virtualizacao","confianca":0.5,"epistemico":"fato","origem":"wiki","rel":"relaciona"},{"alvo":"word","confianca":0.5,"epistemico":"fato","origem":"wiki","rel":"relaciona"},{"alvo":"vontade_de_poder","confianca":0.5,"epistemico":"fato","origem":"wiki","rel":"relaciona"},{"alvo":"while","confianca":0.5,"epistemico":"fato","origem":"wiki","rel":"relaciona"}],"confianca":1.0,"contraexemplos":[],"descricao":"`SETOR INALTERADO; TRANSICAO ATRATOR`","epistemico":"fato","exemplos":[],"id":"veredito_experimental","memoria":"semantica","meta":{"dominio":"manual","nivel":6,"verificacao":"slug idêntico — enriquecer, não duplicar"},"origem":"POW_METAL_SETOR.md","palavras_chave":[],"sinonimos":[],"tipo":"conceito","titulo":"Veredito experimental"}
\section{Veredito experimental}
`SETOR INALTERADO; TRANSICAO ATRATOR`
\prov{relações: parte\_de setor\_algebrico\_descontinuidades\_no\_fecho; relaciona vida-morte; relaciona virtualizacao; relaciona word; relaciona vontade\_de\_poder; relaciona while}
\prov{proveniência: POW\_METAL\_SETOR.md \textperiodcentered{} fato \textperiodcentered{} confiança 1.0 \textperiodcentered{} domínio manual \textperiodcentered{} história 9 versões no jornal}

%CRISTAL {"arestas":[{"alvo":"morte_matematica","confianca":1.0,"epistemico":"fato","origem":"wiki","rel":"contradiz"},{"alvo":"corpo_dual","confianca":1.0,"epistemico":"fato","origem":"wiki","rel":"usa"}],"confianca":1.0,"contraexemplos":[],"descricao":"O dual positivo: um objeto está vivo quando sua identidade é uma dinâmica geradora não trivial (Ut=etL), da qual os estados são manifestações que ressoam (não se igualam) — com frequência fundamental comum, eixo Re s=½ e reversibilidade (Koopman unitário, provado). Custo O(−1): mais dados → mais potência (ganho convexo) — o eixo c² da parábola.","epistemico":"fato","exemplos":[],"id":"vida_matematica","memoria":"semantica","meta":{"autor":"Aarão/broca-so","dominio":"broca_repos","fonte":""},"origem":"wiki","palavras_chave":[],"sinonimos":[],"tipo":"conceito","titulo":"A Vida Matemática (e o custo O(−1))"}
\section{A Vida Matemática (e o custo O(\ensuremath{-}1))}
O dual positivo: um objeto está vivo quando sua identidade é uma dinâmica geradora não trivial (Ut=etL), da qual os estados são manifestações que ressoam (não se igualam) — com frequência fundamental comum, eixo Re s=\ensuremath{\tfrac12} e reversibilidade (Koopman unitário, provado). Custo O(\ensuremath{-}1): mais dados \ensuremath{\to} mais potência (ganho convexo) — o eixo c\ensuremath{^{2}} da parábola.
\prov{relações: contradiz morte\_matematica; usa corpo\_dual}
\prov{proveniência: wiki \textperiodcentered{} fato \textperiodcentered{} confiança 1.0 \textperiodcentered{} domínio broca\_repos \textperiodcentered{} história 4 versões no jornal}

%CRISTAL {"arestas":[{"alvo":"colapso","confianca":0.95,"epistemico":"fato","origem":"paper","rel":"parte_de"}],"confianca":1.0,"contraexemplos":[],"descricao":"| Ficheiro | Conteúdo | |----------|----------| | `atlas/ruido_erro.svg` | crescimento Δ vs ε | | `atlas/ruido_horizonte.svg` | Δ por metal | | `atlas/ruido_trajetoria.svg` | Δ por passo | | `atlas/ruido_parabola.svg` | ideal vs perturbado no plano calor×fase | | `atlas/ruido_desloc.svg` | deslocamento médio vs ε | | `atlas/ruido_trajetorias.csv` | passos exp5 completos |","epistemico":"fato","exemplos":[],"id":"visualizacao","memoria":"semantica","meta":{"dominio":"manual","nivel":6,"verificacao":"slug idêntico — enriquecer, não duplicar"},"origem":"RUÍDO.md","palavras_chave":[],"sinonimos":[],"tipo":"conceito","titulo":"Visualização"}
\section{Visualização}
| Ficheiro | Conteúdo | |----------|----------| | `atlas/ruido\_erro.svg` | crescimento \ensuremath{\Delta} vs \ensuremath{\varepsilon} | | `atlas/ruido\_horizonte.svg` | \ensuremath{\Delta} por metal | | `atlas/ruido\_trajetoria.svg` | \ensuremath{\Delta} por passo | | `atlas/ruido\_parabola.svg` | ideal vs perturbado no plano calor$\times$fase | | `atlas/ruido\_desloc.svg` | deslocamento médio vs \ensuremath{\varepsilon} | | `atlas/ruido\_trajetorias.csv` | passos exp5 completos |
\prov{relações: parte\_de colapso}
\prov{proveniência: RUÍDO.md \textperiodcentered{} fato \textperiodcentered{} confiança 1.0 \textperiodcentered{} domínio manual \textperiodcentered{} história 5 versões no jornal}

%CRISTAL {"arestas":[{"alvo":"koch_atlas_experimento_de_parametrizacao","confianca":1.0,"epistemico":"fato","origem":"KOCH_ATLAS.md","rel":"parte_de"}],"confianca":1.0,"contraexemplos":[],"descricao":"| Figura | Descrição | |--------|-----------| | `atlas/koch_atlas/trajetoria/parametro.svg` | Nonce → parâmetro | | `atlas/koch_atlas/trajetoria/endereco.svg` | Nonce → endereço fractal | | `atlas/koch_atlas/trajetoria/continuidade.svg` | Δ atlas entre nonces consecutivos | | `atlas/koch_atlas/distribuicoes/ramos.svg` | Histograma de ramos |","epistemico":"fato","exemplos":[],"id":"visualizacoes","memoria":"semantica","meta":{"dominio":"manual","nivel":6,"verificacao":"slug idêntico — enriquecer, não duplicar"},"origem":"KOCH_ATLAS.md","palavras_chave":[],"sinonimos":[],"tipo":"conceito","titulo":"Visualizações"}
\section{Visualizações}
| Figura | Descrição | |--------|-----------| | `atlas/koch\_atlas/trajetoria/parametro.svg` | Nonce \ensuremath{\to} parâmetro | | `atlas/koch\_atlas/trajetoria/endereco.svg` | Nonce \ensuremath{\to} endereço fractal | | `atlas/koch\_atlas/trajetoria/continuidade.svg` | \ensuremath{\Delta} atlas entre nonces consecutivos | | `atlas/koch\_atlas/distribuicoes/ramos.svg` | Histograma de ramos |
\prov{relações: parte\_de koch\_atlas\_experimento\_de\_parametrizacao}
\prov{proveniência: KOCH\_ATLAS.md \textperiodcentered{} fato \textperiodcentered{} confiança 1.0 \textperiodcentered{} domínio manual \textperiodcentered{} história 6 versões no jornal}

%CRISTAL {"arestas":[{"alvo":"gato","confianca":1.0,"epistemico":"fato","origem":"wiki","rel":"relaciona"},{"alvo":"lei_da_parabola_e_dinamica_de_anosov","confianca":1.0,"epistemico":"fato","origem":"wiki","rel":"relaciona"},{"alvo":"numeros_de_pisot","confianca":1.0,"epistemico":"fato","origem":"wiki","rel":"relaciona"},{"alvo":"numero_de_salem","confianca":1.0,"epistemico":"fato","origem":"wiki","rel":"relaciona"},{"alvo":"entropia_topologica","confianca":1.0,"epistemico":"fato","origem":"wiki","rel":"relaciona"},{"alvo":"flip_e_rotacao_de_wick","confianca":1.0,"epistemico":"fato","origem":"wiki","rel":"relaciona"},{"alvo":"espectro_irredutivel_floco","confianca":1.0,"epistemico":"fato","origem":"wiki","rel":"relaciona"},{"alvo":"corpo_estelar","confianca":1.0,"epistemico":"fato","origem":"wiki","rel":"relaciona"},{"alvo":"mecanismo_de_la_hire","confianca":1.0,"epistemico":"fato","origem":"wiki","rel":"relaciona"},{"alvo":"chicote_e_o_tensor","confianca":1.0,"epistemico":"fato","origem":"wiki","rel":"relaciona"},{"alvo":"o_chicote_a_limpeza_nao_um_operador","confianca":1.0,"epistemico":"fato","origem":"wiki","rel":"relaciona"},{"alvo":"crivo_de_selberg","confianca":1.0,"epistemico":"fato","origem":"wiki","rel":"relaciona"},{"alvo":"bai_e_a_cosmologia_quantica","confianca":1.0,"epistemico":"fato","origem":"wiki","rel":"relaciona"},{"alvo":"dinamica_de_selberg_o_termo","confianca":1.0,"epistemico":"fato","origem":"wiki","rel":"relaciona"},{"alvo":"beta_numeracao_metalica","confianca":1.0,"epistemico":"fato","origem":"wiki","rel":"relaciona"},{"alvo":"norma_estelar","confianca":1.0,"epistemico":"fato","origem":"wiki","rel":"relaciona"}],"confianca":1.0,"contraexemplos":[],"descricao":"A pedra de Roseta da máquina: para cada nome poético/interno, o termo CONSAGRADO da literatura, declarado como sinônimo. A limpeza geral — os nomes não-consagrados deixam de ser ruído, cada um com seu par estabelecido. 17 termos, todos ancorados no grafo (linguagem/vocabulario.py); o dicionário 'vocabulario' traduz máquina↔consagrado. NÃO renomeia ids (isso quebraria arestas): cristaliza a correspondência. Ex.: gato→automorfismo hiperbólico do toro (Anosov); cristal→Pisot; borda→Salem/linha crítica; estrela→espectro SVD; peneira→crivo de Selberg; dinâmica de Selberg→dinâmica de Anosov (≡ Selberg-Anosov). Não é teoria nova — é a língua da máquina posta em ordem.","epistemico":"fato","exemplos":[],"id":"vocabulario_canonico","memoria":"semantica","meta":{"dominio":"referencia","fonte":"vocabulário canônico"},"origem":"wiki","palavras_chave":[],"sinonimos":[],"tipo":"conceito","titulo":"O Vocabulário Canônico (máquina ↔ literatura)"}
\section{O Vocabulário Canônico (máquina \ensuremath{\leftrightarrow} literatura)}
A pedra de Roseta da máquina: para cada nome poético/interno, o termo CONSAGRADO da literatura, declarado como sinônimo. A limpeza geral — os nomes não-consagrados deixam de ser ruído, cada um com seu par estabelecido. 17 termos, todos ancorados no grafo (linguagem/vocabulario.py); o dicionário 'vocabulario' traduz máquina\ensuremath{\leftrightarrow}consagrado. NÃO renomeia ids (isso quebraria arestas): cristaliza a correspondência. Ex.: gato\ensuremath{\to}automorfismo hiperbólico do toro (Anosov); cristal\ensuremath{\to}Pisot; borda\ensuremath{\to}Salem/linha crítica; estrela\ensuremath{\to}espectro SVD; peneira\ensuremath{\to}crivo de Selberg; dinâmica de Selberg\ensuremath{\to}dinâmica de Anosov (\ensuremath{\equiv} Selberg-Anosov). Não é teoria nova — é a língua da máquina posta em ordem.
\prov{relações: relaciona gato; relaciona lei\_da\_parabola\_e\_dinamica\_de\_anosov; relaciona numeros\_de\_pisot; relaciona numero\_de\_salem; relaciona entropia\_topologica; relaciona flip\_e\_rotacao\_de\_wick; relaciona espectro\_irredutivel\_floco; relaciona corpo\_estelar; relaciona mecanismo\_de\_la\_hire; relaciona chicote\_e\_o\_tensor; relaciona o\_chicote\_a\_limpeza\_nao\_um\_operador; relaciona crivo\_de\_selberg; relaciona bai\_e\_a\_cosmologia\_quantica; relaciona dinamica\_de\_selberg\_o\_termo; relaciona beta\_numeracao\_metalica; relaciona norma\_estelar}
\prov{proveniência: wiki \textperiodcentered{} vocabulário canônico \textperiodcentered{} fato \textperiodcentered{} confiança 1.0 \textperiodcentered{} domínio referencia \textperiodcentered{} história 17 versões no jornal}

%CRISTAL {"arestas":[{"alvo":"mult_prtsu","confianca":1.0,"epistemico":"fato","origem":"wiki","rel":"usa"},{"alvo":"mecanica_quantica","confianca":1.0,"epistemico":"fato","origem":"wiki","rel":"relaciona"}],"confianca":1.0,"contraexemplos":[],"descricao":"Estende o circuito variacional (química quântica) com uma transformação hipercomplexa de 5 parâmetros otimizados junto às rotações; para H₂ (STO-3G, 2 qubits) atinge erro <10⁻⁸ Hartree — sete ordens de magnitude abaixo do ansatz de seis rotações.","epistemico":"fato","exemplos":[],"id":"vqe_ansatz_algebrico","memoria":"semantica","meta":{"autor":"Lopes/Aarão","dominio":"hiper","fonte":""},"origem":"wiki","palavras_chave":[],"sinonimos":[],"tipo":"conceito","titulo":"VQE com Ansatz Algébrico"}
\section{VQE com Ansatz Algébrico}
Estende o circuito variacional (química quântica) com uma transformação hipercomplexa de 5 parâmetros otimizados junto às rotações; para H\ensuremath{_{2}} (STO-3G, 2 qubits) atinge erro <10\ensuremath{^{-8}} Hartree — sete ordens de magnitude abaixo do ansatz de seis rotações.
\prov{relações: usa mult\_prtsu; relaciona mecanica\_quantica}
\prov{proveniência: wiki \textperiodcentered{} fato \textperiodcentered{} confiança 1.0 \textperiodcentered{} domínio hiper \textperiodcentered{} história 4 versões no jornal}

%CRISTAL {"arestas":[{"alvo":"respiracao_dos_programas","confianca":1.0,"epistemico":"fato","origem":"PARABOLA_ISA.md","rel":"parte_de"},{"alvo":"word","confianca":0.5,"epistemico":"fato","origem":"wiki","rel":"relaciona"}],"confianca":1.0,"contraexemplos":[],"descricao":"- calor acumulado: **0.000** - fase acumulada: **9.000** - perda: **0.000** - potência (calor/passo): **0.000** - passos: **9**","epistemico":"fato","exemplos":[],"id":"while","memoria":"semantica","meta":{"dominio":"manual","nivel":6,"verificacao":"slug idêntico — enriquecer, não duplicar"},"origem":"PARABOLA_ISA.md","palavras_chave":[],"sinonimos":[],"tipo":"conceito","titulo":"🌱 while"}
\section{[semente] while}
- calor acumulado: **0.000** - fase acumulada: **9.000** - perda: **0.000** - potência (calor/passo): **0.000** - passos: **9**
\prov{relações: parte\_de respiracao\_dos\_programas; relaciona word}
\prov{proveniência: PARABOLA\_ISA.md \textperiodcentered{} fato \textperiodcentered{} confiança 1.0 \textperiodcentered{} domínio manual \textperiodcentered{} história 7 versões no jornal}

%CRISTAL {"arestas":[{"alvo":"2_semantica_formal","confianca":1.0,"epistemico":"fato","origem":"PROVA_ISA.md","rel":"parte_de"},{"alvo":"ula_da_broca_arquitetura_minima","confianca":1.0,"epistemico":"fato","origem":"README.md","rel":"parte_de"}],"confianca":1.0,"contraexemplos":[],"descricao":"Saída do neurônio: `[total, e]` com `total = e + o`.","epistemico":"fato","exemplos":[],"id":"word","memoria":"semantica","meta":{"dominio":"manual","nivel":6,"verificacao":"slug idêntico — enriquecer, não duplicar"},"origem":"README.md","palavras_chave":[],"sinonimos":[],"tipo":"conceito","titulo":"Word"}
\section{Word}
Saída do neurônio: `[total, e]` com `total = e + o`.
\prov{relações: parte\_de 2\_semantica\_formal; parte\_de ula\_da\_broca\_arquitetura\_minima}
\prov{proveniência: README.md \textperiodcentered{} fato \textperiodcentered{} confiança 1.0 \textperiodcentered{} domínio manual \textperiodcentered{} história 8 versões no jornal}

%CRISTAL {"arestas":[{"alvo":"flip_epsilon_discriminante","confianca":1.0,"epistemico":"fato","origem":"wiki","rel":"usa"},{"alvo":"hopfield","confianca":1.0,"epistemico":"fato","origem":"wiki","rel":"relaciona"}],"confianca":1.0,"contraexemplos":[],"descricao":"As projeções espectrais P_φ,P_ψ (idempotentes ortogonais, divisores de zero) são o yin (vazio) e sua soma I=P_φ+P_ψ é o yang (matéria). No tempo real a antifase é um bit (ℤ₂, o sinal alternante de ψⁿ); sob Wick vira fase (U(1), qubit na esfera de Bloch), com frequência de batimento φ−ψ=√5=√disc.","epistemico":"fato","exemplos":[],"id":"yin_yang_quantico","memoria":"semantica","meta":{"autor":"Aarão/broca-so","dominio":"broca_repos","fonte":""},"origem":"wiki","palavras_chave":[],"sinonimos":[],"tipo":"conceito","titulo":"Yin-Yang Quântico (antifase elevada a fase por Wick)"}
\section{Yin-Yang Quântico (antifase elevada a fase por Wick)}
As projeções espectrais P\_\ensuremath{\varphi},P\_\ensuremath{\psi} (idempotentes ortogonais, divisores de zero) são o yin (vazio) e sua soma I=P\_\ensuremath{\varphi}+P\_\ensuremath{\psi} é o yang (matéria). No tempo real a antifase é um bit (\ensuremath{\mathbb{Z}}\ensuremath{_{2}}, o sinal alternante de \ensuremath{\psi}\ensuremath{^{n}}); sob Wick vira fase (U(1), qubit na esfera de Bloch), com frequência de batimento \ensuremath{\varphi}\ensuremath{-}\ensuremath{\psi}=\ensuremath{\sqrt{\,}}5=\ensuremath{\sqrt{\,}}disc.
\prov{relações: usa flip\_epsilon\_discriminante; relaciona hopfield}
\prov{proveniência: wiki \textperiodcentered{} fato \textperiodcentered{} confiança 1.0 \textperiodcentered{} domínio broca\_repos \textperiodcentered{} história 3 versões no jornal}

\end{document}
